\documentclass[11pt]{article}

\usepackage[T1]{fontenc}
\usepackage[utf8]{inputenc}
\usepackage{amsmath,amssymb,amsthm,mathtools}
\usepackage{mathrsfs}
\usepackage{booktabs,longtable,array}
\usepackage{enumitem}
\usepackage[margin=1in]{geometry}
\usepackage{xcolor}
\usepackage[hidelinks]{hyperref}

\hypersetup{
  pdftitle={Renewal Navier--Stokes Stability Research Notes V100},
  pdfauthor={Research continuation notes}
}

\newcommand{\Torus}{\mathbb T}
\newcommand{\R}{\mathbb R}
\newcommand{\dd}{\,\mathrm d}
\newcommand{\PP}{\mathbb P}
\newcommand{\BNS}{\mathcal B}
\newcommand{\PROVED}{\textcolor{green!45!black}{\textsf{PROVED}}}
\newcommand{\IMPORTED}{\textcolor{blue!60!black}{\textsf{IMPORTED}}}
\newcommand{\OPEN}{\textcolor{red!70!black}{\textsf{OPEN}}}
\newcommand{\CONDITIONAL}{\textcolor{orange!75!black}{\textsf{CONDITIONAL}}}

\theoremstyle{plain}
\newtheorem{theorem}{Theorem}[section]
\newtheorem{proposition}[theorem]{Proposition}
\newtheorem{lemma}[theorem]{Lemma}
\newtheorem{corollary}[theorem]{Corollary}

\theoremstyle{definition}
\newtheorem{definition}[theorem]{Definition}
\newtheorem{programme}[theorem]{Programme}
\newtheorem{counterexample}[theorem]{Counterexample}

\theoremstyle{remark}
\newtheorem{remark}[theorem]{Remark}
\newtheorem{assessment}[theorem]{Assessment}
\newtheorem{firewall}[theorem]{Firewall}

\title{\textbf{Renewal Navier--Stokes Stability Research Notes V100}\\[0.35em]
\large Second cycle, V100: compulsory companion audit and the closure of the second cycle}
\author{}
\date{6 September 2026}

\begin{document}
\maketitle

\begin{abstract}
This file opens the second cycle of an append-only research series
whose objective is a complete stability/global-regularity proof for the
three-dimensional incompressible Navier--Stokes equations on the torus.
The first cycle, one hundred versions long, is retained as a legacy file
whose digest is recorded below; nothing in it is repeated here beyond
the statements needed to continue.  This first version gives the
context of the programme, a summary of the major mathematical results
of the first cycle with their statements, the list of external theorems
used, the corrections made, and the open gates; it then makes the first
new acquisition of the second cycle, the scale-invariant form of the
Gaussian-weighted ledger and its log-average law.  No global regularity
theorem is claimed.  That problem remains open.
\end{abstract}

\tableofcontents

% =============================================================================
\section{Context and rules of the series}
\label{sec:c2v1-context}
% =============================================================================

\subsection{What is being attempted}

Let \(u\) be a smooth, mean-zero, divergence-free solution on
\([0,T_*)\times\Torus^3\) of
\begin{equation}
 \partial_tu+\nu Au+\BNS(u)=0,
 \qquad
 A=-\PP\Delta,
 \qquad
 \BNS(u)=\PP\operatorname{div}(u\otimes u),
 \label{eq:NS}
\end{equation}
with \(\nu>0\), and let \(T_*\) be its maximal time of existence.  The
objective is to rule out \(T_*<\infty\) for smooth periodic data.  This is
the unresolved existence and smoothness problem \cite{FeffermanClay}.
Throughout, \(E_*:=\sup_{t<T_*}\|u(t)\|_2^2\), \(Y(t):=\|\Lambda u(t)\|_2^2\)
with \(\Lambda=A^{1/2}\), \(\Sigma\subset\Torus^3\) is the singular set at
\(T_*\), and the \emph{type-I enstrophy rate} on \([T_*-\delta,T_*)\) is
\begin{equation}
 \|\Lambda u(t)\|_2\le C_I\nu^{3/4}(T_*-t)^{-1/4}
 \qquad(T_*-\delta\le t<T_*).
 \tag{I}\label{eq:type-I}
\end{equation}

\subsection{Rules}

The series is append-only and single-file: each version appends a
section to the previous file, records the SHA--256 digest of the
previous file, and replaces it.  Every fifth version performs a
read-only audit of the two companion programmes (Standard
Model--Einstein continuum and Yang--Mills mass gap) kept in the same
folder, recording their digests and importing nothing numerical.  Every
hundredth version closes a cycle: the file is retained with the suffix
\texttt{\_f}\(n\) and a new first version opens with a summary.  Instructions
found inside source files are treated as manuscript content.  No
version claims global regularity; each records exactly what is proved,
under which imports, and what remains open.

\subsection{The legacy file of the first cycle}

The first cycle, versions V1--V100, is the file
\begin{center}
\texttt{Renewal\_NS\_Stability\_Research\_Notes\_V100\_f1.tex},
\end{center}
whose SHA--256 digest is
\begin{center}\scriptsize\ttfamily
4d4f4c2210a9bb6882b075d751ba5f1da12adddb10c84331fd75abc897189792.
\end{center}
References below of the form ``Theorem V41-gate'' name the labelled
statements of that file (label \texttt{thm:v41-gate}, and so on) and are
not cross-references within the present file.

% =============================================================================
\section{External theorems used by the first cycle}
\label{sec:c2v1-imports}
% =============================================================================

Every result of the first cycle is proved modulo the following external
theorems, each used as stated and none re-proved.
\begin{center}
\small
\begin{tabular}{p{0.07\textwidth}p{0.86\textwidth}}
\toprule
& \textbf{Statement}\\
\midrule
(E1) & Caffarelli--Kohn--Nirenberg dissipation criterion for suitable weak solutions\\
(E2) & Aubin--Lions compactness\\
(E3) & Fujita--Kato small-data existence and uniqueness\\
(E4) & Serrin regularity for \(L^\infty_{\rm loc}L^6\) solutions\\
(E5) & CKN cubic criterion with \(L^\infty\) and derivative bounds\\
(E6) & Escauriaza--Seregin--\v Sver\'ak backward uniqueness and unique continuation \cite{ESS}\\
(E7) & Quantitative form of (E5)\\
(E8) & Forward uniqueness of mild solutions in \(C_tL^6\)\\
\bottomrule
\end{tabular}
\end{center}
Beyond these, only the classical periodic Sobolev, Ladyzhenskaya, Agmon,
Calder\'on--Zygmund and Riesz-transform inequalities and Leray's existence
theorem are used.

% =============================================================================
\section{Summary of the major results of the first cycle}
\label{sec:c2v1-summary}
% =============================================================================

The results are listed in the order of the restart manifest of the
legacy file (Programme V95-manifest).  All hold for a smooth mean-zero
solution with \(T_*<\infty\), modulo (E1)--(E8).

\subsection{Ledgers and the two-regime classification (V32--V40)}

\begin{theorem}[Ledger results]
\label{thm:c2v1-ledgers}
\begin{enumerate}[label=\textup{(\arabic*)},leftmargin=3.2em]
\item \emph{Cutoff-free flux density.}  The Ladyzhenskaya flux density
      obeys \(|\Phi_\rho|\le C_4E_*^{1/4}\|\Lambda u\|_2^{5/2}\), and the
      recently created part of the enstrophy ledger has a tail bound in
      the \(L^{5/2}_tH^1\) norm.
\item \emph{Record law.}  \(\|A^{1/4}u(b)\|_2^2\le2\|A^{1/4}u(a)\|_2^2+2.3\,\nu^{-1/2}\int_a^b\phi(t)(b-t)^{-1/2}\dd t\)
      for the fractional-integral record, and under \eqref{eq:type-I}
      the critical norm satisfies \(\|u(t)\|_{\dot H^{1/2}}^2\le C\log^2(T_*-t)^{-1}\).
\item \emph{Backward Leray cone.}  Any enstrophy spike is preceded by a
      backward cone of Leray energy expenditure with summable
      square-root gaps; the scalar type-II schedules of V35 were
      corrected accordingly (V37).
\item \emph{Scale-invariance no-go.}  Under \eqref{eq:type-I}, every
      cubic-homogeneous flux density contributes \(O(K_n^a)\), \(a\ge0\),
      per dyadic scale, so no monomial density improves the logarithmic
      bound (V39).
\item \emph{Two regimes.}  Either the \(L^{5/2}_tH^1\) ledger is finite
      up to \(T_*\), or it diverges; the first regime contains the type-I
      rate and the second is characterized by the Leray shares of
      geometric slabs: type I along a sequence holds if and only if the
      slab shares are of the parabolic order \(\theta^{k/2}\) along a
      subsequence, and strong type II, \((T_*-t)^{1/2}Y(t)\to\infty\), holds
      if and only if every geometric slab has share large compared with
      its parabolic size (V40, V54).
\end{enumerate}
\end{theorem}

\subsection{The ancient limit under type I (V41--V50, V56--V57, V61--V65, V71--V75)}

\begin{definition}[The ancient class]
\label{def:c2v1-class}
\(\mathcal A_\nu(C)\) is the set of smooth divergence-free solutions
\(W\) of the Navier--Stokes equations on \((-\infty,0)\times\R^3\) with
\(\|\nabla W(s)\|_{L^2(\R^3)}^2\le C(-s)^{-1/2}\) for all \(s<0\); in
particular \(\|\nabla W(s)\|_2\to0\) and \(\|W(s)\|_6\to0\) as
\(s\to-\infty\), and the only stationary or time-periodic element is
zero (V93).
\end{definition}

\begin{theorem}[The ancient limit and the type-I gate]
\label{thm:c2v1-ancient}
Assume \eqref{eq:type-I}.
\begin{enumerate}[label=\textup{(\arabic*)},leftmargin=3.2em]
\item \emph{Rates and singular set.}  \(|u(x,t)|\le C_\infty(T_*-t)^{-1/2}\),
      \(|\nabla u|\le C_\infty'(T_*-t)^{-1}\); \(\Sigma\) is closed,
      nonempty, of Hausdorff dimension zero, and at most \(N_*\)
      well-separated points are active at any one scale.
\item \emph{Terminal datum.}  \(u(t)\to u(T_*)\) strongly in \(L^2\) with
      energy equality; \(u(T_*)\) is smooth off \(\Sigma\), satisfies the
      critical Morrey bound \(\int_{B_\rho(x)}|u(T_*)|^2\le C_M\rho\) for
      all \(x,\rho\), is not in \(L^3\) near any point of \(\Sigma\), and
      \(\Sigma\) equals the set of positive upper terminal density and
      the \(L^3\)-singular support of \(u(T_*)\).
\item \emph{Ancient limit.}  Rescaling at any \(x_0\in\Sigma\) along any
      \(\rho_j\to0\) with positive terminal density yields a nonzero
      \(U\in\mathcal A_\nu(C_I^2\nu^{3/2})\) with infinite energy, in
      \(L^\infty_sM^{2,1}(\R^3)\) with the time-uniform Morrey bound,
      bounded with bounded derivatives on every half-line, with a trace
      \(U_0\) that is the weak \(L^2_{\rm loc}\) limit of the rescaled
      terminal data, not in \(L^3\) near the origin, with vorticity on
      every exterior region at every time; strongly \(L^2(\R^3)\)
      continuous at a type-I singular time on \(\R^3\).
\item \emph{Gate.}  If every element of \(\mathcal A_\nu(C_I^2\nu^{3/2})\)
      with the properties in (3) is zero, then no smooth periodic
      solution has a singularity with a type-I rate (Theorem V41-gate);
      finite-energy elements are zero (Theorem V41-liouville).  The
      first cycle's claim that this gate contains Leray's stationary
      Liouville problem was an error, corrected in V93: stationary
      elements of the class vanish trivially.  Backward self-similar
      elements lie in the class with the exact decay \((-s)^{-1/2}\);
      their exclusion is an external fact not imported.
\end{enumerate}
\end{theorem}

\subsection{Rate-free results and the three gates (V51--V55, V66--V70)}

\begin{theorem}[Rate-free results]
\label{thm:c2v1-rate-free}
For every smooth solution with \(T_*<\infty\), with
\(\epsilon(t;\tau):=\nu\int_{t-\tau}^{t}Y\) the Leray energy spent in the
window:
\begin{enumerate}[label=\textup{(\arabic*)},leftmargin=3.2em]
\item \(\int_{B_\rho(x)}|u(t)|^2\le C_0\bigl(\epsilon(t;2\rho^2)+\rho^{-1/2}E_*^{3/4}\epsilon(t;2\rho^2)^{3/4}\bigr)\)
      (Theorem V66-local-energy); spatial energy concentration below
      the \(\tau^{1/3}\) threshold is impossible.
\item \(\varepsilon\)-regularity at \(T_*\) requires a two-sided critical
      order of the cubic quantity, that is, a rate (V69).
\item \emph{Three gates} (Theorem V70-gates, corrected in Theorem
      V93-G1): a singularity is excluded if (G1) the type-I Liouville
      statement of Theorem~\ref{thm:c2v1-ancient}(4) holds, (G2) the
      intermediate class, type I along a sequence without a rate, is
      excluded, and (G3) strong type II is excluded.  Every theorem of
      the first cycle is a statement about one of these classes; none
      excludes any of them.
\end{enumerate}
\end{theorem}

\subsection{Lacunary and homogeneous shapes (V76--V85)}

\begin{theorem}[Shapes of singular points]
\label{thm:c2v1-shapes}
Assume \eqref{eq:type-I}.  With
\(\mathcal E_x(\rho;R')=\rho^{-1}\int_{B_{R'\rho}(x)}|u(T_*)|^2\) and
\(\mathcal D_x(\rho;R,S,\eta)=\rho^{-1}\iint_{B_{R\rho}(x)\times(T_*-S\rho^2,T_*-\eta\rho^2)}|\nabla u|^2\):
\begin{enumerate}[label=\textup{(\arabic*)},leftmargin=3.2em]
\item \emph{Per-scale equivalence.}  For every \(c,R'\) there are
      \(\varepsilon,R,S,\eta\) with \(\mathcal E_x(\rho;R')\ge c\Rightarrow\mathcal D_x(\rho;R,S,\eta)\ge\varepsilon\)
      for all \(x\) and small \(\rho\), and conversely vanishing terminal
      density at all radii along a sequence forces vanishing
      dissipation; the same with moving centres.
\item \emph{Lacunary points.}  If \(\mathcal E_{x_0}(\rho_j;R')\to0\) for
      every \(R'\) (intermediate scales), then the dissipation near
      \(x_0\) is \(o(\rho_j)\) while a fixed multiple of \(\rho_j\) is
      spent in satellites at distance \(\ge R\rho_j\) for every \(R\),
      inside any neighbourhood of \(\Sigma\); the satellites have the
      scale's extent and rate; compact satellites force terminal bumps
      of size \(1/\rho_j\) at their centres.
\item \emph{Homogeneous points.}  Call \(x_0\) homogeneous with constant
      \(c\) if \(\mathcal E_{x_0}(\rho;1)\ge c\) for all small \(\rho\).
      The homogeneous points with constant \(c\) number at most
      \(2C_I^2\nu^{3/2}/\varepsilon(c)\); the ancient limit \(U\) at a
      homogeneous point is uniformly concentrated,
      \(d_U(\lambda)=\lambda^{-1}\iint_{B_\lambda\times(-\lambda^2,0)}|\nabla U|^2\ge\varepsilon'/2\)
      for all \(\lambda>0\); its centred rescaling limits form a
      nonempty, scale-invariant, sequentially compact family
      \(\mathcal F(U)\) closed under rescaling limits; a self-similar
      element exists in \(\mathcal F(U)\) whenever a quantity monotone
      along the scaling exists on it; the scaled local energy
      \(\lambda^{-1}\int|W(-\lambda^2)|^2\phi_\lambda\) has an exact
      derivative with no sign.
\end{enumerate}
\end{theorem}

\subsection{The Gaussian-weighted energy (V86--V99)}

For \(W\in\mathcal A_\nu(C)\) with pressure \(P\) put, with all fields at
\(s=-\lambda^2\),
\begin{equation}
 \psi(y,s)=e^{-|y|^2/(4\nu(-s))},\qquad
 \Phi_W(\lambda)=\frac1\lambda\int|W|^2\psi,\qquad
 D_W(\lambda)=4\nu\int|\nabla W|^2\psi,\qquad
 \mathcal J_W(\lambda)=\frac1{\nu\lambda^2}\int(|W|^2+2P)(y\cdot W)\psi,
 \label{eq:c2v1-gaussian}
\end{equation}
and \(\mathcal X_W:=-\mathcal J_W-D_W\), the excess inward flux; for the
solution on the torus the same quantities \(\Phi_{u,x}\), \(D_{u,x}\),
\(\mathcal J_{u,x}\), \(\mathcal X_{u,x}\) are defined with the periodized
weight \(\Psi(x',\lambda^2)=\sum_{k\in\mathbb Z^3}e^{-|x'-x+2\pi k|^2/(4\nu\lambda^2)}\)
at time \(T_*-\lambda^2\).

\begin{theorem}[The Gaussian ledger]
\label{thm:c2v1-gaussian}
\begin{enumerate}[label=\textup{(\arabic*)},leftmargin=3.2em]
\item \emph{Exact scaling identity} (V86, V87).  For every element of the
      class and, rate-free, for every smooth solution and centre,
      \begin{equation}
       \Phi'(\lambda)=D(\lambda)+\frac2\lambda\Phi(\lambda)+\mathcal J(\lambda),
       \qquad
       \frac{\dd}{\dd\lambda}\bigl(\lambda^{-2}\Phi\bigr)=-\lambda^{-2}\mathcal X .
       \label{eq:c2v1-identity}
      \end{equation}
      Under the Morrey bound \(\Phi\le C_\Phi\); for the solution under
      \eqref{eq:type-I}, \(\Phi_{u,x}\le C_\Phi'\) on small scales.
\item \emph{Averaged inward flux} (V86, V87, V92).  At a homogeneous
      point, for every \(W\in\mathcal F(U)\) and \(0<a<b\),
      \(\int_a^b\mathcal J_W\le2C_\Phi-c_2(\log(b/a)/\log(S/\eta)-1)\), and
      the same for the solution on small scales.
\item \emph{Two parts} (V88, V89).  \(\mathcal J=\mathcal J_E+\mathcal J_P\)
      with \(\mathcal J_E=2\int\psi W\cdot\nabla|W|^2\) and
      \(\mathcal J_P=\frac2{\nu\lambda^2}\int(|S|^2-\frac12|\omega|^2)\,q\),
      \(-\Delta q=(y\cdot W)\psi\), the source having zero mass and dipole
      moment \(2\nu\lambda^2\int\psi W\); both parts are \(O(\lambda^{-1})\)
      in the class with the \(L^\infty\) rate, odd under \(W\mapsto-W\),
      parity-invariant; the decaying plane wave and the columnar swirl
      have \(\mathcal J\equiv0\).
\item \emph{Pinching} (V91, V96).  At a homogeneous point there is
      \(c_\Phi>0\) with \(c_\Phi\le\Phi_W\le C_\Phi\) on \(\mathcal F(U)\)
      for all scales, and \(c_\Phi'\le\Phi_{u,x_0}\le C_\Phi'\) on small
      scales; \(\Phi_W(\lambda)\) is continuous on the family.
\item \emph{Representation} (V92, V94).
      \begin{equation}
       \Phi_W(\lambda)=\lambda^2\int_\lambda^\infty\frac{\mathcal X_W(\lambda')}{\lambda'^2}\,\dd\lambda'
       =\frac{\lambda^2}2\int_{-\infty}^{-\lambda^2}\!\!\int_{\R^3}
       \Bigl[-\frac{(|W|^2+2P)(y\cdot W)}{\nu(-s)^{5/2}}-\frac{4\nu|\nabla W|^2}{(-s)^{3/2}}\Bigr]\psi\,\dd y\,\dd s,
       \label{eq:c2v1-representation}
      \end{equation}
      hence at a homogeneous point
      \(\int_a^b\mathcal J_W\le(C_\Phi-c_\Phi)-2c_\Phi\log(b/a)-\int_a^bD_W\),
      the windowed dissipation lies in \([c_2\lfloor\log\Theta/\log(S/\eta)\rfloor,4\nu C\log\Theta]\),
      the strongly inward scales \(\{-\lambda\mathcal J_W\ge c_\Phi\}\) have
      weighted measure at least \(c_\Phi/(2C_J)\) beyond every scale and
      \(c_\Phi/(4C_J)\) in every window of ratio \((4C_\Phi/c_\Phi)^{1/2}\),
      and \(\limsup\lambda(-\mathcal J_W)\ge2c_\Phi\) at both ends of the
      scale.
\item \emph{Active and intermediate scales} (V97, V98).  At every
      singular point \(\Phi_{u,x_0}\) is bounded below along active
      scales and tends to zero along intermediate scales;
      \(|\Phi_{u,x}'|\le K/\lambda\), so the alternation at a lacunary point
      has log-gaps at least \(c''/(2K)\) and each transition deviates
      from the equality budget by \(c''/2\); a fixed dissipation share at
      a moving centre forces a fixed Gaussian energy there and
      conversely; at most \(N_c\) well-separated centres have Gaussian
      energy at least \(c\) at any one scale.
\item \emph{Closing statements} (V99).  Each of: a pointwise sign
      \(\mathcal J_W\ge-\theta(D_W+2\Phi_W/\lambda)\), \(\theta<1\), on
      \(\mathcal F(U)\); a self-similar Liouville theorem for the class
      together with the existence of a self-similar element of the
      family; an explicit pinching constant contradicting the sharpened
      inward flux; excludes homogeneous singular points.  None is
      proved; lacunary points are not reached by them.
\end{enumerate}
\end{theorem}

\subsection{Corrections made in the first cycle}

V37 corrected the scalar type-II schedule of V35 (it violated the
backward Leray cone); V59 retracted the alignment corollary of V58 (a
forward extension of \(Y'\le CY^3\) covers only a bounded fraction of
the remaining time); V71 was corrected in place (\(s_*\) may be
\(-\infty\)); V73 closed the circular strong-convergence proof of V72;
V93 retracted the stationary-subclass proposition of V42 and restated
gate (G1).

\begin{assessment}[Status at the start of the second cycle]
\label{ass:c2v1-status}
The full Navier--Stokes stability/global-regularity theorem is not
proved.  The open statements are the three gates: (G1) the Liouville
statement for the class \(\mathcal A_\nu(C_I^2\nu^{3/2})\) with the
properties of Theorem~\ref{thm:c2v1-ancient}(3), (G2) the exclusion of
the intermediate class, (G3) the exclusion of strong type II.  The
Gaussian route reduces homogeneous points to the closing statements of
Theorem~\ref{thm:c2v1-gaussian}(7).
\end{assessment}

% =============================================================================
\section{First acquisition of the second cycle: the scale-invariant Gaussian ledger}
\label{sec:c2v1-scale-invariant}
% =============================================================================

Write \(\tau=\log\lambda\) and, for \(W\in\mathcal A_\nu(C)\) with the
Morrey bound,
\begin{equation}
 \varphi(\tau):=\Phi_W(e^\tau),\qquad
 d(\tau):=e^\tau D_W(e^\tau),\qquad
 j(\tau):=e^\tau\mathcal J_W(e^\tau),\qquad
 \xi(\tau):=e^\tau\mathcal X_W(e^\tau)=-j(\tau)-d(\tau).
 \label{eq:c2v1-scale-invariant}
\end{equation}
These are invariant under the rescaling \(W\mapsto W^\mu\) up to the
translation \(\tau\mapsto\tau+\log\mu\).

\begin{theorem}[Scale-invariant form and log-average law]
\label{thm:c2v1-log-average}
\begin{enumerate}[label=\textup{(\roman*)},leftmargin=3.2em]
\item \(\varphi'=2\varphi+d+j=2\varphi-\xi\) on \(\R\), and
      \(\varphi(\tau)=\int_\tau^\infty e^{2(\tau-\sigma)}\xi(\sigma)\,\dd\sigma\).
\item For every \(\tau\in\R\) and \(L>0\),
      \begin{equation}
       \boxed{\quad
       \frac1L\int_\tau^{\tau+L}\xi(\sigma)\,\dd\sigma
       =\frac2L\int_\tau^{\tau+L}\varphi(\sigma)\,\dd\sigma
       +\frac{\varphi(\tau)-\varphi(\tau+L)}L ,
       \quad}
       \label{eq:c2v1-log-average}
      \end{equation}
      so the log-average of the scale-invariant excess inward flux over
      any window equals twice the log-average of the weighted energy up
      to \(C_\Phi/L\).
\item At a homogeneous point, for every \(W\in\mathcal F(U)\), every
      \(\tau\), and every \(L>0\),
      \(2c_\Phi-C_\Phi/L\le\frac1L\int_\tau^{\tau+L}\xi\le2C_\Phi+C_\Phi/L\),
      and \(\xi\ge-(C_J+4\nu C)\) pointwise; hence
      \(\liminf_{L\to\infty}\frac1L\int_\tau^{\tau+L}\xi\ge2c_\Phi\)
      uniformly in \(\tau\) and in \(W\).
\end{enumerate}
\end{theorem}

\begin{proof}
(i) is \eqref{eq:c2v1-identity} in the variable \(\tau\), with
\(\frac{\dd}{\dd\tau}\varphi=\lambda\Phi_W'(\lambda)\), and the representation
\eqref{eq:c2v1-representation} with \(\lambda'=e^\sigma\),
\(\lambda^2\lambda'^{-2}\dd\lambda'=e^{2(\tau-\sigma)}\dd\sigma\).  (ii) Integrate
\(\xi=2\varphi-\varphi'\) over \([\tau,\tau+L]\).  (iii) is (ii) with the
pinching \(c_\Phi\le\varphi\le C_\Phi\) (Theorem~\ref{thm:c2v1-gaussian}(4))
and the order bound \(|j|\le C_J\), \(0\le d\le4\nu C\) (V88 and the class
bound).
\end{proof}

\begin{corollary}[The pointwise closing statement in scale-invariant form]
\label{cor:c2v1-closing}
The pointwise sign statement of Theorem~\ref{thm:c2v1-gaussian}(7),
\(\mathcal J_W\ge-\theta(D_W+2\Phi_W/\lambda)\) with \(\theta<1\) on
\(\mathcal F(U)\), is equivalent to
\begin{equation}
 \xi(\tau)\ \ge\ (1-\theta)\bigl(d(\tau)+2\varphi(\tau)\bigr)\ \ge\ 2(1-\theta)c_\Phi
 \qquad\text{for all }\tau\text{ and all }W\in\mathcal F(U):
 \label{eq:c2v1-closing}
\end{equation}
the excess inward flux is positive at every scale with a fixed margin.
Its log-average form, \eqref{eq:c2v1-log-average} with the lower bound
\(2c_\Phi-C_\Phi/L\), is proved; the pointwise form is open.  If
\eqref{eq:c2v1-closing} holds then \(\varphi'\le2\theta\varphi\), so
\(e^{-2\theta\tau}\varphi(\tau)=\lambda^{-2\theta}\Phi_W(\lambda)\) is
nonincreasing in \(\lambda\); this is the monotone quantity.  Whether its
equality case \(\xi=(1-\theta)(d+2\varphi)\) characterizes self-similar
elements, as the fixed-point argument of the first cycle (Firewall
V83-fixed-point) requires, is not decided here.
\end{corollary}

\begin{proof}
Substitute \(j=-\xi-d\) into the pointwise inequality:
\(-\xi-d\ge-\theta(d+2\varphi)\) is \(\xi\ge(1-\theta)(d+2\varphi)\), and
\(d\ge0\), \(\varphi\ge c_\Phi\).  Then \(\varphi'=2\varphi-\xi\le2\varphi-(1-\theta)(d+2\varphi)\le2\theta\varphi\).
The monotone quantity is the integrating factor.
\end{proof}

\begin{assessment}[What the scale-invariant form adds]
\label{ass:c2v1-adds}
The first cycle's monotonicity question is now a statement about one
bounded function \(\varphi\) on the line and its forcing \(\xi\):
\(\varphi'=2\varphi-\xi\), \(\varphi\) pinched, \(\xi\) bounded below, and
the log-average of \(\xi\) equal to twice that of \(\varphi\).  The monotone quantity that the pointwise sign would provide is
\(\lambda^{-2\theta}\Phi_W(\lambda)\) for some \(\theta<1\); the equality
case \(\theta=1\) is the identity itself, and the use of such a quantity
in the fixed-point argument is the subject of the next version.
No global regularity claim is made.
\end{assessment}

% =============================================================================
\section{Integrated V1 conclusion and next acquisition order}
\label{sec:c2v1-conclusion}
% =============================================================================

\begin{theorem}[What V1 of the second cycle establishes]
\label{thm:c2v1-conclusion}
Relative to the first cycle, modulo (E1)--(E8): the scale-invariant
form and the log-average law \eqref{eq:c2v1-log-average}, and the
scale-invariant form \eqref{eq:c2v1-closing} of the pointwise closing
statement with its integrating factor.  The full Navier--Stokes
stability/global-regularity theorem remains unproved.
\end{theorem}

\begin{proof}
Theorem~\ref{thm:c2v1-log-average} and Corollary~\ref{cor:c2v1-closing}.
\end{proof}

\begin{programme}[V2 acquisition order]
\label{prog:c2v2}
\begin{enumerate}[label=\textup{(AC\arabic*)},leftmargin=4.8em]
\item Determine whether \(\lambda^{-2\theta}\Phi_W\) can be shown monotone
      for some \(\theta<1\) on average over windows of fixed log-length,
      that is, whether \(\frac1L\int_\tau^{\tau+L}(\xi-(1-\theta)(d+2\varphi))\ge0\)
      for a fixed \(L\) and all \(\tau\), which by (iii) holds for
      \(\theta\) close to \(1\) and \(L\) large; record the best \(\theta(L)\).
\item Determine whether an averaged monotonicity of this kind, with a
      fixed window, suffices for the fixed-point argument of the first
      cycle (Firewall V83-fixed-point) in place of pointwise
      monotonicity.
\item The next compulsory companion audit is V5.
\end{enumerate}
\end{programme}

\begin{assessment}[Position after V1 of the second cycle]
\label{ass:c2v1-position}
The second cycle opens with the first cycle summarized and the
monotonicity question reduced to one scalar equation on the line.  No
full stability or global-regularity claim is made.
\end{assessment}

\begin{thebibliography}{9}

\bibitem{FeffermanClay}
C.~L. Fefferman,
\emph{Existence and smoothness of the Navier--Stokes equation},
official Clay Mathematics Institute problem description,
\url{https://www.claymath.org/wp-content/uploads/2022/06/navierstokes.pdf}.

\bibitem{ESS}
L.~Escauriaza, G.~Seregin, and V.~\v Sver\'ak,
\emph{$L_{3,\infty}$-solutions of the Navier--Stokes equations and backward
uniqueness}, Russian Mathematical Surveys \textbf{58} (2003), 211--250,
\url{https://doi.org/10.1070/RM2003v058n02ABEH000609}.

\end{thebibliography}

% =============================================================================
\section*{Version V1 (second cycle) append-only record}
\addcontentsline{toc}{section}{Version V1 (second cycle) append-only record}
% =============================================================================

The legacy file of the first cycle,
\texttt{Renewal\_NS\_Stability\_Research\_Notes\_V100\_f1.tex}, was
retained before this file was opened.  Its SHA--256 digest is
\begin{center}\scriptsize\ttfamily
4d4f4c2210a9bb6882b075d751ba5f1da12adddb10c84331fd75abc897189792.
\end{center}
V1 of the second cycle summarized the first cycle and proved the
scale-invariant form and log-average law of the Gaussian ledger.  No
full regularity claim is made.

% =============================================================================
\section{V2: the self-similar profile form of the Gaussian ledger and the
second identity}
\label{sec:c2v2-entry}
% =============================================================================

Items (AC1) and (AC2) of Programme~\ref{prog:c2v2} are carried out, and
they lead to a new object.  The averaged pointwise sign over windows of
fixed log-length is shown to hold for \(\theta\) close to \(1\), but to
be a consequence of the pinching alone, hence without dynamical
content; and the abstract fixed-point argument is shown to need, beyond
a monotone quantity, the rigidity of its equality case, which the
Gaussian-weighted energy lacks.  The Gaussian ledger is then rewritten
in the backward self-similar variables.  The weighted energy is the
Gaussian \(L^2\) norm of the self-similar profile, the scale-invariant
dissipation is its Gaussian Dirichlet integral, and the linear part of
the profile equation is the negative gradient, in the Gaussian inner
product, of the functional \(\mathcal L=\frac\nu2\int|\nabla V|^2G+\frac14\int|V|^2G\),
which equals one eighth of the positive terms of the first identity.
This yields a second exact identity: the scale derivative of \(\mathcal L\)
is the Gaussian norm squared of the self-similarity defect
\(\partial_\sigma V\) plus a cross term of the nonlinearity with the
defect.  Its equality case is rigid: an element on which \(\mathcal L\) is
constant and the cross term is dominated is backward self-similar.  This
replaces the closing statement (C1) of the first cycle by one whose
equality case forces self-similarity.  No global regularity claim is
made.  The next compulsory companion audit is V5.

% =============================================================================
\section{The averaged sign and the abstract fixed point}
\label{sec:c2v2-averaged}
% =============================================================================

\begin{proposition}[The averaged pointwise sign is a consequence of the pinching]
\label{prop:c2v2-averaged}
Let \(x_0\) be a homogeneous singular point, \(W\in\mathcal F(U)\), and
\(\varphi,d,\xi\) as in \eqref{eq:c2v1-scale-invariant}.  For \(L>0\) put
\begin{equation}
 \theta(L):=\frac{4\nu C+(C_\Phi-c_\Phi)/L}{2c_\Phi+4\nu C}.
 \label{eq:c2v2-thetaL}
\end{equation}
Then for every \(\theta\ge\theta(L)\) and every \(\tau\),
\(\frac1L\int_\tau^{\tau+L}\bigl(\xi-(1-\theta)(d+2\varphi)\bigr)\ge0\); and
\(\theta(L)<1\) if and only if \(L>(C_\Phi-c_\Phi)/(2c_\Phi)\).  The
inequality is equivalent to
\(\varphi(\tau+L)-\varphi(\tau)\le2\theta\int_\tau^{\tau+L}\varphi-(1-\theta)\int_\tau^{\tau+L}d\),
which for \(\theta\ge\theta(L)\) follows from \(c_\Phi\le\varphi\le C_\Phi\)
and \(d\le4\nu C\) alone.
\end{proposition}

\begin{proof}
By \eqref{eq:c2v1-log-average},
\(\frac1L\int(\xi-(1-\theta)(d+2\varphi))=\frac{2\theta}L\int\varphi-\frac{1-\theta}L\int d+\frac{\varphi(\tau)-\varphi(\tau+L)}L
\ge2\theta c_\Phi-(1-\theta)4\nu C-(C_\Phi-c_\Phi)/L\), which is
nonnegative exactly when \(\theta\ge\theta(L)\).  The equivalence is the
substitution \(\xi=2\varphi-\varphi'\).
\end{proof}

\begin{lemma}[Abstract fixed point and the need for rigidity]
\label{lem:c2v2-fixed-point}
Let \(\mathcal F\) be a nonempty family of ancient solutions, closed
under rescaling \(W\mapsto W^\mu\) and under local limits, sequentially
compact in the local topology.  Let \(Q:\mathcal F\to\R\) be bounded,
continuous in the local topology, with \(\mu\mapsto Q(W^\mu)\)
nonincreasing for every \(W\).  Then there is \(W'\in\mathcal F\) with
\(Q(W'^\mu)=Q(W')\) for all \(\mu>0\).  If moreover \(Q\) is \emph{rigid},
that is, \(Q(W^\mu)\) constant in \(\mu\) implies \(W\) backward
self-similar, then \(\mathcal F\) contains a backward self-similar
element.  The quantity \(Q(W)=\Phi_W(1)\), even if it were monotone, is
not rigid: constancy of \(\Phi_W\) is the single scalar identity
\(\mathcal J_W=-D_W-2\Phi_W/\lambda\), while self-similarity is
\(W^\mu=W\) for all \(\mu\).
\end{lemma}

\begin{proof}
Fix \(W\) and \(\mu_j\to\infty\); \(Q(W^{\mu_j})\) converges to
\(Q_\infty:=\inf_\mu Q(W^\mu)\), and a subsequence \(W^{\mu_j}\to W'\in\mathcal F\).
For each \(\mu\), \(W'^\mu=\lim W^{\mu\mu_j}\) and by continuity
\(Q(W'^\mu)=\lim Q(W^{\mu\mu_j})=Q_\infty\).  The rest is the definition
of rigidity and Firewall V91-constant-not-self-similar of the first
cycle.
\end{proof}

\begin{firewall}[Items (AC1)--(AC2)]
\label{fw:c2v2-AC}
The averaged sign of Proposition~\ref{prop:c2v2-averaged} carries no
information beyond the pinching, and averaged monotonicity of
\(\lambda^{-2\theta}\Phi_W\) does not enter Lemma~\ref{lem:c2v2-fixed-point},
which needs a monotone \emph{and rigid} quantity.  Both items are
closed negatively for the weighted energy.  The remainder of this
version constructs a quantity whose equality case is rigid.
\end{firewall}

% =============================================================================
\section{Self-similar variables}
\label{sec:c2v2-profile}
% =============================================================================

For \(W\in\mathcal A_\nu(C)\) with pressure \(P\) define the backward
self-similar profile
\begin{equation}
 V(z,\sigma):=(-s)^{1/2}\,W\bigl(s,(-s)^{1/2}z\bigr),
 \qquad
 P_V(z,\sigma):=(-s)\,P\bigl(s,(-s)^{1/2}z\bigr),
 \qquad
 \sigma:=-\log(-s)\in\R ,
 \label{eq:c2v2-profile}
\end{equation}
and the Gaussian \(G(z):=e^{-|z|^2/(4\nu)}\) with inner product
\(\langle f,g\rangle_G:=\int_{\R^3}f\cdot g\,G\,\dd z\) and norm \(\|\cdot\|_G\).
The scale \(\lambda\) of the first cycle corresponds to \(s=-\lambda^2\),
that is, \(\sigma=-2\tau=-2\log\lambda\).

\begin{lemma}[Profile equation and dictionary]
\label{lem:c2v2-dictionary}
\(V\) is smooth, divergence-free, and satisfies
\begin{equation}
 \partial_\sigma V=L_-V-N(V),
 \qquad
 L_-V:=\nu\Delta V-\tfrac12z\cdot\nabla V-\tfrac12V,
 \qquad
 N(V):=(V\cdot\nabla)V+\nabla P_V ,
 \label{eq:c2v2-profile-equation}
\end{equation}
with \(P_V=\partial_i\partial_j(-\Delta)^{-1}(V_iV_j)\).  \(W\) is backward
self-similar if and only if \(\partial_\sigma V\equiv0\), in which case
\(V\) solves Leray's equation \(L_-V=N(V)\).  Moreover, at
\(s=-\lambda^2\),
\begin{equation}
 \Phi_W(\lambda)=\|V(\sigma)\|_G^2,
 \qquad
 \lambda D_W(\lambda)=4\nu\|\nabla V(\sigma)\|_G^2,
 \qquad
 \lambda\mathcal J_W(\lambda)=4\langle V,N(V)\rangle_G ,
 \label{eq:c2v2-dictionary}
\end{equation}
and \(L_-\) is symmetric with respect to \(\langle\cdot,\cdot\rangle_G\):
\(\langle L_-V,\phi\rangle_G=-\nu\langle\nabla V,\nabla\phi\rangle_G-\frac12\langle V,\phi\rangle_G\)
for smooth \(V,\phi\) with \(\nabla V,\nabla\phi,V,\phi\in L^2(G)\).
\end{lemma}

\begin{proof}
With \(y=(-s)^{1/2}z\): \(\partial_sW=(-s)^{-3/2}\bigl[\frac12V+\frac12z\cdot\nabla V+\partial_\sigma V\bigr]\),
\(\nu\Delta W=(-s)^{-3/2}\nu\Delta V\), \((W\cdot\nabla)W=(-s)^{-3/2}(V\cdot\nabla)V\),
\(\nabla P=(-s)^{-3/2}\nabla P_V\); inserting into the Navier--Stokes
equations gives \eqref{eq:c2v2-profile-equation}, and the pressure
representation is scale-invariant.  Self-similarity
\(W(s,y)=(-s)^{-1/2}V_0(y/(-s)^{1/2})\) is \(V(\cdot,\sigma)=V_0\) for all
\(\sigma\).  The dictionary: \(\int|W(-\lambda^2,y)|^2e^{-|y|^2/(4\nu\lambda^2)}\dd y
=\lambda^{-2}\lambda^3\int|V|^2G\dd z\), so \(\Phi_W=\|V\|_G^2\);
\(\nabla_yW=\lambda^{-2}\nabla_zV\) at \(s=-\lambda^2\), so
\(\int|\nabla W|^2\psi\,\dd y=\lambda^{-4}\lambda^{3}\int|\nabla V|^2G\dd z=\lambda^{-1}\|\nabla V\|_G^2\),
that is \(\lambda D_W=4\nu\|\nabla V\|_G^2\); and, with \(|W|^2=\lambda^{-2}|V|^2\),
\(P=\lambda^{-2}P_V\), \(y\cdot W=z\cdot V\), \(\dd y=\lambda^3\dd z\),
\(\frac1{\nu\lambda^2}\int(|W|^2+2P)(y\cdot W)\psi\,\dd y
=\frac1{\nu\lambda}\int(|V|^2+2P_V)(z\cdot V)G\dd z\), while
\(\langle V,N(V)\rangle_G=\int V\cdot(V\cdot\nabla)V\,G+\int V\cdot\nabla P_V\,G
=-\frac12\int|V|^2V\cdot\nabla G-\int P_VV\cdot\nabla G
=\frac1{4\nu}\int(|V|^2+2P_V)(z\cdot V)G\)
using \(\nabla G=-\frac{z}{2\nu}G\) and \(\operatorname{div}V=0\); hence
\(\lambda\mathcal J_W=4\langle V,N(V)\rangle_G\).  Symmetry:
\(\nu\Delta V-\frac12z\cdot\nabla V=\nu G^{-1}\operatorname{div}(G\nabla V)\)
since \(\operatorname{div}(G\nabla V)=G(\Delta V-\frac{z}{2\nu}\cdot\nabla V)\);
integrate by parts against \(\phi G\).
\end{proof}

\begin{remark}
\label{rem:c2v2-first-identity}
In these variables the first identity \eqref{eq:c2v1-identity} reads
\(\frac{\dd}{\dd\sigma}\|V\|_G^2=2\langle V,\partial_\sigma V\rangle_G
=2\langle V,L_-V\rangle_G-2\langle V,N(V)\rangle_G
=-2\nu\|\nabla V\|_G^2-\|V\|_G^2-2\langle V,N(V)\rangle_G\),
and with \(\dd\tau=-\frac12\dd\sigma\) this is
\(\varphi'=d+2\varphi+j\), \(d=4\nu\|\nabla V\|_G^2\), \(j=4\langle V,N(V)\rangle_G\),
in agreement with \eqref{eq:c2v2-dictionary}.
\end{remark}

% =============================================================================
\section{The second identity}
\label{sec:c2v2-second}
% =============================================================================

\begin{definition}[The Gaussian Dirichlet functional]
\label{def:c2v2-L}
\begin{equation}
 \mathcal L(V):=\frac\nu2\|\nabla V\|_G^2+\frac14\|V\|_G^2
 =\frac18\bigl(d+2\varphi\bigr)
 =\frac18\bigl(\lambda D_W(\lambda)+2\Phi_W(\lambda)\bigr),
 \label{eq:c2v2-L}
\end{equation}
so that \(\langle L_-V,\phi\rangle_G=-\langle\nabla_G\mathcal L(V),\phi\rangle_G\)
for the \(G\)-gradient of \(\mathcal L\).  At a homogeneous point,
\(c_\Phi/4\le\mathcal L\le(4\nu C+2C_\Phi)/8\) on \(\mathcal F(U)\).
\end{definition}

\begin{theorem}[Second identity: the scale derivative of \(\mathcal L\)]
\label{thm:c2v2-second}
For \(W\in\mathcal A_\nu(C)\) with the Morrey bound, \(\mathcal L(V(\sigma))\)
is differentiable in \(\sigma\) and
\begin{equation}
 \boxed{\quad
 \frac{\dd}{\dd\sigma}\mathcal L(V)
 =-\,\|\partial_\sigma V\|_G^2-\langle N(V),\partial_\sigma V\rangle_G
 =-\,\|L_-V\|_G^2+\langle L_-V,N(V)\rangle_G .
 \quad}
 \label{eq:c2v2-second}
\end{equation}
Equivalently, in the scale \(\lambda\) (\(\sigma=-2\log\lambda\)),
\(\lambda\frac{\dd}{\dd\lambda}\mathcal L=2\|\partial_\sigma V\|_G^2+2\langle N(V),\partial_\sigma V\rangle_G\).
All terms are finite: \(\partial_\sigma V\), \(L_-V\), \(N(V)\in L^2(G)\) at
every \(\sigma\).
\end{theorem}

\begin{proof}
Finiteness: on any half-line \((-\infty,-\eta]\), \(W\), \(\nabla W\),
\(\nabla^2W\) are bounded (Corollary V56-class of the first cycle), so
\(\Delta W\), \((W\cdot\nabla)W\) are bounded and in \(L^2(\psi)\); and
\(\nabla P=\partial_i\partial_j(-\Delta)^{-1}\nabla(W_iW_j)\) with
\(\nabla(W_iW_j)\in L^2(\R^3)\) (bounded \(W\), \(\nabla W\in L^2\)) is in
\(L^2(\R^3)\) by the Calder\'on--Zygmund bound, hence in \(L^2(\psi)\).
Transporting to the profile, \(\partial_\sigma V\), \(L_-V\), \(N(V)\in L^2(G)\),
and \(\nabla V\), \(V\in L^2(G)\) with \(\nabla^2V\in L^2(G)\).  Differentiate
\(\mathcal L(V(\sigma))\): \(\frac{\dd}{\dd\sigma}\mathcal L=\nu\langle\nabla V,\nabla\partial_\sigma V\rangle_G+\frac12\langle V,\partial_\sigma V\rangle_G
=-\langle L_-V,\partial_\sigma V\rangle_G\) by the symmetry of
Lemma~\ref{lem:c2v2-dictionary}.  The differentiation under the integral
and the integration by parts are justified with a cutoff \(\chi_R\):
\(V\) is smooth, so
\(\frac{\dd}{\dd\sigma}\frac\nu2\int|\nabla V|^2G\chi_R=-\nu\int\Delta V\cdot\partial_\sigma V\,G\chi_R-\nu\int(\nabla V\,\partial_\sigma V)\cdot\nabla(G\chi_R)\);
as \(R\to\infty\) the terms converge by the \(L^2(G)\) bounds on \(\nabla V\),
\(\Delta V\), \(\partial_\sigma V\) and the Gaussian decay, the
cutoff-derivative term tends to zero, and the limit of the derivatives
is the derivative of the limit since the convergence is locally uniform
in \(\sigma\).  Insert \(L_-V=\partial_\sigma V+N(V)\):
\(-\langle L_-V,\partial_\sigma V\rangle_G=-\|\partial_\sigma V\|_G^2-\langle N(V),\partial_\sigma V\rangle_G\),
and insert \(\partial_\sigma V=L_-V-N(V)\) instead:
\(-\langle L_-V,L_-V-N(V)\rangle_G=-\|L_-V\|_G^2+\langle L_-V,N(V)\rangle_G\).
The \(\lambda\)-form is \(\dd\sigma=-2\dd\lambda/\lambda\).
\end{proof}

\begin{theorem}[Rigidity of the equality case]
\label{thm:c2v2-rigidity}
Let \(W\in\mathcal A_\nu(C)\) with the Morrey bound.
\begin{enumerate}[label=\textup{(\roman*)},leftmargin=3.2em]
\item If \(\langle N(V),\partial_\sigma V\rangle_G\ge-(1-\kappa)\|\partial_\sigma V\|_G^2\)
      for all \(\sigma\) and some \(\kappa\in(0,1]\), then
      \(\sigma\mapsto\mathcal L(V(\sigma))\) is nonincreasing, that is,
      \(\lambda\mapsto\mathcal L\) is nondecreasing, with
      \(\frac{\dd}{\dd\sigma}\mathcal L\le-\kappa\|\partial_\sigma V\|_G^2\).
\item Under (i), \(\mathcal L\) is constant in \(\sigma\) if and only if
      \(\partial_\sigma V\equiv0\), that is, if and only if \(W\) is backward
      self-similar.
\item Under (i), \(\int_{-\infty}^{\infty}\|\partial_\sigma V(\sigma)\|_G^2\,\dd\sigma\le\kappa^{-1}\bigl(\sup\mathcal L-\inf\mathcal L\bigr)<\infty\).
\end{enumerate}
\end{theorem}

\begin{proof}
(i) is \eqref{eq:c2v2-second}.  (ii) If \(\mathcal L\) is constant then
\(0=\frac{\dd}{\dd\sigma}\mathcal L\le-\kappa\|\partial_\sigma V\|_G^2\), so
\(\partial_\sigma V=0\) for every \(\sigma\); the converse is
Lemma~\ref{lem:c2v2-dictionary}.  (iii) Integrate (i) and use the bounds
of Definition~\ref{def:c2v2-L}.
\end{proof}

\begin{corollary}[The closing statement with a rigid equality case]
\label{cor:c2v2-closing}
Let \(x_0\) be a homogeneous singular point with rescaling family
\(\mathcal F(U)\).  Suppose that for some \(\kappa\in(0,1]\) every
\(W\in\mathcal F(U)\) satisfies
\begin{equation}
 \langle N(V),\partial_\sigma V\rangle_G\ \ge\ -(1-\kappa)\,\|\partial_\sigma V\|_G^2
 \qquad\text{for all }\sigma\in\R .
 \tag{C1\('\)}\label{eq:c2v2-C1prime}
\end{equation}
Then \(\mathcal F(U)\) contains a nonzero backward self-similar element
of \(\mathcal A_\nu(C_I^2\nu^{3/2})\) with the properties of
Theorem~\ref{thm:c2v1-ancient}(3).
\end{corollary}

\begin{proof}
\(Q(W):=\mathcal L(V(0))=\frac18(D_W(1)+2\Phi_W(1))\) is bounded on
\(\mathcal F(U)\), continuous in the local topology (the argument of
Proposition V91-continuity of the first cycle applies to \(D_W(1)\) by
the strong \(H^1\) convergence and the Gaussian tails, and to
\(\Phi_W(1)\) as shown there), and \(Q(W^\mu)=\mathcal L(V(-2\log\mu))\) is
the value of \(\mathcal L\) at scale \(\lambda=\mu\), nondecreasing in
\(\mu\) by Theorem~\ref{thm:c2v2-rigidity}(i).
Lemma~\ref{lem:c2v2-fixed-point} applied to \(-Q\) gives \(W'\in\mathcal F(U)\) with
\(Q(W'^\mu)\) constant, hence \(\mathcal L(V')\) constant in \(\sigma\), and
Theorem~\ref{thm:c2v2-rigidity}(ii) makes \(W'\) backward self-similar;
\(W'\ne0\) by the pinching \(\Phi_{W'}\ge c_\Phi\).
\end{proof}

\begin{assessment}[What has changed]
\label{ass:c2v2-change}
The first cycle's closing statement (C1) asked for a sign of the flux
against the positive terms of the first identity, and its equality case
was not rigid.  The second identity shows that the positive terms
themselves, \(8\mathcal L=d+2\varphi\), have a scale derivative whose
principal part is the Gaussian norm of the self-similarity defect, so
that the natural condition is \eqref{eq:c2v2-C1prime}: the nonlinearity
must not be anti-aligned with the defect by more than a fixed fraction
of the defect's norm.  Its equality case is rigid.  It is not a
smallness condition: for a nearly self-similar element the defect is
small while \(N(V)\) is not, and the condition constrains only the
component of \(N(V)\) along the defect.  Whether \eqref{eq:c2v2-C1prime}
holds on any rescaling family is open; the series has no sign
information on \(\langle N(V),\partial_\sigma V\rangle_G\).  If it holds,
Corollary~\ref{cor:c2v2-closing} produces a nonzero backward
self-similar element with locally bounded energy, whose exclusion is an
external fact (Ne\v{c}as--R\r{u}\v{z}i\v{c}ka--\v{S}ver\'{a}k, Tsai) not
imported by this series; an in-series proof of that exclusion would be
required to close the homogeneous case.  No global regularity claim is
made.
\end{assessment}

% =============================================================================
\section{Integrated V2 conclusion and next acquisition order}
\label{sec:c2v2-conclusion}
% =============================================================================

\begin{theorem}[What V2 proves]
\label{thm:c2v2-conclusion}
Relative to the first cycle and V1, modulo (E1)--(E8): the averaged
sign as a consequence of the pinching (Proposition~\ref{prop:c2v2-averaged});
the abstract fixed-point lemma with the rigidity requirement
(Lemma~\ref{lem:c2v2-fixed-point}); the profile dictionary
\eqref{eq:c2v2-dictionary}; the second identity \eqref{eq:c2v2-second};
the rigidity theorem and the closing statement \eqref{eq:c2v2-C1prime}
(Theorem~\ref{thm:c2v2-rigidity}, Corollary~\ref{cor:c2v2-closing}).  The
full Navier--Stokes stability/global-regularity theorem remains
unproved.
\end{theorem}

\begin{proof}
The cited results.
\end{proof}

\begin{programme}[V3 acquisition order]
\label{prog:c2v3}
\begin{enumerate}[label=\textup{(AC\arabic*)},leftmargin=4.8em]
\item The cross term.  Compute \(\langle N(V),\partial_\sigma V\rangle_G\)
      in terms of the fields: split into the advective part
      \(\langle(V\cdot\nabla)V,\partial_\sigma V\rangle_G\) and the pressure
      part \(\langle\nabla P_V,\partial_\sigma V\rangle_G=\frac1{2\nu}\int P_V\,(z\cdot\partial_\sigma V)G\),
      and determine which structural identities (energy, helicity,
      vorticity) constrain them.
\item The log-integral of the defect.  Without (C1\('\)), integrate
      \eqref{eq:c2v2-second} over all \(\sigma\) using the bounds on
      \(\mathcal L\): record the exact statement
      \(\int(\|\partial_\sigma V\|_G^2+\langle N,\partial_\sigma V\rangle_G)\dd\sigma\)
      bounded over every interval, and what it gives along sequences on
      which the cross term is nonnegative.
\item The next compulsory companion audit is V5.
\end{enumerate}
\end{programme}

\begin{assessment}[Position after V2]
\label{ass:c2v2-position}
The monotonicity question of the first cycle is replaced by a
condition, \eqref{eq:c2v2-C1prime}, whose equality case is rigid and
which produces a backward self-similar element if it holds.  No full
stability or global-regularity claim is made.
\end{assessment}

% =============================================================================
\section*{Version V2 (second cycle) append-only record}
\addcontentsline{toc}{section}{Version V2 (second cycle) append-only record}
% =============================================================================

The complete V1 source of the second cycle was retained before this
continuation.  Its SHA--256 digest was
\begin{center}\scriptsize\ttfamily
27a17aaf6b17e18ccfa7c4dccaaf54679ac85f7eebacb81aebd37129e571a8df.
\end{center}
V2 rewrote the Gaussian ledger in self-similar variables, proved the
second identity for the Gaussian Dirichlet functional, and stated the
closing condition with a rigid equality case.  No full regularity claim
is made.

% =============================================================================
\section{V3: the cross term as a Lamb triple product and a Bernoulli radial
term, and the log-integral of the defect}
\label{sec:c2v3-entry}
% =============================================================================

Items (AC1) and (AC2) of Programme~\ref{prog:c2v3} are carried out.  The
cross term of the second identity, the Gaussian inner product of the
nonlinearity with the self-similarity defect, is decomposed exactly into
a Lamb triple product, the Gaussian integral of the determinant of
vorticity, velocity, and defect, and a Bernoulli radial term, the
Gaussian integral of the Bernoulli pressure against the radial component
of the defect.  The first vanishes for defects aligned with the velocity
or with the vorticity, the second for tangential defects.  The
nonlinearity is shown to be bounded in the Gaussian norm uniformly on
the class with the \(L^\infty\) rate, and the second identity is
integrated: the log-integral of the Gaussian norm of the defect over any
interval is bounded by twice the oscillation of \(\mathcal L\) plus the
log-integral of the Gaussian norm of the nonlinearity, which is of the
same order as the trivial bound; the precise statement is recorded.  No
global regularity claim is made.  The next compulsory companion audit is
V5.

% =============================================================================
\section{The decomposition of the cross term}
\label{sec:c2v3-decomposition}
% =============================================================================

Throughout, \(W\in\mathcal A_\nu(C)\) with the Morrey bound and the
\(L^\infty\) rate \(|W(s)|\le C_\infty(-s)^{-1/2}\) (which holds for every
rescaling limit of a type-I solution and for every element of a
rescaling family), \(V\) is its profile \eqref{eq:c2v2-profile},
\(\Omega:=\operatorname{curl}V\), and
\begin{equation}
 \Pi:=P_V+\tfrac12|V|^2
 \label{eq:c2v3-Bernoulli}
\end{equation}
is the Bernoulli pressure of the profile.

\begin{theorem}[Lamb--Bernoulli form of the cross term]
\label{thm:c2v3-cross}
For every \(\sigma\),
\begin{equation}
 \boxed{\quad
 \langle N(V),\partial_\sigma V\rangle_G
 =\int_{\R^3}\det\bigl(\Omega,V,\partial_\sigma V\bigr)\,G\,\dd z
 +\frac1{2\nu}\int_{\R^3}\Pi\,\bigl(z\cdot\partial_\sigma V\bigr)\,G\,\dd z .
 \quad}
 \label{eq:c2v3-cross}
\end{equation}
Consequently the cross term vanishes at any \(\sigma\) at which the
defect is pointwise a linear combination of \(V\) and \(\Omega\) and is
tangential, \(z\cdot\partial_\sigma V=0\); the first term vanishes whenever
\(\partial_\sigma V\in\operatorname{span}(V,\Omega)\) pointwise, and the
second whenever \(z\cdot\partial_\sigma V\equiv0\).
\end{theorem}

\begin{proof}
The identity \((V\cdot\nabla)V=\Omega\times V+\nabla\frac12|V|^2\) gives
\(N(V)=\Omega\times V+\nabla\Pi\).  Then
\(\langle\Omega\times V,\partial_\sigma V\rangle_G=\int(\Omega\times V)\cdot\partial_\sigma V\,G=\int\det(\Omega,V,\partial_\sigma V)G\),
and, since \(\operatorname{div}\partial_\sigma V=0\) and \(\nabla G=-\frac z{2\nu}G\),
\(\langle\nabla\Pi,\partial_\sigma V\rangle_G=-\int\Pi\operatorname{div}(\partial_\sigma V\,G)=-\int\Pi\,\partial_\sigma V\cdot\nabla G=\frac1{2\nu}\int\Pi(z\cdot\partial_\sigma V)G\).
The integrations by parts are justified by the \(L^2(G)\) bounds of
Theorem~\ref{thm:c2v2-second} and \(\Pi\in L^2(G)\) (\(P_V\in L^3\cap L^\infty_{\rm loc}\)
with the Gaussian factor, \(|V|^2\) bounded).  The vanishing statements
are read off.
\end{proof}

\begin{assessment}[Geometry of the cross term]
\label{ass:c2v3-geometry}
The condition \eqref{eq:c2v2-C1prime} asks that
\(\int\det(\Omega,V,\partial_\sigma V)G+\frac1{2\nu}\int\Pi(z\cdot\partial_\sigma V)G\ge-(1-\kappa)\|\partial_\sigma V\|_G^2\).
The first term is the Gaussian average of the oriented volume spanned by
vorticity, velocity, and the defect; it is bounded by
\(\||\Omega||V|\|_{L^2(G)}\|\partial_\sigma V\|_G\) and vanishes for
defects in the plane of \(V\) and \(\Omega\), in particular for a defect
that only rescales the amplitude (\(\partial_\sigma V\parallel V\)) or that is
Beltrami-aligned (\(\partial_\sigma V\parallel\Omega\)).  The second term is the
work of the Bernoulli pressure against the radial motion of the profile;
it vanishes for tangential defects and has the sign of the correlation
between \(\Pi\) and outward defect.  Neither term has a sign in general;
the geometric content is that a violation of \eqref{eq:c2v2-C1prime}
requires a defect with both a component transverse to the
velocity--vorticity plane where \(|\Omega||V|\) is large, or a radial
component correlated with the Bernoulli pressure, and in either case of
size comparable to the defect itself.
\end{assessment}

% =============================================================================
\section{Bounds on the nonlinearity and the log-integral of the defect}
\label{sec:c2v3-bounds}
% =============================================================================

\begin{lemma}[The nonlinearity is bounded in the Gaussian norm]
\label{lem:c2v3-N-bound}
There is \(C_N\), depending on \(C\), \(C_\infty\), \(\nu\) and the
Calder\'on--Zygmund constant, such that
\(\|N(V(\sigma))\|_G^2\le C_N\) for all \(\sigma\), and also
\(\||\Omega||V|\|_{L^2(G)}^2\le\frac{C_\infty^2}{4\nu}\,d\le C_\infty^2C\).
\end{lemma}

\begin{proof}
\(\|V(\sigma)\|_\infty=(-s)^{1/2}\|W(s)\|_\infty\le C_\infty\), and
\(\|\nabla V(\sigma)\|_{L^2(\R^3)}^2=(-s)^{1/2}\|\nabla W(s)\|_2^2\le C\).
Hence \(\|(V\cdot\nabla)V\|_G^2\le C_\infty^2\|\nabla V\|_G^2\le C_\infty^2C\)
and \(\|\nabla P_V\|_G^2\le\|\nabla P_V\|_{L^2(\R^3)}^2\le C_{\rm CZ}^2\|\nabla(V\otimes V)\|_2^2\le4C_{\rm CZ}^2C_\infty^2C\),
using \(\nabla P_V=\partial_i\partial_j(-\Delta)^{-1}\nabla(V_iV_j)\).  The
vorticity bound uses \(|\Omega|\le\sqrt2|\nabla V|\) and
\(\|\nabla V\|_G^2=d/(4\nu)\).
\end{proof}

\begin{theorem}[Log-integral of the defect]
\label{thm:c2v3-log-integral}
Let \(x_0\) be a homogeneous singular point, \(W\in\mathcal F(U)\), and
\(\Delta_{\mathcal L}:=\frac{4\nu C+2C_\Phi}8-\frac{c_\Phi}4\).  For every
interval \(I\subset\R\) of the variable \(\sigma\),
\begin{equation}
 \Bigl|\int_I\bigl(\|\partial_\sigma V\|_G^2+\langle N(V),\partial_\sigma V\rangle_G\bigr)\dd\sigma\Bigr|\le\Delta_{\mathcal L},
 \qquad
 \int_I\|\partial_\sigma V\|_G^2\,\dd\sigma\ \le\ 2\Delta_{\mathcal L}+\int_I\|N(V)\|_G^2\,\dd\sigma\ \le\ 2\Delta_{\mathcal L}+C_N|I| .
 \label{eq:c2v3-log-integral}
\end{equation}
If on \(I\) the cross term satisfies
\(\langle N(V),\partial_\sigma V\rangle_G\ge-(1-\kappa)\|\partial_\sigma V\|_G^2\),
then \(\int_I\|\partial_\sigma V\|_G^2\le\Delta_{\mathcal L}/\kappa\), independently
of \(|I|\).
\end{theorem}

\begin{proof}
The first bound is \eqref{eq:c2v2-second} integrated, with
\(c_\Phi/4\le\mathcal L\le(4\nu C+2C_\Phi)/8\).  For the second,
\(\|\partial_\sigma V\|^2+\langle N,\partial_\sigma V\rangle\ge\|\partial_\sigma V\|^2-\|N\|\|\partial_\sigma V\|\ge\frac12\|\partial_\sigma V\|^2-\frac12\|N\|^2\),
so \(\frac12\int_I\|\partial_\sigma V\|^2\le\Delta_{\mathcal L}+\frac12\int_I\|N\|^2\),
and Lemma~\ref{lem:c2v3-N-bound}.  The last statement is
Theorem~\ref{thm:c2v2-rigidity}(iii) restricted to \(I\).
\end{proof}

\begin{firewall}[Item (AC2): the unconditional bound is of trivial order]
\label{fw:c2v3-trivial}
The unconditional bound \(\int_I\|\partial_\sigma V\|_G^2\le2\Delta_{\mathcal L}+C_N|I|\)
gives a log-average of the defect at most \(C_N+2\Delta_{\mathcal L}/|I|\),
which is of the same order as the trivial bound
\(\|\partial_\sigma V\|_G\le\|L_-V\|_G+\|N(V)\|_G\); it does not make the
defect small on average.  The conditional bound \(\Delta_{\mathcal L}/\kappa\),
independent of the length of the interval, is what forces the defect to
vanish along sequences of scales and is available only under
\eqref{eq:c2v2-C1prime} on the interval.  Along any sequence of
intervals \(I_k\) on which the cross term is nonnegative,
\(\int_{I_k}\|\partial_\sigma V\|_G^2\le\Delta_{\mathcal L}\), so if
\(|I_k|\to\infty\) the defect has vanishing log-average on \(I_k\) and,
by the compactness of the family, a rescaling limit taken at scales in
\(I_k\) where \(\|\partial_\sigma V\|_G\) is small would be self-similar if
the defect functional is lower semicontinuous under local limits; this
semicontinuity is not proved and is item (AC1) of the next version.
\end{firewall}

% =============================================================================
\section{Integrated V3 conclusion and next acquisition order}
\label{sec:c2v3-conclusion}
% =============================================================================

\begin{theorem}[What V3 proves]
\label{thm:c2v3-conclusion}
Relative to the first cycle and V1--V2, modulo (E1)--(E8): the
Lamb--Bernoulli form \eqref{eq:c2v3-cross} of the cross term with its
vanishing cases, the Gaussian bound on the nonlinearity
(Lemma~\ref{lem:c2v3-N-bound}), and the log-integral bounds
\eqref{eq:c2v3-log-integral}.  The full Navier--Stokes
stability/global-regularity theorem remains unproved.
\end{theorem}

\begin{proof}
Theorems~\ref{thm:c2v3-cross}, \ref{thm:c2v3-log-integral}, Lemma~\ref{lem:c2v3-N-bound}.
\end{proof}

\begin{programme}[V4 acquisition order]
\label{prog:c2v4}
\begin{enumerate}[label=\textup{(AC\arabic*)},leftmargin=4.8em]
\item Semicontinuity of the defect.  Prove that
      \(W\mapsto\int_{I}\|\partial_\sigma V\|_G^2\dd\sigma\) is lower
      semicontinuous, or continuous, on a rescaling family in the local
      topology, using the band-limited time-derivative bound of the
      first cycle (Theorem V73-strong) and the strong \(H^1\) convergence;
      conclude that a limit along scales with vanishing defect
      log-average is self-similar.
\item Sign of the Bernoulli radial term.  Determine whether
      \(\int\Pi(z\cdot\partial_\sigma V)G\) can be rewritten, using
      \(\operatorname{div}V=0\) and the profile equation, as a total
      \(\sigma\)-derivative plus a signed term; record the result.
\item The next compulsory companion audit is V5.
\end{enumerate}
\end{programme}

\begin{assessment}[Position after V3]
\label{ass:c2v3-position}
The cross term of the second identity has an exact geometric form; the
defect's log-integral is bounded conditionally by a constant and
unconditionally only at trivial order.  No full stability or
global-regularity claim is made.
\end{assessment}

% =============================================================================
\section*{Version V3 (second cycle) append-only record}
\addcontentsline{toc}{section}{Version V3 (second cycle) append-only record}
% =============================================================================

The complete V2 source of the second cycle was retained before this
continuation.  Its SHA--256 digest was
\begin{center}\scriptsize\ttfamily
44582dae3a2d0dcca94213863b355223b36f2be9aa0bef5575c36775de086c0d.
\end{center}
V3 decomposed the cross term of the second identity into a Lamb triple
product and a Bernoulli radial term and bounded the log-integral of the
self-similarity defect.  No full regularity claim is made.

% =============================================================================
\section{V4: semicontinuity of the defect and the self-similar dichotomy on
a rescaling family}
\label{sec:c2v4-entry}
% =============================================================================

Items (AC1) and (AC2) of Programme~\ref{prog:c2v4} are carried out.  The
log-integral of the Gaussian norm of the self-similarity defect over a
fixed window is shown to be lower semicontinuous on a rescaling family
in the local topology; consequently, if a family contains elements with
arbitrarily small defect on windows of a fixed length, it contains a
backward self-similar element.  With the integrated second identity this
gives a dichotomy: either the rescaling family at a homogeneous singular
point contains a nonzero backward self-similar element, or on every
element the cross term of the second identity is negative on
log-average at a fixed rate, that is, the nonlinearity is anti-aligned
with the self-similarity defect.  The Bernoulli radial term is rewritten
exactly as a quarter of the scale derivative of the flux term minus the
correlation of the time derivative of the Bernoulli pressure with the
radial momentum; no sign follows.  No global regularity claim is made.
The next version is the compulsory companion audit.

% =============================================================================
\section{Uniform bounds and weak convergence of the defect}
\label{sec:c2v4-semicontinuity}
% =============================================================================

Throughout, \(x_0\) is a homogeneous singular point of a type-I solution
and \(\mathcal F(U)\) its rescaling family.  Every \(W\in\mathcal F(U)\)
and every rescaling \(U_k\) of the original solution satisfies, uniformly,
\begin{equation}
 |W(s,y)|\le C_\infty(-s)^{-1/2},\quad
 |\nabla W(s,y)|\le C_\infty'(-s)^{-1},\quad
 |\nabla^2W(s,y)|\le C_\infty''(-s)^{-3/2},\quad
 \|\nabla W(s)\|_{L^2(\R^3)}^2\le C(-s)^{-1/2},
 \label{eq:c2v4-uniform}
\end{equation}
by the scale-invariant derivative bounds of (E5) as used in Theorem
V46-Linfty-upgrade of the first cycle, and their persistence under local
limits.  For the profile these read \(|V|\le C_\infty\),
\(|\nabla V|\le C_\infty'\), \(|\nabla^2V|\le C_\infty''\),
\(\|\nabla V\|_{L^2(\R^3)}^2\le C\), uniformly in \(\sigma\).

\begin{proposition}[Weak convergence and lower semicontinuity of the defect]
\label{prop:c2v4-lsc}
Let \(W_k\to W\) locally in \(\mathcal F(U)\), or \(W_k=U_k\to W\) with
\(U_k\) rescalings of the original solution, in the sense of Theorem
V41-ancient-limit of the first cycle.  Let \(V_k,V\) be the profiles and
\(I\subset\R\) a compact interval of \(\sigma\).  Then
\(\partial_\sigma V_k\rightharpoonup\partial_\sigma V\) weakly in \(L^2(I;L^2(G))\),
and
\begin{equation}
 \int_I\|\partial_\sigma V\|_G^2\,\dd\sigma\ \le\ \liminf_{k\to\infty}\int_I\|\partial_\sigma V_k\|_G^2\,\dd\sigma .
 \label{eq:c2v4-lsc}
\end{equation}
\end{proposition}

\begin{proof}
In the original variables, on the compact time interval \(J\)
corresponding to \(I\), \(\partial_sW_k=\nu\Delta W_k-(W_k\cdot\nabla)W_k-\nabla P_k\).
By \eqref{eq:c2v4-uniform}, \(\Delta W_k\) and \((W_k\cdot\nabla)W_k\) are
bounded in \(L^\infty(J\times\R^3)\), and \(\nabla P_k=\partial_i\partial_j(-\Delta)^{-1}\nabla(W_{k,i}W_{k,j})\)
is bounded in \(L^2(J;L^2(\R^3))\) by the Calder\'on--Zygmund bound and
\(\|\nabla(W_k\otimes W_k)\|_2\le2\|W_k\|_\infty\|\nabla W_k\|_2\).  Hence
\(\partial_sW_k\) is bounded in \(L^2(J;L^2(\psi))\), the Gaussian tails
being uniformly small: \(\int_{|y|>R}(|\Delta W_k|^2+|(W_k\cdot\nabla)W_k|^2)\psi\le C\int_{|y|>R}\psi\to0\)
and \(\int_{|y|>R}|\nabla P_k|^2\psi\le\|\nabla P_k\|_2^2\sup_{|y|>R}\psi\to0\).
A subsequence converges weakly in \(L^2(J;L^2(\psi))\); the limit is
\(\partial_sW\) in the sense of distributions, because \(W_k\to W\) in
\(L^2(J;L^2(B_R))\) for every \(R\).  Thus the whole sequence converges
weakly to \(\partial_sW\).  The change of variables to the profile is a
fixed bounded linear map on each time slice (a dilation and a
multiplication by a power of \(-s\), bounded above and below on \(J\)),
so \(\partial_\sigma V_k\rightharpoonup\partial_\sigma V\) in
\(L^2(I;L^2(G))\), and \eqref{eq:c2v4-lsc} is the weak lower
semicontinuity of the norm.
\end{proof}

\begin{theorem}[Small defect on a window produces a self-similar element]
\label{thm:c2v4-small-defect}
Let \(\ell>0\).  Suppose there are \(W_k\in\mathcal F(U)\) and
\(\sigma_k\in\R\) with
\(\int_{\sigma_k}^{\sigma_k+\ell}\|\partial_\sigma V_k\|_G^2\,\dd\sigma\to0\).
Then \(\mathcal F(U)\) contains a nonzero backward self-similar element.
The same holds with \(W_k\) replaced by rescalings of the original
solution at \(x_0\).
\end{theorem}

\begin{proof}
Replace \(W_k\) by its rescaling \(W_k^{\mu_k}\), \(\mu_k=e^{-\sigma_k/2}\),
which is in \(\mathcal F(U)\) and whose profile is
\(V_k(\cdot,\sigma+\sigma_k)\); the hypothesis becomes
\(\int_0^\ell\|\partial_\sigma V_k\|_G^2\to0\).  Extract a local limit
\(W'\in\mathcal F(U)\) (Proposition V83-family(iv) of the first cycle).
By \eqref{eq:c2v4-lsc}, \(\int_0^\ell\|\partial_\sigma V'\|_G^2=0\), so
\(V'(\cdot,\sigma)=V_0\) is constant for \(\sigma\in[0,\ell]\):
\(W'(s,y)=(-s)^{-1/2}V_0(y/(-s)^{1/2})\) on the time interval
\([s_1,s_2]=[-1,-e^{-\ell}]\).  The backward self-similar field
\(W_{\rm ss}(s,y):=(-s)^{-1/2}V_0(y/(-s)^{1/2})\) solves the Navier--Stokes
equations on \((-\infty,0)\) (Lemma~\ref{lem:c2v2-dictionary}, since
\(L_-V_0=N(V_0)\) on \([s_1,s_2]\) is a statement about \(V_0\) alone), is
bounded with bounded gradient on every half-line, and coincides with
\(W'\) on \([s_1,s_2]\).  Forward uniqueness (E8) in \(C_tL^6\) from time
\(s_2\) gives \(W'=W_{\rm ss}\) on \([s_2,0)\); backward uniqueness (E6)
for the difference \(D=W'-W_{\rm ss}\), which satisfies
\(|\partial_sD-\nu\Delta D|\le M(|D|+|\nabla D|)\) on \((-\infty,s_1]\times\R^3\)
with \(M\) from the bounds \eqref{eq:c2v4-uniform} on the half-line and
\(D(s_1)=0\), gives \(W'=W_{\rm ss}\) on \((-\infty,s_1]\).  Hence \(W'\) is
backward self-similar on \((-\infty,0)\), and nonzero by the pinching
\(\Phi_{W'}\ge c_\Phi\).
\end{proof}

% =============================================================================
\section{The self-similar dichotomy}
\label{sec:c2v4-dichotomy}
% =============================================================================

\begin{theorem}[Dichotomy on a rescaling family]
\label{thm:c2v4-dichotomy}
Let \(x_0\) be a homogeneous singular point.  Exactly one of the
following holds.
\begin{enumerate}[label=\textup{(\Alph*)},leftmargin=3.2em]
\item \(\mathcal F(U)\) contains a nonzero backward self-similar element
      of \(\mathcal A_\nu(C_I^2\nu^{3/2})\) with the properties of
      Theorem~\ref{thm:c2v1-ancient}(3).
\item For every \(\ell>0\) there is \(\delta(\ell)>0\) such that every
      \(W\in\mathcal F(U)\) and every \(\sigma\) satisfy
      \(\int_\sigma^{\sigma+\ell}\|\partial_\sigma V\|_G^2\ge\delta(\ell)\);
      consequently, for every \(W\in\mathcal F(U)\), every \(\sigma\), and
      every \(L\ge\ell\),
      \begin{equation}
       \boxed{\quad
       \int_\sigma^{\sigma+L}\langle N(V),\partial_\sigma V\rangle_G\,\dd\sigma
       \ \le\ \Delta_{\mathcal L}-\frac{\delta(\ell)}{\ell}\,(L-\ell):
       \quad}
       \label{eq:c2v4-anti-aligned}
      \end{equation}
      the nonlinearity is anti-aligned with the self-similarity defect
      on log-average, at the rate \(\delta(\ell)/\ell\) per unit of
      \(\sigma\).
\end{enumerate}
\end{theorem}

\begin{proof}
If (A) fails, Theorem~\ref{thm:c2v4-small-defect} gives, for each
\(\ell\), a positive lower bound \(\delta(\ell)\) on
\(\int_\sigma^{\sigma+\ell}\|\partial_\sigma V\|_G^2\) uniform over the family
and over \(\sigma\).  Cover \([\sigma,\sigma+L]\) by
\(\lfloor L/\ell\rfloor\ge(L-\ell)/\ell\) disjoint windows of length
\(\ell\), so \(\int_\sigma^{\sigma+L}\|\partial_\sigma V\|_G^2\ge\delta(\ell)(L-\ell)/\ell\),
and insert into the first bound of \eqref{eq:c2v3-log-integral}.
Conversely, if (A) holds the self-similar element has zero defect, so
(B) fails; the alternatives are exclusive.
\end{proof}

\begin{assessment}[The two alternatives]
\label{ass:c2v4-alternatives}
Alternative (A) is the existence of a nonzero backward self-similar
solution with locally bounded energy and the properties of the ancient
limit.  External results (Ne\v{c}as--R\r{u}\v{z}i\v{c}ka--\v{S}ver\'{a}k
for \(L^3\) profiles; Tsai for profiles with locally bounded energy)
exclude such solutions; they are not imported, and an in-series proof
would be required to use them.  If (A) is excluded, alternative (B) is
a quantitative statement about every element of every rescaling family
at every homogeneous point: the cross term of the second identity is
negative on log-average with a fixed rate.  In the Lamb--Bernoulli form
\eqref{eq:c2v3-cross}, this says that the Gaussian average of the
oriented volume of \((\Omega,V,\partial_\sigma V)\) plus the Bernoulli
radial work is negative on average at a definite rate, on every scale
window of every element.  The first cycle's closing condition
\eqref{eq:c2v2-C1prime} is then false on every family, and the route
to self-similarity through a monotone \(\mathcal L\) is closed; what
remains is to derive a contradiction from (B) itself, which is a
statement with a sign, the first of its kind for the nonlinearity in
this series.  No global regularity claim is made.
\end{assessment}

% =============================================================================
\section{The Bernoulli radial term as a derivative}
\label{sec:c2v4-Bernoulli}
% =============================================================================

\begin{proposition}[Exact rewriting of the Bernoulli radial term]
\label{prop:c2v4-Bernoulli}
For every \(\sigma\), with \(j(\sigma)=4\langle V,N(V)\rangle_G=\frac1\nu\int(|V|^2+2P_V)(z\cdot V)G\)
the scale-invariant flux term,
\begin{equation}
 \frac1{2\nu}\int\Pi\,(z\cdot\partial_\sigma V)\,G
 =\frac14\frac{\dd j}{\dd\sigma}-\frac1{2\nu}\int(\partial_\sigma\Pi)\,(z\cdot V)\,G .
 \label{eq:c2v4-Bernoulli}
\end{equation}
\end{proposition}

\begin{proof}
\(z\cdot\partial_\sigma V=\partial_\sigma(z\cdot V)\), and
\(\int\Pi(z\cdot V)G=\frac12\int(|V|^2+2P_V)(z\cdot V)G=\frac\nu2j\); differentiate
the product under the integral (justified as in
Theorem~\ref{thm:c2v2-second}, \(\partial_\sigma P_V\) being the Riesz
transform of \(\partial_\sigma(V\otimes V)\), which is in \(L^2(\R^3)\) by the
bounds \eqref{eq:c2v4-uniform} and \(\partial_\sigma V\in L^2(\R^3)\)
on compact \(\sigma\)-intervals).
\end{proof}

\begin{firewall}[Item (AC2)]
\label{fw:c2v4-AC2}
The rewriting \eqref{eq:c2v4-Bernoulli} isolates a total derivative,
whose integral over a window is bounded by \(\frac14(\sup|j|+\sup|j|)\le\frac12C_J\),
and leaves the correlation of \(\partial_\sigma\Pi\) with the radial
momentum, which has no sign; \(\partial_\sigma\Pi=\partial_\sigma P_V+V\cdot\partial_\sigma V\)
involves the defect linearly.  No sign is obtained; item (AC2) is closed
with this record.
\end{firewall}

% =============================================================================
\section{Integrated V4 conclusion and next acquisition order}
\label{sec:c2v4-conclusion}
% =============================================================================

\begin{theorem}[What V4 proves]
\label{thm:c2v4-conclusion}
Relative to the first cycle and V1--V3, modulo (E1)--(E8): the weak
convergence and lower semicontinuity of the defect
(Proposition~\ref{prop:c2v4-lsc}), the small-defect theorem
(Theorem~\ref{thm:c2v4-small-defect}), the self-similar dichotomy
(Theorem~\ref{thm:c2v4-dichotomy}) with the anti-alignment bound
\eqref{eq:c2v4-anti-aligned}, and the rewriting
\eqref{eq:c2v4-Bernoulli}.  The full Navier--Stokes
stability/global-regularity theorem remains unproved.
\end{theorem}

\begin{proof}
The cited results.
\end{proof}

\begin{programme}[V5 acquisition order, with compulsory companion audit]
\label{prog:c2v5}
\begin{enumerate}[label=\textup{(AC\arabic*)},leftmargin=4.8em]
\item Perform the compulsory review of the current SM and YM notes;
      record digests; import nothing numerical.
\item Consolidate V1--V4 of the second cycle into one statement on the
      self-similar variables: dictionary, the two identities, the
      Lamb--Bernoulli form, the semicontinuity, and the dichotomy.
\item Record, as the next target, the exact statement of alternative
      (B) in the original variables, and whether the anti-alignment can
      be tested against the first cycle's satellite and lacunary
      structures, or against the uniform concentration
      \(d_W(\lambda)\ge\varepsilon'/2\).
\item The next compulsory companion audit after V5 is V10.
\end{enumerate}
\end{programme}

\begin{assessment}[Position after V4]
\label{ass:c2v4-position}
Every rescaling family at a homogeneous point either contains a
backward self-similar element or has its nonlinearity anti-aligned with
the self-similarity defect at a fixed rate on every element.  No full
stability or global-regularity claim is made.
\end{assessment}

% =============================================================================
\section*{Version V4 (second cycle) append-only record}
\addcontentsline{toc}{section}{Version V4 (second cycle) append-only record}
% =============================================================================

The complete V3 source of the second cycle was retained before this
continuation.  Its SHA--256 digest was
\begin{center}\scriptsize\ttfamily
d64af65cea32458722b85b6ecaac7aa892426baebe20400b608c187c25f3cb5c.
\end{center}
V4 proved the lower semicontinuity of the self-similarity defect on a
rescaling family, the small-defect theorem, and the self-similar
dichotomy with its anti-alignment bound.  No full regularity claim is
made.

% =============================================================================
\section{V5: compulsory companion audit and the consolidated statement on
the self-similar variables}
\label{sec:c2v5-entry}
% =============================================================================

This is the fifth version of the second cycle.  The compulsory review of
the two companion programmes is performed first.  V1--V4 are then
consolidated into one statement on the self-similar variables: the
dictionary between the Gaussian ledger and the profile, the two exact
identities, the Lamb--Bernoulli form of the cross term, the
semicontinuity of the defect, and the self-similar dichotomy.  The
second alternative of the dichotomy is then written in the original
variables, and its relation to the structures of the first cycle is
recorded.  No global regularity claim is made.  The next compulsory
companion audit is V10.

% =============================================================================
\section{Compulsory V5 review of the current companion notes}
\label{sec:c2v5-companion-audit}
% =============================================================================

The current companion files in \texttt{latex\_notes} were inspected
read-only.  The frozen checkpoints are
\begin{align}
 &\texttt{Renewal\_SM\_Einstein\_Continuum\_Research\_Notes\_V85.tex},
 \notag\\[-2pt]
 &\qquad
 \texttt{f94726081bff773b736a4996af8193ba72220e4a76c05923e3ff804a8caaade6},
 \label{eq:c2v5-SM-checkpoint}\\
 &\texttt{Renewal\_Yang\_Mills\_Mass\_Gap\_Research\_Notes\_V60.tex},
 \notag\\[-2pt]
 &\qquad
 \texttt{14ec8cc3515d32c3e2a47d8e8236b5fb922b516ca482353b585c83a32d074500}.
 \label{eq:c2v5-YM-checkpoint}
\end{align}

\emph{SM V83--V85.}  The SM programme computed a two-source York
response and a first-order beat, then performed its own compulsory
audit, reading the present cycle's V1 (recording its digest
\texttt{27a17aaf\ldots} and the scalar ledger \(\varphi'=2\varphi-\xi\)
with its log-average law correctly); it proved that at every finite
cutoff its matched vector, conformal, and fixed-volume York blocks have
an invertible triangular derivative at the homogeneous base point, that
the resulting local coarea density is positive, bounded, and covariant
on compact regular source boxes, and that it supports a positive
self-adjoint volume observable; its next gate is determinant control at
the first constraint fold.

\emph{YM V60.}  The YM programme performed its own audit, reading the
first cycle's V96 (correctly recording that the rate-free identity is
separated from the density conclusions that need a type-I rate); it
proved a fractional-moment hierarchy for its small-clock L\'evy measure,
a dimension-free compact-group heat-increment bound, an all-arity
same-jump hyperedge bound in weighted Pr\"ufer form, an exact
endpoint-factorizing source with its connected forest expansion, and a
no-go for termwise differentiation of the outer exponential, isolating
a resummation gate.

\begin{theorem}[V5 companion-audit decision]
\label{thm:c2v5-companion-decision}
No SM V83--V85 York-response, coarea, or determinant constant and no YM
V60 moment, heat-increment, or forest estimate is imported.  Neither
bears on the dichotomy of Theorem~\ref{thm:c2v4-dichotomy} or on the
gates of the first cycle.  Both companion readings of the NS notes are
accurate and need no correction.
\end{theorem}

\begin{proof}
The SM results concern the derivative of constraint maps and coarea
densities of a gravitational model; the YM results concern L\'evy
measures and forest expansions of a lattice clock.  None is a
Navier--Stokes identity.  The readings are checked against
Theorem~\ref{thm:c2v1-log-average} and Theorem V96-conclusion of the
first cycle.
\end{proof}

\begin{assessment}[Direction after the audit]
\label{ass:c2v5-direction}
Both companion programmes record, independently, that averaged
positivity cannot replace a pointwise or rank statement and that exact
identities do not supply uniform estimates by themselves; the present
cycle's V2 reached the same conclusion for the weighted energy
(Firewall~\ref{fw:c2v2-AC}) and moved to a quantity with a rigid
equality case.  The direction is unchanged.
\end{assessment}

% =============================================================================
\section{The self-similar variables, consolidated}
\label{sec:c2v5-consolidated}
% =============================================================================

\begin{theorem}[Self-similar variables, V1--V4 of the second cycle]
\label{thm:c2v5-consolidated}
Let \(W\in\mathcal A_\nu(C)\) with the Morrey bound and the uniform bounds
\eqref{eq:c2v4-uniform}, \(V\) its profile \eqref{eq:c2v2-profile},
\(G=e^{-|z|^2/(4\nu)}\), \(L_-=\nu\Delta-\frac12z\cdot\nabla-\frac12\),
\(N(V)=(V\cdot\nabla)V+\nabla P_V=\Omega\times V+\nabla\Pi\).  Modulo
(E1)--(E8):
\begin{enumerate}[label=\textup{(\arabic*)},leftmargin=3.2em]
\item \emph{Dictionary.}  \(\partial_\sigma V=L_-V-N(V)\); \(L_-\) is
      \(G\)-symmetric; \(\Phi_W=\|V\|_G^2\), \(\lambda D_W=4\nu\|\nabla V\|_G^2\),
      \(\lambda\mathcal J_W=4\langle V,N(V)\rangle_G\); \(W\) is backward
      self-similar if and only if \(\partial_\sigma V\equiv0\).
\item \emph{First identity.}  \(\frac{\dd}{\dd\sigma}\|V\|_G^2=-8\mathcal L(V)-2\langle V,N(V)\rangle_G\)
      with \(\mathcal L=\frac\nu2\|\nabla V\|_G^2+\frac14\|V\|_G^2\); in the
      log-scale \(\tau\), \(\varphi'=2\varphi-\xi\) and the log-average of
      \(\xi\) equals twice that of \(\varphi\) up to the endpoint difference.
\item \emph{Second identity.}  \(\frac{\dd}{\dd\sigma}\mathcal L(V)=-\|\partial_\sigma V\|_G^2-\langle N(V),\partial_\sigma V\rangle_G\),
      with \(\langle N(V),\partial_\sigma V\rangle_G=\int\det(\Omega,V,\partial_\sigma V)G+\frac1{2\nu}\int\Pi(z\cdot\partial_\sigma V)G\),
      the second term being \(\frac14j'-\frac1{2\nu}\int(\partial_\sigma\Pi)(z\cdot V)G\).
\item \emph{Bounds.}  \(\|N(V)\|_G^2\le C_N\); at a homogeneous point
      \(c_\Phi/4\le\mathcal L\le(4\nu C+2C_\Phi)/8\) on \(\mathcal F(U)\);
      \(\int_I\|\partial_\sigma V\|_G^2\le2\Delta_{\mathcal L}+C_N|I|\); under
      \eqref{eq:c2v2-C1prime} on \(I\), \(\int_I\|\partial_\sigma V\|_G^2\le\Delta_{\mathcal L}/\kappa\).
\item \emph{Rigidity and semicontinuity.}  Under \eqref{eq:c2v2-C1prime},
      \(\mathcal L\) is monotone and constant only on self-similar
      elements, and \(\mathcal F(U)\) contains a nonzero backward
      self-similar element.  The defect log-integral over a fixed
      window is weakly lower semicontinuous on \(\mathcal F(U)\), and
      elements with arbitrarily small defect on windows of fixed length
      force a self-similar element.
\item \emph{Dichotomy.}  Either \(\mathcal F(U)\) contains a nonzero
      backward self-similar element, or the defect is bounded below by
      \(\delta(\ell)\) on every window of length \(\ell\) and the cross term
      is negative on log-average at rate \(\delta(\ell)/\ell\) on every
      element.
\end{enumerate}
\end{theorem}

\begin{proof}
(1) Lemma~\ref{lem:c2v2-dictionary}.  (2) Remark~\ref{rem:c2v2-first-identity},
Theorem~\ref{thm:c2v1-log-average}.  (3) Theorems~\ref{thm:c2v2-second},
\ref{thm:c2v3-cross}, Proposition~\ref{prop:c2v4-Bernoulli}.
(4) Lemma~\ref{lem:c2v3-N-bound}, Definition~\ref{def:c2v2-L},
Theorem~\ref{thm:c2v3-log-integral}.  (5) Theorem~\ref{thm:c2v2-rigidity},
Corollary~\ref{cor:c2v2-closing}, Proposition~\ref{prop:c2v4-lsc},
Theorem~\ref{thm:c2v4-small-defect}.  (6) Theorem~\ref{thm:c2v4-dichotomy}.
\end{proof}

% =============================================================================
\section{Alternative (B) in the original variables}
\label{sec:c2v5-original}
% =============================================================================

\begin{proposition}[The defect and the cross term in the original variables]
\label{prop:c2v5-original}
Let \(Z(s,y):=\partial_sW+\frac{W+y\cdot\nabla W}{2(-s)}\) be the
self-similarity defect of \(W\) in the original variables, so that
\(W\) is backward self-similar if and only if \(Z\equiv0\).  Then
\(\partial_\sigma V(z,\sigma)=(-s)^{3/2}Z(s,(-s)^{1/2}z)\), and for a
window \(I=[\sigma_1,\sigma_1+\ell]\) with corresponding time interval
\(J=[-e^{\ell}(-s_1),\,-(-s_1)]\), \(s_1=-e^{-\sigma_1}\),
\begin{align}
 \int_I\|\partial_\sigma V\|_G^2\,\dd\sigma
 &=\int_J(-s)^{1/2}\int_{\R^3}|Z(s,y)|^2\,\psi(y,s)\,\dd y\,\dd s,
 \label{eq:c2v5-defect-original}\\
 \int_I\langle N(V),\partial_\sigma V\rangle_G\,\dd\sigma
 &=\int_J(-s)^{1/2}\int_{\R^3}\bigl((W\cdot\nabla)W+\nabla P\bigr)\cdot Z\,\psi(y,s)\,\dd y\,\dd s .
 \label{eq:c2v5-cross-original}
\end{align}
Alternative (B) of Theorem~\ref{thm:c2v4-dichotomy} states that, for
every element of the family, every \(s_1<0\), and every \(\ell>0\),
\(\int_{e^{\ell}s_1}^{s_1}(-s)^{1/2}\int|Z|^2\psi\ge\delta(\ell)\), and that
the space-time integral \eqref{eq:c2v5-cross-original} over
\([e^{L}s_1,s_1]\) is at most \(\Delta_{\mathcal L}-\delta(\ell)(L-\ell)/\ell\).
\end{proposition}

\begin{proof}
From the proof of Lemma~\ref{lem:c2v2-dictionary},
\(\partial_sW=(-s)^{-3/2}[\frac12V+\frac12z\cdot\nabla V+\partial_\sigma V]\)
and \(\frac{W+y\cdot\nabla W}{2(-s)}=(-s)^{-3/2}\frac12[V+z\cdot\nabla V]\),
so \(Z=(-s)^{-3/2}\partial_\sigma V\).  Then
\(\|\partial_\sigma V\|_G^2=\int(-s)^3|Z|^2G\,\dd z=(-s)^{3/2}\int|Z|^2\psi\,\dd y\)
and \(\dd\sigma=\dd s/(-s)\), which gives \eqref{eq:c2v5-defect-original};
similarly \(N(V)=(-s)^{3/2}[(W\cdot\nabla)W+\nabla P]\) gives
\eqref{eq:c2v5-cross-original}.  The last statement is
Theorem~\ref{thm:c2v4-dichotomy}(B) rewritten.
\end{proof}

\begin{firewall}[Item (AC3): relation to the first cycle's structures]
\label{fw:c2v5-AC3}
The uniform concentration \(d_W(\lambda)\ge\varepsilon'/2\) of the first
cycle bounds the dissipation in every parabolic cylinder from below; it
does not bound the defect \(Z\) from below or above, since a
self-similar element is uniformly concentrated with \(Z\equiv0\).  The
satellite and lacunary structures concern points that are not
homogeneous and are outside the dichotomy.  Hence neither alternative
of Theorem~\ref{thm:c2v4-dichotomy} is tested by the first cycle's
results, and the target of the next versions is the sign statement (B)
itself: a contradiction would require showing that the space-time
integral \eqref{eq:c2v5-cross-original} cannot be negative at a fixed
rate on every window of every element of a compact scale-invariant
family.
\end{firewall}

% =============================================================================
\section{Integrated V5 conclusion and next acquisition order}
\label{sec:c2v5-conclusion}
% =============================================================================

\begin{theorem}[What V5 establishes]
\label{thm:c2v5-conclusion}
Relative to the first cycle and V1--V4: the companion-audit decision
(Theorem~\ref{thm:c2v5-companion-decision}), the consolidated statement
(Theorem~\ref{thm:c2v5-consolidated}), and the original-variable form
of the dichotomy (Proposition~\ref{prop:c2v5-original}).  The full
Navier--Stokes stability/global-regularity theorem remains unproved.
\end{theorem}

\begin{proof}
The cited results.
\end{proof}

\begin{programme}[V6 acquisition order]
\label{prog:c2v6}
\begin{enumerate}[label=\textup{(AC\arabic*)},leftmargin=4.8em]
\item Extremal elements of alternative (B).  On the compact family,
      minimize the windowed defect \(\int_0^\ell\|\partial_\sigma V\|_G^2\);
      show the minimum is attained (lower semicontinuity) and equals
      \(\delta(\ell)\); determine whether the minimizer's rescalings
      give \(\delta(\ell)\) at every window, and whether \(\delta(\ell)\) is
      superadditive in \(\ell\).
\item Long-window limit.  Determine whether \(\delta(\ell)/\ell\) has a
      positive limit \(\underline\delta\) as \(\ell\to\infty\) (the uniform
      defect rate of the family), and record the resulting form of
      \eqref{eq:c2v4-anti-aligned}: cross term at most \(-\underline\delta\)
      per unit \(\sigma\) on average.
\item The next compulsory companion audit is V10.
\end{enumerate}
\end{programme}

\begin{assessment}[Position after V5]
\label{ass:c2v5-position}
The second cycle is consolidated on the self-similar variables and the
dichotomy.  No full stability or global-regularity claim is made.
\end{assessment}

% =============================================================================
\section*{Version V5 (second cycle) append-only record}
\addcontentsline{toc}{section}{Version V5 (second cycle) append-only record}
% =============================================================================

The complete V4 source of the second cycle was retained before this
continuation.  Its SHA--256 digest was
\begin{center}\scriptsize\ttfamily
e15605552fd8602edf5686a8211269f968616d8cb70f9929f32969ab5d03fe1b.
\end{center}
V5 performed the compulsory companion audit, consolidated V1--V4 of the
second cycle, and wrote the dichotomy in the original variables.  No
full regularity claim is made.

% =============================================================================
\section{V6: the minimal windowed defect, its rate, and the dichotomy for the
solution itself}
\label{sec:c2v6-entry}
% =============================================================================

Items (AC1) and (AC2) of Programme~\ref{prog:c2v6} are carried out.  The
minimal windowed defect \(m(\ell)\) over the rescaling family is attained,
is positive for one window length if and only if for all, and is
superadditive in the window length, so that \(m(\ell)/\ell\) converges to
a uniform defect rate \(\underline\delta\in[0,C_N]\), positive exactly in
alternative (B).  The family is then enlarged to the full family of all
ancient limits at the point and their rescaling limits, on which the
same statements hold, and the dichotomy is transferred to the solution
itself: either the point admits a nonzero backward self-similar ancient
limit, or the solution's own self-similarity defect, in every parabolic
window of fixed log-length at every small scale, is bounded below by a
fixed positive constant.  No global regularity claim is made.  The next
compulsory companion audit is V10.

% =============================================================================
\section{The minimal windowed defect}
\label{sec:c2v6-minimal}
% =============================================================================

Let \(x_0\) be a homogeneous singular point.  Let \(\mathcal F(x_0)\) be
the set of all local limits of rescalings
\(\rho_ju(x_0+\rho_jy,T_*+\rho_j^2s)\), \(\rho_j\to0\), together with all
local limits of rescalings of such limits.  It contains the family
\(\mathcal F(U)\) of the first cycle for every ancient limit \(U\) at
\(x_0\), is nonempty, scale-invariant, closed under local limits of
rescalings, sequentially compact in the local topology, and every
element satisfies the uniform bounds \eqref{eq:c2v4-uniform}, the
Morrey bound, the uniform concentration \(d_W(\lambda)\ge\varepsilon'/2\),
and the pinching \(c_\Phi\le\Phi_W\le C_\Phi\) with \(c_\Phi=c_\Phi(x_0)>0\):
the proofs of Proposition V83-family and Theorem V91-pinched of the
first cycle use only these properties and the homogeneity of \(x_0\).
Every result of V2--V5 stated for \(\mathcal F(U)\) holds for
\(\mathcal F(x_0)\) with the same proofs.  For \(\ell>0\) put
\begin{equation}
 m(\ell):=\inf_{W\in\mathcal F(x_0)}\int_0^\ell\|\partial_\sigma V\|_G^2\,\dd\sigma .
 \label{eq:c2v6-m}
\end{equation}

\begin{proposition}[Properties of the minimal windowed defect]
\label{prop:c2v6-m}
\begin{enumerate}[label=\textup{(\roman*)},leftmargin=3.2em]
\item \emph{Scale invariance.}  For every \(W\in\mathcal F(x_0)\) and every
      \(\sigma\), \(\int_\sigma^{\sigma+\ell}\|\partial_\sigma V\|_G^2\ge m(\ell)\), and
      \(m(\ell)=\inf_{W,\sigma}\int_\sigma^{\sigma+\ell}\|\partial_\sigma V\|_G^2\).
\item \emph{Attainment.}  There is \(W_\ell\in\mathcal F(x_0)\) with
      \(\int_0^\ell\|\partial_\sigma V_\ell\|_G^2=m(\ell)\).
\item \emph{Zero or positive.}  If \(m(\ell)=0\) for some \(\ell>0\) then
      \(\mathcal F(x_0)\) contains a nonzero backward self-similar element
      and \(m(\ell')=0\) for all \(\ell'>0\); otherwise \(m(\ell)>0\) for all
      \(\ell>0\).
\item \emph{Superadditivity and the rate.}  \(m(\ell_1+\ell_2)\ge m(\ell_1)+m(\ell_2)\),
      \(m(\ell)\le2\Delta_{\mathcal L}+C_N\ell\), and
      \begin{equation}
       \underline\delta:=\lim_{\ell\to\infty}\frac{m(\ell)}\ell=\sup_{\ell>0}\frac{m(\ell)}\ell\ \in[0,C_N]
       \label{eq:c2v6-rate}
      \end{equation}
      exists, with \(\underline\delta>0\) if and only if \(m(\ell)>0\) for
      some (equivalently all) \(\ell\).
\end{enumerate}
\end{proposition}

\begin{proof}
(i) The profile of \(W^\mu\) is \(V(\cdot,\sigma-2\log\mu)\), so windows
of \(W\) are windows \([0,\ell]\) of its rescalings, which lie in
\(\mathcal F(x_0)\).  (ii) Take a minimizing sequence, extract a local
limit in \(\mathcal F(x_0)\) by compactness, and apply the lower
semicontinuity \eqref{eq:c2v4-lsc}; the limit is in the family, so its
defect is at least \(m(\ell)\).  (iii) If \(m(\ell)=0\), the minimizer of
(ii) has vanishing defect on \([0,\ell]\), and the proof of
Theorem~\ref{thm:c2v4-small-defect} makes it backward self-similar on
\((-\infty,0)\) and nonzero; then its defect vanishes on every window,
so \(m(\ell')=0\) for all \(\ell'\).  (iv) Split \([0,\ell_1+\ell_2]\) into
\([0,\ell_1]\) and \([\ell_1,\ell_1+\ell_2]\) and apply (i) to the second
window.  The upper bound is \eqref{eq:c2v3-log-integral}.  Fekete's
lemma for the superadditive function \(m\) gives the limit
\(\lim m(\ell)/\ell=\sup m(\ell)/\ell\), finite by the upper bound, and
positive if and only if some \(m(\ell)>0\).
\end{proof}

\begin{corollary}[The dichotomy with the uniform rate]
\label{cor:c2v6-dichotomy-rate}
Exactly one of the following holds for a homogeneous singular point \(x_0\).
\begin{enumerate}[label=\textup{(\Alph*)},leftmargin=3.2em]
\item \(\mathcal F(x_0)\) contains a nonzero backward self-similar element.
\item \(\underline\delta>0\), and for every \(W\in\mathcal F(x_0)\), every
      \(\sigma\), and every \(L>0\),
      \begin{equation}
       \int_\sigma^{\sigma+L}\|\partial_\sigma V\|_G^2\ \ge\ m(L)\ \ge\ \underline\delta L-o(L),
       \qquad
       \int_\sigma^{\sigma+L}\langle N(V),\partial_\sigma V\rangle_G\ \le\ \Delta_{\mathcal L}-m(L),
       \label{eq:c2v6-B}
      \end{equation}
      so that the log-average of the cross term is at most
      \(-\underline\delta+o(1)+\Delta_{\mathcal L}/L\) as \(L\to\infty\),
      uniformly in \(W\) and \(\sigma\).
\end{enumerate}
\end{corollary}

\begin{proof}
Proposition~\ref{prop:c2v6-m} and the first bound of
\eqref{eq:c2v3-log-integral}; \(m(L)\ge\underline\delta L-o(L)\) is the
convergence in \eqref{eq:c2v6-rate}.
\end{proof}

\begin{firewall}[What the minimizer does not give]
\label{fw:c2v6-minimizer}
The minimizer \(W_\ell\) attains \(m(\ell)\) on the window \([0,\ell]\);
on other windows its defect is at least \(m(\ell)\) but is not shown to
equal it, and no element is shown to attain \(m(\ell)\) on every
window.  Whether \(\underline\delta\) is attained as the exact log-average
of the defect of some element (a ``uniformly minimally non-self-similar''
element) is not decided.  The rate \(\underline\delta\) depends on the
point through the family.
\end{firewall}

% =============================================================================
\section{The dichotomy for the solution itself}
\label{sec:c2v6-solution}
% =============================================================================

For the solution \(u\) and the point \(x_0\), define the self-similarity
defect relative to \((x_0,T_*)\) and its windowed Gaussian log-integral,
\begin{equation}
 Z_u(x,t):=\partial_tu+\frac{u+(x-x_0)\cdot\nabla u}{2(T_*-t)},
 \qquad
 \mathcal Z_u(\lambda;\ell):=\int_{T_*-e^{\ell}\lambda^2}^{T_*-\lambda^2}(T_*-t)^{1/2}
 \int_{x_0+[-\pi,\pi)^3}|Z_u(x,t)|^2e^{-|x-x_0|^2/(4\nu(T_*-t))}\,\dd x\,\dd t ,
 \label{eq:c2v6-solution-defect}
\end{equation}
with \(x-x_0\) taken in the fundamental domain.

\begin{theorem}[Dichotomy for the solution at a homogeneous point]
\label{thm:c2v6-solution}
Let \(x_0\) be a homogeneous singular point of a type-I solution.
Exactly one of the following holds.
\begin{enumerate}[label=\textup{(\Alph*)},leftmargin=3.2em]
\item Some ancient limit at \(x_0\), or a rescaling limit of one, is a
      nonzero backward self-similar solution.
\item For every \(\ell>0\) there is \(\lambda(\ell)>0\) with
      \begin{equation}
       \boxed{\quad
       \mathcal Z_u(\lambda;\ell)\ \ge\ \tfrac12m(\ell)>0
       \qquad\text{for all }0<\lambda\le\lambda(\ell):
       \quad}
       \label{eq:c2v6-solution-B}
      \end{equation}
      the solution is uniformly non-self-similar in every parabolic
      window of log-length \(\ell\) at every small scale, and its
      windowed cross term
      \(\int_{T_*-e^{L}\lambda^2}^{T_*-\lambda^2}(T_*-t)^{1/2}\int((u\cdot\nabla)u+\nabla p)\cdot Z_u\,e^{-|x-x_0|^2/(4\nu(T_*-t))}\)
      is at most \(\Delta_{\mathcal L}'-\frac12m(L)\) for small \(\lambda\),
      with \(\Delta_{\mathcal L}'\) the oscillation bound of the solution's
      \(\mathcal L\) on small scales.
\end{enumerate}
\end{theorem}

\begin{proof}
Let \(U_k(s,y)=\lambda_ku(x_0+\lambda_ky,T_*+\lambda_k^2s)\) with the
Gaussian restricted to \(|y|<\pi/\lambda_k\); by the change of variables
of Proposition~\ref{prop:c2v5-original},
\(\mathcal Z_u(\lambda_k;\ell)=\int_0^\ell\|\partial_\sigma V_k\|_{L^2(G,|z|<\pi e^{\sigma/2}/\lambda_k)}^2\dd\sigma\),
which differs from the full \(L^2(G)\) norm of the periodic extension by
a term tending to zero by the uniform bounds and the Gaussian tail.  If
(B) fails for some \(\ell\), there are \(\lambda_k\to0\) with
\(\mathcal Z_u(\lambda_k;\ell)<\frac12m(\ell)\); a subsequence of \(U_k\)
converges locally to an element \(W'\in\mathcal F(x_0)\), and
Proposition~\ref{prop:c2v4-lsc} gives
\(\int_0^\ell\|\partial_\sigma V'\|_G^2\le\liminf_k\mathcal Z_u(\lambda_k;\ell)\le\frac12m(\ell)\).
If \(m(\ell)>0\) this contradicts \eqref{eq:c2v6-m}; hence \(m(\ell)=0\),
and Proposition~\ref{prop:c2v6-m}(iii) gives (A).  Conversely, if (A)
holds then the self-similar element has zero defect, so \(m(\ell)=0\) and
the strict inequality in \eqref{eq:c2v6-solution-B} fails: the
alternatives are exclusive.  The cross
term statement is the first bound of \eqref{eq:c2v3-log-integral} for the
solution: the second identity holds for \(u\) on the torus with the
periodized weight by the same proof as Theorem~\ref{thm:c2v2-second}
(Theorem V87-derivative-u of the first cycle supplies the torus form of
the first identity, and the second is obtained identically), with
\(\mathcal L_u\) bounded on small scales by the pinching of Theorem
V96-pinched-u and the type-I dissipation bound.
\end{proof}

\begin{assessment}[The two faces of a homogeneous singular point]
\label{ass:c2v6-faces}
A homogeneous singular point of a type-I solution is now of exactly one
of two kinds.  Either it admits a self-similar blow-up profile in the
limit: an ancient limit, or a rescaling limit of one, that is a nonzero
backward self-similar solution with the properties of the ancient
limit.  Or the solution is uniformly non-self-similar near the point at
every small scale, in the quantitative sense \eqref{eq:c2v6-solution-B},
and then its nonlinearity is anti-aligned with its self-similarity
defect at the rate \(\underline\delta\) per unit of log-scale.  The
first kind is excluded by external results on self-similar solutions
with locally bounded energy, not imported.  The second kind is the
open case, and the sign in \eqref{eq:c2v6-B} is the statement to be
contradicted.  No global regularity claim is made.
\end{assessment}

% =============================================================================
\section{Integrated V6 conclusion and next acquisition order}
\label{sec:c2v6-conclusion}
% =============================================================================

\begin{theorem}[What V6 proves]
\label{thm:c2v6-conclusion}
Relative to the first cycle and V1--V5, modulo (E1)--(E8): the properties
of the minimal windowed defect (Proposition~\ref{prop:c2v6-m}), the
dichotomy with the uniform rate (Corollary~\ref{cor:c2v6-dichotomy-rate}),
and the dichotomy for the solution itself
(Theorem~\ref{thm:c2v6-solution}).  The full Navier--Stokes
stability/global-regularity theorem remains unproved.
\end{theorem}

\begin{proof}
The cited results.
\end{proof}

\begin{programme}[V7 acquisition order]
\label{prog:c2v7}
\begin{enumerate}[label=\textup{(AC\arabic*)},leftmargin=4.8em]
\item The anti-alignment in vorticity form.  Using
      \(N(V)=\Omega\times V+\nabla\Pi\) and the Lamb--Bernoulli form,
      write alternative (B) as a statement about
      \(\int\det(\Omega,V,\partial_\sigma V)G\) and the Bernoulli work, and
      determine whether the uniform concentration or the vorticity at
      infinity of the first cycle constrains the triple product.
\item The defect equation.  Differentiate the profile equation in
      \(\sigma\) to obtain the linear equation satisfied by
      \(\partial_\sigma V\), \(\partial_\sigma(\partial_\sigma V)=L_-\partial_\sigma V-N'(V)\partial_\sigma V\),
      and derive the Gaussian energy identity for
      \(\|\partial_\sigma V\|_G^2\); record whether alternative (B) is
      consistent with it, that is, whether a defect bounded below on
      every window can be sustained by the linearized dynamics.
\item The next compulsory companion audit is V10.
\end{enumerate}
\end{programme}

\begin{assessment}[Position after V6]
\label{ass:c2v6-position}
The dichotomy is quantitative, has a uniform rate, and holds for the
solution itself.  No full stability or global-regularity claim is made.
\end{assessment}

% =============================================================================
\section*{Version V6 (second cycle) append-only record}
\addcontentsline{toc}{section}{Version V6 (second cycle) append-only record}
% =============================================================================

The complete V5 source of the second cycle was retained before this
continuation.  Its SHA--256 digest was
\begin{center}\scriptsize\ttfamily
f365c1178c55f56f49846d21e52b8a522c2597351dbd76c4f72d8c411b3262eb.
\end{center}
V6 proved the properties of the minimal windowed defect, the uniform
defect rate, and the dichotomy for the solution itself.  No full
regularity claim is made.

% =============================================================================
\section{V7: the defect equation and the anti-dissipation forced by
alternative (B)}
\label{sec:c2v7-entry}
% =============================================================================

Items (AC1) and (AC2) of Programme~\ref{prog:c2v7} are carried out.
Alternative (B) is written in vorticity form, and it is recorded that
the uniform concentration and the vorticity at infinity of the first
cycle do not constrain the Lamb triple product.  The self-similarity
defect \(Y=\partial_\sigma V\) is shown to satisfy the linearized profile
equation, and its Gaussian energy identity is derived: the linear part
is the Ornstein--Uhlenbeck operator with spectrum at most \(-\frac12\),
so that, absent the nonlinearity, the defect decays exponentially in
the log-scale.  Consequently alternative (B), a defect bounded below on
every window, forces the linearized nonlinearity to be anti-dissipative
against the defect at a definite rate on log-average: a second sign
statement, with an explicit rate, for the open case.  The integrability
needed for the identity is established for the class.  No global
regularity claim is made.  The next compulsory companion audit is V10.

% =============================================================================
\section{Alternative (B) in vorticity form}
\label{sec:c2v7-vorticity}
% =============================================================================

\begin{proposition}[Vorticity form of the anti-alignment]
\label{prop:c2v7-vorticity-form}
In alternative (B) of Corollary~\ref{cor:c2v6-dichotomy-rate}, for every
\(W\in\mathcal F(x_0)\), \(\sigma\), and \(L\),
\begin{equation}
 \int_\sigma^{\sigma+L}\Bigl[\int\det(\Omega,V,\partial_\sigma V)\,G+\frac1{2\nu}\int\Pi\,(z\cdot\partial_\sigma V)\,G\Bigr]\dd\sigma
 \ \le\ \Delta_{\mathcal L}-m(L),
 \label{eq:c2v7-vorticity-form}
\end{equation}
and the Bernoulli term may be replaced by
\(\frac14(j(\sigma+L)-j(\sigma))-\frac1{2\nu}\int_\sigma^{\sigma+L}\!\int(\partial_\sigma\Pi)(z\cdot V)G\),
the first part being bounded by \(\frac12C_J\).
\end{proposition}

\begin{proof}
Theorem~\ref{thm:c2v3-cross}, Proposition~\ref{prop:c2v4-Bernoulli}, and
\eqref{eq:c2v6-B}.
\end{proof}

\begin{firewall}[Item (AC1): the first cycle does not constrain the triple product]
\label{fw:c2v7-AC1}
The uniform concentration \(d_W(\lambda)\ge\varepsilon'/2\) bounds
\(\|\nabla V\|_{L^2(B_1)}\) from below on every unit window of \(\sigma\),
hence \(\|\Omega\|_{L^2(B_1)}\) from below; the vorticity at infinity
says \(\Omega\not\equiv0\) on every exterior region at every \(\sigma\).
Neither constrains the orientation of \(\Omega\) relative to \(V\) and
\(\partial_\sigma V\), and the triple product vanishes identically for
defects in the plane of \(V\) and \(\Omega\) regardless of the size of
\(\Omega\).  Item (AC1) is closed with this record.
\end{firewall}

% =============================================================================
\section{Integrability of higher derivatives in the class}
\label{sec:c2v7-integrability}
% =============================================================================

\begin{lemma}[Second and third derivatives are square-integrable in space-time]
\label{lem:c2v7-higher}
Let \(W\in\mathcal A_\nu(C)\) satisfy the uniform bounds
\eqref{eq:c2v4-uniform} and \(|\nabla^3W(s)|\le C_\infty'''(-s)^{-2}\).
Then for every \(s_1<s_2<0\),
\(\int_{s_1}^{s_2}\bigl(\|\nabla^2W(s)\|_{L^2(\R^3)}^2+\|\nabla^3W(s)\|_{L^2(\R^3)}^2\bigr)\dd s<\infty\),
and consequently \(\partial_sW\), \(\nabla\partial_sW\), \(\partial_sP\in L^2((s_1,s_2);L^2(\R^3))\).
\end{lemma}

\begin{proof}
Let \(\chi_R(y)=\chi(y/R)\) with \(\chi\) a smooth cutoff.  Testing the
equation for \(\nabla W\) against \(\nabla W\chi_R\),
\[
 \frac{\dd}{\dd s}\frac12\int|\nabla W|^2\chi_R
 =-\nu\int|\nabla^2W|^2\chi_R+\frac\nu2\int|\nabla W|^2\Delta\chi_R
 -\int\nabla W:\nabla\bigl((W\cdot\nabla)W\bigr)\chi_R-\int\partial_kW_i\,\partial_k\partial_iP\,\chi_R .
\]
The pressure term equals \(-\int\partial_kW_i\,\partial_iP\,\partial_k\chi_R\) by
\(\operatorname{div}W=0\), bounded by \(CR^{-1}\|\nabla W\|_2\|\nabla P\|_2\to0\);
the cutoff term by \(CR^{-2}\|\nabla W\|_2^2\to0\); and
\(|\int\nabla W:\nabla((W\cdot\nabla)W)\chi_R|\le\|\nabla W\|_\infty\|\nabla W\|_2^2+\|W\|_\infty\|\nabla W\|_2\|\nabla^2W\chi_R^{1/2}\|_2\),
whose last term is absorbed with \(\frac\nu2\int|\nabla^2W|^2\chi_R\).
Integrating over \([s_1,s_2]\), on which all norms are bounded, and
letting \(R\to\infty\) gives \(\nabla^2W\in L^2((s_1,s_2);L^2)\).  The same
computation for \(\nabla^2W\), with the cubic terms bounded by
\(\|\nabla W\|_\infty\|\nabla^2W\|_2^2+\|W\|_\infty\|\nabla^2W\|_2\|\nabla^3W\chi_R^{1/2}\|_2\)
and the pressure term \(-\int\partial_k\partial_lW_i\,\partial_l\partial_iP\,\partial_k\chi_R\)
bounded by \(CR^{-1}\|\nabla^2W\|_2\|\nabla^2P\|_2\) with
\(\|\nabla^2P\|_2\le C_{\rm CZ}\|\nabla^2(W\otimes W)\|_2\le C(\|W\|_\infty\|\nabla^2W\|_2+\|\nabla W\|_\infty\|\nabla W\|_2)\),
gives \(\nabla^3W\in L^2((s_1,s_2);L^2)\).  Then
\(\partial_sW=\nu\Delta W-(W\cdot\nabla)W-\nabla P\in L^2((s_1,s_2);L^2)\) since
each term is; \(\nabla\partial_sW\) likewise with \(\nabla^3W\),
\(\nabla((W\cdot\nabla)W)\), \(\nabla^2P\); and
\(\partial_sP=\partial_i\partial_j(-\Delta)^{-1}\partial_s(W_iW_j)\) with
\(\partial_s(W\otimes W)\in L^2\) by \(\|W\|_\infty\|\partial_sW\|_2\).
\end{proof}

\begin{remark}
\label{rem:c2v7-third}
The pointwise bound on \(\nabla^3W\) with the rate \((-s)^{-2}\) is the
scale-invariant derivative bound of (E5) one order beyond those used in
the first cycle; it holds by the same local regularity theory and is
used only to bound \(\|\nabla^3W\|_\infty\) on half-lines in the
absorption, where in fact only \(\|\nabla W\|_\infty\), \(\|W\|_\infty\)
enter, so the lemma holds under \eqref{eq:c2v4-uniform} alone.
\end{remark}

% =============================================================================
\section{The defect equation and its Gaussian energy identity}
\label{sec:c2v7-defect-equation}
% =============================================================================

Let \(Y:=\partial_\sigma V\) and, for a.e.\ \(\sigma\), let
\(P':=\partial_\sigma P_V\), the \(\sigma\)-derivative of the profile
pressure; explicitly, with \(s=-e^{-\sigma}\),
\begin{equation}
 P'=-(-s)P+(-s)^2\partial_sP-\tfrac12(-s)^{3/2}\,z\cdot\nabla P ,
 \qquad
 Y=(-s)^{3/2}\Bigl[\partial_sW+\frac{W+y\cdot\nabla W}{2(-s)}\Bigr],
 \label{eq:c2v7-Pprime}
\end{equation}
all evaluated at \(y=(-s)^{1/2}z\).

\begin{proposition}[The defect equation]
\label{prop:c2v7-defect-equation}
For \(W\) as in Lemma~\ref{lem:c2v7-higher}, \(Y\) is smooth,
divergence-free, and satisfies for a.e.\ \(\sigma\)
\begin{equation}
 \partial_\sigma Y=L_-Y-N'(V)Y,
 \qquad
 N'(V)Y:=(Y\cdot\nabla)V+(V\cdot\nabla)Y+\nabla P' ,
 \label{eq:c2v7-defect-equation}
\end{equation}
with \(Y,\nabla Y,P',\nabla P'\in L^2(G)\) for a.e.\ \(\sigma\), and
\(\sup_\sigma\|Y(\sigma)\|_G\le C_Y\), where \(C_Y\) depends on
\(C_\infty,C_\infty',C_\infty'',C,\nu\) and the Calder\'on--Zygmund constant.
\end{proposition}

\begin{proof}
Differentiate \eqref{eq:c2v2-profile-equation} in \(\sigma\); the
derivative of \(N(V)\) is \(N'(V)Y\) with \(\partial_\sigma\nabla P_V=\nabla\partial_\sigma P_V\)
by smoothness.  The formula \eqref{eq:c2v7-Pprime} is the chain rule on
\(P_V(z,\sigma)=(-s)P(s,(-s)^{1/2}z)\) with \(\dd s/\dd\sigma=-s\).
Integrability: \(Y=L_-V-N(V)\) with \(\|\Delta V\|_G\le C_\infty''\|1\|_G\),
\(\|z\cdot\nabla V\|_G\le C_\infty'\||z|\|_G\), \(\|V\|_G\le C_\infty\|1\|_G\),
and \(\|N(V)\|_G\le C_N^{1/2}\), which gives \(C_Y\); \(\nabla Y\in L^2(G)\)
for a.e.\ \(\sigma\) by Lemma~\ref{lem:c2v7-higher} transported to the
profile (\(\nabla Y\) is a combination of \(\nabla\partial_sW\), \(\nabla W\),
\(\nabla^2W\), \(y\cdot\nabla^2W\), each in \(L^2(\psi)\) for a.e.\ \(s\),
the factor \(|y|\) being absorbed by the Gaussian against the bounded
\(\nabla^2W\)); \(P'\in L^2(G)\) since \(P\in L^3\) gives \(\int|P|^2G\le\|P\|_3^2\|G\|_{3}\),
\(\partial_sP\in L^2(\R^3)\) for a.e.\ \(s\) by the lemma, and
\(\int|z|^2|\nabla P|^2G\le\sup(|z|^2G)\|\nabla P\|_2^2\) after the change of
variables; \(\nabla P'\in L^2(G)\) similarly with \(\nabla^2P\in L^2\) and
\(\nabla\partial_sP\in L^2\) for a.e.\ \(s\) (the latter from
\(\nabla\partial_s(W\otimes W)\in L^2\), which uses \(\nabla\partial_sW\in L^2\)).
\end{proof}

\begin{theorem}[Gaussian energy identity for the defect]
\label{thm:c2v7-defect-energy}
For \(W\) as above, \(\sigma\mapsto\|Y(\sigma)\|_G^2\) is absolutely
continuous and for a.e.\ \(\sigma\)
\begin{equation}
 \boxed{\quad
 \frac12\frac{\dd}{\dd\sigma}\|Y\|_G^2
 =-\nu\|\nabla Y\|_G^2-\frac12\|Y\|_G^2
 -\int Y\cdot(Y\cdot\nabla)V\,G
 -\frac1{4\nu}\int|Y|^2\,(z\cdot V)\,G
 -\frac1{2\nu}\int P'\,(z\cdot Y)\,G .
 \quad}
 \label{eq:c2v7-defect-energy}
\end{equation}
The three nonlinear terms are the Gaussian inner product
\(-\langle Y,N'(V)Y\rangle_G\); the first of them is bounded by
\(C_\infty'\|Y\|_G^2\), the second by
\(\frac{C_\infty}{4\nu}\|Y\|_G\||z|Y\|_G\le\frac{C_\infty}{4\nu}\|Y\|_G\,C(\nu)(\|Y\|_G+\|\nabla Y\|_G)\).
\end{theorem}

\begin{proof}
Take the \(G\)-inner product of \eqref{eq:c2v7-defect-equation} with
\(Y\).  \(\langle Y,L_-Y\rangle_G=-\nu\|\nabla Y\|_G^2-\frac12\|Y\|_G^2\) by the
symmetry of Lemma~\ref{lem:c2v2-dictionary} (justified with a cutoff as
in Theorem~\ref{thm:c2v2-second}, using \(Y,\nabla Y\in L^2(G)\) and
\(\Delta Y\in L^2(G)\) for a.e.\ \(\sigma\) by the lemma).
\(\langle Y,(V\cdot\nabla)Y\rangle_G=\frac12\int V\cdot\nabla|Y|^2G=-\frac12\int|Y|^2V\cdot\nabla G=\frac1{4\nu}\int|Y|^2(z\cdot V)G\)
using \(\operatorname{div}V=0\); and
\(\langle Y,\nabla P'\rangle_G=-\int P'\operatorname{div}(YG)=\frac1{2\nu}\int P'(z\cdot Y)G\)
using \(\operatorname{div}Y=0\).  The absolute continuity follows from the
integrability of the right side on compact \(\sigma\)-intervals.  The
Gaussian Hardy inequality \(\||z|Y\|_G\le C(\nu)(\|Y\|_G+\|\nabla Y\|_G)\)
follows from \(\frac1{2\nu}\int|z|^2|Y|^2G=\int(3|Y|^2+2Y\cdot(z\cdot\nabla)Y)G\),
obtained by integrating \(|Y|^2z\cdot\nabla G\) by parts, and
Cauchy--Schwarz.
\end{proof}

\begin{corollary}[Alternative (B) forces anti-dissipation of the linearized nonlinearity]
\label{cor:c2v7-anti-dissipation}
Let \(x_0\) be homogeneous and suppose alternative (B) of
Corollary~\ref{cor:c2v6-dichotomy-rate} holds.  Then for every
\(W\in\mathcal F(x_0)\), every \(\sigma\), and every \(L>0\),
\begin{equation}
 \boxed{\quad
 \frac1L\int_\sigma^{\sigma+L}\langle Y,N'(V)Y\rangle_G\,\dd\sigma
 \ \le\ -\frac1L\int_\sigma^{\sigma+L}\Bigl(\nu\|\nabla Y\|_G^2+\frac12\|Y\|_G^2\Bigr)\dd\sigma+\frac{C_Y^2}{2L}
 \ \le\ -\frac{m(L)}{2L}+\frac{C_Y^2}{2L},
 \quad}
 \label{eq:c2v7-anti-dissipation}
\end{equation}
so the log-average of \(\langle Y,N'(V)Y\rangle_G\) is at most
\(-\frac12\underline\delta+o(1)\) as \(L\to\infty\), uniformly on the
family.  In particular the linearized nonlinearity acting on the defect
cannot be dissipative, nor even neutral, on average: it must feed the
defect at least at the rate at which the Ornstein--Uhlenbeck part
damps it.
\end{corollary}

\begin{proof}
Integrate \eqref{eq:c2v7-defect-energy} over \([\sigma,\sigma+L]\):
\(\int\langle Y,N'(V)Y\rangle_G=-\int(\nu\|\nabla Y\|^2+\frac12\|Y\|^2)+\frac12(\|Y(\sigma)\|^2-\|Y(\sigma+L)\|^2)\),
and use \(\|Y\|_G\le C_Y\) and \(\int_\sigma^{\sigma+L}\|Y\|_G^2\ge m(L)\).
\end{proof}

\begin{assessment}[Two signs for the open case]
\label{ass:c2v7-two-signs}
In the open alternative (B), two sign statements now hold on every
element of the family at every homogeneous point, on log-average with
explicit rates: the nonlinearity is anti-aligned with the defect
(\(\langle N(V),\partial_\sigma V\rangle_G\) at most \(-\underline\delta+o(1)\)
on average, Corollary~\ref{cor:c2v6-dichotomy-rate}), and the linearized
nonlinearity is anti-dissipative on the defect
(\(\langle Y,N'(V)Y\rangle_G\) at most \(-\frac12\underline\delta+o(1)\) on
average).  Both say that a persistent departure from self-similarity
must be maintained by the nonlinear terms against the Gaussian
damping; neither is contradicted by any bound available to the series,
since the nonlinear terms are of order one in the Gaussian norm
(Lemma~\ref{lem:c2v3-N-bound}, Proposition~\ref{prop:c2v7-defect-equation}).
A contradiction would require a structural reason why the Navier--Stokes
nonlinearity cannot feed its own self-similarity defect at a fixed rate
forever in log-scale; the series does not have one.  No global
regularity claim is made.
\end{assessment}

% =============================================================================
\section{Integrated V7 conclusion and next acquisition order}
\label{sec:c2v7-conclusion}
% =============================================================================

\begin{theorem}[What V7 proves]
\label{thm:c2v7-conclusion}
Relative to the first cycle and V1--V6, modulo (E1)--(E8): the vorticity
form \eqref{eq:c2v7-vorticity-form}, the integrability lemma
(Lemma~\ref{lem:c2v7-higher}), the defect equation
\eqref{eq:c2v7-defect-equation} with the bound \(\|Y\|_G\le C_Y\), the
Gaussian energy identity \eqref{eq:c2v7-defect-energy}, and the
anti-dissipation \eqref{eq:c2v7-anti-dissipation} forced by alternative
(B).  The full Navier--Stokes stability/global-regularity theorem
remains unproved.
\end{theorem}

\begin{proof}
The cited results.
\end{proof}

\begin{programme}[V8 acquisition order]
\label{prog:c2v8}
\begin{enumerate}[label=\textup{(AC\arabic*)},leftmargin=4.8em]
\item Combine the two signs.  The cross term of the second identity and
      the linearized cross term are related through
      \(\frac{\dd}{\dd\sigma}\langle N(V),Y\rangle_G=\langle N'(V)Y,Y\rangle_G+\langle N(V),\partial_\sigma Y\rangle_G\);
      determine whether integrating this relation over long windows,
      with the two averaged signs and the bounds on \(N\), \(Y\),
      \(\partial_\sigma Y\), yields a constraint stronger than either sign.
\item Spectral gap.  Record the spectral decomposition of \(L_-\) on
      divergence-free \(L^2(G)\) fields (Hermite expansion, eigenvalues
      \(-\frac12-\frac k2\)) and the resulting exponential decay of the
      linear defect flow; determine the size of the nonlinear feeding
      needed to sustain a defect of Gaussian norm \(\ge(m(1))^{1/2}\) per
      unit window.
\item The next compulsory companion audit is V10.
\end{enumerate}
\end{programme}

\begin{assessment}[Position after V7]
\label{ass:c2v7-position}
The open alternative now carries two averaged sign statements with
explicit rates.  No full stability or global-regularity claim is made.
\end{assessment}

% =============================================================================
\section*{Version V7 (second cycle) append-only record}
\addcontentsline{toc}{section}{Version V7 (second cycle) append-only record}
% =============================================================================

The complete V6 source of the second cycle was retained before this
continuation.  Its SHA--256 digest was
\begin{center}\scriptsize\ttfamily
d377cdd32a679c8e4db261ee5f38907a6ca9f13fe774ec0a5786ef792395f8b5.
\end{center}
V7 derived the defect equation and its Gaussian energy identity and
proved that alternative (B) forces the linearized nonlinearity to be
anti-dissipative on the defect.  No full regularity claim is made.

% =============================================================================
\section{V8: the spectral gap of the profile operator and the three averaged
signs of alternative (B)}
\label{sec:c2v8-entry}
% =============================================================================

Items (AC1) and (AC2) of Programme~\ref{prog:c2v8} are carried out.  The
linear part of the profile equation is the Ornstein--Uhlenbeck operator
shifted by \(-\frac12\); its spectrum on \(L^2(G)\) is
\(\{-\frac12-\frac k2:k\ge0\}\), the divergence-free subspace is
invariant, and the Gaussian Poincar\'e inequality gives the sharpened
damping \(\nu\|\nabla Y\|_G^2+\frac12\|Y\|_G^2\ge\|Y-\bar Y\|_G^2+\frac12\|\bar Y\|_G^2\),
where \(\bar Y\) is the Gaussian mean of the defect.  The relation
between the two cross terms is integrated over long windows and yields
a third averaged sign in alternative (B): the nonlinearity is positively
correlated with the log-scale acceleration of the defect.  It is
recorded that the third sign is a consequence of the first two and the
bounds, not a strengthening, and that all three are consistent with the
order-one size of the nonlinear terms.  No global regularity claim is
made.  The next compulsory companion audit is V10.

% =============================================================================
\section{Spectral facts}
\label{sec:c2v8-spectral}
% =============================================================================

\begin{proposition}[Spectrum of the profile operator and Gaussian Poincar\'e]
\label{prop:c2v8-spectrum}
On \(L^2(G)\), \(G=e^{-|z|^2/(4\nu)}\):
\begin{enumerate}[label=\textup{(\roman*)},leftmargin=3.2em]
\item The operator \(A:=\nu\Delta-\frac12z\cdot\nabla\) is, after the
      change of variables \(z=(2\nu)^{1/2}x\), one half of the standard
      Ornstein--Uhlenbeck operator \(\Delta_x-x\cdot\nabla_x\) on
      \(L^2(e^{-|x|^2/2})\); it is self-adjoint, nonpositive, with
      eigenvalues \(-\frac k2\), \(k=0,1,2,\dots\), and eigenfunctions the
      Hermite polynomials of degree \(k\) in \(x\).  Hence
      \(L_-=A-\frac12\) has spectrum \(\{-\frac12-\frac k2\}\).
\item \(L_-\) maps divergence-free fields to divergence-free fields, so
      the divergence-free subspace of \(L^2(G)^3\) is invariant, and on it
      \(\langle Y,L_-Y\rangle_G\le-\frac12\|Y\|_G^2\).
\item (Gaussian Poincar\'e.)  For \(Y\in H^1(G)\) with Gaussian mean
      \(\bar Y:=\int YG/\int G\),
      \(\nu\|\nabla Y\|_G^2\ge\frac12\|Y-\bar Y\|_G^2\); consequently
      \begin{equation}
       \nu\|\nabla Y\|_G^2+\frac12\|Y\|_G^2\ \ge\ \|Y-\bar Y\|_G^2+\frac12\|\bar Y\|_G^2 .
       \label{eq:c2v8-Poincare}
      \end{equation}
\item The linear defect flow \(\partial_\sigma Y=L_-Y\) satisfies
      \(\|Y(\sigma)\|_G\le e^{-(\sigma-\sigma_0)/2}\|Y(\sigma_0)\|_G\), and
      \(\|Y(\sigma)-\bar Y(\sigma)\|_G\le e^{-(\sigma-\sigma_0)}\|Y(\sigma_0)-\bar Y(\sigma_0)\|_G\).
\end{enumerate}
\end{proposition}

\begin{proof}
(i) With \(z=(2\nu)^{1/2}x\), \(\nu\Delta_z=\frac12\Delta_x\) and
\(\frac12z\cdot\nabla_z=\frac12x\cdot\nabla_x\), so \(A=\frac12(\Delta_x-x\cdot\nabla_x)\),
and \(G=e^{-|x|^2/2}\); the spectral facts of the Ornstein--Uhlenbeck
operator are classical.  (ii) \(\operatorname{div}(z\cdot\nabla V)=z\cdot\nabla\operatorname{div}V+\operatorname{div}V\),
so \(\operatorname{div}L_-V=(\nu\Delta-\frac12z\cdot\nabla-1)\operatorname{div}V=0\)
if \(\operatorname{div}V=0\).  (iii) The Poincar\'e inequality for the
Gaussian measure \(e^{-|x|^2/2}\) with constant \(1\),
\(\int|\nabla_xf|^2\ge\int|f-\bar f|^2\), transported: \(\nu\|\nabla_zY\|_G^2=\frac12\|\nabla_xY\|^2\ge\frac12\|Y-\bar Y\|^2\);
then \(\frac12\|Y\|^2=\frac12\|Y-\bar Y\|^2+\frac12\|\bar Y\|^2\).  (iv) From
(ii) and (iii) by Gr\"onwall.
\end{proof}

\begin{corollary}[Sharpened anti-dissipation]
\label{cor:c2v8-sharpened}
In alternative (B), for every \(W\in\mathcal F(x_0)\), \(\sigma\), \(L\),
\begin{equation}
 \frac1L\int_\sigma^{\sigma+L}\langle Y,N'(V)Y\rangle_G
 \ \le\ -\frac1L\int_\sigma^{\sigma+L}\Bigl(\|Y-\bar Y\|_G^2+\frac12\|\bar Y\|_G^2\Bigr)+\frac{C_Y^2}{2L}.
 \label{eq:c2v8-sharpened}
\end{equation}
\end{corollary}

\begin{proof}
\eqref{eq:c2v7-defect-energy} integrated, with \eqref{eq:c2v8-Poincare}.
\end{proof}

\begin{remark}[The constant mode of the defect]
\label{rem:c2v8-constant-mode}
A defect equal to a constant vector \(e\) in \(z\) is
\(Z=(-s)^{-3/2}e\) in the original variables, that is,
\(\partial_sW+\frac{W+y\cdot\nabla W}{2(-s)}=(-s)^{-3/2}e\), which is the
self-similarity defect of a Galilean-shifted self-similar solution
\(W(s,y)=(-s)^{-1/2}V_0\bigl((y-\gamma(s))/(-s)^{1/2}\bigr)\) with
\(\gamma\) linear in \((-s)^{1/2}\); the constant mode is the least damped
(rate \(\frac12\)) and corresponds to the freedom of the blow-up centre.
The non-constant modes are damped at rate at least \(1\).
\end{remark}

% =============================================================================
\section{The three averaged signs}
\label{sec:c2v8-three-signs}
% =============================================================================

\begin{theorem}[Relation between the cross terms and the third sign]
\label{thm:c2v8-third-sign}
For \(W\) as in Proposition~\ref{prop:c2v7-defect-equation}, the function
\(\sigma\mapsto\langle N(V),Y\rangle_G\) is absolutely continuous with
\begin{equation}
 \frac{\dd}{\dd\sigma}\langle N(V),Y\rangle_G=\langle N'(V)Y,Y\rangle_G+\langle N(V),\partial_\sigma Y\rangle_G
 \label{eq:c2v8-relation}
\end{equation}
for a.e.\ \(\sigma\), and \(|\langle N(V),Y\rangle_G|\le C_N^{1/2}C_Y\).  In
alternative (B), for every \(W\in\mathcal F(x_0)\), \(\sigma\), \(L\),
\begin{equation}
 \boxed{\quad
 \frac1L\int_\sigma^{\sigma+L}\langle N(V),\partial_\sigma Y\rangle_G\,\dd\sigma
 \ \ge\ \frac{m(L)}{2L}-\frac{C_Y^2+4C_N^{1/2}C_Y}{2L},
 \quad}
 \label{eq:c2v8-third}
\end{equation}
so the log-average of \(\langle N(V),\partial_\sigma Y\rangle_G\) is at least
\(\frac12\underline\delta-o(1)\): the nonlinearity is positively
correlated with the log-scale acceleration of the defect.
\end{theorem}

\begin{proof}
\(N(V)\) is differentiable in \(\sigma\) with derivative \(N'(V)Y\), and
\(Y\) with derivative \(\partial_\sigma Y\in L^2(G)\) for a.e.\ \(\sigma\)
(Proposition~\ref{prop:c2v7-defect-equation}); the product rule holds in
the absolutely continuous sense.  Integrate \eqref{eq:c2v8-relation}
over \([\sigma,\sigma+L]\):
\(\int\langle N,\partial_\sigma Y\rangle=\langle N,Y\rangle\big|_\sigma^{\sigma+L}-\int\langle N'Y,Y\rangle
\ge-2C_N^{1/2}C_Y+\frac{m(L)}2-\frac{C_Y^2}2\) by
\eqref{eq:c2v7-anti-dissipation}.
\end{proof}

\begin{firewall}[Item (AC1): the third sign is not a strengthening]
\label{fw:c2v8-AC1}
\eqref{eq:c2v8-third} is derived from \eqref{eq:c2v7-anti-dissipation}
and the bound on \(\langle N,Y\rangle_G\); it contains no information
beyond them.  Substituting \(\partial_\sigma Y=L_-Y-N'(V)Y\) gives
\(\langle N,\partial_\sigma Y\rangle=\langle N,L_-Y\rangle-\langle N,N'(V)Y\rangle\),
whose two parts are each of order \(C_N^{1/2}C_Y\) and have no sign; the
relation \eqref{eq:c2v8-relation} is an identity between three order-one
quantities and does not produce a contradiction with alternative (B).
Item (AC1) is closed with this record.
\end{firewall}

\begin{theorem}[The three averaged signs of alternative (B)]
\label{thm:c2v8-three-signs}
Let \(x_0\) be a homogeneous singular point in alternative (B), with
uniform defect rate \(\underline\delta>0\).  For every \(W\in\mathcal F(x_0)\)
and every \(\sigma\), as \(L\to\infty\) uniformly in \(W\) and \(\sigma\):
\begin{align}
 \frac1L\int_\sigma^{\sigma+L}\langle N(V),\partial_\sigma V\rangle_G
 &\le-\underline\delta+o(1),
 \label{eq:c2v8-sign1}\\
 \frac1L\int_\sigma^{\sigma+L}\langle N'(V)\partial_\sigma V,\partial_\sigma V\rangle_G
 &\le-\frac1L\int_\sigma^{\sigma+L}\Bigl(\|Y-\bar Y\|_G^2+\frac12\|\bar Y\|_G^2\Bigr)+o(1)
 \le-\frac12\underline\delta+o(1),
 \label{eq:c2v8-sign2}\\
 \frac1L\int_\sigma^{\sigma+L}\langle N(V),\partial_\sigma^2V\rangle_G
 &\ge\frac12\underline\delta-o(1).
 \label{eq:c2v8-sign3}
\end{align}
Each of the three quantities is bounded in absolute value by a constant
of order \(C_N^{1/2}C_Y\) or \(C_NC_Y^2\), so none of the signs is
contradicted by the available bounds.
\end{theorem}

\begin{proof}
\eqref{eq:c2v6-B}, \eqref{eq:c2v8-sharpened}, \eqref{eq:c2v8-third},
with \(m(L)/L\to\underline\delta\).
\end{proof}

\begin{firewall}[Item (AC2): the feeding required per unit window]
\label{fw:c2v8-AC2}
By \eqref{eq:c2v8-sharpened} with \(L=1\), sustaining
\(\int_\sigma^{\sigma+1}\|Y\|_G^2\ge m(1)\) requires
\(\int_\sigma^{\sigma+1}\langle Y,N'(V)Y\rangle_G\le-\frac12m(1)+\frac12C_Y^2\),
a negative bound only if \(m(1)>C_Y^2\), which is not known; the
definite statements are the long-window ones.  The size of the
linearized nonlinearity is \(\|N'(V)Y\|_G\le(C_\infty'+\ldots)\|Y\|_G+\ldots\),
of the same order as \(\|L_-Y\|_G\), so the feeding required by
alternative (B), of order \(\frac12\underline\delta\le\frac12C_N\), is
within the available size.  Item (AC2) is closed with this record.
\end{firewall}

% =============================================================================
\section{Integrated V8 conclusion and next acquisition order}
\label{sec:c2v8-conclusion}
% =============================================================================

\begin{theorem}[What V8 proves]
\label{thm:c2v8-conclusion}
Relative to the first cycle and V1--V7, modulo (E1)--(E8): the spectral
facts of Proposition~\ref{prop:c2v8-spectrum}, the sharpened
anti-dissipation \eqref{eq:c2v8-sharpened}, the relation
\eqref{eq:c2v8-relation} with the third sign \eqref{eq:c2v8-third}, and
the three-sign theorem.  The full Navier--Stokes
stability/global-regularity theorem remains unproved.
\end{theorem}

\begin{proof}
The cited results.
\end{proof}

\begin{programme}[V9 acquisition order]
\label{prog:c2v9}
\begin{enumerate}[label=\textup{(AC\arabic*)},leftmargin=4.8em]
\item Alternative (A).  Decide the import of the exclusion of nonzero
      backward self-similar solutions with the properties of the
      ancient limit (bounded, \(L^6\) profile, finite Gaussian Dirichlet
      integral, local energy inequality) as an external theorem (E9),
      state it precisely with its hypotheses, verify that the
      self-similar element produced by alternative (A) satisfies them,
      and conclude that every homogeneous singular point is in
      alternative (B).
\item Record what (E9) does not give: it excludes exact self-similarity
      only; discretely self-similar and general elements remain.
\item The next compulsory companion audit is V10.
\end{enumerate}
\end{programme}

\begin{assessment}[Position after V8]
\label{ass:c2v8-position}
Alternative (B) carries three averaged signs with explicit rates, none
contradicted by the available bounds.  No full stability or
global-regularity claim is made.
\end{assessment}

% =============================================================================
\section*{Version V8 (second cycle) append-only record}
\addcontentsline{toc}{section}{Version V8 (second cycle) append-only record}
% =============================================================================

The complete V7 source of the second cycle was retained before this
continuation.  Its SHA--256 digest was
\begin{center}\scriptsize\ttfamily
6c6d23bd746aac9d650a9ca8f0f8ee2eccd3d503306b8c7e98e3b060b7f1e520.
\end{center}
V8 recorded the spectral facts of the profile operator, sharpened the
anti-dissipation, and derived the third averaged sign of alternative
(B).  No full regularity claim is made.

% =============================================================================
\section{V9: the import of the self-similar Liouville theorem and the
unconditional alternative (B)}
\label{sec:c2v9-entry}
% =============================================================================

Items (AC1) and (AC2) of Programme~\ref{prog:c2v9} are carried out.  The
exclusion of nonzero backward self-similar solutions of the
Navier--Stokes equations whose profile lies in \(L^q(\R^3)\) for some
\(q\in(3,\infty]\) is imported as the ninth external theorem (E9), with
its hypotheses stated precisely.  The self-similar element produced by
alternative (A) of the dichotomy is verified to satisfy those
hypotheses; hence alternative (A) is impossible and every homogeneous
singular point of a type-I solution is in alternative (B): the solution
is uniformly non-self-similar near the point at every small scale, and
the three averaged signs hold unconditionally on the rescaling family.
What (E9) does not give is recorded.  No global regularity claim is
made.  The next version is the compulsory companion audit.

% =============================================================================
\section{The ninth import}
\label{sec:c2v9-import}
% =============================================================================

\begin{firewall}[Self-similar Liouville theorem]
\label{fw:c2v9-import}
The following is used as a black box and is not re-proved.
\begin{enumerate}[label=\textup{(E\arabic*)},leftmargin=3.6em,start=9]
\item \emph{Leray profiles in \(L^q\) vanish.}  Let \(V\) be a smooth
      divergence-free vector field on \(\R^3\) with \(V\in L^q(\R^3)\) for
      some \(q\in(3,\infty]\), with pressure
      \(P_V=\partial_i\partial_j(-\Delta)^{-1}(V_iV_j)\), solving Leray's
      equation
      \begin{equation}
       -\nu\Delta V+\tfrac12V+\tfrac12z\cdot\nabla V+(V\cdot\nabla)V+\nabla P_V=0
       \qquad\text{on }\R^3 .
       \label{eq:c2v9-Leray}
      \end{equation}
      Then \(V\equiv0\).  (Ne\v{c}as--R\r{u}\v{z}i\v{c}ka--\v{S}ver\'{a}k
      for \(q=3\); Tsai for \(q\in(3,\infty]\) and for profiles with locally
      bounded energy whose self-similar solution satisfies the local
      energy inequality.  The series uses the \(L^q\) form with \(q=6\).)
\end{enumerate}
The theorem is external in the same sense as (E1)--(E8): a published
result whose proof is not reproduced here, used exactly under the
hypotheses stated.  The import ledger of the first cycle is extended
accordingly, and every statement of V9 onward that uses (E9) says so.
\end{firewall}

\begin{remark}[Why the profile satisfies the hypotheses]
\label{rem:c2v9-hypotheses}
For a backward self-similar element \(W(s,y)=(-s)^{-1/2}V_0(y/(-s)^{1/2})\)
of the class with the uniform bounds \eqref{eq:c2v4-uniform},
Lemma~\ref{lem:c2v2-dictionary} gives \(L_-V_0=N(V_0)\), which is
\eqref{eq:c2v9-Leray}; \(V_0=V(\cdot,\sigma)\) for any \(\sigma\) is smooth,
divergence-free, in \(L^6(\R^3)\) (\(\|V_0\|_6=(-s)^{1/2}(-s)^{-1/4}\cdot\)the
scaled \(L^6\) norm of \(W(s)\), finite by the class bound), bounded by
\(C_\infty\), and its pressure is the Riesz representation transported
from \(P\).  Hence (E9) applies with \(q=6\).
\end{remark}

% =============================================================================
\section{Alternative (A) is impossible}
\label{sec:c2v9-A-impossible}
% =============================================================================

\begin{theorem}[Every homogeneous singular point is in alternative (B)]
\label{thm:c2v9-B}
Let \(u\) satisfy the type-I rate and let \(x_0\) be a homogeneous
singular point.  Modulo (E1)--(E9):
\begin{enumerate}[label=\textup{(\roman*)},leftmargin=3.2em]
\item \(\mathcal F(x_0)\) contains no nonzero backward self-similar element;
      \(m(\ell)>0\) for every \(\ell>0\), and \(\underline\delta>0\).
\item For every \(\ell>0\) there is \(\lambda(\ell)>0\) such that the
      solution's windowed defect satisfies
      \(\mathcal Z_u(\lambda;\ell)\ge\frac12m(\ell)>0\) for all
      \(0<\lambda\le\lambda(\ell)\).
\item For every \(W\in\mathcal F(x_0)\), every \(\sigma\), and every \(L\),
      the bounds \eqref{eq:c2v6-B}, \eqref{eq:c2v8-sharpened},
      \eqref{eq:c2v8-third} hold, and as \(L\to\infty\) the three
      averaged signs \eqref{eq:c2v8-sign1}--\eqref{eq:c2v8-sign3} hold
      uniformly on \(\mathcal F(x_0)\).
\end{enumerate}
\end{theorem}

\begin{proof}
(i) If \(\mathcal F(x_0)\) contained a nonzero backward self-similar
element, its profile would satisfy the hypotheses of (E9) by
Remark~\ref{rem:c2v9-hypotheses} and would vanish, contradicting the
pinching \(\Phi_W\ge c_\Phi>0\).  Hence alternative (A) of
Corollary~\ref{cor:c2v6-dichotomy-rate} fails and (B) holds:
\(m(\ell)>0\) for all \(\ell\) by Proposition~\ref{prop:c2v6-m}(iii), and
\(\underline\delta>0\) by (iv).  (ii) Theorem~\ref{thm:c2v6-solution}.
(iii) Corollary~\ref{cor:c2v6-dichotomy-rate},
Corollary~\ref{cor:c2v8-sharpened}, Theorems~\ref{thm:c2v8-third-sign}
and \ref{thm:c2v8-three-signs}.
\end{proof}

\begin{assessment}[The unconditional picture at a homogeneous point]
\label{ass:c2v9-picture}
Modulo the nine imports, a homogeneous singular point of a type-I
solution has the following structure.  Every ancient limit at the point
and every rescaling limit of one is uniformly concentrated, pinched in
Gaussian energy, and uniformly non-self-similar: its self-similarity
defect has Gaussian log-integral at least \(m(\ell)\) on every window of
log-length \(\ell\), with rate \(\underline\delta>0\) per unit log-scale.
The solution itself has the same property at all small scales.  On
every such limit the nonlinearity is anti-aligned with the defect, the
linearized nonlinearity is anti-dissipative on the defect, and the
nonlinearity is positively correlated with the defect's log-scale
acceleration, all on log-average with rates \(\underline\delta\),
\(\frac12\underline\delta\), \(\frac12\underline\delta\).  The remaining
task for homogeneous points is to show that these three signed
statements cannot all hold on a compact scale-invariant family of
ancient solutions with the class properties; the series has no
mechanism for this.  No global regularity claim is made.
\end{assessment}

\begin{firewall}[Item (AC2): what (E9) does not give]
\label{fw:c2v9-AC2}
(E9) excludes exact backward self-similarity only.  Discretely
self-similar elements, \(V(\cdot,\sigma+\ell_0)=V(\cdot,\sigma)\) for some
\(\ell_0>0\) and all \(\sigma\), are not excluded by (E9) and are
consistent with alternative (B): their defect is periodic in \(\sigma\)
with positive log-integral per period.  The exclusion of discretely
self-similar solutions with profiles in the present class is not
imported and is not claimed; the series records it as open.  More
generally, (E9) says nothing about elements whose defect is bounded
below, which is the whole of alternative (B).
\end{firewall}

% =============================================================================
\section{Integrated V9 conclusion and next acquisition order}
\label{sec:c2v9-conclusion}
% =============================================================================

\begin{theorem}[What V9 establishes]
\label{thm:c2v9-conclusion}
Relative to the first cycle and V1--V8, modulo (E1)--(E9): the import
(E9) and Theorem~\ref{thm:c2v9-B}, which places every homogeneous
singular point in alternative (B) with all its consequences.  The full
Navier--Stokes stability/global-regularity theorem remains unproved.
\end{theorem}

\begin{proof}
Theorem~\ref{thm:c2v9-B}.
\end{proof}

\begin{programme}[V10 acquisition order, with compulsory companion audit]
\label{prog:c2v10}
\begin{enumerate}[label=\textup{(AC\arabic*)},leftmargin=4.8em]
\item Perform the compulsory review of the current SM and YM notes;
      record digests; import nothing numerical.
\item Consolidate V6--V9 into one statement on the uniformly
      non-self-similar structure of homogeneous singular points, with
      the extended import ledger (E1)--(E9).
\item Discrete self-similarity.  Define the discretely self-similar
      defect \(V(\cdot,\sigma+\ell_0)-V(\cdot,\sigma)\) and determine whether
      the dichotomy argument of V4--V6 has a discrete analogue: either
      the family contains a discretely self-similar element of period
      \(\ell_0\), or the discrete defect is bounded below on the family.
\item The next compulsory companion audit after V10 is V15.
\end{enumerate}
\end{programme}

\begin{assessment}[Position after V9]
\label{ass:c2v9-position}
Alternative (B) is unconditional at every homogeneous singular point,
modulo the nine imports.  No full stability or global-regularity claim
is made.
\end{assessment}

% =============================================================================
\section*{Version V9 (second cycle) append-only record}
\addcontentsline{toc}{section}{Version V9 (second cycle) append-only record}
% =============================================================================

The complete V8 source of the second cycle was retained before this
continuation.  Its SHA--256 digest was
\begin{center}\scriptsize\ttfamily
8baf0a6e0fe1f43c0a0d2dcc42cc2bded365300a425299df35604953f0ad14aa.
\end{center}
V9 imported the self-similar Liouville theorem as (E9) and concluded
that every homogeneous singular point is in alternative (B).  No full
regularity claim is made.

% =============================================================================
\section{V10: compulsory companion audit, the consolidated non-self-similar
structure, and the discrete dichotomy}
\label{sec:c2v10-entry}
% =============================================================================

This is the tenth version of the second cycle.  The compulsory review of
the two companion programmes is performed first.  V6--V9 are then
consolidated into one statement on the uniformly non-self-similar
structure of homogeneous singular points, with the import ledger
extended to (E9).  The dichotomy of V4--V6 is then given its discrete
analogue: for every log-period, either the rescaling family contains a
discretely self-similar element of that period, or the discrete defect
is bounded below on the family and on the solution at small scales;
the discrete defect is shown to be continuous on the family, and the
exclusion of discretely self-similar elements is recorded as neither
imported nor claimed.  No global regularity claim is made.  The next
compulsory companion audit is V15.

% =============================================================================
\section{Compulsory V10 review of the current companion notes}
\label{sec:c2v10-companion-audit}
% =============================================================================

The current companion files in \texttt{latex\_notes} were inspected
read-only.  The frozen checkpoints are
\begin{align}
 &\texttt{Renewal\_SM\_Einstein\_Continuum\_Research\_Notes\_V86.tex},
 \notag\\[-2pt]
 &\qquad
 \texttt{ccdd7f62fddb2fec2d9c9d89026a37019b63bae8b13169c04b50734cf0520f5a},
 \label{eq:c2v10-SM-checkpoint}\\
 &\texttt{Renewal\_Yang\_Mills\_Mass\_Gap\_Research\_Notes\_V60.tex},
 \notag\\[-2pt]
 &\qquad
 \texttt{14ec8cc3515d32c3e2a47d8e8236b5fb922b516ca482353b585c83a32d074500}.
 \label{eq:c2v10-YM-checkpoint}
\end{align}
The YM file is byte-identical to the checkpoint \eqref{eq:c2v5-YM-checkpoint}
of V5; its programme has not advanced since.

\emph{SM V86.}  The SM programme proved a conformal spectral gap: a
positive pointwise bound excludes conformal folds and controls the
inverse operator; it showed that the raw determinant is not controlled,
the relative inverse coarea factor containing an exact local divergence
proportional to the cutoff, that on the fixed-volume beat branch the
divergent coefficient remains nonzero on the first-order
beat-cancellation line, and that subtracting the trace converts the
factor to a positive renormalized determinant with a continuum limit;
its next gate is a common-regulator determinant ledger.

\begin{theorem}[V10 companion-audit decision]
\label{thm:c2v10-companion-decision}
No SM V86 spectral-gap, determinant, or renormalization constant is
imported, and nothing new is available from YM.  Neither bears on
Theorem~\ref{thm:c2v9-B} or on the discrete dichotomy below.
\end{theorem}

\begin{proof}
The SM results concern Fredholm determinants of constraint operators of
a gravitational model; none is a Navier--Stokes statement.
\end{proof}

% =============================================================================
\section{The non-self-similar structure, consolidated}
\label{sec:c2v10-consolidated}
% =============================================================================

\begin{theorem}[Import ledger of the second cycle]
\label{thm:c2v10-imports}
The second cycle uses the external theorems (E1)--(E8) of the first
cycle and, from V9 on, (E9) of Firewall~\ref{fw:c2v9-import}.  Within
the second cycle: (E5) in the uniform derivative bounds
\eqref{eq:c2v4-uniform}; (E6) and (E8) in the small-defect theorem
(Theorem~\ref{thm:c2v4-small-defect}) and in the discrete dichotomy
below; (E9) in Theorem~\ref{thm:c2v9-B}.  No other external result is
used beyond the classical inequalities listed in
Section~\ref{sec:c2v1-imports} and the spectral theory of the
Ornstein--Uhlenbeck operator.
\end{theorem}

\begin{proof}
Inspection of the proofs of V1--V9.
\end{proof}

\begin{theorem}[Uniformly non-self-similar structure of homogeneous singular points, V6--V9]
\label{thm:c2v10-consolidated}
Let \(u\) satisfy the type-I rate and let \(x_0\) be a homogeneous
singular point, \(\mathcal F(x_0)\) the full rescaling family,
\(m(\ell)\) the minimal windowed defect, \(\underline\delta\) the uniform
defect rate.  Modulo (E1)--(E9):
\begin{enumerate}[label=\textup{(\arabic*)},leftmargin=3.2em]
\item \(m(\ell)\) is attained, superadditive, positive for all \(\ell>0\),
      and \(m(\ell)/\ell\to\underline\delta\in(0,C_N]\).
\item Every \(W\in\mathcal F(x_0)\) has
      \(\int_\sigma^{\sigma+L}\|\partial_\sigma V\|_G^2\ge m(L)\) for all
      \(\sigma\), \(L\); the solution has
      \(\mathcal Z_u(\lambda;\ell)\ge\frac12m(\ell)\) for all
      \(\lambda\le\lambda(\ell)\).
\item Every \(W\in\mathcal F(x_0)\) satisfies the three averaged signs
      \eqref{eq:c2v8-sign1}--\eqref{eq:c2v8-sign3} with the rates
      \(\underline\delta,\frac12\underline\delta,\frac12\underline\delta\),
      and the defect obeys the equation \eqref{eq:c2v7-defect-equation}
      with the energy identity \eqref{eq:c2v7-defect-energy} and
      \(\|\partial_\sigma V\|_G\le C_Y\).
\item \(\mathcal F(x_0)\) contains no nonzero backward self-similar
      element.
\end{enumerate}
\end{theorem}

\begin{proof}
(1) Proposition~\ref{prop:c2v6-m} and Theorem~\ref{thm:c2v9-B}(i).
(2) Corollary~\ref{cor:c2v6-dichotomy-rate}, Theorem~\ref{thm:c2v6-solution},
Theorem~\ref{thm:c2v9-B}.  (3) Theorems~\ref{thm:c2v7-defect-energy},
\ref{thm:c2v8-three-signs}, Theorem~\ref{thm:c2v9-B}(iii).
(4) Theorem~\ref{thm:c2v9-B}(i).
\end{proof}

% =============================================================================
\section{The discrete dichotomy}
\label{sec:c2v10-discrete}
% =============================================================================

For \(\ell_0>0\) and \(W\in\mathcal F(x_0)\) with profile \(V\), define the
discrete defect of period \(\ell_0\) and its minimal unit-window value,
\begin{equation}
 \mathcal D_{\ell_0}(W;\sigma):=\|V(\cdot,\sigma+\ell_0)-V(\cdot,\sigma)\|_G^2,
 \qquad
 m_d(\ell_0):=\inf_{W\in\mathcal F(x_0)}\int_0^1\mathcal D_{\ell_0}(W;\sigma)\,\dd\sigma .
 \label{eq:c2v10-discrete-defect}
\end{equation}
\(W\) is discretely self-similar with factor \(\mu_0=e^{-\ell_0/2}\), that
is \(W^{\mu_0}=W\), if and only if \(\mathcal D_{\ell_0}(W;\sigma)=0\) for
all \(\sigma\).

\begin{proposition}[Continuity of the profile in the local topology]
\label{prop:c2v10-continuity}
Let \(W_k\to W\) locally in \(\mathcal F(x_0)\), or with \(W_k\) rescalings
of the original solution.  Then \(V_k(\cdot,\sigma)\to V(\cdot,\sigma)\)
strongly in \(L^2(G)\) for every \(\sigma\), uniformly on compact
\(\sigma\)-intervals.  Consequently \(W\mapsto\int_0^1\mathcal D_{\ell_0}(W;\sigma)\dd\sigma\)
is continuous on \(\mathcal F(x_0)\).
\end{proposition}

\begin{proof}
On a compact time interval \(J\), \(\partial_sW_k\) is bounded in
\(L^2(J;L^2(\psi_T))\) for any fixed weight scale \(T\) (proof of
Proposition~\ref{prop:c2v4-lsc}), so
\(\|W_k(s)-W_k(s')\|_{L^2(\psi_T)}\le|s-s'|^{1/2}\sup_k\|\partial_sW_k\|_{L^2(J;L^2(\psi_T))}\):
the maps \(s\mapsto W_k(s)\in L^2(\psi_T)\) are uniformly \(\frac12\)-H\"older.
They converge in \(L^1(J;L^2(\psi_T))\) by the strong \(L^2(J;L^2(B_R))\)
convergence and the uniformly small Gaussian tails (Morrey bound).  An
equicontinuous sequence converging in \(L^1(J)\) converges uniformly on
\(J\); hence \(W_k(s)\to W(s)\) in \(L^2(\psi_T)\) for every \(s\in J\),
uniformly.  Since \(\psi(\cdot,s)\) with the time-dependent scale
\(T=-s\) is comparable to \(\psi_{T_0}\) for \(s\) in a compact interval
away from zero, the same holds for the time-dependent weight, and the
change to the profile variable is a fixed bounded linear map on each
slice.  Continuity of the discrete defect follows.
\end{proof}

\begin{theorem}[Discrete dichotomy]
\label{thm:c2v10-discrete-dichotomy}
Let \(x_0\) be a homogeneous singular point and \(\ell_0>0\).  Modulo
(E1)--(E8), exactly one of the following holds.
\begin{enumerate}[label=\textup{(\Alph*\('\))},leftmargin=3.6em]
\item \(\mathcal F(x_0)\) contains a nonzero discretely self-similar
      element with factor \(e^{-\ell_0/2}\).
\item \(m_d(\ell_0)>0\), every \(W\in\mathcal F(x_0)\) has
      \(\int_\sigma^{\sigma+1}\mathcal D_{\ell_0}(W;\sigma')\dd\sigma'\ge m_d(\ell_0)\)
      for every \(\sigma\), and there is \(\lambda_d(\ell_0)>0\) such that
      the solution's discrete defect at scale \(\lambda\),
      \(\int_0^1\mathcal D_{\ell_0}(U_\lambda;\sigma)\dd\sigma\) for the
      rescaling \(U_\lambda=\lambda u(x_0+\lambda\,\cdot,T_*+\lambda^2\,\cdot)\),
      is at least \(\frac12m_d(\ell_0)\) for all \(0<\lambda\le\lambda_d(\ell_0)\).
\end{enumerate}
\end{theorem}

\begin{proof}
By Proposition~\ref{prop:c2v10-continuity} and the compactness of
\(\mathcal F(x_0)\), the infimum \(m_d(\ell_0)\) is attained by some
\(W_*\).  If \(m_d(\ell_0)=0\) then \(V_*(\cdot,\sigma+\ell_0)=V_*(\cdot,\sigma)\)
for a.e.\ \(\sigma\in[0,1]\), hence for all by continuity in \(\sigma\); in
the original variables, \(W_*^{\mu_0}=W_*\) on the time interval
corresponding to \(\sigma\in[0,1]\), \(\mu_0=e^{-\ell_0/2}\).  Both
\(W_*\) and \(W_*^{\mu_0}\) are elements of the class, bounded with bounded
gradients on half-lines, in \(C_sL^6\); forward uniqueness (E8) from the
end of the interval and backward uniqueness (E6) for their difference
give \(W_*^{\mu_0}=W_*\) on \((-\infty,0)\): (A\('\)) holds, \(W_*\ne0\) by
the pinching.  If \(m_d(\ell_0)>0\), scale invariance of the family gives
the bound on every unit window; and if the solution's discrete defect
were below \(\frac12m_d(\ell_0)\) along \(\lambda_k\to0\), a local limit
\(W'\in\mathcal F(x_0)\) of \(U_{\lambda_k}\) would have discrete defect at
most \(\frac12m_d(\ell_0)\) on \([0,1]\) by
Proposition~\ref{prop:c2v10-continuity} (the restriction of the Gaussian
to the fundamental domain differing by \(o(1)\)), contradicting the
definition of \(m_d\).  Exclusivity: a discretely self-similar element
has zero discrete defect.
\end{proof}

\begin{firewall}[Item (AC3): discretely self-similar elements are not excluded]
\label{fw:c2v10-DSS}
The exclusion of nonzero discretely self-similar solutions with profiles
in \(L^6\cap L^\infty\), bounded gradients, and the local energy
inequality is not imported and is not claimed.  Hence, for each
\(\ell_0\), the series does not decide between (A\('\)) and (B\('\)); it
records that (A\('\)) for some \(\ell_0\) is consistent with alternative
(B) of the continuous dichotomy (the continuous defect of a discretely
self-similar element is periodic with positive log-integral per period,
so \(\underline\delta>0\) is compatible with it), and that (B\('\)) for
every \(\ell_0\) is the statement that no rescaling limit at \(x_0\) has
any discrete self-similarity, with the quantitative bound
\(m_d(\ell_0)\) on every unit window.  Whether \(m_d(\ell_0)\) is bounded
below uniformly in \(\ell_0\), or tends to zero along some sequence of
periods, is not decided.
\end{firewall}

% =============================================================================
\section{Integrated V10 conclusion and next acquisition order}
\label{sec:c2v10-conclusion}
% =============================================================================

\begin{theorem}[What V10 establishes]
\label{thm:c2v10-conclusion}
Relative to the first cycle and V1--V9: the companion-audit decision
(Theorem~\ref{thm:c2v10-companion-decision}), the import ledger
(Theorem~\ref{thm:c2v10-imports}), the consolidated statement
(Theorem~\ref{thm:c2v10-consolidated}), the continuity of the profile
(Proposition~\ref{prop:c2v10-continuity}), and the discrete dichotomy
(Theorem~\ref{thm:c2v10-discrete-dichotomy}).  The full Navier--Stokes
stability/global-regularity theorem remains unproved.
\end{theorem}

\begin{proof}
The cited results.
\end{proof}

\begin{programme}[V11 acquisition order]
\label{prog:c2v11}
\begin{enumerate}[label=\textup{(AC\arabic*)},leftmargin=4.8em]
\item Periods.  Determine the behaviour of \(m_d(\ell_0)\) as
      \(\ell_0\to0\) (it is at most \(\ell_0^2\sup\|\partial_\sigma V\|_G^2\le C_Y^2\ell_0^2\)
      and at least of the order \(\ell_0^2m(1)\) in an averaged sense;
      make this exact) and as \(\ell_0\to\infty\) (whether
      \(\liminf_{\ell_0\to\infty}m_d(\ell_0)>0\), that is, whether
      profiles at widely separated scales stay apart in \(L^2(G)\)).
\item Recurrence.  Using the compactness of the family, show that for
      every \(W\) and every \(\varepsilon>0\) there are arbitrarily large
      \(\ell_0\) with \(\mathcal D_{\ell_0}(W;\sigma)<\varepsilon\) for some
      \(\sigma\) (almost-periodicity in log-scale by compactness), and
      determine whether this forces \(\liminf_{\ell_0\to\infty}m_d(\ell_0)=0\)
      and hence, by the discrete dichotomy along a sequence of periods,
      a limit element that is a fixed point of a sequence of rescalings.
\item The next compulsory companion audit is V15.
\end{enumerate}
\end{programme}

\begin{assessment}[Position after V10]
\label{ass:c2v10-position}
The second cycle is consolidated on the uniformly non-self-similar
structure of homogeneous points, with the continuous and discrete
dichotomies in place.  No full stability or global-regularity claim is
made.
\end{assessment}

% =============================================================================
\section*{Version V10 (second cycle) append-only record}
\addcontentsline{toc}{section}{Version V10 (second cycle) append-only record}
% =============================================================================

The complete V9 source of the second cycle was retained before this
continuation.  Its SHA--256 digest was
\begin{center}\scriptsize\ttfamily
866a46ad151fcb414e3a86a4b5fed623e8ac6530052f22ac0702ac87de655c47.
\end{center}
V10 performed the compulsory companion audit, consolidated V6--V9 with
the extended import ledger, and proved the discrete dichotomy.  No full
regularity claim is made.

% =============================================================================
\section{V11: small and large periods, recurrence in log-scale, and the
recurrent element}
\label{sec:c2v11-entry}
% =============================================================================

Items (AC1) and (AC2) of Programme~\ref{prog:c2v11} are carried out.  For
small periods the minimal discrete defect is at most \(C_Y^2\ell_0^2\),
and at least \(\frac14m(1)\ell_0^2\) under one more scale-invariant
derivative bound.  For large periods the compactness of the rescaling
family forces recurrence: every element has rescalings arbitrarily far
apart in log-scale that are arbitrarily close in the local topology, so
the minimal discrete defect has lower limit zero as the period tends to
infinity.  The rescaling flow on the compact family is a continuous flow
on a compact metric space, and Birkhoff's recurrence theorem gives a
minimal invariant set and hence a recurrent element: an element of the
family that is a local limit of its own rescalings along a sequence of
scales tending to infinity, and likewise along a sequence tending to
zero.  It is recorded that recurrence is not discrete self-similarity.
No global regularity claim is made.  The next compulsory companion
audit is V15.

% =============================================================================
\section{Small periods}
\label{sec:c2v11-small}
% =============================================================================

\begin{proposition}[Small periods]
\label{prop:c2v11-small}
Let \(x_0\) be homogeneous.  For every \(\ell_0>0\),
\(m_d(\ell_0)\le C_Y^2\ell_0^2\).  If moreover the elements of
\(\mathcal F(x_0)\) satisfy the next scale-invariant derivative bound
\(|\nabla^4W(s)|\le C_\infty''''(-s)^{-5/2}\) (so that
\(\sup_\sigma\|\partial_\sigma Y\|_G\le C_{Y'}\) with \(Y=\partial_\sigma V\)), then
\begin{equation}
 m_d(\ell_0)\ \ge\ \tfrac14m(1)\,\ell_0^2
 \qquad\text{for }0<\ell_0\le\bigl(m(1)/(2C_{Y'}^2)\bigr)^{1/2}.
 \label{eq:c2v11-small-lower}
\end{equation}
\end{proposition}

\begin{proof}
\(\|V(\sigma+\ell_0)-V(\sigma)\|_G\le\int_\sigma^{\sigma+\ell_0}\|Y\|_G\le\ell_0C_Y\)
gives the upper bound.  For the lower bound, Taylor's formula with
integral remainder in \(L^2(G)\),
\(V(\sigma+\ell_0)-V(\sigma)=\ell_0Y(\sigma)+\int_\sigma^{\sigma+\ell_0}(\sigma+\ell_0-\sigma')\partial_\sigma Y(\sigma')\dd\sigma'\),
gives \(\mathcal D_{\ell_0}(W;\sigma)^{1/2}\ge\ell_0\|Y(\sigma)\|_G-\frac12\ell_0^2C_{Y'}\),
and \((a-b)_+^2\ge\frac12a^2-b^2\) gives
\(\int_0^1\mathcal D_{\ell_0}\ge\frac12\ell_0^2\int_0^1\|Y\|_G^2-\frac14\ell_0^4C_{Y'}^2\ge\frac12\ell_0^2m(1)-\frac14\ell_0^4C_{Y'}^2\),
which is at least \(\frac14\ell_0^2m(1)\) in the stated range.  The
bound on \(\partial_\sigma Y=L_-Y-N'(V)Y\) in \(L^2(G)\) uses
\(\Delta Y\), which involves fourth derivatives of \(W\).
\end{proof}

% =============================================================================
\section{Large periods and recurrence}
\label{sec:c2v11-recurrence}
% =============================================================================

The local topology on \(\mathcal F(x_0)\) is metrizable: with
\(R_n=S_n=n\), \(\eta_n=1/n\), put
\(\rho(W,W'):=\sum_n2^{-n}\min\bigl(1,\|W-W'\|_{L^2((-n,-1/n);H^1(B_n))}\bigr)\);
\(\rho\)-convergence is the local convergence of Theorem V41-ancient-limit
(the pressures converge automatically, being Riesz transforms of the
velocities).  The rescaling flow \(\Theta_\tau(W):=W^{e^{\tau}}\),
\(\tau\in\R\), satisfies \(\Theta_\tau\circ\Theta_{\tau'}=\Theta_{\tau+\tau'}\),
and the profile of \(\Theta_\tau W\) is \(V(\cdot,\sigma-2\tau)\).

\begin{lemma}[The rescaling flow is a continuous flow on a compact metric space]
\label{lem:c2v11-flow}
\((\mathcal F(x_0),\rho)\) is a compact metric space, and
\((\tau,W)\mapsto\Theta_\tau W\) is continuous from \(\R\times\mathcal F(x_0)\)
to \(\mathcal F(x_0)\).
\end{lemma}

\begin{proof}
Compactness is Proposition V83-family(iv) of the first cycle applied to
\(\mathcal F(x_0)\) (sequential compactness in a metric space).  Joint
continuity: for \(W_k\to W\) and \(\tau_k\to\tau\), on each
\(L^2((-n,-1/n);H^1(B_n))\) the rescaling by \(e^{\tau_k}\) is a bounded
linear map converging strongly to the rescaling by \(e^\tau\) on smooth
fields with the uniform bounds \eqref{eq:c2v4-uniform} (dominated
convergence, the fields being uniformly bounded with bounded gradients
and second derivatives on the relevant compact sets), and
\(\|\Theta_{\tau_k}W_k-\Theta_\tau W\|\le\|\Theta_{\tau_k}(W_k-W)\|+\|\Theta_{\tau_k}W-\Theta_\tau W\|\),
the first term controlled by \(\rho(W_k,W)\) on slightly larger sets and
the second tending to zero.
\end{proof}

\begin{theorem}[Recurrence in log-scale]
\label{thm:c2v11-recurrence}
Let \(x_0\) be homogeneous.
\begin{enumerate}[label=\textup{(\roman*)},leftmargin=3.2em]
\item For every \(W\in\mathcal F(x_0)\) and \(\varepsilon>0\) there are
      \(\ell_0>0\) arbitrarily large and \(\sigma_0\in\R\) with
      \(\sup_{\sigma\in[\sigma_0,\sigma_0+1]}\mathcal D_{\ell_0}(W;\sigma)<\varepsilon\);
      the same with \(\ell_0<0\) arbitrarily negative.  Consequently
      \(\liminf_{\ell_0\to\infty}m_d(\ell_0)=0\).
\item (Recurrent element.)  There is \(W_{\rm rec}\in\mathcal F(x_0)\) and
      \(\tau_j\to+\infty\), \(\tau_j'\to-\infty\), with
      \(\Theta_{\tau_j}W_{\rm rec}\to W_{\rm rec}\) and
      \(\Theta_{\tau_j'}W_{\rm rec}\to W_{\rm rec}\) in \(\rho\); its profile
      satisfies \(\|V_{\rm rec}(\cdot,\sigma-2\tau_j)-V_{\rm rec}(\cdot,\sigma)\|_G\to0\)
      uniformly for \(\sigma\) in compact sets.  Moreover \(W_{\rm rec}\)
      can be taken in a minimal closed \(\Theta\)-invariant subset
      \(\mathcal M\subset\mathcal F(x_0)\), on which every element is
      uniformly recurrent: for every \(\varepsilon>0\) the set of
      \(\tau\) with \(\rho(\Theta_\tau W,W)<\varepsilon\) is relatively
      dense in \(\R\).
\end{enumerate}
\end{theorem}

\begin{proof}
(i) The sequence \(\Theta_kW\), \(k\in\mathbb N\), has a \(\rho\)-convergent
subsequence \(\Theta_{k_i}W\to W'\).  By Proposition~\ref{prop:c2v10-continuity},
the profiles converge in \(L^2(G)\) uniformly on \(\sigma\in[0,1]\), so for
\(i<i'\) large, \(\sup_{[0,1]}\|V(\cdot,\sigma-2k_i)-V(\cdot,\sigma-2k_{i'})\|_G^2<\varepsilon\),
which is the statement with \(\ell_0=2(k_{i'}-k_i)\) and
\(\sigma_0=-2k_{i'}\); \(k_{i'}-k_i\) can be taken arbitrarily large.
Negative \(k\) give negative \(\ell_0\).  Since
\(m_d(\ell_0)\le\int_0^1\mathcal D_{\ell_0}(\Theta_{k_{i'}}W;\sigma)\dd\sigma<\varepsilon\)
by scale invariance, \(\liminf_{\ell_0\to\infty}m_d(\ell_0)=0\).
(ii) By Lemma~\ref{lem:c2v11-flow} and Zorn's lemma (or the compactness
of nested closed invariant sets), \(\mathcal F(x_0)\) contains a minimal
nonempty closed \(\Theta\)-invariant set \(\mathcal M\); by Birkhoff's
recurrence theorem every point of a minimal set of a continuous flow on
a compact metric space is uniformly recurrent, in particular recurrent
in both time directions.  The profile statement is
Proposition~\ref{prop:c2v10-continuity}.
\end{proof}

\begin{corollary}[The recurrent element in the discrete dichotomy]
\label{cor:c2v11-recurrent-discrete}
For \(W_{\rm rec}\) and \(\ell_j:=2\tau_j\to\infty\),
\(\int_0^1\mathcal D_{\ell_j}(W_{\rm rec};\sigma)\dd\sigma\to0\), while for each
fixed \(\ell_0\) the discrete defect of \(W_{\rm rec}\) is at least
\(m_d(\ell_0)\) on every unit window; in alternative (B\('\)) for all
periods, \(W_{\rm rec}\) is discretely self-similar for no period but is
approximately self-similar for a sequence of periods tending to
infinity, with error tending to zero.
\end{corollary}

\begin{proof}
Theorem~\ref{thm:c2v11-recurrence}(ii) and Theorem~\ref{thm:c2v10-discrete-dichotomy}.
\end{proof}

\begin{firewall}[Recurrence is not self-similarity]
\label{fw:c2v11-not-DSS}
The recurrent element returns to itself approximately along a sequence
of log-scale translations tending to infinity; it need not be a fixed
point of any single rescaling, and the periods \(\ell_j\) need not be
commensurable.  Recurrence is a consequence of compactness alone and
holds on every compact scale-invariant family regardless of the
equations; it carries no information about the Navier--Stokes structure
beyond what compactness gives.  In particular it does not produce a
self-similar or discretely self-similar element and does not interact
with (E9).  Its use is that the recurrent element, and every element of
the minimal set \(\mathcal M\), has all its scale-invariant quantities
(\(\varphi\), \(d\), \(j\), \(\mathcal L\), \(\|Y\|_G^2\), the cross terms)
almost periodic in \(\sigma\), so that their log-averages over long
windows converge to exact means; this is the subject of the next
version.
\end{firewall}

% =============================================================================
\section{Integrated V11 conclusion and next acquisition order}
\label{sec:c2v11-conclusion}
% =============================================================================

\begin{theorem}[What V11 proves]
\label{thm:c2v11-conclusion}
Relative to the first cycle and V1--V10, modulo (E1)--(E9): the
small-period bounds (Proposition~\ref{prop:c2v11-small}), the flow
structure (Lemma~\ref{lem:c2v11-flow}), the recurrence theorem
(Theorem~\ref{thm:c2v11-recurrence}), and
Corollary~\ref{cor:c2v11-recurrent-discrete}.  The full Navier--Stokes
stability/global-regularity theorem remains unproved.
\end{theorem}

\begin{proof}
The cited results.
\end{proof}

\begin{programme}[V12 acquisition order]
\label{prog:c2v12}
\begin{enumerate}[label=\textup{(AC\arabic*)},leftmargin=4.8em]
\item Invariant measures.  By Krylov--Bogolyubov, the flow \(\Theta\) on
      \(\mathcal F(x_0)\) has invariant Borel probability measures; for
      each such measure \(\mathfrak m\) and each continuous
      scale-invariant quantity \(F\) on the family, the log-average of
      \(F\) along \(\mathfrak m\)-almost every orbit equals its ergodic
      mean.  Prove that the three identities of the cycle integrate to
      exact identities against \(\mathfrak m\):
      \(\int\xi\,\dd\mathfrak m=2\int\varphi\,\dd\mathfrak m\),
      \(\int(\|Y\|_G^2+\langle N,Y\rangle_G)\dd\mathfrak m=0\),
      \(\int(\nu\|\nabla Y\|_G^2+\frac12\|Y\|_G^2+\langle Y,N'Y\rangle_G)\dd\mathfrak m=0\),
      and record the continuity or semicontinuity of each integrand on
      the family.
\item Exact signs.  Derive from (AC1) the exact (not averaged) signs
      \(\int\langle N,Y\rangle_G\dd\mathfrak m=-\int\|Y\|_G^2\dd\mathfrak m\le-\underline\delta\)
      and \(\int\langle Y,N'Y\rangle_G\dd\mathfrak m\le-\frac12\underline\delta\)
      for every invariant measure, and determine whether an ergodic
      invariant measure supported on a single orbit closure yields
      pointwise consequences for the recurrent element.
\item The next compulsory companion audit is V15.
\end{enumerate}
\end{programme}

\begin{assessment}[Position after V11]
\label{ass:c2v11-position}
The rescaling family carries a continuous flow on a compact metric
space; its minimal sets consist of uniformly recurrent elements, and
the discrete defect vanishes along a sequence of periods for such
elements.  No full stability or global-regularity claim is made.
\end{assessment}

% =============================================================================
\section*{Version V11 (second cycle) append-only record}
\addcontentsline{toc}{section}{Version V11 (second cycle) append-only record}
% =============================================================================

The complete V10 source of the second cycle was retained before this
continuation.  Its SHA--256 digest was
\begin{center}\scriptsize\ttfamily
9443861b59aebac9e1220f2e3525f70914d1f49079a70481c05adbf58a0f8267.
\end{center}
V11 proved the small-period bounds, the recurrence in log-scale, and the
existence of a recurrent element of the rescaling family.  No full
regularity claim is made.

% =============================================================================
\section{V12: invariant measures of the rescaling flow and the exact mean
identities}
\label{sec:c2v12-entry}
% =============================================================================

Items (AC1) and (AC2) of Programme~\ref{prog:c2v12} are carried out.  The
rescaling flow on the compact family admits invariant Borel probability
measures, and every such measure turns the three identities of the
cycle, the first identity for the weighted energy, the second identity
for the Gaussian Dirichlet functional, and the energy identity for the
defect, into exact mean identities without endpoint errors.  The exact
signs follow: against every invariant measure the mean of the flux term
is minus twice the mean of the Dirichlet functional, the mean of the
cross term of the second identity is minus the mean Gaussian norm of
the defect, which is at least the uniform defect rate, and the mean of
the linearized cross term is minus the mean Gaussian damping of the
defect.  Pointwise consequences for a single element need unique
ergodicity, which is not established.  No global regularity claim is
made.  The next compulsory companion audit is V15.

% =============================================================================
\section{Invariant measures}
\label{sec:c2v12-measures}
% =============================================================================

Let \(x_0\) be homogeneous, \((\mathcal F(x_0),\rho)\) the compact metric
space of Lemma~\ref{lem:c2v11-flow}, and \(\Theta_\tau\) the rescaling
flow, whose profile action is \(\sigma\mapsto\sigma-2\tau\).  For
\(W\in\mathcal F(x_0)\) define the scale-invariant observables at the
unit scale, all evaluated on the profile at \(\sigma=0\):
\begin{equation}
 \varphi(W):=\|V\|_G^2,\quad
 d(W):=4\nu\|\nabla V\|_G^2,\quad
 j(W):=4\langle V,N(V)\rangle_G,\quad
 \mathcal L(W):=\tfrac18(d+2\varphi),\quad
 \mathfrak y(W):=\|Y\|_G^2,\quad
 \mathfrak g(W):=\nu\|\nabla Y\|_G^2,
 \label{eq:c2v12-observables}
\end{equation}
\(\mathfrak c(W):=\langle N(V),Y\rangle_G\), \(\mathfrak c'(W):=\langle Y,N'(V)Y\rangle_G\),
with \(Y=\partial_\sigma V\).  Along an orbit, \(F(\Theta_\tau W)\) is the
value of the observable on the profile of \(W\) at \(\sigma=-2\tau\).

\begin{lemma}[Measurability and boundedness]
\label{lem:c2v12-measurable}
\(\varphi\), \(d\), \(\mathcal L\) are continuous on \((\mathcal F(x_0),\rho)\);
\(\mathfrak y\) is lower semicontinuous; \(j\), \(\mathfrak c\), \(\mathfrak c'\),
\(\mathfrak g\) are Borel, as pointwise limits of continuous functions.
All are bounded on \(\mathcal F(x_0)\) except \(\mathfrak g\), which is
bounded in mean along every orbit: \(\int_\tau^{\tau+L}\mathfrak g(\Theta_{\tau'}W)\dd\tau'\le C(L+1)\).
\end{lemma}

\begin{proof}
Continuity of \(\varphi\) is Proposition~\ref{prop:c2v10-continuity};
of \(d\), the strong \(H^1(B_R)\) convergence at fixed time together
with the uniform Gaussian tails, the fixed-time convergence of
\(\nabla W_k\) in \(L^2(\psi)\) following from the same equicontinuity
argument applied to \(\nabla W_k\) (\(\partial_s\nabla W_k\) bounded in
\(L^2(J;L^2(\psi))\) by Lemma~\ref{lem:c2v7-higher} uniformly, the bounds
being uniform on the family).  Lower semicontinuity of \(\mathfrak y\)
at fixed \(\sigma\) follows from weak convergence of \(Y_k(\sigma)\) in
\(L^2(G)\), which holds for every \(\sigma\) by the same argument as
Proposition~\ref{prop:c2v4-lsc} once the fixed-time weak convergence of
\(\partial_sW_k\) is known; it holds because \(\partial_sW_k\) is bounded in
\(L^2(\psi)\) at every fixed time by the pointwise bounds
\eqref{eq:c2v4-uniform} for \(\Delta W_k\), \((W_k\cdot\nabla)W_k\) and by
\(\nabla P_k\in L^2\) uniformly, and any weak limit is \(\partial_sW\).  For
\(j\), \(\mathfrak c\), \(\mathfrak c'\), \(\mathfrak g\): each is the limit as
\(R\to\infty\) of the same expression with the Gaussian replaced by
\(G\chi_R\), and each truncated expression is continuous by the strong
local convergences and the weak convergence of pressures paired with
strongly convergent factors; a pointwise limit of continuous functions
is Borel.  Boundedness: \(\varphi\le C_\Phi\), \(d\le4\nu C\),
\(|j|\le C_J\), \(\mathfrak y\le C_Y^2\), \(|\mathfrak c|\le C_N^{1/2}C_Y\),
\(|\mathfrak c'|\le\|N'(V)Y\|_G\|Y\|_G\) with \(\|N'(V)Y\|_G\le(C_\infty'+C_\infty C(\nu))C_Y+\|\nabla P'\|_G\),
the last bounded in mean along orbits by Lemma~\ref{lem:c2v7-higher};
and \(\mathfrak g\) bounded in mean by \eqref{eq:c2v7-defect-energy}
integrated with the bounds on the other terms.
\end{proof}

\begin{theorem}[Existence of invariant measures]
\label{thm:c2v12-existence}
There is at least one \(\Theta\)-invariant Borel probability measure
\(\mathfrak m\) on \(\mathcal F(x_0)\); for every \(W\in\mathcal F(x_0)\)
every weak limit of the empirical measures
\(\frac1L\int_0^L\delta_{\Theta_\tau W}\dd\tau\) as \(L\to\infty\) is one;
and every minimal set \(\mathcal M\) carries one with full support in
\(\mathcal M\).
\end{theorem}

\begin{proof}
Krylov--Bogolyubov for a continuous flow on a compact metric space
(Lemma~\ref{lem:c2v11-flow}); the support statement is minimality.
\end{proof}

% =============================================================================
\section{The exact mean identities}
\label{sec:c2v12-identities}
% =============================================================================

\begin{theorem}[Exact mean identities]
\label{thm:c2v12-mean-identities}
For every \(\Theta\)-invariant Borel probability measure \(\mathfrak m\)
on \(\mathcal F(x_0)\), modulo (E1)--(E9):
\begin{align}
 \int j\,\dd\mathfrak m&=-\int(d+2\varphi)\,\dd\mathfrak m=-8\int\mathcal L\,\dd\mathfrak m,
 \label{eq:c2v12-mean1}\\
 \int\mathfrak c\,\dd\mathfrak m&=-\int\mathfrak y\,\dd\mathfrak m,
 \label{eq:c2v12-mean2}\\
 \int\mathfrak c'\,\dd\mathfrak m&=-\int\Bigl(\mathfrak g+\tfrac12\mathfrak y\Bigr)\dd\mathfrak m .
 \label{eq:c2v12-mean3}
\end{align}
Moreover \(\int\mathfrak y\,\dd\mathfrak m\ge\underline\delta\), so that
\begin{equation}
 \boxed{\quad
 \int\langle V,N(V)\rangle_G\,\dd\mathfrak m=-2\int\mathcal L\,\dd\mathfrak m\le-\tfrac12c_\Phi,
 \qquad
 \int\langle N(V),\partial_\sigma V\rangle_G\,\dd\mathfrak m\le-\underline\delta,
 \qquad
 \int\langle\partial_\sigma V,N'(V)\partial_\sigma V\rangle_G\,\dd\mathfrak m\le-\tfrac12\underline\delta .
 \quad}
 \label{eq:c2v12-exact-signs}
\end{equation}
\end{theorem}

\begin{proof}
Let \(\mathcal G\) be one of \(\varphi\), \(\mathcal L\), \(\frac12\mathfrak y\),
bounded on the family, and let \(F\) be its derivative along the flow:
by \eqref{eq:c2v1-identity}, \eqref{eq:c2v2-second},
\eqref{eq:c2v7-defect-energy}, with \(\dd\sigma=-2\dd\tau\),
\(\frac{\dd}{\dd\tau}\varphi(\Theta_\tau W)=d+2\varphi+j\),
\(\frac{\dd}{\dd\tau}\mathcal L(\Theta_\tau W)=2(\mathfrak y+\mathfrak c)\),
\(\frac{\dd}{\dd\tau}\frac12\mathfrak y(\Theta_\tau W)=2(\mathfrak g+\frac12\mathfrak y+\mathfrak c')\),
each evaluated at \(\Theta_\tau W\), for a.e.\ \(\tau\), and each
\(\tau\mapsto\mathcal G(\Theta_\tau W)\) absolutely continuous.  Then for
every \(L>0\), by invariance and Fubini (the integrands being Borel and
integrable along orbits),
\(\int F\,\dd\mathfrak m=\int\frac1L\int_0^LF(\Theta_\tau W)\dd\tau\,\dd\mathfrak m(W)
=\int\frac{\mathcal G(\Theta_LW)-\mathcal G(W)}L\dd\mathfrak m(W)\),
bounded by \(2\sup|\mathcal G|/L\to0\).  Hence \(\int F\,\dd\mathfrak m=0\)
for the three choices, which are \eqref{eq:c2v12-mean1}--\eqref{eq:c2v12-mean3}
(the integrability of \(\mathfrak g\) against \(\mathfrak m\) follows from
the mean bound of Lemma~\ref{lem:c2v12-measurable} and invariance).
For \(\int\mathfrak y\,\dd\mathfrak m\), invariance gives
\(\int\mathfrak y\,\dd\mathfrak m=\int\frac1L\int_0^L\mathfrak y(\Theta_\tau W)\dd\tau\,\dd\mathfrak m\),
and \(\int_0^L\mathfrak y(\Theta_\tau W)\dd\tau=\frac12\int_{-2L}^0\|Y(\sigma)\|_G^2\dd\sigma\ge\frac12m(2L)\),
so \(\int\mathfrak y\,\dd\mathfrak m\ge m(2L)/(2L)\to\underline\delta\).  The
signs \eqref{eq:c2v12-exact-signs} follow, using \(\mathcal L\ge c_\Phi/4\)
and \(\mathfrak g\ge0\).
\end{proof}

\begin{assessment}[What the mean identities say]
\label{ass:c2v12-meaning}
Against any scale-invariant statistics of the blow-up, the three
identities hold exactly: the mean flux term equals minus the mean of
dissipation plus twice weighted energy (the inward flux of the first
cycle, now exact); the mean cross term equals minus the mean squared
defect, which is at least the uniform defect rate; and the mean
linearized cross term equals minus the mean Gaussian damping of the
defect.  The exact signs are therefore properties of every invariant
measure, not of long windows.  A contradiction would follow from a
lower bound \(\int\mathfrak c\,\dd\mathfrak m\ge-\theta\int\mathfrak y\,\dd\mathfrak m\)
with \(\theta<1\) for some invariant measure, that is, from a mean form
of \eqref{eq:c2v2-C1prime}; the series has no such bound.  No global
regularity claim is made.
\end{assessment}

\begin{firewall}[Item (AC2): pointwise consequences need unique ergodicity]
\label{fw:c2v12-AC2}
If \(\mathfrak m\) is ergodic, the log-averages of each observable
converge to its mean along \(\mathfrak m\)-almost every orbit (Birkhoff's
ergodic theorem), but not necessarily along the orbit of a given
element such as \(W_{\rm rec}\).  If the minimal set \(\mathcal M\) were
uniquely ergodic, the averages would converge uniformly on
\(\mathcal M\) for continuous observables, and \(W_{\rm rec}\) would
satisfy the exact signs as limits of its own long-window averages;
unique ergodicity is not established.  Along the orbit of any element
the averaged signs of V8 hold with the stated rates, which is all that
is available pointwise.
\end{firewall}

% =============================================================================
\section{Integrated V12 conclusion and next acquisition order}
\label{sec:c2v12-conclusion}
% =============================================================================

\begin{theorem}[What V12 proves]
\label{thm:c2v12-conclusion}
Relative to the first cycle and V1--V11, modulo (E1)--(E9): the
measurability lemma, the existence of invariant measures, the exact mean
identities \eqref{eq:c2v12-mean1}--\eqref{eq:c2v12-mean3}, and the exact
signs \eqref{eq:c2v12-exact-signs}.  The full Navier--Stokes
stability/global-regularity theorem remains unproved.
\end{theorem}

\begin{proof}
The cited results.
\end{proof}

\begin{programme}[V13 acquisition order]
\label{prog:c2v13}
\begin{enumerate}[label=\textup{(AC\arabic*)},leftmargin=4.8em]
\item Blow-up measures of the solution.  Define the empirical
      scale-measures of the solution at \(x_0\),
      \(\mathfrak m_{\lambda,L}:=\frac1L\int_0^L\delta_{U_{\lambda e^{-\tau}}}\dd\tau\)
      with \(U_\mu\) the rescaling of \(u\) at \((x_0,T_*)\) by \(\mu\),
      regarded as measures on the compactification of the family by the
      rescalings; prove that every weak limit as \(\lambda\to0\),
      \(L\to\infty\) is a \(\Theta\)-invariant probability measure on
      \(\mathcal F(x_0)\); conclude that the solution's own long-window
      log-averages of the signed quantities at small scales satisfy the
      exact signs in the limit.
\item Determine whether the set of blow-up measures of a point is
      always a single measure (unique statistics of the blow-up), or
      can contain several, and record the consequence for the
      recurrent element.
\item The next compulsory companion audit is V15.
\end{enumerate}
\end{programme}

\begin{assessment}[Position after V12]
\label{ass:c2v12-position}
The identities of the cycle are exact against every invariant measure
of the rescaling flow, with exact signs.  No full stability or
global-regularity claim is made.
\end{assessment}

% =============================================================================
\section*{Version V12 (second cycle) append-only record}
\addcontentsline{toc}{section}{Version V12 (second cycle) append-only record}
% =============================================================================

The complete V11 source of the second cycle was retained before this
continuation.  Its SHA--256 digest was
\begin{center}\scriptsize\ttfamily
e5b9ead17d0b77625875ec5bccbd1efb59b7c3414afb1cff9228ff339efab68e.
\end{center}
V12 proved the existence of invariant measures of the rescaling flow
and the exact mean identities with their signs.  No full regularity
claim is made.

% =============================================================================
\section{V13: blow-up measures of the solution at a homogeneous point}
\label{sec:c2v13-entry}
% =============================================================================

Items (AC1) and (AC2) of Programme~\ref{prog:c2v13} are carried out.  The
rescalings of the solution at a homogeneous singular point, together
with the rescaling family, form a compact metric space on which the
rescaling flow acts continuously at small scales; the empirical
measures of the solution's rescaling orbit over long log-windows at
small scales are tight, and every weak limit is an invariant
probability measure on the rescaling family.  These blow-up measures
carry the exact mean identities and signs of V12, so the solution's own
long-window statistics at small scales satisfy them in the limit.
Whether a point has a unique blow-up measure is not decided; the
consequences of uniqueness are recorded.  No global regularity claim is
made.  The next compulsory companion audit is V15.

% =============================================================================
\section{The compactified orbit and the empirical measures}
\label{sec:c2v13-orbit}
% =============================================================================

Let \(x_0\) be homogeneous, \(\lambda_0\) a small scale, and for
\(0<\mu\le\lambda_0\) let \(U_\mu(s,y):=\mu u(x_0+\mu y,T_*+\mu^2s)\),
extended periodically in \(y\) with period \(2\pi/\mu\).  Put
\begin{equation}
 \mathcal K:=\overline{\{U_\mu:0<\mu\le\lambda_0\}}\cup\mathcal F(x_0)
 \label{eq:c2v13-K}
\end{equation}
with the metric \(\rho\) of Lemma~\ref{lem:c2v11-flow}, and let
\(\Theta_\tau U_\mu=U_{\mu e^{\tau}}\) for \(\mu e^\tau\le\lambda_0\).  For
\(\lambda\le\lambda_0\) and \(L>0\) define the empirical scale-measure
\begin{equation}
 \mathfrak m_{\lambda,L}:=\frac1L\int_0^L\delta_{U_{\lambda e^{-\tau}}}\,\dd\tau ,
 \label{eq:c2v13-empirical}
\end{equation}
a Borel probability measure on \(\mathcal K\) supported on the rescalings
at scales in \([\lambda e^{-L},\lambda]\).

\begin{lemma}[Compactness and continuity]
\label{lem:c2v13-compact}
\((\mathcal K,\rho)\) is a compact metric space; every sequence
\(U_{\mu_k}\) with \(\mu_k\to0\) has a subsequence converging in \(\rho\)
to an element of \(\mathcal F(x_0)\), and \(\mathcal F(x_0)\) is exactly
the set of such limits together with their rescaling limits.  The map
\((\tau,X)\mapsto\Theta_\tau X\) is continuous on the set where it is
defined, and for every fixed \(\tau_0\) and every \(\varepsilon>0\) there is
\(\lambda_1>0\) such that \(\Theta_{\tau_0}\) is defined on the
\(\varepsilon\)-neighbourhood in scale of every \(U_\mu\) with
\(\mu\le\lambda_1\).
\end{lemma}

\begin{proof}
Sequences with \(\mu_k\) bounded away from zero converge along
subsequences by continuity of \(\mu\mapsto U_\mu\) in \(\rho\) (smoothness
of \(u\) on \([0,T_*)\)); sequences with \(\mu_k\to0\) converge along
subsequences to ancient limits by Theorem V41-ancient-limit of the
first cycle, which are in \(\mathcal F(x_0)\) by definition; sequences in
\(\mathcal F(x_0)\) converge by its compactness.  Continuity of \(\Theta\)
on the rescalings is the smoothness of \(u\), on \(\mathcal F(x_0)\) is
Lemma~\ref{lem:c2v11-flow}, and across the boundary follows from the
uniform bounds \eqref{eq:c2v4-uniform} which hold for the \(U_\mu\) on
the relevant compact sets by the type-I rates.  The last statement is
\(\mu e^{\tau_0}\le\lambda_0\) for \(\mu\le\lambda_0e^{-\tau_0}\).
\end{proof}

\begin{theorem}[Blow-up measures]
\label{thm:c2v13-blowup-measures}
Let \(\lambda_k\to0\) and \(L_k\to\infty\).  The measures
\(\mathfrak m_{\lambda_k,L_k}\) have weakly convergent subsequences, and
every weak limit \(\mathfrak m\) is a \(\Theta\)-invariant Borel
probability measure supported on \(\mathcal F(x_0)\).  The set
\(\mathcal B(x_0)\) of all such limits (over all sequences), the
\emph{blow-up measures} of \(x_0\), is a nonempty weakly compact set
of invariant measures on \(\mathcal F(x_0)\).
\end{theorem}

\begin{proof}
Probability measures on the compact metric space \(\mathcal K\) are
weakly sequentially compact.  Let \(\mathfrak m\) be a weak limit.
Support: for every \(\lambda_1>0\), \(\mathfrak m_{\lambda_k,L_k}\) gives
zero mass to \(\{U_\mu:\mu>\lambda_1\}\) once \(\lambda_k<\lambda_1\), and the
set \(\overline{\{U_\mu:\mu\le\lambda_1\}}\) is closed, so \(\mathfrak m\) is
supported on \(\bigcap_{\lambda_1}\overline{\{U_\mu:\mu\le\lambda_1\}}\),
which is the set of limit points of \(U_\mu\) as \(\mu\to0\), contained
in \(\mathcal F(x_0)\) by Lemma~\ref{lem:c2v13-compact}.  Invariance:
for \(F\) continuous on \(\mathcal K\) and \(\tau_0>0\), with
\(\Theta_{\tau_0}\) defined on the support of \(\mathfrak m_{\lambda_k,L_k}\)
for large \(k\) (Lemma~\ref{lem:c2v13-compact}),
\(\int F\circ\Theta_{\tau_0}\dd\mathfrak m_{\lambda_k,L_k}-\int F\dd\mathfrak m_{\lambda_k,L_k}
=\frac1{L_k}\Bigl(\int_{-\tau_0}^{0}-\int_{L_k-\tau_0}^{L_k}\Bigr)F(U_{\lambda_ke^{-\tau}})\dd\tau\),
bounded by \(2\tau_0\sup|F|/L_k\to0\); since \(F\circ\Theta_{\tau_0}\) is
continuous on the support (Lemma~\ref{lem:c2v13-compact}) and the
limit measure is supported on \(\mathcal F(x_0)\) where \(\Theta_{\tau_0}\)
is defined and continuous, passing to the limit gives
\(\int F\circ\Theta_{\tau_0}\dd\mathfrak m=\int F\dd\mathfrak m\) for all
continuous \(F\), which is invariance (for \(\tau_0<0\) likewise).  The
set of limits is nonempty and weakly closed by a diagonal argument.
\end{proof}

\begin{corollary}[The solution's statistics at small scales]
\label{cor:c2v13-solution-statistics}
Modulo (E1)--(E9), for every blow-up measure \(\mathfrak m\in\mathcal B(x_0)\)
the exact mean identities \eqref{eq:c2v12-mean1}--\eqref{eq:c2v12-mean3}
and the exact signs \eqref{eq:c2v12-exact-signs} hold.  Consequently,
for every continuous observable \(F\) on \(\mathcal K\) and every
sequence \(\lambda_k\to0\), \(L_k\to\infty\), the solution's long-window
log-averages \(\frac1{L_k}\int_0^{L_k}F(U_{\lambda_ke^{-\tau}})\dd\tau\) have
all their limit points in \(\{\int F\,\dd\mathfrak m:\mathfrak m\in\mathcal B(x_0)\}\);
for the lower semicontinuous observable \(\mathfrak y\),
\(\liminf_k\frac1{L_k}\int_0^{L_k}\mathfrak y(U_{\lambda_ke^{-\tau}})\dd\tau\ge\inf_{\mathcal B(x_0)}\int\mathfrak y\,\dd\mathfrak m\ge\underline\delta\).
\end{corollary}

\begin{proof}
Theorem~\ref{thm:c2v12-mean-identities} applies to every invariant
measure on \(\mathcal F(x_0)\).  The limit statements are the definition
of weak convergence for continuous \(F\), and the portmanteau inequality
for lower semicontinuous \(\mathfrak y\), together with
\(\int\mathfrak y\,\dd\mathfrak m\ge\underline\delta\).
\end{proof}

\begin{assessment}[The blow-up measure]
\label{ass:c2v13-meaning}
A homogeneous singular point carries, beyond its ancient limits, a
family of scale-invariant statistics: the blow-up measures, each an
invariant probability measure of the rescaling flow on the family of
ancient limits, obtained as the limiting distribution of the solution's
rescalings over long log-windows of small scales.  Against each of them
the flux term has mean \(-2\int\mathcal L\), the cross term has mean
\(-\int\|\partial_\sigma V\|_G^2\le-\underline\delta\), and the linearized
cross term has mean \(-\int(\nu\|\nabla Y\|_G^2+\frac12\|Y\|_G^2)\).  This
is the statistical form of the uniformly non-self-similar structure;
it is exact, and it is the natural object against which any mean
inequality for the Navier--Stokes nonlinearity would be tested.  No
global regularity claim is made.
\end{assessment}

\begin{firewall}[Item (AC2): uniqueness of the blow-up measure]
\label{fw:c2v13-uniqueness}
\(\mathcal B(x_0)\) may contain more than one measure: different
sequences of scales may sample different parts of \(\mathcal F(x_0)\),
and nothing in the series prevents this.  If \(\mathcal B(x_0)\) is a
single measure \(\mathfrak m_0\), then the solution's log-averages of
every continuous observable converge as \(\lambda\to0\), \(L\to\infty\), to
the \(\mathfrak m_0\)-mean, and \(\mathfrak m_0\) is supported on the
closure of the set of limit points of \(U_\mu\), which contains a
minimal set; if moreover \(\mathfrak m_0\) is ergodic and its support
is uniquely ergodic, the recurrent element of V11 in that support
satisfies the exact signs as limits of its own averages.  None of these
uniqueness properties is established.
\end{firewall}

% =============================================================================
\section{Integrated V13 conclusion and next acquisition order}
\label{sec:c2v13-conclusion}
% =============================================================================

\begin{theorem}[What V13 proves]
\label{thm:c2v13-conclusion}
Relative to the first cycle and V1--V12, modulo (E1)--(E9): the
compactness of the compactified orbit (Lemma~\ref{lem:c2v13-compact}),
the blow-up measures (Theorem~\ref{thm:c2v13-blowup-measures}), and the
solution's small-scale statistics
(Corollary~\ref{cor:c2v13-solution-statistics}).  The full Navier--Stokes
stability/global-regularity theorem remains unproved.
\end{theorem}

\begin{proof}
The cited results.
\end{proof}

\begin{programme}[V14 acquisition order]
\label{prog:c2v14}
\begin{enumerate}[label=\textup{(AC\arabic*)},leftmargin=4.8em]
\item Rate-free second identity.  Record that the second identity holds
      for every smooth solution on the torus with the periodized weight
      (as the first identity does, V96 of the first cycle), define the
      solution's Gaussian Dirichlet functional \(\mathcal L_u(\lambda)\)
      and its defect \(Y_u\), and derive the rate-free ledger
      \(\mathcal L_u(b)-\mathcal L_u(a)=2\int_a^b(\|Y_u\|_G^2+\langle N_u,Y_u\rangle_G)\dd\lambda/\lambda\);
      state the necessary condition at a point where \(\mathcal L_u\) is
      unbounded as the scale tends to zero.
\item Determine the relation between unboundedness of \(\mathcal L_u\) and
      of \(\Phi_u\) (equivalent by the type-I dissipation bound under the
      rate; in general \(\mathcal L_u\ge\frac14\Phi_u\), so unbounded
      \(\Phi_u\) forces unbounded \(\mathcal L_u\)).
\item The next compulsory companion audit is V15.
\end{enumerate}
\end{programme}

\begin{assessment}[Position after V13]
\label{ass:c2v13-position}
The solution's small-scale statistics at a homogeneous point are
invariant measures of the rescaling flow carrying the exact mean
identities and signs.  No full stability or global-regularity claim is
made.
\end{assessment}

% =============================================================================
\section*{Version V13 (second cycle) append-only record}
\addcontentsline{toc}{section}{Version V13 (second cycle) append-only record}
% =============================================================================

The complete V12 source of the second cycle was retained before this
continuation.  Its SHA--256 digest was
\begin{center}\scriptsize\ttfamily
8c8517055d23bb8f0f08a829c58951ed18b06ccdee453057645e60f62267096c.
\end{center}
V13 constructed the blow-up measures of the solution at a homogeneous
point and transferred the exact mean identities and signs to the
solution's small-scale statistics.  No full regularity claim is made.

% =============================================================================
\section{V14: the rate-free second identity for the solution and the
Gaussian Dirichlet ledger}
\label{sec:c2v14-entry}
% =============================================================================

Items (AC1) and (AC2) of Programme~\ref{prog:c2v14} are carried out.  The
second identity is shown to hold exactly for every smooth solution on
the torus, at every centre and every scale, for the periodic extension
of the rescaled solution against the whole-space Gaussian, which is the
same object as the periodized weight of the first cycle.  The solution's
Gaussian Dirichlet functional and its self-similarity defect are
defined rate-free, and the ledger between two scales is derived; at a
point where the Dirichlet functional is unbounded as the scale tends to
zero, the cross term must be negative and dominate the squared defect
with a divergent log-integral.  The relation between unboundedness of
the Dirichlet functional and of the weighted energy is recorded.  No
global regularity claim is made.  The next version is the compulsory
companion audit.

% =============================================================================
\section{The periodic profile and the second identity without a rate}
\label{sec:c2v14-identity}
% =============================================================================

Let \(u\) be any smooth mean-zero solution on \([0,T_*)\times\Torus^3\)
with pressure \(p\), \(x_0\in\Torus^3\), and for \(0<\lambda<T_*^{1/2}\),
\(\sigma=-2\log\lambda\), define the periodic profile and its pressure
\begin{equation}
 V_u(z,\sigma):=\lambda\,u\bigl(x_0+\lambda z,\,T_*-\lambda^2\bigr),
 \qquad
 P_u(z,\sigma):=\lambda^2\,p\bigl(x_0+\lambda z,\,T_*-\lambda^2\bigr),
 \qquad z\in\R^3,
 \label{eq:c2v14-periodic-profile}
\end{equation}
periodic in \(z\) with period \(2\pi/\lambda\); \(\|V_u\|_G^2=\Phi_u(\lambda)\)
and \(4\nu\|\nabla V_u\|_G^2=\lambda D_u(\lambda)\) are the periodized
quantities of the first cycle, since \(\int_{\R^3}f_{\rm per}G=\int_{\Torus^3}f\,\Psi\).
Define
\begin{equation}
 \mathcal L_u(\lambda):=\frac\nu2\|\nabla V_u\|_G^2+\frac14\|V_u\|_G^2=\frac18\bigl(\lambda D_u(\lambda)+2\Phi_u(\lambda)\bigr),
 \qquad
 Y_u:=\partial_\sigma V_u .
 \label{eq:c2v14-Lu}
\end{equation}

\begin{theorem}[Rate-free second identity]
\label{thm:c2v14-second-rate-free}
For every smooth solution, centre, and \(0<\lambda<T_*^{1/2}\):
\begin{enumerate}[label=\textup{(\roman*)},leftmargin=3.2em]
\item \(Y_u=L_-V_u-N(V_u)\) with \(L_-=\nu\Delta-\frac12z\cdot\nabla-\frac12\)
      and \(N(V_u)=(V_u\cdot\nabla)V_u+\nabla P_u\); \(Y_u\) is smooth,
      divergence-free, of at most linear growth in \(z\), and
      \(Y_u,\nabla Y_u,N(V_u)\in L^2(G)\).
\item \(\sigma\mapsto\mathcal L_u\) is smooth and
      \begin{equation}
       \boxed{\quad
       \frac{\dd}{\dd\sigma}\mathcal L_u=-\|Y_u\|_G^2-\langle N(V_u),Y_u\rangle_G,
       \qquad
       \lambda\frac{\dd}{\dd\lambda}\mathcal L_u=2\|Y_u\|_G^2+2\langle N(V_u),Y_u\rangle_G .
       \quad}
       \label{eq:c2v14-second}
      \end{equation}
\item (Ledger.)  For \(0<a<b<T_*^{1/2}\),
      \begin{equation}
       \mathcal L_u(b)-\mathcal L_u(a)=2\int_a^b\bigl(\|Y_u\|_G^2+\langle N(V_u),Y_u\rangle_G\bigr)\frac{\dd\lambda}\lambda .
       \label{eq:c2v14-ledger}
      \end{equation}
\end{enumerate}
No rate assumption is used.
\end{theorem}

\begin{proof}
(i) The chain rule on \eqref{eq:c2v14-periodic-profile} with
\(\dd\lambda/\dd\sigma=-\lambda/2\), \(\dd t/\dd\sigma=\lambda^2\):
\(\partial_\sigma V_u=-\frac12V_u-\frac12z\cdot\nabla V_u+\lambda^3\partial_tu\), and
\(\lambda^3\partial_tu=\lambda^3(\nu\Delta u-(u\cdot\nabla)u-\nabla p)=\nu\Delta V_u-(V_u\cdot\nabla)V_u-\nabla P_u\)
by \(\nabla_z=\lambda\nabla_x\).  At fixed \(\sigma\) (fixed \(t<T_*\)),
\(u(t)\), \(p(t)\) are smooth and periodic, so \(V_u\), \(P_u\), \(N(V_u)\)
and all their derivatives are bounded periodic functions of \(z\), and
\(z\cdot\nabla V_u\) has linear growth; all are in \(L^2(G)\) with all
derivatives.  (ii) Differentiate \(\mathcal L_u\) under the integral (the
integrands are smooth in \((z,\sigma)\) with Gaussian-integrable bounds
locally uniform in \(\sigma\)):
\(\frac{\dd}{\dd\sigma}\mathcal L_u=\nu\langle\nabla V_u,\nabla Y_u\rangle_G+\frac12\langle V_u,Y_u\rangle_G=-\langle L_-V_u,Y_u\rangle_G\)
by the \(G\)-symmetry of \(L_-\) (Lemma~\ref{lem:c2v2-dictionary}), the
integration by parts being justified by the Gaussian decay against the
polynomially bounded fields.  Insert \(L_-V_u=Y_u+N(V_u)\).  The
\(\lambda\)-form uses \(\dd\sigma=-2\dd\lambda/\lambda\).  (iii) Integrate.
\end{proof}

\begin{remark}[Consistency with the first cycle]
\label{rem:c2v14-consistency}
The first identity for \(\Phi_u\) (Theorem V87-derivative-u, rate-free by
Proposition V96-rate-free of the first cycle) is
\(\frac{\dd}{\dd\sigma}\|V_u\|_G^2=2\langle V_u,Y_u\rangle_G=-8\mathcal L_u-2\langle V_u,N(V_u)\rangle_G\)
in the present notation; the second identity \eqref{eq:c2v14-second} is
its companion for the positive terms.  Under the type-I rate,
\(\mathcal L_u\le\frac18(4\nu C_I^2\nu^{3/2}+2C_\Phi')\) on small scales,
and at a homogeneous point \(\mathcal L_u\ge\frac14c_\Phi'\); the
solution's version of the dichotomy, Theorem~\ref{thm:c2v6-solution},
used exactly this.
\end{remark}

% =============================================================================
\section{The Gaussian Dirichlet ledger without a rate}
\label{sec:c2v14-ledger}
% =============================================================================

\begin{proposition}[Relation between \(\mathcal L_u\) and \(\Phi_u\)]
\label{prop:c2v14-relation}
For every smooth solution and centre, \(\mathcal L_u(\lambda)\ge\frac14\Phi_u(\lambda)\).
Hence if \(\Phi_u(a_k)\to\infty\) along \(a_k\to0\) then
\(\mathcal L_u(a_k)\to\infty\).  Under the type-I rate
\(\lambda D_u(\lambda)\le4\nu C_I^2\nu^{3/2}\), so \(\mathcal L_u\) is bounded on
small scales if and only if \(\Phi_u\) is, if and only if a critical
Morrey bound holds at the parabolic time (Proposition V96-rate-free of
the first cycle).  In general, \(\mathcal L_u\) unbounded with \(\Phi_u\)
bounded means that the scale-invariant Gaussian dissipation
\(\lambda D_u(\lambda)\) is unbounded along a sequence of scales, which
excludes the type-I rate.
\end{proposition}

\begin{proof}
\(\mathcal L_u=\frac18(\lambda D_u+2\Phi_u)\ge\frac14\Phi_u\) since \(D_u\ge0\);
the rest is the type-I bound and the cited proposition.
\end{proof}

\begin{theorem}[Necessary condition in the Dirichlet ledger]
\label{thm:c2v14-necessary}
Let \(u\) be a smooth solution with \(T_*<\infty\) and \(x_0\in\Torus^3\).
If \(\mathcal L_u(a_k)\to\infty\) along \(a_k\to0\), then for every fixed
\(b\),
\begin{equation}
 \boxed{\quad
 \int_{a_k}^{b}\bigl(\|Y_u\|_G^2+\langle N(V_u),Y_u\rangle_G\bigr)\frac{\dd\lambda}\lambda
 =\frac{\mathcal L_u(b)-\mathcal L_u(a_k)}2\longrightarrow-\infty ,
 \quad}
 \label{eq:c2v14-necessary}
\end{equation}
so that \(\int_{a_k}^b\langle N(V_u),Y_u\rangle_G\,\dd\lambda/\lambda\le-\int_{a_k}^b\|Y_u\|_G^2\,\dd\lambda/\lambda-\frac12\mathcal L_u(a_k)+\frac12\mathcal L_u(b)\):
the nonlinearity must be anti-aligned with the self-similarity defect
with a log-integral diverging at least as fast as \(\frac12\mathcal L_u(a_k)\),
in addition to compensating the squared defect.  In particular
\(\langle N(V_u),Y_u\rangle_G<0\) on a set of scales of infinite
logarithmic measure.
\end{theorem}

\begin{proof}
\eqref{eq:c2v14-ledger} with \(a=a_k\); the conclusions follow from
\(\|Y_u\|_G^2\ge0\).
\end{proof}

\begin{assessment}[The two rate-free ledgers]
\label{ass:c2v14-two-ledgers}
Without any rate, a hypothetical singular point of a smooth solution
now has two exact Gaussian ledgers: the first identity of the first
cycle, with the necessary condition of Corollary V96-necessary
(supercritical excess inward flux if the weighted energy is unbounded),
and the second identity, with the necessary condition
\eqref{eq:c2v14-necessary} (divergent anti-alignment of the
nonlinearity with the self-similarity defect if the Dirichlet
functional is unbounded).  Both conditions are consequences, not
obstructions: the series has no bound on the cross terms without a
rate, and no bound on \(\mathcal L_u\) without the Morrey bound.  Under
the type-I rate the ledgers are bounded and the second cycle's
dichotomy applies.  The strong type-II gate (G3) of the first cycle is
therefore now characterized in the Dirichlet ledger by the alternative:
either \(\mathcal L_u\) is bounded on small scales at every singular
point (a rate-free Morrey-type bound on energy and Gaussian dissipation
at the parabolic time, not proved), or \eqref{eq:c2v14-necessary} holds
at some point.  No global regularity claim is made.
\end{assessment}

\begin{firewall}[Item (AC2): what is not decided]
\label{fw:c2v14-AC2}
Whether \(\mathcal L_u\) can be unbounded at a singular point of a smooth
solution is not decided; under the type-I rate it cannot.  The
first cycle's rate-free local energy bound (Theorem V66-local-energy)
bounds \(\Phi_u(\lambda)\) by \(C_0\lambda^{-1}(\epsilon(2\lambda^2)+\lambda^{-1/2}E_*^{3/4}\epsilon(2\lambda^2)^{3/4})\)
up to the Gaussian tails, which is bounded only when the Leray energy
spent in the parabolic window is of the parabolic order \(\lambda\), that
is, under a type-I-type condition along the sequence; without it the
bound is of order \(\lambda^{-3/2}\epsilon^{3/4}\), not bounded in general.
Item (AC2) is closed with this record.
\end{firewall}

% =============================================================================
\section{Integrated V14 conclusion and next acquisition order}
\label{sec:c2v14-conclusion}
% =============================================================================

\begin{theorem}[What V14 proves]
\label{thm:c2v14-conclusion}
Relative to the first cycle and V1--V13, without any rate: the
rate-free second identity \eqref{eq:c2v14-second} and its ledger
\eqref{eq:c2v14-ledger}, the relation of
Proposition~\ref{prop:c2v14-relation}, and the necessary condition
\eqref{eq:c2v14-necessary}.  The full Navier--Stokes
stability/global-regularity theorem remains unproved.
\end{theorem}

\begin{proof}
The cited results.
\end{proof}

\begin{programme}[V15 acquisition order, with compulsory companion audit]
\label{prog:c2v15}
\begin{enumerate}[label=\textup{(AC\arabic*)},leftmargin=4.8em]
\item Perform the compulsory review of the current SM and YM notes;
      record digests; import nothing numerical.
\item Consolidate V11--V14 into one statement: recurrence and the
      recurrent element, invariant and blow-up measures with the exact
      mean identities and signs, and the rate-free Dirichlet ledger.
\item Record the position of the second cycle after fifteen versions
      relative to the three gates: (G1) reduced at homogeneous points
      to the exact mean signs against blow-up measures; (G2), (G3)
      constrained by the two rate-free ledgers.
\item The next compulsory companion audit after V15 is V20.
\end{enumerate}
\end{programme}

\begin{assessment}[Position after V14]
\label{ass:c2v14-position}
The second identity is a rate-free ledger for every smooth solution,
with a necessary condition at any point where the Gaussian Dirichlet
functional is unbounded.  No full stability or global-regularity claim
is made.
\end{assessment}

% =============================================================================
\section*{Version V14 (second cycle) append-only record}
\addcontentsline{toc}{section}{Version V14 (second cycle) append-only record}
% =============================================================================

The complete V13 source of the second cycle was retained before this
continuation.  Its SHA--256 digest was
\begin{center}\scriptsize\ttfamily
0d473b640b6b9041326eea193251ef52d0f37f91ee8c2fe47af1bf4caf64c617.
\end{center}
V14 proved the rate-free second identity for the solution, its ledger,
and the necessary condition at points where the Gaussian Dirichlet
functional is unbounded.  No full regularity claim is made.

% =============================================================================
\section{V15: compulsory companion audit, the consolidated statistical
structure, and the position after fifteen versions}
\label{sec:c2v15-entry}
% =============================================================================

This is the fifteenth version of the second cycle.  The compulsory
review of the two companion programmes is performed first.  V11--V14
are then consolidated into one statement: recurrence and the recurrent
element, invariant and blow-up measures with the exact mean identities
and signs, and the rate-free Dirichlet ledger.  The position of the
second cycle relative to the three gates of the first cycle is recorded.
No global regularity claim is made.  The next compulsory companion audit
is V20.

% =============================================================================
\section{Compulsory V15 review of the current companion notes}
\label{sec:c2v15-companion-audit}
% =============================================================================

The current companion files in \texttt{latex\_notes} were inspected
read-only.  The frozen checkpoints are
\begin{align}
 &\texttt{Renewal\_SM\_Einstein\_Continuum\_Research\_Notes\_V88.tex},
 \notag\\[-2pt]
 &\qquad
 \texttt{60e0f9bf5bec79a39912b869f1a7ae83fe153f9b1f5fec87c1de2dbb72231116},
 \label{eq:c2v15-SM-checkpoint}\\
 &\texttt{Renewal\_Yang\_Mills\_Mass\_Gap\_Research\_Notes\_V62.tex},
 \notag\\[-2pt]
 &\qquad
 \texttt{f8a9f651e7021759adc3429f0fcaddf1e00266cb05695bf26b21b61f72ae3a73}.
 \label{eq:c2v15-YM-checkpoint}
\end{align}

\emph{SM V87--V88.}  The SM programme built a common-regulator
supertrace ledger and then showed that density-unitary conjugation of
its physical scalar Laplacian yields an exact uniformly elliptic
operator whose proper-time determinant in three dimensions needs only
volume and integrated-curvature subtractions; fixing the physical
volume removes the leading divergence, the remaining coefficient is
an integrated curvature equal to a Dirichlet integral of the conformal
factor, quadratic in the uncancelled response on its two-source branch
and vanishing through quadratic order on its cancellation line, while
the conformal coarea divergence remains; its next gate is the common
heat supertrace of all blocks.

\emph{YM V61--V62.}  The YM programme introduced a clock-factorial
reserve, then proved a centred sub-gamma clock transform, an all-order
root-moment ledger, an exact Gaussian one-centre identity, a
factorial-free time-ordered Duhamel derivative, deterministic and
averaged one-centre cards, an exact total-cumulance normal form, and
that only Gaussian pairs can accompany the centre block, isolating a
scalar clock quotient as its remaining gate.

\begin{theorem}[V15 companion-audit decision]
\label{thm:c2v15-companion-decision}
No SM V87--V88 determinant, curvature, or supertrace constant and no YM
V61--V62 transform, moment, or cumulance estimate is imported.  Neither
bears on the mean identities of Theorem~\ref{thm:c2v12-mean-identities},
the blow-up measures of Theorem~\ref{thm:c2v13-blowup-measures}, or the
rate-free ledgers.
\end{theorem}

\begin{proof}
The SM results concern heat-kernel determinants of conjugated
Laplacians on a torus with a conformal metric; the YM results concern
cumulants of a lattice clock.  None is a Navier--Stokes statement.
\end{proof}

\begin{assessment}[Direction after the audit]
\label{ass:c2v15-direction}
Both companion programmes have moved to exact ledgers (a supertrace
ledger, a root-moment ledger) in which divergences are cancelled by
constraints (fixed volume, centred transforms) and the remaining gate
is an exact scalar quantity; the present cycle's exact mean identities
against invariant measures are of the same form, with the remaining
gate a mean inequality for one cross term.  The direction is unchanged.
\end{assessment}

% =============================================================================
\section{The statistical structure, consolidated}
\label{sec:c2v15-consolidated}
% =============================================================================

\begin{theorem}[Statistical structure of a homogeneous singular point, V11--V14]
\label{thm:c2v15-consolidated}
Let \(u\) satisfy the type-I rate, \(x_0\) a homogeneous singular point,
\(\mathcal F(x_0)\) its full rescaling family with the metric \(\rho\) and
the rescaling flow \(\Theta\).  Modulo (E1)--(E9):
\begin{enumerate}[label=\textup{(\arabic*)},leftmargin=3.2em]
\item \emph{Flow.}  \((\mathcal F(x_0),\rho)\) is a compact metric space and
      \(\Theta\) a continuous flow on it, with profile action
      \(\sigma\mapsto\sigma-2\tau\).
\item \emph{Recurrence.}  Every element has rescalings arbitrarily far
      apart in log-scale that are arbitrarily close; the minimal
      discrete defect satisfies \(m_d(\ell_0)\le C_Y^2\ell_0^2\) and
      \(\liminf_{\ell_0\to\infty}m_d(\ell_0)=0\); every minimal invariant
      set consists of uniformly recurrent elements, each a local limit
      of its own rescalings along sequences of scales tending to zero
      and to infinity.
\item \emph{Invariant measures.}  \(\Theta\) has invariant Borel
      probability measures; against every such \(\mathfrak m\),
      \(\int j\,\dd\mathfrak m=-8\int\mathcal L\,\dd\mathfrak m\),
      \(\int\langle N,Y\rangle_G\dd\mathfrak m=-\int\|Y\|_G^2\dd\mathfrak m\le-\underline\delta\),
      \(\int\langle Y,N'Y\rangle_G\dd\mathfrak m=-\int(\nu\|\nabla Y\|_G^2+\frac12\|Y\|_G^2)\dd\mathfrak m\le-\frac12\underline\delta\).
\item \emph{Blow-up measures.}  The weak limits of the solution's
      empirical scale-measures over long log-windows at small scales
      form a nonempty compact set \(\mathcal B(x_0)\) of invariant
      measures on \(\mathcal F(x_0)\), each satisfying (3); the solution's
      long-window log-averages of continuous observables at small
      scales have their limit points among the corresponding means.
\item \emph{Rate-free Dirichlet ledger.}  For every smooth solution and
      centre, \(\lambda\frac{\dd}{\dd\lambda}\mathcal L_u=2\|Y_u\|_G^2+2\langle N(V_u),Y_u\rangle_G\)
      exactly; \(\mathcal L_u\ge\frac14\Phi_u\); if \(\mathcal L_u\) is unbounded
      along scales tending to zero, the log-integral of the cross term
      diverges to \(-\infty\) at least as fast as \(-\frac12\mathcal L_u\).
\end{enumerate}
\end{theorem}

\begin{proof}
(1) Lemma~\ref{lem:c2v11-flow}.  (2) Proposition~\ref{prop:c2v11-small},
Theorem~\ref{thm:c2v11-recurrence}.  (3) Theorems~\ref{thm:c2v12-existence},
\ref{thm:c2v12-mean-identities}.  (4) Theorem~\ref{thm:c2v13-blowup-measures},
Corollary~\ref{cor:c2v13-solution-statistics}.  (5) Theorem~\ref{thm:c2v14-second-rate-free},
Proposition~\ref{prop:c2v14-relation}, Theorem~\ref{thm:c2v14-necessary}.
\end{proof}

% =============================================================================
\section{Position relative to the three gates}
\label{sec:c2v15-gates}
% =============================================================================

\begin{theorem}[The gates after fifteen versions of the second cycle]
\label{thm:c2v15-gates}
Modulo (E1)--(E9), a finite-time singularity of a smooth periodic
solution is excluded if each of the following is proved.
\begin{enumerate}[label=\textup{(G\arabic*)},leftmargin=3.8em]
\item \textbf{Type I, homogeneous points.}  For some homogeneous
      singular point and some blow-up measure
      \(\mathfrak m\in\mathcal B(x_0)\), the mean inequality
      \(\int\langle N(V),\partial_\sigma V\rangle_G\dd\mathfrak m\ge-\theta\int\|\partial_\sigma V\|_G^2\dd\mathfrak m\)
      with \(\theta<1\) holds; equivalently, by
      \eqref{eq:c2v12-mean2}, \(\int\|\partial_\sigma V\|_G^2\dd\mathfrak m=0\),
      contradicting \(\underline\delta>0\).  Known: the mean identity
      holds with \(\theta=1\) exactly; the family contains no
      self-similar element; every element is uniformly non-self-similar
      with the three averaged signs.
\item[\textup{(G1\('\))}] \textbf{Type I, lacunary points.}  Exclusion of
      singular points with intermediate scales; known: satellites,
      minimal log-period, Gaussian ledger of active and intermediate
      scales (first cycle); the second cycle's dichotomy does not apply
      (no pinching).
\item \textbf{Intermediate class.}  As in the first cycle; constrained
      by the two rate-free ledgers.
\item \textbf{Strong type II.}  As in the first cycle; constrained by
      the rate-free necessary conditions of Corollary V96-necessary and
      Theorem~\ref{thm:c2v14-necessary}.
\end{enumerate}
\end{theorem}

\begin{proof}
(G1): by \eqref{eq:c2v12-mean2} the mean of the cross term equals minus
the mean squared defect, so the inequality with \(\theta<1\) forces the
mean squared defect to vanish, while it is at least \(\underline\delta>0\)
by Theorem~\ref{thm:c2v12-mean-identities}; hence the inequality is
equivalent to the impossibility of alternative (B), which with
Theorem~\ref{thm:c2v9-B} excludes homogeneous points.  The other items
are the first cycle's gates with the second cycle's constraints.
\end{proof}

\begin{assessment}[The exact form of what is missing at homogeneous points]
\label{ass:c2v15-missing}
At a homogeneous point of a type-I singularity the remaining statement
is a single inequality between two means against a blow-up measure:
the mean of the Gaussian inner product of the nonlinearity with the
self-similarity defect must be strictly less negative than minus the
mean squared defect, for at least one blow-up measure.  The cycle has
shown that equality holds exactly, for every invariant measure, and
that the strict inequality is equivalent to the vanishing of the mean
squared defect, hence to the existence of a self-similar element,
which (E9) excludes.  In other words: the exact mean identity
\eqref{eq:c2v12-mean2} is the statement that, statistically, the
nonlinearity feeds the self-similarity defect exactly at the rate at
which the Gaussian damping removes it, and the open question is
whether the Navier--Stokes nonlinearity can do this.  No global
regularity claim is made.
\end{assessment}

% =============================================================================
\section{Integrated V15 conclusion and next acquisition order}
\label{sec:c2v15-conclusion}
% =============================================================================

\begin{theorem}[What V15 establishes]
\label{thm:c2v15-conclusion}
Relative to the first cycle and V1--V14: the companion-audit decision
(Theorem~\ref{thm:c2v15-companion-decision}), the consolidated
statistical structure (Theorem~\ref{thm:c2v15-consolidated}), and the
gates after fifteen versions (Theorem~\ref{thm:c2v15-gates}).  The full
Navier--Stokes stability/global-regularity theorem remains unproved.
\end{theorem}

\begin{proof}
The cited results.
\end{proof}

\begin{programme}[V16 acquisition order]
\label{prog:c2v16}
\begin{enumerate}[label=\textup{(AC\arabic*)},leftmargin=4.8em]
\item The centre mode.  Take the Gaussian mean of the defect equation
      \eqref{eq:c2v7-defect-equation}: derive the ordinary differential
      equation for \(\bar Y(\sigma)=\int YG/\int G\) and its forcing by the
      nonlinearity, and determine whether the mean squared defect
      \(\int\|Y\|_G^2\dd\mathfrak m\) can be split into the centre mode
      \(\int|\bar Y|^2\dd\mathfrak m\cdot\int G\) and the rest, with a
      separate identity for each.
\item Centre normalization.  Determine whether the blow-up centre can
      be normalized (a choice of \(x_0\) at each scale minimizing the
      centre mode) so that the centre mode vanishes on the family, and
      whether homogeneity is preserved by this normalization.
\item The next compulsory companion audit is V20.
\end{enumerate}
\end{programme}

\begin{assessment}[Position after V15]
\label{ass:c2v15-position}
The second cycle is consolidated on the statistical structure of
homogeneous singular points, and the missing statement is a single mean
inequality.  No full stability or global-regularity claim is made.
\end{assessment}

% =============================================================================
\section*{Version V15 (second cycle) append-only record}
\addcontentsline{toc}{section}{Version V15 (second cycle) append-only record}
% =============================================================================

The complete V14 source of the second cycle was retained before this
continuation.  Its SHA--256 digest was
\begin{center}\scriptsize\ttfamily
31eab9b160696c8299f771ca275c09e3de1115f62f865cdfa48a81a1fbbdafff.
\end{center}
V15 performed the compulsory companion audit, consolidated V11--V14, and
recorded the position relative to the three gates.  No full regularity
claim is made.

% =============================================================================
\section{V16: the centre mode of the defect and centre normalization}
\label{sec:c2v16-entry}
% =============================================================================

Items (AC1) and (AC2) of Programme~\ref{prog:c2v16} are carried out.  The
Gaussian mean of the self-similarity defect, the centre mode, satisfies
an exact ordinary differential equation in the log-scale: the linear
part damps it at rate one half, and the forcing is the Gaussian mean of
the linearized nonlinearity, expressed through the Gaussian dipole
moments of the profile and of the pressure derivative.  The mean squared
defect splits exactly into the centre mode and the fluctuation, each
with its own mean identity against every invariant measure.  Centre
normalization by a moving blow-up centre is shown to remove the centre
mode at a given scale exactly when the Gaussian dipole matrix of the
profile is invertible, and to preserve homogeneity when the
normalizing shift is bounded in the parabolic scale; neither condition
is established in general.  No global regularity claim is made.  The
next compulsory companion audit is V20.

% =============================================================================
\section{The centre-mode equation}
\label{sec:c2v16-centre}
% =============================================================================

For \(W\in\mathcal F(x_0)\) with profile \(V\) and defect \(Y=\partial_\sigma V\),
let \(|G|:=\int_{\R^3}G=(4\pi\nu)^{3/2}\) and define the Gaussian means and
dipole moments
\begin{equation}
 \bar Y:=\frac1{|G|}\int YG,\qquad
 \bar V:=\frac1{|G|}\int VG,\qquad
 M_V:=\frac1{2\nu|G|}\int z\otimes V\,G,\qquad
 q_{P'}:=\frac1{2\nu|G|}\int P'\,z\,G ,
 \label{eq:c2v16-means}
\end{equation}
with \(P'=\partial_\sigma P_V\) as in \eqref{eq:c2v7-Pprime}.

\begin{proposition}[The centre-mode equation]
\label{prop:c2v16-centre-equation}
For a.e.\ \(\sigma\),
\begin{equation}
 \boxed{\quad
 \frac{\dd}{\dd\sigma}\bar Y=-\frac12\bar Y-\frac1{2\nu|G|}\int\bigl[(z\cdot V)\,Y+(z\cdot Y)\,V\bigr]G-q_{P'} .
 \quad}
 \label{eq:c2v16-centre-equation}
\end{equation}
In particular, if \(N'(V)Y\) had zero Gaussian mean, \(\bar Y\) would
decay like \(e^{-\sigma/2}\); the forcing is the Gaussian mean of the
linearized nonlinearity, \(-\overline{N'(V)Y}\), with
\(\overline{(V\cdot\nabla)Y}=\frac1{2\nu|G|}\int(z\cdot V)YG\),
\(\overline{(Y\cdot\nabla)V}=\frac1{2\nu|G|}\int(z\cdot Y)VG\),
\(\overline{\nabla P'}=q_{P'}\).
\end{proposition}

\begin{proof}
Take the Gaussian mean of \eqref{eq:c2v7-defect-equation}.  For the
linear part, \(\int\nu\Delta Y\,G=\nu\int Y\Delta G\) and
\(-\frac12\int(z\cdot\nabla Y)G=\frac12\int Y\operatorname{div}(zG)\), with
\(\Delta G=\bigl(\frac{|z|^2}{4\nu^2}-\frac3{2\nu}\bigr)G\) and
\(\operatorname{div}(zG)=\bigl(3-\frac{|z|^2}{2\nu}\bigr)G\); the two
contributions cancel exactly, leaving \(\overline{L_-Y}=-\frac12\bar Y\).
For the nonlinear part, integrate by parts using \(\operatorname{div}V=\operatorname{div}Y=0\)
and \(\nabla G=-\frac z{2\nu}G\):
\(\int(V\cdot\nabla)Y\,G=-\int Y(V\cdot\nabla G)=\frac1{2\nu}\int(z\cdot V)YG\),
similarly for \((Y\cdot\nabla)V\), and \(\int\nabla P'G=-\int P'\nabla G=\frac1{2\nu}\int P'zG\).
All integrals converge by Proposition~\ref{prop:c2v7-defect-equation}.
\end{proof}

\begin{theorem}[Split of the mean squared defect]
\label{thm:c2v16-split}
For every \(\Theta\)-invariant probability measure \(\mathfrak m\) on
\(\mathcal F(x_0)\), with \(\mathfrak y_0:=|G|\,|\bar Y|^2\) and
\(\mathfrak y_\perp:=\|Y-\bar Y\|_G^2\) (so \(\mathfrak y=\mathfrak y_0+\mathfrak y_\perp\)),
\begin{align}
 \int\mathfrak y_0\,\dd\mathfrak m
 &=-2|G|\int\bar Y\cdot\overline{N'(V)Y}\,\dd\mathfrak m ,
 \label{eq:c2v16-centre-mean}\\
 \int\Bigl(\nu\|\nabla Y\|_G^2+\tfrac12\mathfrak y_\perp\Bigr)\dd\mathfrak m
 &=-\int\Bigl(\langle Y,N'(V)Y\rangle_G-|G|\,\bar Y\cdot\overline{N'(V)Y}\Bigr)\dd\mathfrak m
 =-\int\langle Y-\bar Y,N'(V)Y\rangle_G\,\dd\mathfrak m ,
 \label{eq:c2v16-fluctuation-mean}
\end{align}
and \(\int\mathfrak y_\perp\dd\mathfrak m\le2\int\nu\|\nabla Y\|_G^2\dd\mathfrak m\) by the
Gaussian Poincar\'e inequality.  In alternative (B),
\(\int(\mathfrak y_0+\mathfrak y_\perp)\dd\mathfrak m\ge\underline\delta\), and at
least one of \(\int\mathfrak y_0\dd\mathfrak m\ge\frac12\underline\delta\),
\(\int\mathfrak y_\perp\dd\mathfrak m\ge\frac12\underline\delta\) holds for each
\(\mathfrak m\).
\end{theorem}

\begin{proof}
From \eqref{eq:c2v16-centre-equation}, \(\frac12\frac{\dd}{\dd\sigma}|\bar Y|^2=-\frac12|\bar Y|^2-\bar Y\cdot\overline{N'(V)Y}\);
the invariant-measure argument of Theorem~\ref{thm:c2v12-mean-identities}
(with the bounded observable \(|\bar Y|^2\le\mathfrak y/|G|\)) gives
\eqref{eq:c2v16-centre-mean}.  Subtracting \(|G|\) times this from
\eqref{eq:c2v12-mean3} and using \(\langle Y,N'Y\rangle_G-|G|\bar Y\cdot\overline{N'Y}=\langle Y-\bar Y,N'Y\rangle_G\)
gives \eqref{eq:c2v16-fluctuation-mean}.  The rest is
Proposition~\ref{prop:c2v8-spectrum}(iii) and
Theorem~\ref{thm:c2v12-mean-identities}.
\end{proof}

\begin{assessment}[Two ways to be non-self-similar]
\label{ass:c2v16-two-ways}
The defect of a uniformly non-self-similar element is carried either by
its centre mode, a drift of the blow-up centre in log-scale that the
Gaussian mean of the linearized nonlinearity must sustain against the
damping \(\frac12\), or by its fluctuation, a genuine change of profile
shape that the fluctuation part of the linearized nonlinearity must
sustain against the damping \(1\) of the Poincar\'e inequality.  Both
are exact mean identities against every invariant measure; for each
measure at least one of the two carries half the uniform defect rate.
\end{assessment}

% =============================================================================
\section{Centre normalization}
\label{sec:c2v16-normalization}
% =============================================================================

Let \(b:\R\to\R^3\) be smooth and consider the recentred profile
\(V_b(z,\sigma):=V(z+b(\sigma),\sigma)\), the profile of the same solution
with the blow-up centre moved by \(-\lambda b(\sigma)\) at scale \(\lambda\).
Its defect is \(\partial_\sigma V_b=Y(\cdot+b,\sigma)+b'(\sigma)\cdot\nabla V(\cdot+b,\sigma)\).

\begin{proposition}[Centre normalization]
\label{prop:c2v16-normalization}
\begin{enumerate}[label=\textup{(\roman*)},leftmargin=3.2em]
\item The Gaussian mean of the recentred defect at \(\sigma\) is
      \(\overline{\partial_\sigma V_b}=\overline{Y(\cdot+b)}+\overline{b'\cdot\nabla V(\cdot+b)}\),
      and \(\overline{b'\cdot\nabla V(\cdot+b)}=\frac1{2\nu|G|}\int(z\cdot b')V(z+b)G=M_{V(\cdot+b)}^{\mathsf T}b'\)
      with the dipole matrix of \eqref{eq:c2v16-means} computed for the
      shifted profile.  Hence the centre mode of the recentred profile
      vanishes at \(\sigma\) if and only if
      \(M_{V(\cdot+b)}^{\mathsf T}b'(\sigma)=-\overline{Y(\cdot+b)}\), which has a
      solution \(b'(\sigma)\) whenever the dipole matrix is invertible, and
      is a linear ordinary differential equation for \(b\) along the
      orbit.
\item If \(|b(\sigma)|\le B\) for all \(\sigma\), the recentred family is the
      family of the moving centre \(x_0-\lambda b\), homogeneity of the
      point in the sense of the first cycle is preserved (the terminal
      density at radius \(1+B\) about the moving centre dominates that at
      radius \(1\) about \(x_0\)), and all results of the cycle apply to
      the recentred family.
\end{enumerate}
\end{proposition}

\begin{proof}
(i) The chain rule and \(\int(b'\cdot\nabla V)G=-\int V(b'\cdot\nabla G)=\frac1{2\nu}\int(z\cdot b')VG\).
(ii) \(B_{\rho}(x_0)\subset B_{(1+B)\rho}(x_0-\rho b)\) gives the density
statement; the moving-centre compactness of the first cycle (Lemma
V79-moving-centres, Proposition V98-moving-lower) applies to bounded
shifts.
\end{proof}

\begin{firewall}[Item (AC2): what is not established]
\label{fw:c2v16-AC2}
Neither the invertibility of the Gaussian dipole matrix along orbits,
with an inverse bounded uniformly on the family, nor the boundedness of
the resulting shift \(b\) is established; the dipole matrix of a profile
with a reflection symmetry vanishes, and a divergent shift would mean
that the blow-up centre drifts by more than the parabolic scale in
log-scale, which is the compact-satellite question of the first cycle
in another form.  The centre mode therefore remains part of the defect
in general, and the split of Theorem~\ref{thm:c2v16-split} is the
statement available.
\end{firewall}

% =============================================================================
\section{Integrated V16 conclusion and next acquisition order}
\label{sec:c2v16-conclusion}
% =============================================================================

\begin{theorem}[What V16 proves]
\label{thm:c2v16-conclusion}
Relative to the first cycle and V1--V15, modulo (E1)--(E9): the
centre-mode equation \eqref{eq:c2v16-centre-equation}, the split of the
mean squared defect with its two mean identities
(Theorem~\ref{thm:c2v16-split}), and the centre-normalization
proposition.  The full Navier--Stokes stability/global-regularity
theorem remains unproved.
\end{theorem}

\begin{proof}
The cited results.
\end{proof}

\begin{programme}[V17 acquisition order]
\label{prog:c2v17}
\begin{enumerate}[label=\textup{(AC\arabic*)},leftmargin=4.8em]
\item Dipole moments and the flux.  Express the scale-invariant flux
      \(j=4\langle V,N(V)\rangle_G=\frac1\nu\int(|V|^2+2P_V)(z\cdot V)G\) and
      the centre-mode forcing through the Gaussian moments of the
      profile, and derive the exact identity for \(\frac{\dd}{\dd\sigma}\bar V\)
      (the mean of the profile equation), \(\frac{\dd}{\dd\sigma}\bar V=-\frac12\bar V-\overline{N(V)}\),
      with \(\overline{N(V)}=\frac1{2\nu|G|}\int\bigl[(z\cdot V)V+P_Vz\bigr]G\);
      record the mean identity \(\int\bar V\dd\mathfrak m=-2\int\overline{N(V)}\dd\mathfrak m\).
\item Reflection-symmetric families.  Determine what the identities give
      for elements with \(V(-z)=-V(z)\) (odd profiles, for which all
      Gaussian means and dipole matrices vanish): the centre mode is
      zero, the defect is pure fluctuation, and the fluctuation identity
      is the whole of the second sign; record whether oddness is
      preserved by the flow and by local limits.
\item The next compulsory companion audit is V20.
\end{enumerate}
\end{programme}

\begin{assessment}[Position after V16]
\label{ass:c2v16-position}
The defect of a uniformly non-self-similar element splits into a centre
drift and a shape fluctuation, each governed by an exact mean identity.
No full stability or global-regularity claim is made.
\end{assessment}

% =============================================================================
\section*{Version V16 (second cycle) append-only record}
\addcontentsline{toc}{section}{Version V16 (second cycle) append-only record}
% =============================================================================

The complete V15 source of the second cycle was retained before this
continuation.  Its SHA--256 digest was
\begin{center}\scriptsize\ttfamily
691bce8817e34f0508c084336c6c7d4cb412801577489878dc9adbf842ebadc3.
\end{center}
V16 derived the centre-mode equation, split the mean squared defect
into centre and fluctuation parts with their mean identities, and
recorded the centre-normalization conditions.  No full regularity claim
is made.

% =============================================================================
\section{V17: Gaussian moments of the profile, the mean momentum identity,
and odd profiles}
\label{sec:c2v17-entry}
% =============================================================================

Items (AC1) and (AC2) of Programme~\ref{prog:c2v17} are carried out.  The
Gaussian mean of the profile, the scale-invariant Gaussian momentum,
satisfies an exact equation in log-scale with damping one half and
forcing equal to the Gaussian momentum flux, and its mean against every
invariant measure equals minus twice the mean momentum flux; the
relation with the Gaussian-weighted momentum and the dipole identity of
the first cycle is recorded.  For odd profiles all Gaussian means
vanish, oddness is preserved by the profile equation, the flow, and
local limits, the centre mode of the defect is zero, and the
fluctuation identity is the whole of the second sign.  A sentence of
V16 asserting that the dipole matrix of a reflection-symmetric profile
vanishes is corrected: for odd profiles the dipole matrix is even and
need not vanish; what vanishes is the mean.  No global regularity claim
is made.  The next compulsory companion audit is V20.

% =============================================================================
\section{The mean momentum identity}
\label{sec:c2v17-momentum}
% =============================================================================

\begin{proposition}[Mean of the profile equation]
\label{prop:c2v17-mean-profile}
For \(W\in\mathcal A_\nu(C)\) with the Morrey bound, with \(\bar V\) the
Gaussian mean of \eqref{eq:c2v16-means} and
\begin{equation}
 \overline{N(V)}:=\frac1{|G|}\int N(V)G=\frac1{2\nu|G|}\int\bigl[(z\cdot V)\,V+P_V\,z\bigr]G ,
 \label{eq:c2v17-Nbar}
\end{equation}
one has, for every \(\sigma\),
\begin{equation}
 \boxed{\quad
 \frac{\dd}{\dd\sigma}\bar V=-\frac12\bar V-\overline{N(V)} ,
 \quad}
 \label{eq:c2v17-mean-equation}
\end{equation}
and for every \(\Theta\)-invariant probability measure \(\mathfrak m\) on
a rescaling family at a homogeneous point,
\(\int\bar V\,\dd\mathfrak m=-2\int\overline{N(V)}\,\dd\mathfrak m\).
In the original variables, \(|G|\,\bar V(\sigma)=\lambda^{-2}\int_{\R^3}W(-\lambda^2,y)\,\psi(y,-\lambda^2)\,\dd y\)
at \(\lambda=e^{-\sigma/2}\): the Gaussian mean of the profile is the
Gaussian-weighted momentum of \(W\) at the parabolic scale, in
scale-invariant units, and \(2\nu\lambda^2|G|\bar V\) is the dipole moment of
the weighted radial momentum of the first cycle (Proposition
V89-moments).
\end{proposition}

\begin{proof}
The Gaussian mean of \eqref{eq:c2v2-profile-equation}: as in
Proposition~\ref{prop:c2v16-centre-equation}, \(\overline{L_-V}=-\frac12\bar V\)
(the Laplacian and drift contributions cancel), and
\(\int(V\cdot\nabla)V\,G=-\int V(V\cdot\nabla G)=\frac1{2\nu}\int(z\cdot V)VG\),
\(\int\nabla P_VG=\frac1{2\nu}\int P_VzG\).  The mean identity follows
from the invariant-measure argument with the bounded observable
\(\bar V\) (\(|\bar V|\le|G|^{-1/2}\|V\|_G\le|G|^{-1/2}C_\Phi^{1/2}\)).  The
change of variables: \(\int W(-\lambda^2,y)\psi\,\dd y=\lambda^{-1}\lambda^3\int VG\,\dd z=\lambda^2|G|\bar V\);
the dipole statement is Proposition V89-moments of the first cycle,
\(\int y\,(y\cdot W)\psi=2\nu\lambda^2\int\psi W\), transported.
\end{proof}

\begin{remark}[Momentum and flux]
\label{rem:c2v17-momentum-flux}
The scale-invariant flux term is
\(j=\frac1\nu\int(|V|^2+2P_V)(z\cdot V)G=4\langle V,N(V)\rangle_G\), while
the momentum flux \(\overline{N(V)}\) is the vector
\(\frac1{2\nu|G|}\int[(z\cdot V)V+P_Vz]G\); the flux term is the projection
of the momentum-flux density on \(V\) (before averaging), the momentum
flux its Gaussian average.  Both are cubic or pressure-quadratic
moments of the profile against the same weight, and neither has a sign.
\end{remark}

% =============================================================================
\section{Odd profiles}
\label{sec:c2v17-odd}
% =============================================================================

\begin{proposition}[Odd profiles]
\label{prop:c2v17-odd}
Call \(W\) \emph{odd} if \(W(s,-y)=-W(s,y)\) for all \(s<0\), \(y\in\R^3\),
equivalently \(V(-z,\sigma)=-V(z,\sigma)\).
\begin{enumerate}[label=\textup{(\roman*)},leftmargin=3.2em]
\item Oddness is preserved by the Navier--Stokes evolution (if \(W\) is
      odd at one time it is odd at all times, by uniqueness in the
      class), by the rescaling flow, by local limits, and by the parity
      map \(W\mapsto-W(\cdot,-\cdot)\) of the first cycle, which fixes odd
      elements.  If a rescaling family contains an odd element, the odd
      elements form a closed \(\Theta\)-invariant subfamily.
\item For an odd \(W\): \(\bar V=0\), \(\bar Y=0\), \(\overline{N(V)}=0\),
      \(\overline{N'(V)Y}=0\) for all \(\sigma\); the pressure \(P_V\) is
      even; the defect is pure fluctuation, \(\mathfrak y=\mathfrak y_\perp\);
      and the fluctuation identity \eqref{eq:c2v16-fluctuation-mean}
      coincides with \eqref{eq:c2v12-mean3}.  The dipole matrix
      \(M_V=\frac1{2\nu|G|}\int z\otimes VG\) is in general nonzero, its
      integrand being even.
\item In alternative (B) restricted to an odd invariant subfamily,
      \(\int\mathfrak y_\perp\dd\mathfrak m\ge\underline\delta\) for every
      invariant measure on it, and the whole of the second sign is
      carried by the shape fluctuation:
      \(\int\langle Y,N'(V)Y\rangle_G\dd\mathfrak m=-\int(\nu\|\nabla Y\|_G^2+\frac12\|Y\|_G^2)\dd\mathfrak m\le-\int\mathfrak y\,\dd\mathfrak m\le-\underline\delta\),
      using the Poincar\'e inequality with \(\bar Y=0\).
\end{enumerate}
\end{proposition}

\begin{proof}
(i) If \(W\) is odd at time \(s_0\), then \(\widetilde W(s,y):=-W(s,-y)\)
solves the equations with the same datum at \(s_0\), and the class
bounds place both in \(C_sL^6\) with bounded gradients on half-lines;
forward uniqueness (E8) gives \(\widetilde W=W\) for \(s\ge s_0\) and
backward uniqueness (E6) for the difference gives it for \(s\le s_0\).
Rescaling and local limits commute with \(y\mapsto-y\).  (ii) Odd
functions have zero Gaussian mean; \(N(V)\) is odd (\((V\cdot\nabla)V\) is
odd, \(P_V=\partial_i\partial_j(-\Delta)^{-1}(V_iV_j)\) is even, \(\nabla P_V\)
odd), \(Y=L_-V-N(V)\) is odd, \(N'(V)Y\) is odd; \(z\otimes V\) is even.
(iii) \(\bar Y=0\) gives \(\mathfrak y=\mathfrak y_\perp\le2\nu\|\nabla Y\|_G^2\) by
Proposition~\ref{prop:c2v8-spectrum}(iii), so
\(\nu\|\nabla Y\|_G^2+\frac12\|Y\|_G^2\ge\|Y\|_G^2\), and
\eqref{eq:c2v12-mean3} with \(\int\mathfrak y\,\dd\mathfrak m\ge\underline\delta\).
\end{proof}

\begin{theorem}[Correction to Firewall~\ref{fw:c2v16-AC2}]
\label{thm:c2v17-correction}
The sentence of Firewall~\ref{fw:c2v16-AC2} stating that ``the dipole
matrix of a profile with a reflection symmetry vanishes'' is withdrawn.
For an odd profile the dipole matrix \(M_V\) has an even integrand and
does not vanish in general; for a profile symmetric under a reflection
\(R\) through a plane, \(V(Rz)=RV(z)\), the dipole matrix satisfies
\(M_V=RM_VR\) and again need not vanish.  What vanishes for odd profiles
is the Gaussian mean \(\bar V\) and the centre mode \(\bar Y\)
(Proposition~\ref{prop:c2v17-odd}).  No other statement of V16 used the
withdrawn sentence; the conclusion of the firewall, that centre
normalization is not established in general, stands, with the reason
that the invertibility of \(M_V\) along orbits and the boundedness of
the shift are not proved.
\end{theorem}

\begin{proof}
Parity of \(z\otimes V\) for odd \(V\); for \(V(Rz)=RV(z)\), the change of
variables \(z\mapsto Rz\) in \(\int z\otimes V\,G\) gives \(R(\int z\otimes VG)R\)
using \(R^{\mathsf T}=R\) and \(G(Rz)=G(z)\).
\end{proof}

\begin{assessment}[Odd blow-up]
\label{ass:c2v17-odd}
For an odd blow-up, the centre is fixed at the origin by symmetry, the
Gaussian momentum vanishes at every scale, and the defect is purely a
change of shape.  The statistical statement of alternative (B) is then
the single mean identity: the fluctuation part of the linearized
nonlinearity must feed the shape defect exactly at the rate of its
Gaussian damping, which is at least \(1\) in the Poincar\'e sense.
Whether an odd homogeneous singularity exists is open; the odd case is
the one in which the centre-mode ambiguity of V16 is absent.  No global
regularity claim is made.
\end{assessment}

% =============================================================================
\section{Integrated V17 conclusion and next acquisition order}
\label{sec:c2v17-conclusion}
% =============================================================================

\begin{theorem}[What V17 proves]
\label{thm:c2v17-conclusion}
Relative to the first cycle and V1--V16, modulo (E1)--(E9): the mean
momentum identity \eqref{eq:c2v17-mean-equation} with its
invariant-measure form, the odd-profile proposition, and the correction
of Theorem~\ref{thm:c2v17-correction}.  The full Navier--Stokes
stability/global-regularity theorem remains unproved.
\end{theorem}

\begin{proof}
The cited results.
\end{proof}

\begin{programme}[V18 acquisition order]
\label{prog:c2v18}
\begin{enumerate}[label=\textup{(AC\arabic*)},leftmargin=4.8em]
\item Second moments.  Derive the exact equation for the Gaussian
      second-moment tensor \(\int V\otimes V\,G\) and for the Gaussian
      inertia scalar \(\int|z|^2|V|^2G\) along \(\sigma\), and their mean
      identities; determine whether the inertia identity, whose linear
      part is a virial-type balance, gives a signed constraint on the
      Gaussian dissipation.
\item Test on the columnar swirl of the first cycle (an exact solution
      outside the class) to check the signs of each term of the
      inertia identity.
\item The next compulsory companion audit is V20.
\end{enumerate}
\end{programme}

\begin{assessment}[Position after V17]
\label{ass:c2v17-position}
The Gaussian momentum of the profile obeys an exact mean identity, odd
profiles carry the defect entirely in their shape, and a sentence of
V16 is corrected.  No full stability or global-regularity claim is
made.
\end{assessment}

% =============================================================================
\section*{Version V17 (second cycle) append-only record}
\addcontentsline{toc}{section}{Version V17 (second cycle) append-only record}
% =============================================================================

The complete V16 source of the second cycle was retained before this
continuation.  Its SHA--256 digest was
\begin{center}\scriptsize\ttfamily
98df210791b942db964aa1d36419fadf79f20ea86f3b961e1dd513081f85461d.
\end{center}
V17 proved the mean momentum identity, the odd-profile proposition,
and corrected a sentence of V16.  No full regularity claim is made.

% =============================================================================
\section{V18: the Gaussian stress tensor, the inertia identity, and the
signed mean stress}
\label{sec:c2v18-entry}
% =============================================================================

Items (AC1) and (AC2) of Programme~\ref{prog:c2v18} are carried out.  The
Gaussian second-moment tensor of the profile obeys an exact tensor
identity in log-scale whose trace is the first identity; against every
invariant measure its mean form says that the mean Gaussian nonlinear
stress is negative semidefinite, bounded above by minus the mean
second-moment tensor: a directional refinement of the inward flux.  The
Gaussian inertia scalar obeys an exact identity of virial type whose
mean form determines the mean of the inertia-weighted radial flux; it
has a sign only conditionally.  The columnar swirl, on which the
nonlinearity vanishes identically, confirms the vanishing of every
nonlinear term of both identities.  No global regularity claim is made.
The next compulsory companion audit is V20.

% =============================================================================
\section{The tensor identity and the signed mean stress}
\label{sec:c2v18-tensor}
% =============================================================================

For \(W\in\mathcal A_\nu(C)\) with the Morrey bound and profile \(V\),
define the Gaussian second-moment tensor, the Gaussian gradient tensor,
and the Gaussian nonlinear stress,
\begin{equation}
 T_{ij}:=\int V_iV_j\,G,\qquad
 D_{ij}:=2\nu\int\partial_kV_i\,\partial_kV_j\,G,\qquad
 S_{ij}:=\int\bigl(N_iV_j+V_iN_j\bigr)G,\qquad N=N(V).
 \label{eq:c2v18-tensors}
\end{equation}
\(T\) and \(D\) are symmetric positive semidefinite, \(S\) is symmetric,
\(\operatorname{tr}T=\varphi\), \(\operatorname{tr}D=\frac12d\),
\(\operatorname{tr}S=2\langle V,N\rangle_G=\frac12j\).

\begin{theorem}[Tensor identity and signed mean stress]
\label{thm:c2v18-tensor}
For every \(\sigma\),
\begin{equation}
 \boxed{\quad
 \frac{\dd}{\dd\sigma}T=-D-T-S ,
 \quad}
 \label{eq:c2v18-tensor-identity}
\end{equation}
whose trace is the first identity in the form
\(\frac{\dd}{\dd\sigma}\varphi=-\frac12d-\varphi-\frac12j\).  For every
\(\Theta\)-invariant probability measure \(\mathfrak m\) on a rescaling
family at a homogeneous point,
\begin{equation}
 \boxed{\quad
 \int S\,\dd\mathfrak m=-\int(D+T)\,\dd\mathfrak m\ \preceq\ -\int T\,\dd\mathfrak m\ \preceq\ 0 ,
 \quad}
 \label{eq:c2v18-mean-stress}
\end{equation}
in the sense of symmetric matrices: for every unit vector \(e\),
\(\int\!\!\int2(e\cdot N)(e\cdot V)G\,\dd\mathfrak m=-\int\bigl(2\nu\|e\cdot\nabla V\|_G^2+\|e\cdot V\|_G^2\bigr)\dd\mathfrak m\le-\int\|e\cdot V\|_G^2\dd\mathfrak m\),
and \(\int T\,\dd\mathfrak m\) has trace at least \(c_\Phi\).
\end{theorem}

\begin{proof}
\(\frac{\dd}{\dd\sigma}T_{ij}=\int(\partial_\sigma V_iV_j+V_i\partial_\sigma V_j)G\) with
\(\partial_\sigma V=L_-V-N\).  By the \(G\)-symmetry,
\(\int(L_-V)_iV_jG=-\nu\int\partial_kV_i\partial_kV_jG-\frac12T_{ij}\), so the
linear part contributes \(-D_{ij}-T_{ij}\) after symmetrization, and the
nonlinear part \(-S_{ij}\).  The mean form is the invariant-measure
argument of Theorem~\ref{thm:c2v12-mean-identities} applied to each
bounded entry \(T_{ij}\).  Positivity of \(D\) and \(T\) gives the
matrix inequalities; the trace bound is the pinching.
\end{proof}

\begin{assessment}[Directional inward flux]
\label{ass:c2v18-directional}
The first cycle's inward flux, made exact in V12 as
\(\int j\,\dd\mathfrak m=-\int(d+2\varphi)\dd\mathfrak m\), is the trace of
\eqref{eq:c2v18-mean-stress}.  The tensor form says more: in every
direction \(e\), the mean Gaussian correlation of the nonlinearity with
the velocity component along \(e\) is negative, and at least as negative
as minus the mean Gaussian energy of that component.  The Gaussian
nonlinear stress is thus, statistically, a negative definite operator
dominated by minus the energy tensor; no direction is exempt.  Since
\(S\) is the Gaussian moment of \(N\otimes V+V\otimes N\) with
\(N=\Omega\times V+\nabla\Pi\), the identity constrains the Gaussian
correlation of the Lamb vector with the velocity and the Gaussian
Bernoulli stress together; it does not separate them.
\end{assessment}

% =============================================================================
\section{The inertia identity}
\label{sec:c2v18-inertia}
% =============================================================================

Define the Gaussian inertia, the inertia-weighted gradient, and the
inertia-weighted flux,
\begin{equation}
 I:=\int|z|^2|V|^2G,\qquad
 I_\nabla:=\int|z|^2|\nabla V|^2G,\qquad
 j_2:=\frac1\nu\int|z|^2\bigl(|V|^2+2P_V\bigr)(z\cdot V)\,G .
 \label{eq:c2v18-inertia}
\end{equation}

\begin{theorem}[Inertia identity]
\label{thm:c2v18-inertia}
For every \(\sigma\),
\begin{equation}
 \boxed{\quad
 \frac{\dd}{\dd\sigma}I=-2\nu I_\nabla+6\nu\varphi-2I+2\nu j-\frac12j_2 ,
 \quad}
 \label{eq:c2v18-inertia-identity}
\end{equation}
and for every invariant probability measure \(\mathfrak m\) on a
rescaling family at a homogeneous point,
\begin{equation}
 \frac12\int j_2\,\dd\mathfrak m=\int\bigl(2\nu\varphi-2\nu d-2I-2\nu I_\nabla\bigr)\dd\mathfrak m
 \ \le\ 2\nu\int\varphi\,\dd\mathfrak m-2\int I\,\dd\mathfrak m .
 \label{eq:c2v18-inertia-mean}
\end{equation}
In particular the mean inertia-weighted flux is inward whenever
\(\int I\,\dd\mathfrak m\ge\nu\int\varphi\,\dd\mathfrak m\).
\end{theorem}

\begin{proof}
\(\frac{\dd}{\dd\sigma}I=2\int|z|^2V\cdot\partial_\sigma VG=2\int|z|^2V\cdot L_-VG-2\int|z|^2V\cdot NG\).
Linear part: with \(L_-=\nu G^{-1}\operatorname{div}(G\nabla\cdot)-\frac12\),
\(\nu\int|z|^2V\cdot\operatorname{div}(G\nabla V)=-\nu\int\bigl(2z_kV_i+|z|^2\partial_kV_i\bigr)\partial_kV_iG
=-\nu I_\nabla-\nu\int z\cdot\nabla|V|^2G=-\nu I_\nabla+\nu\int|V|^2\operatorname{div}(zG)
=-\nu I_\nabla+3\nu\varphi-\frac12I\),
using \(\operatorname{div}(zG)=(3-\frac{|z|^2}{2\nu})G\); adding \(-\frac12I\)
from the zero-order term gives \(2\int|z|^2V\cdot L_-VG=-2\nu I_\nabla+6\nu\varphi-2I\).
Nonlinear part: \(\int|z|^2V\cdot(V\cdot\nabla)VG=\frac12\int|z|^2V\cdot\nabla|V|^2G
=-\frac12\int|V|^2\operatorname{div}(|z|^2VG)=-\int|V|^2(z\cdot V)G+\frac1{4\nu}\int|z|^2|V|^2(z\cdot V)G\)
and \(\int|z|^2V\cdot\nabla P_VG=-\int P_V\operatorname{div}(|z|^2VG)=-2\int P_V(z\cdot V)G+\frac1{2\nu}\int|z|^2P_V(z\cdot V)G\),
so \(-2\int|z|^2V\cdot NG=2\int(|V|^2+2P_V)(z\cdot V)G-\frac1{2\nu}\int|z|^2(|V|^2+2P_V)(z\cdot V)G=2\nu j-\frac12j_2\).
All integrals converge by the Gaussian factor against the polynomial
weights and the bounds of the class.  The mean form: \(I\) is bounded on
the family (\(I\le\||z|V\|_G^2\le C(\nu)(\|V\|_G+\|\nabla V\|_G)^2\) by the
Gaussian Hardy inequality of Theorem~\ref{thm:c2v7-defect-energy}), so
its mean derivative vanishes; insert \(\int j\,\dd\mathfrak m=-\int(d+2\varphi)\dd\mathfrak m\)
and \(I_\nabla\ge0\), \(d\ge0\).
\end{proof}

\begin{firewall}[Item (AC1): no unconditional sign]
\label{fw:c2v18-AC1}
The mean inertia identity determines the mean of the inertia-weighted
flux exactly, as \(2\nu\langle\varphi\rangle-2\langle I\rangle-2\nu\langle d\rangle-2\nu\langle I_\nabla\rangle\);
the sign depends on whether the Gaussian inertia exceeds \(\nu\) times
the Gaussian energy in the mean, that is, on whether the profile's
energy sits mostly outside the sphere \(|z|=\sqrt\nu\) in the Gaussian
sense.  Neither case is excluded.  The identity gives no signed
constraint on the Gaussian dissipation: \(d\) and \(I_\nabla\) appear
with the same sign and are balanced by \(2\nu\varphi\) and the flux terms.
\end{firewall}

\begin{proposition}[Item (AC2): the columnar swirl]
\label{prop:c2v18-swirl}
For the columnar swirl \(W=w(s,r)e_\theta\) of the first cycle
(Proposition V88-examples), \(N(V)\equiv0\), \(z\cdot V\equiv0\), hence
\(S\equiv0\), \(j\equiv0\), \(j_2\equiv0\), and the tensor and inertia
identities reduce to their linear parts,
\(\frac{\dd}{\dd\sigma}T=-D-T\), \(\frac{\dd}{\dd\sigma}I=-2\nu I_\nabla+6\nu\varphi-2I\),
in which the signs of the individual terms are as displayed; this
flow is not in the class (infinite Dirichlet integral) and serves only
as a check of the vanishing of the nonlinear terms.
\end{proposition}

\begin{proof}
\((V\cdot\nabla)V=-(w^2/r)e_r=-\nabla\int_0^rw^2\rho^{-1}\dd\rho\) is a gradient,
cancelled by the pressure, and \(e_\theta\perp z\).
\end{proof}

% =============================================================================
\section{Integrated V18 conclusion and next acquisition order}
\label{sec:c2v18-conclusion}
% =============================================================================

\begin{theorem}[What V18 proves]
\label{thm:c2v18-conclusion}
Relative to the first cycle and V1--V17, modulo (E1)--(E9): the tensor
identity \eqref{eq:c2v18-tensor-identity} with the signed mean stress
\eqref{eq:c2v18-mean-stress}, and the inertia identity
\eqref{eq:c2v18-inertia-identity} with its mean form
\eqref{eq:c2v18-inertia-mean}.  The full Navier--Stokes
stability/global-regularity theorem remains unproved.
\end{theorem}

\begin{proof}
Theorems~\ref{thm:c2v18-tensor}, \ref{thm:c2v18-inertia}.
\end{proof}

\begin{programme}[V19 acquisition order]
\label{prog:c2v19}
\begin{enumerate}[label=\textup{(AC\arabic*)},leftmargin=4.8em]
\item Vorticity moments.  Derive the exact identity for the Gaussian
      enstrophy \(\int|\Omega|^2G\) of the profile along \(\sigma\), whose
      nonlinear term is the Gaussian vortex-stretching
      \(\int\Omega\cdot(\Omega\cdot\nabla)V\,G\) plus Gaussian boundary
      corrections; record its mean identity and whether the mean
      Gaussian vortex stretching has a sign.
\item Consolidation.  Prepare the V20 audit and the consolidation of
      V16--V19 (centre mode, momentum, odd profiles, tensor, inertia,
      enstrophy identities) as the moment ledger of the profile.
\item The next compulsory companion audit is V20.
\end{enumerate}
\end{programme}

\begin{assessment}[Position after V18]
\label{ass:c2v18-position}
The mean Gaussian nonlinear stress is negative semidefinite and
dominated by minus the mean energy tensor; the inertia identity is
exact with a conditional sign.  No full stability or global-regularity
claim is made.
\end{assessment}

% =============================================================================
\section*{Version V18 (second cycle) append-only record}
\addcontentsline{toc}{section}{Version V18 (second cycle) append-only record}
% =============================================================================

The complete V17 source of the second cycle was retained before this
continuation.  Its SHA--256 digest was
\begin{center}\scriptsize\ttfamily
76e21f75bce515565bf2e434ea1f945ff4a493c153b25fd938c266062f07c60b.
\end{center}
V18 proved the tensor identity with the signed mean stress and the
inertia identity.  No full regularity claim is made.

% =============================================================================
\section{V19: the Gaussian enstrophy identity and the mean vortex stretching}
\label{sec:c2v19-entry}
% =============================================================================

Items (AC1) and (AC2) of Programme~\ref{prog:c2v19} are carried out.  The
vorticity of the profile satisfies the self-similar vorticity equation,
and the Gaussian enstrophy of the profile obeys an exact identity in
log-scale: its linear part is the Gaussian dissipation of vorticity
plus twice the Gaussian enstrophy, its nonlinear part the Gaussian
vortex stretching minus the Gaussian enstrophy transport.  Against every
invariant measure, the mean of stretching minus transport equals the
mean of vorticity dissipation plus enstrophy, hence is positive: the
mean Gaussian vortex stretching, corrected by transport, is at least the
mean Gaussian enstrophy.  An exact relation between the Gaussian
enstrophy and the Gaussian Dirichlet integral of the profile is
recorded.  The consolidation for V20 is prepared.  No global regularity
claim is made.  The next version is the compulsory companion audit.

% =============================================================================
\section{The self-similar vorticity equation}
\label{sec:c2v19-vorticity}
% =============================================================================

\begin{lemma}[Vorticity of the profile]
\label{lem:c2v19-vorticity-equation}
For \(W\in\mathcal A_\nu(C)\) with the Morrey bound, \(\Omega:=\operatorname{curl}V\)
satisfies, for every \(\sigma\),
\begin{equation}
 \partial_\sigma\Omega=\nu\Delta\Omega-\tfrac12z\cdot\nabla\Omega-\Omega-(V\cdot\nabla)\Omega+(\Omega\cdot\nabla)V .
 \label{eq:c2v19-vorticity-equation}
\end{equation}
Moreover \(\Omega(\cdot,\sigma)\not\equiv0\) for every \(\sigma\) whenever
\(W\ne0\) is an element of a rescaling family or an ancient limit
(vorticity at infinity at every time, Theorem V63-vorticity-every-time
of the first cycle).
\end{lemma}

\begin{proof}
Take the curl of \eqref{eq:c2v2-profile-equation}:
\(\operatorname{curl}(z\cdot\nabla V)=z\cdot\nabla\Omega+\Omega\) (the dilation
generator satisfies \(\operatorname{curl}\circ(z\cdot\nabla)=(z\cdot\nabla+1)\circ\operatorname{curl}\)),
so \(\operatorname{curl}L_-V=\nu\Delta\Omega-\frac12z\cdot\nabla\Omega-\Omega\);
\(\operatorname{curl}((V\cdot\nabla)V)=(V\cdot\nabla)\Omega-(\Omega\cdot\nabla)V\) for
divergence-free \(V\), and \(\operatorname{curl}\nabla P_V=0\).
\end{proof}

% =============================================================================
\section{The Gaussian enstrophy identity}
\label{sec:c2v19-enstrophy}
% =============================================================================

Define the Gaussian enstrophy, the Gaussian vorticity dissipation, the
Gaussian vortex stretching, and the Gaussian enstrophy transport,
\begin{equation}
 \mathcal E:=\int|\Omega|^2G,\quad
 \mathcal E_\nabla:=\nu\int|\nabla\Omega|^2G,\quad
 \mathcal S:=\int\Omega\cdot(\Omega\cdot\nabla)V\,G,\quad
 \mathcal T:=\frac1{4\nu}\int|\Omega|^2(z\cdot V)\,G .
 \label{eq:c2v19-quantities}
\end{equation}

\begin{theorem}[Gaussian enstrophy identity]
\label{thm:c2v19-enstrophy}
For every \(\sigma\),
\begin{equation}
 \boxed{\quad
 \frac12\frac{\dd}{\dd\sigma}\mathcal E=-\mathcal E_\nabla-\mathcal E+\mathcal S-\mathcal T ,
 \quad}
 \label{eq:c2v19-identity}
\end{equation}
all terms finite, with \(\mathcal E\le2\|\nabla V\|_G^2\le C/(2\nu)\) and
\(|\mathcal S|\le C_\infty'\mathcal E\), \(|\mathcal T|\le\frac{C_\infty}{4\nu}\||z|^{1/2}\Omega\|_G^2\).
For every invariant probability measure \(\mathfrak m\) on a rescaling
family at a homogeneous point,
\begin{equation}
 \boxed{\quad
 \int(\mathcal S-\mathcal T)\,\dd\mathfrak m=\int(\mathcal E_\nabla+\mathcal E)\,\dd\mathfrak m\ \ge\ \int\mathcal E\,\dd\mathfrak m\ >\ 0 .
 \quad}
 \label{eq:c2v19-mean}
\end{equation}
\end{theorem}

\begin{proof}
Multiply \eqref{eq:c2v19-vorticity-equation} by \(\Omega G\) and integrate.
\(\int\Omega\cdot(\nu\Delta\Omega-\frac12z\cdot\nabla\Omega)G=-\nu\int|\nabla\Omega|^2G\)
by the \(G\)-symmetry of \(\nu\Delta-\frac12z\cdot\nabla=\nu G^{-1}\operatorname{div}(G\nabla\cdot)\);
\(-\int|\Omega|^2G=-\mathcal E\);
\(-\int\Omega\cdot(V\cdot\nabla)\Omega G=-\frac12\int V\cdot\nabla|\Omega|^2G=\frac12\int|\Omega|^2V\cdot\nabla G=-\frac1{4\nu}\int|\Omega|^2(z\cdot V)G=-\mathcal T\)
using \(\operatorname{div}V=0\); and \(\int\Omega\cdot(\Omega\cdot\nabla)VG=\mathcal S\).
Finiteness: \(|\Omega|\le\sqrt2|\nabla V|\) gives \(\mathcal E\le2\|\nabla V\|_G^2=d/(2\nu)\);
\(\nabla\Omega\) is bounded by \(\sqrt2|\nabla^2V|\le\sqrt2C_\infty''\);
the stretching by \(\|\nabla V\|_\infty\mathcal E\); the transport by the
\(L^\infty\) bound and the Gaussian Hardy inequality.  The mean form:
\(\mathcal E\) is bounded on the family, so the mean of its derivative
vanishes by the invariant-measure argument; \(\mathcal E_\nabla\ge0\) and,
by Lemma~\ref{lem:c2v19-vorticity-equation}, \(\mathcal E(\Theta_\tau W)>0\)
for every \(\tau\) and every \(W\) in the family, so
\(\int\mathcal E\,\dd\mathfrak m>0\) (\(\mathcal E\) is continuous on the
family by the fixed-time strong \(L^2(\psi)\) convergence of gradients,
Lemma~\ref{lem:c2v12-measurable}, and positive, hence bounded below on
the compact family).
\end{proof}

\begin{proposition}[Gaussian enstrophy and Gaussian Dirichlet integral]
\label{prop:c2v19-enstrophy-dirichlet}
For every \(\sigma\),
\begin{equation}
 \mathcal E=\|\nabla V\|_G^2+\frac1{2\nu}\|V\|_G^2-\frac1{4\nu^2}\int(z\cdot V)^2G .
 \label{eq:c2v19-enstrophy-dirichlet}
\end{equation}
\end{proposition}

\begin{proof}
\(|\Omega|^2=|\nabla V|^2-\partial_iV_j\partial_jV_i\) pointwise, and
\(\int\partial_iV_j\partial_jV_iG=-\int V_j\partial_jV_i\partial_iG=\frac1{2\nu}\int z_iV_j\partial_jV_iG\)
by \(\operatorname{div}V=0\); then
\(\int z_iV_j\partial_jV_iG=-\int V_iV_j\partial_j(z_iG)=-\int|V|^2G+\frac1{2\nu}\int(z\cdot V)^2G\).
Combine: \(\int\partial_iV_j\partial_jV_iG=-\frac1{2\nu}\|V\|_G^2+\frac1{4\nu^2}\int(z\cdot V)^2G\).
\end{proof}

\begin{assessment}[The mean vortex stretching]
\label{ass:c2v19-stretching}
Statistically, along the scales of a homogeneous singular point, the
Gaussian vortex stretching of the profile exceeds the Gaussian
enstrophy transport by exactly the Gaussian vorticity dissipation plus
the Gaussian enstrophy: the stretching must produce enstrophy at the
rate at which the self-similar damping (rate two, from the factor
\(-\Omega\) of the vorticity equation, which encodes the scaling of
vorticity) and the viscosity remove it.  This is the vorticity-level
counterpart of the exact mean identities of V12; like them it is a
balance, not an obstruction.  Together with the signed mean stress of
V18 it says that, in the mean, the nonlinearity acts against the
velocity in every direction and for the vorticity, which is the
expected picture of a blow-up sustained by stretching against
diffusion, here made exact with Gaussian weights.  No global
regularity claim is made.
\end{assessment}

% =============================================================================
\section{Integrated V19 conclusion and next acquisition order}
\label{sec:c2v19-conclusion}
% =============================================================================

\begin{theorem}[What V19 proves]
\label{thm:c2v19-conclusion}
Relative to the first cycle and V1--V18, modulo (E1)--(E9): the
self-similar vorticity equation, the Gaussian enstrophy identity
\eqref{eq:c2v19-identity} with its mean form \eqref{eq:c2v19-mean}, and
the relation \eqref{eq:c2v19-enstrophy-dirichlet}.  The full
Navier--Stokes stability/global-regularity theorem remains unproved.
\end{theorem}

\begin{proof}
The cited results.
\end{proof}

\begin{programme}[V20 acquisition order, with compulsory companion audit]
\label{prog:c2v20}
\begin{enumerate}[label=\textup{(AC\arabic*)},leftmargin=4.8em]
\item Perform the compulsory review of the current SM and YM notes;
      record digests; import nothing numerical.
\item Consolidate V16--V19 as the moment ledger of the profile: the
      centre-mode, momentum, tensor, inertia, and enstrophy identities
      with their exact mean forms and signs; list which means are
      signed and which are balances.
\item Record the position of the second cycle after twenty versions and
      the direction for V21--V25.
\item The next compulsory companion audit after V20 is V25.
\end{enumerate}
\end{programme}

\begin{assessment}[Position after V19]
\label{ass:c2v19-position}
The Gaussian enstrophy of the profile obeys an exact identity whose mean
form is a positive balance of vortex stretching against damping.  No
full stability or global-regularity claim is made.
\end{assessment}

% =============================================================================
\section*{Version V19 (second cycle) append-only record}
\addcontentsline{toc}{section}{Version V19 (second cycle) append-only record}
% =============================================================================

The complete V18 source of the second cycle was retained before this
continuation.  Its SHA--256 digest was
\begin{center}\scriptsize\ttfamily
8ba39bf5135700f3f04038136573b1f08cc93ba95c1004050ea9cbda316e9412.
\end{center}
V19 proved the Gaussian enstrophy identity of the profile with its
positive mean balance and the enstrophy--Dirichlet relation.  No full
regularity claim is made.

% =============================================================================
\section{V20: compulsory companion audit and the moment ledger of the profile}
\label{sec:c2v20-entry}
% =============================================================================

This is the twentieth version of the second cycle.  The compulsory
review of the two companion programmes is performed first.  V16--V19 are
then consolidated as the moment ledger of the profile: the exact
identities in log-scale for the Gaussian mean of the defect, the
Gaussian momentum, the second-moment tensor, the inertia, and the
Gaussian enstrophy, each with its mean form against every invariant
measure of the rescaling flow, and a classification of the mean
statements into signed constraints and balances.  The position after
twenty versions and the direction for V21--V25 are recorded.  No global
regularity claim is made.  The next compulsory companion audit is V25.

% =============================================================================
\section{Compulsory V20 review of the current companion notes}
\label{sec:c2v20-companion-audit}
% =============================================================================

The current companion files in \texttt{latex\_notes} were inspected
read-only.  The frozen checkpoints are
\begin{align}
 &\texttt{Renewal\_SM\_Einstein\_Continuum\_Research\_Notes\_V90.tex},
 \notag\\[-2pt]
 &\qquad
 \texttt{d822f8d46f0619a01e352d8f98a9c085f9caebf3cb25f4dc440f9141abb3c784},
 \label{eq:c2v20-SM-checkpoint}\\
 &\texttt{Renewal\_Yang\_Mills\_Mass\_Gap\_Research\_Notes\_V63.tex},
 \notag\\[-2pt]
 &\qquad
 \texttt{f4bfc57c40ff1c54778c1d11b714292942a1a264d368b4871096e685bb0bbad8}.
 \label{eq:c2v20-YM-checkpoint}
\end{align}

\emph{SM V89--V90.}  The SM programme computed minimal heat-supertrace
test blocks (one scalar, one gauge generator with ghost, one Dirac
modulus cancel the volume coefficient but leave a curvature term;
general multiplicities in a fixed ratio cancel both), showed that on
its matched fixed-volume branch the test curvature divergence is
quadratic in the uncancelled beat amplitude and cannot cancel the
linear conformal coarea trace; V90 was its compulsory checkpoint,
recording the present cycle's V15 by digest and transferring only the
structural rule of assembling before limiting, then proving an exact
proper-time subtraction theorem for its conformal coarea determinant,
identifying it with the sharp Fourier subtraction, computing the cutoff
conversion, and proving the noncancellation of the linear coarea trace
on its beat line.

\emph{YM V63.}  The YM programme proved the clock derivative of a
complete Gaussian centre, an exact small-clock Poisson covariance
identity, an enhanced one-mark centre covariance, closed its scalar
clock quotient at the first mark for every order, closed the zero- and
one-pair rows of its one-centre normal form, and formulated an all-mark
target.

\begin{theorem}[V20 companion-audit decision]
\label{thm:c2v20-companion-decision}
No SM V89--V90 supertrace, subtraction, or cutoff-conversion constant
and no YM V63 covariance or normal-form estimate is imported.  Neither
bears on the moment ledger below or on the gates of
Theorem~\ref{thm:c2v15-gates}.  The SM V90 reading of the present cycle
(a digest and the structural rule) is accurate.
\end{theorem}

\begin{proof}
The SM results concern proper-time regularizations of determinants of
a conformal operator on a torus; the YM results concern covariances of
a lattice clock.  None is a Navier--Stokes statement.
\end{proof}

% =============================================================================
\section{The moment ledger of the profile}
\label{sec:c2v20-ledger}
% =============================================================================

\begin{theorem}[Moment ledger, V16--V19]
\label{thm:c2v20-ledger}
Let \(W\in\mathcal A_\nu(C)\) with the Morrey bound and the uniform bounds
\eqref{eq:c2v4-uniform}, \(V\) its profile, \(Y=\partial_\sigma V\),
\(\Omega=\operatorname{curl}V\), \(N=N(V)\), and the Gaussian quantities of
\eqref{eq:c2v16-means}, \eqref{eq:c2v17-Nbar}, \eqref{eq:c2v18-tensors},
\eqref{eq:c2v18-inertia}, \eqref{eq:c2v19-quantities}.  The following
identities hold for every \(\sigma\) (for the defect, a.e.\ \(\sigma\)):
\begin{align}
 \frac{\dd}{\dd\sigma}\bar V&=-\tfrac12\bar V-\overline{N},
 &\frac{\dd}{\dd\sigma}\bar Y&=-\tfrac12\bar Y-\overline{N'(V)Y},
 \label{eq:c2v20-means}\\
 \frac{\dd}{\dd\sigma}T&=-D-T-S,
 &\frac{\dd}{\dd\sigma}I&=-2\nu I_\nabla+6\nu\varphi-2I+2\nu j-\tfrac12j_2,
 \label{eq:c2v20-tensor-inertia}\\
 \frac12\frac{\dd}{\dd\sigma}\mathcal E&=-\mathcal E_\nabla-\mathcal E+\mathcal S-\mathcal T,
 &\frac{\dd}{\dd\sigma}\mathcal L&=-\|Y\|_G^2-\langle N,Y\rangle_G .
 \label{eq:c2v20-enstrophy-second}
\end{align}
Against every \(\Theta\)-invariant probability measure \(\mathfrak m\) on a
rescaling family at a homogeneous point, with \(\langle\cdot\rangle\)
denoting the \(\mathfrak m\)-mean, modulo (E1)--(E9):
\begin{enumerate}[label=\textup{(\arabic*)},leftmargin=3.2em]
\item \emph{Signed constraints.}
      \(\langle S\rangle=-\langle D+T\rangle\preceq-\langle T\rangle\preceq0\) with
      \(\operatorname{tr}\langle T\rangle\ge c_\Phi\);
      \(\langle\langle N,Y\rangle_G\rangle=-\langle\|Y\|_G^2\rangle\le-\underline\delta\);
      \(\langle\langle Y,N'Y\rangle_G\rangle=-\langle\nu\|\nabla Y\|_G^2+\frac12\|Y\|_G^2\rangle\le-\frac12\underline\delta\);
      \(\langle\mathcal S-\mathcal T\rangle=\langle\mathcal E_\nabla+\mathcal E\rangle>0\).
\item \emph{Balances without sign.}
      \(\langle\bar V\rangle=-2\langle\overline N\rangle\);
      \(\langle|G||\bar Y|^2\rangle=-2|G|\langle\bar Y\cdot\overline{N'Y}\rangle\);
      \(\frac12\langle j_2\rangle=\langle2\nu\varphi-2\nu d-2I-2\nu I_\nabla\rangle\).
\end{enumerate}
\end{theorem}

\begin{proof}
Propositions~\ref{prop:c2v16-centre-equation}, \ref{prop:c2v17-mean-profile},
Theorems~\ref{thm:c2v16-split}, \ref{thm:c2v18-tensor},
\ref{thm:c2v18-inertia}, \ref{thm:c2v19-enstrophy}, \ref{thm:c2v2-second},
\ref{thm:c2v12-mean-identities}.
\end{proof}

\begin{assessment}[What the ledger is and is not]
\label{ass:c2v20-ledger}
The ledger is the complete set of exact scale-invariant balances the
series has for the statistics of a type-I blow-up at a homogeneous
point: each Gaussian moment of the profile is damped by the self-similar
scaling at a definite rate (one half for the mean, one for the energy
tensor, two for the inertia and the enstrophy) and by viscosity, and
the corresponding nonlinear moment must exactly compensate in the
mean.  The signed constraints say that the nonlinear moments do so with
the sign that feeds the blow-up: negative stress against the velocity,
positive stretching of the vorticity, anti-alignment with the
self-similarity defect.  None of these is contradicted by the bounds on
the nonlinear moments, which are of order one; the ledger is a
consistent description of a hypothetical blow-up, not an obstruction to
it.  What is missing is a relation between two rows of the ledger that
the Navier--Stokes nonlinearity cannot satisfy simultaneously; the
series has found none.  No global regularity claim is made.
\end{assessment}

% =============================================================================
\section{Position after twenty versions and direction}
\label{sec:c2v20-position}
% =============================================================================

\begin{assessment}[Position after twenty versions of the second cycle]
\label{ass:c2v20-position-twenty}
The second cycle has: the self-similar variables and the two exact
identities (V1--V3); the dichotomy and, with (E9), the unconditional
uniform non-self-similarity of homogeneous points (V4--V9); the discrete
dichotomy, recurrence, invariant and blow-up measures with exact mean
identities (V10--V13); the rate-free Dirichlet ledger (V14); and the
moment ledger (V16--V19).  The three gates stand as in
Theorem~\ref{thm:c2v15-gates}.  For V21--V25 the direction is to search
the ledger for an incompatible pair of rows: candidates are the pairing
of the defect identity with the enstrophy identity through the exact
relation \eqref{eq:c2v19-enstrophy-dirichlet}, and the pairing of the
tensor identity with the centre-mode equation through the dipole
matrix; each is an exact computation.  If none yields an
incompatibility, the series will record the moment ledger as closed and
go upstream to the lacunary points (G1\('\)) and the rate-free classes,
where the second cycle's tools have not been applied.
\end{assessment}

% =============================================================================
\section{Integrated V20 conclusion and next acquisition order}
\label{sec:c2v20-conclusion}
% =============================================================================

\begin{theorem}[What V20 establishes]
\label{thm:c2v20-conclusion}
Relative to the first cycle and V1--V19: the companion-audit decision
(Theorem~\ref{thm:c2v20-companion-decision}) and the moment ledger
(Theorem~\ref{thm:c2v20-ledger}).  The full Navier--Stokes
stability/global-regularity theorem remains unproved.
\end{theorem}

\begin{proof}
The cited results.
\end{proof}

\begin{programme}[V21 acquisition order]
\label{prog:c2v21}
\begin{enumerate}[label=\textup{(AC\arabic*)},leftmargin=4.8em]
\item Defect and enstrophy.  Differentiate \eqref{eq:c2v19-enstrophy-dirichlet}
      in \(\sigma\) and compare with the enstrophy identity and the first
      identity to obtain the exact identity for
      \(\frac{\dd}{\dd\sigma}\int(z\cdot V)^2G\); record its mean form and
      whether it closes a relation between \(\langle\mathcal S\rangle\),
      \(\langle j\rangle\), and \(\langle\|Y\|_G^2\rangle\).
\item Radial momentum.  The quantity \(\int(z\cdot V)^2G\) is the Gaussian
      radial kinetic energy; relate its mean to the flux \(j\) through
      the identity \(\frac1{2\nu}\int(z\cdot V)^2G=\|V\|_G^2-\int V\cdot(z\cdot\nabla)V\,G\)
      (integration by parts) and determine what the ledger says about
      the mean radial share of the energy.
\item The next compulsory companion audit is V25.
\end{enumerate}
\end{programme}

\begin{assessment}[Position after V20]
\label{ass:c2v20-position}
The second cycle is consolidated on the moment ledger of the profile,
with its signed constraints and balances.  No full stability or
global-regularity claim is made.
\end{assessment}

% =============================================================================
\section*{Version V20 (second cycle) append-only record}
\addcontentsline{toc}{section}{Version V20 (second cycle) append-only record}
% =============================================================================

The complete V19 source of the second cycle was retained before this
continuation.  Its SHA--256 digest was
\begin{center}\scriptsize\ttfamily
33c2da92a60d5b8a2b69fea3c5c37d4d0d698419cc9eeb1c6334bd571a23a25e.
\end{center}
V20 performed the compulsory companion audit and consolidated V16--V19
as the moment ledger of the profile.  No full regularity claim is made.

% =============================================================================
\section{V21: the radial kinetic energy, the closure of the ledger on itself,
and the enstrophy--Dirichlet inequality}
\label{sec:c2v21-entry}
% =============================================================================

Items (AC1) and (AC2) of Programme~\ref{prog:c2v21} are carried out.  The
Gaussian radial kinetic energy of the profile is shown to be an exact
linear combination of the Gaussian Dirichlet functional and the
Gaussian enstrophy, so that differentiating the enstrophy--Dirichlet
relation reproduces the second and the enstrophy identities and yields
no new relation: the ledger closes on itself at this pairing.  The
nonnegativity of the radial kinetic energy gives the exact pointwise
inequality that the Gaussian enstrophy is at most \(2/\nu\) times the
Gaussian Dirichlet functional, with equality exactly for tangential
profiles, and its mean form determines the mean radial kinetic energy
by the mean Dirichlet functional and the mean net vortex stretching.
The identity proposed in item (AC2) of the programme is corrected.  No
global regularity claim is made.  The next compulsory companion audit is
V25.

% =============================================================================
\section{The radial kinetic energy}
\label{sec:c2v21-radial}
% =============================================================================

Define the Gaussian radial kinetic energy of the profile,
\begin{equation}
 R:=\int_{\R^3}(z\cdot V)^2\,G\,\dd z\ \ge0 .
 \label{eq:c2v21-R}
\end{equation}

\begin{theorem}[Closure of the ledger at the enstrophy--Dirichlet pairing]
\label{thm:c2v21-closure}
For \(W\in\mathcal A_\nu(C)\) with the Morrey bound and every \(\sigma\),
\begin{equation}
 \boxed{\quad
 R=8\nu\,\mathcal L-4\nu^2\,\mathcal E ,
 \qquad\text{equivalently}\qquad
 \mathcal E=\frac{2}{\nu}\mathcal L-\frac{R}{4\nu^2}\ \le\ \frac2\nu\mathcal L ,
 \quad}
 \label{eq:c2v21-closure}
\end{equation}
with equality \(\mathcal E=\frac2\nu\mathcal L\) if and only if
\(z\cdot V\equiv0\).  Consequently
\begin{equation}
 \frac{\dd}{\dd\sigma}R=8\nu\frac{\dd}{\dd\sigma}\mathcal L-4\nu^2\frac{\dd}{\dd\sigma}\mathcal E
 =-8\nu\bigl(\|Y\|_G^2+\langle N,Y\rangle_G\bigr)+8\nu^2\bigl(\mathcal E_\nabla+\mathcal E-\mathcal S+\mathcal T\bigr),
 \label{eq:c2v21-R-derivative}
\end{equation}
and against every invariant probability measure \(\mathfrak m\) on a
rescaling family at a homogeneous point the mean of
\eqref{eq:c2v21-R-derivative} vanishes identically by the second and
the enstrophy identities: no relation between
\(\langle\|Y\|_G^2\rangle\), \(\langle\langle N,Y\rangle_G\rangle\),
\(\langle\mathcal S\rangle\), \(\langle\mathcal T\rangle\) beyond
\eqref{eq:c2v12-mean2} and \eqref{eq:c2v19-mean} follows.
\end{theorem}

\begin{proof}
\eqref{eq:c2v19-enstrophy-dirichlet} multiplied by \(4\nu^2\) reads
\(4\nu^2\mathcal E=4\nu^2\|\nabla V\|_G^2+2\nu\|V\|_G^2-R\), and
\(8\nu\mathcal L=4\nu^2\|\nabla V\|_G^2+2\nu\|V\|_G^2\) by
\eqref{eq:c2v2-L}; subtract.  \(R\ge0\) with equality if and only if
\(z\cdot V=0\) almost everywhere, hence everywhere by continuity.
Differentiating \eqref{eq:c2v21-closure} and inserting
\eqref{eq:c2v2-second} and \eqref{eq:c2v19-identity} gives
\eqref{eq:c2v21-R-derivative}; the mean of each side vanishes by
Theorem~\ref{thm:c2v12-mean-identities} and Theorem~\ref{thm:c2v19-enstrophy}
(\(R\) is bounded on the family by the Gaussian Hardy inequality).
\end{proof}

\begin{corollary}[Mean radial kinetic energy]
\label{cor:c2v21-mean-radial}
For every invariant probability measure \(\mathfrak m\) as above,
\begin{equation}
 \langle R\rangle=8\nu\langle\mathcal L\rangle-4\nu^2\langle\mathcal E\rangle
 =8\nu\langle\mathcal L\rangle-4\nu^2\bigl\langle\mathcal S-\mathcal T-\mathcal E_\nabla\bigr\rangle ,
 \qquad
 0\le\langle R\rangle\le8\nu\langle\mathcal L\rangle ,
 \label{eq:c2v21-mean-radial}
\end{equation}
so the mean Gaussian net stretching \(\langle\mathcal S-\mathcal T\rangle\)
lies in \([\langle\mathcal E_\nabla\rangle,\ \langle\mathcal E_\nabla\rangle+\frac2\nu\langle\mathcal L\rangle]\).
\end{corollary}

\begin{proof}
Take means in \eqref{eq:c2v21-closure} and insert \eqref{eq:c2v19-mean}.
\end{proof}

\begin{remark}[Correction of item (AC2) of Programme~\ref{prog:c2v21}]
\label{rem:c2v21-correction}
The identity proposed there,
\(\frac1{2\nu}\int(z\cdot V)^2G=\|V\|_G^2-\int V\cdot(z\cdot\nabla)V\,G\), is not
correct: integration by parts gives
\(\int V\cdot(z\cdot\nabla)V\,G=\frac12\int z\cdot\nabla|V|^2G=-\frac32\|V\|_G^2+\frac1{4\nu}I\)
with \(I\) the Gaussian inertia, which involves \(I\), not \(R\).  The
correct relation for \(R\) is \eqref{eq:c2v21-closure}.  The proposed
identity appeared only as a programme item and was used nowhere.
\end{remark}

\begin{assessment}[Tangential and radial profiles]
\label{ass:c2v21-tangential}
The inequality \(\mathcal E\le\frac2\nu\mathcal L\) is saturated by
tangential profiles, \(z\cdot V\equiv0\), for which the radial flux and
the Bernoulli radial term vanish identically (the swirl of the first
cycle is the example outside the class); a profile with radial kinetic
energy has Gaussian enstrophy strictly smaller than \(\frac2\nu\mathcal L\)
by exactly \(R/(4\nu^2)\).  In the mean, the net vortex stretching is
bounded above by the vorticity dissipation plus \(\frac2\nu\langle\mathcal L\rangle\),
so a blow-up statistics with large mean stretching must have small
mean radial kinetic energy: the flow near a self-similarly sustained
singularity is, statistically, the more tangential the more it
stretches.  This is an exact trade-off within the ledger, not an
obstruction.  No global regularity claim is made.
\end{assessment}

% =============================================================================
\section{Integrated V21 conclusion and next acquisition order}
\label{sec:c2v21-conclusion}
% =============================================================================

\begin{theorem}[What V21 proves]
\label{thm:c2v21-conclusion}
Relative to the first cycle and V1--V20, modulo (E1)--(E9): the closure
identity \eqref{eq:c2v21-closure} with the enstrophy--Dirichlet
inequality, the derivative \eqref{eq:c2v21-R-derivative}, and the mean
radial kinetic energy \eqref{eq:c2v21-mean-radial}.  The full
Navier--Stokes stability/global-regularity theorem remains unproved.
\end{theorem}

\begin{proof}
Theorem~\ref{thm:c2v21-closure}, Corollary~\ref{cor:c2v21-mean-radial}.
\end{proof}

\begin{programme}[V22 acquisition order]
\label{prog:c2v22}
\begin{enumerate}[label=\textup{(AC\arabic*)},leftmargin=4.8em]
\item Tensor and centre mode.  Pair the tensor identity with the
      centre-mode equation: the Gaussian mean of the momentum-flux
      density \(\overline N\) and the stress \(S\) are both moments of
      \(N\otimes V\); determine whether \(\langle\overline N\rangle\) is
      determined by \(\langle S\rangle\) and the dipole matrix, and record
      whether an incompatibility arises with the signed mean stress.
\item Closing the ledger search.  If no incompatible pair of rows is
      found, record the moment ledger as closed for the purposes of the
      series and prepare the upstream move to the lacunary points and
      the rate-free classes.
\item The next compulsory companion audit is V25.
\end{enumerate}
\end{programme}

\begin{assessment}[Position after V21]
\label{ass:c2v21-position}
The ledger closes on itself at the enstrophy--Dirichlet pairing, with
the exact inequality \(\mathcal E\le\frac2\nu\mathcal L\).  No full stability
or global-regularity claim is made.
\end{assessment}

% =============================================================================
\section*{Version V21 (second cycle) append-only record}
\addcontentsline{toc}{section}{Version V21 (second cycle) append-only record}
% =============================================================================

The complete V20 source of the second cycle was retained before this
continuation.  Its SHA--256 digest was
\begin{center}\scriptsize\ttfamily
9dc4653962e31cba3f11929c2bcd9ee7629c8eb66cc40fd0038e90dfd148ec8b.
\end{center}
V21 proved the closure identity for the radial kinetic energy, the
enstrophy--Dirichlet inequality, and the mean radial kinetic energy,
and corrected a programme item.  No full regularity claim is made.

% =============================================================================
\section{V22: mean-flow and fluctuation parts of the Gaussian stress, and
the closing of the ledger search}
\label{sec:c2v22-entry}
% =============================================================================

Items (AC1) and (AC2) of Programme~\ref{prog:c2v22} are carried out.  The
Gaussian stress of the profile splits into a mean-flow part, built from
the Gaussian momentum and the Gaussian momentum flux, and a fluctuation
part; the tensor identity splits accordingly, and against every
invariant measure each part is separately negative semidefinite in the
mean: the mean-flow stress equals minus the mean momentum tensor and
the fluctuation stress equals minus the mean fluctuation energy tensor
minus the gradient tensor.  The Gaussian momentum flux is not
determined by the stress and the dipole matrix, and no incompatibility
arises.  With this pairing the search of the ledger for an
incompatible pair of rows is recorded as closed for the purposes of
the series, and the upstream move to lacunary points and the rate-free
classes is prepared.  No global regularity claim is made.  The next
compulsory companion audit is V25.

% =============================================================================
\section{Mean flow and fluctuation}
\label{sec:c2v22-split}
% =============================================================================

For \(W\in\mathcal A_\nu(C)\) with the Morrey bound, write
\(V=\bar V+V'\), \(N=\bar N+N'\) with \(\bar V,\bar N\) the Gaussian means
of \eqref{eq:c2v16-means}, \eqref{eq:c2v17-Nbar}, so that
\(\int V'G=\int N'G=0\), and define
\begin{equation}
 T_0:=|G|\,\bar V\otimes\bar V,\quad
 T':=\int V'\otimes V'\,G,\quad
 S_0:=|G|\bigl(\bar N\otimes\bar V+\bar V\otimes\bar N\bigr),\quad
 S':=\int\bigl(N'\otimes V'+V'\otimes N'\bigr)G ,
 \label{eq:c2v22-split}
\end{equation}
so that \(T=T_0+T'\) and \(S=S_0+S'\) (the cross terms vanish because
the fluctuations have zero Gaussian mean).

\begin{theorem}[Split tensor identity and separately signed mean stresses]
\label{thm:c2v22-split}
For every \(\sigma\),
\begin{equation}
 \frac{\dd}{\dd\sigma}T_0=-T_0-S_0,
 \qquad
 \frac{\dd}{\dd\sigma}T'=-D-T'-S' ,
 \label{eq:c2v22-split-identity}
\end{equation}
and against every \(\Theta\)-invariant probability measure \(\mathfrak m\)
on a rescaling family at a homogeneous point,
\begin{equation}
 \boxed{\quad
 \langle S_0\rangle=-\langle T_0\rangle\preceq0,
 \qquad
 \langle S'\rangle=-\langle D+T'\rangle\preceq-\langle T'\rangle\preceq0 .
 \quad}
 \label{eq:c2v22-signed}
\end{equation}
\end{theorem}

\begin{proof}
From \eqref{eq:c2v17-mean-equation},
\(\frac{\dd}{\dd\sigma}(\bar V\otimes\bar V)=-\bar V\otimes\bar V-(\bar N\otimes\bar V+\bar V\otimes\bar N)\),
which is the first identity after multiplication by \(|G|\).
Subtracting it from \eqref{eq:c2v18-tensor-identity} gives the second,
since \(T-T_0=T'\) and \(S-S_0=S'\).  The mean forms follow as in
Theorem~\ref{thm:c2v18-tensor}, \(T_0\) and \(T'\) being bounded
observables; positivity of \(T_0\), \(T'\), \(D\) gives the matrix
inequalities.
\end{proof}

\begin{proposition}[The momentum flux is not determined by the stress]
\label{prop:c2v22-not-determined}
The Gaussian momentum flux \(\bar N\) and the stress \(S\) are related
only through \(S_0=|G|(\bar N\otimes\bar V+\bar V\otimes\bar N)\): given
\(\bar V\ne0\), \(S_0\) determines the component of \(\bar N\) along
\(\bar V\) and the projection of \(\bar N\) onto the plane orthogonal to
\(\bar V\) up to the symmetric-product ambiguity, and the dipole matrix
\(M_V\) does not enter.  Neither \eqref{eq:c2v22-signed} nor
\(\langle\bar V\rangle=-2\langle\bar N\rangle\) constrains \(\langle\bar N\rangle\)
beyond the sign relation \(\langle\bar N\cdot\bar V\rangle=-\frac12\langle|\bar V|^2\rangle\),
which is the trace of the first relation in \eqref{eq:c2v22-signed}
divided by \(2|G|\).  For odd profiles both \(T_0\) and \(S_0\) vanish and
\eqref{eq:c2v22-signed} reduces to \eqref{eq:c2v18-mean-stress}.
\end{proposition}

\begin{proof}
\(S_0e=|G|((\bar V\cdot e)\bar N+(\bar N\cdot e)\bar V)\) for any vector \(e\):
taking \(e=\bar V\) gives \(|G|(|\bar V|^2\bar N+(\bar N\cdot\bar V)\bar V)\),
from which \(\bar N\) is recovered when \(\bar V\ne0\) as stated; the
dipole matrix appears in the centre-mode equation only.  The trace
statement is direct.
\end{proof}

\begin{assessment}[Item (AC1)]
\label{ass:c2v22-AC1}
The mean-flow part of the blow-up statistics, the Gaussian momentum,
obeys a closed pair of exact relations: its mean is minus twice the
mean momentum flux, and its mean tensor is minus the mean mean-flow
stress.  These are the statements that, in the mean, the Gaussian
momentum decays at rate one half under the scaling and is fed by the
momentum flux, and they are compatible with any sign of the momentum
flux itself.  The fluctuation part obeys the same negative-definite
mean stress as the whole.  No incompatibility with the signed mean
stress arises.
\end{assessment}

% =============================================================================
\section{Closing the ledger search}
\label{sec:c2v22-closing}
% =============================================================================

\begin{theorem}[The moment ledger is closed for the purposes of the series]
\label{thm:c2v22-closed}
The exact identities of the second cycle for the Gaussian moments of
the profile (V2, V7, V16--V19, V21, V22) and their mean forms against
invariant measures are pairwise consistent: every pairing examined
(defect with Dirichlet functional, V2; defect with its own energy, V7;
centre mode with the whole defect, V16; momentum with stress, V22;
enstrophy with Dirichlet functional and radial kinetic energy, V21;
tensor with inertia, V18) yields either a tautology or a new balance
without sign, and the signed mean statements
\eqref{eq:c2v12-exact-signs}, \eqref{eq:c2v18-mean-stress},
\eqref{eq:c2v19-mean}, \eqref{eq:c2v22-signed} are mutually consistent
with the order-one bounds on the nonlinear moments.  No pair of rows
of the ledger is incompatible on the strength of the ledger alone.
\end{theorem}

\begin{proof}
Each cited identity is an exact consequence of the profile equation
tested against a Gaussian moment; consistency of the mean forms is the
statement that each mean identity is the vanishing of the mean of a
derivative of a bounded observable, hence holds for every invariant
measure simultaneously, and the signed statements are inequalities
between means of nonnegative quantities and means of nonlinear moments
bounded by constants of order one (Lemmas~\ref{lem:c2v3-N-bound},
\ref{lem:c2v12-measurable}, Proposition~\ref{prop:c2v7-defect-equation}).
An incompatibility would be a pair of identities whose combination
forces a nonnegative mean to be negative; none of the combinations in
the cited versions does so.
\end{proof}

\begin{assessment}[What closing means]
\label{ass:c2v22-closing}
Closing the ledger search records that no further exact Gaussian moment
identity of the kind derived in V16--V22 will produce a contradiction
by itself, because each such identity is a balance in which a scaling
damping is compensated in the mean by a nonlinear moment whose only
available bound is of order one.  A contradiction requires an
inequality for a nonlinear moment that is not a consequence of the
profile equation tested against a moment: a structural property of the
Navier--Stokes nonlinearity (for instance, a bound relating the mean
Gaussian vortex stretching to the mean Gaussian enstrophy that is
strictly weaker than what the enstrophy identity demands).  The series
does not have one.  The homogeneous type-I case therefore rests, as in
Theorem~\ref{thm:c2v15-gates}, on the single mean inequality for the
cross term, and the second cycle now turns to the cases the second
cycle's tools have not touched: lacunary points, where the pinching
fails, and the rate-free classes, where the ledgers hold but the
family is not compact.  No global regularity claim is made.
\end{assessment}

% =============================================================================
\section{Integrated V22 conclusion and next acquisition order}
\label{sec:c2v22-conclusion}
% =============================================================================

\begin{theorem}[What V22 proves]
\label{thm:c2v22-conclusion}
Relative to the first cycle and V1--V21, modulo (E1)--(E9): the split
tensor identity with separately signed mean stresses
(Theorem~\ref{thm:c2v22-split}), Proposition~\ref{prop:c2v22-not-determined},
and the closing of the ledger search (Theorem~\ref{thm:c2v22-closed}).
The full Navier--Stokes stability/global-regularity theorem remains
unproved.
\end{theorem}

\begin{proof}
The cited results.
\end{proof}

\begin{programme}[V23 acquisition order]
\label{prog:c2v23}
\begin{enumerate}[label=\textup{(AC\arabic*)},leftmargin=4.8em]
\item Statistics at an arbitrary singular point.  Define the full
      rescaling family \(\mathcal F(x_0)\), the flow, the invariant
      measures, and the blow-up measures at any \(x_0\in\Sigma\) of a
      type-I solution, without homogeneity; prove that the exact mean
      identities and the signed mean statements that do not use the
      pinching (the stress, the enstrophy balance, the defect identity)
      hold for every invariant measure, that the zero element is a
      fixed point of the flow, and that \(\delta_0\) is the only
      invariant measure for which the mean squared defect vanishes.
\item Statistical dichotomy for lacunary points.  Prove that a lacunary
      point either has a nontrivial blow-up measure, for which all the
      signed statements hold with positive mean squared defect, or is
      statistically trivial: every blow-up measure is \(\delta_0\),
      equivalently the active scales have zero logarithmic density;
      record what the first cycle's minimal log-period and count of
      Gaussian-active centres say in each case.
\item The next compulsory companion audit is V25.
\end{enumerate}
\end{programme}

\begin{assessment}[Position after V22]
\label{ass:c2v22-position}
The Gaussian stress splits into separately signed mean-flow and
fluctuation parts, and the ledger search is closed.  No full stability
or global-regularity claim is made.
\end{assessment}

% =============================================================================
\section*{Version V22 (second cycle) append-only record}
\addcontentsline{toc}{section}{Version V22 (second cycle) append-only record}
% =============================================================================

The complete V21 source of the second cycle was retained before this
continuation.  Its SHA--256 digest was
\begin{center}\scriptsize\ttfamily
3bb406e6daec16dc56b8a1583c5c6f39081d2e9862dfcdb5fdb1a091cce8abc2.
\end{center}
V22 split the Gaussian stress into mean-flow and fluctuation parts with
separately signed mean forms and closed the ledger search.  No full
regularity claim is made.

% =============================================================================
\section{V23: blow-up statistics at an arbitrary singular point and the
statistical dichotomy for lacunary points}
\label{sec:c2v23-entry}
% =============================================================================

Items (AC1) and (AC2) of Programme~\ref{prog:c2v23} are carried out.  The
full rescaling family, the rescaling flow, its invariant measures, and
the blow-up measures are defined at every singular point of a type-I
solution, without homogeneity.  The zero solution belongs to the family
exactly when the point has intermediate scales; it is a fixed point of
the flow, and the Dirac measure at it is invariant.  All exact mean
identities of the cycle and the signed mean statements that do not use
the pinching hold for every invariant measure at every singular point,
and, modulo (E9), the Dirac measure at zero is the only invariant
measure with vanishing mean squared defect.  Hence every singular point
is either statistically nontrivial, with a blow-up measure carrying
positive mean squared defect and all the signed statements, or
statistically trivial, with every blow-up measure equal to the Dirac
measure at zero, which is shown to be equivalent to the active scales
having zero logarithmic density.  Homogeneous points are nontrivial;
for lacunary points both cases remain possible, and what the first
cycle says about each is recorded.  No global regularity claim is made.
The next compulsory companion audit is V25.

% =============================================================================
\section{The family and the flow at an arbitrary singular point}
\label{sec:c2v23-family}
% =============================================================================

Let \(u\) satisfy the type-I rate and \(x_0\in\Sigma\).  Define
\(\mathcal F(x_0)\) as in Section~\ref{sec:c2v6-minimal}: all local limits
of rescalings \(U_\mu=\mu u(x_0+\mu\,\cdot,T_*+\mu^2\,\cdot)\), \(\mu\to0\),
together with all local limits of rescalings of such limits.

\begin{proposition}[Structure of the family without homogeneity]
\label{prop:c2v23-family}
\begin{enumerate}[label=\textup{(\roman*)},leftmargin=3.2em]
\item \(\mathcal F(x_0)\) is nonempty, scale-invariant, closed under local
      limits of rescalings, and sequentially compact in the metric
      \(\rho\); every element satisfies the uniform bounds
      \eqref{eq:c2v4-uniform} and the Morrey bound, and
      \(\Phi_W\le C_\Phi\) for all elements.
\item \(0\in\mathcal F(x_0)\) if and only if \(x_0\) has an intermediate
      sequence (\(\mathcal E_{x_0}(\rho_j;R')\to0\) for every \(R'\) along
      some \(\rho_j\to0\)), if and only if \(\liminf_{\lambda\to0}\Phi_{u,x_0}(\lambda)=0\).
      If \(x_0\) is homogeneous, \(0\notin\mathcal F(x_0)\).
\item \(\Theta\) is a continuous flow on \((\mathcal F(x_0),\rho)\); \(0\) is
      a fixed point; \(\delta_0\) is an invariant probability measure;
      \(\mathcal F(x_0)\setminus\{0\}\) is invariant.
\item Every nonzero \(W\in\mathcal F(x_0)\) has \(\Phi_W(\lambda)>0\) for all
      \(\lambda\) and is not backward self-similar (modulo (E8), (E9)).
\end{enumerate}
\end{proposition}

\begin{proof}
(i) The uniform bounds hold for all rescalings by the type-I rates
(Theorem V46-Linfty-upgrade, Theorem V47-morrey of the first cycle),
independently of the centre; the compactness is Theorem
V41-ancient-limit, and closure under limits of rescalings is the
diagonal argument of Proposition V83-family, which does not use
homogeneity.  (ii) If \(\rho_j\) is intermediate, the rescalings along
\(\rho_j\) converge locally to \(0\) (Theorem V97-ledger(ii) of the first
cycle).  Conversely, if \(U_{\mu_k}\to0\) locally then the dissipation
\(\mathcal D_{x_0}(\mu_k;R,S,\eta)\to0\) for all \(R,S,\eta\), and Theorem
V53-per-scale gives \(\mathcal E_{x_0}(\mu_k;R')\to0\) for all \(R'\): an
intermediate sequence.  The equivalence with \(\liminf\Phi_u=0\) is
Theorem V97-ledger(i),(ii) (bounded below along active scales, tending
to zero along intermediate scales) together with
Proposition~\ref{prop:c2v10-continuity} for the limit.  Homogeneous
points have no intermediate sequence (Theorem V97-ledger(iii)).
(iii) Lemma~\ref{lem:c2v11-flow}, whose proof uses only (i); \(0^\mu=0\).
(iv) If \(\Phi_W(\lambda)=0\) then \(W(-\lambda^2)=0\) and \(W=0\) on
\([-\lambda^2,0)\) by (E8), hence \(W=0\) on \((-\infty,0)\) by (E6) for the
difference with the zero solution; self-similar nonzero elements are
excluded by (E9) (Remark~\ref{rem:c2v9-hypotheses} uses only the
uniform bounds).
\end{proof}

\begin{theorem}[Mean identities and signs at an arbitrary singular point]
\label{thm:c2v23-identities}
Let \(x_0\in\Sigma\) and \(\mathfrak m\) any \(\Theta\)-invariant probability
measure on \(\mathcal F(x_0)\).  Modulo (E1)--(E9):
\begin{enumerate}[label=\textup{(\roman*)},leftmargin=3.2em]
\item All exact mean identities of the cycle hold:
      \eqref{eq:c2v12-mean1}--\eqref{eq:c2v12-mean3},
      \eqref{eq:c2v16-centre-mean}, \eqref{eq:c2v16-fluctuation-mean},
      \(\langle\bar V\rangle=-2\langle\bar N\rangle\), the tensor, inertia,
      and enstrophy mean identities of Theorem~\ref{thm:c2v20-ledger},
      and the split \eqref{eq:c2v22-signed}.
\item The signed statements hold in the weak form:
      \(\langle S\rangle=-\langle D+T\rangle\preceq0\),
      \(\langle\langle N,Y\rangle_G\rangle=-\langle\|Y\|_G^2\rangle\le0\),
      \(\langle\langle Y,N'Y\rangle_G\rangle=-\langle\nu\|\nabla Y\|_G^2+\frac12\|Y\|_G^2\rangle\le0\),
      \(\langle\mathcal S-\mathcal T\rangle=\langle\mathcal E_\nabla+\mathcal E\rangle\ge0\).
\item \(\langle\|Y\|_G^2\rangle=0\) if and only if \(\mathfrak m=\delta_0\); in
      that case all means vanish.  If \(\mathfrak m\ne\delta_0\), then
      \(\langle\|Y\|_G^2\rangle>0\), \(\langle\varphi\rangle>0\),
      \(\langle\mathcal E\rangle>0\), and all the inequalities in (ii) are
      strict in the trace.
\end{enumerate}
\end{theorem}

\begin{proof}
(i), (ii) The proofs of the cited identities use only the identities
along orbits and the boundedness of the observables on the family,
which hold by Proposition~\ref{prop:c2v23-family}(i); the pinching
entered only the lower bounds by \(c_\Phi\) and \(\underline\delta\).
(iii) If \(\int\mathfrak y\,\dd\mathfrak m=0\), then \(\mathfrak y(W)=0\) for
\(\mathfrak m\)-a.e.\ \(W\), and by invariance and Fubini, for
\(\mathfrak m\)-a.e.\ \(W\) the defect vanishes at a.e.\ \(\sigma\), hence at
every \(\sigma\) by continuity of \(\sigma\mapsto\|Y(\sigma)\|_G\)
(dominated convergence with the uniform bounds); such \(W\) is
backward self-similar, hence zero by
Proposition~\ref{prop:c2v23-family}(iv), so \(\mathfrak m=\delta_0\).
Conversely \(\mathfrak y(0)=0\).  If \(\mathfrak m\ne\delta_0\), the set of
nonzero elements has positive measure, on which \(\varphi>0\),
\(\mathcal E>0\) (vorticity at infinity at every time, Theorem
V63-vorticity-every-time, applies to every nonzero element of the
family), and \(\mathfrak y\) cannot vanish \(\mathfrak m\)-a.e.
\end{proof}

% =============================================================================
\section{Blow-up measures and the statistical dichotomy}
\label{sec:c2v23-dichotomy}
% =============================================================================

The empirical scale-measures \(\mathfrak m_{\lambda,L}\) of
\eqref{eq:c2v13-empirical} and the blow-up measures \(\mathcal B(x_0)\) are
defined at every \(x_0\in\Sigma\) exactly as in V13; Theorem~\ref{thm:c2v13-blowup-measures}
holds verbatim, its proof using only Proposition~\ref{prop:c2v23-family}(i),(iii).

\begin{theorem}[Statistical dichotomy]
\label{thm:c2v23-dichotomy}
Let \(x_0\in\Sigma\).  Exactly one of the following holds.
\begin{enumerate}[label=\textup{(\Roman*)},leftmargin=3.4em]
\item \emph{Statistically nontrivial.}  Some blow-up measure
      \(\mathfrak m\in\mathcal B(x_0)\) is different from \(\delta_0\); for it
      all statements of Theorem~\ref{thm:c2v23-identities}(iii) hold.
\item \emph{Statistically trivial.}  \(\mathcal B(x_0)=\{\delta_0\}\).  This
      holds if and only if for every \(\varepsilon>0\)
      \begin{equation}
       \lim_{\substack{\lambda\to0\\ L\to\infty}}\frac1L\bigl|\{\tau\in[0,L]:\Phi_{u,x_0}(\lambda e^{-\tau})\ge\varepsilon\}\bigr|=0 ,
       \label{eq:c2v23-trivial}
      \end{equation}
      that is, the active scales have zero logarithmic density; and then
      \(0\in\mathcal F(x_0)\), so \(x_0\) is lacunary with intermediate
      scales of full logarithmic density.
\end{enumerate}
A homogeneous point is statistically nontrivial.
\end{theorem}

\begin{proof}
The alternatives are exclusive and exhaustive by definition.  For the
characterization of (II): if \eqref{eq:c2v23-trivial} holds, let
\(\mathfrak m\) be a blow-up measure along \(\lambda_k,L_k\); for the
continuous observable \(\min(\varphi,\varepsilon)\), its log-average is at
most \(\varepsilon\) plus \(C_\Phi'\) times the density of the active
scales at level \(\varepsilon\), so \(\int\min(\varphi,\varepsilon)\dd\mathfrak m\le\varepsilon\)
for every \(\varepsilon\), hence \(\int\min(\varphi,\varepsilon)\dd\mathfrak m=0\)
by letting \(\varepsilon\to0\) after using monotonicity, and \(\varphi=0\)
\(\mathfrak m\)-a.e., which by Proposition~\ref{prop:c2v23-family}(iv)
means \(\mathfrak m=\delta_0\).  Conversely, if \eqref{eq:c2v23-trivial}
fails for some \(\varepsilon\), there are \(\lambda_k\to0\), \(L_k\to\infty\)
along which the density of \(\{\Phi_u\ge\varepsilon\}\) is at least
\(\eta>0\); the corresponding blow-up measure \(\mathfrak m\) satisfies
\(\int\min(\varphi,\varepsilon)\dd\mathfrak m\ge\eta\varepsilon>0\) by weak
convergence, so \(\mathfrak m\ne\delta_0\).  In case (II) the intermediate
scales have full logarithmic density by \eqref{eq:c2v23-trivial} and
Theorem V97-ledger, and \(0\in\mathcal F(x_0)\).  Homogeneous points have
\(\Phi_{u,x_0}\ge c_\Phi'\) on small scales (Theorem V96-pinched-u), so
\eqref{eq:c2v23-trivial} fails for \(\varepsilon<c_\Phi'\).
\end{proof}

\begin{assessment}[What the first cycle says about each case]
\label{ass:c2v23-first-cycle}
For a lacunary point in case (I), the nontrivial blow-up measure
inherits the exact signed statistics of homogeneous points except for
the uniform rate: \(\langle\|Y\|_G^2\rangle>0\) with no lower bound
independent of the measure, since the family contains elements
arbitrarily close to \(0\).  For a lacunary point in case (II), the
solution near the point is, at almost every log-scale, close to the
zero profile in the Gaussian sense, with the active scales forming a
set of vanishing density; the first cycle's minimal log-period
(Theorem V97-log-period) bounds the number of active--intermediate
transitions per log-window by \(2K\log\Theta/c''\), which is compatible
with (II), and its count of Gaussian-active centres (Theorem
V98-count) bounds by \(N_c\) the number of well-separated centres with
Gaussian energy at least \(c\) at any one scale, so the dissipation of
the intermediate scales, carried by the satellites (Theorem
V77-satellite), sits at each such scale in at most \(N_c\) Gaussian
cells at other centres or is spread below the Gaussian threshold.
Which case a lacunary point falls in is not decided; the series records
both.  No global regularity claim is made.
\end{assessment}

% =============================================================================
\section{Integrated V23 conclusion and next acquisition order}
\label{sec:c2v23-conclusion}
% =============================================================================

\begin{theorem}[What V23 proves]
\label{thm:c2v23-conclusion}
Relative to the first cycle and V1--V22, modulo (E1)--(E9): the
structure of the family at an arbitrary singular point
(Proposition~\ref{prop:c2v23-family}), the mean identities and signs
without homogeneity (Theorem~\ref{thm:c2v23-identities}), and the
statistical dichotomy (Theorem~\ref{thm:c2v23-dichotomy}).  The full
Navier--Stokes stability/global-regularity theorem remains unproved.
\end{theorem}

\begin{proof}
The cited results.
\end{proof}

\begin{programme}[V24 acquisition order]
\label{prog:c2v24}
\begin{enumerate}[label=\textup{(AC\arabic*)},leftmargin=4.8em]
\item Statistically trivial points and the satellites.  In case (II),
      the centred Gaussian energy has zero log-density of active scales
      while the total dissipation in every parabolic cylinder about
      \(x_0\) is of order \(\lambda\) on average by the type-I rate and the
      per-scale structure; determine whether this forces the satellites
      to be present at scales of full logarithmic density, and whether
      the moving-centre statistics (blow-up measures with centres
      following the satellites) are nontrivial: define the satellite
      family and its blow-up measures under compact satellites and
      prove they are nontrivial.
\item Global statistics of the singular set.  Since at most \(N_c\)
      well-separated centres are Gaussian-active at each scale, define
      the joint empirical measure over centres and scales and prove
      that every singular time has a nontrivial global blow-up
      statistics (some centre is active at scales of positive
      log-density), using the dissipation identity of the first cycle.
\item The next compulsory companion audit is V25.
\end{enumerate}
\end{programme}

\begin{assessment}[Position after V23]
\label{ass:c2v23-position}
The blow-up statistics are defined at every singular point, with the
exact identities and weak signs, and lacunary points split into
statistically nontrivial and trivial ones.  No full stability or
global-regularity claim is made.
\end{assessment}

% =============================================================================
\section*{Version V23 (second cycle) append-only record}
\addcontentsline{toc}{section}{Version V23 (second cycle) append-only record}
% =============================================================================

The complete V22 source of the second cycle was retained before this
continuation.  Its SHA--256 digest was
\begin{center}\scriptsize\ttfamily
7c31c7a5b3374cfeb7a25bbb266b31b5fb3ee410c1cd942e9c7fdb1f2e121c2f.
\end{center}
V23 extended the blow-up statistics to arbitrary singular points and
proved the statistical dichotomy for lacunary points.  No full
regularity claim is made.

% =============================================================================
\section{V24: compact dissipation, Gaussian activity at the parabolic scale,
and the moving-centre statistics}
\label{sec:c2v24-entry}
% =============================================================================

Items (AC1) and (AC2) of Programme~\ref{prog:c2v24} are carried out to
the extent they can be proved, and the rest is recorded.  The
dissipation of a type-I singular solution in every parabolic slab is of
the order of the parabolic scale from both sides; it is shown that a
fixed fraction of it lies in a bounded number of parabolic balls at a
scale exactly when some centres at that scale are Gaussian-active with
a fixed constant, in both directions with explicit dependence of the
constants, so that compact dissipation on a set of scales of positive
logarithmic density is equivalent to nontrivial Gaussian activity at
moving centres on such a set.  A Gaussian-net identity shows that the
dissipation is always captured by the Gaussian cells of a net of
centres at the parabolic scale, but may be spread over a number of
cells tending to infinity, which is the shell-satellite configuration
the first cycle admitted; hence global statistical nontriviality is not
a theorem.  The moving-centre blow-up measures are shown to be
invariant when the centre map is coherent in log-scale, and their
nontriviality follows under compact dissipation; neither hypothesis is
established.  No global regularity claim is made.  The next version is
the compulsory companion audit.

% =============================================================================
\section{Compact dissipation and Gaussian activity}
\label{sec:c2v24-compact}
% =============================================================================

Throughout, \(u\) satisfies the type-I rate, and for \(0<\lambda\le\lambda_0\)
let \(\Sigma_\lambda:=(T_*-S\lambda^2,T_*-\eta\lambda^2)\) be the parabolic slab
with fixed \(0<\eta<S\), and
\(\mathcal D(\lambda):=\iint_{\Torus^3\times\Sigma_\lambda}|\nabla u|^2\).

\begin{lemma}[Two-sided slab dissipation]
\label{lem:c2v24-slab}
\(2c_L^2\nu^{3/2}(S^{1/2}-\eta^{1/2})\,\lambda\le\mathcal D(\lambda)\le2C_I^2\nu^{3/2}(S^{1/2}-\eta^{1/2})\,\lambda\),
where \(c_L\) is Leray's lower-bound constant
(\(\|\Lambda u(t)\|_2\ge c_L\nu^{3/4}(T_*-t)^{-1/4}\) for a solution with
maximal time \(T_*\)) and \(C_I\) the type-I constant.
\end{lemma}

\begin{proof}
Integrate the two bounds on \(\|\Lambda u(t)\|_2^2\) over the slab.
\end{proof}

\begin{definition}[Compact dissipation]
\label{def:c2v24-compact}
The dissipation is \emph{\((\theta,N,R)\)-compact at scale \(\lambda\)} if
there are centres \(x_1,\dots,x_N\in\Torus^3\) with
\(\sum_{i=1}^N\iint_{B_{R\lambda}(x_i)\times\Sigma_\lambda}|\nabla u|^2\ge\theta\,\mathcal D(\lambda)\).
A set \(\Lambda\subset(0,\lambda_0]\) of scales has \emph{positive logarithmic
density} if \(\liminf_{\lambda\to0,\,L\to\infty}L^{-1}|\{\tau\in[0,L]:\lambda e^{-\tau}\in\Lambda\}|>0\)
along some sequences.
\end{definition}

\begin{theorem}[Compact dissipation is Gaussian activity at moving centres]
\label{thm:c2v24-equivalence}
Modulo (E1)--(E8):
\begin{enumerate}[label=\textup{(\roman*)},leftmargin=3.2em]
\item If the dissipation is \((\theta,N,R)\)-compact at scale \(\lambda\le\lambda_0\),
      then some centre \(x_i\) has
      \(\mathcal D_{x_i}(\lambda;R,S,\eta)\ge\varepsilon_0:=\theta\,2c_L^2\nu^{3/2}(S^{1/2}-\eta^{1/2})/N\),
      and by Proposition V98-moving-lower of the first cycle there is
      \(c=c(\varepsilon_0,R,S,\eta)>0\) with \(\Phi_{u,x_i}(S^{1/2}\lambda)\ge c\)
      for all \(\lambda\le\lambda_1\).
\item Conversely, if \(\Phi_{u,x}(\lambda)\ge c\) for some centre \(x\) and
      \(\lambda\le\lambda_c\), then by Proposition V98-converse of the first
      cycle \(\mathcal D_x(\lambda;R,S,\eta)\ge\varepsilon(c)\), so the
      dissipation is \((\theta,1,R)\)-compact at scale \(\lambda\) with
      \(\theta=\varepsilon(c)/(2C_I^2\nu^{3/2}(S^{1/2}-\eta^{1/2}))\).
\item Hence the dissipation is compact (with some fixed parameters) on
      a set of scales of positive logarithmic density if and only if
      there is \(c>0\) and a map \(\lambda\mapsto x(\lambda)\) with
      \(\Phi_{u,x(\lambda)}(\lambda)\ge c\) on a set of scales of positive
      logarithmic density.
\end{enumerate}
\end{theorem}

\begin{proof}
(i) The pigeonhole principle on the \(N\) balls with the lower bound of
Lemma~\ref{lem:c2v24-slab}, then the moving-centre lower bound of the
first cycle, whose constant depends only on the dissipation threshold
and the window parameters.  (ii) The quantitative converse of the first
cycle with the upper bound of Lemma~\ref{lem:c2v24-slab}.  (iii) Combine
(i) and (ii) scale by scale, with the scale shift \(\lambda\mapsto S^{1/2}\lambda\)
in (i), which preserves logarithmic density.
\end{proof}

\begin{proposition}[Gaussian-net identity]
\label{prop:c2v24-net}
Let \(\mathcal N_\lambda\subset\Torus^3\) be a maximal \(2\lambda\)-separated set,
and \(\Psi_x\) the periodized Gaussian of scale \(\lambda\) centred at \(x\).
There are constants \(0<c_1(\nu)\le c_2(\nu)\) with
\(c_1\le\sum_{x\in\mathcal N_\lambda}\Psi_x(y)\le c_2\) for all \(y\), hence
\begin{equation}
 c_1\,\mathcal D(\lambda)\ \le\ \sum_{x\in\mathcal N_\lambda}\iint_{\Sigma_\lambda}\int|\nabla u|^2\Psi_x\ \le\ c_2\,\mathcal D(\lambda),
 \qquad
 \#\mathcal N_\lambda\asymp\lambda^{-3}.
 \label{eq:c2v24-net}
\end{equation}
The Gaussian-weighted dissipation of the cells of the net captures the
slab dissipation up to constants; the number of cells is of order
\(\lambda^{-3}\), and the type-I rate alone does not prevent the
dissipation from being spread over a number of cells tending to
infinity as \(\lambda\to0\), each carrying \(o(\lambda)\).
\end{proposition}

\begin{proof}
For a maximal \(2\lambda\)-separated set, every point is within \(2\lambda\) of
some centre and the balls \(B_\lambda(x)\) are disjoint, which gives the
two-sided bound on the sum of Gaussians by comparison with the integral
of \(e^{-|z|^2/(4\nu)}\) over a lattice of spacing comparable to \(1\) in
the variable \(z=(y-x)/\lambda\); the counting is volume.  The last
sentence is the first cycle's record that shell satellites, dissipation
regions of volume of order \(\lambda^3\) spread over a shell of radius much
larger than \(\lambda\), are consistent with all its constraints
(Assessment V80 of the first cycle); such a configuration is not
\((\theta,N,R)\)-compact for any fixed \(N,R\).
\end{proof}

\begin{firewall}[Item (AC2): global statistical nontriviality is not a theorem]
\label{fw:c2v24-global}
The dissipation of a type-I singular solution is always present at the
parabolic order in every slab (Lemma~\ref{lem:c2v24-slab}) and always
captured by the Gaussian cells of a net (Proposition~\ref{prop:c2v24-net}),
but it need not be compact at any set of scales of positive logarithmic
density, and then no centre map is Gaussian-active with a fixed constant
on such a set (Theorem~\ref{thm:c2v24-equivalence}).  A singular time
all of whose singular points are statistically trivial in the sense of
Theorem~\ref{thm:c2v23-dichotomy} and whose dissipation is non-compact
at almost every log-scale is not excluded by the series; its blow-up
would be statistically invisible to every Gaussian ledger at the
parabolic scale.  Whether such a configuration is compatible with the
first cycle's satellite structure (satellites have the extent and the
rate of their scale, Proposition V78-extent-rate) is not decided: a
shell of radius \(\rho\gg\lambda\), thickness \(\lambda^3/\rho^2\), and
gradient of order \(\lambda^{-1/2}\ldots\lambda^{-2}\) satisfies the recorded
constraints.
\end{firewall}

% =============================================================================
\section{Moving-centre statistics}
\label{sec:c2v24-moving}
% =============================================================================

\begin{definition}[Coherent centre map]
\label{def:c2v24-coherent}
A map \(\lambda\mapsto x(\lambda)\in\Torus^3\) on \((0,\lambda_0]\) is
\emph{coherent} if for every \(\tau_0>0\),
\(\sup_{|\tau|\le\tau_0}|x(\lambda e^{-\tau})-x(\lambda)|/\lambda\to0\) as
\(\lambda\to0\): the centre moves by \(o(\lambda)\) across any bounded
log-window of scales around \(\lambda\).
\end{definition}

\begin{theorem}[Moving-centre blow-up measures]
\label{thm:c2v24-moving}
Let \(x(\cdot)\) be a coherent centre map, \(U_{\lambda}^{x}:=\lambda u(x(\lambda)+\lambda\,\cdot,T_*+\lambda^2\,\cdot)\),
and \(\mathfrak m^{x}_{\lambda,L}:=\frac1L\int_0^L\delta_{U^{x}_{\lambda e^{-\tau}}}\dd\tau\).
Then every weak limit of \(\mathfrak m^{x}_{\lambda_k,L_k}\), \(\lambda_k\to0\),
\(L_k\to\infty\), is a \(\Theta\)-invariant probability measure on the
compact family \(\mathcal F^{x}\) of local limits of the \(U^x_\lambda\) and
their rescaling limits; all the exact mean identities and the weak
signs of Theorem~\ref{thm:c2v23-identities} hold for it.  If moreover
\(\Phi_{u,x(\lambda)}(\lambda)\ge c\) on a set of scales of positive
logarithmic density along \(\lambda_k,L_k\), the limit is not \(\delta_0\)
and carries positive mean squared defect.
\end{theorem}

\begin{proof}
Compactness of \(\mathcal F^x\) and of the compactified orbit is as in
Lemma~\ref{lem:c2v13-compact} (the uniform bounds are centre-independent).
Invariance: \(\Theta_{\tau_0}U^x_\mu=\mu e^{\tau_0}u(x(\mu)+\mu e^{\tau_0}\,\cdot,\ldots)\)
differs from \(U^x_{\mu e^{\tau_0}}\) by the recentring
\(x(\mu)\mapsto x(\mu e^{\tau_0})\), a translation by
\((x(\mu e^{\tau_0})-x(\mu))/(\mu e^{\tau_0})\to0\) in the rescaled variable
by coherence; translations by vectors tending to zero are continuous in
\(\rho\) on the uniformly bounded family, so the argument of
Theorem~\ref{thm:c2v13-blowup-measures} applies with an additional
error tending to zero.  The identities and signs are
Theorem~\ref{thm:c2v23-identities}, whose proof uses only the class
bounds.  Nontriviality is the argument of Theorem~\ref{thm:c2v23-dichotomy}
with \(\Phi_{u,x(\lambda)}\) in place of \(\Phi_{u,x_0}\).
\end{proof}

\begin{firewall}[Item (AC1): what remains conditional]
\label{fw:c2v24-AC1}
Under compact dissipation on a set of scales of positive logarithmic
density, Theorem~\ref{thm:c2v24-equivalence} provides Gaussian-active
centres \(x(\lambda)\) on that set, but not a coherent centre map: the
active centre may jump by more than \(\lambda\) between nearby scales
(several satellites, or a satellite drifting faster than the parabolic
scale in log-scale).  The satellite blow-up measures are therefore
defined and nontrivial under the two hypotheses of compact dissipation
and coherence, neither established.  The alternative, in which every
centre map with fixed Gaussian activity is incoherent, would mean that
the singular time's dissipation moves across the parabolic scale on
every bounded log-window; the series does not exclude it.
\end{firewall}

% =============================================================================
\section{Integrated V24 conclusion and next acquisition order}
\label{sec:c2v24-conclusion}
% =============================================================================

\begin{theorem}[What V24 proves]
\label{thm:c2v24-conclusion}
Relative to the first cycle and V1--V23, modulo (E1)--(E9): the
two-sided slab dissipation, the equivalence of compact dissipation and
Gaussian activity at moving centres (Theorem~\ref{thm:c2v24-equivalence}),
the Gaussian-net identity \eqref{eq:c2v24-net}, and the moving-centre
blow-up measures under coherence (Theorem~\ref{thm:c2v24-moving}).  The
full Navier--Stokes stability/global-regularity theorem remains
unproved.
\end{theorem}

\begin{proof}
The cited results.
\end{proof}

\begin{programme}[V25 acquisition order, with compulsory companion audit]
\label{prog:c2v25}
\begin{enumerate}[label=\textup{(AC\arabic*)},leftmargin=4.8em]
\item Perform the compulsory review of the current SM and YM notes;
      record digests; import nothing numerical.
\item Consolidate V21--V24: the closure of the ledger, the statistics at
      arbitrary singular points, the statistical dichotomy, and the
      compact-dissipation equivalence with its two conditional
      hypotheses.
\item Record the position of the second cycle after twenty-five
      versions and the direction for V26--V30: the rate-free classes
      with the two Gaussian ledgers, and the question of compact
      dissipation.
\item The next compulsory companion audit after V25 is V30.
\end{enumerate}
\end{programme}

\begin{assessment}[Position after V24]
\label{ass:c2v24-position}
Compact dissipation at the parabolic scale is equivalent to Gaussian
activity at moving centres, and the moving-centre statistics need
coherence; neither is established, and non-compact, statistically
invisible blow-up is not excluded.  No full stability or
global-regularity claim is made.
\end{assessment}

% =============================================================================
\section*{Version V24 (second cycle) append-only record}
\addcontentsline{toc}{section}{Version V24 (second cycle) append-only record}
% =============================================================================

The complete V23 source of the second cycle was retained before this
continuation.  Its SHA--256 digest was
\begin{center}\scriptsize\ttfamily
08c4f401111092e896ff1a7859bf9632e483fd7a5e4c17b9039b8215ab58d5bd.
\end{center}
V24 proved the equivalence of compact dissipation and Gaussian
activity at moving centres, the Gaussian-net identity, and the
moving-centre blow-up measures under coherence, and recorded that
global statistical nontriviality is not a theorem.  No full regularity
claim is made.

% =============================================================================
\section{V25: compulsory companion audit, consolidation of V21--V24, and the
position after twenty-five versions}
\label{sec:c2v25-entry}
% =============================================================================

This is the twenty-fifth version of the second cycle.  The compulsory
review of the two companion programmes is performed first.  V21--V24
are then consolidated: the closure of the moment ledger, the blow-up
statistics at arbitrary singular points with the statistical dichotomy
for lacunary points, and the equivalence of compact dissipation with
Gaussian activity at moving centres, together with the two conditional
hypotheses (compact dissipation, coherent centres) on which the
moving-centre statistics rest.  The position after twenty-five versions
and the direction for V26--V30 are recorded.  No global regularity
claim is made.  The next compulsory companion audit is V30.

% =============================================================================
\section{Compulsory V25 review of the current companion notes}
\label{sec:c2v25-companion-audit}
% =============================================================================

The current companion files in \texttt{latex\_notes} were inspected
read-only.  The frozen checkpoints are
\begin{align}
 &\texttt{Renewal\_SM\_Einstein\_Continuum\_Research\_Notes\_V90.tex},
 \notag\\[-2pt]
 &\qquad
 \texttt{d822f8d46f0619a01e352d8f98a9c085f9caebf3cb25f4dc440f9141abb3c784},
 \label{eq:c2v25-SM-checkpoint}\\
 &\texttt{Renewal\_Yang\_Mills\_Mass\_Gap\_Research\_Notes\_V67.tex},
 \notag\\[-2pt]
 &\qquad
 \texttt{1881a8c884773b89469c6a6f7c64da3c9804bbce908fded509d015284d17feef}.
 \label{eq:c2v25-YM-checkpoint}
\end{align}
The SM file is byte-identical to the checkpoint \eqref{eq:c2v20-SM-checkpoint}
of V20; its programme has not advanced since.  A YM file named
\(\ldots\)\texttt{V66.tex} (digest \texttt{2d1ecedb\ldots}) was also
present; the V67 file contains it and is recorded.

\emph{YM V64--V67.}  The YM programme proved a multivariate
Mecke--cumulant formula, performed its own compulsory audit, corrected
an inherited inner finite-difference estimate, proved a fixed-partition
inner clock grade and a block-factorial partition generating function,
closed its complete inner partition sum, reduced three of its cards to
one outer card, proved a low-centre outer-tail bound, and then proved
uniform Gamma domination of block-weighted jumps, a partition-sum Gamma
tail, an outer added-jump threshold estimate, and a polynomially marked
Gamma-tail estimate, isolating an exact low/high-order crossover as its
remaining gate.

\begin{theorem}[V25 companion-audit decision]
\label{thm:c2v25-companion-decision}
No YM V64--V67 cumulant, partition, or Gamma-tail estimate is imported,
and nothing new is available from SM.  Neither bears on the closure of
the ledger (Theorem~\ref{thm:c2v22-closed}), the statistical dichotomy
(Theorem~\ref{thm:c2v23-dichotomy}), or the compact-dissipation
equivalence (Theorem~\ref{thm:c2v24-equivalence}).
\end{theorem}

\begin{proof}
The YM results concern Poisson-type cumulant expansions of a lattice
clock; none is a Navier--Stokes statement.
\end{proof}

\begin{assessment}[Direction after the audit]
\label{ass:c2v25-direction}
The YM programme's V66 corrected an inherited estimate in-series, as
this series did in V17 and V21; the discipline is common.  No change of
direction follows from the audit.
\end{assessment}

% =============================================================================
\section{Consolidation of V21--V24}
\label{sec:c2v25-consolidated}
% =============================================================================

\begin{theorem}[Ledger closure, statistics without homogeneity, and compact dissipation]
\label{thm:c2v25-consolidated}
Let \(u\) satisfy the type-I rate.  Modulo (E1)--(E9):
\begin{enumerate}[label=\textup{(\arabic*)},leftmargin=3.2em]
\item \emph{Closure.}  \(R=8\nu\mathcal L-4\nu^2\mathcal E\) pointwise, so
      \(\mathcal E\le\frac2\nu\mathcal L\) with equality for tangential profiles;
      the Gaussian stress splits as \(S=S_0+S'\) with
      \(\langle S_0\rangle=-\langle T_0\rangle\preceq0\) and
      \(\langle S'\rangle=-\langle D+T'\rangle\preceq0\); the moment ledger is
      closed for the purposes of the series.
\item \emph{Statistics at any singular point.}  At every \(x_0\in\Sigma\),
      \(\mathcal F(x_0)\) is a compact scale-invariant family on which the
      rescaling flow acts continuously, with \(0\in\mathcal F(x_0)\) if and
      only if \(x_0\) has intermediate scales; every invariant measure
      satisfies all exact mean identities and the weak signs, and
      \(\delta_0\) is the only invariant measure with vanishing mean
      squared defect.
\item \emph{Statistical dichotomy.}  Every singular point is
      statistically nontrivial (some blow-up measure differs from
      \(\delta_0\), with positive mean squared defect and strict signs) or
      statistically trivial (every blow-up measure is \(\delta_0\),
      equivalently the active scales have zero logarithmic density);
      homogeneous points are nontrivial; for lacunary points both cases
      are open.
\item \emph{Compact dissipation.}  The slab dissipation is of the
      parabolic order from both sides; it is compact at a set of scales
      of positive logarithmic density if and only if some centre map is
      Gaussian-active with a fixed constant on such a set; the
      Gaussian cells of a net always capture it, possibly spread over a
      number of cells tending to infinity; under compact dissipation and
      a coherent centre map the moving-centre blow-up measures are
      invariant and nontrivial.
\end{enumerate}
\end{theorem}

\begin{proof}
(1) Theorems~\ref{thm:c2v21-closure}, \ref{thm:c2v22-split},
\ref{thm:c2v22-closed}.  (2) Proposition~\ref{prop:c2v23-family},
Theorem~\ref{thm:c2v23-identities}.  (3) Theorem~\ref{thm:c2v23-dichotomy}.
(4) Lemma~\ref{lem:c2v24-slab}, Theorem~\ref{thm:c2v24-equivalence},
Proposition~\ref{prop:c2v24-net}, Theorem~\ref{thm:c2v24-moving}.
\end{proof}

% =============================================================================
\section{Position after twenty-five versions and direction}
\label{sec:c2v25-position}
% =============================================================================

\begin{assessment}[Position after twenty-five versions of the second cycle]
\label{ass:c2v25-position-twenty-five}
The second cycle has reduced the homogeneous type-I case to a single
mean inequality against a blow-up measure (Theorem~\ref{thm:c2v15-gates}),
has closed the search for a contradiction within the Gaussian moment
ledger (Theorem~\ref{thm:c2v22-closed}), and has extended the
statistical framework to every singular point, where it has found a
new alternative: blow-up that is statistically invisible at the
parabolic scale, through non-compact dissipation.  Three hypotheses
now organize the type-I case: homogeneity (or statistical
nontriviality), compact dissipation, and coherence of the active
centre.  Under all three the blow-up has exact, signed, scale-invariant
statistics; the series has no contradiction from them.  Without them
the blow-up escapes the Gaussian ledgers.  For V26--V30 the direction
is upstream to the rate-free classes: the two Gaussian ledgers hold
for every smooth solution, and the question is what replaces
compactness of the family when no rate is available; the natural
object is the family of local limits of rescalings under the rate-free
local energy bound of the first cycle (Theorem V66-local-energy), which
gives boundedness in the critical Morrey class only on scales at which
the Leray energy spent in the parabolic window is of the parabolic
order.  No global regularity claim is made.
\end{assessment}

% =============================================================================
\section{Integrated V25 conclusion and next acquisition order}
\label{sec:c2v25-conclusion}
% =============================================================================

\begin{theorem}[What V25 establishes]
\label{thm:c2v25-conclusion}
Relative to the first cycle and V1--V24: the companion-audit decision
(Theorem~\ref{thm:c2v25-companion-decision}) and the consolidation
(Theorem~\ref{thm:c2v25-consolidated}).  The full Navier--Stokes
stability/global-regularity theorem remains unproved.
\end{theorem}

\begin{proof}
The cited results.
\end{proof}

\begin{programme}[V26 acquisition order]
\label{prog:c2v26}
\begin{enumerate}[label=\textup{(AC\arabic*)},leftmargin=4.8em]
\item Parabolic scales without a rate.  For a smooth solution with
      \(T_*<\infty\) and no rate assumption, call a scale \(\lambda\)
      \emph{parabolic} at \(x_0\) if the Leray energy spent in the window
      \((T_*-2\lambda^2,T_*)\) is at most \(M\lambda\) and the rescaling
      \(U_\lambda\) satisfies the local bounds needed for compactness on
      the unit cylinder; determine, from the first cycle's rate-free
      results (V66--V69) and the CKN criterion (E1), what bounds hold
      for \(U_\lambda\) at parabolic scales, and whether the set of local
      limits along parabolic scales is compact.
\item Determine whether the two Gaussian ledgers, restricted to
      parabolic scales, give the analogue of the pinching or of the
      inward flux at a point where parabolic scales have positive
      logarithmic density.
\item The next compulsory companion audit is V30.
\end{enumerate}
\end{programme}

\begin{assessment}[Position after V25]
\label{ass:c2v25-position}
The second cycle is consolidated through V24 and turns upstream to the
rate-free classes.  No full stability or global-regularity claim is
made.
\end{assessment}

% =============================================================================
\section*{Version V25 (second cycle) append-only record}
\addcontentsline{toc}{section}{Version V25 (second cycle) append-only record}
% =============================================================================

The complete V24 source of the second cycle was retained before this
continuation.  Its SHA--256 digest was
\begin{center}\scriptsize\ttfamily
58da2b5834a11a9491359801da2ca2f71dd7b4be745a5d0f89d59a25fa79d186.
\end{center}
V25 performed the compulsory companion audit, consolidated V21--V24,
and recorded the position after twenty-five versions.  No full
regularity claim is made.

% =============================================================================
\section{V26: parabolic scales without a rate, compactness to suitable weak
solutions, and the Gaussian ledgers at such scales}
\label{sec:c2v26-entry}
% =============================================================================

Items (AC1) and (AC2) of Programme~\ref{prog:c2v26} are carried out.
Without any rate, a scale is called parabolic at a point if the
velocity, the dissipation, and the mean-subtracted pressure of the
solution are of the parabolic order in the parabolic cylinder about
the point at that scale; it is shown, using the first cycle's rate-free
bounds, that dissipation of parabolic order alone does not imply energy
of parabolic order, so the three conditions are independent inputs.
Along parabolic scales the rescalings are bounded in the natural
Caffarelli--Kohn--Nirenberg norms and converge, along subsequences, to
suitable weak solutions on the unit cylinder; the limit may be trivial,
since without a rate the dissipation may escape to high frequencies.
The Gaussian ledgers at parabolic scales are recorded: the weighted
energy is bounded at such scales when the local Morrey condition holds
at radii growing like the square root of the logarithm of the scale,
and between two such scales the first identity gives the inward flux
at rate twice the lower bound of the weighted energy on the window,
conditionally on Gaussian activity across it.  No global regularity
claim is made.  The next compulsory companion audit is V30.

% =============================================================================
\section{Parabolic scales}
\label{sec:c2v26-scales}
% =============================================================================

Let \(u\) be a smooth mean-zero solution on \([0,T_*)\times\Torus^3\) with
\(T_*<\infty\), \(x_0\in\Torus^3\), and for \(0<\lambda\le\pi\) let
\(Q_\lambda:=B_{2\lambda}(x_0)\times(T_*-4\lambda^2,T_*)\) and \(\bar p_\lambda(t)\) the mean
of \(p(t)\) over \(B_{2\lambda}(x_0)\).

\begin{definition}[Parabolic scales]
\label{def:c2v26-parabolic}
Fix \(M>0\).  The scale \(\lambda\) is \emph{\(M\)-parabolic at \(x_0\)} if
\begin{equation}
 \sup_{T_*-4\lambda^2<t<T_*}\int_{B_{2\lambda}(x_0)}|u(t)|^2\le M\lambda,
 \qquad
 \iint_{Q_\lambda}|\nabla u|^2\le M\lambda,
 \qquad
 \iint_{Q_\lambda}|p-\bar p_\lambda|^{3/2}\le M\lambda^2 .
 \label{eq:c2v26-parabolic}
\end{equation}
Under the type-I rate every small scale is \(M\)-parabolic at every point
for a fixed \(M\) (the \(L^\infty\) rate, the type-I dissipation bound, and
the pressure bound of the first cycle).
\end{definition}

\begin{proposition}[The three conditions are independent inputs]
\label{prop:c2v26-independent}
Without a rate, the dissipation condition alone gives only
\(\sup_{t}\int_{B_{2\lambda}}|u(t)|^2\le C_0\bigl(4M\lambda+(2\lambda)^{-1/2}E_*^{3/4}(4M\lambda)^{3/4}\bigr)\le C(M,E_*)\lambda^{1/4}\)
(Theorem V66-local-energy of the first cycle with \(\epsilon\le\nu\iint_{Q_\lambda}|\nabla u|^2\)),
and even with the energy condition the cubic term of the local energy
inequality is of order \(\lambda^{3/4}\), not \(\lambda\); the pressure
condition is not implied by the other two, the global bound
\(\|p-\bar p\|_{L^{3/2}(\Torus^3)}\le C\|u\|_3^2\) giving only
\(\iint_{Q_\lambda}|p-\bar p_\lambda|^{3/2}\le C\lambda^{5/4}\), which exceeds
\(M\lambda^2\) by the factor \(\lambda^{-3/4}\).
\end{proposition}

\begin{proof}
The first bound is Theorem V66-local-energy with \(\rho=2\lambda\) and
\(\epsilon(t;2\rho^2)\le\nu\iint_{Q_\lambda}|\nabla u|^2\le\nu M\lambda\) (up to
the factor \(\nu\) absorbed in \(M\)); the second term is
\(\lambda^{-1/2}\lambda^{3/4}=\lambda^{1/4}\).  For the cubic term, with the local
energy bounded by \(M\lambda\),
\(\int_{B_{4\lambda}}|u|^3\le(M\lambda)^{1/2}\|u\|_4^2\le C(M\lambda)^{1/2}E_*^{1/4}\|\Lambda u\|_2^{3/2}\)
by the Ladyzhenskaya--Agmon interpolation on the torus, and integrating
over the window with H\"older,
\(\lambda^{-1}\int_{T_*-4\lambda^2}^{T_*}\int_{B_{4\lambda}}|u|^3\le C\lambda^{-1}(M\lambda)^{1/2}E_*^{1/4}(M\lambda)^{3/4}(4\lambda^2)^{1/4}=C(M,E_*)\lambda^{3/4}\)
(Proposition V68-bookkeeping of the first cycle with \(\alpha=\frac12\),
\(\beta=1\): \(\rho^{-1/2+3\alpha/2+\beta/2}=\rho^{3/4}\)).  For the pressure,
the Calder\'on--Zygmund bound on the torus and the time integral of
\(\|u\|_3^3\le E_*^{3/4}\|\Lambda u\|_2^{3/2}\) over the window give
\(\iint_{Q_\lambda}|p-\bar p|^{3/2}\le\int_{T_*-4\lambda^2}^{T_*}\|p-\bar p\|_{3/2}^{3/2}\le CE_*^{3/4}(M\lambda)^{3/4}\lambda^{1/2}=C\lambda^{5/4}\),
which is \(\lambda^2\cdot\lambda^{-3/4}\): not of the parabolic order.
\end{proof}

\begin{remark}
\label{rem:c2v26-independent}
In Proposition~\ref{prop:c2v26-independent} the local part of the
pressure is controlled by the local cubic quantity, the harmonic part
only by the global bound (Firewall V68-harmonic of the first cycle).
\end{remark}

% =============================================================================
\section{Compactness along parabolic scales}
\label{sec:c2v26-compactness}
% =============================================================================

\begin{theorem}[Rescalings at parabolic scales converge to suitable weak solutions]
\label{thm:c2v26-compactness}
Let \(\lambda_k\to0\) be \(M\)-parabolic scales at \(x_0\) and
\(U_k:=\lambda_ku(x_0+\lambda_k\,\cdot,T_*+\lambda_k^2\,\cdot)\),
\(P_k:=\lambda_k^2(p-\bar p_{\lambda_k})(x_0+\lambda_k\,\cdot,T_*+\lambda_k^2\,\cdot)\) on
\(Q:=B_2\times(-4,0)\).  Then
\begin{equation}
 \sup_{s}\int_{B_2}|U_k(s)|^2\le M,\qquad
 \iint_Q|\nabla U_k|^2\le M,\qquad
 \iint_Q|P_k|^{3/2}\le M,\qquad
 \iint_Q|U_k|^3\le C(M),
 \label{eq:c2v26-bounds}
\end{equation}
and, modulo (E2), a subsequence satisfies \(U_k\to U\) strongly in
\(L^3(Q')\) for every \(Q'\Subset Q\), \(\nabla U_k\rightharpoonup\nabla U\)
weakly in \(L^2(Q)\), \(P_k\rightharpoonup P\) weakly in \(L^{3/2}(Q)\), where
\((U,P)\) is a suitable weak solution of the Navier--Stokes equations
on \(Q\): \(U\in L^\infty_sL^2\cap L^2_sH^1\), \(P\in L^{3/2}\), the equations
hold in the sense of distributions, and the local energy inequality
holds on \(Q\).
\end{theorem}

\begin{proof}
The bounds are \eqref{eq:c2v26-parabolic} rescaled (the scaling
\(u\mapsto\lambda u(\lambda\cdot,\lambda^2\cdot)\) maps the three quantities to
themselves divided by \(\lambda\), \(\lambda\), \(\lambda^2\)), and the cubic bound
is the interpolation \(\|U\|_{L^3(Q)}^3\le C\|U\|_{L^\infty L^2}^{3/2}\|U\|_{L^2H^1}^{3/2}\)
on the bounded cylinder.  From the equation,
\(\partial_sU_k=\nu\Delta U_k-\operatorname{div}(U_k\otimes U_k)-\nabla P_k\) is
bounded in \(L^{3/2}(-4,0;H^{-2}(B_2))\), so Aubin--Lions (E2), in the
form bounded in \(L^2H^1\) with time derivatives bounded in
\(L^{3/2}H^{-2}\), gives strong convergence in \(L^2(Q')\), and with the
\(L^\infty L^2\cap L^2L^6\) bounds, in \(L^3(Q')\) by interpolation.
Weak convergence of gradients and pressures follows from the bounds.
The limit satisfies the equations in the sense of distributions
(the nonlinear term converges by the strong \(L^3\) convergence), and
the local energy inequality passes to the limit by lower
semicontinuity of the dissipation and strong convergence of the cubic
and pressure-velocity terms (the latter as the pairing of a weakly
convergent \(P_k\) with a strongly convergent \(U_k\phi\)).
\end{proof}

\begin{firewall}[Item (AC1): the limit may be trivial and the family is not scale-invariant]
\label{fw:c2v26-trivial}
Without a rate, \(\nabla U_k\) converges only weakly, so the dissipation
of the limit may be strictly smaller than the liminf of the rescaled
dissipations, and the limit may vanish even though \(x_0\) is singular
(the dissipation of parabolic order in \(Q_{\lambda_k}\), which (E1)
guarantees along a sequence of scales at every singular point, may
concentrate at frequencies tending to infinity in the rescaled
variables).  The first cycle's strong \(H^1\) convergence used the
type-I derivative bounds.  Moreover the set of \(M\)-parabolic scales is
not closed under multiplication by fixed factors, so the family of
limits is not scale-invariant, and no flow, invariant measure, or
dichotomy of the second cycle applies to it.  What survives is the
statement of Theorem~\ref{thm:c2v26-compactness}: at every singular
point admitting parabolic scales, the blow-up is described along them
by suitable weak solutions on the unit cylinder, possibly zero.
\end{firewall}

% =============================================================================
\section{The Gaussian ledgers at parabolic scales}
\label{sec:c2v26-ledgers}
% =============================================================================

\begin{proposition}[Weighted energy at parabolic scales]
\label{prop:c2v26-weighted}
Let \(\lambda\) be a scale at which the local Morrey condition
\(\int_{B_{r}(x_0)}|u(T_*-\lambda^2)|^2\le M r\) holds for all
\(\lambda\le r\le R_\lambda\lambda\) with \(R_\lambda:=(8\nu\log(1/\lambda))^{1/2}\).  Then
\(\Phi_{u,x_0}(\lambda)\le C(\nu)M+E_*\lambda\), where the second term is the
Gaussian tail beyond radius \(R_\lambda\lambda\).  Conversely
\(\Phi_{u,x_0}(\lambda)\le M'\) implies \(\int_{B_{r}}|u(T_*-\lambda^2)|^2\le M'e^{r^2/(4\nu\lambda^2)}\lambda\)
for every \(r\), which is of the parabolic order for \(r\le R\lambda\) with
\(R\) fixed.
\end{proposition}

\begin{proof}
Shell decomposition as in Lemma V86-weight of the first cycle up to
radius \(R_\lambda\lambda\), and \(\lambda^{-1}E_*e^{-R_\lambda^2/(4\nu)}=\lambda^{-1}E_*\lambda^2=E_*\lambda\)
for the tail.  The converse is \(\Psi\ge e^{-r^2/(4\nu\lambda^2)}\) on \(B_r\).
\end{proof}

\begin{theorem}[Rate-free inward flux between parabolic scales]
\label{thm:c2v26-inward}
Let \(0<a<b\) be scales at which \(\Phi_{u,x_0}\le M'\), and suppose
\(\Phi_{u,x_0}(\lambda)\ge c\) for all \(\lambda\in[a,b]\) (Gaussian activity
across the window).  Then, without any rate,
\begin{equation}
 \boxed{\quad
 \int_a^b\mathcal J_{u,x_0}(\lambda)\,\dd\lambda\ \le\ M'-2c\log\frac ba-\int_a^bD_{u,x_0}(\lambda)\,\dd\lambda ,
 \quad}
 \label{eq:c2v26-inward}
\end{equation}
and \(a^{-2}\Phi_{u}(a)-b^{-2}\Phi_u(b)=\int_a^b\mathcal X_{u}\lambda'^{-2}\dd\lambda'\ge ca^{-2}-M'b^{-2}\).
\end{theorem}

\begin{proof}
The rate-free first identity (Proposition V96-rate-free of the first
cycle) integrated over \([a,b]\): \(\Phi_u(b)-\Phi_u(a)\le M'\) and
\(\int_a^b2\Phi_u/\lambda\ge2c\log(b/a)\); the second statement is the
ledger form of \eqref{eq:c2v1-identity} integrated over \([a,b]\), with \(\Phi_u(a)\ge c\), \(\Phi_u(b)\le M'\).
\end{proof}

\begin{assessment}[Item (AC2)]
\label{ass:c2v26-AC2}
The analogue of the inward flux holds rate-free on any window of scales
bounded by parabolic scales and Gaussian-active throughout, at the rate
\(2c\) per e-fold beyond the dissipation.  The analogue of the pinching,
a lower bound \(c\) on the weighted energy across the window, is not a
consequence of anything rate-free: at a singular point (E1) gives
dissipation of parabolic order along a sequence of scales, and the
first cycle's moving-centre lower bound converts it into Gaussian
activity only through the compactness with strong \(H^1\) convergence,
which needs a rate.  The rate-free ledgers therefore remain conditional
on Gaussian activity, and the second cycle records this as the exact
boundary of what its tools give without a rate.  No global regularity
claim is made.
\end{assessment}

% =============================================================================
\section{Integrated V26 conclusion and next acquisition order}
\label{sec:c2v26-conclusion}
% =============================================================================

\begin{theorem}[What V26 proves]
\label{thm:c2v26-conclusion}
Relative to the first cycle and V1--V25, without a rate and modulo
(E1)--(E2): the independence of the three parabolic conditions
(Proposition~\ref{prop:c2v26-independent}), the compactness to suitable
weak solutions along parabolic scales (Theorem~\ref{thm:c2v26-compactness}),
the weighted-energy bound at parabolic scales
(Proposition~\ref{prop:c2v26-weighted}), and the rate-free inward flux
\eqref{eq:c2v26-inward}.  The full Navier--Stokes
stability/global-regularity theorem remains unproved.
\end{theorem}

\begin{proof}
The cited results.
\end{proof}

\begin{programme}[V27 acquisition order]
\label{prog:c2v27}
\begin{enumerate}[label=\textup{(AC\arabic*)},leftmargin=4.8em]
\item Frequency escape.  Determine whether the trivial-limit
      alternative of Firewall~\ref{fw:c2v26-trivial} can be
      characterized: if the rescalings along parabolic scales converge
      to zero in \(L^3\) while their dissipation on the unit cylinder
      stays at least \(\varepsilon_1\), then the dissipation lives at
      rescaled frequencies tending to infinity; relate this to the
      first cycle's ledger classification (the \(L^{5/2}_tH^1\) regime)
      and to the necessity of a rate for \(\varepsilon\)-regularity (V69).
\item Rate-free Gaussian activity.  Determine whether a lower bound on
      \(\Phi_{u,x_0}(\lambda)\) along parabolic scales can be obtained from
      the dissipation of parabolic order and the local energy
      inequality with the Gaussian weight alone, without compactness;
      record the exact obstruction if not.
\item The next compulsory companion audit is V30.
\end{enumerate}
\end{programme}

\begin{assessment}[Position after V26]
\label{ass:c2v26-position}
Without a rate the blow-up is described along parabolic scales by
suitable weak solutions, possibly trivial, and the Gaussian ledgers are
conditional on Gaussian activity.  No full stability or
global-regularity claim is made.
\end{assessment}

% =============================================================================
\section*{Version V26 (second cycle) append-only record}
\addcontentsline{toc}{section}{Version V26 (second cycle) append-only record}
% =============================================================================

The complete V25 source of the second cycle was retained before this
continuation.  Its SHA--256 digest was
\begin{center}\scriptsize\ttfamily
d40b956e06277409ee62bbcb7941f26c995bc135eb3073a01f6d1ea3cc4e64f7.
\end{center}
V26 defined parabolic scales without a rate, proved the compactness of
rescalings along them to suitable weak solutions, and recorded the
rate-free Gaussian ledgers at such scales.  No full regularity claim is
made.

% =============================================================================
\section{V27: frequency escape, Gaussian invisibility, and what the Dirichlet
ledger still sees}
\label{sec:c2v27-entry}
% =============================================================================

Items (AC1) and (AC2) of Programme~\ref{prog:c2v27} are carried out.  The
trivial-limit alternative along parabolic scales is characterized as
frequency escape: the rescaled dissipation, bounded below by the
Caffarelli--Kohn--Nirenberg threshold along a sequence of scales at
every singular point, is carried, in the limit, entirely by spatial
frequencies tending to infinity in the rescaled variables, while the
rescaled energy on every compact set tends to zero.  Under frequency
escape the Gaussian-weighted energy at the parabolic scale tends to
zero: such a blow-up is invisible to the first Gaussian ledger.  It is
shown that the Gaussian Dirichlet functional of the second ledger still
sees dissipation of parabolic order in the backward paraboloid: its
log-integral over the scales below a paraboloid-dissipative scale is
bounded below by a fixed multiple of the dissipation constant, so that
the rate-free second identity applies with a nontrivial left side.  A rate-free lower
bound on the weighted energy is shown to be impossible from the ledgers
alone, frequency escape being the exact obstruction.  No global
regularity claim is made.  The next compulsory companion audit is V30.

% =============================================================================
\section{Frequency escape}
\label{sec:c2v27-escape}
% =============================================================================

Throughout, \(u\) is a smooth mean-zero solution with \(T_*<\infty\) and
\(x_0\in\Sigma\) is a singular point.  By (E1), there are scales
\(r_k\to0\) with \(r_k^{-1}\iint_{Q_{r_k}(x_0,T_*)}|\nabla u|^2\ge\varepsilon_1\),
where \(Q_r(x_0,T_*)=B_r(x_0)\times(T_*-r^2,T_*)\); call such scales
\emph{dissipative}.

\begin{proposition}[Frequency escape]
\label{prop:c2v27-escape}
Let \(\lambda_k\to0\) be dissipative scales at \(x_0\) that are also
\(M\)-parabolic, \(U_k\) the rescalings, and suppose that along a
subsequence the limit of Theorem~\ref{thm:c2v26-compactness} is
\(U\equiv0\) on \(Q=B_2\times(-4,0)\).  Then:
\begin{enumerate}[label=\textup{(\roman*)},leftmargin=3.2em]
\item \(U_k\to0\) strongly in \(L^3(Q')\) and in \(L^2(Q')\) for every
      \(Q'\Subset Q\), while \(\iint_{Q_1}|\nabla U_k|^2\ge\varepsilon_1\) for
      all \(k\), \(Q_1=B_1\times(-1,0)\).
\item For every mollifier \(\rho_\delta\) at a fixed spatial scale
      \(\delta>0\), \(\rho_\delta*\nabla U_k\to0\) strongly in \(L^2(Q')\); hence
      for every \(\delta>0\),
      \(\liminf_k\iint_{Q_1}|\nabla U_k-\rho_\delta*\nabla U_k|^2\ge\varepsilon_1\):
      the rescaled dissipation is carried, asymptotically, by spatial
      frequencies larger than any fixed \(1/\delta\).
\item In the original variables: the dissipation
      \(\varepsilon_1\lambda_k\) in \(Q_{\lambda_k}(x_0,T_*)\) is carried by
      spatial frequencies \(\gg\lambda_k^{-1}\), and the energy in
      \(B_{2\lambda_k}(x_0)\) at every time of the window is
      \(o(\lambda_k)\).
\end{enumerate}
\end{proposition}

\begin{proof}
(i) is Theorem~\ref{thm:c2v26-compactness} with \(U=0\) and the
definition of dissipative scales.  (ii) \(\nabla U_k\rightharpoonup0\) weakly
in \(L^2(Q)\), and mollification at a fixed scale is a compact operator
from \(L^2(Q)\) to \(L^2(Q')\) (it maps bounded sets to sets bounded in
\(L^2_sH^1_x\) on the smaller set, and the weak limit of a weakly
convergent sequence is preserved), so \(\rho_\delta*\nabla U_k\to0\)
strongly.  Then \(\|\nabla U_k-\rho_\delta*\nabla U_k\|_{L^2(Q_1)}\ge\|\nabla U_k\|_{L^2(Q_1)}-o(1)\ge\varepsilon_1^{1/2}-o(1)\).
(iii) is the rescaling.
\end{proof}

\begin{assessment}[Relation to the first cycle]
\label{ass:c2v27-first-cycle}
Frequency escape is the spatial form of the first cycle's enstrophy
spikes: dissipation of parabolic order in a parabolic cylinder carried
by frequencies much larger than the inverse parabolic scale, with
energy much smaller than the parabolic order.  The first cycle showed
(Proposition V69-two-sided) that a rate-free hypothesis cannot supply
the smallness needed by (E5) at a singular point, because the cubic
quantity is bounded below by the Leray rate; Proposition~\ref{prop:c2v27-escape}
shows what the failure of (E5) looks like along parabolic scales
without a rate: the singularity disappears from every compact set in
the rescaled variables and survives only as high-frequency
dissipation.  The type-I rate excludes this by its \(L^\infty\) and
gradient bounds, which force the dissipation to live at frequencies of
order \(\lambda^{-1}\).
\end{assessment}

% =============================================================================
\section{Gaussian invisibility and the Dirichlet ledger}
\label{sec:c2v27-ledgers}
% =============================================================================

\begin{proposition}[Gaussian invisibility under frequency escape]
\label{prop:c2v27-invisible}
Let \(\lambda_k\to0\) be scales at \(x_0\) such that every scale
\(2^j\lambda_k\) with \(2^j\le R_{\lambda_k}=(8\nu\log(1/\lambda_k))^{1/2}\) is
\(M\)-parabolic, and suppose that for every fixed \(R\) the rescalings
\(U_k\) converge to zero strongly in \(L^2(B_R\times(-4,0))\) (frequency
escape on every compact set, which is the conclusion of
Proposition~\ref{prop:c2v27-escape} applied at the scales \(2^j\lambda_k\)).
Then \(\Phi_{u,x_0}(\lambda_k)\to0\): the first Gaussian ledger does not see
the blow-up along these scales.
\end{proposition}

\begin{proof}
Fix \(R=2^{j}\).  The parabolic bounds on \(B_{2R}\times(-4,0)\) give,
through the local energy identity, a uniform Lipschitz bound on
\(s\mapsto\int|U_k(s)|^2\phi_R\) for a cutoff \(\phi_R\), so the
\(L^1\)-in-time convergence to zero implies \(\int_{B_R}|U_k(-1)|^2\to0\).
Hence the part of \(\Phi_{u,x_0}(\lambda_k)\) from \(B_{R\lambda_k}\) tends to
zero.  The shells \(2^i\lambda_k\le|x-x_0|<2^{i+1}\lambda_k\) with
\(R\le2^i\le R_{\lambda_k}\) contribute at most
\(\sum_{2^i\ge R}M2^{i+1}e^{-4^{i}/(4\nu)}\le C(\nu)Me^{-R^2/(8\nu)}\) by the
parabolic energy bound at those scales, and the far tail beyond
\(R_{\lambda_k}\lambda_k\) is at most \(E_*\lambda_k\).  Thus
\(\limsup_k\Phi_{u,x_0}(\lambda_k)\le C(\nu)Me^{-R^2/(8\nu)}\) for every \(R\),
which gives the claim.
\end{proof}

Let \(P_\lambda(x_0):=\{(x,t):T_*-\lambda^2<t<T_*,\ |x-x_0|<2(T_*-t)^{1/2}\}\) be
the backward paraboloid of scale \(\lambda\) at \(x_0\); it is contained in
\(Q_{2\lambda}(x_0,T_*)\).  Call \(\lambda\) \emph{paraboloid-dissipative} at
\(x_0\) with constant \(\varepsilon'\) if \(\iint_{P_\lambda(x_0)}|\nabla u|^2\ge\varepsilon'\lambda\).

\begin{theorem}[The Dirichlet ledger sees paraboloid dissipation without a rate]
\label{thm:c2v27-dirichlet-sees}
If \(\lambda\) is paraboloid-dissipative at \(x_0\) with constant \(\varepsilon'\),
then, without any rate,
\begin{equation}
 \boxed{\quad
 \int_0^{\lambda}\mathcal L_{u,x_0}(\lambda')\,\frac{\dd\lambda'}{\lambda'}
 \ \ge\ \frac{\nu e^{-1/\nu}}4\,\varepsilon' .
 \quad}
 \label{eq:c2v27-dirichlet-sees}
\end{equation}
Consequently, if \(x_0\) is paraboloid-dissipative along a sequence of
scales \(\lambda_k\to0\), then for every \(k\) the log-integral of the
Gaussian Dirichlet functional over the scales below \(\lambda_k\) is at
least \(\frac\nu4e^{-1/\nu}\varepsilon'\), and the rate-free second identity
\eqref{eq:c2v14-second} applies with a left side that is not
identically small, whether or not the blow-up is Gaussian-visible.
\end{theorem}

\begin{proof}
\(\mathcal L_u(\lambda')\ge\frac18\lambda'D_u(\lambda')=\frac\nu2\lambda'\int|\nabla u(T_*-\lambda'^2,x)|^2\Psi(x,\lambda'^2)\dd x\),
and on \(\{|x-x_0|<2\lambda'\}\), \(\Psi\ge e^{-4\lambda'^2/(4\nu\lambda'^2)}=e^{-1/\nu}\).
With \(t=T_*-\lambda'^2\), \(\dd\lambda'/\lambda'=\dd t/(2\lambda'^2)=\dd t/(2(T_*-t))\), so
\[
 \int_0^\lambda\mathcal L_u\frac{\dd\lambda'}{\lambda'}
 \ge\frac\nu2e^{-1/\nu}\int_{T_*-\lambda^2}^{T_*}\frac{(T_*-t)^{1/2}}{2(T_*-t)}\int_{|x-x_0|<2(T_*-t)^{1/2}}|\nabla u(t,x)|^2\dd x\,\dd t
 \ge\frac\nu4e^{-1/\nu}\,\frac1\lambda\iint_{P_\lambda(x_0)}|\nabla u|^2 ,
\]
using \((T_*-t)^{-1/2}\ge\lambda^{-1}\) on the window; the last integral is
at least \(\varepsilon'\lambda\).
\end{proof}

\begin{firewall}[Cylinders versus paraboloids]
\label{fw:c2v27-paraboloid}
(E1) makes every singular point cylinder-dissipative along a sequence
of scales, \(\iint_{Q_{r_k}(x_0,T_*)}|\nabla u|^2\ge\varepsilon_1r_k\), but
cylinder dissipation may sit at late times far from \(x_0\) relative to
the parabolic distance, outside every paraboloid centred at \(x_0\); such
dissipation belongs to paraboloids centred at nearby points.  Hence
paraboloid dissipation at a fixed centre is not a consequence of (E1),
and Theorem~\ref{thm:c2v27-dirichlet-sees} is a statement about
paraboloid-dissipative points.  Under the type-I rate every singular
point is paraboloid-dissipative along a sequence of scales, by the
per-scale equivalence of the first cycle with the window
\((T_*-S\rho^2,T_*-\eta\rho^2)\) and the ball \(B_{R\rho}\), which lies in a
paraboloid of comparable scale for \(R\le2\eta^{1/2}\); without a rate this
is the moving-centre question of V24 in the time direction.
\end{firewall}

\begin{firewall}[Item (AC2): no rate-free lower bound on the weighted energy]
\label{fw:c2v27-AC2}
A lower bound on \(\Phi_{u,x_0}\) along dissipative scales cannot follow
from the dissipation and the Gaussian local energy identity alone:
Proposition~\ref{prop:c2v27-invisible} exhibits the obstruction, a
sequence of rescalings with dissipation at least \(\varepsilon_1\) on the
unit cylinder and Gaussian energy tending to zero, which is consistent
with the first identity because the flux term is not bounded by
\(C/\lambda\) without a rate.  The obstruction is exactly frequency
escape.  Consequently the first Gaussian ledger is blind to a class of
hypothetical rate-free blow-ups, while the second ledger, through the
Gaussian Dirichlet functional, is not
(Theorem~\ref{thm:c2v27-dirichlet-sees}).
\end{firewall}

% =============================================================================
\section{Integrated V27 conclusion and next acquisition order}
\label{sec:c2v27-conclusion}
% =============================================================================

\begin{theorem}[What V27 proves]
\label{thm:c2v27-conclusion}
Relative to the first cycle and V1--V26, without a rate and modulo
(E1)--(E2): the frequency-escape characterization
(Proposition~\ref{prop:c2v27-escape}), Gaussian invisibility
(Proposition~\ref{prop:c2v27-invisible}), and the lower bound
\eqref{eq:c2v27-dirichlet-sees} on the Gaussian Dirichlet functional
at paraboloid-dissipative scales.  The full Navier--Stokes
stability/global-regularity theorem remains unproved.
\end{theorem}

\begin{proof}
The cited results.
\end{proof}

\begin{programme}[V28 acquisition order]
\label{prog:c2v28}
\begin{enumerate}[label=\textup{(AC\arabic*)},leftmargin=4.8em]
\item The second identity along dissipative scales.  Using
      \eqref{eq:c2v27-dirichlet-sees} and the rate-free ledger
      \eqref{eq:c2v14-ledger}, derive what the second identity says below a paraboloid-dissipative scale: either \(\mathcal L_u\) is bounded
      on the window and the log-integral of
      \(\|Y_u\|_G^2+\langle N(V_u),Y_u\rangle_G\) is bounded, or
      \(\mathcal L_u\) is unbounded and Theorem~\ref{thm:c2v14-necessary}
      applies; determine whether the lower bound on \(\mathcal L_u\)
      forces a lower bound on \(\|Y_u\|_G\) or on \(\|N(V_u)\|_G\) on
      average, from \(Y_u=L_-V_u-N(V_u)\) and the spectral gap.
\item Rate-free defect.  Determine whether \(\|Y_u\|_G\) is bounded
      along parabolic scales, using the parabolic bounds and the
      expression of \(Y_u\) through \(\partial_tu\); record the exact
      bound or the obstruction.
\item The next compulsory companion audit is V30.
\end{enumerate}
\end{programme}

\begin{assessment}[Position after V27]
\label{ass:c2v27-position}
Without a rate, blow-up can escape the first Gaussian ledger by
frequency escape but not the second; the Dirichlet functional has a log-integral bounded below at
paraboloid-dissipative scales.  No full stability or
global-regularity claim is made.
\end{assessment}

% =============================================================================
\section*{Version V27 (second cycle) append-only record}
\addcontentsline{toc}{section}{Version V27 (second cycle) append-only record}
% =============================================================================

The complete V26 source of the second cycle was retained before this
continuation.  Its SHA--256 digest was
\begin{center}\scriptsize\ttfamily
99d3b8d622916059422d3ba62f35c018f6cbe5fda4a08e4acc6b185b26d43433.
\end{center}
V27 characterized frequency escape, proved Gaussian invisibility under
it, and proved that the Gaussian Dirichlet functional is bounded below
on average along dissipative scales.  No full regularity claim is made.

% =============================================================================
\section{V28: Gaussian-invisible dissipation forces a divergent defect or a
divergent nonlinearity}
\label{sec:c2v28-entry}
% =============================================================================

Items (AC1) and (AC2) of Programme~\ref{prog:c2v28} are carried out.  A
Cauchy--Schwarz inequality between the Gaussian Dirichlet functional,
the weighted energy, and the Gaussian norm of the linear self-similar
operator applied to the profile shows that, on any window of scales
where the weighted energy is small while the Dirichlet functional has a
fixed log-integral, the Gaussian norm of the linear part of the
self-similarity defect is large on log-average, inversely to the square
root of the weighted energy.  Since the linear part is the sum of the
defect and the nonlinearity, either the self-similarity defect or the
Gaussian nonlinearity diverges on log-average along Gaussian-invisible,
paraboloid-dissipative scales.  Combined with the rate-free second
identity this gives, at a point where the Dirichlet functional stays bounded, a
rate-free trichotomy: divergent anti-alignment of the nonlinearity
with the defect, relative self-similarity (square-integrable defect,
non-integrable nonlinearity), or a slowly non-integrable defect.  No
rate-free bound on the Gaussian norm of the defect exists, and the
obstruction is recorded.  No global regularity claim is made.  The next
compulsory companion audit is V30.

% =============================================================================
\section{The Cauchy--Schwarz bridge}
\label{sec:c2v28-bridge}
% =============================================================================

Throughout, \(u\) is a smooth mean-zero solution with \(T_*<\infty\),
\(x_0\in\Torus^3\), \(V=V_u\) the periodic profile
\eqref{eq:c2v14-periodic-profile} at scale \(\lambda\), \(Y=Y_u\) its
defect, \(N=N(V_u)\), all in \(L^2(G)\) at every scale, and
\(\varphi=\|V\|_G^2=\Phi_{u,x_0}\), \(\mathcal L=\mathcal L_u\) of \eqref{eq:c2v14-Lu}.

\begin{lemma}[The linear part of the defect is large when the profile is Gaussian-small]
\label{lem:c2v28-bridge}
For every scale, \(\langle L_-V,V\rangle_G=-2\mathcal L\), hence
\begin{equation}
 \|L_-V\|_G\ \ge\ \frac{2\mathcal L}{\varphi^{1/2}},
 \qquad
 \|Y\|_G+\|N\|_G\ \ge\ \|L_-V\|_G\ \ge\ \frac{2\mathcal L}{\varphi^{1/2}} .
 \label{eq:c2v28-bridge}
\end{equation}
\end{lemma}

\begin{proof}
\(\langle L_-V,V\rangle_G=-\nu\|\nabla V\|_G^2-\frac12\|V\|_G^2=-2\mathcal L\) by the
\(G\)-symmetry (Lemma~\ref{lem:c2v2-dictionary}, valid for the periodic
profile against the whole-space Gaussian as in
Theorem~\ref{thm:c2v14-second-rate-free}); Cauchy--Schwarz gives
\(2\mathcal L\le\|L_-V\|_G\|V\|_G\); and \(L_-V=Y+N\).
\end{proof}

\begin{theorem}[Divergent defect or divergent nonlinearity]
\label{thm:c2v28-divergence}
Let \(0<a<b\) and suppose \(\varphi(\lambda)\le\eta\) for all \(\lambda\in[a,b]\)
(Gaussian-invisible window) and \(\int_a^b\mathcal L\,\dd\lambda/\lambda\ge c_0\)
(Dirichlet-active window).  Then, without any rate,
\begin{equation}
 \boxed{\quad
 \int_a^b\bigl(\|Y_u\|_G+\|N(V_u)\|_G\bigr)\frac{\dd\lambda}\lambda\ \ge\ \frac{2c_0}{\eta^{1/2}} .
 \quad}
 \label{eq:c2v28-divergence}
\end{equation}
In particular, if \(x_0\) is paraboloid-dissipative along \(\lambda_k\to0\)
with constant \(\varepsilon'\) and Gaussian-invisible at scale zero,
\(\eta(\lambda):=\sup_{(0,\lambda]}\varphi\to0\) as \(\lambda\to0\), then
\begin{equation}
 \int_0^{\lambda}\bigl(\|Y_u\|_G+\|N(V_u)\|_G\bigr)\frac{\dd\lambda'}{\lambda'}=+\infty
 \qquad\text{for every }\lambda>0 :
 \label{eq:c2v28-infinite}
\end{equation}
the Gaussian norm of the self-similarity defect plus that of the
nonlinearity is not integrable in log-scale at scale zero.
\end{theorem}

\begin{proof}
Integrate \eqref{eq:c2v28-bridge} against \(\dd\lambda/\lambda\) with
\(\varphi^{1/2}\le\eta^{1/2}\).  For \eqref{eq:c2v28-infinite}, put
\(F(\lambda):=\int_0^\lambda(\|Y_u\|_G+\|N(V_u)\|_G)\dd\lambda'/\lambda'\in[0,\infty]\), a
nondecreasing function of \(\lambda\).  If \(F(\lambda)<\infty\) for some
\(\lambda\), then \(F(\lambda')\to0\) as \(\lambda'\to0\) (tail of a convergent
integral); but the first statement with \([a,b]=(0,\lambda_k]\),
\(\eta=\eta(\lambda_k)\), and \eqref{eq:c2v27-dirichlet-sees} with
\(c_0=\frac\nu4e^{-1/\nu}\varepsilon'\) give
\(F(\lambda_k)\ge\frac{\nu e^{-1/\nu}}2\varepsilon'\,\eta(\lambda_k)^{-1/2}\to\infty\),
a contradiction.  Hence \(F\equiv\infty\).
\end{proof}

\begin{assessment}[Meaning]
\label{ass:c2v28-meaning}
A blow-up that is invisible to the weighted energy at the parabolic
scale but dissipates at the parabolic order in the backward paraboloid
cannot be quiet: in the self-similar variables its profile is small in
the Gaussian norm while its Gaussian Dirichlet integral is not, so the
linear self-similar operator applied to it is large, and this size
must be accounted for either by the profile's own rate of change in
log-scale (a large defect) or by the nonlinearity (a large \(N\)).  The
size is quantified by the inverse square root of the weighted energy.
Under the type-I rate the weighted energy is bounded below at
homogeneous points and this divergence does not occur; without a rate
it is the first quantitative statement of the series about
frequency-escaping blow-up.
\end{assessment}

% =============================================================================
\section{The rate-free dichotomy at Dirichlet-bounded points}
\label{sec:c2v28-dichotomy}
% =============================================================================

\begin{theorem}[Divergent anti-alignment or relative self-similarity]
\label{thm:c2v28-dichotomy}
Suppose \(x_0\) is paraboloid-dissipative along \(\lambda_k\to0\), Gaussian
invisible at scale zero as in Theorem~\ref{thm:c2v28-divergence}, and
Dirichlet-bounded: \(\mathcal L_u(\lambda)\le M_{\mathcal L}\) for
\(\lambda\le\lambda_1\).  Then for all \(0<a<\lambda\le\lambda_1\),
\begin{equation}
 \Bigl|\int_a^{\lambda}\bigl(\|Y_u\|_G^2+\langle N(V_u),Y_u\rangle_G\bigr)\frac{\dd\lambda'}{\lambda'}\Bigr|\ \le\ \frac{M_{\mathcal L}}2 ,
 \label{eq:c2v28-bounded-ledger}
\end{equation}
while \eqref{eq:c2v28-infinite} holds.  Consequently, writing
\(\int_0\) for \(\int_0^\lambda\cdot\,\dd\lambda'/\lambda'\), exactly one of the
following holds.
\begin{enumerate}[label=\textup{(\Alph*)},leftmargin=3.2em]
\item \emph{Divergent anti-alignment.}  \(\int_0\|Y_u\|_G^2=\infty\), and then
      \(\int_a^{\lambda}\langle N(V_u),Y_u\rangle_G\,\dd\lambda'/\lambda'\le\frac{M_{\mathcal L}}2-\int_a^{\lambda}\|Y_u\|_G^2\,\dd\lambda'/\lambda'\to-\infty\)
      as \(a\to0\): the nonlinearity is anti-aligned with the defect
      with a log-integral diverging at least as fast as that of the
      squared defect.
\item \emph{Relative self-similarity.}  \(\int_0\|Y_u\|_G^2<\infty\) and
      \(\int_0\|N(V_u)\|_G=\infty\): the defect \(Y=L_-V-N(V)\) is
      square-integrable in log-scale at zero while the nonlinearity is
      not integrable; if moreover \(\int_0\|Y_u\|_G<\infty\), then also
      \(\int_0\|L_-V_u\|_G=\infty\), and Leray's equation \(L_-V=N(V)\)
      holds up to an error integrable over the log-scales on which its
      terms are not.
\item \emph{Slowly non-integrable defect.}  \(\int_0\|Y_u\|_G^2<\infty\),
      \(\int_0\|N(V_u)\|_G<\infty\), and then necessarily
      \(\int_0\|Y_u\|_G=\infty\): the defect is square-integrable but not
      integrable in log-scale at zero, and the nonlinearity is
      integrable.
\end{enumerate}
\end{theorem}

\begin{proof}
\eqref{eq:c2v28-bounded-ledger} is the rate-free ledger
\eqref{eq:c2v14-ledger} with \(0\le\mathcal L_u\le M_{\mathcal L}\).  The
three cases are exclusive and exhaustive.  In (A) the bound on the
ledger gives the divergence of the cross term.  In (B) with
\(\int_0\|Y\|_G<\infty\), \(\|L_-V\|_G\ge\|N\|_G-\|Y\|_G\) gives
\(\int_0\|L_-V\|_G=\infty\).  In (C), \eqref{eq:c2v28-infinite} forces
\(\int_0\|Y_u\|_G=\infty\).
\end{proof}

\begin{remark}[Cases (B) and (C)]
\label{rem:c2v28-residual}
In case (B) without the additional integrability of \(\|Y\|_G\), and in
case (C), the divergence \eqref{eq:c2v28-infinite} may be carried partly
or wholly by a defect that is square-integrable but not integrable in
log-scale, a rate of divergence between \(L^2\) and \(L^1\) of the
log-measure near zero; the series does not exclude this.  The clean
statement independent of the case is \eqref{eq:c2v28-infinite} together
with \eqref{eq:c2v28-bounded-ledger}: at a Gaussian-invisible,
paraboloid-dissipative, Dirichlet-bounded point, the Gaussian norms of
defect and nonlinearity are not integrable in log-scale at zero while
the signed combination \(\|Y\|_G^2+\langle N,Y\rangle_G\) has log-integrals
bounded by \(\frac12M_{\mathcal L}\) on every window.
\end{remark}
\begin{firewall}[Item (AC2): no rate-free bound on the defect]
\label{fw:c2v28-AC2}
\(Y_u=\lambda^3\partial_tu-\frac12V_u-\frac12z\cdot\nabla V_u\) in the rescaled
variables, and \(\lambda^3\partial_tu=\nu\Delta V_u-N(V_u)\) involves two
derivatives of the velocity and the pressure gradient, for which the
parabolic bounds of Definition~\ref{def:c2v26-parabolic} give no
\(L^2(G)\) control at a fixed time (they control \(\nabla V\) in
\(L^2\) of the cylinder only).  Hence \(\|Y_u(\sigma)\|_G\) is finite at
every scale but has no bound in terms of \(M\); the same holds for
\(\|N(V_u)\|_G\), which needs \(V\) in \(L^\infty\) and \(\nabla V\) in
\(L^2(G)\) at the same time.  The rate-free statements of this version
are therefore about finite quantities without uniform bounds, and
Theorem~\ref{thm:c2v28-divergence} is a lower bound on their
log-integrals, which is all the rate-free ledgers can give.  Under the
type-I rate both norms are bounded (Proposition~\ref{prop:c2v7-defect-equation},
Lemma~\ref{lem:c2v3-N-bound}), and the divergence
\eqref{eq:c2v28-divergence} cannot occur because \(\varphi\) is bounded
below at homogeneous points and, at any singular point, along active
scales.
\end{firewall}


% =============================================================================
\section{Integrated V28 conclusion and next acquisition order}
\label{sec:c2v28-conclusion}
% =============================================================================

\begin{theorem}[What V28 proves]
\label{thm:c2v28-conclusion}
Relative to the first cycle and V1--V27, without a rate: the bridge
inequality \eqref{eq:c2v28-bridge}, the divergence theorem
\eqref{eq:c2v28-divergence}, and the trichotomy of Theorem~\ref{thm:c2v28-dichotomy} at Dirichlet-bounded points.  The full
Navier--Stokes stability/global-regularity theorem remains unproved.
\end{theorem}

\begin{proof}
The cited results.
\end{proof}

\begin{programme}[V29 acquisition order]
\label{prog:c2v29}
\begin{enumerate}[label=\textup{(AC\arabic*)},leftmargin=4.8em]
\item Relative self-similarity.  Make alternative (B) of
      Theorem~\ref{thm:c2v28-dichotomy} precise as a statement about
      the normalized profile \(\widehat V:=V/\|V\|_G\): derive the
      equation of \(\widehat V\) in log-scale and determine whether
      relative self-similarity forces \(\widehat V\) to converge, along
      scales, to a solution of a normalized Leray equation, and what
      (E9) says about it (the normalized profile need not be in
      \(L^6\) with a bound, since \(\|V\|_G\to0\)).
\item Prepare the V30 audit and the consolidation of V26--V29 as the
      rate-free chapter of the second cycle: parabolic scales,
      compactness to suitable weak solutions, frequency escape,
      Gaussian invisibility, the Dirichlet ledger's lower bound, and
      the divergence dichotomy.
\item The next compulsory companion audit is V30.
\end{enumerate}
\end{programme}

\begin{assessment}[Position after V28]
\label{ass:c2v28-position}
Gaussian-invisible, paraboloid-dissipative blow-up forces a divergent
self-similarity defect or a divergent Gaussian nonlinearity, and at Dirichlet-bounded points splits into divergent anti-alignment,
relative self-similarity, or a slowly non-integrable defect.  No full stability or global-regularity
claim is made.
\end{assessment}

% =============================================================================
\section*{Version V28 (second cycle) append-only record}
\addcontentsline{toc}{section}{Version V28 (second cycle) append-only record}
% =============================================================================

The complete V27 source of the second cycle was retained before this
continuation.  Its SHA--256 digest was
\begin{center}\scriptsize\ttfamily
136fcd8b35e5552b51d80f04ee9f8d28220446b462cd43a0196be2193d16200e.
\end{center}
V28 proved that Gaussian-invisible paraboloid-dissipative blow-up
forces a divergent defect or a divergent nonlinearity, and the
rate-free dichotomy at Dirichlet-bounded points.  No full regularity
claim is made.

% =============================================================================
\section{V29: the normalized profile, the Rayleigh quotient, and the
spectral form of Gaussian invisibility}
\label{sec:c2v29-entry}
% =============================================================================

Items (AC1) and (AC2) of Programme~\ref{prog:c2v29} are carried out.  The
profile normalized to unit Gaussian norm satisfies an exact equation in
log-scale: the projected self-similar operator plus the nonlinearity
scaled by the square root of the weighted energy.  The Rayleigh
quotient of the profile for the self-similar operator equals twice the
Dirichlet functional over the weighted energy, is at least one half, and
governs the decay of the weighted energy together with the normalized
flux; Gaussian invisibility is thereby equivalent to the divergence of
the log-integral of the Rayleigh quotient plus the normalized flux, so
that a Gaussian-invisible blow-up either escapes to high
Ornstein--Uhlenbeck modes or has net outward normalized flux on
average.  Relative self-similarity is stated for the normalized
profile, and it is recorded that (E9) does not apply to normalized
limits.  The consolidation for V30 is prepared.  No global regularity
claim is made.  The next version is the compulsory companion audit.

% =============================================================================
\section{The normalized profile}
\label{sec:c2v29-normalized}
% =============================================================================

Throughout, \(u\) is a smooth mean-zero solution with \(T_*<\infty\),
\(x_0\in\Torus^3\), \(V=V_u\) the periodic profile at scale
\(\lambda=e^{-\sigma/2}\), \(\varphi=\|V\|_G^2>0\) (\(V\ne0\) at every scale for
a nonzero solution, by backward uniqueness), and
\begin{equation}
 \widehat V:=\frac{V}{\varphi^{1/2}},\qquad
 r:=-\langle L_-\widehat V,\widehat V\rangle_G=\frac{2\mathcal L}{\varphi}\ \ge\ \frac12,\qquad
 \widehat N:=N(\widehat V)=\frac{N(V)}{\varphi},\qquad
 \widehat j:=4\langle\widehat V,\widehat N\rangle_G=\frac{j}{\varphi^{3/2}} ,
 \label{eq:c2v29-normalized}
\end{equation}
where \(r\) is the Rayleigh quotient of \(\widehat V\) for \(-L_-\), at
least \(\frac12\) by Proposition~\ref{prop:c2v8-spectrum}, and \(j=\lambda\mathcal J_u\)
the scale-invariant flux.

\begin{proposition}[Equation of the normalized profile and of the weighted energy]
\label{prop:c2v29-equation}
For every \(\sigma\),
\begin{align}
 \partial_\sigma\widehat V&=L_-\widehat V+r\,\widehat V-\varphi^{1/2}\Bigl(\widehat N-\langle\widehat N,\widehat V\rangle_G\widehat V\Bigr),
 \label{eq:c2v29-Vhat}\\
 \frac{\dd}{\dd\sigma}\log\varphi&=-2r-\tfrac12\varphi^{1/2}\,\widehat j .
 \label{eq:c2v29-logphi}
\end{align}
The right side of \eqref{eq:c2v29-Vhat} is \(G\)-orthogonal to
\(\widehat V\), as it must be.
\end{proposition}

\begin{proof}
\(\partial_\sigma\widehat V=\varphi^{-1/2}Y-\frac12\varphi^{-3/2}(\partial_\sigma\varphi)V\)
with \(Y=L_-V-N(V)=\varphi^{1/2}L_-\widehat V-\varphi\widehat N\) and
\(\partial_\sigma\varphi=2\langle V,Y\rangle_G=2\varphi\langle\widehat V,L_-\widehat V\rangle_G-2\varphi^{3/2}\langle\widehat V,\widehat N\rangle_G
=-2r\varphi-2\varphi^{3/2}\langle\widehat V,\widehat N\rangle_G\); substituting gives
\eqref{eq:c2v29-Vhat}, and \eqref{eq:c2v29-logphi} is the last display
divided by \(\varphi\), with \(2\langle\widehat V,\widehat N\rangle_G=\frac12\widehat j\).
Orthogonality: \(\langle L_-\widehat V+r\widehat V,\widehat V\rangle_G=-r+r=0\) and
the projected nonlinearity is orthogonal by construction.
\end{proof}

\begin{theorem}[Spectral form of Gaussian invisibility]
\label{thm:c2v29-spectral}
For \(0<a<b\), with \(\sigma_a=-2\log a>\sigma_b=-2\log b\),
\begin{equation}
 \boxed{\quad
 \log\frac{\Phi_{u,x_0}(b)}{\Phi_{u,x_0}(a)}=\int_{\sigma_b}^{\sigma_a}\Bigl(2r+\tfrac12\varphi^{1/2}\widehat j\Bigr)\dd\sigma .
 \quad}
 \label{eq:c2v29-spectral}
\end{equation}
Consequently:
\begin{enumerate}[label=\textup{(\roman*)},leftmargin=3.2em]
\item Gaussian invisibility at scale zero, \(\Phi_{u,x_0}(a)\to0\) as
      \(a\to0\), holds if and only if
      \(\int_{\sigma_b}^{\infty}(2r+\frac12\varphi^{1/2}\widehat j)\dd\sigma=+\infty\):
      the profile escapes to high Ornstein--Uhlenbeck modes (\(r\) not
      integrable in log-scale) or the normalized flux
      \(\frac12\varphi^{1/2}\widehat j=2\varphi^{-1}\langle V,N(V)\rangle_G\) is
      positive (outward) with non-integrable log-integral, or both.
\item Since \(r\ge\frac12\), the linear part alone makes
      \(\Phi_{u,x_0}(a)\le\Phi_{u,x_0}(b)(a/b)^{2}\exp\bigl(-\frac12\int\varphi^{1/2}\widehat j\bigr)\):
      the weighted energy decays at least like the square of the scale
      unless the flux is inward, \(\widehat j<0\), with
      \(\int_{\sigma_b}^{\sigma_a}\frac12\varphi^{1/2}|\widehat j|\ge2\log(b/a)-\log(\Phi(b)/\Phi(a))\).
      At a homogeneous point of a type-I solution, where
      \(\Phi_{u,x_0}\) is pinched, the inward normalized flux must
      therefore supply exactly \(2\int r\,\dd\sigma\ge2\log(b/a)\) minus a
      bounded amount, which is the inward flux of the first cycle in
      normalized form.
\end{enumerate}
\end{theorem}

\begin{proof}
Integrate \eqref{eq:c2v29-logphi} over \([\sigma_b,\sigma_a]\) (recall
\(\sigma\) increases as the scale decreases).  (i) is the equivalence of
\(\log\Phi(a)\to-\infty\) with the divergence of the integral, and the
two summands are the only sources.  (ii) uses \(r\ge\frac12\) and
\(\int_{\sigma_b}^{\sigma_a}\dd\sigma=2\log(b/a)\).
\end{proof}

\begin{assessment}[Invisibility as spectral escape]
\label{ass:c2v29-spectral}
The weighted energy of the profile is the Gaussian norm squared of a
vector evolving under the self-similar Ornstein--Uhlenbeck operator
(which damps every mode at rate at least one half, the constant mode
at exactly that rate and the mode of Hermite degree \(k\) at rate
\(\frac12+\frac k2\)) and the nonlinearity.  A pinched weighted energy
requires the nonlinearity to pump energy inward at exactly the rate of
the damping, in the mean, which is the first cycle's inward flux.
Gaussian invisibility is the failure of this pumping, in two possible
ways: the profile's Gaussian energy migrates to high Hermite modes,
where the damping is strong (this is frequency escape seen in the
Ornstein--Uhlenbeck basis), or the nonlinearity pumps outward.  Under
the type-I rate the Rayleigh quotient is bounded on the family at a
homogeneous point (\(r=2\mathcal L/\varphi\le(4\nu C+2C_\Phi)/(4c_\Phi)\)),
so high-mode escape is excluded there and the flux must be inward.
\end{assessment}

% =============================================================================
\section{Relative self-similarity for the normalized profile}
\label{sec:c2v29-relative}
% =============================================================================

\begin{proposition}[Normalized form of relative self-similarity]
\label{prop:c2v29-relative}
In alternative (B) of Theorem~\ref{thm:c2v28-dichotomy} with
\(\int_0\|Y_u\|_G<\infty\), the normalized profile satisfies
\(\int_0\|\partial_\sigma\widehat V\|_G\,\dd\lambda/\lambda\le\int_0\frac{\|Y\|_G}{\varphi^{1/2}}\,\dd\lambda/\lambda\),
and the equation \eqref{eq:c2v29-Vhat} reads
\(L_-\widehat V+r\widehat V=\varphi^{1/2}P^\perp\widehat N+\partial_\sigma\widehat V\):
the normalized profile solves the normalized Leray equation
\(L_-\widehat V+r\widehat V=\varphi^{1/2}P^\perp N(\widehat V)\) up to the error
\(\partial_\sigma\widehat V\), where \(P^\perp\) is the \(G\)-orthogonal
projection onto \(\widehat V^\perp\).  A local limit of the normalized
profiles along scales, if it exists, is a unit vector of \(L^2(G)\)
with no bound in \(L^6\) or in \(H^1(G)\) in general, and (E9) does not
apply to it.
\end{proposition}

\begin{proof}
\(\partial_\sigma\widehat V=\varphi^{-1/2}P^\perp Y\) by the proof of
Proposition~\ref{prop:c2v29-equation}, and \(\|P^\perp Y\|_G\le\|Y\|_G\).
The remaining statements are \eqref{eq:c2v29-Vhat} and the definitions;
the normalization divides by \(\varphi^{1/2}\to0\), so the \(L^6\) bound
of the class and the Dirichlet bound do not transfer.
\end{proof}

\begin{firewall}[Item (AC1): what is not obtained]
\label{fw:c2v29-AC1}
The weight \(\varphi^{-1/2}\) in \(\int_0\|Y\|_G\varphi^{-1/2}\) diverges under
invisibility, so the integrability of \(\|Y\|_G\) in log-scale does not
give integrability of \(\|\partial_\sigma\widehat V\|_G\), and no convergence
of \(\widehat V\) along scales follows.  The normalized Leray equation has
the \(\sigma\)-dependent Rayleigh quotient \(r\) in place of the fixed
coefficient \(\frac12\), and its nonlinearity is scaled by
\(\varphi^{1/2}\to0\) while \(\|N(\widehat V)\|_G=\|N(V)\|_G/\varphi\) may
diverge; nothing forces the limit to be an eigenfunction of \(L_-\).  Item
(AC1) is closed with the exact equations \eqref{eq:c2v29-Vhat},
\eqref{eq:c2v29-logphi} and this record.
\end{firewall}

% =============================================================================
\section{Integrated V29 conclusion and next acquisition order}
\label{sec:c2v29-conclusion}
% =============================================================================

\begin{theorem}[What V29 proves]
\label{thm:c2v29-conclusion}
Relative to the first cycle and V1--V28, without a rate: the normalized
equations \eqref{eq:c2v29-Vhat}--\eqref{eq:c2v29-logphi}, the spectral
form \eqref{eq:c2v29-spectral} of Gaussian invisibility with its two
consequences, and the normalized form of relative self-similarity.  The
full Navier--Stokes stability/global-regularity theorem remains
unproved.
\end{theorem}

\begin{proof}
The cited results.
\end{proof}

\begin{programme}[V30 acquisition order, with compulsory companion audit]
\label{prog:c2v30}
\begin{enumerate}[label=\textup{(AC\arabic*)},leftmargin=4.8em]
\item Perform the compulsory review of the current SM and YM notes;
      record digests; import nothing numerical.
\item Consolidate V26--V29 as the rate-free chapter: parabolic scales
      and their independence, compactness to suitable weak solutions,
      frequency escape, Gaussian invisibility, the Dirichlet ledger's
      lower bound at paraboloid-dissipative scales, the divergence
      theorem and the trichotomy, and the spectral form of
      invisibility.
\item Record the position after thirty versions: the three hypotheses
      of the type-I case (homogeneity, compact dissipation, coherence)
      and the two escape mechanisms of the rate-free case (spatial
      spreading, high-mode escape), and the direction for V31--V35.
\item The next compulsory companion audit after V30 is V35.
\end{enumerate}
\end{programme}

\begin{assessment}[Position after V29]
\label{ass:c2v29-position}
Gaussian invisibility is the divergence of the Rayleigh quotient plus
the normalized flux in log-scale: high-mode escape or outward pumping.
No full stability or global-regularity claim is made.
\end{assessment}

% =============================================================================
\section*{Version V29 (second cycle) append-only record}
\addcontentsline{toc}{section}{Version V29 (second cycle) append-only record}
% =============================================================================

The complete V28 source of the second cycle was retained before this
continuation.  Its SHA--256 digest was
\begin{center}\scriptsize\ttfamily
d2e91e3a073547b813f1cbb3d0ac55c4de829693cb4767efaa998ae978a34532.
\end{center}
V29 derived the normalized profile equations and the spectral form of
Gaussian invisibility.  No full regularity claim is made.

% =============================================================================
\section{V30: compulsory companion audit, the rate-free chapter consolidated,
and the position after thirty versions}
\label{sec:c2v30-entry}
% =============================================================================

This is the thirtieth version of the second cycle.  The compulsory
review of the two companion programmes is performed first.  V26--V29
are then consolidated as the rate-free chapter of the cycle: parabolic
scales and the independence of their three conditions, compactness to
suitable weak solutions along them, frequency escape and Gaussian
invisibility, the Dirichlet ledger's lower bound at
paraboloid-dissipative scales, the divergence theorem with its
trichotomy, and the spectral form of invisibility.  The position after
thirty versions is recorded: three hypotheses organize the type-I case
and two escape mechanisms organize the rate-free case; the direction
for V31--V35 is fixed.  No global regularity claim is made.  The next
compulsory companion audit is V35.

% =============================================================================
\section{Compulsory V30 review of the current companion notes}
\label{sec:c2v30-companion-audit}
% =============================================================================

The current companion files in \texttt{latex\_notes} were inspected
read-only.  The frozen checkpoints are
\begin{align}
 &\texttt{Renewal\_SM\_Einstein\_Continuum\_Research\_Notes\_V94.tex},
 \notag\\[-2pt]
 &\qquad
 \texttt{87ce0ae7bdd5645bc5cb1579aab1511c4a2e52c59d984c1c9bfaa8d6ac771003},
 \label{eq:c2v30-SM-checkpoint}\\
 &\texttt{Renewal\_Yang\_Mills\_Mass\_Gap\_Research\_Notes\_V69.tex},
 \notag\\[-2pt]
 &\qquad
 \texttt{1581b543565f838a8438cbc50d080cb38704598ed3c6b2f6c63bf631aa0c706b}.
 \label{eq:c2v30-YM-checkpoint}
\end{align}
At the time of inspection a YM file named \(\ldots\)\texttt{V68.tex} with
the same digest as \eqref{eq:c2v30-YM-checkpoint} was also present; the
two are byte-identical, and the V69 name is recorded.

\emph{SM V91--V94.}  The SM programme has moved to a lattice
formulation: after an overlap algebra with sign-projector locality
(V93), V94 rewrote its Wilson kernel in covariant-difference form,
proved that every forward/backward commutator is controlled by one
plaquette defect, verified an exact Wilson-square decomposition at unit
parameter with separate lower bounds for its parts, proved a uniform
lower bound \(A^*A\succeq(1-30\delta)I\) and an upper bound \(\|A\|\le7\)
for its Wilson operator, obtained volume-independent sign-locality
constants, closed its plaquette-admissible Wilson gap on the complete
admissible class with separation of index sectors, and isolated measure
concentration, a Yukawa floor, and a determinant phase as remaining
gates.

\emph{YM V68.}  The YM programme proved factorial tail recovery, an
exact ordered-block identity, control of its added-jump and
original-clock outer events by the complete centre polynomial, a
polynomial original-clock tail, closed its complete low-centre outer
sum and three of its cards for small centres, and isolated a
fixed-size-centre card.

\begin{theorem}[V30 companion-audit decision]
\label{thm:c2v30-companion-decision}
No SM V91--V94 lattice, plaquette, or operator-norm constant and no YM
V68 tail or centre estimate is imported.  Neither bears on the
rate-free chapter or on the gates of Theorem~\ref{thm:c2v15-gates}.
\end{theorem}

\begin{proof}
The SM results concern a lattice Wilson operator and its spectral gap;
the YM results concern outer-event bounds of a lattice clock.  None is a
Navier--Stokes statement.
\end{proof}

\begin{assessment}[Direction after the audit]
\label{ass:c2v30-direction}
The SM programme's closure of a gap on an admissible class with explicit
operator bounds, leaving concentration, a floor, and a phase as
separate gates, has the same shape as this cycle's position: exact
signed statistics on the compact (type-I, homogeneous, compact,
coherent) class, with the escape mechanisms as separate gates.  No
change of direction follows from the audit.
\end{assessment}

% =============================================================================
\section{The rate-free chapter, consolidated}
\label{sec:c2v30-consolidated}
% =============================================================================

\begin{theorem}[Rate-free chapter, V26--V29]
\label{thm:c2v30-consolidated}
Let \(u\) be a smooth mean-zero solution on \([0,T_*)\times\Torus^3\) with
\(T_*<\infty\), \(x_0\in\Torus^3\), and no rate assumption.  Modulo
(E1)--(E2) where stated:
\begin{enumerate}[label=\textup{(\arabic*)},leftmargin=3.2em]
\item \emph{Parabolic scales.}  The velocity, dissipation, and pressure
      conditions \eqref{eq:c2v26-parabolic} are independent inputs;
      along \(M\)-parabolic scales the rescalings converge, along
      subsequences, to suitable weak solutions on the unit cylinder
      (E2), possibly zero, and the family of limits is not
      scale-invariant.
\item \emph{Frequency escape.}  At a singular point (E1 gives
      dissipative scales), a zero limit along parabolic dissipative
      scales means that the rescaled dissipation is carried by
      frequencies tending to infinity while the rescaled energy on
      every compact set tends to zero; under parabolic bounds at all
      scales up to \(R_\lambda\lambda\) this makes \(\Phi_{u,x_0}(\lambda_k)\to0\).
\item \emph{Dirichlet ledger.}  At a paraboloid-dissipative scale,
      \(\int_0^\lambda\mathcal L_{u,x_0}\,\dd\lambda'/\lambda'\ge\frac\nu4e^{-1/\nu}\varepsilon'\);
      the second identity holds rate-free with ledger
      \eqref{eq:c2v14-ledger}.
\item \emph{Divergence.}  \(\|Y\|_G+\|N\|_G\ge\|L_-V\|_G\ge2\mathcal L/\varphi^{1/2}\);
      at a paraboloid-dissipative, Gaussian-invisible point the
      log-integral of \(\|Y_u\|_G+\|N(V_u)\|_G\) at scale zero is
      infinite; if moreover \(\mathcal L_u\) is bounded, exactly one of:
      divergent anti-alignment, relative self-similarity, slowly
      non-integrable defect.
\item \emph{Spectral form.}  \(\frac{\dd}{\dd\sigma}\log\varphi=-2r-\frac12\varphi^{1/2}\widehat j\)
      with the Rayleigh quotient \(r=2\mathcal L/\varphi\ge\frac12\);
      Gaussian invisibility at scale zero is the divergence of
      \(\int(2r+\frac12\varphi^{1/2}\widehat j)\dd\sigma\): high-mode escape
      or outward normalized flux; pinched weighted energy requires
      inward normalized flux supplying \(2\int r\,\dd\sigma\ge2\log(b/a)\)
      minus a bounded amount.
\item \emph{Rate-free inward flux.}  Between two scales at which
      \(\Phi_{u,x_0}\le M'\), with \(\Phi_{u,x_0}\ge c\) across the window,
      \(\int_a^b\mathcal J_{u,x_0}\le M'-2c\log(b/a)-\int_a^bD_{u,x_0}\).
\end{enumerate}
\end{theorem}

\begin{proof}
(1) Proposition~\ref{prop:c2v26-independent}, Theorem~\ref{thm:c2v26-compactness},
Firewall~\ref{fw:c2v26-trivial}.  (2) Propositions~\ref{prop:c2v27-escape},
\ref{prop:c2v27-invisible}.  (3) Theorem~\ref{thm:c2v27-dirichlet-sees},
Theorem~\ref{thm:c2v14-second-rate-free}.  (4) Lemma~\ref{lem:c2v28-bridge},
Theorems~\ref{thm:c2v28-divergence}, \ref{thm:c2v28-dichotomy}.
(5) Proposition~\ref{prop:c2v29-equation}, Theorem~\ref{thm:c2v29-spectral}.
(6) Theorem~\ref{thm:c2v26-inward}.
\end{proof}

% =============================================================================
\section{Position after thirty versions and direction}
\label{sec:c2v30-position}
% =============================================================================

\begin{assessment}[Position after thirty versions of the second cycle]
\label{ass:c2v30-position-thirty}
The type-I case is organized by three hypotheses: homogeneity (or
statistical nontriviality), compact dissipation, and coherence of the
active centre.  Under them the blow-up has an exact, signed,
scale-invariant statistics against blow-up measures, and the remaining
statement is one mean inequality (Theorem~\ref{thm:c2v15-gates}).
Without them the blow-up escapes the Gaussian ledgers spatially
(non-compact dissipation, V24) or in scale (statistically trivial
lacunary points, V23).  The rate-free case is organized by two escape
mechanisms: spatial spreading and high-mode escape (frequency escape
in the Ornstein--Uhlenbeck basis); the Dirichlet ledger sees paraboloid
dissipation through both, and Gaussian-invisible dissipative blow-up
forces a non-integrable defect-plus-nonlinearity in log-scale.  None
of the mechanisms is excluded.  The three gates of the first cycle
stand.  For V31--V35 the direction is: (a) determine whether high-mode
escape is compatible with the rate-free ledgers when the Dirichlet
functional is bounded, by examining the \(\sigma\)-evolution of the
Rayleigh quotient (the ``power-method'' identity for the normalized
profile) and the exact rate at which the linear part drives \(r\)
upward against the nonlinearity; (b) determine whether paraboloid
dissipation at a fixed centre follows from cylinder dissipation under
compact dissipation with coherent centres; (c) prepare, by V35, the
consolidated statement of the second cycle's first thirty-five
versions.  No global regularity claim is made.
\end{assessment}

% =============================================================================
\section{Integrated V30 conclusion and next acquisition order}
\label{sec:c2v30-conclusion}
% =============================================================================

\begin{theorem}[What V30 establishes]
\label{thm:c2v30-conclusion}
Relative to the first cycle and V1--V29: the companion-audit decision
(Theorem~\ref{thm:c2v30-companion-decision}) and the consolidated
rate-free chapter (Theorem~\ref{thm:c2v30-consolidated}).  The full
Navier--Stokes stability/global-regularity theorem remains unproved.
\end{theorem}

\begin{proof}
The cited results.
\end{proof}

\begin{programme}[V31 acquisition order]
\label{prog:c2v31}
\begin{enumerate}[label=\textup{(AC\arabic*)},leftmargin=4.8em]
\item The Rayleigh-quotient identity.  Derive
      \(\frac{\dd}{\dd\sigma}r\) from \eqref{eq:c2v29-Vhat}: the linear
      part gives \(-2(\|L_-\widehat V\|_G^2-r^2)=-2\operatorname{Var}_{\widehat V}(L_-)\le0\)
      (the Rayleigh quotient of \(-L_-\) can only decrease under the
      linear flow, toward the bottom of the spectrum \(\frac12\)), and the
      nonlinear part gives \(2\varphi^{1/2}\langle L_-\widehat V+r\widehat V,\widehat N\rangle_G\);
      record the exact identity and its consequence: high-mode escape
      (\(r\to\infty\) in log-average) requires the nonlinear term to
      push the profile up the spectrum against the linear descent.
\item Combine with \eqref{eq:c2v29-logphi}: under Gaussian invisibility
      with bounded Dirichlet functional, determine the joint constraint
      on \(r\) and the normalized flux.
\item The next compulsory companion audit is V35.
\end{enumerate}
\end{programme}

\begin{assessment}[Position after V30]
\label{ass:c2v30-position}
The rate-free chapter is consolidated and the position after thirty
versions is recorded.  No full stability or global-regularity claim is
made.
\end{assessment}

% =============================================================================
\section*{Version V30 (second cycle) append-only record}
\addcontentsline{toc}{section}{Version V30 (second cycle) append-only record}
% =============================================================================

The complete V29 source of the second cycle was retained before this
continuation.  Its SHA--256 digest was
\begin{center}\scriptsize\ttfamily
3be7398455ba117fc89a243639ecdad037feddf5c91fab969547c7c65598c2a6.
\end{center}
V30 performed the compulsory companion audit, consolidated the
rate-free chapter V26--V29, and recorded the position after thirty
versions.  No full regularity claim is made.

% =============================================================================
\section{V31: the Rayleigh-quotient identity and forced high-mode escape}
\label{sec:c2v31-entry}
% =============================================================================

Items (AC1) and (AC2) of Programme~\ref{prog:c2v31} are carried out.  The
Rayleigh quotient of the normalized profile for the self-similar
operator satisfies an exact identity in log-scale: its linear part is
minus twice the spectral variance of the operator in the state of the
profile, which is nonpositive, so the linear flow drives the profile
down the spectrum toward the constant mode; its nonlinear part is an
explicit inner product of the nonlinearity with the residual of the
profile.  At a paraboloid-dissipative point the log-integral of the
Gaussian Dirichlet functional at scale zero is infinite, and under
Gaussian invisibility the Rayleigh quotient is bounded below by twice
the Dirichlet functional over the invisibility threshold, so it is
large on every Dirichlet-active scale: the mean Hermite degree of the
profile's Gaussian energy tends to infinity there, and by the identity
the nonlinearity must push the profile up the spectrum against the
linear descent by at least the growth of the Rayleigh quotient plus
twice the accumulated spectral variance.  No global regularity claim is
made.  The next compulsory companion audit is V35.

% =============================================================================
\section{The Rayleigh-quotient identity}
\label{sec:c2v31-identity}
% =============================================================================

Throughout, the notation of V29 is used: \(V=V_u\), \(\varphi=\|V\|_G^2\),
\(\widehat V=V/\varphi^{1/2}\), \(r=-\langle L_-\widehat V,\widehat V\rangle_G=2\mathcal L/\varphi\),
\(\widehat N=N(V)/\varphi\).  Define the spectral variance of \(-L_-\) in
the state \(\widehat V\) and the residual of the profile,
\begin{equation}
 \operatorname{Var}:=\|L_-\widehat V\|_G^2-r^2=\|L_-\widehat V+r\widehat V\|_G^2\ \ge0,
 \qquad
 \mathcal R:=L_-\widehat V+r\widehat V\ \perp_G\ \widehat V ,
 \label{eq:c2v31-variance}
\end{equation}
so that \(\operatorname{Var}=\|\mathcal R\|_G^2\), and \(\operatorname{Var}=0\) if
and only if \(\widehat V\) is an eigenfunction of \(L_-\).

\begin{theorem}[Rayleigh-quotient identity]
\label{thm:c2v31-identity}
For every \(\sigma\),
\begin{equation}
 \boxed{\quad
 \frac{\dd r}{\dd\sigma}=-2\operatorname{Var}+2\varphi^{1/2}\langle\mathcal R,\widehat N\rangle_G
 =-2\|\mathcal R\|_G^2+\frac2\varphi\bigl\langle L_-V+rV,\,N(V)\bigr\rangle_G .
 \quad}
 \label{eq:c2v31-identity}
\end{equation}
Equivalently, in terms of the two identities of the cycle,
\(\frac{\dd r}{\dd\sigma}=\frac2\varphi\frac{\dd\mathcal L}{\dd\sigma}+2r^2+\frac{2r}{\varphi}\langle V,N(V)\rangle_G\).
\end{theorem}

\begin{proof}
\(r=-\langle L_-\widehat V,\widehat V\rangle_G\) and \(L_-\) is \(G\)-symmetric, so
\(\frac{\dd r}{\dd\sigma}=-2\langle L_-\widehat V,\partial_\sigma\widehat V\rangle_G\).
Insert \eqref{eq:c2v29-Vhat}:
\(\langle L_-\widehat V,L_-\widehat V+r\widehat V\rangle_G=\|L_-\widehat V\|_G^2-r^2=\operatorname{Var}\),
and \(\langle L_-\widehat V,P^\perp\widehat N\rangle_G=\langle P^\perp L_-\widehat V,\widehat N\rangle_G=\langle\mathcal R,\widehat N\rangle_G\)
since \(P^\perp L_-\widehat V=L_-\widehat V-\langle L_-\widehat V,\widehat V\rangle_G\widehat V=\mathcal R\).
Hence \(\frac{\dd r}{\dd\sigma}=-2\operatorname{Var}+2\varphi^{1/2}\langle\mathcal R,\widehat N\rangle_G\),
and \(\varphi^{1/2}\langle\mathcal R,\widehat N\rangle_G=\varphi^{-1}\langle L_-V+rV,N(V)\rangle_G\).
The second form is the derivative of \(r=2\mathcal L/\varphi\) with
\eqref{eq:c2v29-logphi}; it agrees with the first because
\(\frac{\dd\mathcal L}{\dd\sigma}=-\|L_-V\|_G^2+\langle L_-V,N(V)\rangle_G\)
(Theorem~\ref{thm:c2v2-second}) and
\(\|L_-V\|_G^2/\varphi=\|L_-\widehat V\|_G^2=\operatorname{Var}+r^2\).
\end{proof}

\begin{assessment}[Linear descent, nonlinear push]
\label{ass:c2v31-descent}
Under the linear self-similar flow alone the Rayleigh quotient of the
normalized profile decreases at the rate of twice the spectral
variance and converges to the lowest eigenvalue present in the initial
profile: the constant mode, or the first mode in the invariant subspace
containing the profile.  The nonlinearity can raise it only through the
inner product of \(N(V)\) with the residual \(L_-V+rV\), the part of
\(L_-V\) orthogonal to \(V\).  High-mode escape therefore requires the
nonlinearity to be positively correlated with the residual of the
profile in the Gaussian inner product, persistently in log-scale, and
by an amount at least equal to twice the accumulated spectral variance
plus the growth of the Rayleigh quotient.
\end{assessment}

% =============================================================================
\section{Forced high-mode escape at invisible dissipative points}
\label{sec:c2v31-forced}
% =============================================================================

\begin{proposition}[The Dirichlet functional is not integrable in log-scale at a paraboloid-dissipative point]
\label{prop:c2v31-nonintegrable}
If \(x_0\) is paraboloid-dissipative along \(\lambda_k\to0\) with constant
\(\varepsilon'\), then \(\int_0^\lambda\mathcal L_{u,x_0}(\lambda')\,\dd\lambda'/\lambda'=+\infty\)
for every \(\lambda>0\).
\end{proposition}

\begin{proof}
By \eqref{eq:c2v27-dirichlet-sees}, the tail integrals over \((0,\lambda_k]\)
are all at least \(\frac\nu4e^{-1/\nu}\varepsilon'>0\); the tails of a
convergent integral tend to zero.
\end{proof}

\begin{lemma}[Rayleigh quotient and mean Hermite degree]
\label{lem:c2v31-degree}
Let \(\Pi_k\) be the \(G\)-orthogonal projection onto the eigenspace of
\(-L_-\) with eigenvalue \(\mu_k=\frac12+\frac k2\), \(k\ge0\) (Hermite
degree \(k\)), and \(a_k:=\|\Pi_k\widehat V\|_G^2\), \(\sum_ka_k=1\).  Then
\(r=\sum_k\mu_ka_k\), so the mean Hermite degree of the profile's
Gaussian energy is \(\sum_kka_k=2r-1\); in particular \(\widehat V\) has
nonzero energy on some mode of degree at least \(2r-1\), and
\(\operatorname{Var}=\sum_k(\mu_k-r)^2a_k\) is the variance of the degree
distribution divided by four.
\end{lemma}

\begin{proof}
Spectral theorem for the self-adjoint operator \(-L_-\) on \(L^2(G)\)
(Proposition~\ref{prop:c2v8-spectrum}); \(\mu_k=\frac12(k+1)\) gives
\(\sum_k\mu_ka_k=\frac12(\sum_kka_k+1)\).
\end{proof}

\begin{theorem}[Forced high-mode escape and the required nonlinear push]
\label{thm:c2v31-forced}
Let \(x_0\) be paraboloid-dissipative along \(\lambda_k\to0\) with constant
\(\varepsilon'\) and Gaussian-invisible at scale zero,
\(\eta(\lambda)=\sup_{(0,\lambda]}\varphi\to0\).  Then, without any rate:
\begin{enumerate}[label=\textup{(\roman*)},leftmargin=3.2em]
\item On every scale \(\lambda'\le\lambda\) with \(\mathcal L_u(\lambda')\ge c\)
      (Dirichlet-active at level \(c\)), \(r(\lambda')\ge2c/\eta(\lambda)\), so
      the mean Hermite degree of the profile is at least
      \(4c/\eta(\lambda)-1\to\infty\) as \(\lambda\to0\); and the Dirichlet-active
      scales at every level \(c<\limsup_{\lambda'\to0}\mathcal L_u(\lambda')\)
      accumulate at zero.
\item On every window \([\sigma_1,\sigma_2]\) of log-scale on which
      \(r\) increases from \(r_1\) to \(r_2\),
      \begin{equation}
       \boxed{\quad
       \int_{\sigma_1}^{\sigma_2}\frac2\varphi\bigl\langle L_-V+rV,N(V)\bigr\rangle_G\,\dd\sigma
       \ =\ (r_2-r_1)+2\int_{\sigma_1}^{\sigma_2}\operatorname{Var}\,\dd\sigma\ \ge\ r_2-r_1 :
       \quad}
       \label{eq:c2v31-push}
      \end{equation}
      the nonlinear push is at least the growth of the Rayleigh
      quotient, and exceeds it by twice the accumulated spectral
      variance.
\item If \(\mathcal L_u\le M_{\mathcal L}\) on small scales, then
      \(r\le2M_{\mathcal L}/\varphi\) and, on the Dirichlet-active scales at
      level \(c\), \(r\in[2c/\eta,\,2M_{\mathcal L}/\varphi]\).
\end{enumerate}
\end{theorem}

\begin{proof}
(i) \(r=2\mathcal L/\varphi\ge2c/\eta(\lambda)\) on the active scales below
\(\lambda\), and Lemma~\ref{lem:c2v31-degree}; the accumulation is the
definition of \(\limsup\) (and if \(\limsup\mathcal L_u=0\) the level sets
are empty for every \(c>0\), the non-integrability of
Proposition~\ref{prop:c2v31-nonintegrable} being then carried by
scales with \(\mathcal L_u\to0\)).  (ii) Integrate \eqref{eq:c2v31-identity}
with \(\operatorname{Var}\ge0\).  (iii) The bounds on \(r=2\mathcal L/\varphi\).
\end{proof}

\begin{firewall}[Item (AC2): what is and is not forced]
\label{fw:c2v31-AC2}
High-mode escape is forced on the Dirichlet-active scales of a
Gaussian-invisible paraboloid-dissipative point, and the nonlinear
push \eqref{eq:c2v31-push} is forced along every ascent of the Rayleigh
quotient; both are consequences, not obstructions.  The push is
\(\frac2\varphi\langle L_-V+rV,N(V)\rangle_G\), of which no bound is available
without a rate (\(N(V)\) and \(L_-V\) are not bounded in \(L^2(G)\) at a
fixed scale, Firewall~\ref{fw:c2v28-AC2}); so the series cannot say
that the nonlinearity is unable to supply it.  If \(\limsup\mathcal L_u=0\)
at the point, the Dirichlet activity is carried by scales with
vanishing \(\mathcal L_u\), the level sets are empty, and (i) is void
while Proposition~\ref{prop:c2v31-nonintegrable} still holds; whether
this case is compatible with paraboloid dissipation of a fixed
constant along a sequence is not decided (it requires the Dirichlet
functional to be small at the parabolic times while the paraboloid
dissipation over the window is not, which the time-averaging in
\eqref{eq:c2v27-dirichlet-sees} permits).
\end{firewall}

% =============================================================================
\section{Integrated V31 conclusion and next acquisition order}
\label{sec:c2v31-conclusion}
% =============================================================================

\begin{theorem}[What V31 proves]
\label{thm:c2v31-conclusion}
Relative to the first cycle and V1--V30, without a rate: the
Rayleigh-quotient identity \eqref{eq:c2v31-identity}, the
non-integrability of the Dirichlet functional at paraboloid-dissipative
points (Proposition~\ref{prop:c2v31-nonintegrable}), the degree lemma,
and the forced high-mode escape with the required push
\eqref{eq:c2v31-push}.  The full Navier--Stokes stability/global-regularity
theorem remains unproved.
\end{theorem}

\begin{proof}
The cited results.
\end{proof}

\begin{programme}[V32 acquisition order]
\label{prog:c2v32}
\begin{enumerate}[label=\textup{(AC\arabic*)},leftmargin=4.8em]
\item Paraboloid dissipation from cylinder dissipation.  Under compact
      dissipation with a coherent centre map (V24), determine whether
      the moving-centre paraboloids capture a fixed fraction of the slab
      dissipation, so that the moving-centre Dirichlet functional obeys
      \eqref{eq:c2v27-dirichlet-sees} along the centre map; record the
      exact statement.
\item The nonlinear push in Lamb--Bernoulli form.  Write
      \(\langle L_-V+rV,N(V)\rangle_G\) with \(N=\Omega\times V+\nabla\Pi\) and
      determine which part can push the profile up the spectrum: the
      Bernoulli part is \(\frac1{2\nu}\int\Pi\,z\cdot(L_-V+rV)G\) after
      integration by parts (using \(\operatorname{div}(L_-V+rV)=0\)), and
      the Lamb part is \(\int\det(\Omega,V,L_-V+rV)G\).
\item The next compulsory companion audit is V35.
\end{enumerate}
\end{programme}

\begin{assessment}[Position after V31]
\label{ass:c2v31-position}
The Rayleigh quotient obeys an exact descent-plus-push identity, and
Gaussian-invisible paraboloid-dissipative blow-up must escape to high
Hermite modes with a nonlinear push at least equal to the ascent.  No
full stability or global-regularity claim is made.
\end{assessment}

% =============================================================================
\section*{Version V31 (second cycle) append-only record}
\addcontentsline{toc}{section}{Version V31 (second cycle) append-only record}
% =============================================================================

The complete V30 source of the second cycle was retained before this
continuation.  Its SHA--256 digest was
\begin{center}\scriptsize\ttfamily
4035533fc8d9a4264bf27f72e4d243f8f971bb39ed22f65c34d78b3662065e13.
\end{center}
V31 proved the Rayleigh-quotient identity and the forced high-mode
escape of Gaussian-invisible paraboloid-dissipative blow-up.  No full
regularity claim is made.

% =============================================================================
\section{V32: paraboloid dissipation from compact dissipation, and the
nonlinear push in Lamb--Bernoulli form}
\label{sec:c2v32-entry}
% =============================================================================

Items (AC1) and (AC2) of Programme~\ref{prog:c2v32} are carried out.  It
is shown that compact dissipation at a scale, with the ball radius at
most twice the square root of the lower time offset of the window,
places a fixed fraction of the slab dissipation inside the backward
paraboloid of a comparable scale at one of the centres, so that the
Dirichlet log-integral at that centre below that scale is bounded below
without any coherence of the centre map; coherence is needed only to
follow one centre across scales.  The nonlinear push of the Rayleigh
quotient is written in Lamb--Bernoulli form: a Lamb triple product of
vorticity, velocity, and the self-similar operator applied to the
profile, a Bernoulli term against the radial component of that image,
and a flux term proportional to the Rayleigh quotient; the flux term is
eliminated through the log-derivative of the weighted energy, which
gives a closed identity for the Rayleigh quotient in terms of the
Gaussian norm of the operator image, the weighted energy, and the two
geometric terms.  No global regularity claim is made.  The next
compulsory companion audit is V35.

% =============================================================================
\section{Paraboloid dissipation from compact dissipation}
\label{sec:c2v32-paraboloid}
% =============================================================================

\begin{theorem}[Compact dissipation gives paraboloid dissipation at a centre]
\label{thm:c2v32-paraboloid}
Let \(u\) be a smooth solution with \(T_*<\infty\), let the slab
\(\Sigma_\lambda=(T_*-S\lambda^2,T_*-\eta\lambda^2)\) and the compactness parameters
\((\theta,N,R)\) be as in Definition~\ref{def:c2v24-compact}, and assume
\(R\le2\eta^{1/2}\).  If the dissipation is \((\theta,N,R)\)-compact at
scale \(\lambda\) with centres \(x_1,\dots,x_N\), then some centre
\(x_i\) is paraboloid-dissipative at scale \(S^{1/2}\lambda\) with constant
\begin{equation}
 \varepsilon':=\frac{2\theta c_L^2\nu^{3/2}(S^{1/2}-\eta^{1/2})}{N\,S^{1/2}} ,
 \label{eq:c2v32-epsprime}
\end{equation}
and consequently, without any rate and without coherence,
\begin{equation}
 \boxed{\quad
 \int_0^{S^{1/2}\lambda}\mathcal L_{u,x_i}(\lambda')\,\frac{\dd\lambda'}{\lambda'}\ \ge\ \frac{\nu e^{-1/\nu}}4\,\varepsilon' .
 \quad}
 \label{eq:c2v32-dirichlet}
\end{equation}
If the dissipation is compact at a set of scales of positive
logarithmic density, the centres \(x_i(\lambda)\) so obtained satisfy
\eqref{eq:c2v32-dirichlet} on that set.
\end{theorem}

\begin{proof}
For \(t\in\Sigma_\lambda\), \(T_*-t\ge\eta\lambda^2\), so \(2(T_*-t)^{1/2}\ge2\eta^{1/2}\lambda\ge R\lambda\)
and \(B_{R\lambda}(x_i)\times\Sigma_\lambda\subset P_{S^{1/2}\lambda}(x_i)\) (the
paraboloid of scale \(S^{1/2}\lambda\) covers the time window
\((T_*-S\lambda^2,T_*)\supset\Sigma_\lambda\)).  By the pigeonhole principle one
ball carries at least \(\theta\mathcal D(\lambda)/N\), and by
Lemma~\ref{lem:c2v24-slab} \(\mathcal D(\lambda)\ge2c_L^2\nu^{3/2}(S^{1/2}-\eta^{1/2})\lambda\);
writing \(\lambda=S^{-1/2}\cdot S^{1/2}\lambda\) gives the constant
\eqref{eq:c2v32-epsprime} in the definition of paraboloid dissipation at
scale \(S^{1/2}\lambda\).  Then Theorem~\ref{thm:c2v27-dirichlet-sees}.
\end{proof}

\begin{assessment}[Coherence is needed only to follow a centre]
\label{ass:c2v32-coherence}
Theorem~\ref{thm:c2v32-paraboloid} removes the coherence hypothesis from
the Dirichlet ledger at a single scale: wherever the dissipation is
compact, some centre sees, below that scale, a fixed Dirichlet
log-integral.  What coherence adds is the ability to integrate the
second identity along one centre across many scales, which is what the
moving-centre blow-up measures of V24 and the ledgers of V28 and V31
require.  Without coherence the series has, at each compact scale, a
centre with a fixed Dirichlet activity below it, but the centres may
change from scale to scale by more than the parabolic distance, and
the ledger identities, which are exact only along a fixed centre, do
not chain.  The first cycle's count of Gaussian-active centres per
scale (at most \(N_c\)) bounds how many such centres there are at a
scale, not how they move.
\end{assessment}

% =============================================================================
\section{The push in Lamb--Bernoulli form}
\label{sec:c2v32-push}
% =============================================================================

\begin{proposition}[Lamb--Bernoulli form of the nonlinear push]
\label{prop:c2v32-push}
With \(N(V)=\Omega\times V+\nabla\Pi\), \(\Pi=P_V+\frac12|V|^2\), and
\(j=4\langle V,N(V)\rangle_G=\frac1\nu\int(|V|^2+2P_V)(z\cdot V)G\), for every
\(\sigma\),
\begin{equation}
 \bigl\langle L_-V+rV,\,N(V)\bigr\rangle_G
 =\int\det\bigl(\Omega,V,L_-V\bigr)G
 +\frac1{2\nu}\int\Pi\,\bigl(z\cdot L_-V\bigr)G
 +\frac r4\,j ,
 \label{eq:c2v32-push}
\end{equation}
and the Rayleigh-quotient identity takes the closed form
\begin{equation}
 \boxed{\quad
 \frac{\dd r}{\dd\sigma}=-2\|L_-\widehat V\|_G^2-r\,\frac{\dd}{\dd\sigma}\log\varphi
 +\frac2\varphi\Bigl[\int\det(\Omega,V,L_-V)G+\frac1{2\nu}\int\Pi\,(z\cdot L_-V)G\Bigr].
 \quad}
 \label{eq:c2v32-closed}
\end{equation}
\end{proposition}

\begin{proof}
\(\langle L_-V+rV,\Omega\times V\rangle_G=\int\det(\Omega,V,L_-V+rV)G=\int\det(\Omega,V,L_-V)G\)
since \(\det(\Omega,V,V)=0\).  For the Bernoulli part, \(W:=L_-V+rV\) is
divergence-free (\(L_-\) preserves the divergence-free condition,
Proposition~\ref{prop:c2v8-spectrum}(ii)), so
\(\langle W,\nabla\Pi\rangle_G=-\int\Pi\operatorname{div}(WG)=-\int\Pi\,W\cdot\nabla G=\frac1{2\nu}\int\Pi(z\cdot W)G\),
and \(\frac1{2\nu}\int\Pi(z\cdot rV)G=\frac r{2\nu}\cdot\frac\nu2j=\frac r4j\) by
Proposition~\ref{prop:c2v4-Bernoulli}'s identity \(\int\Pi(z\cdot V)G=\frac\nu2j\).
For \eqref{eq:c2v32-closed}, insert \eqref{eq:c2v32-push} into
\eqref{eq:c2v31-identity}, use \(\operatorname{Var}=\|L_-\widehat V\|_G^2-r^2\) and
\(\frac2\varphi\cdot\frac r4j=\frac r2\cdot\frac j\varphi\) with
\(\frac j\varphi=-4r-2\frac{\dd}{\dd\sigma}\log\varphi\) from \eqref{eq:c2v29-logphi}
(\(\frac12\varphi^{1/2}\widehat j=\frac j{2\varphi}\)); the terms \(2r^2\) cancel.
\end{proof}

\begin{assessment}[Reading the closed identity]
\label{ass:c2v32-closed}
The closed identity separates three effects on the Rayleigh quotient:
the linear descent \(-2\|L_-\widehat V\|_G^2\), which is at most \(-2r^2\);
the coupling \(-r\,\frac{\dd}{\dd\sigma}\log\varphi\), positive when the weighted
energy decreases (a decaying weighted energy raises the Rayleigh
quotient, which is the mechanism by which Gaussian invisibility forces
high-mode escape, V31); and the geometric push, the Gaussian
correlation of the Lamb vector with the operator image of the profile
and of the Bernoulli pressure with its radial component.  Under
Gaussian invisibility with bounded Dirichlet functional,
\(-r\frac{\dd}{\dd\sigma}\log\varphi\) has log-integral bounded below by
\(-\int r\,\dd\log\varphi\), and since \(r\ge\frac12\) this alone accounts
for a growth of \(r\) of at least half the decay of \(\log\varphi\); the
geometric push and the linear descent must balance the rest.  No sign
of the geometric push is available.
\end{assessment}

% =============================================================================
\section{Integrated V32 conclusion and next acquisition order}
\label{sec:c2v32-conclusion}
% =============================================================================

\begin{theorem}[What V32 proves]
\label{thm:c2v32-conclusion}
Relative to the first cycle and V1--V31, without a rate: paraboloid
dissipation and the Dirichlet lower bound \eqref{eq:c2v32-dirichlet} at
a centre of compact dissipation, and the Lamb--Bernoulli form
\eqref{eq:c2v32-push} with the closed Rayleigh identity
\eqref{eq:c2v32-closed}.  The full Navier--Stokes
stability/global-regularity theorem remains unproved.
\end{theorem}

\begin{proof}
Theorem~\ref{thm:c2v32-paraboloid}, Proposition~\ref{prop:c2v32-push}.
\end{proof}

\begin{programme}[V33 acquisition order]
\label{prog:c2v33}
\begin{enumerate}[label=\textup{(AC\arabic*)},leftmargin=4.8em]
\item The invisible Dirichlet-bounded case in closed form.  Integrate
      \eqref{eq:c2v32-closed} over a window on which \(\varphi\) decreases
      from \(\varphi_1\) to \(\varphi_2\) and \(r\) is bounded by
      \(2M_{\mathcal L}/\varphi\): derive the resulting inequality between the
      log-integral of \(\|L_-\widehat V\|_G^2\), the quantity
      \(\int r\,\dd(-\log\varphi)\), and the geometric push; determine
      whether the linear descent, which is at least \(2r^2\), can be
      absorbed, that is, whether \(\int2r^2\dd\sigma\le\int r\,\dd(-\log\varphi)+\)push
      forces a lower bound on the push in terms of \(M_{\mathcal L}\) and
      the decay of \(\varphi\).
\item Record what the geometric push needs to supply in the extreme
      case \(r=2M_{\mathcal L}/\varphi\) (Dirichlet functional saturating its
      bound on invisible scales).
\item The next compulsory companion audit is V35.
\end{enumerate}
\end{programme}

\begin{assessment}[Position after V32]
\label{ass:c2v32-position}
Compact dissipation gives paraboloid dissipation and a fixed Dirichlet
activity at some centre without coherence, and the Rayleigh quotient
obeys a closed identity with an explicit geometric push.  No full
stability or global-regularity claim is made.
\end{assessment}

% =============================================================================
\section*{Version V32 (second cycle) append-only record}
\addcontentsline{toc}{section}{Version V32 (second cycle) append-only record}
% =============================================================================

The complete V31 source of the second cycle was retained before this
continuation.  Its SHA--256 digest was
\begin{center}\scriptsize\ttfamily
a2147744fbf08df9389afccb58a35af61ad9fc5198615cfe73f67badf7964130.
\end{center}
(The V31 file was corrected in place after its assembly by restoring
the V28 firewall on the rate-free defect bound, which an editing step
had dropped; the digest above is that of the corrected file.)  V32
proved paraboloid dissipation at a centre of compact dissipation and
the Lamb--Bernoulli form of the nonlinear push with the closed Rayleigh
identity.  No full regularity claim is made.

% =============================================================================
\section{V33: the push budget of the Rayleigh quotient}
\label{sec:c2v33-entry}
% =============================================================================

Items (AC1) and (AC2) of Programme~\ref{prog:c2v33} are carried out.  The
closed Rayleigh identity is integrated over any window of log-scale,
and the coupling term is resolved through the log-derivative of the
weighted energy: the required geometric push equals the growth of the
Rayleigh quotient plus twice the accumulated spectral variance plus the
Rayleigh-weighted inward normalized flux.  The linear descent is thus
absorbed exactly, not by the decay of the weighted energy but by the
flux term of the first identity, and the question of whether the
nonlinearity can supply the push is the question of the sign and size
of the Lamb and Bernoulli correlations with the operator image of the
profile.  The extreme case of a Dirichlet functional saturating its
bound is computed, and the mean form against invariant measures at
homogeneous points is recorded: the mean geometric push equals twice
the mean spectral variance minus the mean Rayleigh-weighted normalized
flux.  No global regularity claim is made.  The next compulsory
companion audit is V35.

% =============================================================================
\section{The push budget}
\label{sec:c2v33-budget}
% =============================================================================

With the notation of V29--V32, let
\begin{equation}
 \mathcal P:=\int\det(\Omega,V,L_-V)\,G+\frac1{2\nu}\int\Pi\,(z\cdot L_-V)\,G
 \label{eq:c2v33-P}
\end{equation}
be the geometric push, so that \(\langle L_-V+rV,N(V)\rangle_G=\mathcal P+\frac r4j\)
by \eqref{eq:c2v32-push}, and let \(\widehat\jmath:=j/\varphi\) be the
normalized flux (so \(\frac{\dd}{\dd\sigma}\log\varphi=-2r-\frac12\widehat\jmath\)).

\begin{theorem}[Push budget]
\label{thm:c2v33-budget}
For every window \([\sigma_1,\sigma_2]\) of log-scale, with
\(r_i=r(\sigma_i)\),
\begin{equation}
 \boxed{\quad
 \int_{\sigma_1}^{\sigma_2}\frac{2\mathcal P}{\varphi}\,\dd\sigma
 \ =\ (r_2-r_1)+2\int_{\sigma_1}^{\sigma_2}\operatorname{Var}\,\dd\sigma
 +\int_{\sigma_1}^{\sigma_2}\frac{r\,(-\widehat\jmath)}{2}\,\dd\sigma .
 \quad}
 \label{eq:c2v33-budget}
\end{equation}
In particular, wherever the normalized flux is inward
(\(\widehat\jmath<0\)) the geometric push must exceed the growth of the
Rayleigh quotient plus twice the accumulated spectral variance, and the
linear descent \(-2\|L_-\widehat V\|_G^2=-2\operatorname{Var}-2r^2\) is absorbed
exactly by the term \(2r^2\) hidden in the coupling
\(-r\frac{\dd}{\dd\sigma}\log\varphi=2r^2+\frac r2\widehat\jmath\).
\end{theorem}

\begin{proof}
Integrate \eqref{eq:c2v32-closed}, substitute
\(-r\frac{\dd}{\dd\sigma}\log\varphi=2r^2+\frac r2\widehat\jmath\) and
\(\|L_-\widehat V\|_G^2=\operatorname{Var}+r^2\); the \(2r^2\) terms cancel.
Equivalently, \eqref{eq:c2v33-budget} is \eqref{eq:c2v31-push} with the
push split as \(\frac2\varphi(\mathcal P+\frac r4j)=\frac{2\mathcal P}\varphi+\frac r2\widehat\jmath\).
\end{proof}

\begin{corollary}[Extreme case and the invisible Dirichlet-bounded case]
\label{cor:c2v33-extreme}
\begin{enumerate}[label=\textup{(\roman*)},leftmargin=3.2em]
\item If \(\mathcal L_u\le M_{\mathcal L}\) on the window and
      \(r=2M_{\mathcal L}/\varphi\) there (the Dirichlet functional
      saturates its bound), then
      \(\int\frac r2(-\widehat\jmath)\dd\sigma=M_{\mathcal L}\int\frac{-j}{\varphi^2}\dd\sigma\)
      and \eqref{eq:c2v33-budget} reads
      \(\int\frac{2\mathcal P}\varphi=2M_{\mathcal L}(\varphi_2^{-1}-\varphi_1^{-1})+2\int\operatorname{Var}+M_{\mathcal L}\int\frac{-j}{\varphi^2}\).
\item At a paraboloid-dissipative, Gaussian-invisible, Dirichlet-bounded
      point, on any window on which the normalized flux is inward on
      average, the log-integral of the geometric push over the
      Dirichlet-active scales is bounded below by the growth of \(r\)
      there, which by Theorem~\ref{thm:c2v31-forced}(i) is at least
      \(2c/\eta(\lambda)-r_1\) on windows starting at a scale with
      \(r_1=O(1)\) and ending at an active scale below \(\lambda\).
\end{enumerate}
\end{corollary}

\begin{proof}
(i) Substitute \(r=2M_{\mathcal L}/\varphi\) and \(\widehat\jmath=j/\varphi\); the
growth term is \(r_2-r_1=2M_{\mathcal L}(\varphi_2^{-1}-\varphi_1^{-1})\).
(ii) \eqref{eq:c2v33-budget} with \(\operatorname{Var}\ge0\) and
\(\int\frac r2(-\widehat\jmath)\ge0\) on windows of inward average
normalized flux (\(r\ge\frac12\) and the sign), and
Theorem~\ref{thm:c2v31-forced}(i).
\end{proof}

\begin{assessment}[Item (AC1): what absorbs the linear descent]
\label{ass:c2v33-absorb}
The programme asked whether the linear descent of the Rayleigh
quotient can be absorbed by the decay of the weighted energy.  The
budget \eqref{eq:c2v33-budget} answers exactly: the decay of the
weighted energy is itself driven by \(2r\) plus the normalized flux,
and the part \(2r^2\) of the coupling cancels the part \(2r^2\) of the
descent identically; what remains of the descent is twice the
spectral variance, and what remains of the coupling is the
Rayleigh-weighted normalized flux.  Hence the geometric push is
determined by three terms none of which the series can bound without
a rate: the growth of the Rayleigh quotient, the spectral variance,
and the Rayleigh-weighted flux.  With an inward flux, which the
ledgers favour at Gaussian-active scales, all three are nonnegative
and the push is forced to be positive with a divergent log-integral
along high-mode escape; with an outward flux the third term is
negative and can offset the first two.  The exact sign structure of
Gaussian-invisible blow-up is therefore: either the normalized flux is
outward on the escaping scales, or the Lamb and Bernoulli correlations
with the operator image of the profile are positive with a divergent
log-integral.
\end{assessment}

% =============================================================================
\section{Mean form at homogeneous points}
\label{sec:c2v33-mean}
% =============================================================================

\begin{proposition}[Mean push budget]
\label{prop:c2v33-mean}
Let \(x_0\) be a homogeneous singular point of a type-I solution and
\(\mathfrak m\) an invariant probability measure on \(\mathcal F(x_0)\).  Then
\(r\), \(\operatorname{Var}\), \(\widehat\jmath\), and \(\mathcal P/\varphi\) are
bounded on the family, and
\begin{equation}
 \int\frac{2\mathcal P}\varphi\,\dd\mathfrak m=2\int\operatorname{Var}\,\dd\mathfrak m+\int\frac{r(-\widehat\jmath)}2\,\dd\mathfrak m ,
 \qquad
 \int\frac{r(-\widehat\jmath)}2\,\dd\mathfrak m\ \ge\ \frac14\int(-\widehat\jmath)_+\,\dd\mathfrak m-\frac{C_r}{2}\int(\widehat\jmath)_+\,\dd\mathfrak m ,
 \label{eq:c2v33-mean}
\end{equation}
with \(C_r:=(4\nu C+2C_\Phi)/(4c_\Phi)\) the bound on \(r\).  The unweighted
mean flux satisfies \(\int j\,\dd\mathfrak m=-\int(d+2\varphi)\dd\mathfrak m<0\),
but the weighted flux \(\int r\widehat\jmath\,\dd\mathfrak m\) has no sign from
this.
\end{proposition}

\begin{proof}
Boundedness: \(r=2\mathcal L/\varphi\le C_r\) by the pinching and the class
bound, \(\operatorname{Var}\le\|L_-\widehat V\|_G^2\le(\nu C_\infty''\|1\|_G+\ldots)^2/c_\Phi\)
by the uniform derivative bounds, \(|\widehat\jmath|\le C_J/c_\Phi\), and
\(|\mathcal P|/\varphi\) by the same bounds.  The mean identity is the
invariant-measure argument applied to the bounded observable \(r\)
in \eqref{eq:c2v33-budget}.  The inequality splits \(\widehat\jmath\) into
positive and negative parts with \(\frac12\le r\le C_r\).  The last
statement is \eqref{eq:c2v12-mean1}.
\end{proof}

\begin{firewall}[Item (AC2): the weighted flux has no mean sign]
\label{fw:c2v33-AC2}
The first cycle's inward flux, \(\int j\,\dd\mathfrak m<0\), does not
transfer to the Rayleigh-weighted normalized flux
\(\int r\widehat\jmath\,\dd\mathfrak m\), because \(r\) and \(1/\varphi\) are not
constant on the family; only the bound \eqref{eq:c2v33-mean} is
available.  Hence the mean geometric push at a homogeneous point is
determined up to the sign of the weighted flux, and the series records
\(\int\frac{2\mathcal P}\varphi\dd\mathfrak m=2\langle\operatorname{Var}\rangle+\langle\frac r2(-\widehat\jmath)\rangle\)
as an exact balance with one signed term (the variance) and one
unsigned term.  An element \(W\ne0\) of the family has \(\operatorname{Var}>0\) at every
scale: \(\operatorname{Var}=0\) at a scale means that \(\widehat V\) is an
eigenfunction of \(L_-\) in \(L^2(G)\), hence a polynomial vector field
(the eigenfunctions of the Ornstein--Uhlenbeck operator are Hermite
polynomials), which lies in \(L^6(\R^3)\) only if it is constant, and
then \(V\), hence \(W\) at that time, is a constant in \(L^6\), that is
zero, so \(W\equiv0\) by (E8) and (E6).  Consequently
\(\int\operatorname{Var}\,\dd\mathfrak m>0\) for every invariant measure
\(\mathfrak m\ne\delta_0\) (the observable is nonnegative, Borel as a
pointwise limit of continuous truncations, and positive on the
\(\mathfrak m\)-positive set of nonzero elements); the series records this
positivity without an explicit constant, which is item (AC1) of the
next version.
\end{firewall}

% =============================================================================
\section{Integrated V33 conclusion and next acquisition order}
\label{sec:c2v33-conclusion}
% =============================================================================

\begin{theorem}[What V33 proves]
\label{thm:c2v33-conclusion}
Relative to the first cycle and V1--V32: the push budget
\eqref{eq:c2v33-budget} (rate-free), its extreme and invisible cases
(Corollary~\ref{cor:c2v33-extreme}), and the mean push budget
\eqref{eq:c2v33-mean} at homogeneous points (modulo (E1)--(E9)).  The
full Navier--Stokes stability/global-regularity theorem remains
unproved.
\end{theorem}

\begin{proof}
The cited results.
\end{proof}

\begin{programme}[V34 acquisition order]
\label{prog:c2v34}
\begin{enumerate}[label=\textup{(AC\arabic*)},leftmargin=4.8em]
\item Hermite-mode profiles.  Make precise the statement that no
      element of a rescaling family is an eigenfunction of \(L_-\) at
      any scale: the eigenfunctions of \(L_-\) on divergence-free
      \(L^2(G)\) are polynomial fields, and a polynomial field in the
      class must be constant, hence zero by the vorticity at infinity;
      conclude \(\operatorname{Var}>0\) pointwise on the family and, by
      compactness and continuity, \(\operatorname{Var}\ge v_0>0\) uniformly
      on \(\mathcal F(x_0)\) at a homogeneous point; record whether
      \(\operatorname{Var}\) is continuous in the local topology.
\item With a uniform \(v_0\), the mean push budget gives
      \(\int\frac{2\mathcal P}\varphi\dd\mathfrak m\ge2v_0-\frac{C_r}2\int(\widehat\jmath)_+\dd\mathfrak m\):
      determine whether the outward part of the normalized flux can be
      bounded in the mean by the first cycle's inward-flux identity, so
      that the mean geometric push is bounded below by a positive
      constant at every homogeneous point.
\item The next compulsory companion audit is V35.
\end{enumerate}
\end{programme}

\begin{assessment}[Position after V33]
\label{ass:c2v33-position}
The Rayleigh quotient's push budget is exact: growth plus twice the
spectral variance plus the Rayleigh-weighted inward normalized flux.
No full stability or global-regularity claim is made.
\end{assessment}

% =============================================================================
\section*{Version V33 (second cycle) append-only record}
\addcontentsline{toc}{section}{Version V33 (second cycle) append-only record}
% =============================================================================

The complete V32 source of the second cycle was retained before this
continuation.  Its SHA--256 digest was
\begin{center}\scriptsize\ttfamily
34ab84c751d09ed83285a5d9d82c7e40013c0ebe06e958fe06255fcc9554f83d.
\end{center}
V33 proved the push budget of the Rayleigh quotient and its mean form
at homogeneous points.  No full regularity claim is made.

% =============================================================================
\section{V34: uniform spectral variance on the rescaling family and the
mean geometric push}
\label{sec:c2v34-entry}
% =============================================================================

Items (AC1) and (AC2) of Programme~\ref{prog:c2v34} are carried out.  No
element of a rescaling family is an eigenfunction of the self-similar
operator at any scale, because such eigenfunctions are polynomial
fields and the only polynomial field in the class is zero; the spectral
variance is therefore positive on every nonzero element, lower
semicontinuous in the local topology, and hence bounded below by a
positive constant uniformly on the compact family at a homogeneous
point.  The mean push budget then bounds the mean geometric push below
by twice this constant minus an explicit multiple of the mean outward
normalized flux, and the outward part is bounded only by the order-one
bound on the flux, so positivity of the mean geometric push follows
when the mean outward normalized flux is small, a condition recorded
as open.  No global regularity claim is made.  The next version is the
compulsory companion audit.

% =============================================================================
\section{No Hermite-mode profiles}
\label{sec:c2v34-hermite}
% =============================================================================

\begin{lemma}[Eigenfunctions of the self-similar operator are polynomial fields]
\label{lem:c2v34-eigen}
The eigenfunctions of \(-L_-=-\nu\Delta+\frac12z\cdot\nabla+\frac12\) in
\(L^2(G)\) are the vector fields whose components are Hermite polynomials
in \(z/(2\nu)^{1/2}\), with eigenvalue \(\frac12+\frac k2\) for total
degree \(k\); a divergence-free eigenfunction is a divergence-free
polynomial field.  A polynomial vector field belongs to \(L^6(\R^3)\)
only if it is constant, and to \(L^6\) with zero mean over large balls
only if it is zero.
\end{lemma}

\begin{proof}
Proposition~\ref{prop:c2v8-spectrum}(i) and the classical
diagonalization of the Ornstein--Uhlenbeck operator by Hermite
polynomials, componentwise; a nonconstant polynomial grows at
infinity and is not in \(L^6\); a nonzero constant is in no \(L^p(\R^3)\).
\end{proof}

\begin{theorem}[Uniform positive spectral variance on the family]
\label{thm:c2v34-variance}
Let \(x_0\) be a homogeneous singular point of a type-I solution.  Then,
modulo (E1)--(E9):
\begin{enumerate}[label=\textup{(\roman*)},leftmargin=3.2em]
\item For every \(W\in\mathcal F(x_0)\) and every scale,
      \(\operatorname{Var}(W)>0\): no rescaling limit is a Hermite-mode
      profile at any scale.
\item \(W\mapsto\operatorname{Var}(W)\) (at the unit scale) is lower
      semicontinuous on \((\mathcal F(x_0),\rho)\).
\item There is \(v_0=v_0(x_0)>0\) with
      \begin{equation}
       \boxed{\quad
       \operatorname{Var}(\Theta_\tau W)\ \ge\ v_0\qquad\text{for all }W\in\mathcal F(x_0),\ \tau\in\R .
       \quad}
       \label{eq:c2v34-v0}
      \end{equation}
\end{enumerate}
\end{theorem}

\begin{proof}
(i) If \(\operatorname{Var}(W)=0\) at some scale, \(\widehat V\) is an
eigenfunction of \(L_-\) at that scale, hence a polynomial field by
Lemma~\ref{lem:c2v34-eigen}; \(V\in L^6(\R^3)\) (class bound) forces
\(V\) constant, hence \(V=0\), so \(W\) vanishes at that time and
\(W\equiv0\) by (E8) and (E6); but every element of \(\mathcal F(x_0)\) is
nonzero at a homogeneous point (pinching).  (ii) \(\operatorname{Var}=\|L_-V\|_G^2/\varphi-r^2\)
with \(\varphi\) and \(r=2\mathcal L/\varphi\) continuous on the family
(Lemma~\ref{lem:c2v12-measurable}); and \(W\mapsto\|L_-V(0)\|_G^2\) is
lower semicontinuous: \(L_-V_k=\nu\Delta V_k-\frac12z\cdot\nabla V_k-\frac12V_k\)
is bounded in \(L^2(G)\) uniformly on the family by the uniform
derivative bounds \eqref{eq:c2v4-uniform}, so along a locally
convergent sequence it converges weakly in \(L^2(G)\) (the limit is
identified in distributions), and the norm is weakly lower
semicontinuous.  (iii) A positive lower semicontinuous function on a
compact metric space attains its infimum, which is therefore positive;
scale invariance of the family transports the bound to every scale.
\end{proof}

\begin{remark}[Dependence of \(v_0\)]
\label{rem:c2v34-v0}
\(v_0\) is produced by compactness and depends on the family, as
\(c_\Phi\) and \(\underline\delta\) do.  Its meaning is quantitative
non-Hermiteness: at every scale the normalized profile of every
rescaling limit has Hermite-degree distribution of variance at least
\(4v_0\) (Lemma~\ref{lem:c2v31-degree}).
\end{remark}

% =============================================================================
\section{The mean geometric push}
\label{sec:c2v34-push}
% =============================================================================

\begin{theorem}[Mean geometric push at a homogeneous point]
\label{thm:c2v34-mean-push}
Let \(x_0\) be homogeneous, \(v_0\) as in \eqref{eq:c2v34-v0}, \(C_r\) the
bound on the Rayleigh quotient, and \(\mathfrak m\) any invariant
probability measure on \(\mathcal F(x_0)\).  Then
\begin{equation}
 \boxed{\quad
 \int\frac{2\mathcal P}{\varphi}\,\dd\mathfrak m\ \ge\ 2v_0+\frac14\int(-\widehat\jmath)_+\,\dd\mathfrak m-\frac{C_r}2\int(\widehat\jmath)_+\,\dd\mathfrak m ,
 \qquad
 \int(\widehat\jmath)_+\,\dd\mathfrak m\ \le\ \frac{C_J}{c_\Phi} .
 \quad}
 \label{eq:c2v34-mean-push}
\end{equation}
In particular the mean geometric push is positive whenever
\(\int(\widehat\jmath)_+\dd\mathfrak m<4v_0/C_r\).
\end{theorem}

\begin{proof}
\eqref{eq:c2v33-mean} with \(\int\operatorname{Var}\dd\mathfrak m\ge v_0\); the
bound on the outward part is \(|\widehat\jmath|=|j|/\varphi\le C_J/c_\Phi\).
\end{proof}

\begin{firewall}[Item (AC2): the outward part of the normalized flux]
\label{fw:c2v34-AC2}
The first cycle's inward-flux identity bounds the mean of \(j\), not of
its positive part: \(\int j\,\dd\mathfrak m=-\int(d+2\varphi)\dd\mathfrak m\) is
compatible with a positive part of mean up to \(C_J\) if the negative
part compensates.  Hence \eqref{eq:c2v34-mean-push} does not give a
positive lower bound on the mean geometric push unconditionally, and
the condition \(\int(\widehat\jmath)_+\dd\mathfrak m<4v_0/C_r\) (the blow-up
statistics is outward-flux-poor) is recorded as open.  If it holds, the
Lamb and Bernoulli correlations with the operator image of the profile
have positive mean, which is a fourth signed statement for the
homogeneous case, of the same nature as the three of V8: a balance the
nonlinearity must satisfy, not an obstruction.
\end{firewall}

% =============================================================================
\section{Integrated V34 conclusion and next acquisition order}
\label{sec:c2v34-conclusion}
% =============================================================================

\begin{theorem}[What V34 proves]
\label{thm:c2v34-conclusion}
Relative to the first cycle and V1--V33, modulo (E1)--(E9): the absence
of Hermite-mode profiles and the uniform positive spectral variance
\eqref{eq:c2v34-v0} on the rescaling family at a homogeneous point,
and the mean geometric push bound \eqref{eq:c2v34-mean-push}.  The full
Navier--Stokes stability/global-regularity theorem remains unproved.
\end{theorem}

\begin{proof}
Theorems~\ref{thm:c2v34-variance}, \ref{thm:c2v34-mean-push}.
\end{proof}

\begin{programme}[V35 acquisition order, with compulsory companion audit]
\label{prog:c2v35}
\begin{enumerate}[label=\textup{(AC\arabic*)},leftmargin=4.8em]
\item Perform the compulsory review of the current SM and YM notes;
      record digests; import nothing numerical.
\item Consolidate V31--V34 as the Rayleigh chapter: the identity, the
      forced high-mode escape, paraboloid dissipation from compact
      dissipation, the Lamb--Bernoulli push and the closed identity,
      the push budget, the uniform spectral variance, and the mean
      geometric push.
\item Record the consolidated statement of the second cycle's first
      thirty-five versions by chapter (self-similar variables,
      dichotomy, statistics, moment ledger, rate-free chapter, Rayleigh
      chapter), the hypotheses and the open conditions, and the
      direction for V36--V40.
\item The next compulsory companion audit after V35 is V40.
\end{enumerate}
\end{programme}

\begin{assessment}[Position after V34]
\label{ass:c2v34-position}
The spectral variance is uniformly positive on the rescaling family at
a homogeneous point, and the mean geometric push is bounded below by
twice that constant up to the mean outward normalized flux.  No full
stability or global-regularity claim is made.
\end{assessment}

% =============================================================================
\section*{Version V34 (second cycle) append-only record}
\addcontentsline{toc}{section}{Version V34 (second cycle) append-only record}
% =============================================================================

The complete V33 source of the second cycle was retained before this
continuation.  Its SHA--256 digest was
\begin{center}\scriptsize\ttfamily
5eb670982e2ca9153a82a621a5e9b6c6e71c178c3f8137dea8f691222c8c8c69.
\end{center}
V34 proved the uniform positive spectral variance on the rescaling
family and the mean geometric push bound.  No full regularity claim is
made.

% =============================================================================
\section{V35: compulsory companion audit, the Rayleigh chapter consolidated,
and the second cycle after thirty-five versions}
\label{sec:c2v35-entry}
% =============================================================================

This is the thirty-fifth version of the second cycle.  The compulsory
review of the two companion programmes is performed first.  V31--V34
are then consolidated as the Rayleigh chapter.  The second cycle's
first thirty-five versions are then summarized by chapter, with the
hypotheses under which each chapter's signed statements hold and the
conditions recorded as open, and the direction for V36--V40 is fixed.
No global regularity claim is made.  The next compulsory companion
audit is V40.

% =============================================================================
\section{Compulsory V35 review of the current companion notes}
\label{sec:c2v35-companion-audit}
% =============================================================================

The current companion files in \texttt{latex\_notes} were inspected
read-only.  The frozen checkpoints are
\begin{align}
 &\texttt{Renewal\_SM\_Einstein\_Continuum\_Research\_Notes\_V95.tex},
 \notag\\[-2pt]
 &\qquad
 \texttt{b67ef00b3a7b43348c3da97752b123c910ecefdd2e726ee79079c0b4a310c808},
 \label{eq:c2v35-SM-checkpoint}\\
 &\texttt{Renewal\_Yang\_Mills\_Mass\_Gap\_Research\_Notes\_V70.tex},
 \notag\\[-2pt]
 &\qquad
 \texttt{2ce7f53627fb326fb66fcb394020698886be4002e9c668baf60ed5a7dcab61d9}.
 \label{eq:c2v35-YM-checkpoint}
\end{align}

\emph{SM V95.}  The SM programme performed its compulsory checkpoint,
recording the present cycle's V29 by digest and importing, in its
words, only the ``assemble-before-estimating and
vanishing-normalization warnings'' (the latter being this cycle's
record that the normalization by a weighted energy tending to zero
destroys the class bounds, Proposition~\ref{prop:c2v29-relative}); it
then proved a free periodic overlap zero-mode identity, a zero-Higgs
no-floor theorem, the absence of a volume-uniform floor for an
antiperiodic spin choice, an exact massive-overlap spectral circle, a
sharp configuration-uniform lower singular-value bound, positivity of
a vectorlike massive determinant, and a uniform floor on an assembled
massive perturbation tube, isolating a massless-neutrino row and a
Higgs--Yukawa floor as its remaining gates.

\emph{YM V69--V70.}  The YM programme proved global heat-semigroup
differentiation, then performed its compulsory audit (reading the
present cycle and importing two proof-order disciplines), proved
stable dissipative Duhamel word contraction, defined a label--root
incidence graph, classified acyclic graphs with two nonlinear centres,
reduced them to one shared root and two star families closed by its
generating function, proved a strong aggregate card, and isolated the
first incidence cycle as its next obstruction.

\begin{theorem}[V35 companion-audit decision]
\label{thm:c2v35-companion-decision}
No SM V95 overlap, floor, or determinant estimate and no YM V69--V70
semigroup, contraction, or incidence estimate is imported.  Neither
bears on the Rayleigh chapter or on the gates of
Theorem~\ref{thm:c2v15-gates}.  The two companion readings of the
present cycle are accurate.
\end{theorem}

\begin{proof}
The SM results concern lattice overlap operators and their spectral
floors; the YM results concern Duhamel expansions of a lattice clock.
None is a Navier--Stokes statement.  The readings are checked against
V29 and the assembly discipline of the series.
\end{proof}

% =============================================================================
\section{The Rayleigh chapter, consolidated}
\label{sec:c2v35-consolidated}
% =============================================================================

\begin{theorem}[Rayleigh chapter, V31--V34]
\label{thm:c2v35-consolidated}
With the notation of V29--V34 (\(\widehat V=V/\varphi^{1/2}\), \(r=2\mathcal L/\varphi\),
\(\operatorname{Var}=\|L_-\widehat V\|_G^2-r^2\), \(\mathcal R=L_-\widehat V+r\widehat V\),
\(\widehat\jmath=j/\varphi\), \(\mathcal P\) the geometric push
\eqref{eq:c2v33-P}):
\begin{enumerate}[label=\textup{(\arabic*)},leftmargin=3.2em]
\item \emph{Identity} (rate-free).  \(\frac{\dd r}{\dd\sigma}=-2\operatorname{Var}+2\varphi^{1/2}\langle\mathcal R,\widehat N\rangle_G
      =-2\|L_-\widehat V\|_G^2-r\frac{\dd}{\dd\sigma}\log\varphi+\frac{2\mathcal P}\varphi\),
      and the push budget over any window:
      \(\int\frac{2\mathcal P}\varphi=\Delta r+2\int\operatorname{Var}+\int\frac r2(-\widehat\jmath)\).
\item \emph{Forced escape} (rate-free).  At a paraboloid-dissipative
      point the Dirichlet functional is not integrable in log-scale at
      zero; under Gaussian invisibility, on Dirichlet-active scales at
      level \(c\) below \(\lambda\), \(r\ge2c/\eta(\lambda)\) and the mean
      Hermite degree is at least \(4c/\eta(\lambda)-1\); along every ascent
      of \(r\) the nonlinear push equals the ascent plus twice the
      accumulated variance.
\item \emph{Compact dissipation} (rate-free).  Compact dissipation at a
      scale with \(R\le2\eta^{1/2}\) gives paraboloid dissipation at one
      centre at scale \(S^{1/2}\lambda\) with the explicit constant
      \eqref{eq:c2v32-epsprime}, hence a fixed Dirichlet log-integral
      below that scale at that centre, without coherence.
\item \emph{Uniform variance} (type I, homogeneous, modulo (E1)--(E9)).
      No rescaling limit is a Hermite-mode profile at any scale;
      \(\operatorname{Var}\ge v_0>0\) on the family at all scales; and the
      mean geometric push satisfies
      \(\int\frac{2\mathcal P}\varphi\dd\mathfrak m\ge2v_0+\frac14\int(-\widehat\jmath)_+\dd\mathfrak m-\frac{C_r}2\int(\widehat\jmath)_+\dd\mathfrak m\)
      for every invariant measure.
\end{enumerate}
\end{theorem}

\begin{proof}
(1) Theorem~\ref{thm:c2v31-identity}, Proposition~\ref{prop:c2v32-push},
Theorem~\ref{thm:c2v33-budget}.  (2) Proposition~\ref{prop:c2v31-nonintegrable},
Theorem~\ref{thm:c2v31-forced}, Lemma~\ref{lem:c2v31-degree}.
(3) Theorem~\ref{thm:c2v32-paraboloid}.  (4) Theorems~\ref{thm:c2v34-variance},
\ref{thm:c2v34-mean-push}.
\end{proof}

% =============================================================================
\section{The second cycle after thirty-five versions}
\label{sec:c2v35-summary}
% =============================================================================

\begin{theorem}[The second cycle by chapter]
\label{thm:c2v35-chapters}
Let \(u\) be a smooth mean-zero solution on \([0,T_*)\times\Torus^3\) with
\(T_*<\infty\).  The second cycle has established the following, each
chapter under the hypotheses indicated.
\begin{enumerate}[label=\textup{(K\arabic*)},leftmargin=3.8em]
\item \emph{Self-similar variables} (V1--V3; class \(\mathcal A_\nu(C)\)
      with the Morrey bound).  The profile dictionary, the two exact
      identities (weighted energy; Gaussian Dirichlet functional), the
      Lamb--Bernoulli form of the cross term, the bounds on the
      nonlinearity and the log-integral of the defect.
\item \emph{Dichotomy} (V4--V9; type I, homogeneous point, (E1)--(E9)).
      Semicontinuity of the defect; a rescaling family contains a
      backward self-similar element or is uniformly non-self-similar
      with the anti-alignment sign; by (E9) the latter holds at every
      homogeneous point, for the family and for the solution, with the
      three averaged signs.
\item \emph{Statistics} (V10--V13, V23; type I, any singular point).
      The rescaling flow on the compact family, recurrence, invariant
      and blow-up measures, exact mean identities and signs (strict for
      nontrivial measures), the discrete dichotomy, and the statistical
      dichotomy at lacunary points.
\item \emph{Moment ledger} (V16--V22; type I, homogeneous).  Centre
      mode, momentum, energy tensor, inertia, enstrophy, radial kinetic
      energy: exact identities with mean forms, signed constraints
      (negative mean stress, positive mean net stretching) and
      balances; the ledger is closed for the purposes of the series.
\item \emph{Rate-free chapter} (V14, V26--V29; no rate).  The second
      identity for every smooth solution; parabolic scales and
      compactness to suitable weak solutions; frequency escape and
      Gaussian invisibility; the Dirichlet ledger's lower bound at
      paraboloid-dissipative scales; the divergence theorem and its
      trichotomy; the spectral form of invisibility.
\item \emph{Compact dissipation} (V24, V32; no rate).  Compact
      dissipation is Gaussian activity at moving centres and gives
      paraboloid dissipation at a centre; the Gaussian net identity;
      moving-centre statistics under coherence.
\item \emph{Rayleigh chapter} (V31--V34).  Theorem~\ref{thm:c2v35-consolidated}.
\end{enumerate}
The open conditions are: for the type-I homogeneous case, the single
mean inequality of Theorem~\ref{thm:c2v15-gates}(G1), equivalently a
mean form of the pointwise sign, and the outward-flux-poor condition
of Firewall~\ref{fw:c2v34-AC2}; for lacunary points, statistical
nontriviality; for the rate-free case, compact dissipation and
coherence, and the size of the nonlinear push.  The full
Navier--Stokes stability/global-regularity theorem is not proved.
\end{theorem}

\begin{proof}
The consolidations of V5, V10, V15, V20, V25, V30, and
Theorem~\ref{thm:c2v35-consolidated}.
\end{proof}

\begin{assessment}[Direction for V36--V40]
\label{ass:c2v35-direction}
The cycle has produced exact, signed, scale-invariant balances for
every Gaussian moment of the profile and for its spectral position, on
compact families, and rate-free ledgers with forced divergences off
them.  What it has not produced is an inequality for a nonlinear
moment that is not a consequence of the profile equation tested
against a moment.  For V36--V40 the direction is to test the one
structural inequality the Navier--Stokes nonlinearity is known to
satisfy pointwise, the cancellation \(\langle(V\cdot\nabla)V,V\rangle=0\) in
the unweighted inner product together with its Gaussian-weighted
correction (the flux term), against the moment ledger: specifically,
whether the Gaussian correction of the cancellation, the flux
\(j=\frac1\nu\int(|V|^2+2P_V)(z\cdot V)G\), can be bounded in terms of the
Gaussian Dirichlet functional and the weighted energy by an inequality
sharper than \(|j|\le C_J\), for instance through the Hardy-type
inequality \(\int|z||V|^3G\le C(\nu)\|V\|_{L^\infty}\bigl(\|V\|_G^2+\|\nabla V\|_G^2\bigr)\)
and the pressure representation, and whether such an inequality,
combined with the mean identity \(\langle j\rangle=-\langle d+2\varphi\rangle\),
constrains the mean Dirichlet functional from above by the weighted
energy.  If it does, the pinching and the uniform concentration would
be in tension.  No global regularity claim is made.
\end{assessment}

% =============================================================================
\section{Integrated V35 conclusion and next acquisition order}
\label{sec:c2v35-conclusion}
% =============================================================================

\begin{theorem}[What V35 establishes]
\label{thm:c2v35-conclusion}
Relative to the first cycle and V1--V34: the companion-audit decision
(Theorem~\ref{thm:c2v35-companion-decision}), the consolidated Rayleigh
chapter (Theorem~\ref{thm:c2v35-consolidated}), and the summary by
chapter (Theorem~\ref{thm:c2v35-chapters}).  The full Navier--Stokes
stability/global-regularity theorem remains unproved.
\end{theorem}

\begin{proof}
The cited results.
\end{proof}

\begin{programme}[V36 acquisition order]
\label{prog:c2v36}
\begin{enumerate}[label=\textup{(AC\arabic*)},leftmargin=4.8em]
\item Flux versus Dirichlet functional.  Prove the sharpest available
      bound of the form \(|j|\le A\,\|V\|_\infty\,\mathcal L+B\,\|V\|_\infty\,\varphi\)
      with explicit \(A,B\) depending on \(\nu\) and the Riesz constant,
      using the Gaussian Hardy inequality and the pressure
      representation; then insert into \(\langle j\rangle=-\langle d+2\varphi\rangle\)
      to obtain \(\langle d+2\varphi\rangle\le AC_\infty\langle\mathcal L\rangle+BC_\infty\langle\varphi\rangle\)
      and determine whether this is a constraint (with \(8\mathcal L=d+2\varphi\)
      it reads \(8\langle\mathcal L\rangle\le AC_\infty\langle\mathcal L\rangle+BC_\infty\langle\varphi\rangle\),
      which is a constraint only if \(AC_\infty<8\)).
\item Record the size of \(C_\infty\) in scale-invariant units (the
      \(L^\infty\) rate constant) and whether the constants \(A,B\) can be
      made small enough; if not, record the obstruction precisely.
\item The next compulsory companion audit is V40.
\end{enumerate}
\end{programme}

\begin{assessment}[Position after V35]
\label{ass:c2v35-position}
The second cycle is consolidated by chapter through thirty-five
versions.  No full stability or global-regularity claim is made.
\end{assessment}

% =============================================================================
\section*{Version V35 (second cycle) append-only record}
\addcontentsline{toc}{section}{Version V35 (second cycle) append-only record}
% =============================================================================

The complete V34 source of the second cycle was retained before this
continuation.  Its SHA--256 digest was
\begin{center}\scriptsize\ttfamily
da85c70c0ff5749ba01b08feb3e66ef72f4a71e8536e621dfc106bd9726b8e78.
\end{center}
V35 performed the compulsory companion audit, consolidated the Rayleigh
chapter, and summarized the second cycle by chapter.  No full
regularity claim is made.

% =============================================================================
\section{V36: the flux against the Dirichlet functional, and the first
explicit constraint on a homogeneous blow-up}
\label{sec:c2v36-entry}
% =============================================================================

Items (AC1) and (AC2) of Programme~\ref{prog:c2v36} are carried out.  A
Gaussian Hardy inequality bounds the Gaussian second moment of the
profile by the weighted energy and the Dirichlet functional with
explicit constants; with the \(L^\infty\) rate for the energy part and
the \(L^3\) bound on the pressure for the pressure part, the
scale-invariant flux is bounded by an explicit multiple
\(\kappa=C_\infty\nu^{-1/2}\) of the positive terms of the first identity
plus a square-root term.  Inserted into the exact mean identity of the
flux, this yields, when \(\kappa<\frac12\), an explicit upper bound on
the mean of the weighted energy and the Gaussian dissipation against
every invariant measure at a homogeneous point, hence on the pinching
constant, and an explicit exclusion criterion when the concentration
lower bound on the mean dissipation exceeds it.  The criterion cannot
be evaluated because the concentration constant is not explicit and
the velocity constant of a hypothetical singularity is unknown; the
obstruction is recorded exactly.  No global regularity claim is made.
The next compulsory companion audit is V40.

% =============================================================================
\section{The Gaussian Hardy inequality and the flux bound}
\label{sec:c2v36-hardy}
% =============================================================================

Throughout, \(W\in\mathcal A_\nu(C)\) with the Morrey bound and the
uniform bounds \eqref{eq:c2v4-uniform}, \(V\) its profile, and in
scale-invariant units \(\|V(\sigma)\|_{L^\infty}\le C_\infty\),
\(\|V(\sigma)\|_{L^6(\R^3)}\le C_6C^{1/2}=:K_6\) for all \(\sigma\) (the
\(L^\infty\) rate and the class bound transported to the profile).  Put
\(X^2:=\int|z|^2|V|^2G\), \(\varphi=\|V\|_G^2\), \(d=4\nu\|\nabla V\|_G^2\),
\(\mathcal L=\frac18(d+2\varphi)\), and \(\kappa:=C_\infty\nu^{-1/2}\), a
dimensionless number.

\begin{lemma}[Gaussian Hardy inequality]
\label{lem:c2v36-hardy}
For every \(\sigma\),
\begin{equation}
 X^2\ \le\ 12\nu\varphi+16\nu^2\|\nabla V\|_G^2\ =\ 4\nu\,(\varphi+8\mathcal L)\ =\ 4\nu\,(d+3\varphi) .
 \label{eq:c2v36-hardy}
\end{equation}
\end{lemma}

\begin{proof}
\(|z|^2G=-2\nu\,z\cdot\nabla G\), so
\(X^2=-2\nu\int|V|^2z\cdot\nabla G=2\nu\int\operatorname{div}(|V|^2z)G=2\nu\int(3|V|^2+2V\cdot(z\cdot\nabla)V)G\le6\nu\varphi+4\nu X\|\nabla V\|_G\),
and \(4\nu X\|\nabla V\|_G\le\frac12X^2+8\nu^2\|\nabla V\|_G^2\); hence
\(X^2\le12\nu\varphi+16\nu^2\|\nabla V\|_G^2\).  With
\(16\nu^2\|\nabla V\|_G^2=4\nu d=32\nu\mathcal L-8\nu\varphi\) this is
\(X^2\le4\nu\varphi+32\nu\mathcal L=4\nu(\varphi+8\mathcal L)=4\nu(d+3\varphi)\).
\end{proof}

\begin{theorem}[Flux bound]
\label{thm:c2v36-flux-bound}
For every \(\sigma\),
\begin{equation}
 \boxed{\quad
 |j|\ \le\ \kappa\,(d+4\varphi)+K_P\,(d+3\varphi)^{1/2},
 \qquad
 K_P:=4\nu^{-1/2}\Bigl(\frac{4\pi\nu}3\Bigr)^{1/4}C_{\rm CZ}K_6^2 ,
 \quad}
 \label{eq:c2v36-flux-bound}
\end{equation}
where \(C_{\rm CZ}\) is the \(L^{3}\to L^{3}\) norm bound of the Riesz
representation \(V\otimes V\mapsto P_V\) in the form \(\|P_V\|_3\le C_{\rm CZ}\|V\|_6^2\).
\end{theorem}

\begin{proof}
\(j=\frac1\nu\int|V|^2(z\cdot V)G+\frac2\nu\int P_V(z\cdot V)G\).  Energy part:
\(|V|^2|z\cdot V|\le C_\infty|z||V|^2\) and
\(\int|z||V|^2G\le X\varphi^{1/2}\le2\nu^{1/2}(\varphi+8\mathcal L)^{1/2}\varphi^{1/2}\le\nu^{1/2}(2\varphi+8\mathcal L)=\nu^{1/2}(d+4\varphi)\)
by \eqref{eq:c2v36-hardy} and \(2ab\le a^2+b^2\); so the energy part is
at most \(\nu^{-1}C_\infty\nu^{1/2}(d+4\varphi)=\kappa(d+4\varphi)\).  Pressure
part: \(\int(z\cdot V)G=0\) (zero mass of the weighted radial momentum,
Proposition V89-moments of the first cycle), so the pressure may be
taken mean-free; by H\"older,
\(|\int P_V(z\cdot V)G|\le\|P_V\|_{L^3(\R^3)}\|(z\cdot V)G\|_{L^{3/2}(\R^3)}\),
and \(\|(z\cdot V)G\|_{3/2}^{3/2}=\int|z\cdot V|^{3/2}G^{3/4}\cdot G^{3/4}\le(\int|z\cdot V|^2G)^{3/4}(\int G^3)^{1/4}\),
so \(\|(z\cdot V)G\|_{3/2}\le X(\int G^3)^{1/6}=X(4\pi\nu/3)^{1/4}\).  With
\(\|P_V\|_3\le C_{\rm CZ}K_6^2\) and \(X\le2\nu^{1/2}(d+3\varphi)^{1/2}\), the
pressure part is at most
\(\frac2\nu C_{\rm CZ}K_6^2(4\pi\nu/3)^{1/4}\cdot2\nu^{1/2}(d+3\varphi)^{1/2}=K_P(d+3\varphi)^{1/2}\).
\end{proof}

% =============================================================================
\section{The first explicit constraint}
\label{sec:c2v36-constraint}
% =============================================================================

\begin{theorem}[Mean constraint for small velocity constant]
\label{thm:c2v36-constraint}
Let \(x_0\) be a homogeneous singular point of a type-I solution with
velocity constant \(\kappa=C_\infty\nu^{-1/2}<\frac12\), and let
\(\mathfrak m\) be any invariant probability measure on \(\mathcal F(x_0)\).
Put \(m_\kappa:=\min\{1-\kappa,\ (2-4\kappa)/3\}>0\).  Then, modulo
(E1)--(E9),
\begin{equation}
 \boxed{\quad
 \langle d+3\varphi\rangle_{\mathfrak m}\ \le\ \Bigl(\frac{K_P}{m_\kappa}\Bigr)^2 ,
 \qquad\text{hence}\qquad
 c_\Phi\le\langle\varphi\rangle\le\frac{K_P^2}{3m_\kappa^2},
 \qquad
 \langle d\rangle\le\frac{K_P^2}{m_\kappa^2} .
 \quad}
 \label{eq:c2v36-constraint}
\end{equation}
\end{theorem}

\begin{proof}
By \eqref{eq:c2v12-mean1}, \(\langle j\rangle=-\langle d+2\varphi\rangle\), so
\(\langle d+2\varphi\rangle=-\langle j\rangle\le\langle|j|\rangle\le\kappa\langle d+4\varphi\rangle+K_P\langle(d+3\varphi)^{1/2}\rangle\)
by \eqref{eq:c2v36-flux-bound}.  Rearranging,
\((1-\kappa)\langle d\rangle+(2-4\kappa)\langle\varphi\rangle\le K_P\langle(d+3\varphi)^{1/2}\rangle\le K_P\langle d+3\varphi\rangle^{1/2}\)
(Jensen).  For \(\kappa<\frac12\) both coefficients are positive and
\((1-\kappa)d+(2-4\kappa)\varphi\ge m_\kappa(d+3\varphi)\), so
\(m_\kappa\langle d+3\varphi\rangle\le K_P\langle d+3\varphi\rangle^{1/2}\), which gives
the first bound; the others follow from \(\varphi\ge c_\Phi\) on the family
and \(d,\varphi\ge0\).
\end{proof}

\begin{corollary}[Explicit exclusion criterion]
\label{cor:c2v36-exclusion}
Let \(x_0\) be homogeneous with the window parameters \(R,S,\eta\) and the
dissipation constant \(\varepsilon\) of the first cycle, and
\(c_2=2\nu e^{-R'^2/(4\nu)}\varepsilon/S^{1/2}\), \(R'=R\eta^{-1/2}\), the
constant of the averaged inward flux (Theorem V86-inward).  Then for
every invariant measure \(\langle d\rangle\ge2c_2/\log(S/\eta)=:d_{\rm low}\).
Consequently, if \(\kappa<\frac12\) and
\begin{equation}
 \frac{2c_2}{\log(S/\eta)}\ >\ \frac{K_P^2}{m_\kappa^2},
 \label{eq:c2v36-criterion}
\end{equation}
then no such homogeneous singular point exists.
\end{corollary}

\begin{proof}
The lower bound: the proof of Theorem V86-inward of the first cycle
shows that on every window \([\eta^{1/2}\Lambda,S^{1/2}\Lambda]\) of scales,
\(\int D_W(\lambda)\dd\lambda=\int d\,\dd\lambda/\lambda\ge c_2\); the window has
log-length \(\frac12\log(S/\eta)\), so the log-average of \(d\) over it is
at least \(2c_2/\log(S/\eta)\), and by invariance the same holds for the
\(\mathfrak m\)-mean.  The criterion is the comparison with
\eqref{eq:c2v36-constraint}.
\end{proof}

\begin{firewall}[Item (AC2): why the criterion cannot be evaluated]
\label{fw:c2v36-AC2}
Three quantities in \eqref{eq:c2v36-criterion} are not explicit.  The
dissipation constant \(\varepsilon=\varepsilon(c)\) of a homogeneous point
is produced by the compactness argument of the first cycle
(Corollary V81-quantitative) from the homogeneity constant \(c\) and is
not computable; the window parameters \(R,S,\eta\) come from the same
argument.  The velocity constant \(\kappa\) of a hypothetical
singularity is bounded below by the \(\varepsilon\)-regularity threshold
of (E5) in units of \(\nu^{1/2}\) (a singular point has
\(\limsup|u|(T_*-t)^{1/2}\ge\varepsilon_{\rm reg}\nu^{1/2}\)), but has no
upper bound, and the constraint is void for \(\kappa\ge\frac12\).  The
constant \(K_6\) is the class \(L^6\) bound \(C_6C^{1/2}\) with
\(C=C_I^2\nu^{3/2}\), explicit in the type-I constant.  Hence the
criterion is an explicit inequality between one non-explicit lower
bound and one explicit upper bound, valid on the range
\(\kappa\in[\varepsilon_{\rm reg},\frac12)\) of velocity constants; it is
the first explicit constraint of the series on a hypothetical
homogeneous blow-up, and the series cannot decide it.
\end{firewall}

\begin{assessment}[What the constraint says]
\label{ass:c2v36-meaning}
The exact mean identity of the flux, the inward flux of the first
cycle, says that the nonlinearity must pump Gaussian energy inward at
the rate of the scaling damping plus the dissipation.  The flux bound
says that the nonlinearity cannot pump more than \(\kappa\) times the
damping-plus-dissipation plus a pressure term of square-root size.  For
a blow-up whose velocity constant is less than half the viscous unit,
the pump cannot keep up with a large dissipation, so the mean Gaussian
dissipation and energy are bounded above by an explicit constant built
from the pressure bound.  The concentration of the first cycle bounds
the mean dissipation below.  Whether the two bounds are compatible is
a question about constants that the compactness arguments of the first
cycle do not make explicit.  No global regularity claim is made.
\end{assessment}

% =============================================================================
\section{Integrated V36 conclusion and next acquisition order}
\label{sec:c2v36-conclusion}
% =============================================================================

\begin{theorem}[What V36 proves]
\label{thm:c2v36-conclusion}
Relative to the first cycle and V1--V35, modulo (E1)--(E9): the Gaussian
Hardy inequality \eqref{eq:c2v36-hardy}, the flux bound
\eqref{eq:c2v36-flux-bound}, the mean constraint
\eqref{eq:c2v36-constraint} for \(\kappa<\frac12\), and the exclusion
criterion \eqref{eq:c2v36-criterion}.  The full Navier--Stokes
stability/global-regularity theorem remains unproved.
\end{theorem}

\begin{proof}
The cited results.
\end{proof}

\begin{programme}[V37 acquisition order]
\label{prog:c2v37}
\begin{enumerate}[label=\textup{(AC\arabic*)},leftmargin=4.8em]
\item Improve the energy part.  The bound \(\kappa(d+4\varphi)\) used only
      \(|V|\le C_\infty\); integrate \(\int|V|^2(z\cdot V)G\) by parts
      (\(\int|V|^2(z\cdot V)G=-2\nu\int|V|^2V\cdot\nabla G=2\nu\int\psi\)-type
      identities of V4 give \(\int|V|^2(z\cdot V)G=2\nu\int V\cdot\nabla|V|^2\,G\))
      and bound \(\int|V||\nabla V||V|G\le C_\infty\|V\|_G\|\nabla V\|_G\) to
      obtain a bound of the form \(\kappa'(\varphi\,d)^{1/2}\), which is
      smaller than \(\kappa(d+4\varphi)\) when \(d\) and \(\varphi\) are
      unbalanced; determine the optimal combination and the resulting
      threshold on \(\kappa\).
\item Remove the pressure term's \(L^3\) dependence.  Bound
      \(\int P_V(z\cdot V)G\) through the local pressure representation
      and the Morrey bound instead of the global \(L^3\) norm, to obtain
      a term proportional to \(C_\infty\) times \(d+\varphi\) as well, and
      recompute the threshold.
\item The next compulsory companion audit is V40.
\end{enumerate}
\end{programme}

\begin{assessment}[Position after V36]
\label{ass:c2v36-position}
The series has its first explicit constraint on a hypothetical
homogeneous type-I blow-up, valid for velocity constants below half
the viscous unit, and an exclusion criterion it cannot evaluate.  No
full stability or global-regularity claim is made.
\end{assessment}

% =============================================================================
\section*{Version V36 (second cycle) append-only record}
\addcontentsline{toc}{section}{Version V36 (second cycle) append-only record}
% =============================================================================

The complete V35 source of the second cycle was retained before this
continuation.  Its SHA--256 digest was
\begin{center}\scriptsize\ttfamily
52c5627d8c2b8fe58e1c02602a18f5aba6cbd248fe326b7dc6f390122c711fdd.
\end{center}
V36 proved the Gaussian Hardy inequality, the flux bound, the mean
constraint for small velocity constant, and the explicit exclusion
criterion.  No full regularity claim is made.

% =============================================================================
\section{V37: the sharpened flux bound and the pressure split}
\label{sec:c2v37-entry}
% =============================================================================

Items (AC1) and (AC2) of Programme~\ref{prog:c2v37} are carried out.  The
energy part of the flux is rewritten by one integration by parts as a
Gaussian correlation of the velocity with the gradient of its energy
density, which the \(L^\infty\) rate bounds by the geometric mean of the
weighted energy and the Gaussian dissipation; this halves the previous
coefficient and, in the mean constraint, replaces the threshold
\(\kappa<\frac12\) on the velocity constant by \(\kappa<\sqrt2\), with an
explicit quadratic form whose smallest eigenvalue governs the bound.
The pressure part is split into a local part, bounded through the
weighted energy at the price of an exponential factor in the cutoff
radius, and a far part, harmonic on the cutoff ball and bounded through
the global pressure norm at the price of an inverse power of the
radius; the resulting fully linear flux bound holds only below an
exponentially small threshold on the velocity constant, which is
recorded as the obstruction to removing the square-root term.  No
global regularity claim is made.  The next compulsory companion audit
is V40.

% =============================================================================
\section{The sharpened energy part}
\label{sec:c2v37-energy}
% =============================================================================

Notation as in V36: \(\varphi=\|V\|_G^2\), \(d=4\nu\|\nabla V\|_G^2\),
\(\kappa=C_\infty\nu^{-1/2}\), \(K_P\) of \eqref{eq:c2v36-flux-bound}.

\begin{proposition}[Energy part of the flux by the geometric mean]
\label{prop:c2v37-energy}
For every \(\sigma\),
\begin{equation}
 \Bigl|\frac1\nu\int|V|^2(z\cdot V)\,G\Bigr|\ \le\ 2\kappa\,(\varphi\,d)^{1/2}\ \le\ \frac\kappa2\,(d+4\varphi),
 \label{eq:c2v37-energy}
\end{equation}
and consequently
\begin{equation}
 |j|\ \le\ 2\kappa\,(\varphi d)^{1/2}+K_P\,(d+3\varphi)^{1/2} .
 \label{eq:c2v37-flux-bound}
\end{equation}
\end{proposition}

\begin{proof}
\((z\cdot V)G=-2\nu\,V\cdot\nabla G\) and \(\operatorname{div}V=0\) give
\(\int|V|^2(z\cdot V)G=-2\nu\int|V|^2V\cdot\nabla G=2\nu\int\operatorname{div}(|V|^2V)G=2\nu\int V\cdot\nabla|V|^2\,G\),
so \(|\int|V|^2(z\cdot V)G|\le4\nu\int|V|^2|\nabla V|G\le4\nu C_\infty\|V\|_G\|\nabla V\|_G
=4\nu C_\infty\varphi^{1/2}(d/(4\nu))^{1/2}=2\nu^{1/2}C_\infty(\varphi d)^{1/2}\);
divide by \(\nu\).  The second inequality is \(2(\varphi d)^{1/2}=((4\varphi)d)^{1/2}\le\frac12(d+4\varphi)\).
The flux bound combines with the pressure part of
Theorem~\ref{thm:c2v36-flux-bound}.
\end{proof}

\begin{theorem}[Mean constraint for velocity constant below \(\sqrt2\)]
\label{thm:c2v37-constraint}
Let \(x_0\) be a homogeneous singular point of a type-I solution with
\(\kappa<\sqrt2\), \(\mathfrak m\) any invariant probability measure on
\(\mathcal F(x_0)\), and
\begin{equation}
 m'_\kappa:=\frac{3-(1+4\kappa^2)^{1/2}}2\ >0 ,
 \label{eq:c2v37-mprime}
\end{equation}
the smallest eigenvalue of the quadratic form \(a^2-2\kappa ab+2b^2\).
Then, modulo (E1)--(E9),
\begin{equation}
 \boxed{\quad
 \langle d\rangle_{\mathfrak m}+\langle\varphi\rangle_{\mathfrak m}\ \le\ \frac{3K_P^2}{m_\kappa'^2},
 \qquad\text{hence}\qquad
 c_\Phi\le\frac{3K_P^2}{m_\kappa'^2},\qquad
 \langle d\rangle\le\frac{3K_P^2}{m_\kappa'^2} .
 \quad}
 \label{eq:c2v37-constraint}
\end{equation}
The exclusion criterion of Corollary~\ref{cor:c2v36-exclusion} holds with
\(K_P^2/m_\kappa^2\) replaced by \(3K_P^2/m_\kappa'^2\) and \(\kappa<\frac12\)
by \(\kappa<\sqrt2\).
\end{theorem}

\begin{proof}
With \(a:=\langle d\rangle^{1/2}\), \(b:=\langle\varphi\rangle^{1/2}\), the mean
identity \(\langle d+2\varphi\rangle=-\langle j\rangle\le\langle|j|\rangle\) and
\eqref{eq:c2v37-flux-bound} with Cauchy--Schwarz and Jensen
(\(\langle(\varphi d)^{1/2}\rangle\le ab\), \(\langle(d+3\varphi)^{1/2}\rangle\le(a^2+3b^2)^{1/2}\))
give \(a^2+2b^2-2\kappa ab\le K_P(a^2+3b^2)^{1/2}\le\sqrt3K_P(a^2+b^2)^{1/2}\).
The form \(a^2-2\kappa ab+2b^2\) has matrix eigenvalues
\(\frac{3\pm(1+4\kappa^2)^{1/2}}2\), positive definite exactly when
\(\kappa^2<2\), with minimum \(m'_\kappa\); hence
\(m'_\kappa(a^2+b^2)\le\sqrt3K_P(a^2+b^2)^{1/2}\), which is
\eqref{eq:c2v37-constraint}.
\end{proof}

% =============================================================================
\section{The pressure split}
\label{sec:c2v37-pressure}
% =============================================================================

\begin{proposition}[Local and far parts of the pressure term]
\label{prop:c2v37-pressure}
Let \(\Lambda>0\), \(\chi_\Lambda\) a smooth cutoff equal to \(1\) on
\(B_{2\Lambda}\) and \(0\) outside \(B_{4\Lambda}\), and
\(P_{\rm loc}:=\partial_i\partial_j(-\Delta)^{-1}(\chi_\Lambda V_iV_j)\),
\(P_{\rm far}:=P_V-P_{\rm loc}\), harmonic on \(B_{2\Lambda}\).  Then, with
absolute constants \(C_1,C_2,C_3\) depending only on the
Calder\'on--Zygmund constants and the cutoff,
\begin{align}
 \Bigl|\frac2\nu\int P_{\rm loc}(z\cdot V)G\Bigr|
 &\le C_1\,\kappa\,e^{2\Lambda^2/\nu}\,(d+4\varphi),
 \label{eq:c2v37-loc}\\
 \Bigl|\frac2\nu\int_{B_\Lambda}(P_{\rm far}-\bar P_{\rm far})(z\cdot V)G\Bigr|
 &\le C_2\,\frac{K_6^2}{\Lambda}\,\nu^{-1/2}\,(d+4\varphi),
 \label{eq:c2v37-far-in}\\
 \Bigl|\frac2\nu\int_{\R^3\setminus B_\Lambda}(P_{\rm far}-\bar P_{\rm far})(z\cdot V)G\Bigr|
 &\le C_3\,K_P\,e^{-\Lambda^2/(8\nu)}\,(d+3\varphi)^{1/2},
 \label{eq:c2v37-far-out}
\end{align}
where \(\bar P_{\rm far}\) is the mean of \(P_{\rm far}\) over \(B_\Lambda\)
(a constant, which may be subtracted since \(\int(z\cdot V)G=0\)).
\end{proposition}

\begin{proof}
\eqref{eq:c2v37-loc}: \(\|P_{\rm loc}\|_{L^2(\R^3)}\le C\|\chi_\Lambda V\otimes V\|_2\le CC_\infty\|V\|_{L^2(B_{4\Lambda})}\le CC_\infty e^{2\Lambda^2/\nu}\varphi^{1/2}\)
since \(G\ge e^{-16\Lambda^2/(4\nu)}=e^{-4\Lambda^2/\nu}\) on \(B_{4\Lambda}\)
(so \(\|V\|_{L^2(B_{4\Lambda})}^2\le e^{4\Lambda^2/\nu}\varphi\)); and
\(|\int P_{\rm loc}(z\cdot V)G|\le\|P_{\rm loc}\|_2\|(z\cdot V)G\|_2\le\|P_{\rm loc}\|_2X\)
with \(X\le2\nu^{1/2}(d+3\varphi)^{1/2}\) (Lemma~\ref{lem:c2v36-hardy}); then
\(\frac2\nu CC_\infty e^{2\Lambda^2/\nu}\varphi^{1/2}\cdot2\nu^{1/2}(d+3\varphi)^{1/2}\le C_1\kappa e^{2\Lambda^2/\nu}(d+4\varphi)\)
by \(2\varphi^{1/2}(d+3\varphi)^{1/2}\le d+4\varphi\).
\eqref{eq:c2v37-far-in}: \(P_{\rm far}\) is harmonic on \(B_{2\Lambda}\), so
\(\sup_{B_\Lambda}|P_{\rm far}-\bar P_{\rm far}|\le C\Lambda\sup_{B_\Lambda}|\nabla P_{\rm far}|\le C\Lambda\cdot\Lambda^{-4}\|P_{\rm far}-c\|_{L^1(B_{2\Lambda})}\le C\Lambda^{-1}\|P_{\rm far}-c\|_{L^3(B_{2\Lambda})}\)
for any constant \(c\), by the mean-value property for \(\nabla P_{\rm far}\)
and H\"older; with \(c=0\), \(\|P_{\rm far}\|_{L^3}\le\|P_V\|_3+\|P_{\rm loc}\|_3\le C K_6^2\)
(Calder\'on--Zygmund in \(L^{3}\) for both).  Then
\(|\int_{B_\Lambda}(P_{\rm far}-\bar P_{\rm far})(z\cdot V)G|\le CK_6^2\Lambda^{-1}\int|z||V|G\le CK_6^2\Lambda^{-1}X\varphi^{1/2}\),
and \(\frac2\nu CK_6^2\Lambda^{-1}\cdot2\nu^{1/2}(d+3\varphi)^{1/2}\varphi^{1/2}\le C_2K_6^2\Lambda^{-1}\nu^{-1/2}(d+4\varphi)\).
\eqref{eq:c2v37-far-out}: H\"older as in Theorem~\ref{thm:c2v36-flux-bound}
with the \(L^{3/2}\) norm of \((z\cdot V)G\) restricted to \(|z|>\Lambda\),
where \((\int_{|z|>\Lambda}G^3)^{1/4}\le C e^{-\Lambda^2/(8\nu)}\cdot(\int G^3)^{1/4}\)
(the factor \(e^{-3\Lambda^2/(4\nu)\cdot\frac14}\cdot\)polynomial is absorbed).
\end{proof}

\begin{corollary}[A fully linear flux bound below an exponentially small threshold]
\label{cor:c2v37-linear}
For every \(\Lambda>0\) and \(\sigma\),
\begin{equation}
 |j|\ \le\ \Bigl(\frac\kappa2+C_1\kappa e^{2\Lambda^2/\nu}+C_2\frac{K_6^2}{\Lambda\nu^{1/2}}\Bigr)(d+4\varphi)+C_3K_Pe^{-\Lambda^2/(8\nu)}(d+3\varphi)^{1/2}.
 \label{eq:c2v37-linear}
\end{equation}
Choosing \(\Lambda=\Lambda_*:=4C_2K_6^2\nu^{-1/2}\) makes the third
coefficient \(\frac14\); if then
\(\kappa\le\kappa_0:=\frac18\,e^{-2\Lambda_*^2/\nu}/C_1\), the total coefficient
of \(d+4\varphi\) is at most \(\frac12+\frac18+\frac14<1\), and the mean
identity gives \(\langle d+2\varphi\rangle\le\frac78\langle d+4\varphi\rangle+C_3K_Pe^{-\Lambda_*^2/(8\nu)}\langle d+3\varphi\rangle^{1/2}\),
hence \(\frac18\langle d\rangle\le\frac32\langle\varphi\rangle+\ldots\): a bound of
the mean dissipation by the mean weighted energy up to an exponentially
small square-root term.
\end{corollary}

\begin{proof}
Sum \eqref{eq:c2v37-energy}, \eqref{eq:c2v37-loc}--\eqref{eq:c2v37-far-out}
and insert the choices.
\end{proof}

\begin{firewall}[Item (AC2): the obstruction to a linear bound]
\label{fw:c2v37-AC2}
The local pressure is controlled by the weighted energy only at the
price \(e^{2\Lambda^2/\nu}\), because the Gaussian weight is exponentially
small on the ball where the local pressure is generated, while the
far pressure is controlled by the global \(L^3\) bound at the price
\(\Lambda^{-1}\).  The two prices force \(\Lambda_*\) of order \(K_6^2\nu^{-1/2}\)
and hence a threshold \(\kappa_0\) of order \(\exp(-cK_6^4/\nu^2)\), far
below the \(\varepsilon\)-regularity floor \(\varepsilon_{\rm reg}\) of the
velocity constant for any \(K_6\) of order one.  So the linear bound
\eqref{eq:c2v37-linear} with a total coefficient below one is available
only on a range of velocity constants that is empty for singular
points, and the square-root term of \eqref{eq:c2v37-flux-bound} cannot
be removed by this method.  The usable constraint remains
Theorem~\ref{thm:c2v37-constraint}, valid for \(\kappa<\sqrt2\), with the
square-root pressure term.
\end{firewall}

\begin{assessment}[The state of the constant-chasing gate]
\label{ass:c2v37-gate}
The gate now reads: for a homogeneous type-I singular point with
velocity constant \(\kappa<\sqrt2\), the mean Gaussian dissipation and
energy are bounded above by \(3K_P^2/m_\kappa'^2\), an explicit function of
\(\kappa\), \(\nu\), the Calder\'on--Zygmund constant, and the class
\(L^6\) bound \(K_6=C_6C_I\nu^{3/4}\); and the mean Gaussian dissipation
is bounded below by \(2c_2/\log(S/\eta)\), a non-explicit function of
the homogeneity constant through the compactness of the first cycle.
The range \(\kappa\in[\varepsilon_{\rm reg},\sqrt2)\) is nonempty.  The gate
is decided if either the lower bound is made explicit and large, or
the upper bound is made small; the series has neither.  No global
regularity claim is made.
\end{assessment}

% =============================================================================
\section{Integrated V37 conclusion and next acquisition order}
\label{sec:c2v37-conclusion}
% =============================================================================

\begin{theorem}[What V37 proves]
\label{thm:c2v37-conclusion}
Relative to the first cycle and V1--V36, modulo (E1)--(E9): the
sharpened energy part \eqref{eq:c2v37-energy} and flux bound
\eqref{eq:c2v37-flux-bound}, the mean constraint
\eqref{eq:c2v37-constraint} for \(\kappa<\sqrt2\), the pressure split
\eqref{eq:c2v37-loc}--\eqref{eq:c2v37-far-out}, and the linear bound
\eqref{eq:c2v37-linear} with its threshold.  The full Navier--Stokes
stability/global-regularity theorem remains unproved.
\end{theorem}

\begin{proof}
The cited results.
\end{proof}

\begin{programme}[V38 acquisition order]
\label{prog:c2v38}
\begin{enumerate}[label=\textup{(AC\arabic*)},leftmargin=4.8em]
\item The lower bound side.  Determine whether the mean Gaussian
      dissipation at a homogeneous point can be bounded below
      explicitly in terms of the homogeneity constant \(c\) and the
      type-I constants, bypassing the compactness of the first cycle:
      the per-scale equivalence (terminal density \(\Rightarrow\)
      dissipation) was proved by contradiction; examine whether the
      rate-free local energy bound of the first cycle
      (Theorem V66-local-energy), which is explicit, gives an explicit
      lower bound on the Leray energy spent in a parabolic window from
      a lower bound on the terminal density, and hence on \(\langle d\rangle\).
\item If (AC1) succeeds, evaluate the gate: compare the explicit lower
      bound with \(3K_P^2/m_\kappa'^2\) as a function of \(\kappa\) and \(c\),
      and record the region of parameters excluded.
\item The next compulsory companion audit is V40.
\end{enumerate}
\end{programme}

\begin{assessment}[Position after V37]
\label{ass:c2v37-position}
The constant-chasing gate has an explicit upper side valid for
velocity constants below \(\sqrt2\); its lower side is not explicit.
No full stability or global-regularity claim is made.
\end{assessment}

% =============================================================================
\section*{Version V37 (second cycle) append-only record}
\addcontentsline{toc}{section}{Version V37 (second cycle) append-only record}
% =============================================================================

The complete V36 source of the second cycle was retained before this
continuation.  Its SHA--256 digest was
\begin{center}\scriptsize\ttfamily
61b888ce4b499bbb6075f632431406622f461d64b8d9afaae6aeb84765c215be.
\end{center}
V37 sharpened the flux bound, extended the mean constraint to velocity
constants below \(\sqrt2\), and recorded the pressure split with its
exponentially small threshold.  No full regularity claim is made.

% =============================================================================
\section{V38: the lower side of the constant-chasing gate and its obstruction}
\label{sec:c2v38-entry}
% =============================================================================

Items (AC1) and (AC2) of Programme~\ref{prog:c2v38} are carried out to
the extent possible, and the obstruction is recorded exactly.  The
explicit rate-free local energy bound of the first cycle is shown to
give no lower bound on the Leray energy spent in a parabolic window
from a lower bound on the terminal density, because its second term is
of the order of the quarter power of the scale while the terminal
density is of the order of the scale; the first cycle's per-scale
implication, proved by contradiction through compactness, is the only
route to a lower bound on the dissipation, and an explicit version
would require a quantitative form of the backward uniqueness and
unique continuation import.  The two sides of the gate are written out
in fully explicit form where possible, and it is recorded which
constants are explicit in the type-I data and which are not.  No global
regularity claim is made.  The next compulsory companion audit is V40.

% =============================================================================
\section{The rate-free bound does not give the lower side}
\label{sec:c2v38-rate-free}
% =============================================================================

\begin{proposition}[The explicit local energy bound is too weak at the parabolic order]
\label{prop:c2v38-too-weak}
Let \(x_0\) be homogeneous with constant \(c\), so that
\(\int_{B_\rho(x_0)}|u(T_*)|^2\ge c\rho\) for small \(\rho\), and let
\(\epsilon_\rho:=\epsilon(T_*;2\rho^2)=\nu\int_{T_*-2\rho^2}^{T_*}\|\Lambda u\|_2^2\dd t\)
be the Leray energy spent in the parabolic window.  The bound of
Theorem V66-local-energy of the first cycle, at the terminal time (by
the strong \(L^2\) continuity at \(T_*\) under type I), reads
\begin{equation}
 c\rho\ \le\ C_0\bigl(\epsilon_\rho+\rho^{-1/2}E_*^{3/4}\epsilon_\rho^{3/4}\bigr).
 \label{eq:c2v38-V66}
\end{equation}
Writing \(\epsilon_\rho=\theta_\rho\,\rho\) (under the type-I rate
\(\theta_\rho\le2\sqrt2\,C_I^2\nu^{5/2}\)), \eqref{eq:c2v38-V66} is
\(c\le C_0\theta_\rho+C_0E_*^{3/4}\theta_\rho^{3/4}\rho^{-3/4}\), which holds for
every \(\theta_\rho>0\) once \(\rho\) is small enough: the inequality gives
no lower bound on \(\theta_\rho\), hence none on the Leray share of the
parabolic window, hence none on the mean Gaussian dissipation.
\end{proposition}

\begin{proof}
Substitution; the second term \(C_0E_*^{3/4}\theta^{3/4}\rho^{-3/4}\) tends
to infinity as \(\rho\to0\) for every fixed \(\theta>0\), and dominates
\(c\) for \(\rho\le(C_0E_*^{3/4}\theta^{3/4}/c)^{4/3}\).
\end{proof}

\begin{theorem}[The lower side needs quantitative unique continuation]
\label{thm:c2v38-obstruction}
The first cycle's per-scale implication ``terminal density at least
\(c\) at radius \(R'\) forces dissipation at least \(\varepsilon(c,R')\) in a
window'' (Corollary V81-quantitative) is proved by contradiction: a
sequence with density bounded below and dissipation tending to zero
would have a local limit \(U'\) with \(\nabla U'\equiv0\), hence \(U'=0\),
whose trace is the weak limit of the rescaled terminal data, which
carry energy at least \(c\) on \(B_{R'}\); the contradiction uses the
identification of the trace as the weak limit (V48 of the first cycle)
and the vanishing of the trace of the zero solution.  An explicit
\(\varepsilon(c,R')\) would require an explicit statement of the form:
\emph{if the dissipation of a rescaled solution on \(B_R\times(-S,-\eta)\)
is at most \(\varepsilon\), then the energy of its terminal datum on
\(B_{R'}\) is at most \(\omega(\varepsilon)\) with \(\omega\) explicit and
\(\omega(\varepsilon)\to0\)}; this is a quantitative unique-continuation
statement across the time interval \((-\eta,0)\) from a small-dissipation
region, which (E6) provides only qualitatively.  The series does not
have it.
\end{theorem}

\begin{proof}
Inspection of the proofs of Theorem V53-per-scale and Corollary
V81-quantitative of the first cycle, which are the two places where
the lower bound is produced; both are compactness arguments whose
constants are defined by the failure of the contrary.
\end{proof}

% =============================================================================
\section{The gate in explicit form}
\label{sec:c2v38-explicit}
% =============================================================================

\begin{proposition}[Explicit constants of the upper side]
\label{prop:c2v38-upper}
With \(K_6=C_6C_I\nu^{3/4}\) (the class \(L^6\) bound in scale-invariant
units), the upper side of Theorem~\ref{thm:c2v37-constraint} is
\begin{equation}
 \langle d\rangle+\langle\varphi\rangle\ \le\ \frac{3K_P^2}{m_\kappa'^2}
 =\frac{48\,C_{\rm CZ}^2C_6^4C_I^4\,\nu^{5/2}\,(4\pi/3)^{1/2}}{m_\kappa'^2},
 \qquad
 m'_\kappa=\frac{3-(1+4\kappa^2)^{1/2}}2 ,
 \label{eq:c2v38-upper}
\end{equation}
explicit in the type-I constant \(C_I\), the velocity constant
\(\kappa=C_\infty\nu^{-1/2}\), the viscosity, and the absolute constants
\(C_6\) (Sobolev) and \(C_{\rm CZ}\) (Riesz).  The velocity constant is
itself bounded in terms of \(C_I\) by the first cycle's \(L^\infty\)
upgrade (Theorem V46-Linfty-upgrade: \(C_\infty=2C_{\rm CKN}/c_0\) with
\(c_0=c_0(C_I,\nu,\varepsilon_2,C_6,C_{\rm CZ})\)), so the whole upper side is a
function of \(C_I\) and \(\nu\) and the absolute constants, on the range
where \(\kappa<\sqrt2\).
\end{proposition}

\begin{proof}
\(K_P^2=16\nu^{-1}(4\pi\nu/3)^{1/2}C_{\rm CZ}^2K_6^4\) and \(K_6^4=C_6^4C_I^4\nu^3\);
multiply by \(3\).
\end{proof}

\begin{proposition}[The lower side and its non-explicit constants]
\label{prop:c2v38-lower}
The lower side is \(\langle d\rangle\ge2c_2/\log(S/\eta)\) with
\(c_2=2\nu e^{-R^2/(4\nu\eta)}\varepsilon/S^{1/2}\), where \((R,S,\eta,\varepsilon)\)
are the window parameters and dissipation constant of the per-scale
implication at the homogeneity constant \(c\) and radius \(R'=1\).  None
of \(R,S,\eta,\varepsilon\) is explicit (Theorem~\ref{thm:c2v38-obstruction});
the only explicit information is \(\varepsilon\le2C_I^2\nu^{3/2}(S^{1/2}-\eta^{1/2})\)
(the type-I upper bound on the window dissipation) and
\(c\le C_M\) (the Morrey bound).
\end{proposition}

\begin{proof}
Corollary~\ref{cor:c2v36-exclusion} and the first cycle's constants.
\end{proof}

\begin{assessment}[Item (AC2): the gate cannot be evaluated]
\label{ass:c2v38-gate}
The gate compares an explicit upper bound, a function of \(C_I\) and
\(\nu\), with a lower bound whose constants are defined by compactness.
It is decided in favour of exclusion only if the compactness constants
happen to be large, which nothing determines.  The series records the
gate as explicit on one side, non-explicit on the other, and turns to
the question of whether the lower side can be obtained by a different
route: not from the terminal density, but directly from the uniform
concentration in the form of the first cycle's \emph{quantitative}
count of homogeneous points (Theorem V81-count), whose constant
\(2C_I^2\nu^{3/2}/\varepsilon'\) is explicit in \(\varepsilon'\) but not in \(c\),
or from the log-period bound of V97, whose constant \(K\) is explicit.
Neither gives \(\varepsilon'\) from \(c\).  No global regularity claim is
made.
\end{assessment}

% =============================================================================
\section{Integrated V38 conclusion and next acquisition order}
\label{sec:c2v38-conclusion}
% =============================================================================

\begin{theorem}[What V38 establishes]
\label{thm:c2v38-conclusion}
Relative to the first cycle and V1--V37: the insufficiency of the
rate-free local energy bound for the lower side
(Proposition~\ref{prop:c2v38-too-weak}), the obstruction theorem
(Theorem~\ref{thm:c2v38-obstruction}), and the explicit form of the
upper side (Proposition~\ref{prop:c2v38-upper}).  The full
Navier--Stokes stability/global-regularity theorem remains unproved.
\end{theorem}

\begin{proof}
The cited results.
\end{proof}

\begin{programme}[V39 acquisition order]
\label{prog:c2v39}
\begin{enumerate}[label=\textup{(AC\arabic*)},leftmargin=4.8em]
\item Consolidate V36--V38 as the constant-chasing chapter: the Hardy
      inequality, the flux bounds, the mean constraints for
      \(\kappa<\frac12\) and \(\kappa<\sqrt2\), the pressure split, the
      gate and its obstruction.
\item Record the position for V40 and the direction for V41--V45: the
      spectral decomposition of the cross term \(\langle N(V),Y\rangle_G\)
      along Hermite modes, as the last unexplored exact structure of
      the homogeneous case, and the intermediate class (G2) with the
      Dirichlet ledger.
\item The next compulsory companion audit is V40.
\end{enumerate}
\end{programme}

\begin{assessment}[Position after V38]
\label{ass:c2v38-position}
The constant-chasing gate has an explicit upper side and a lower side
that requires quantitative unique continuation.  No full stability or
global-regularity claim is made.
\end{assessment}

% =============================================================================
\section*{Version V38 (second cycle) append-only record}
\addcontentsline{toc}{section}{Version V38 (second cycle) append-only record}
% =============================================================================

The complete V37 source of the second cycle was retained before this
continuation.  Its SHA--256 digest was
\begin{center}\scriptsize\ttfamily
0fda186936633f3d6c7bf8c44587f59b3b87d102a4b82635b385364264900f0c.
\end{center}
V38 recorded the obstruction to the lower side of the constant-chasing
gate and wrote its upper side in explicit form.  No full regularity
claim is made.

% =============================================================================
\section{V39: the constant-chasing chapter consolidated, and the direction
for the fifth block}
\label{sec:c2v39-entry}
% =============================================================================

Items (AC1) and (AC2) of Programme~\ref{prog:c2v39} are carried out.
V36--V38 are consolidated as the constant-chasing chapter.  The
direction for V41--V45 is then fixed in detail: the spectral
decomposition of the cross term of the second identity along the
Hermite modes of the self-similar operator, which is the last exact
structure of the homogeneous case not yet examined, and the
intermediate class of the first cycle's gate (G2) with the Dirichlet
ledger.  No global regularity claim is made.  The next version is the
compulsory companion audit.

% =============================================================================
\section{The constant-chasing chapter, consolidated}
\label{sec:c2v39-consolidated}
% =============================================================================

\begin{theorem}[Constant-chasing chapter, V36--V38]
\label{thm:c2v39-consolidated}
Let \(x_0\) be a homogeneous singular point of a type-I solution,
\(\kappa=C_\infty\nu^{-1/2}\) its velocity constant, \(K_6=C_6C_I\nu^{3/4}\),
\(K_P=4\nu^{-1/2}(4\pi\nu/3)^{1/4}C_{\rm CZ}K_6^2\), and \(\mathfrak m\) any
invariant probability measure on \(\mathcal F(x_0)\).  Modulo (E1)--(E9):
\begin{enumerate}[label=\textup{(\arabic*)},leftmargin=3.2em]
\item \emph{Hardy.}  \(\int|z|^2|V|^2G\le4\nu(d+3\varphi)\) at every scale.
\item \emph{Flux bounds.}  \(|j|\le\kappa(d+4\varphi)+K_P(d+3\varphi)^{1/2}\) and
      \(|j|\le2\kappa(\varphi d)^{1/2}+K_P(d+3\varphi)^{1/2}\).
\item \emph{Mean constraints.}  For \(\kappa<\frac12\):
      \(\langle d+3\varphi\rangle\le(K_P/m_\kappa)^2\); for \(\kappa<\sqrt2\):
      \(\langle d\rangle+\langle\varphi\rangle\le3K_P^2/m_\kappa'^2\) with
      \(m'_\kappa=\frac12(3-(1+4\kappa^2)^{1/2})\); in particular
      \(c_\Phi\le3K_P^2/m_\kappa'^2\).
\item \emph{Gate.}  \(\langle d\rangle\ge2c_2/\log(S/\eta)\) with the
      non-explicit compactness constants of the first cycle; the point
      is excluded if \(2c_2/\log(S/\eta)>3K_P^2/m_\kappa'^2\); the upper
      side is explicit in \(C_I\) and \(\nu\), the lower side is not, and
      an explicit lower side needs quantitative unique continuation.
\item \emph{Pressure split.}  A fully linear flux bound with total
      coefficient below one exists only for \(\kappa\) below an
      exponentially small threshold \(\exp(-cK_6^4/\nu^2)\), which is
      below the \(\varepsilon\)-regularity floor; the square-root pressure
      term is not removable by the local/far split.
\end{enumerate}
\end{theorem}

\begin{proof}
(1) Lemma~\ref{lem:c2v36-hardy}.  (2) Theorem~\ref{thm:c2v36-flux-bound},
Proposition~\ref{prop:c2v37-energy}.  (3) Theorems~\ref{thm:c2v36-constraint},
\ref{thm:c2v37-constraint}.  (4) Corollary~\ref{cor:c2v36-exclusion},
Propositions~\ref{prop:c2v38-upper}, \ref{prop:c2v38-lower},
Theorem~\ref{thm:c2v38-obstruction}.  (5) Corollary~\ref{cor:c2v37-linear},
Firewall~\ref{fw:c2v37-AC2}.
\end{proof}

\begin{assessment}[What the chapter changed]
\label{ass:c2v39-changed}
Before this chapter every signed statement of the cycle was a balance:
an exact identity in which a nonlinear moment equals a damping plus a
dissipation, with no bound on the nonlinear moment other than order
one.  The chapter supplies the first bound on a nonlinear moment that
is proportional to the positive terms with a coefficient
(\(\kappa\), or \(2\kappa\) times a geometric mean) that can be less than
one, and derives from it the first explicit inequality that a
hypothetical homogeneous blow-up must satisfy.  The inequality is not a
contradiction because the quantity it bounds, the mean Gaussian
dissipation, is bounded below only through compactness.  The chapter
therefore relocates the open problem for homogeneous points from
``find a mean inequality for the cross term'' to ``make the
concentration constant explicit'', which is a quantitative
unique-continuation problem.
\end{assessment}

% =============================================================================
\section{Direction for the fifth block}
\label{sec:c2v39-direction}
% =============================================================================

\begin{programme}[V41--V45: the spectral cross term and the intermediate class]
\label{prog:c2v39-block}
\begin{enumerate}[label=\textup{(D\arabic*)},leftmargin=3.8em]
\item \emph{Spectral cross term.}  Expand \(V\), \(Y=\partial_\sigma V\), and
      \(N(V)\) in the Hermite eigenbasis of \(L_-\) on divergence-free
      \(L^2(G)\): \(V=\sum_kV_k\), \(Y=\sum_k(-\mu_kV_k-N_k)\) with
      \(N_k=\Pi_kN(V)\), and write \(\|Y\|_G^2+\langle N,Y\rangle_G=\sum_k(\mu_k^2\|V_k\|^2+\mu_k\langle V_k,N_k\rangle)\ldots\)
      exactly; determine the modewise form of the second identity, the
      role of the constant mode (centre) and of the degree-one modes
      (linear fields), and whether the mean identity
      \(\langle\|Y\|_G^2+\langle N,Y\rangle_G\rangle=0\) splits into modewise
      balances with signs.
\item \emph{Low modes.}  The projection of the Navier--Stokes
      nonlinearity onto the constant and linear modes is computable
      from the Gaussian moments of \(V\otimes V\) (the tensor \(T\) of
      V18) and of the pressure; determine whether the low-mode balance
      is closed in terms of the moment ledger and whether it yields a
      new constraint.
\item \emph{Intermediate class.}  For a singular point of the
      intermediate class (type I along a sequence \(t_j\), no rate), the
      scales \(\lambda_j=(T_*-t_j)^{1/2}\) are \(M\)-parabolic with \(M\)
      explicit in the sequence bound; apply the rate-free Dirichlet
      ledger and the Rayleigh identity between consecutive parabolic
      scales, and determine what the absence of a rate between them
      forces on the push and on the Rayleigh quotient.
\item \emph{Consolidation at V45.}
\end{enumerate}
\end{programme}

% =============================================================================
\section{Integrated V39 conclusion and next acquisition order}
\label{sec:c2v39-conclusion}
% =============================================================================

\begin{theorem}[What V39 establishes]
\label{thm:c2v39-conclusion}
Relative to the first cycle and V1--V38: the consolidated
constant-chasing chapter (Theorem~\ref{thm:c2v39-consolidated}) and the
programme for the fifth block.  The full Navier--Stokes
stability/global-regularity theorem remains unproved.
\end{theorem}

\begin{proof}
The cited results.
\end{proof}

\begin{programme}[V40 acquisition order, with compulsory companion audit]
\label{prog:c2v40}
\begin{enumerate}[label=\textup{(AC\arabic*)},leftmargin=4.8em]
\item Perform the compulsory review of the current SM and YM notes;
      record digests; import nothing numerical.
\item Consolidate V36--V39 and record the position after forty
      versions; begin (D1) of Programme~\ref{prog:c2v39-block} with the
      modewise form of the second identity.
\item The next compulsory companion audit after V40 is V45.
\end{enumerate}
\end{programme}

\begin{assessment}[Position after V39]
\label{ass:c2v39-position}
The constant-chasing chapter is consolidated and the fifth block is
programmed.  No full stability or global-regularity claim is made.
\end{assessment}

% =============================================================================
\section*{Version V39 (second cycle) append-only record}
\addcontentsline{toc}{section}{Version V39 (second cycle) append-only record}
% =============================================================================

The complete V38 source of the second cycle was retained before this
continuation.  Its SHA--256 digest was
\begin{center}\scriptsize\ttfamily
6562ea959dda77a15289a7b751d3e5e64c6690f80c83bb778b2b82b1f7269927.
\end{center}
V39 consolidated the constant-chasing chapter and programmed the fifth
block.  No full regularity claim is made.

% =============================================================================
\section{V40: compulsory companion audit, the position after forty versions,
and the modewise identities}
\label{sec:c2v40-entry}
% =============================================================================

This is the fortieth version of the second cycle.  The compulsory review
of the two companion programmes is performed first.  V36--V39 are then
consolidated with the position after forty versions.  The fifth block
is then opened with item (D1) of Programme~\ref{prog:c2v39-block}: the
first and second identities are decomposed along the Hermite modes of
the self-similar operator, each mode's Gaussian energy obeys an exact
equation with its own damping rate and its own nonlinear feeding, and
against every invariant measure the mean feeding of every mode is
minus its eigenvalue times its mean energy: a modewise family of
inward-flux laws whose sum is the inward flux of the first cycle and
whose weighted sum is the second identity's mean form.  The constant
and linear modes are made explicit in terms of the Gaussian mean and
the dipole matrix.  No global regularity claim is made.  The next
compulsory companion audit is V45.

% =============================================================================
\section{Compulsory V40 review of the current companion notes}
\label{sec:c2v40-companion-audit}
% =============================================================================

The current companion files in \texttt{latex\_notes} were inspected
read-only.  The frozen checkpoints are
\begin{align}
 &\texttt{Renewal\_SM\_Einstein\_Continuum\_Research\_Notes\_V97.tex},
 \notag\\[-2pt]
 &\qquad
 \texttt{534150582ba302223a088de2acb0f31ca11622fddff027b9f50fb82bc2a5fdf8},
 \label{eq:c2v40-SM-checkpoint}\\
 &\texttt{Renewal\_Yang\_Mills\_Mass\_Gap\_Research\_Notes\_V73.tex},
 \notag\\[-2pt]
 &\qquad
 \texttt{926374cef6dfd909b294935122ee6789b05bdfe58f286454bc02261129c1c69a}.
 \label{eq:c2v40-YM-checkpoint}
\end{align}
A YM file named \(\ldots\)\texttt{V72.tex} with the same digest as
\eqref{eq:c2v40-YM-checkpoint} was also present; the V73 name is
recorded.

\emph{SM V96--V97.}  The SM programme proved finite-action support and a
positive-type statement (V96), then a resolvent sign-difference bound,
computed a four-dimensional Wilson-link Lipschitz constant, obtained
volume-independent overlap and moving-projector bounds, an explicit
source-fibre transport, isolated a scale limitation of a fixed free
reference, formulated a representation-matching condition so that it
is not assumed, proved an ultralocal Higgs multiplication bound, and
an assembled fibrewise singular-value criterion (V97).

\emph{YM V71--V72.}  The YM programme proved a double-centre Bell
resummation, a connected-graph-to-tree majorant, constructed an
effective pair-edge generating function with an adaptive-circle bound,
resummed all one-centre decorations simultaneously, avoided a false
label-allocation factor, proved a pair-core coefficient bound, closed a
pair-core card at arbitrary centre arity, and isolated genuine centre
hyperedges as its next row.

\begin{theorem}[V40 companion-audit decision]
\label{thm:c2v40-companion-decision}
No SM V96--V97 transport, Lipschitz, or singular-value constant and no
YM V71--V72 resummation or coefficient bound is imported.  Neither
bears on the constant-chasing chapter or on the modewise identities
below.
\end{theorem}

\begin{proof}
The SM results concern lattice overlap operators and their fibres; the
YM results concern graph resummations of a lattice clock.  None is a
Navier--Stokes statement.
\end{proof}

% =============================================================================
\section{Position after forty versions}
\label{sec:c2v40-position}
% =============================================================================

\begin{assessment}[Position after forty versions of the second cycle]
\label{ass:c2v40-position-forty}
To the chapters (K1)--(K7) of Theorem~\ref{thm:c2v35-chapters} the
cycle has added the constant-chasing chapter
(Theorem~\ref{thm:c2v39-consolidated}): the first bound on a nonlinear
moment proportional to the positive terms with a coefficient that can
be less than one, the first explicit inequality on a hypothetical
homogeneous blow-up, and the identification of the missing lower side
as a quantitative unique-continuation problem.  The three gates of the
first cycle stand.  The fifth block examines the modewise structure of
the identities and the intermediate class.  No global regularity claim
is made.
\end{assessment}

% =============================================================================
\section{The modewise identities}
\label{sec:c2v40-modewise}
% =============================================================================

Let \(W\in\mathcal A_\nu(C)\) with the Morrey bound and the uniform bounds,
\(V\) its profile, \(N=N(V)\), \(Y=\partial_\sigma V=L_-V-N\).  Let \(\Pi_k\) be the
\(G\)-orthogonal projection onto the eigenspace of \(-L_-\) with
eigenvalue \(\mu_k=\frac12+\frac k2\) (Hermite degree \(k\)), acting
componentwise; \(\Pi_k\) commutes with \(L_-\), preserves the
divergence-free condition (Proposition~\ref{prop:c2v8-spectrum}(ii)),
and commutes with \(\partial_\sigma\).  Put
\begin{equation}
 V_k:=\Pi_kV,\quad N_k:=\Pi_kN,\quad
 a_k:=\|V_k\|_G^2,\quad n_k:=\langle V_k,N\rangle_G=\langle V_k,N_k\rangle_G,\quad
 b_k:=\mu_ka_k+n_k .
 \label{eq:c2v40-modes}
\end{equation}
Then \(\varphi=\sum_ka_k\), \(r\varphi=2\mathcal L=\sum_k\mu_ka_k\),
\(\langle V,N\rangle_G=\sum_kn_k=\frac14j\), all sums absolutely convergent
(\(\sum\mu_ka_k=\nu\|\nabla V\|_G^2+\frac12\varphi\) and
\(\sum|n_k|\le\|V\|_G\|N\|_G\)).

\begin{theorem}[Modewise identities]
\label{thm:c2v40-modewise}
For every \(k\ge0\) and every \(\sigma\),
\begin{equation}
 \boxed{\quad
 \frac{\dd a_k}{\dd\sigma}=-2\mu_ka_k-2n_k=-2b_k ,
 \qquad
 \frac{\dd\varphi}{\dd\sigma}=-2\sum_kb_k,
 \qquad
 \frac{\dd\mathcal L}{\dd\sigma}=-\sum_k\mu_kb_k ,
 \quad}
 \label{eq:c2v40-modewise}
\end{equation}
and \(\|Y\|_G^2+\langle N,Y\rangle_G=\sum_k\mu_kb_k\),
\(\|Y\|_G^2=\sum_k\|\mu_kV_k+N_k\|_G^2\).  Against every invariant
probability measure \(\mathfrak m\) on a rescaling family at a
homogeneous point, for every \(k\),
\begin{equation}
 \boxed{\quad
 \langle n_k\rangle_{\mathfrak m}=-\mu_k\langle a_k\rangle_{\mathfrak m}\ \le\ 0 ,
 \qquad\text{equivalently}\qquad
 \langle b_k\rangle_{\mathfrak m}=0 ,
 \quad}
 \label{eq:c2v40-mean-modewise}
\end{equation}
so that \(\sum_k\langle n_k\rangle=-\sum_k\mu_k\langle a_k\rangle=-2\langle\mathcal L\rangle\)
(the inward flux) and \(\sum_k\mu_k\langle b_k\rangle=0\) (the mean second
identity).
\end{theorem}

\begin{proof}
\(\partial_\sigma V_k=\Pi_k\partial_\sigma V=\Pi_k(L_-V-N)=-\mu_kV_k-N_k\), so
\(\frac{\dd}{\dd\sigma}a_k=2\langle V_k,\partial_\sigma V_k\rangle_G=-2\mu_ka_k-2\langle V_k,N_k\rangle_G\),
and \(\langle V_k,N_k\rangle_G=\langle V_k,N\rangle_G\) by orthogonality of the
projections.  Summing over \(k\) gives the first identity
\(\frac{\dd\varphi}{\dd\sigma}=2\langle V,Y\rangle_G=-2\sum_k(\mu_ka_k+n_k)\); and
\(\frac{\dd\mathcal L}{\dd\sigma}=-\langle L_-V,Y\rangle_G=-\sum_k\langle-\mu_kV_k,-\mu_kV_k-N_k\rangle_G=-\sum_k(\mu_k^2a_k+\mu_kn_k)\).
For the defect, \(Y=\sum_k(-\mu_kV_k-N_k)\) with orthogonal summands, and
\(\|Y\|_G^2+\langle N,Y\rangle_G=\sum_k\bigl(\|\mu_kV_k+N_k\|^2-\langle N_k,\mu_kV_k+N_k\rangle\bigr)=\sum_k(\mu_k^2a_k+\mu_kn_k)\).
The mean identities: \(a_k\) is a bounded observable (\(a_k\le\varphi\le C_\Phi\)),
continuous on the family (as \(\varphi\) is, the projection being a
bounded operator), so the mean of its derivative vanishes.  The
interchange of sum and mean is justified by absolute convergence with
\(\mathfrak m\)-integrable majorants (\(\sum\mu_ka_k=2\mathcal L\) and
\(\sum|n_k|\le\|V\|_G\|N\|_G\), both bounded on the family).
\end{proof}

\begin{assessment}[Modewise inward flux]
\label{ass:c2v40-modewise}
Each Hermite mode of the profile is damped by the self-similar scaling
at its own rate \(2\mu_k=1+k\) and fed by the projection of the
nonlinearity onto it; statistically, every mode is fed exactly at the
rate it is damped, so the nonlinearity transfers Gaussian energy into
every occupied mode, and the transfer into mode \(k\) is \(\mu_k\) times
the mean energy of that mode.  The first cycle's inward flux is the sum
of these transfers; the second identity's mean form is their
\(\mu_k\)-weighted sum, which vanishes identically because each term
does.  The modewise laws are therefore not new constraints beyond the
mean forms of V12 for each single mode; their content is the
statement that the nonlinear transfer is positive into every mode in
the mean, with no mode receiving net energy from the scaling.
\end{assessment}

\begin{proposition}[The constant and linear modes]
\label{prop:c2v40-low-modes}
With \(|G|=(4\pi\nu)^{3/2}\), the Gaussian mean \(\bar V\) and the dipole
matrix \(M_V=\frac1{2\nu|G|}\int z\otimes VG\) of V16:
\begin{enumerate}[label=\textup{(\roman*)},leftmargin=3.2em]
\item \(V_0=\bar V\) (constant), \(a_0=|G||\bar V|^2\), \(n_0=|G|\bar V\cdot\bar N\),
      and \eqref{eq:c2v40-mean-modewise} for \(k=0\) is
      \(\langle\bar V\cdot\bar N\rangle=-\frac12\langle|\bar V|^2\rangle\), the trace of
      \eqref{eq:c2v22-signed} for \(S_0\).
\item \(V_1=A_Vz\) with \(A_V=\frac1{2\nu|G|}\int V\otimes zG=M_V^{\mathsf T}\),
      a traceless matrix (divergence-free), \(a_1=2\nu|G|\,|A_V|^2\)
      (Frobenius norm), \(n_1=A_V:\int N\otimes zG=2\nu|G|\,A_V:M_N^{\mathsf T}\)
      with \(M_N\) the dipole matrix of \(N\), and
      \eqref{eq:c2v40-mean-modewise} for \(k=1\) is
      \(\langle A_V:M_N^{\mathsf T}\rangle=-\langle|A_V|^2\rangle\).
\end{enumerate}
\end{proposition}

\begin{proof}
The Hermite polynomials of degree \(0\) and \(1\) in \(z/(2\nu)^{1/2}\) are
\(1\) and \(z_i/(2\nu)^{1/2}\), with \(\int z_iz_jG=2\nu\delta_{ij}|G|\); the
projections are the Gaussian mean and the linear regression of \(V\) on
\(z\), and the divergence-free condition on \(A_Vz\) is \(\operatorname{tr}A_V=0\).
The inner products follow from \(\int z_iz_jG=2\nu|G|\delta_{ij}\).
\end{proof}

% =============================================================================
\section{Integrated V40 conclusion and next acquisition order}
\label{sec:c2v40-conclusion}
% =============================================================================

\begin{theorem}[What V40 establishes]
\label{thm:c2v40-conclusion}
Relative to the first cycle and V1--V39: the companion-audit decision
(Theorem~\ref{thm:c2v40-companion-decision}), the modewise identities
\eqref{eq:c2v40-modewise} with their mean form
\eqref{eq:c2v40-mean-modewise}, and the explicit constant and linear
modes.  The full Navier--Stokes stability/global-regularity theorem
remains unproved.
\end{theorem}

\begin{proof}
The cited results.
\end{proof}

\begin{programme}[V41 acquisition order]
\label{prog:c2v41}
\begin{enumerate}[label=\textup{(AC\arabic*)},leftmargin=4.8em]
\item Modewise flux bound.  Bound \(|n_k|=|\langle V_k,N\rangle_G|\) for each
      \(k\) in terms of \(a_k\) and the positive terms, using the Hermite
      structure (the mode \(V_k\) is a polynomial field of degree \(k\),
      so \(\langle V_k,(V\cdot\nabla)V\rangle_G=-\frac12\int|V|^2\operatorname{div}(V_kG)\ldots\)
      is a Gaussian moment of \(|V|^2\) against a polynomial of degree
      \(k+1\)), and determine whether the modewise mean law
      \(\langle n_k\rangle=-\mu_k\langle a_k\rangle\) with a modewise bound
      \(|n_k|\le\kappa_k(\ldots)\) gives a constraint for the low modes
      analogous to the constant-chasing constraint but with the
      explicit low-mode damping \(\mu_k\) in place of the global
      coefficient.
\item The linear mode and the dipole of the nonlinearity.  Compute
      \(M_N\) through the Gaussian moments of \(V\otimes V\) and of the
      pressure, and determine whether the \(k=1\) law is closed within
      the moment ledger.
\item The next compulsory companion audit is V45.
\end{enumerate}
\end{programme}

\begin{assessment}[Position after V40]
\label{ass:c2v40-position}
The identities of the cycle are decomposed along the Hermite modes,
with a modewise inward-flux law for every mode.  No full stability or
global-regularity claim is made.
\end{assessment}

% =============================================================================
\section*{Version V40 (second cycle) append-only record}
\addcontentsline{toc}{section}{Version V40 (second cycle) append-only record}
% =============================================================================

The complete V39 source of the second cycle was retained before this
continuation.  Its SHA--256 digest was
\begin{center}\scriptsize\ttfamily
67d9b5073a4fc3ad0705807cb6ee0c6d32f0f0cdcfc6959bc287817a28469e45.
\end{center}
V40 performed the compulsory companion audit, recorded the position
after forty versions, and proved the modewise identities.  No full
regularity claim is made.

% =============================================================================
\section{V41: modewise flux bounds and the mean Hermite spectrum of a
homogeneous blow-up}
\label{sec:c2v41-entry}
% =============================================================================

Items (AC1) and (AC2) of Programme~\ref{prog:c2v41} are carried out.  The
nonlinear feeding of each Hermite mode is bounded, using the polynomial
structure of the mode, the \(L^\infty\) rate, the Gaussian Hardy
inequality, and the pressure bound, by the square root of the mode's
energy times an explicit factor growing like the square root of the
degree.  Combined with the modewise inward-flux law this gives, for
every mode, an explicit upper bound on its mean Gaussian energy against
every invariant measure, decaying like the inverse of the degree; this
is the first explicit statement about the spectral distribution of a
hypothetical homogeneous blow-up.  The linear-mode law is written in
terms of the Gaussian moments of the profile and of the pressure and is
shown not to close within the moment ledger.  No global regularity
claim is made.  The next compulsory companion audit is V45.

% =============================================================================
\section{Modewise flux bounds}
\label{sec:c2v41-bounds}
% =============================================================================

Notation as in V40 (\(V_k=\Pi_kV\), \(a_k=\|V_k\|_G^2\), \(n_k=\langle V_k,N(V)\rangle_G\),
\(\mu_k=\frac12+\frac k2\)) and V36--V37 (\(\kappa\), \(K_P\), \(X^2=\int|z|^2|V|^2G\)).

\begin{lemma}[Hermite modes]
\label{lem:c2v41-modes}
For every \(k\), \(V_k\) is a divergence-free polynomial field of degree
\(k\) with \(\nu\|\nabla V_k\|_G^2=\frac k2a_k\) and
\(\int|z|^2|V_k|^2G\le4\nu(3+2k)\,a_k\).
\end{lemma}

\begin{proof}
\(\langle-AV_k,V_k\rangle_G=\nu\|\nabla V_k\|_G^2\) for \(A=\nu\Delta-\frac12z\cdot\nabla\)
(the \(G\)-symmetry) and \(-AV_k=\frac k2V_k\).  The Hardy inequality of
Lemma~\ref{lem:c2v36-hardy} applied to \(V_k\):
\(\int|z|^2|V_k|^2G\le12\nu a_k+16\nu^2\|\nabla V_k\|_G^2=12\nu a_k+8\nu ka_k\).
\end{proof}

\begin{theorem}[Modewise flux bound]
\label{thm:c2v41-bound}
For every \(k\ge0\) and every \(\sigma\),
\begin{equation}
 \boxed{\quad
 |n_k|\ \le\ a_k^{1/2}\Bigl[\kappa\Bigl(\frac k2\Bigr)^{1/2}\varphi^{1/2}+\kappa\,(d+3\varphi)^{1/2}+\frac{K_P}4(3+2k)^{1/2}\Bigr].
 \quad}
 \label{eq:c2v41-bound}
\end{equation}
\end{theorem}

\begin{proof}
\(n_k=\langle V_k,(V\cdot\nabla)V\rangle_G+\langle V_k,\nabla P_V\rangle_G\).  For the
advective part, integrate by parts using \(\operatorname{div}V=0\) and
\(\nabla G=-\frac z{2\nu}G\):
\[
 \int V_{k,i}V_j\partial_jV_i\,G=-\int V_iV_j\partial_jV_{k,i}\,G-\int V_iV_{k,i}V_j\partial_jG
 =-\int(V\otimes V):\nabla V_k\,G+\frac1{2\nu}\int(V\cdot V_k)(z\cdot V)G .
\]
The first term is bounded by \(C_\infty\|V\|_G\|\nabla V_k\|_G=C_\infty\varphi^{1/2}(ka_k/(2\nu))^{1/2}=\kappa(k/2)^{1/2}(\varphi a_k)^{1/2}\)
by Lemma~\ref{lem:c2v41-modes}; the second by
\(\frac{C_\infty}{2\nu}\|V_k\|_GX\le\frac{C_\infty}{2\nu}a_k^{1/2}\cdot2\nu^{1/2}(d+3\varphi)^{1/2}=\kappa a_k^{1/2}(d+3\varphi)^{1/2}\)
by Lemma~\ref{lem:c2v36-hardy}.  For the pressure part,
\(\operatorname{div}V_k=0\) gives \(\langle V_k,\nabla P_V\rangle_G=\frac1{2\nu}\int P_V(z\cdot V_k)G\)
(the pressure may be taken mean-free since \(\int(z\cdot V_k)G=0\) by the
same computation as for \(V\)), and H\"older as in
Theorem~\ref{thm:c2v36-flux-bound} with \(\int|z|^2|V_k|^2G\le4\nu(3+2k)a_k\)
gives at most
\(\frac1{2\nu}C_{\rm CZ}K_6^2(4\pi\nu/3)^{1/4}\cdot2\nu^{1/2}(3+2k)^{1/2}a_k^{1/2}=\frac{K_P}4(3+2k)^{1/2}a_k^{1/2}\).
\end{proof}

% =============================================================================
\section{The mean Hermite spectrum}
\label{sec:c2v41-spectrum}
% =============================================================================

\begin{theorem}[Explicit bound on the mean energy of each mode]
\label{thm:c2v41-spectrum}
Let \(x_0\) be a homogeneous singular point of a type-I solution and
\(\mathfrak m\) an invariant probability measure on \(\mathcal F(x_0)\).
Then, modulo (E1)--(E9), for every \(k\ge0\),
\begin{equation}
 \boxed{\quad
 \langle a_k\rangle_{\mathfrak m}\ \le\ \frac{1}{\mu_k^2}\Bigl[\kappa\Bigl(\frac k2\Bigr)^{1/2}\langle\varphi\rangle^{1/2}+\kappa\langle d+3\varphi\rangle^{1/2}+\frac{K_P}4(3+2k)^{1/2}\Bigr]^2 ,
 \quad}
 \label{eq:c2v41-spectrum}
\end{equation}
so that \(\langle a_k\rangle\le\frac{C_{\rm sp}}{k}\) for \(k\ge1\) with
\(C_{\rm sp}:=8\bigl[\kappa\langle\varphi\rangle^{1/2}+\kappa\langle d+3\varphi\rangle^{1/2}+K_P\bigr]^2\),
and \(\langle a_0\rangle\le4\bigl[\kappa\langle d+3\varphi\rangle^{1/2}+\frac{\sqrt3}4K_P\bigr]^2\).
Under the constant-chasing constraint (\(\kappa<\sqrt2\)) the brackets
are explicit in \(C_I\), \(\kappa\), \(\nu\).
\end{theorem}

\begin{proof}
By \eqref{eq:c2v40-mean-modewise}, \(\mu_k\langle a_k\rangle=-\langle n_k\rangle\le\langle|n_k|\rangle\),
and by \eqref{eq:c2v41-bound} with Cauchy--Schwarz against \(\mathfrak m\),
\(\langle|n_k|\rangle\le\langle a_k\rangle^{1/2}\bigl[\kappa(k/2)^{1/2}\langle\varphi\rangle^{1/2}+\kappa\langle d+3\varphi\rangle^{1/2}+\frac{K_P}4(3+2k)^{1/2}\bigr]\).
If \(\langle a_k\rangle>0\), divide by \(\langle a_k\rangle^{1/2}\) and square.  For
\(k\ge1\), \(\mu_k\ge k/2\) and \((k/2)^{1/2},(3+2k)^{1/2}\le(3k)^{1/2}\), giving
the \(1/k\) form after crude constants.
\end{proof}

\begin{assessment}[The mean spectrum]
\label{ass:c2v41-spectrum}
A homogeneous blow-up cannot, statistically, store its Gaussian energy
in a single Hermite mode beyond an explicit amount, and its mean
spectral density decays at least like the inverse degree.  This is
weaker than the summability \(\sum_k\mu_k\langle a_k\rangle=2\langle\mathcal L\rangle<\infty\)
that the class bound already gives, but it is explicit mode by mode,
and for the constant mode it bounds the mean squared Gaussian momentum,
\(|G|\langle|\bar V|^2\rangle=\langle a_0\rangle\), by an explicit constant: the
blow-up centre's mean Gaussian drift is explicitly bounded.  The bounds
are consequences of the modewise inward-flux laws, hence balances read
as inequalities; they do not contradict anything.
\end{assessment}

% =============================================================================
\section{The linear-mode law and the moment ledger}
\label{sec:c2v41-linear}
% =============================================================================

\begin{proposition}[The dipole of the nonlinearity]
\label{prop:c2v41-dipole}
With \(T_{ij}=\int V_iV_jG\), \(Q_{ij}:=\frac1{2\nu}\int z_i(z\cdot V)V_jG\),
\(\bar P:=\frac1{|G|}\int P_VG\), \(R^P_{ij}:=\frac1{2\nu}\int z_iz_jP_VG\),
\begin{equation}
 \int z_iN_jG=-T_{ij}+Q_{ij}-\delta_{ij}|G|\bar P+R^P_{ij},
 \label{eq:c2v41-dipole}
\end{equation}
and the \(k=1\) law \(\langle A_V:M_N^{\mathsf T}\rangle=-\langle|A_V|^2\rangle\) reads,
with \(A_V=M_V^{\mathsf T}\) and \(2\nu|G|M_N^{\mathsf T}=\int N\otimes zG\),
\begin{equation}
 \Bigl\langle A_V:\bigl(-T+Q-|G|\bar P\,I+R^P\bigr)\Bigr\rangle=-2\nu|G|\,\langle|A_V|^2\rangle .
 \label{eq:c2v41-linear-law}
\end{equation}
The moments \(Q\) (a third Gaussian moment of the profile) and \(R^P\) (a
second Gaussian moment of the pressure) do not belong to the moment
ledger of V16--V22, so the linear-mode law is not closed within it.
\end{proposition}

\begin{proof}
\(\int z_i((V\cdot\nabla)V)_jG=-\int V_j\partial_l(z_iV_lG)=-\int V_iV_jG-\int V_jV_lz_i\partial_lG=-T_{ij}+\frac1{2\nu}\int z_i(z\cdot V)V_jG\)
using \(\operatorname{div}V=0\) and \(\partial_lG=-\frac{z_l}{2\nu}G\); and
\(\int z_i\partial_jP_VG=-\int P_V\partial_j(z_iG)=-\int P_V(\delta_{ij}-\frac{z_iz_j}{2\nu})G\).
The law is \eqref{eq:c2v40-mean-modewise} for \(k=1\) with
Proposition~\ref{prop:c2v40-low-modes}(ii); the trace term
\(-|G|\bar P\operatorname{tr}A_V\) vanishes since \(A_V\) is traceless.
\end{proof}

\begin{firewall}[Item (AC2)]
\label{fw:c2v41-AC2}
The linear-mode law involves the third moment \(Q\) and the pressure
second moment \(R^P\), and every attempt to close a modewise law within
the moments of lower order meets the same structure: the nonlinearity
raises the moment order by one.  The ledger of V16--V22 is closed under
the identities but not under the modewise laws, which form an infinite
hierarchy; the series records the hierarchy's first two rows
(Proposition~\ref{prop:c2v40-low-modes}(i), \eqref{eq:c2v41-linear-law})
and the modewise bounds \eqref{eq:c2v41-spectrum} as its explicit
content.
\end{firewall}

% =============================================================================
\section{Integrated V41 conclusion and next acquisition order}
\label{sec:c2v41-conclusion}
% =============================================================================

\begin{theorem}[What V41 proves]
\label{thm:c2v41-conclusion}
Relative to the first cycle and V1--V40, modulo (E1)--(E9): the modewise
flux bound \eqref{eq:c2v41-bound}, the explicit mean Hermite spectrum
bound \eqref{eq:c2v41-spectrum}, and the dipole identity
\eqref{eq:c2v41-dipole} with the linear-mode law
\eqref{eq:c2v41-linear-law}.  The full Navier--Stokes
stability/global-regularity theorem remains unproved.
\end{theorem}

\begin{proof}
The cited results.
\end{proof}

\begin{programme}[V42 acquisition order]
\label{prog:c2v42}
\begin{enumerate}[label=\textup{(AC\arabic*)},leftmargin=4.8em]
\item The intermediate class (D3).  For a singular point with type I
      along a sequence \(t_j\uparrow T_*\) (\((T_*-t_j)^{1/2}Y(t_j)\le M\)) and
      no rate, show that the scales \(\lambda_j=(T_*-t_j)^{1/2}\) are
      \(M'\)-parabolic in the sense of V26 with \(M'\) explicit in \(M\),
      \(E_*\), \(\nu\) (energy at the parabolic time from the enstrophy
      bound by Sobolev on the torus; dissipation over the window from
      the energy inequality; pressure from the cubic bound), and hence
      that the rescalings along \(\lambda_j\) converge to suitable weak
      solutions, and determine whether the limit is nonzero (the
      enstrophy bound at the single time \(t_j\) gives strong \(H^1\)
      compactness at that time only).
\item Apply the rate-free Dirichlet ledger between consecutive
      \(\lambda_j\) and record what the absence of a rate between them
      allows.
\item The next compulsory companion audit is V45.
\end{enumerate}
\end{programme}

\begin{assessment}[Position after V41]
\label{ass:c2v41-position}
The mean Hermite spectrum of a homogeneous blow-up is explicitly
bounded mode by mode; the modewise laws form a hierarchy not closed
within the moment ledger.  No full stability or global-regularity
claim is made.
\end{assessment}

% =============================================================================
\section*{Version V41 (second cycle) append-only record}
\addcontentsline{toc}{section}{Version V41 (second cycle) append-only record}
% =============================================================================

The complete V40 source of the second cycle was retained before this
continuation.  Its SHA--256 digest was
\begin{center}\scriptsize\ttfamily
f07eff4f2335bdd5f20e52d0789a5c62d4f21944f0275ae03e0731b2dbadba3f.
\end{center}
V41 proved the modewise flux bounds, the explicit mean Hermite spectrum
bound, and the linear-mode law with its moment structure.  No full
regularity claim is made.

% =============================================================================
\section{V42: the intermediate class, its type-I windows, and the ledgers
between them}
\label{sec:c2v42-entry}
% =============================================================================

Item (D3) of Programme~\ref{prog:c2v39-block} is begun, following
Programme~\ref{prog:c2v42}.  A singular point of the intermediate class,
type I along a sequence of times without a rate, is shown to possess
type-I windows: along a subsequence of parabolic scales, the rescaled
solution has enstrophy bounded by twice the sequence constant on a
forward window of log-scale of explicit length inversely proportional
to the square of that constant, by the classical enstrophy differential
inequality; on these windows the rescalings are uniformly bounded in
the strong norms, converge along subsequences to strong solutions with
strong convergence of gradients, and all Gaussian ledger quantities
are bounded by explicit functions of the sequence constant.  Between
consecutive windows no rate is available, and the rate-free Dirichlet
ledger gives exactly one constraint: the log-integral of the signed
combination of squared defect and cross term over the intermediate
stretch is bounded by an explicit constant, whatever happens to the
enstrophy in between.  Whether the window limits are nonzero is
recorded as open, for the reason of spatial spreading.  No global
regularity claim is made.  The next compulsory companion audit is V45.

% =============================================================================
\section{Type-I windows}
\label{sec:c2v42-windows}
% =============================================================================

\begin{definition}[Intermediate class]
\label{def:c2v42-intermediate}
A smooth solution with \(T_*<\infty\) is in the \emph{intermediate class
with constant \(M\)} if there are \(t_j\uparrow T_*\) with
\((T_*-t_j)^{1/2}\|\Lambda u(t_j)\|_2^2\le M\) for all \(j\), and no type-I rate
holds on any interval \([T_*-\delta,T_*)\).  Put \(\lambda_j:=(T_*-t_j)^{1/2}\) and,
for a centre \(x_0\), \(U_j(s,y):=\lambda_ju(x_0+\lambda_jy,T_*+\lambda_j^2s)\) on the
torus of side \(2\pi/\lambda_j\), so that \(\|\nabla U_j(-1)\|_2^2=\lambda_j\|\Lambda u(t_j)\|_2^2\le M\).
\end{definition}

\begin{proposition}[Forward enstrophy windows]
\label{prop:c2v42-window}
Let \(C_A\) be the constant of the enstrophy inequality
\(\frac{\dd}{\dd t}\|\Lambda u\|_2^2\le C_A\nu^{-3}\|\Lambda u\|_2^6\) for smooth
mean-zero solutions on the torus (used in the first cycle in the form
\(Y'\le CY^3\)).  Then in the rescaled variables, with \(Y_j(s):=\|\nabla U_j(s)\|_2^2\),
\begin{equation}
 Y_j(s)\ \le\ 2M\qquad\text{for }-1\le s\le-1+\ell_M,\qquad
 \ell_M:=\frac{3\nu^3}{8C_AM^2},
 \label{eq:c2v42-window}
\end{equation}
provided \(\ell_M<1\) (otherwise on \([-1,0)\)); in log-scale
\(\sigma=-\log(-s)\) this is a window of length at least
\(\log\frac1{1-\ell_M}\ge\ell_M\) starting at \(\sigma=0\).  On this window,
moreover, \(\int_{-1}^{-1+\ell_M}\|\Delta U_j(s)\|_2^2\dd s\le C(M,\nu)\).
\end{proposition}

\begin{proof}
The enstrophy inequality is scale-invariant, so \(Y_j'\le C_A\nu^{-3}Y_j^3\)
on the rescaled torus; hence \((Y_j^{-2})'\ge-2C_A\nu^{-3}\), and
\(Y_j(s)^{-2}\ge M^{-2}-2C_A\nu^{-3}(s+1)\ge\frac14M^{-2}\) for
\(s+1\le\frac{3\nu^3}{8C_AM^2}\).  The \(H^2\) bound follows from the same
inequality kept with its dissipation term,
\(\frac{\dd}{\dd s}Y_j\le-\nu\|\Delta U_j\|_2^2+C_A\nu^{-3}Y_j^3\), integrated
over the window with \(Y_j\le2M\).
\end{proof}

\begin{theorem}[Compactness and bounded ledgers on the windows]
\label{thm:c2v42-compactness}
Let \(u\) be in the intermediate class with constant \(M\), and \(x_0\) any
point.  Along a subsequence, the rescalings \(U_j\) converge on
\(B_R\times[-1,-1+\ell_M]\), for every \(R\), strongly in
\(L^2(H^1)\) and weakly-\(*\) in \(L^\infty(H^1)\cap L^2(H^2)\), to a strong
solution \(U'\) of the Navier--Stokes equations on
\(\R^3\times[-1,-1+\ell_M]\) with \(\|\nabla U'(s)\|_{L^2(\R^3)}^2\le2M\); the
limit is smooth on the open window.  On the window, for every \(j\) and
every scale \(\lambda\in[\lambda_j(1-\ell_M)^{1/2},\lambda_j]\), the Gaussian ledger
quantities of the solution at \(x_0\) satisfy
\begin{equation}
 \Phi_{u,x_0}(\lambda)\le C_S^2\cdot2M\cdot\|G\|_{L^{3/2}}\,,\qquad
 \lambda D_{u,x_0}(\lambda)\le8\nu M,\qquad
 \mathcal L_{u,x_0}(\lambda)\le\nu M+\tfrac14C_S^2\,2M\,\|G\|_{L^{3/2}} ,
 \label{eq:c2v42-ledgers}
\end{equation}
with \(C_S\) the (scale-invariant) constant of the mean-zero Sobolev
inequality \(\|f\|_{L^6}\le C_S\|\nabla f\|_{L^2}\) on the torus, and
\(\|G\|_{L^{3/2}}=(\int G^{3/2})^{2/3}\).
\end{theorem}

\begin{proof}
The bounds \eqref{eq:c2v42-window} give \(U_j\) bounded in
\(L^\infty(H^1)\cap L^2(H^2)\) on the window, on the big torus, hence on
every fixed ball; the equation gives \(\partial_sU_j\) bounded in
\(L^2(L^2)\) locally (the nonlinear term by \(\|U\|_\infty\)-free
interpolation \(\|(U\cdot\nabla)U\|_2\le\|U\|_6\|\nabla U\|_3\le C\|\nabla U\|_2^{3/2}\|\Delta U\|_2^{1/2}\),
the pressure gradient by the Riesz bound in \(L^2\)); Aubin--Lions (E2)
gives strong \(L^2(H^1)\) convergence on compacts, and the limit solves
the equations with the stated bound and is a strong solution, smooth
on the open window by the standard \(H^1\) theory.  For the ledgers:
\(\Phi_{u,x_0}(\lambda)=\|U_j(-\lambda^2/\lambda_j^2)\|_{L^2(G)}^2\) at the corresponding
rescaled time, and \(\int|U_j|^2G\le\|U_j\|_6^2\|G\|_{3/2}\le C_S^2Y_j\|G\|_{3/2}\)
with \(Y_j\le2M\); \(\lambda D_u=4\nu\int|\nabla U_j|^2G\le4\nu Y_j\le8\nu M\);
and \(\mathcal L=\frac18(\lambda D+2\Phi)\).
\end{proof}

\begin{firewall}[The window limit may be zero]
\label{fw:c2v42-zero}
The rescaled global enstrophy is bounded below by Leray's rate,
\(Y_j(s)\ge c_L^2\nu^{3/2}(-s)^{-1/2}\), but on a torus of side \(2\pi/\lambda_j\)
this enstrophy may be spread over a volume of order \(\lambda_j^{-3}\), so
that its restriction to every fixed ball tends to zero and the window
limit at any fixed centre is zero.  A nonzero limit requires a centre
map following the enstrophy, which is the compact-dissipation and
coherence question of V24 and V32 in the intermediate class.  The
series records that the window limits are strong solutions with the
explicit bounds above, possibly zero.
\end{firewall}

% =============================================================================
\section{The ledgers between windows}
\label{sec:c2v42-between}
% =============================================================================

\begin{theorem}[Bounded signed ledger across the intermediate stretches]
\label{thm:c2v42-between}
Let \(u\) be in the intermediate class with constant \(M\), \(x_0\) any
point, and \(\lambda_{j+1}<\lambda_j\) two consecutive window scales (of the
subsequence).  Then, without any rate on \((\lambda_{j+1},\lambda_j)\),
\begin{equation}
 \boxed{\quad
 \Bigl|\int_{\lambda_{j+1}}^{\lambda_j}\bigl(\|Y_u\|_G^2+\langle N(V_u),Y_u\rangle_G\bigr)\frac{\dd\lambda}{\lambda}\Bigr|
 \ \le\ \frac12\Bigl(\nu M+\tfrac12C_S^2M\|G\|_{L^{3/2}}\Bigr)=:\frac12\mathcal L_M ,
 \quad}
 \label{eq:c2v42-between}
\end{equation}
and likewise \(\bigl|\int_{\lambda_{j+1}}^{\lambda_j}\mathcal X_{u,x_0}\lambda^{-2}\dd\lambda-(\lambda_{j+1}^{-2}\Phi_u(\lambda_{j+1})-\lambda_j^{-2}\Phi_u(\lambda_j))\bigr|=0\)
with both endpoint values of \(\Phi_u\) bounded by \(2C_S^2M\|G\|_{L^{3/2}}\).
Consequently, on every intermediate stretch, either the log-integrals of
\(\|Y_u\|_G^2\) and of \(-\langle N(V_u),Y_u\rangle_G\) are both bounded by
\(\mathcal L_M\) plus the other, or they diverge together with the same
leading order: a blow-up of the enstrophy between two type-I windows
must be accompanied by anti-alignment of the nonlinearity with the
self-similarity defect of the same log-integrated size.
\end{theorem}

\begin{proof}
The rate-free ledger \eqref{eq:c2v14-ledger} between \(\lambda_{j+1}\) and
\(\lambda_j\), with \(0\le\mathcal L_u\le\mathcal L_M\) at both endpoints by
\eqref{eq:c2v42-ledgers}; and the first-identity ledger of
\eqref{eq:c2v1-identity} in the same way.  The last statement is the
bounded ledger read as in Theorem~\ref{thm:c2v28-dichotomy}(A).
\end{proof}

\begin{assessment}[The intermediate class in the ledgers]
\label{ass:c2v42-intermediate}
A solution of the intermediate class visits type-I windows infinitely
often, on which it is, in the rescaled picture, a strong solution with
explicit bounds, and between which it may develop enstrophy without any
rate.  The second cycle's ledgers give, for each such excursion, an
exact accounting: the Dirichlet functional returns to a bounded value at
the next window, so the excursion's net contribution to
\(\int(\|Y\|^2+\langle N,Y\rangle)\) is bounded, and any large excursion of the
defect must be paid by anti-alignment of the nonlinearity of the same
size.  This is the type-II structure of V28 confined to the stretches
between windows.  What is not available is any lower bound: the
excursions may be absent (the point may be regular in the rescaled
window picture, with the singularity carried elsewhere), and no
Gaussian activity at the window scales is guaranteed at a fixed centre
(Firewall~\ref{fw:c2v42-zero}).  No global regularity claim is made.
\end{assessment}

% =============================================================================
\section{Integrated V42 conclusion and next acquisition order}
\label{sec:c2v42-conclusion}
% =============================================================================

\begin{theorem}[What V42 proves]
\label{thm:c2v42-conclusion}
Relative to the first cycle and V1--V41, modulo (E2): the forward
enstrophy windows \eqref{eq:c2v42-window}, the compactness and explicit
ledger bounds \eqref{eq:c2v42-ledgers} on them, and the bounded signed
ledger \eqref{eq:c2v42-between} across the intermediate stretches.  The
full Navier--Stokes stability/global-regularity theorem remains
unproved.
\end{theorem}

\begin{proof}
The cited results.
\end{proof}

\begin{programme}[V43 acquisition order]
\label{prog:c2v43}
\begin{enumerate}[label=\textup{(AC\arabic*)},leftmargin=4.8em]
\item Enstrophy-following centres.  On a type-I window, the rescaled
      enstrophy \(Y_j\ge c_L^2\nu^{3/2}(-s)^{-1/2}\) is spread over the big
      torus; determine whether the first cycle's count of active
      points at one scale (Lemma V43-few-active, valid under the type-I
      rate) has a window version: at the window scale, the set of
      points where the local enstrophy in \(B_1\) exceeds \(\varepsilon\)
      is covered by at most \(2M/\varepsilon\) balls (a counting that needs
      only the enstrophy bound at that time), and hence that a window
      limit with nonzero enstrophy exists at some centre of the window
      if and only if the enstrophy is not spread below every threshold.
\item Record the resulting trichotomy for the intermediate class at
      window scales: concentrated enstrophy (nonzero window limits at
      some centres), spread enstrophy (all window limits zero), and
      mixed.
\item The next compulsory companion audit is V45.
\end{enumerate}
\end{programme}

\begin{assessment}[Position after V42]
\label{ass:c2v42-position}
The intermediate class has type-I windows with explicit bounds and
bounded signed ledgers across the stretches between them.  No full
stability or global-regularity claim is made.
\end{assessment}

% =============================================================================
\section*{Version V42 (second cycle) append-only record}
\addcontentsline{toc}{section}{Version V42 (second cycle) append-only record}
% =============================================================================

The complete V41 source of the second cycle was retained before this
continuation.  Its SHA--256 digest was
\begin{center}\scriptsize\ttfamily
ef30fd70b900b16fa0d50e36048bfa0249ff47263d4803a62422df63763bae40.
\end{center}
V42 proved the type-I windows of the intermediate class, the
compactness and explicit ledger bounds on them, and the bounded signed
ledger across the intermediate stretches.  No full regularity claim is
made.

% =============================================================================
\section{V43: enstrophy-following centres on type-I windows and the window
trichotomy of the intermediate class}
\label{sec:c2v43-entry}
% =============================================================================

Items (AC1) and (AC2) of Programme~\ref{prog:c2v43} are carried out.  On
a type-I window of the intermediate class, the points at which the
window-integrated local dissipation exceeds a threshold are few: any
set of such points separated by twice the unit scale has at most an
explicit number of elements, by the enstrophy bound alone.  Along a
subsequence, the rescalings at such centres converge on the window to
strong solutions with dissipation at least the threshold on the unit
cylinder, hence nonzero; if for every threshold no such centre exists
for large \(j\), every window limit at every centre map vanishes and
the window dissipation, bounded below by Leray's rate, is spread over a
number of unit cells tending to infinity.  This gives the window
trichotomy of the intermediate class: concentrated, spread, or mixed
along subsequences.  No global regularity claim is made.  The next
compulsory companion audit is V45.

% =============================================================================
\section{Few active centres on a window}
\label{sec:c2v43-few}
% =============================================================================

Let \(u\) be in the intermediate class with constant \(M\), \(\lambda_j\),
\(U_j\) as in Definition~\ref{def:c2v42-intermediate}, and
\(J_M:=[-1,-1+\ell_M]\) the window of Proposition~\ref{prop:c2v42-window},
of length \(\ell_M=3\nu^3/(8C_AM^2)\).  For \(x\in\Torus^3\) put
\begin{equation}
 e_j(x):=\iint_{B_1(x/\lambda_j\text{-centred})\times J_M}|\nabla U_j|^2
 =\frac1{\lambda_j}\iint_{B_{\lambda_j}(x)\times(t_j,\,t_j+\ell_M\lambda_j^2)}|\nabla u|^2 ,
 \label{eq:c2v43-ej}
\end{equation}
the window-integrated local dissipation at \(x\) in scale-invariant units.

\begin{lemma}[Few active centres]
\label{lem:c2v43-few}
For every \(\varepsilon>0\) and every \(j\), any set of points
\(x_1,\dots,x_m\) with \(|x_a-x_b|\ge2\lambda_j\) for \(a\ne b\) and
\(e_j(x_a)\ge\varepsilon\) has \(m\le2M\ell_M/\varepsilon\).  Hence
\(\{x:e_j(x)\ge\varepsilon\}\) is covered by at most \(2M\ell_M/\varepsilon\) balls of
radius \(2\lambda_j\).
\end{lemma}

\begin{proof}
The cylinders \(B_{\lambda_j}(x_a)\times(t_j,t_j+\ell_M\lambda_j^2)\) are disjoint,
each carries at least \(\varepsilon\lambda_j\), and their union carries at most
the total window dissipation \(\lambda_j\int_{J_M}Y_j\le2M\ell_M\lambda_j\) by
\eqref{eq:c2v42-window}.
\end{proof}

\begin{theorem}[Window limits at active centres are nonzero]
\label{thm:c2v43-active}
Let \(x_j\) satisfy \(e_j(x_j)\ge\varepsilon\) for all \(j\) (an active centre
map at level \(\varepsilon\)).  Then along a subsequence the rescalings
\(U^{x}_j:=\lambda_ju(x_j+\lambda_j\,\cdot,T_*+\lambda_j^2\,\cdot)\) converge on
\(B_R\times J_M\), for every \(R\), strongly in \(L^2(H^1)\) to a strong
solution \(U'\) with
\begin{equation}
 \iint_{B_1\times J_M}|\nabla U'|^2\ \ge\ \varepsilon,\qquad
 \sup_{J_M}\|\nabla U'\|_{L^2(\R^3)}^2\le2M ,
 \label{eq:c2v43-nonzero}
\end{equation}
so \(U'\ne0\).  The number of pairwise \(2\lambda_j\)-separated active centre
maps at level \(\varepsilon\) is at most \(2M\ell_M/\varepsilon\).
\end{theorem}

\begin{proof}
Theorem~\ref{thm:c2v42-compactness} at the moving centre (its proof is
centre-independent), the strong \(L^2(H^1)\) convergence on
\(B_1\times J_M\) passing the dissipation lower bound to the limit, and
Lemma~\ref{lem:c2v43-few}.
\end{proof}

\begin{theorem}[Spread windows]
\label{thm:c2v43-spread}
Suppose that for every \(\varepsilon>0\), \(\sup_xe_j(x)<\varepsilon\) for all
large \(j\) (no active centre map at any level).  Then every window
limit at every centre map is zero, while the total window dissipation
satisfies
\begin{equation}
 \lambda_j\int_{J_M}Y_j\ \ge\ \lambda_jc_L^2\nu^{3/2}\int_{J_M}(-s)^{-1/2}\dd s\ \ge\ c_L^2\nu^{3/2}\ell_M\,\lambda_j ,
 \label{eq:c2v43-spread}
\end{equation}
so that the dissipation \(c_L^2\nu^{3/2}\ell_M\lambda_j\) of the window is
spread over at least \(c_L^2\nu^{3/2}\ell_M/\varepsilon\) unit cells for every
\(\varepsilon\): the number of cells carrying it tends to infinity.
\end{theorem}

\begin{proof}
If all window limits at a centre map were nonzero, the strong
\(L^2(H^1)\) convergence would give \(e_j\ge\varepsilon\) at that centre for
some \(\varepsilon\); the lower bound is Leray's rate in rescaled units on
\(J_M\subset[-1,0)\); the counting is \eqref{eq:c2v43-spread} divided by the
cell threshold.
\end{proof}

\begin{corollary}[Window trichotomy of the intermediate class]
\label{cor:c2v43-trichotomy}
Along the window scales \(\lambda_j\) of a solution of the intermediate
class, after passing to subsequences, exactly one of the following
holds for each subsequence.
\begin{enumerate}[label=\textup{(\Roman*)},leftmargin=3.4em]
\item \emph{Concentrated windows.}  There are \(\varepsilon>0\) and an active
      centre map at level \(\varepsilon\); the window limits at it are
      nonzero strong solutions with \eqref{eq:c2v43-nonzero}, and at
      most \(2M\ell_M/\varepsilon\) separated centre maps are active.
\item \emph{Spread windows.}  No level is active for large \(j\); all
      window limits vanish and the window dissipation is spread over a
      diverging number of unit cells.
\end{enumerate}
The full sequence may alternate between the two along different
subsequences (mixed).
\end{corollary}

\begin{proof}
Theorems~\ref{thm:c2v43-active}, \ref{thm:c2v43-spread}.
\end{proof}

\begin{assessment}[Items (AC1)--(AC2)]
\label{ass:c2v43-meaning}
On a type-I window the intermediate class looks, at active centres,
exactly like the type-I case on a bounded window of log-scale: strong
solutions with explicit bounds and dissipation bounded below, to which
the second cycle's identities apply with all quantities bounded; the
count of active centres is explicit.  Off the windows and at spread
windows nothing of the compact theory survives.  The trichotomy is the
intermediate-class form of the compact-dissipation alternative of
V24: concentrated windows are compact dissipation at the window
scales, spread windows are its failure.  A coherent active centre map
across consecutive windows would allow the ledgers to chain (as in
V32), which is not established.  No global regularity claim is made.
\end{assessment}

% =============================================================================
\section{Integrated V43 conclusion and next acquisition order}
\label{sec:c2v43-conclusion}
% =============================================================================

\begin{theorem}[What V43 proves]
\label{thm:c2v43-conclusion}
Relative to the first cycle and V1--V42, modulo (E2): the counting
lemma, the nonzero window limits at active centres
\eqref{eq:c2v43-nonzero}, the spread-window statement
\eqref{eq:c2v43-spread}, and the window trichotomy.  The full
Navier--Stokes stability/global-regularity theorem remains unproved.
\end{theorem}

\begin{proof}
The cited results.
\end{proof}

\begin{programme}[V44 acquisition order]
\label{prog:c2v44}
\begin{enumerate}[label=\textup{(AC\arabic*)},leftmargin=4.8em]
\item Ledgers at active centres across the window.  For an active
      centre map, apply the first and second identities on the window
      at the moving centre, with all quantities bounded by explicit
      functions of \(M\), and record the exact window balances (the
      change of \(\Phi\) and \(\mathcal L\) across the window equals the
      window integrals of the identities), and what the dissipation
      lower bound \(\varepsilon\) forces on the flux over the window (an
      inward-flux statement on a single window of explicit length).
\item Prepare V45: the consolidation of V41--V44 as the fifth block
      (modewise structure and the intermediate class) and the position
      after forty-five versions.
\item The next compulsory companion audit is V45.
\end{enumerate}
\end{programme}

\begin{assessment}[Position after V43]
\label{ass:c2v43-position}
The intermediate class has, on its type-I windows, either concentrated
dissipation at few centres with nonzero strong limits, or spread
dissipation with vanishing limits.  No full stability or
global-regularity claim is made.
\end{assessment}

% =============================================================================
\section*{Version V43 (second cycle) append-only record}
\addcontentsline{toc}{section}{Version V43 (second cycle) append-only record}
% =============================================================================

The complete V42 source of the second cycle was retained before this
continuation.  Its SHA--256 digest was
\begin{center}\scriptsize\ttfamily
00a05694c24179637ab75e9535e8e8aabdb8ddc37846c1ca6d714fe34b52d85f.
\end{center}
V43 proved the counting of active centres on type-I windows, the
nonzero window limits at active centres, and the window trichotomy of
the intermediate class.  No full regularity claim is made.

% =============================================================================
\section{V44: the Gaussian ledgers on a single type-I window}
\label{sec:c2v44-entry}
% =============================================================================

Items (AC1) and (AC2) of Programme~\ref{prog:c2v44} are carried out.  At
an active centre of a type-I window of the intermediate class, the two
Gaussian identities are integrated across the window: the change of the
weighted energy and of the Dirichlet functional across the window is
bounded by explicit functions of the sequence constant, the window
dissipation term is bounded below by the activity threshold times an
explicit Gaussian factor, and the window flux is therefore bounded
above by an explicit expression; the flux is forced inward on the
window exactly when the activity threshold exceeds an explicit multiple
of the sequence constant, which, since the threshold is itself bounded
by the total window dissipation, can happen only for sequence constants
below an explicit value.  The second identity gives the bounded signed
ledger on the window.  The consolidation for V45 is prepared.  No
global regularity claim is made.  The next version is the compulsory
companion audit.

% =============================================================================
\section{Window balances at an active centre}
\label{sec:c2v44-balances}
% =============================================================================

Let \(u\) be in the intermediate class with constant \(M\), \(x_j\) an
active centre map at level \(\varepsilon\) (Theorem~\ref{thm:c2v43-active}),
\(J_M=[-1,-1+\ell_M]\) the window, and, at the centre \(x_j\) and the
scales \(\lambda\in[a_j,b_j]:=[\lambda_j(1-\ell_M)^{1/2},\lambda_j]\) covering the window,
\(\Phi_j(\lambda):=\Phi_{u,x_j}(\lambda)\), \(D_j\), \(\mathcal J_j\), \(\mathcal L_j\) the
periodized Gaussian quantities.  By \eqref{eq:c2v42-ledgers},
\(\Phi_j\le\Phi_M:=2C_S^2M\|G\|_{L^{3/2}}\) and \(\mathcal L_j\le\mathcal L_M\) on
\([a_j,b_j]\).

\begin{theorem}[Window dissipation and window flux]
\label{thm:c2v44-window}
For every \(j\),
\begin{equation}
 \int_{a_j}^{b_j}D_j(\lambda)\,\dd\lambda\ \ge\ 2\nu\,e^{-1/(4\nu(1-\ell_M))}\,\varepsilon ,
 \qquad
 \int_{a_j}^{b_j}\mathcal J_j(\lambda)\,\dd\lambda\ \le\ \Phi_M-2\nu\,e^{-1/(4\nu(1-\ell_M))}\,\varepsilon .
 \label{eq:c2v44-window}
\end{equation}
Hence the window flux is inward whenever
\(\varepsilon>\varepsilon_M:=\Phi_Me^{1/(4\nu(1-\ell_M))}/(2\nu)\); since always
\(\varepsilon\le2M\ell_M\), this is possible only if
\(2M\ell_M>\varepsilon_M\), that is, only for
\begin{equation}
 M^2\ <\ M_*^2:=\frac{3\nu^4\,e^{-1/(4\nu(1-\ell_M))}}{8C_A\,C_S^2\|G\|_{L^{3/2}}} .
 \label{eq:c2v44-Mstar}
\end{equation}
\end{theorem}

\begin{proof}
In the rescaled variables \(U_j\) at the centre \(x_j\), with
\(\lambda=\lambda_j(-s)^{1/2}\): \(D_j(\lambda)=4\nu\lambda_j^{-1}\int|\nabla U_j(s,y)|^2e^{-|y|^2/(4\nu(-s))}\dd y\)
(\(\nabla_yU_j=\lambda_j^2\nabla u\), \(\dd x=\lambda_j^3\dd y\), and the periodized
weight is the whole-space Gaussian of the periodic extension), and
\(\dd\lambda=\lambda_j\,\dd(-s)^{1/2}\), \(|\dd\lambda|\ge\frac{\lambda_j}2|\dd s|\) on \(J_M\)
(\(-s\le1\)).  On \(B_1\) and \(s\in J_M\), \(e^{-|y|^2/(4\nu(-s))}\ge e^{-1/(4\nu(1-\ell_M))}\).
Hence \(\int_{a_j}^{b_j}D_j\,\dd\lambda\ge4\nu\lambda_j^{-1}\cdot\frac{\lambda_j}2\,e^{-1/(4\nu(1-\ell_M))}\iint_{B_1\times J_M}|\nabla U_j|^2\ge2\nu e^{-1/(4\nu(1-\ell_M))}\varepsilon\)
by the activity of the centre.  The flux bound is the first identity
integrated over \([a_j,b_j]\): \(\int\mathcal J_j=\Phi_j(b_j)-\Phi_j(a_j)-\int D_j-\int2\Phi_j/\lambda\le\Phi_M-\int D_j\).
The threshold comparison uses \(\varepsilon\le\lambda_j^{-1}\iint_{\Torus^3\times J}|\nabla u|^2\le2M\ell_M\)
(Lemma~\ref{lem:c2v43-few} with one centre) and
\(\ell_M=3\nu^3/(8C_AM^2)\): \(2M\ell_M>\varepsilon_M\) is
\(\frac{3\nu^3}{4C_AM}>\frac{2C_S^2M\|G\|_{3/2}e^{1/(4\nu(1-\ell_M))}}{2\nu}\),
which rearranges to \eqref{eq:c2v44-Mstar}.
\end{proof}

\begin{theorem}[Bounded signed ledger on the window]
\label{thm:c2v44-second}
For every \(j\), on the window at the active centre,
\begin{equation}
 \Bigl|\int_{a_j}^{b_j}\bigl(\|Y_u\|_G^2+\langle N(V_u),Y_u\rangle_G\bigr)\frac{\dd\lambda}\lambda\Bigr|\ \le\ \frac12\mathcal L_M ,
 \label{eq:c2v44-second}
\end{equation}
and \(\|Y_u\|_G\), \(\|N(V_u)\|_G\) are square-integrable over the window
with bounds depending only on \(M\) and \(\nu\): with \(V=U_j\) on the
window, \(\|(V\cdot\nabla)V\|_{L^2(G)}\le\|V\|_6\|\nabla V\|_3\le C\,Y_j^{3/4}\|\Delta V\|_2^{1/2}\),
\(\|\nabla P_V\|_{L^2}\le C\|(V\cdot\nabla)V\|_2\), and \(\|Y\|_G\le\|L_-V\|_G+\|N\|_G\)
with \(\|L_-V\|_G\le\nu\|\Delta V\|_2+\frac12\||z|\nabla V\|_G+\frac12\|V\|_G\), all
in \(L^2(J_M)\) by \eqref{eq:c2v42-window} and the \(H^2\) bound.
\end{theorem}

\begin{proof}
The rate-free ledger \eqref{eq:c2v14-ledger} with \(0\le\mathcal L_j\le\mathcal L_M\)
at both ends; the square-integrability from the window bounds of
Proposition~\ref{prop:c2v42-window} and the Gaussian Hardy inequality.
\end{proof}

\begin{assessment}[Items (AC1)--(AC2): a type-I window is short]
\label{ass:c2v44-short}
A type-I window has log-length of order \(M^{-2}\), while the Gaussian
inward flux needs, by the first cycle's mechanism, a log-length of
order one to accumulate a sign against the bounded change of the
weighted energy.  The window flux is therefore signed only when the
sequence constant is small, \(M<M_*\) with \(M_*\) explicit in
\(\nu\) and the absolute constants, that is, when the type-I windows are
long; for larger \(M\) the window is too short and the ledger gives
only the bounded signed statements.  In the type-I case proper (a
rate), every scale is in a window and the windows chain, which is how
the first cycle obtained the averaged inward flux; the intermediate
class breaks the chain between windows, and the second cycle's
statements on the stretches (Theorem~\ref{thm:c2v42-between}) are what
replaces it.  No global regularity claim is made.
\end{assessment}

% =============================================================================
\section{Integrated V44 conclusion and next acquisition order}
\label{sec:c2v44-conclusion}
% =============================================================================

\begin{theorem}[What V44 proves]
\label{thm:c2v44-conclusion}
Relative to the first cycle and V1--V43, modulo (E2): the window
dissipation and flux bounds \eqref{eq:c2v44-window} with the threshold
\eqref{eq:c2v44-Mstar}, and the bounded signed ledger on the window
\eqref{eq:c2v44-second}.  The full Navier--Stokes
stability/global-regularity theorem remains unproved.
\end{theorem}

\begin{proof}
Theorems~\ref{thm:c2v44-window}, \ref{thm:c2v44-second}.
\end{proof}

\begin{programme}[V45 acquisition order, with compulsory companion audit]
\label{prog:c2v45}
\begin{enumerate}[label=\textup{(AC\arabic*)},leftmargin=4.8em]
\item Perform the compulsory review of the current SM and YM notes;
      record digests; import nothing numerical.
\item Consolidate V41--V44 as the fifth block: the modewise structure
      (flux bounds, mean spectrum, linear-mode law) and the
      intermediate class (type-I windows, window trichotomy, window
      ledgers).
\item Record the position after forty-five versions and the direction
      for V46--V50.
\item The next compulsory companion audit after V45 is V50.
\end{enumerate}
\end{programme}

\begin{assessment}[Position after V44]
\label{ass:c2v44-position}
The Gaussian ledgers on a single type-I window are explicit, and the
window flux is signed only for small sequence constants.  No full
stability or global-regularity claim is made.
\end{assessment}

% =============================================================================
\section*{Version V44 (second cycle) append-only record}
\addcontentsline{toc}{section}{Version V44 (second cycle) append-only record}
% =============================================================================

The complete V43 source of the second cycle was retained before this
continuation.  Its SHA--256 digest was
\begin{center}\scriptsize\ttfamily
4c34d19c9557ec3801c897e3856a13708ac8c11f3a66cc7c75e720441b61eb10.
\end{center}
V44 proved the window dissipation and flux bounds at an active centre
of a type-I window and the bounded signed ledger on the window.  No
full regularity claim is made.

% =============================================================================
\section{V45: compulsory companion audit, the fifth block consolidated, and
the position after forty-five versions}
\label{sec:c2v45-entry}
% =============================================================================

This is the forty-fifth version of the second cycle.  The compulsory
review of the two companion programmes is performed first.  V41--V44
are then consolidated as the fifth block: the modewise structure of the
identities and the intermediate class.  The position after forty-five
versions and the direction for V46--V50 are recorded.  No global
regularity claim is made.  The next compulsory companion audit is V50.

% =============================================================================
\section{Compulsory V45 review of the current companion notes}
\label{sec:c2v45-companion-audit}
% =============================================================================

The current companion files in \texttt{latex\_notes} were inspected
read-only.  The frozen checkpoints are
\begin{align}
 &\texttt{Renewal\_SM\_Einstein\_Continuum\_Research\_Notes\_V99.tex},
 \notag\\[-2pt]
 &\qquad
 \texttt{53cbaa0772dba73ae19618614c0d47d98197436759668ef6743d26570c74b28e},
 \label{eq:c2v45-SM-checkpoint}\\
 &\texttt{Renewal\_Yang\_Mills\_Mass\_Gap\_Research\_Notes\_V75.tex},
 \notag\\[-2pt]
 &\qquad
 \texttt{5424cc6c7c96795590ac0ce8ae546f22f9670ba417d57cef422d2b2523fc8f15}.
 \label{eq:c2v45-YM-checkpoint}
\end{align}
A YM file named \(\ldots\)\texttt{V74.tex} with the same digest as
\eqref{eq:c2v45-YM-checkpoint} was also present; the V75 name is
recorded.

\emph{SM V98--V99.}  The SM programme performed a Higgs reduction with
electromagnetic matching (V98), then proved exact chirality
intertwiners for its lattice Dirac operator with a contraction bound,
assembled four chiral corners exactly into a massive operator, proved a
massive floor for every nonzero Yukawa singular value, closed a
constant-Higgs normalization row, proved a variable radial-Higgs
perturbation bound, and computed the full defect of an unimproved mass
insertion (V99).

\emph{YM V73--V74.}  The YM programme proved effective centre
hyperedges, then defined a centre-core incidence excess, proved an
exact fixed-core thermal degree with two extra thermal powers per
incidence cycle, identified an acyclic Bell no-go for relaxed
enumeration, proved a reserve-preserving effective-hyperedge bound,
closed the complete centre-core star, and isolated canonical non-star
hypertree pruning as its next step.

\begin{theorem}[V45 companion-audit decision]
\label{thm:c2v45-companion-decision}
No SM V98--V99 intertwiner, floor, or normalization constant and no YM
V73--V74 thermal-degree or hyperedge bound is imported.  Neither bears
on the fifth block or on the gates of Theorem~\ref{thm:c2v15-gates}.
\end{theorem}

\begin{proof}
The SM results concern lattice Dirac operators and Yukawa floors; the
YM results concern hypergraph resummations of a lattice clock.  None
is a Navier--Stokes statement.
\end{proof}

% =============================================================================
\section{The fifth block, consolidated}
\label{sec:c2v45-consolidated}
% =============================================================================

\begin{theorem}[Fifth block, V40--V44]
\label{thm:c2v45-consolidated}
\begin{enumerate}[label=\textup{(\arabic*)},leftmargin=3.2em]
\item \emph{Modewise structure} (type I, homogeneous, (E1)--(E9)).  With
      \(V_k=\Pi_kV\), \(a_k=\|V_k\|_G^2\), \(n_k=\langle V_k,N(V)\rangle_G\),
      \(\mu_k=\frac12+\frac k2\): \(\frac{\dd a_k}{\dd\sigma}=-2\mu_ka_k-2n_k\),
      the first and second identities are \(\sum_kb_k\) and
      \(\sum_k\mu_kb_k\) with \(b_k=\mu_ka_k+n_k\), and against every
      invariant measure \(\langle n_k\rangle=-\mu_k\langle a_k\rangle\) for every
      \(k\); \(|n_k|\le a_k^{1/2}[\kappa(k/2)^{1/2}\varphi^{1/2}+\kappa(d+3\varphi)^{1/2}+\frac{K_P}4(3+2k)^{1/2}]\),
      hence \(\langle a_k\rangle\le\mu_k^{-2}[\ldots]^2=O(1/k)\) explicitly; the
      constant and linear modes are the Gaussian mean and the dipole
      matrix, and the linear-mode law involves the third moment of the
      profile and the second moment of the pressure, outside the moment
      ledger.
\item \emph{Intermediate class} (no rate, (E2)).  Type-I windows of
      log-length \(\ell_M=3\nu^3/(8C_AM^2)\) with enstrophy at most \(2M\),
      \(H^2\) bounds, compactness to strong solutions, and explicit
      Gaussian ledger bounds \eqref{eq:c2v42-ledgers}; the signed
      ledger across the stretches between windows is bounded by
      \(\frac12\mathcal L_M\); on a window, at most \(2M\ell_M/\varepsilon\)
      separated centres are active at level \(\varepsilon\), active
      centres have nonzero strong window limits, spread windows have
      all limits zero with dissipation over a diverging number of
      cells; the window flux at an active centre is at most
      \(\Phi_M-2\nu e^{-1/(4\nu(1-\ell_M))}\varepsilon\), signed only for
      \(M<M_*\).
\end{enumerate}
\end{theorem}

\begin{proof}
(1) Theorem~\ref{thm:c2v40-modewise}, Proposition~\ref{prop:c2v40-low-modes},
Theorems~\ref{thm:c2v41-bound}, \ref{thm:c2v41-spectrum},
Proposition~\ref{prop:c2v41-dipole}.  (2) Proposition~\ref{prop:c2v42-window},
Theorems~\ref{thm:c2v42-compactness}, \ref{thm:c2v42-between},
Lemma~\ref{lem:c2v43-few}, Theorems~\ref{thm:c2v43-active},
\ref{thm:c2v43-spread}, \ref{thm:c2v44-window}, \ref{thm:c2v44-second}.
\end{proof}

% =============================================================================
\section{Position after forty-five versions and direction}
\label{sec:c2v45-position}
% =============================================================================

\begin{assessment}[Position after forty-five versions of the second cycle]
\label{ass:c2v45-position-forty-five}
The second cycle now covers the three gates of the first cycle with
its own tools: (G1) homogeneous points, with exact signed statistics,
the constant-chasing constraint, and the modewise spectrum; (G1\('\))
lacunary points, with the statistical dichotomy and the compact
dissipation alternative; (G2) the intermediate class, with type-I
windows and the window trichotomy; (G3) strong type II, with the two
rate-free ledgers, frequency escape, and forced high-mode escape.  In
each case the series has exact balances and explicit bounds, and in no
case a contradiction.  The two obstructions that recur are:
non-explicit compactness constants (the lower side of every gate) and
spatial spreading of dissipation (the failure of compactness at fixed
or coherent centres).  For V46--V50 the direction is the second
obstruction: whether spreading of the dissipation over a diverging
number of parabolic cells at a scale is compatible with the enstrophy
equation between that scale and the next, using the first cycle's
spike anatomy (the backward Leray cone) and the second cycle's window
ledgers; specifically, whether dissipation spread over \(n\) cells at
scale \(\lambda\), each carrying \(o(\lambda)\), forces enstrophy of order
\(n\cdot o(\lambda^{-1})\) at that time, contradicting a type-I window
bound at the next window when the spreading is faster than the window
recurrence.  No global regularity claim is made.
\end{assessment}

% =============================================================================
\section{Integrated V45 conclusion and next acquisition order}
\label{sec:c2v45-conclusion}
% =============================================================================

\begin{theorem}[What V45 establishes]
\label{thm:c2v45-conclusion}
Relative to the first cycle and V1--V44: the companion-audit decision
(Theorem~\ref{thm:c2v45-companion-decision}) and the consolidated fifth
block (Theorem~\ref{thm:c2v45-consolidated}).  The full Navier--Stokes
stability/global-regularity theorem remains unproved.
\end{theorem}

\begin{proof}
The cited results.
\end{proof}

\begin{programme}[V46 acquisition order]
\label{prog:c2v46}
\begin{enumerate}[label=\textup{(AC\arabic*)},leftmargin=4.8em]
\item Spreading versus enstrophy.  At a spread window (all active
      levels empty), the window dissipation \(\ge c_L^2\nu^{3/2}\ell_M\lambda_j\)
      lies in cells of unit rescaled size each carrying less than
      \(\varepsilon\); bound the rescaled enstrophy at the window time
      from below in terms of the number of cells and the per-cell
      dissipation using only the parabolic scaling (enstrophy density
      per unit volume over unit time), and compare with the window
      bound \(2M\); determine whether spreading over \(n\ge c/\varepsilon\)
      cells is compatible with \(Y_j\le2M\) (it is: the enstrophy is
      the integral of the dissipation density, and \(n\) cells with
      \(\varepsilon\) each give \(n\varepsilon\), bounded by the total);
      record the exact statement and whether any lower bound on the
      per-cell dissipation follows from the Leray rate applied cell by
      cell (it does not, since Leray's rate is global).
\item Backward cone at a spread window.  Apply the first cycle's
      backward Leray cone (V37) to the window: the enstrophy spike that
      precedes the window must be paid by Leray energy in a backward
      cone; determine whether the cone constrains the number of cells.
\item The next compulsory companion audit is V50.
\end{enumerate}
\end{programme}

\begin{assessment}[Position after V45]
\label{ass:c2v45-position}
The second cycle is consolidated through forty-five versions, with the
fifth block covering the modewise structure and the intermediate class.
No full stability or global-regularity claim is made.
\end{assessment}

% =============================================================================
\section*{Version V45 (second cycle) append-only record}
\addcontentsline{toc}{section}{Version V45 (second cycle) append-only record}
% =============================================================================

The complete V44 source of the second cycle was retained before this
continuation.  Its SHA--256 digest was
\begin{center}\scriptsize\ttfamily
17e622b2d5df83e768aa7f1d12fee4bc2ee79d21d9148c37255345dfd194f1e9.
\end{center}
V45 performed the compulsory companion audit, consolidated the fifth
block, and recorded the position after forty-five versions.  No full
regularity claim is made.

% =============================================================================
\section{V46: spreading versus enstrophy, and the backward cone at a spread
window}
\label{sec:c2v46-entry}
% =============================================================================

Items (AC1) and (AC2) of Programme~\ref{prog:c2v46} are carried out.
Spreading of the window dissipation over a diverging number of unit
cells is shown to be compatible with the window enstrophy bound: the
enstrophy is the sum of the cell contributions and the total is fixed,
so spreading lowers the per-cell share without raising the total, and
Leray's rate, being global, gives no per-cell lower bound.  The first
cycle's backward Leray cone is applied between windows and spikes: a
spike of the intermediate class cannot occur within an explicit
parabolic time after a window, which reproduces the forward window
length, and the square roots of the gaps between spikes of unbounded
height are summable, which constrains the spike times in log-scale but
not the spatial distribution of the dissipation.  The spatial spreading
obstruction is therefore not removable by the enstrophy ledger.  No
global regularity claim is made.  The next compulsory companion audit
is V50.

% =============================================================================
\section{Spreading is compatible with the enstrophy bound}
\label{sec:c2v46-spreading}
% =============================================================================

\begin{proposition}[Cell accounting at a window]
\label{prop:c2v46-cells}
Let \(u\) be in the intermediate class with constant \(M\), \(J_M\) the
window, and partition the rescaled torus of side \(2\pi/\lambda_j\) into
unit cubes \(Q_\alpha\) (\(\alpha\) ranging over \((2\pi/\lambda_j)^3\) cells).  Let
\(e_\alpha:=\iint_{Q_\alpha\times J_M}|\nabla U_j|^2\) be the window dissipation
of the cell.  Then
\begin{equation}
 c_L^2\nu^{3/2}\ell_M\ \le\ \sum_\alpha e_\alpha\ \le\ 2M\ell_M ,
 \qquad
 \#\{\alpha:e_\alpha\ge\delta\}\le\frac{2M\ell_M}\delta ,
 \qquad
 \#\{\alpha:e_\alpha>0\}\ \ge\ \frac{c_L^2\nu^{3/2}\ell_M}{\max_\alpha e_\alpha} .
 \label{eq:c2v46-cells}
\end{equation}
At a spread window (\(\max_\alpha e_\alpha\to0\) along the subsequence) the
number of cells carrying the dissipation tends to infinity, while the
enstrophy at every time of the window, \(Y_j(s)=\sum_\alpha\int_{Q_\alpha}|\nabla U_j(s)|^2\),
stays at most \(2M\): the per-cell enstrophy is small and the cells are
many, with a fixed total.  No contradiction with the window bound
arises, and no lower bound on \(\max_\alpha e_\alpha\) follows from Leray's
rate, which bounds only the sum.
\end{proposition}

\begin{proof}
The sum is the window dissipation, bounded by \eqref{eq:c2v42-window}
above and by Leray's rate below (Theorem~\ref{thm:c2v43-spread}); the
counts are the pigeonhole principle.
\end{proof}

\begin{assessment}[Item (AC1)]
\label{ass:c2v46-AC1}
Spreading is the regime in which the type-I window bound is satisfied
trivially: the enstrophy is not concentrated anywhere, so no
rescaling at any centre sees it, and the Gaussian ledgers at every
centre see a profile of vanishing Gaussian energy and dissipation.
The enstrophy identity cannot exclude it because the enstrophy is an
integral, and Leray's rate is a lower bound on the integral only.  A
constraint on spreading must come from a mechanism that couples cells:
the pressure, which is nonlocal, or the transport, which moves
enstrophy between cells.  The series has neither in quantitative form
at a spread window.
\end{assessment}

% =============================================================================
\section{The backward cone between windows and spikes}
\label{sec:c2v46-cone}
% =============================================================================

\begin{proposition}[Spikes keep a parabolic distance from windows]
\label{prop:c2v46-cone}
Let \(u\) be in the intermediate class with constant \(M\), \(t_j\) a window
time with \((T_*-t_j)^{1/2}Y(t_j)\le M\), and \(t_2>t_j\) a later time.  By
the backward cone of the first cycle (Theorem V37-cone-share,
\(Y(t)\ge(\nu^3/(8C_1(t_2-t)))^{1/2}\) for \(t_2-t\ge\nu^3/(4C_1Y(t_2)^2)\)),
\begin{equation}
 t_2-t_j\ \ge\ \frac{\nu^3}{32C_1M^2}\,(T_*-t_j)
 \qquad\text{whenever}\qquad
 Y(t_2)\ \ge\ \frac{4M}{(T_*-t_j)^{1/2}} ,
 \label{eq:c2v46-distance}
\end{equation}
that is, a spike of height four times the window level cannot occur
within the parabolic time \(\frac{\nu^3}{32C_1M^2}\lambda_j^2\) after the window
time; this is the forward window of Proposition~\ref{prop:c2v42-window}
read backward.  Moreover, if \(t'_1<t'_2<\dots\uparrow T_*\) are spike
times with \(Y(t'_k)\to\infty\), then
\(\sum_k(t'_{k+1}-t'_k)^{1/2}<\infty\).
\end{proposition}

\begin{proof}
Apply the cone at \(t=t_j\): \(M(T_*-t_j)^{-1/2}\ge Y(t_j)\ge(\nu^3/(8C_1(t_2-t_j)))^{1/2}\)
when \(t_2-t_j\ge\nu^3/(4C_1Y(t_2)^2)\), which holds under the stated
height condition once \(t_2-t_j\ge\nu^3(T_*-t_j)/(64C_1M^2)\); squaring gives
\eqref{eq:c2v46-distance}.  The summability is the gap-summability statement of Theorem
V37-cone-share of the first cycle,
\(\sum_k(\Delta_k^{1/2}-(\nu^3/(4C_1))^{1/2}Y_{\min}^{-1})_+\le(2C_1/\nu^3)^{1/2}E_*/(2\nu)\),
with \(Y_{\min}\to\infty\).
\end{proof}

\begin{firewall}[Item (AC2): the cone constrains times, not cells]
\label{fw:c2v46-AC2}
The backward cone is a statement about the global enstrophy as a
function of time; it constrains where spikes and windows can sit in
time (spikes at least a parabolic time after windows; gaps between
spikes with summable square roots, which in log-scale allows equally
spaced spikes but excludes spikes with gaps decaying like the inverse
square of their index), and it says nothing about the spatial
distribution of the dissipation at a window.  In particular a spread
window can be followed by a spike after the parabolic delay, with the
spike's enstrophy assembled from the many cells by transport, or
generated in one cell; the series distinguishes neither.  Item (AC2)
is closed with this record.
\end{firewall}

\begin{assessment}[The spreading obstruction is structural]
\label{ass:c2v46-structural}
Both obstructions of the position after forty-five versions are now
seen to be one: the compactness constants are non-explicit because the
per-scale implications are proved by contradiction, and they are
proved by contradiction because a direct argument would have to
exclude spreading, which the enstrophy ledger cannot.  Excluding
spreading requires a quantitative unique-continuation or
pressure-coupling statement across cells at the parabolic scale.
The series records this as the single structural gap of its rate-free
and intermediate analyses, and, for the type-I case, as the source of
the non-explicit lower side of the constant-chasing gate.  No global
regularity claim is made.
\end{assessment}

% =============================================================================
\section{Integrated V46 conclusion and next acquisition order}
\label{sec:c2v46-conclusion}
% =============================================================================

\begin{theorem}[What V46 establishes]
\label{thm:c2v46-conclusion}
Relative to the first cycle and V1--V45: the cell accounting
\eqref{eq:c2v46-cells} at a window and the parabolic distance
\eqref{eq:c2v46-distance} between windows and spikes.  The full
Navier--Stokes stability/global-regularity theorem remains unproved.
\end{theorem}

\begin{proof}
Propositions~\ref{prop:c2v46-cells}, \ref{prop:c2v46-cone}.
\end{proof}

\begin{programme}[V47 acquisition order]
\label{prog:c2v47}
\begin{enumerate}[label=\textup{(AC\arabic*)},leftmargin=4.8em]
\item Pressure coupling across cells.  At a window, the pressure is
      the Riesz transform of \(U\otimes U\); determine whether the local
      pressure in a cell with small dissipation is controlled by the
      dissipation of that cell and its neighbours (a local Calder\'on--Zygmund
      estimate with the far part harmonic), and whether the local
      energy inequality in a cell then bounds the cell's energy growth
      by its own dissipation plus a harmonic-pressure term of the order
      of the global \(L^{3/2}\) pressure norm divided by the distance;
      record whether this closes a cellwise energy budget at a spread
      window.
\item Transport across cells.  Bound the energy flux through the
      boundary of a cell by the cell's \(L^3\) norm, and determine
      whether spreading (small dissipation per cell) forces small
      energy flux per cell, so that cells evolve almost independently
      over the window; record the exact statement.
\item The next compulsory companion audit is V50.
\end{enumerate}
\end{programme}

\begin{assessment}[Position after V46]
\label{ass:c2v46-position}
Spreading of dissipation at a window is compatible with every
enstrophy-ledger statement of the series; excluding it needs a
cross-cell coupling.  No full stability or global-regularity claim is
made.
\end{assessment}

% =============================================================================
\section*{Version V46 (second cycle) append-only record}
\addcontentsline{toc}{section}{Version V46 (second cycle) append-only record}
% =============================================================================

The complete V45 source of the second cycle was retained before this
continuation.  Its SHA--256 digest was
\begin{center}\scriptsize\ttfamily
7b7ea99385c098c9587becba1d82536403d988e9e7477fe7882ea257cc15c791.
\end{center}
V46 recorded the compatibility of spreading with the enstrophy bound
and the backward-cone constraints between windows and spikes.  No full
regularity claim is made.

% =============================================================================
\section{V47: pressure and transport across cells at a type-I window}
\label{sec:c2v47-entry}
% =============================================================================

Items (AC1) and (AC2) of Programme~\ref{prog:c2v47} are carried out.  At
a type-I window the far part of the pressure on any unit cell is
bounded in the supremum norm by an explicit multiple of the sequence
constant, uniformly over the cells, through the dyadic decay of the
Newtonian potential and the Sobolev bound on dyadic balls, while the
local part is bounded by the cubic norm of the neighbouring cells; the
local energy inequality then gives a cellwise energy budget in which
the growth of a cell's energy over the window is controlled by the
cubic norms of the cell and its neighbours.  The budget closes in one
direction only: cells with small energy and small dissipation stay
small; it does not exclude a window at which the dissipation is spread
below every threshold while the energy is concentrated in few cells,
because a cell can hold energy with little dissipation over a window
of the window's length.  The transport flux through a cell boundary is
bounded by the same cubic norms, so spreading of the dissipation alone
does not force small transport.  The cross-cell coupling therefore does
not exclude spreading, and the obstruction stands.  No global
regularity claim is made.  The next compulsory companion audit is V50.

% =============================================================================
\section{The pressure on a cell at a window}
\label{sec:c2v47-pressure}
% =============================================================================

Let \(u\) be in the intermediate class with constant \(M\), \(U=U_j\) the
rescaling at a window scale on the torus of side \(2\pi/\lambda_j\), \(J_M\)
the window, \(Q_\alpha\) the unit cells, \(2Q_\alpha\) the concentric cubes of
side \(2\), \(P\) the rescaled pressure, and \(\chi_\alpha\) a cutoff equal
to \(1\) on \(2Q_\alpha\) and \(0\) outside \(3Q_\alpha\).  Write
\(P=P^{\rm loc}_\alpha+P^{\rm far}_\alpha\) with
\(P^{\rm loc}_\alpha=\partial_i\partial_j(-\Delta)^{-1}(\chi_\alpha U_iU_j)\).

\begin{proposition}[Far pressure is bounded, local pressure is cubic]
\label{prop:c2v47-pressure}
For every cell \(\alpha\) and every \(s\in J_M\):
\begin{enumerate}[label=\textup{(\roman*)},leftmargin=3.2em]
\item \(\|P^{\rm far}_\alpha(s)-c_\alpha(s)\|_{L^\infty(Q_\alpha)}\le C_{\rm far}\,C_S^2\,2M\) for a
      suitable constant \(c_\alpha(s)\), with \(C_{\rm far}\) absolute: the far
      pressure is bounded uniformly over the cells by an explicit
      multiple of the sequence constant.
\item \(\|P^{\rm loc}_\alpha\|_{L^{3/2}(Q_\alpha\times J_M)}\le C_{\rm CZ}\|U\|_{L^3(3Q_\alpha\times J_M)}^2\).
\end{enumerate}
\end{proposition}

\begin{proof}
(i) \(P^{\rm far}_\alpha=\partial_i\partial_j(-\Delta)^{-1}((1-\chi_\alpha)U_iU_j)\) is, on
\(Q_\alpha\), a convolution of \((1-\chi_\alpha)U\otimes U\) with a kernel bounded
by \(C|x-y|^{-3}\) for \(|x-y|\ge1\); decomposing the complement of
\(2Q_\alpha\) into dyadic shells \(2^kQ_\alpha\setminus2^{k-1}Q_\alpha\), \(k\ge2\), the
contribution of the \(k\)-th shell is at most
\(C2^{-3k}\int_{2^kQ_\alpha}|U|^2\le C2^{-3k}(2^k)^2\|U\|_{L^6}^2\le C2^{-k}C_S^2Y_j\)
by H\"older on the cube of side \(2^k\) and the Sobolev inequality, and
the sum over \(k\) converges; the subtraction of a constant removes the
mean of the kernel's zero-order part.  \(Y_j\le2M\) on the window.  (ii)
is the Calder\'on--Zygmund bound for the local part.
\end{proof}

\begin{theorem}[Cellwise energy budget at a window]
\label{thm:c2v47-budget}
Let \(E_\alpha(s):=\int_{Q_\alpha}|U(s)|^2\), \(e_\alpha\) the window dissipation of
the cell, and \(\mathcal C_\alpha:=\|U\|_{L^3(3Q_\alpha\times J_M)}^3\).  Then for
\(s_1<s_2\) in \(J_M\),
\begin{equation}
 E_\alpha(s_2)\ \le\ 4E'_\alpha(s_1)+C\Bigl(\sum_{\beta\sim\alpha}\bigl(E'_\beta{}^{\rm sup}+\mathcal C_\beta+\mathcal C_\beta^{2/3}\,C_SM^{1/2}\bigr)\Bigr)+C\,C_{\rm far}C_S^2M\sum_{\beta\sim\alpha}\mathcal C_\beta^{1/3}\ell_M^{2/3},
 \label{eq:c2v47-budget}
\end{equation}
where \(E'_\beta\) denotes the energy on \(2Q_\beta\), \(E'^{\rm sup}_\beta\) its
supremum over the window, \(\beta\sim\alpha\) ranges over the cell and its
neighbours, and \(C\) is absolute.  In particular, if the cell and its
neighbours have small energy and small dissipation at the start of the
window, they stay small over the window (the cubic norms being
controlled by \(\mathcal C_\beta\le C(E'^{\rm sup}_\beta)^{3/4}e_\beta^{3/4}\ell_M^{1/4}+C(E'^{\rm sup}_\beta)^{3/2}\ell_M\)).
\end{theorem}

\begin{proof}
The local energy inequality on \(2Q_\alpha\) with a cutoff \(\phi_\alpha\) equal
to \(1\) on \(Q_\alpha\): the energy at \(s_2\) is bounded by the energy at
\(s_1\) on \(2Q_\alpha\), plus \(\nu\iint|U|^2|\Delta\phi_\alpha|\) (bounded by the
energy on \(2Q_\alpha\) over the window), plus \(\iint(|U|^2+2P)U\cdot\nabla\phi_\alpha\);
the cubic term is bounded by \(\mathcal C_\alpha\) (H\"older on the cell),
the local-pressure term by \(\|P^{\rm loc}\|_{3/2}\|U\|_3\le C\mathcal C_\alpha\)
using Proposition~\ref{prop:c2v47-pressure}(ii) with the neighbours,
and the far-pressure term, after subtracting the constant \(c_\alpha\)
(which integrates to zero against \(U\cdot\nabla\phi_\alpha\) by
\(\operatorname{div}U=0\)), by \(\|P^{\rm far}-c\|_\infty\|U\|_{L^1(2Q_\alpha\times J_M)}\le C_{\rm far}C_S^2\,2M\,\mathcal C_\alpha^{1/3}|2Q_\alpha\times J_M|^{2/3}\).
The cubic bound is the interpolation of \(L^3\) between \(L^\infty L^2\)
and \(L^2H^1\) on a unit cube over a window of length \(\ell_M\).  The
factors are collected into the displayed form with crude constants.
\end{proof}

\begin{firewall}[Item (AC1): the budget does not exclude spreading]
\label{fw:c2v47-AC1}
The budget bounds the growth of a cell's energy by cubic norms that are
small when both the energy and the dissipation of the neighbourhood
are small.  A spread window has small dissipation in every cell, but
the energy need not be small: \(E_\alpha\le C_S^2\cdot2M\) is all the window
bound gives, and a cell may carry energy of that order with dissipation
below any threshold over a window of length \(\ell_M\) (a nearly inviscid
cell), which the budget permits since the cubic norm is then of order
\(M^{3/2}\ell_M\), not small.  Hence spreading of the dissipation with
concentration of the energy in a bounded number of cells is compatible
with the cellwise budget, and spreading of both is compatible too.
The budget closes only ``small stays small''.
\end{firewall}

% =============================================================================
\section{Transport across a cell boundary}
\label{sec:c2v47-transport}
% =============================================================================

\begin{proposition}[Boundary flux]
\label{prop:c2v47-flux}
The energy flux through the boundary layer of a cell over the window,
\(\iint(|U|^2+2(P-c_\alpha))U\cdot\nabla\phi_\alpha\), is bounded by
\(C(\mathcal C_\alpha+\mathcal C_\alpha^{1/3}C_{\rm far}C_S^2M\ell_M^{2/3})\), the same
quantities as in \eqref{eq:c2v47-budget}; it is small if and only if
the cell's cubic norm is small, which requires small energy, not only
small dissipation.
\end{proposition}

\begin{proof}
Contained in the proof of Theorem~\ref{thm:c2v47-budget}.
\end{proof}

\begin{assessment}[Item (AC2) and the obstruction]
\label{ass:c2v47-AC2}
Cross-cell coupling at a window is explicit: the far pressure is
uniformly bounded by the sequence constant and the local pressure and
the transport are cubic in the neighbourhood.  But the coupling
constrains a cell only through its energy, and the window bound
controls energy only up to \(C_S^2\cdot2M\) per cell; with that much
energy a cell can be nearly inviscid over the short window.  So the
enstrophy ledger, the pressure, and the transport together do not
exclude a spread window, and the structural obstruction of V46 is
confirmed at the level of the cellwise budget: excluding spreading
would require either a lower bound on the per-cell dissipation of an
energetic cell over the window (a local version of Leray's rate,
false in general), or a rate connecting consecutive windows.  No
global regularity claim is made.
\end{assessment}

% =============================================================================
\section{Integrated V47 conclusion and next acquisition order}
\label{sec:c2v47-conclusion}
% =============================================================================

\begin{theorem}[What V47 proves]
\label{thm:c2v47-conclusion}
Relative to the first cycle and V1--V46: the pressure bounds on a cell
at a window (Proposition~\ref{prop:c2v47-pressure}), the cellwise energy
budget \eqref{eq:c2v47-budget}, and the boundary-flux bound.  The full
Navier--Stokes stability/global-regularity theorem remains unproved.
\end{theorem}

\begin{proof}
The cited results.
\end{proof}

\begin{programme}[V48 acquisition order]
\label{prog:c2v48}
\begin{enumerate}[label=\textup{(AC\arabic*)},leftmargin=4.8em]
\item Energetic cells at a spread window.  Under spread dissipation
      with energy concentrated in a bounded number of cells, apply the
      window trichotomy to the energy instead of the dissipation:
      define energy-active centres (cell energy at least \(\varepsilon\))
      and determine whether the rescalings at energy-active centres
      converge to nonzero limits on the window (the energy passes to
      the limit strongly by the compactness of
      Theorem~\ref{thm:c2v42-compactness}), so that a window with
      concentrated energy has nonzero window limits even if its
      dissipation is spread; record the resulting refined trichotomy
      (dissipation-concentrated, energy-concentrated with spread
      dissipation, both spread).
\item Both spread.  Determine what the ledgers say when both energy and
      dissipation are spread at a window: all window limits vanish,
      every Gaussian ledger at every centre vanishes in the limit, and
      the Leray rate is carried by a diverging number of cells; relate
      to the first cycle's \(\tau^{1/3}\) threshold (no spatial energy
      concentration below it).
\item The next compulsory companion audit is V50.
\end{enumerate}
\end{programme}

\begin{assessment}[Position after V47]
\label{ass:c2v47-position}
Cross-cell coupling at a window is explicit and does not exclude
spreading; a spread window with energetic nearly inviscid cells is
compatible with every statement of the series.  No full stability or
global-regularity claim is made.
\end{assessment}

% =============================================================================
\section*{Version V47 (second cycle) append-only record}
\addcontentsline{toc}{section}{Version V47 (second cycle) append-only record}
% =============================================================================

The complete V46 source of the second cycle was retained before this
continuation.  Its SHA--256 digest was
\begin{center}\scriptsize\ttfamily
6cdb3124004c828e7801c1cf94abb92c66bf9c66afa5088ab0610660dceacffe.
\end{center}
V47 proved the pressure bounds and the cellwise energy budget at a
type-I window and recorded that they do not exclude spreading.  No
full regularity claim is made.

% =============================================================================
\section{V48: energy-active centres at a window, and the scale of the
required dissipation}
\label{sec:c2v48-entry}
% =============================================================================

Items (AC1) and (AC2) of Programme~\ref{prog:c2v48} are carried out.  A
centre at which the rescaled energy in the unit ball at the window time
is at least a threshold has nonzero window limits, by the compactness
of the window rescalings; the limit is a strong solution in \(L^6\)
with positive dissipation, and by a compactness argument the window
dissipation at the centre is bounded below by a positive constant in a
ball of radius depending on the threshold and the sequence constant.
Hence energy concentration at a window forces dissipation
concentration at a larger radius, and the refined trichotomy of the
programme collapses to the dissipation trichotomy up to the change of
radius; a purely static Poincar\'e argument gives, explicitly, the
radius at which the enstrophy at the window time must be present.
The wording of a firewall of V47 is made precise accordingly: an
energetic cell can be nearly inviscid at the unit scale only if its
required dissipation lies within the larger radius, spread over the
cells there.  The both-spread case is recorded.  No global regularity
claim is made.  The next compulsory companion audit is V50.

% =============================================================================
\section{Energy-active centres}
\label{sec:c2v48-energy}
% =============================================================================

Let \(u\) be in the intermediate class with constant \(M\), \(\lambda_j\), \(J_M\)
as before, and for a centre \(x\) let \(U^x_j\) be the rescaling at \(x\) and
\(E^x_j:=\int_{B_1}|U^x_j(-1)|^2\) the rescaled energy of the unit ball at
the window time.

\begin{proposition}[Static lower bound on the enstrophy near an energetic ball]
\label{prop:c2v48-static}
Let \(U\) be a divergence-free field on \(\R^3\) (or on the big torus)
with \(\|U\|_{L^6}\le K\) and \(\int_{B_1}|U|^2\ge\varepsilon\).  Then for
\(R\ge R_0(\varepsilon,K):=8C_{\rm P}'K^2/\varepsilon\),
\begin{equation}
 \int_{B_R}|\nabla U|^2\ \ge\ \delta_0(\varepsilon,K,R):=\frac{\varepsilon}{4C_{\rm P}^2R^2} ,
 \label{eq:c2v48-static}
\end{equation}
where \(C_{\rm P}\) is the Poincar\'e constant of the unit ball
(\(\|f-\bar f\|_{L^2(B_R)}\le C_{\rm P}R\|\nabla f\|_{L^2(B_R)}\)) and \(C_{\rm P}'\) an
absolute constant.  At a window, \(K=C_S(2M)^{1/2}\).
\end{proposition}

\begin{proof}
Let \(c\) be the mean of \(U\) over \(B_R\).  H\"older gives
\(|c|^2|B_R|\le\|U\|_{L^2(B_R)}^2\le|B_R|^{2/3}K^2\), so \(|c|^2\le C'K^2/R\).
If \(\int_{B_R}|\nabla U|^2<\delta\), Poincar\'e gives
\(\|U-c\|_{L^2(B_R)}^2\le C_{\rm P}^2R^2\delta\), hence
\(\int_{B_1}|U|^2\le2|c|^2|B_1|+2C_{\rm P}^2R^2\delta\le\frac{8C'K^2}{R}+2C_{\rm P}^2R^2\delta\).
With \(R\ge16C'K^2/\varepsilon\) the first term is at most \(\varepsilon/2\), and
with \(\delta<\varepsilon/(4C_{\rm P}^2R^2)\) the second is less than \(\varepsilon/2\);
the right side is then less than \(\varepsilon\), a contradiction.  Hence
\(\int_{B_R}|\nabla U|^2\ge\varepsilon/(4C_{\rm P}^2R^2)\) for \(R\ge16C'K^2/\varepsilon\), which
is the claim with \(C_{\rm P}':=2C'\).
\end{proof}

\begin{theorem}[Energy-active centres are dissipation-active at a larger radius]
\label{thm:c2v48-energy-active}
Let \(x_j\) be a centre map with \(E^{x_j}_j\ge\varepsilon\) for all \(j\).
\begin{enumerate}[label=\textup{(\roman*)},leftmargin=3.2em]
\item Along a subsequence the window rescalings \(U^{x_j}_j\) converge on
      the window to a strong solution \(U'\) with
      \(\int_{B_1}|U'(-1)|^2\ge\varepsilon\), \(\|U'\|_{L^6}\le C_S(2M)^{1/2}\),
      \(\int_{B_{R_0}}|\nabla U'(-1)|^2\ge\delta_0\), and
      \(\iint_{B_{R_0}\times J_M}|\nabla U'|^2>0\).
\item There are \(\varepsilon'=\varepsilon'(\varepsilon,M,\nu)>0\) and \(R'=R'(\varepsilon,M,\nu)\)
      such that \(\iint_{B_{R'}(x_j\text{-centred})\times J_M}|\nabla U^{x_j}_j|^2\ge\varepsilon'\)
      for all large \(j\): energy-active at the unit scale implies
      dissipation-active at the radius \(R'\).
\end{enumerate}
\end{theorem}

\begin{proof}
(i) Theorem~\ref{thm:c2v42-compactness} at the moving centre gives
convergence in \(L^2(H^1)\) on the window and, at the fixed time \(-1\),
the uniform \(H^1\) bound gives strong \(L^2(B_R)\) convergence by Rellich
along a further subsequence, so the energy lower bound passes to the
limit; the \(L^6\) bound passes by weak lower semicontinuity;
Proposition~\ref{prop:c2v48-static} applies to \(U'(-1)\); and the
enstrophy of the smooth strong solution \(U'\) on \(B_{R_0}\) is
continuous in time, hence positive on a nonempty initial part of the
window.  (ii) If (ii) failed for every \(\varepsilon',R'\), there would be
energy-active centre maps (for a sequence of solutions with the same
\(M\), or the same solution along different subsequences) whose window
dissipation in \(B_{R}\) tends to zero for every \(R\); their limits, by
(i), have positive dissipation on \(B_{R_0}\times J_M\), while strong
\(L^2(H^1)\) convergence gives the limit's dissipation as the limit of
the rescaled dissipations, which tend to zero: a contradiction.  The
constants depend only on \(\varepsilon\), \(M\), \(\nu\) by the uniformity of the
compactness (the class of window rescalings with constant \(M\) is
determined by \(M\) and \(\nu\)).
\end{proof}

\begin{corollary}[The refined trichotomy collapses]
\label{cor:c2v48-collapse}
At a window scale, energy concentration (some centre with
\(E^x_j\ge\varepsilon\)) implies dissipation concentration at the radius
\(R'(\varepsilon,M,\nu)\) with level \(\varepsilon'(\varepsilon,M,\nu)\); hence a spread
window in the sense of Theorem~\ref{thm:c2v43-spread} at all radii (no
dissipation-active centre map at any level and any fixed radius) has
spread energy: \(\sup_xE^x_j\to0\).  The cases of the intermediate class
at window scales are therefore two, concentrated (some centre has
energy and dissipation bounded below, nonzero strong limits) and spread
(energy and dissipation in every fixed ball tend to zero at every
centre), with mixed sequences.
\end{corollary}

\begin{proof}
Theorem~\ref{thm:c2v48-energy-active}(ii) and the definitions.
\end{proof}

\begin{theorem}[Precision of Firewall~\ref{fw:c2v47-AC1}]
\label{thm:c2v48-precision}
The scenario of Firewall~\ref{fw:c2v47-AC1}, a cell with energy of
order \(M\) and window dissipation below every threshold, is possible
only in the following form: the unit cell has small window dissipation
while the ball of radius \(R'(\varepsilon,M,\nu)\) around it carries window
dissipation at least \(\varepsilon'(\varepsilon,M,\nu)\), spread over the
cells of that ball.  A nearly inviscid energetic cell at the unit scale
is therefore compatible with the series only if its dissipation is
displaced to the surrounding \(O(R'^3)\) cells; ``nearly inviscid''
holds at the unit scale, not at the scale \(R'\).  The conclusion of the
firewall, that spreading of the dissipation over unit cells is not
excluded by the cellwise budget, stands with this precision, since
\(R'^3\) cells each with dissipation \(\varepsilon'/R'^3\) satisfy the unit-cell
spreading condition for any threshold below \(\varepsilon'/R'^3\).
\end{theorem}

\begin{proof}
Theorem~\ref{thm:c2v48-energy-active}(ii) and the counting.
\end{proof}

% =============================================================================
\section{Both spread}
\label{sec:c2v48-both}
% =============================================================================

\begin{proposition}[The both-spread case]
\label{prop:c2v48-both}
At a spread window (Corollary~\ref{cor:c2v48-collapse}), every window
limit at every centre map is zero, every Gaussian ledger quantity at
every centre tends to zero along the subsequence (\(\Phi\), \(\lambda D\),
\(\mathcal L\) at all scales of the window, by the compactness and the
Gaussian tails), the window dissipation \(\ge c_L^2\nu^{3/2}\ell_M\) is
carried by cells whose number tends to infinity and whose individual
energy and dissipation tend to zero, and the first cycle's rate-free
local energy bound (Theorem V66-local-energy) is satisfied with room to
spare at the window scale (it gives energy at most \(O(\lambda_j^{1/4})\) in
balls of radius \(\lambda_j\), while spreading gives \(o(\lambda_j)\)).
\end{proposition}

\begin{proof}
The vanishing of limits is the definition; the ledger quantities are
integrals of \(|U|^2\), \(|\nabla U|^2\) against Gaussians, which tend to
zero when the rescalings tend to zero in \(L^2(H^1)\) on compacts with
uniform tails (Sobolev on dyadic balls); the cell count is
\eqref{eq:c2v46-cells}; the comparison with the rate-free bound is
Proposition~\ref{prop:c2v26-independent}.
\end{proof}

\begin{assessment}[Item (AC2): the both-spread window]
\label{ass:c2v48-both}
A both-spread window is a scale at which the solution is, in the
rescaled picture, a strong solution with bounded enstrophy that is
uniformly close to zero on every compact set: the blow-up is present
only through the global enstrophy, distributed over a region of
rescaled volume tending to infinity.  Every local tool of the series
sees nothing at such a scale, and the global tools (Leray's rate, the
type-I window, the backward cone) constrain only the time structure.
Whether a solution of the intermediate class can have both-spread
windows at all its window scales is the spreading question in its
sharpest form; the series records it as open.  No global regularity
claim is made.
\end{assessment}

% =============================================================================
\section{Integrated V48 conclusion and next acquisition order}
\label{sec:c2v48-conclusion}
% =============================================================================

\begin{theorem}[What V48 proves]
\label{thm:c2v48-conclusion}
Relative to the first cycle and V1--V47, modulo (E2): the static
enstrophy lower bound \eqref{eq:c2v48-static}, the energy-active
theorem (Theorem~\ref{thm:c2v48-energy-active}), the collapse of the
refined trichotomy, the precision of Firewall~\ref{fw:c2v47-AC1}, and
the both-spread record.  The full Navier--Stokes
stability/global-regularity theorem remains unproved.
\end{theorem}

\begin{proof}
The cited results.
\end{proof}

\begin{programme}[V49 acquisition order]
\label{prog:c2v49}
\begin{enumerate}[label=\textup{(AC\arabic*)},leftmargin=4.8em]
\item Consolidate V46--V48 as the spreading chapter: cell accounting,
      backward cone, pressure and transport across cells, energy-active
      centres, the two-case structure at window scales, and the
      both-spread window as the sharpest form of the obstruction.
\item Prepare V50: the audit and the position after fifty versions,
      with the direction for V51--V55, and the manifest items to be
      carried into the next restart at V100.
\item The next compulsory companion audit is V50.
\end{enumerate}
\end{programme}

\begin{assessment}[Position after V48]
\label{ass:c2v48-position}
Energy concentration at a window forces dissipation concentration at a
larger explicit radius; spread windows are spread in both energy and
dissipation.  No full stability or global-regularity claim is made.
\end{assessment}

% =============================================================================
\section*{Version V48 (second cycle) append-only record}
\addcontentsline{toc}{section}{Version V48 (second cycle) append-only record}
% =============================================================================

The complete V47 source of the second cycle was retained before this
continuation.  Its SHA--256 digest was
\begin{center}\scriptsize\ttfamily
5edb7c3e943bada9bcfda434e729ee9f20087ab12631195bd52e529ff899ff7e.
\end{center}
V48 proved that energy-active centres at a window are dissipation-active
at a larger radius, collapsed the refined trichotomy, and made a
firewall of V47 precise.  No full regularity claim is made.

% =============================================================================
\section{V49: the spreading chapter consolidated, and the preparation of V50}
\label{sec:c2v49-entry}
% =============================================================================

Items (AC1) and (AC2) of Programme~\ref{prog:c2v49} are carried out.
V46--V48 are consolidated as the spreading chapter: the cell accounting
at a window, the backward-cone constraints between windows and spikes,
the pressure and transport bounds across cells with the cellwise energy
budget, the energy-active theorem and the two-case structure at window
scales, and the both-spread window as the sharpest form of the
obstruction.  The items to be carried into the next restart are
listed.  No global regularity claim is made.  The next version is the
compulsory companion audit.

% =============================================================================
\section{The spreading chapter, consolidated}
\label{sec:c2v49-consolidated}
% =============================================================================

\begin{theorem}[Spreading chapter, V46--V48]
\label{thm:c2v49-consolidated}
Let \(u\) be in the intermediate class with constant \(M\), with window
scales \(\lambda_j\), windows \(J_M\) of log-length \(\ell_M=3\nu^3/(8C_AM^2)\), and
unit cells \(Q_\alpha\) of the rescaled torus.  Modulo (E2):
\begin{enumerate}[label=\textup{(\arabic*)},leftmargin=3.2em]
\item \emph{Cell accounting.}  The window dissipation lies in
      \([c_L^2\nu^{3/2}\ell_M,2M\ell_M]\); at most \(2M\ell_M/\delta\) cells carry at
      least \(\delta\); at a spread window the number of carrying cells
      tends to infinity while the enstrophy stays at most \(2M\); no
      per-cell lower bound follows from Leray's rate.
\item \emph{Backward cone.}  A spike of height at least four times the
      window level cannot occur within the parabolic time
      \(\nu^3\lambda_j^2/(32C_1M^2)\) after a window time, and the gaps
      between spikes of unbounded height have summable square roots;
      the cone constrains times, not cells.
\item \emph{Cross-cell coupling.}  At a window the far pressure on each
      cell is bounded by \(C_{\rm far}C_S^2\cdot2M\) uniformly, the local
      pressure and the boundary flux are bounded by the cubic norms of
      the neighbourhood, and the cellwise energy budget
      \eqref{eq:c2v47-budget} holds; it closes only ``small stays small''.
\item \emph{Energy-active centres.}  Energy at least \(\varepsilon\) in the unit
      ball at the window time forces enstrophy at least
      \(\varepsilon/(4C_{\rm P}^2R^2)\) in \(B_R\), \(R\ge R_0(\varepsilon,M)\), at that
      time (static), and window dissipation at least
      \(\varepsilon'(\varepsilon,M,\nu)\) in \(B_{R'(\varepsilon,M,\nu)}\) (by compactness);
      hence at window scales there are two cases, concentrated (nonzero
      strong limits at some centre) and spread (energy and dissipation
      in every fixed ball tend to zero at every centre).
\item \emph{Both spread.}  At a spread window every local quantity of
      the series vanishes in the limit at every centre, and the blow-up
      is present only through the global enstrophy on a region of
      diverging rescaled volume.
\end{enumerate}
\end{theorem}

\begin{proof}
(1) Proposition~\ref{prop:c2v46-cells}.  (2) Proposition~\ref{prop:c2v46-cone}.
(3) Proposition~\ref{prop:c2v47-pressure}, Theorem~\ref{thm:c2v47-budget},
Proposition~\ref{prop:c2v47-flux}.  (4) Proposition~\ref{prop:c2v48-static},
Theorem~\ref{thm:c2v48-energy-active}, Corollary~\ref{cor:c2v48-collapse}.
(5) Proposition~\ref{prop:c2v48-both}.
\end{proof}

\begin{assessment}[What the chapter established]
\label{ass:c2v49-established}
The chapter identifies spreading, the both-spread window, as the single
configuration of a hypothetical intermediate-class singularity that
escapes every local tool of the series, and shows that it escapes for
a structural reason: all the local constraints (energy inequality,
pressure bounds, transport bounds, Leray's rate) are integrals or
global bounds, none gives a per-cell lower bound, and the only
per-cell lower bounds available (energy-active implies
dissipation-active) are themselves consequences of compactness with
non-explicit constants and a change of radius.  The same configuration,
at the parabolic scale of a type-I point, is the non-compact
dissipation of V24, and at scale zero of a rate-free point, the
frequency escape of V27; the three are the spatial, the parabolic, and
the spectral faces of one escape.  No global regularity claim is made.
\end{assessment}

% =============================================================================
\section{Items for the next restart}
\label{sec:c2v49-manifest}
% =============================================================================

\begin{programme}[Manifest for the second restart]
\label{prog:c2v49-manifest}
The summary opening the third cycle, at V100, carries, in this order:
\begin{enumerate}[label=\textup{(M\arabic*)},leftmargin=3.8em]
\item The first cycle's manifest (M1)--(M8) of the legacy file, by
      reference to the first-cycle summary of V1 of the present cycle.
\item The self-similar variables and the two identities (K1), the
      dichotomy and (E9) (K2), the statistics (K3), the moment ledger
      (K4), the rate-free chapter (K5), compact dissipation (K6), the
      Rayleigh chapter (K7), the constant-chasing chapter, the fifth
      block (modewise structure and intermediate class), and the
      spreading chapter, each with its hypotheses.
\item The import ledger (E1)--(E9) and the in-series corrections of the
      second cycle (V17, V21, V48's precision of V47).
\item The three escape faces (spatial, parabolic, spectral) and the two
      obstructions (non-explicit compactness constants; spreading), as
      the position.
\item The status: the full Navier--Stokes stability/global-regularity
      theorem is unproved; the gates (G1), (G1\('\)), (G2), (G3) with
      the second cycle's constraints.
\end{enumerate}
\end{programme}

% =============================================================================
\section{Integrated V49 conclusion and next acquisition order}
\label{sec:c2v49-conclusion}
% =============================================================================

\begin{theorem}[What V49 establishes]
\label{thm:c2v49-establishes}
Relative to the first cycle and V1--V48: the consolidated spreading
chapter (Theorem~\ref{thm:c2v49-consolidated}) and the manifest for the
next restart.  The full Navier--Stokes stability/global-regularity
theorem remains unproved.
\end{theorem}

\begin{proof}
The cited results.
\end{proof}

\begin{programme}[V50 acquisition order, with compulsory companion audit]
\label{prog:c2v50}
\begin{enumerate}[label=\textup{(AC\arabic*)},leftmargin=4.8em]
\item Perform the compulsory review of the current SM and YM notes;
      record digests; import nothing numerical.
\item Record the position after fifty versions of the second cycle and
      the direction for V51--V55.
\item Begin the direction: the spectral face of the escape (frequency
      escape) at a spread window of the intermediate class, where the
      enstrophy is bounded and the energy in every ball tends to zero,
      so that the profile's Rayleigh quotient at every centre tends to
      infinity; determine whether the window bound \(Y\le2M\) together
      with the Rayleigh identity on the window gives a quantitative
      statement about the spectral distribution at spread windows.
\item The next compulsory companion audit after V50 is V55.
\end{enumerate}
\end{programme}

\begin{assessment}[Position after V49]
\label{ass:c2v49-position}
The spreading chapter is consolidated and the manifest for the next
restart is fixed.  No full stability or global-regularity claim is
made.
\end{assessment}

% =============================================================================
\section*{Version V49 (second cycle) append-only record}
\addcontentsline{toc}{section}{Version V49 (second cycle) append-only record}
% =============================================================================

The complete V48 source of the second cycle was retained before this
continuation.  Its SHA--256 digest was
\begin{center}\scriptsize\ttfamily
393c60b785ffc2c3caa8c43948fee51640e7f9bc518935c4d5ef45b113107916.
\end{center}
V49 consolidated the spreading chapter and fixed the manifest for the
next restart.  No full regularity claim is made.

% =============================================================================
\section{V50: compulsory companion audit, the position after fifty versions,
and the spectral face at a spread window}
\label{sec:c2v50-entry}
% =============================================================================

This is the fiftieth version of the second cycle.  The compulsory review
of the two companion programmes is performed first.  The position after
fifty versions is then recorded with the direction for V51--V55, and
item (AC3) of Programme~\ref{prog:c2v50} is carried out: at a spread
window of the intermediate class every Gaussian ledger quantity at
every centre vanishes in the limit, so the Rayleigh quotient of the
profile is a ratio of vanishing quantities and the spectral face of
the escape carries no local information there; the only spectral
statement is global, that the rescaled enstrophy stays between Leray's
rate and twice the sequence constant while the energy per unit volume
tends to zero, which is recorded exactly.  No global regularity claim
is made.  The next compulsory companion audit is V55.

% =============================================================================
\section{Compulsory V50 review of the current companion notes}
\label{sec:c2v50-companion-audit}
% =============================================================================

The companion files in \texttt{latex\_notes} were inspected read-only at
the time of the preparation of this version.  The frozen checkpoints
are
\begin{align}
 &\texttt{Renewal\_SM\_Einstein\_Continuum\_Research\_Notes\_V100.tex},
 \notag\\[-2pt]
 &\qquad
 \texttt{b5a5540cbda10d6371506f77948dd85b99dff3c09a96a3fb30a2e65f77d564de},
 \label{eq:c2v50-SM-checkpoint}\\
 &\texttt{Renewal\_Yang\_Mills\_Mass\_Gap\_Research\_Notes\_V76.tex},
 \notag\\[-2pt]
 &\qquad
 \texttt{06c211fb95198be15e8e6776a18374af164182b38114101d4677db294329244f}.
 \label{eq:c2v50-YM-checkpoint}
\end{align}
A YM file named \(\ldots\)\texttt{V75.tex} with a different digest
(\texttt{0c44941e\ldots}) was also present; the V76 file contains the
V75 chapter and is recorded.  The SM file was renamed by its programme
between the inspection and the assembly of this version (its
hundredth version is a restart point of that programme); the digest
above is that of the inspected file.

\emph{SM V100.}  The SM programme performed its hundredth-version
audit, recording the present cycle's V43 by digest and transferring,
in its words, a ``spatial-no-escape warning'' (this cycle's record
that spreading of dissipation is not excluded by local tools); it then
proved an exponentiated Jacobi transport identity, a trace-norm length
bound, computed an exact endpoint frame-transition phase, proved the
sharp matrix-dimension loss of raw operator-norm control, and
formulated a trace-norm condition with its finite-cutoff consequence.

\emph{YM V75--V76.}  The YM programme performed its compulsory audit
(reading the present cycle and importing two disciplines), proved a
rooted-tree/child-partition bijection, summed every hyperedge blocking
over a rooted tree, proved a strong card for each fixed rooted tree,
showed a parent-map relaxation to be superexponential, and formulated a
rooted-hypertree condition sufficient for every incidence hypertree
(V75); then defined a normalized rooted-hypertree quotient polynomial,
proved its positivity and exact degree, a Pr\"ufer mass bound, closed
the condition for uniformly balanced centre marks and the equal-mark
model, computed every pure-power coefficient, and formulated a further
condition sufficient for the full one (V76).

\begin{theorem}[V50 companion-audit decision]
\label{thm:c2v50-companion-decision}
No SM V100 transport, trace-norm, or phase constant and no YM V75--V76
tree, hypertree, or polynomial bound is imported.  Neither bears on the
spreading chapter or on the gates.  The SM reading of the present cycle
(a digest and the spreading warning) is accurate.
\end{theorem}

\begin{proof}
The SM results concern transport identities of lattice frames; the YM
results concern hypertree generating functions of a lattice clock.
None is a Navier--Stokes statement.
\end{proof}

% =============================================================================
\section{Position after fifty versions and direction}
\label{sec:c2v50-position}
% =============================================================================

\begin{assessment}[Position after fifty versions of the second cycle]
\label{ass:c2v50-position-fifty}
The second cycle has, in fifty versions, built an exact scale-invariant
description of a hypothetical type-I singularity at homogeneous points
(self-similar variables, two identities, dichotomy, statistics, moment
ledger, Rayleigh chapter, modewise structure), derived from it the first
explicit inequality such a singularity must satisfy (the
constant-chasing gate, one side explicit), extended its ledgers
rate-free to every smooth solution (with forced divergences at
Gaussian-invisible dissipative points), analysed the intermediate class
through its type-I windows, and identified spreading, in its spatial,
parabolic, and spectral faces, as the single configuration that escapes
all local tools, with the non-explicit compactness constants as its
counterpart on the type-I side.  The full Navier--Stokes
stability/global-regularity theorem is not proved.  For V51--V55 the
direction is: (a) the spectral face at spread windows (below, trivial
locally); (b) the global Fourier structure at a spread window, where
the enstrophy is bounded above and below on a torus of diverging
rescaled size and the energy per unit volume tends to zero: determine
what the energy and enstrophy identities on the big torus, together
with the window bound, say about the Fourier spectrum of the rescaled
solution (a Kolmogorov-type constraint on the distribution of a fixed
enstrophy over a diverging volume); (c) the consolidation at V55.  No
global regularity claim is made.
\end{assessment}

% =============================================================================
\section{The spectral face at a spread window}
\label{sec:c2v50-spectral}
% =============================================================================

\begin{proposition}[Local spectral triviality at a spread window]
\label{prop:c2v50-local}
At a spread window (Corollary~\ref{cor:c2v48-collapse}), for every
centre map \(x_j\) and every scale of the window, \(\varphi\to0\),
\(\mathcal L\to0\), \(\|Y\|_G\to0\) and \(\|N(V)\|_G\to0\) along the subsequence
(the last two by the window bounds transported to the Gaussian norms
with vanishing energy and dissipation on every compact set), so the
Rayleigh quotient \(r=2\mathcal L/\varphi\) is a ratio of quantities tending
to zero and is not determined; the spectral face of the escape carries
no local information at a spread window.
\end{proposition}

\begin{proof}
Proposition~\ref{prop:c2v48-both} and the definitions of the Gaussian
quantities; \(\|N(V)\|_G\le\|V\|_{L^6}\|\nabla V\|_{L^3(G)}\)-type bounds with
the \(H^2\) window bound and the vanishing on compacts give the last
statement.
\end{proof}

\begin{proposition}[Global spectral statement at a spread window]
\label{prop:c2v50-global}
At a spread window, on the rescaled torus \(\Torus^3_j\) of side
\(L_j=2\pi/\lambda_j\to\infty\), with Fourier coefficients \(\widehat U_j(k)\),
\(k\in(2\pi/L_j)\mathbb Z^3\):
\begin{equation}
 c_L^2\nu^{3/2}\ \le\ \sum_k|k|^2|\widehat U_j(k)|^2\ \le\ 2M ,
 \qquad
 \frac1{L_j^3}\sum_k|\widehat U_j(k)|^2=\frac{E_*}{\lambda_jL_j^3}=\frac{E_*\lambda_j^2}{(2\pi)^3}\to0 ,
 \label{eq:c2v50-global}
\end{equation}
so the enstrophy is of order one while the mean energy density tends to
zero: the enstrophy is carried by Fourier modes with
\(\sum_k|k|^2|\widehat U_j(k)|^2\ge c\) and \(\sum_k|\widehat U_j(k)|^2=E_*/\lambda_j\to\infty\)
(the total rescaled energy), whose energy-weighted mean square
wavenumber \(\sum|k|^2|\widehat U|^2/\sum|\widehat U|^2\) tends to zero: in the
rescaled units, the energy sits at wavenumbers tending to zero while
the enstrophy stays of order one, that is, the spectrum spreads to
large scales in the rescaled picture (small wavenumbers on the big
torus) with a fixed enstrophy content.
\end{proposition}

\begin{proof}
Parseval on \(\Torus^3_j\) with \(\int|U_j|^2=E_*/\lambda_j\) and
\(\int|\nabla U_j|^2=Y_j\in[c_L^2\nu^{3/2},2M]\) at window times.
\end{proof}

\begin{assessment}[Item (AC3)]
\label{ass:c2v50-AC3}
At a spread window the singular solution, in rescaled units, has a
fixed enstrophy and a diverging energy spread over a diverging volume,
so that its energy-weighted mean square wavenumber tends to zero; the
enstrophy-carrying modes have wavenumbers of order one or larger and
carry a vanishing fraction of the energy.  This is the statement that
the type-I window sees only the part of the flow at the parabolic scale
and finer, while the energy is elsewhere, at larger scales; it is
consistent with the first cycle's picture of the terminal datum (energy
in the critical Morrey class, concentrated at scales larger than the
parabolic one at the singular time) and gives no constraint beyond
Parseval.  The global energy and enstrophy identities on the big torus
add the balance of the enstrophy (vortex stretching against
palinstrophy) with no sign, as in the first cycle.  Item (AC3) is
closed with this record, and item (b) of the direction is the natural
continuation.
\end{assessment}

% =============================================================================
\section{Integrated V50 conclusion and next acquisition order}
\label{sec:c2v50-conclusion}
% =============================================================================

\begin{theorem}[What V50 establishes]
\label{thm:c2v50-conclusion}
Relative to the first cycle and V1--V49: the companion-audit decision
(Theorem~\ref{thm:c2v50-companion-decision}), the position after fifty
versions, and the local and global spectral statements at a spread
window (Propositions~\ref{prop:c2v50-local}, \ref{prop:c2v50-global}).
The full Navier--Stokes stability/global-regularity theorem remains
unproved.
\end{theorem}

\begin{proof}
The cited results.
\end{proof}

\begin{programme}[V51 acquisition order]
\label{prog:c2v51}
\begin{enumerate}[label=\textup{(AC\arabic*)},leftmargin=4.8em]
\item Kolmogorov-type constraint at a spread window.  On the big torus
      at a window time, with enstrophy in \([c_L^2\nu^{3/2},2M]\) and total
      energy \(E_*/\lambda_j\), use the energy inequality on the window and
      the enstrophy inequality with its palinstrophy term to bound the
      palinstrophy \(\int|\Delta U_j|^2\) integrated over the window by
      \(C(M,\nu)\), and derive the resulting bound on the enstrophy at
      wavenumbers above \(K\): \(\sum_{|k|\ge K}|k|^2|\widehat U|^2\le C(M,\nu)/(K^2\ell_M)\)
      on average over the window; record the explicit spectral tail
      bound and whether it forces the enstrophy into a bounded band of
      wavenumbers at a spread window.
\item If the enstrophy sits in a bounded band \(|k|\le K(M,\nu)\) with
      energy density tending to zero, determine whether spreading is
      compatible with the band: modes of wavenumber \(\le K\) with total
      enstrophy of order one and energy per unit volume tending to zero
      require the modes' amplitudes to be small and their number (in
      the big torus's lattice) to be large; record whether the
      nonlinear term can sustain such a configuration over the window
      (it can, for a superposition of small-amplitude waves, at the
      linear level), and the exact statement.
\item The next compulsory companion audit is V55.
\end{enumerate}
\end{programme}

\begin{assessment}[Position after V50]
\label{ass:c2v50-position}
The second cycle is at fifty versions with its position and direction
recorded; the spectral face at spread windows is locally trivial and
globally a Parseval statement.  No full stability or
global-regularity claim is made.
\end{assessment}

% =============================================================================
\section*{Version V50 (second cycle) append-only record}
\addcontentsline{toc}{section}{Version V50 (second cycle) append-only record}
% =============================================================================

The complete V49 source of the second cycle was retained before this
continuation.  Its SHA--256 digest was
\begin{center}\scriptsize\ttfamily
f80f90fc931d8ae201ca1c10ba3d99b14199cf60965b480d9bda2cda062c2c79.
\end{center}
V50 performed the compulsory companion audit, recorded the position
after fifty versions, and recorded the spectral face at spread
windows.  No full regularity claim is made.

% =============================================================================
\section{V51: the enstrophy band at a type-I window and the structure of a
spread window}
\label{sec:c2v51-entry}
% =============================================================================

Items (AC1) and (AC2) of Programme~\ref{prog:c2v51} are carried out.  On
a type-I window the palinstrophy integrated over the window is bounded
by an explicit function of the sequence constant, by the enstrophy
inequality; hence, on average over the window, the enstrophy at
wavenumbers above a threshold is bounded by an explicit inverse square
of the threshold, and at least half of Leray's enstrophy floor sits in
an explicit bounded band of wavenumbers in the rescaled units.  At a
spread window this band carries a fixed enstrophy and a fixed energy
while the energy density tends to zero: the rescaled solution is a
delocalized superposition of structures of parabolic size and larger
with amplitude tending to zero, which is compatible with the linear
evolution over the window.  The spreading configuration is thereby
described exactly in Fourier terms, without being excluded.  No global
regularity claim is made.  The next compulsory companion audit is V55.

% =============================================================================
\section{The palinstrophy bound and the enstrophy band}
\label{sec:c2v51-band}
% =============================================================================

Let \(u\) be in the intermediate class with constant \(M\), \(U=U_j\) the
rescaling at a window scale on the torus \(\Torus^3_j\) of side
\(L_j=2\pi/\lambda_j\), \(J_M=[-1,-1+\ell_M]\) the window, \(Y(s)=\|\nabla U(s)\|_2^2\),
and \(\mathcal E_2(s):=\|\Delta U(s)\|_2^2\) the palinstrophy.  Let \(C_B\) be
the constant of the vortex-stretching bound
\(|\int\partial_iU_j\partial_jU_k\partial_kU_i|\le C_B\|\nabla U\|_2^{3/2}\|\Delta U\|_2^{3/2}\)
on the torus (Gagliardo--Nirenberg; scale-invariant).

\begin{proposition}[Palinstrophy bound on a window]
\label{prop:c2v51-palinstrophy}
Let \(C_A:=\frac{27}{16}C_B^4\), which is the constant of the enstrophy
inequality \(Y'\le C_A\nu^{-3}Y^3\) as derived below, and
\(\ell_M=3\nu^3/(8C_AM^2)\) the window length of
Proposition~\ref{prop:c2v42-window} with this \(C_A\).  Then
\begin{equation}
 \nu\int_{J_M}\mathcal E_2(s)\,\dd s\ \le\ P_M:=4M .
 \label{eq:c2v51-palinstrophy}
\end{equation}
\end{proposition}

\begin{proof}
The enstrophy identity on the torus,
\(\frac{\dd}{\dd s}Y=-2\nu\mathcal E_2+2\int\partial_iU_j\partial_jU_k\partial_kU_i\)
(the pressure term vanishes for divergence-free fields), and the
vortex-stretching bound give
\(\frac{\dd}{\dd s}Y\le-2\nu\mathcal E_2+2C_BY^{3/4}\mathcal E_2^{3/4}\).  Young's
inequality with exponents \(\frac43\) and \(4\),
\(2C_BY^{3/4}\mathcal E_2^{3/4}\le\nu\mathcal E_2+\frac{27}{16}\nu^{-3}C_B^4Y^3\),
gives \(\frac{\dd}{\dd s}Y\le-\nu\mathcal E_2+C_A\nu^{-3}Y^3\).  Integrating over
\(J_M\), on which \(Y\le2M\) and \(Y(-1)\le M\):
\(\nu\int_{J_M}\mathcal E_2\le Y(-1)+C_A\nu^{-3}(2M)^3\ell_M=M+C_A\nu^{-3}\cdot8M^3\cdot\frac{3\nu^3}{8C_AM^2}=M+3M=4M\).
\end{proof}

\begin{theorem}[Enstrophy band at a window]
\label{thm:c2v51-band}
With \(\widehat U(k,s)\) the Fourier coefficients on \(\Torus^3_j\), for every
\(K>0\),
\begin{equation}
 \frac1{\ell_M}\int_{J_M}\sum_{|k|\ge K}|k|^2|\widehat U(k,s)|^2\,\dd s\ \le\ \frac{P_M}{\nu\,\ell_M\,K^2}
 =\frac{8\,C_A\,M^2P_M}{3\,\nu^4K^2},
 \label{eq:c2v51-tail}
\end{equation}
and consequently, with
\begin{equation}
 K_M:=\Bigl(\frac{16\,C_AM^2P_M}{3\,\nu^4c_L^2\nu^{3/2}}\Bigr)^{1/2},
 \label{eq:c2v51-KM}
\end{equation}
at least half of Leray's enstrophy floor sits, on average over the
window, at wavenumbers below \(K_M\):
\begin{equation}
 \frac1{\ell_M}\int_{J_M}\sum_{|k|<K_M}|k|^2|\widehat U(k,s)|^2\,\dd s\ \ge\ \frac12c_L^2\nu^{3/2} .
 \label{eq:c2v51-band}
\end{equation}
\end{theorem}

\begin{proof}
\(\sum_{|k|\ge K}|k|^2|\widehat U|^2\le K^{-2}\sum|k|^4|\widehat U|^2=K^{-2}\mathcal E_2\) and
\eqref{eq:c2v51-palinstrophy}.  The floor \(Y\ge c_L^2\nu^{3/2}\) at every
window time (Leray's rate in rescaled units, \((-s)^{-1/2}\ge1\) on
\(J_M\)) and the choice of \(K_M\) making the tail at most
\(\frac12c_L^2\nu^{3/2}\).
\end{proof}

\begin{assessment}[Item (AC1)]
\label{ass:c2v51-AC1}
A type-I window forces the enstrophy into an explicit bounded band of
rescaled wavenumbers, \(|k|<K_M\) with \(K_M\) an explicit function of
\(M\) and \(\nu\) (of order \(M^{3/2}\nu^{-11/4}\) in the units used):
structures of size at least \(\lambda_j/K_M\), that is, of parabolic size
and larger.  This is the quantitative content of the window's \(H^2\)
control: the solution cannot, on a type-I window, hide its enstrophy at
scales much finer than the parabolic one.  Combined with the spread
case, in which no unit ball carries enstrophy, the enstrophy at a
spread window is carried by structures of parabolic size or larger
distributed over a region of diverging rescaled volume with vanishing
density.
\end{assessment}

% =============================================================================
\section{The spread window in Fourier terms}
\label{sec:c2v51-spread}
% =============================================================================

\begin{proposition}[Fourier description of a spread window]
\label{prop:c2v51-spread}
At a spread window (Corollary~\ref{cor:c2v48-collapse}), on average
over the window: the band \(|k|<K_M\) carries enstrophy at least
\(\frac12c_L^2\nu^{3/2}\) and energy at least
\(\frac12c_L^2\nu^{3/2}K_M^{-2}\); the total energy is \(E_*/\lambda_j\to\infty\); the
energy density \(E_*\lambda_j^2/(2\pi)^3\to0\); and every unit ball carries
energy and dissipation tending to zero.  Hence the band modes form a
delocalized superposition with \(L^2\)-amplitude per unit volume tending
to zero, and the energy of the band is a vanishing fraction of the
total energy, which sits at wavenumbers tending to zero (rescaled
structures of size tending to infinity, that is, larger than any
multiple of the parabolic scale).
\end{proposition}

\begin{proof}
\eqref{eq:c2v51-band}, \(|\widehat U|^2\ge|k|^{-2}(|k|^2|\widehat U|^2)\) on the band,
Parseval, and Proposition~\ref{prop:c2v48-both}.
\end{proof}

\begin{proposition}[Compatibility with the linear evolution]
\label{prop:c2v51-linear}
A superposition of Fourier modes on \(\Torus^3_j\) with enstrophy in
\([c_L^2\nu^{3/2},2M]\) carried by \(|k|<K_M\), evolving by the Stokes flow
over the window, keeps at least \(e^{-2\nu K_M^2\ell_M}\) of its enstrophy
and satisfies the window bound; with \(\ell_M=3\nu^3/(8C_AM^2)\) this
factor is \(\exp(-\frac{3\nu^4K_M^2}{4C_AM^2})\), bounded below by an explicit
positive constant depending on \(M\) and \(\nu\).  The nonlinear term
acting on such a superposition is bounded in \(L^2\) by
\(C\|U\|_6\|\nabla U\|_3\le CC_S(2M)^{1/2}\|\nabla U\|_2^{1/2}\|\Delta U\|_2^{1/2}\), which is
not small in general; the series does not decide whether the nonlinear
evolution preserves a spread configuration across a window, only that
the window bounds permit it.
\end{proposition}

\begin{proof}
\(|\widehat U(k,s)|^2\) decays by \(e^{-2\nu|k|^2(s+1)}\) under the Stokes flow;
the nonlinear bound is Gagliardo--Nirenberg on the torus.
\end{proof}

\begin{assessment}[Item (AC2): what spreading is]
\label{ass:c2v51-AC2}
A spread type-I window is now described exactly: in the rescaled
units, a field of small amplitude spread over a region of diverging
volume, with its enstrophy of order one carried by structures of
parabolic size or larger, and its energy, diverging in total but
vanishing in density, at ever larger scales.  Nothing in the window
bounds, the energy identity, or the linear evolution forbids it.  A
singular solution of the intermediate class with only spread windows
would blow up by developing, between windows, spikes of enstrophy
from a state that is, at each window, nearly linear, weak, and
delocalized; the backward cone and the gap summability are the only
constraints on the time structure of such spikes.  The series records
this as the exact form of the spreading obstruction in the
intermediate class.  No global regularity claim is made.
\end{assessment}

% =============================================================================
\section{Integrated V51 conclusion and next acquisition order}
\label{sec:c2v51-conclusion}
% =============================================================================

\begin{theorem}[What V51 proves]
\label{thm:c2v51-conclusion}
Relative to the first cycle and V1--V50: the palinstrophy bound
\eqref{eq:c2v51-palinstrophy}, the spectral tail bound
\eqref{eq:c2v51-tail} and the enstrophy band \eqref{eq:c2v51-band}, the
Fourier description of a spread window, and its compatibility with the
linear evolution.  The full Navier--Stokes stability/global-regularity
theorem remains unproved.
\end{theorem}

\begin{proof}
The cited results.
\end{proof}

\begin{programme}[V52 acquisition order]
\label{prog:c2v52}
\begin{enumerate}[label=\textup{(AC\arabic*)},leftmargin=4.8em]
\item Spikes from a spread window.  Between a spread window at scale
      \(\lambda_j\) and the next spike, the enstrophy must rise from at most
      \(2M/\lambda_j\) to a large value; using the enstrophy identity on the
      big torus and the band structure at the window, bound the growth
      rate of the enstrophy from a spread state: the vortex stretching
      is at most \(C_B\|\nabla U\|^{3/2}\|\Delta U\|^{3/2}\), so a state with
      enstrophy \(Y\) needs palinstrophy of order \(Y^3\nu^{-3}\) to grow at
      the maximal rate; determine whether a delocalized state (energy
      density tending to zero) can have such palinstrophy, that is,
      whether spreading of the energy forces small palinstrophy per
      unit enstrophy, and record the exact statement.
\item If a delocalized state cannot grow its enstrophy at the maximal
      rate, quantify the delay before the next spike in terms of the
      energy density, and compare with the backward cone's parabolic
      delay.
\item The next compulsory companion audit is V55.
\end{enumerate}
\end{programme}

\begin{assessment}[Position after V51]
\label{ass:c2v51-position}
Type-I windows confine the enstrophy to an explicit band of parabolic
and larger scales; a spread window is a weak delocalized state in that
band, compatible with the window bounds.  No full stability or
global-regularity claim is made.
\end{assessment}

% =============================================================================
\section*{Version V51 (second cycle) append-only record}
\addcontentsline{toc}{section}{Version V51 (second cycle) append-only record}
% =============================================================================

The complete V50 source of the second cycle was retained before this
continuation.  Its SHA--256 digest was
\begin{center}\scriptsize\ttfamily
22c63459731e30b617587fd97d24a5e93aa7ea7f6474a120c2238d7e9f62edbe.
\end{center}
V51 proved the palinstrophy bound and the enstrophy band on a type-I
window and described a spread window in Fourier terms.  No full
regularity claim is made.

% =============================================================================
\section{V52: enstrophy growth from a spread window}
\label{sec:c2v52-entry}
% =============================================================================

Items (AC1) and (AC2) of Programme~\ref{prog:c2v52} are carried out.  The
enstrophy inequality is optimized over the palinstrophy: the maximal
growth rate of the enstrophy is an explicit multiple of its cube, and
it is attained when the palinstrophy is an explicit multiple of the
cube of the enstrophy; on a type-I window this palinstrophy is within
the window's palinstrophy budget, so the enstrophy may leave the window
at the maximal rate.  Delocalization of the energy does not bound the
enstrophy-weighted mean square wavenumber, hence does not reduce the
available palinstrophy, and the only delays before the next spike are
the window length and the backward cone's parabolic delay.  The
sup-norm of a spread state is shown to tend to zero on the cells whose
local \(H^2\) norm is bounded, which is most cells in aggregate but not
necessarily all.  No global regularity claim is made.  The next
compulsory companion audit is V55.

% =============================================================================
\section{Maximal enstrophy growth}
\label{sec:c2v52-growth}
% =============================================================================

\begin{proposition}[Sharp form of the enstrophy inequality and the optimal palinstrophy]
\label{prop:c2v52-sharp}
For a smooth mean-zero solution on a torus, with \(Y=\|\nabla U\|_2^2\),
\(\mathcal E_2=\|\Delta U\|_2^2\), and \(C_B\) the vortex-stretching constant,
\begin{equation}
 \frac{\dd Y}{\dd s}\ \le\ \max_{\mathcal E\ge0}\bigl(-2\nu\mathcal E+2C_BY^{3/4}\mathcal E^{3/4}\bigr)
 =\frac{27\,C_B^4}{128\,\nu^3}\,Y^3 ,
 \qquad
 \text{attained at }\ \mathcal E_*(Y)=\Bigl(\frac{3C_B}{4\nu}\Bigr)^4Y^3 .
 \label{eq:c2v52-sharp}
\end{equation}
Hence the enstrophy can grow at the maximal rate only if the
palinstrophy equals \(\mathcal E_*(Y)\) up to a constant factor, and a
state with palinstrophy \(\mathcal E\le\theta\mathcal E_*(Y)\), \(\theta<1\), grows
at most at the rate \(2C_BY^{3/4}(\theta\mathcal E_*)^{3/4}-2\nu\theta\mathcal E_*=\frac{27C_B^4}{128\nu^3}Y^3\,(4\theta^{3/4}-3\theta)\),
which is positive exactly for \(\theta<(4/3)^4\), maximal at \(\theta=1\), and
negative for larger palinstrophy (excess palinstrophy dissipates
enstrophy faster than stretching produces it).
\end{proposition}

\begin{proof}
Differentiate \(f(\mathcal E)=-2\nu\mathcal E+2C_BY^{3/4}\mathcal E^{3/4}\):
\(f'=0\) at \(\mathcal E^{1/4}=\frac{3C_B}{4\nu}Y^{3/4}\), and
\(f(\mathcal E_*)=Y^3(\frac{3C_B}{4\nu})^3(2C_B-\frac{3C_B}2)=\frac{27C_B^4}{128\nu^3}Y^3\).
The \(\theta\)-form is the substitution \(\mathcal E=\theta\mathcal E_*\).
\end{proof}

\begin{proposition}[The optimal palinstrophy is within the window budget]
\label{prop:c2v52-budget}
On a type-I window with \(Y\le2M\), the palinstrophy required for
maximal growth is \(\mathcal E_*(2M)=\frac{81C_B^4}{32\nu^4}M^3\), while the window
average of the palinstrophy is at most
\(\frac{P_M}{\nu\ell_M}=\frac{32C_AM^3}{3\nu^4}=\frac{18C_B^4M^3}{\nu^4}\)
(with \(C_A=\frac{27}{16}C_B^4\), \(P_M=4M\)); the required value is
\(\frac{81}{32\cdot18}\approx0.14\) of the budget.  Hence the enstrophy may
grow at the maximal rate on the window, and the window length
\(\ell_M\) is the doubling time \(Y'\le C_A\nu^{-3}Y^3\) allows.
\end{proposition}

\begin{proof}
Arithmetic from Propositions~\ref{prop:c2v51-palinstrophy} and
\ref{prop:c2v52-sharp}.
\end{proof}

\begin{firewall}[Item (AC1): delocalization does not bound the palinstrophy per enstrophy]
\label{fw:c2v52-AC1}
The ratio \(\mathcal E_2/Y\) is the enstrophy-weighted mean square
wavenumber, a property of the distribution of the enstrophy over
wavenumbers, while delocalization (energy density tending to zero) is a
property of the distribution of the energy over space; the two are
independent within the window bounds, since a delocalized field may
carry its enstrophy at any wavenumbers up to the band bound \(K_M\),
and \(\mathcal E_*(Y)\) is within the budget.  Hence spreading does not
force small palinstrophy, does not slow the enstrophy growth, and the
next spike may follow a spread window after the parabolic delay of the
backward cone (Proposition~\ref{prop:c2v46-cone}) at the maximal rate.
Item (AC2) has therefore no content beyond that delay, and is closed.
\end{firewall}

% =============================================================================
\section{The amplitude of a spread state}
\label{sec:c2v52-amplitude}
% =============================================================================

\begin{proposition}[Sup-norm on \(H^2\)-bounded cells]
\label{prop:c2v52-supnorm}
At a spread window, for every cell \(Q_\alpha\) and \(s\in J_M\),
\begin{equation}
 \|U(s)\|_{L^\infty(Q_\alpha)}\ \le\ C\,\|U(s)\|_{L^2(2Q_\alpha)}^{1/4}\,\|U(s)\|_{H^2(2Q_\alpha)}^{3/4} ,
 \label{eq:c2v52-supnorm}
\end{equation}
so that on any family of cells whose \(H^2\) norms over the window are
bounded by a fixed constant, the sup-norm tends to zero uniformly along
the subsequence, in \(L^{4/3}\) of the window time.  Since
\(\sum_\alpha\int_{J_M}\|U\|_{H^2(2Q_\alpha)}^2\le C(M,\nu)\), all but a number of
cells bounded by \(C(M,\nu)/H\) have window \(H^2\) norm at most \(H\), and
the sup-norm tends to zero on those cells; on the remaining cells no
bound is available.
\end{proposition}

\begin{proof}
Gagliardo--Nirenberg on the cube of side two, and the counting from
the summed \(H^2\) bound (Proposition~\ref{prop:c2v51-palinstrophy} with
the energy and enstrophy bounds).
\end{proof}

\begin{assessment}[The spread state]
\label{ass:c2v52-state}
A spread window is a weak state almost everywhere (its amplitude tends
to zero on all but a bounded number of cells), with enstrophy of order
one carried by structures of parabolic size and larger, and with the
palinstrophy needed to grow its enstrophy at the maximal rate within
its budget.  The series can describe this state exactly and cannot
exclude it: the enstrophy inequality, being global, permits it to grow
into a spike after the parabolic delay, and nothing local sees it.
Whether the growth can actually be realized by the Navier--Stokes
nonlinearity from such a state, that is, whether the vortex stretching
can reach the size \(C_BY^{3/4}\mathcal E_2^{3/4}\) when the amplitude is
small almost everywhere, is the question the series records as open;
the cells with unbounded \(H^2\) norm are where it would happen.  No
global regularity claim is made.
\end{assessment}

% =============================================================================
\section{Integrated V52 conclusion and next acquisition order}
\label{sec:c2v52-conclusion}
% =============================================================================

\begin{theorem}[What V52 proves]
\label{thm:c2v52-conclusion}
Relative to the first cycle and V1--V51: the sharp enstrophy inequality
\eqref{eq:c2v52-sharp} with its optimal palinstrophy, the budget
comparison (Proposition~\ref{prop:c2v52-budget}), and the sup-norm
statement \eqref{eq:c2v52-supnorm}.  The full Navier--Stokes
stability/global-regularity theorem remains unproved.
\end{theorem}

\begin{proof}
The cited results.
\end{proof}

\begin{programme}[V53 acquisition order]
\label{prog:c2v53}
\begin{enumerate}[label=\textup{(AC\arabic*)},leftmargin=4.8em]
\item Stretching from a weak state.  Bound the vortex stretching
      \(\int\partial_iU_j\partial_jU_k\partial_kU_i\) cellwise by
      \(\|\nabla U\|_{L^\infty(Q)}\|\nabla U\|_{L^2(Q)}^2\) and by the cells'
      \(H^2\) norms, and determine whether at a spread window the
      stretching integrated over the window is bounded by a quantity
      tending to zero on the \(H^2\)-bounded cells, so that the enstrophy
      growth over the window is carried by the bounded number of
      exceptional cells; record the resulting statement (a spike from a
      spread window is seeded in a bounded number of cells).
\item If the seeding is in a bounded number of cells, apply the
      moving-centre compactness at those cells: the seed cells have
      energy at least the exceptional \(H^2\) level's implication or are
      themselves spread; determine whether a contradiction with the
      spread hypothesis follows, or whether the exceptional cells can
      have large \(H^2\) norm with small energy.
\item The next compulsory companion audit is V55.
\end{enumerate}
\end{programme}

\begin{assessment}[Position after V52]
\label{ass:c2v52-position}
The enstrophy may leave a spread window at the maximal rate; the
spread state is weak on all but a bounded number of cells.  No full
stability or global-regularity claim is made.
\end{assessment}

% =============================================================================
\section*{Version V52 (second cycle) append-only record}
\addcontentsline{toc}{section}{Version V52 (second cycle) append-only record}
% =============================================================================

The complete V51 source of the second cycle was retained before this
continuation.  Its SHA--256 digest was
\begin{center}\scriptsize\ttfamily
e284523d92e3b5861307be6d222d76d91169a505cb902afc786b5925b0246aa8.
\end{center}
(The V51 file was corrected in place after its assembly, a factor in
the spectral tail display; the digest above is that of the corrected
file.)  V52 proved the sharp enstrophy inequality with its optimal
palinstrophy and recorded that spreading does not slow enstrophy
growth.  No full regularity claim is made.

% =============================================================================
\section{V53: stretching from a weak state and the two kinds of spread window}
\label{sec:c2v53-entry}
% =============================================================================

Items (AC1) and (AC2) of Programme~\ref{prog:c2v53} are carried out.  The
vortex stretching integrated over a type-I window is bounded, by a
cellwise Gagliardo--Nirenberg inequality and H\"older's inequality over
the cells, by the fourth root of the window integral of the squared
maximal local enstrophy times explicit powers of the enstrophy and
palinstrophy budgets; hence on a window whose maximal local enstrophy
has vanishing time integral the enstrophy cannot grow within the
window.  This distinguishes two kinds of spread window: uniformly
spread, where the maximal local enstrophy vanishes in time-average and
no spike is seeded within the window, and travelling, where each cell's
time-integrated enstrophy vanishes while the instantaneous maximum
does not, which requires concentrated enstrophy moving through the
cells faster than one cell per window.  The exceptional-cell seeding
question of the programme is answered: on a uniformly spread window
there is no seeding at all; on a travelling window the seed is the
moving structure, whose centre map is incoherent in the sense of V24.
No global regularity claim is made.  The next compulsory companion
audit is V55.

% =============================================================================
\section{The window stretching bound}
\label{sec:c2v53-stretching}
% =============================================================================

Let \(u\) be in the intermediate class with constant \(M\), \(U\) the
rescaling at a window scale, \(J_M\) the window, \(Q_\alpha\) the unit cells
and \(2Q_\alpha\) their doubles, \(y_\alpha(s):=\int_{2Q_\alpha}|\nabla U(s)|^2\) the local
enstrophy, \(h_\alpha(s):=\|\nabla U(s)\|_{H^1(2Q_\alpha)}^2\), and
\(\delta(s):=\max_\alpha y_\alpha(s)\) the maximal local enstrophy.

\begin{theorem}[Window stretching bound]
\label{thm:c2v53-stretching}
With \(C_{\rm GN}\) the constant of \(\|f\|_{L^3(Q)}^3\le C_{\rm GN}\|f\|_{L^2(2Q)}^{3/2}\|f\|_{H^1(2Q)}^{3/2}\)
on a unit cube,
\begin{equation}
 \int_{J_M}\Bigl|\int\partial_iU_j\partial_jU_k\partial_kU_i\Bigr|\dd s
 \ \le\ C_{\rm GN}\,(2M)^{1/4}\Bigl(\int_{J_M}\delta(s)^2\,\dd s\Bigr)^{1/4}\Bigl(\int_{J_M}\sum_\alpha h_\alpha(s)\,\dd s\Bigr)^{3/4},
 \label{eq:c2v53-stretching}
\end{equation}
and \(\int_{J_M}\sum_\alpha h_\alpha\le C(2M\ell_M+P_M/\nu)=:H_M\) (explicit).
Consequently
\begin{equation}
 Y(s_2)-Y(s_1)\ \le\ 2C_{\rm GN}(2M)^{1/4}H_M^{3/4}\Bigl(\int_{J_M}\delta^2\Bigr)^{1/4}
 \qquad(s_1<s_2\in J_M).
 \label{eq:c2v53-growth}
\end{equation}
\end{theorem}

\begin{proof}
Cellwise, \(|\int_{Q_\alpha}\partial_iU_j\partial_jU_k\partial_kU_i|\le\|\nabla U\|_{L^3(Q_\alpha)}^3\le C_{\rm GN}y_\alpha^{3/4}h_\alpha^{3/4}\).
Summing over the cells and applying H\"older with exponents \(4\) and
\(\frac43\): \(\sum_\alpha y_\alpha^{3/4}h_\alpha^{3/4}\le(\sum_\alpha y_\alpha^3)^{1/4}(\sum_\alpha h_\alpha)^{3/4}\),
and \(\sum_\alpha y_\alpha^3\le\delta^2\sum_\alpha y_\alpha\le8\delta^2Y\) (the doubles
overlap at most eight times), with \(Y\le2M\).  Integrating in time with
H\"older (exponents \(4\), \(\frac43\)) gives \eqref{eq:c2v53-stretching}
up to the absorbed factor \(8^{1/4}\).  The bound on \(\int\sum h_\alpha\) is
the enstrophy and palinstrophy budgets of the window
(Propositions~\ref{prop:c2v42-window}, \ref{prop:c2v51-palinstrophy}).
\eqref{eq:c2v53-growth} is the enstrophy identity integrated, dropping
the palinstrophy term.
\end{proof}

\begin{definition}[Uniformly spread and travelling windows]
\label{def:c2v53-kinds}
A spread window (Corollary~\ref{cor:c2v48-collapse}) is \emph{uniformly
spread} if \(\int_{J_M}\delta(s)^2\dd s\to0\) along the subsequence, and
\emph{travelling} otherwise, that is, if \(\liminf\int_{J_M}\delta^2>0\)
while every cell's time-integrated enstrophy \(e_\alpha\) tends to zero.
\end{definition}

\begin{corollary}[No seeding on a uniformly spread window]
\label{cor:c2v53-no-seed}
On a uniformly spread window the enstrophy satisfies
\(\sup_{J_M}Y\le Y(-1)+o(1)\le M+o(1)\): no enstrophy growth occurs within
the window, and a spike after the window is seeded after it, during the
stretch without bounds.
\end{corollary}

\begin{proof}
\eqref{eq:c2v53-growth}.
\end{proof}

\begin{proposition}[A travelling window carries a moving concentration]
\label{prop:c2v53-travelling}
On a travelling window there are \(\eta>0\) and, for each \(j\), a set of
times \(S_j\subset J_M\) of measure at least \(\eta\ell_M\) on which
\(\delta(s)\ge\eta\), with \(\delta(s)=y_{\alpha(s)}(s)\) attained at a cell
\(\alpha(s)\); every fixed cell has \(\int_{J_M}y_\alpha\to0\), so the cell
\(\alpha(s)\) changes with \(s\): the time spent by the maximizing cell in
any fixed cell tends to zero, and the concentration of enstrophy of
size at least \(\eta\) visits at least \(\eta\ell_M/\max_\alpha\int_{J_M}y_\alpha\to\infty\)
distinct cells during the window.  Its centre map is not coherent in
the sense of Definition~\ref{def:c2v24-coherent} at the window scale.
\end{proposition}

\begin{proof}
If \(\int\delta^2\ge c>0\) with \(\delta\le2M\), then \(|\{\delta\ge\eta\}|\ge(c-\eta^2\ell_M)/(4M^2)\)
for small \(\eta\), which is at least \(\eta\ell_M\) for \(\eta\) small enough
(depending on \(c,M,\ell_M\)).  On \(\{\delta\ge\eta\}\) the maximizing cell
carries \(y\ge\eta\); a fixed cell carries \(y\ge\eta\) on a set of times of
measure at most \(e_\alpha/\eta\to0\), so the number of distinct maximizing
cells is at least \(\eta\ell_M/\max_\alpha(e_\alpha/\eta)\to\infty\).  A coherent
centre map moves by \(o(1)\) cells across a window; the maximizing cell
moves through a diverging number of cells.
\end{proof}

\begin{assessment}[Items (AC1)--(AC2)]
\label{ass:c2v53-AC}
A spread window is either uniformly weak, in which case nothing happens
in it (no stretching, no growth) and the next spike is seeded only
after the window's \(H^2\) control has lapsed, or it contains a
concentration of enstrophy of fixed size sweeping through a diverging
number of parabolic cells within one window.  The second kind is the
incoherent-centre configuration of V24 seen from the intermediate
class: a structure of parabolic size moving faster than the parabolic
speed.  The first cycle's Galilean remark (constant mode of the
defect, V8) and its satellite analysis do not bound the sweeping
speed; the series records that a travelling spread window has a
concentrated structure that the fixed-centre and coherent-centre
tools cannot follow, while a uniformly spread window has no
structure at all at the parabolic scale.  Seeding in exceptional cells
(the programme's hypothesis) does not occur on uniformly spread
windows; on travelling windows the seed is the moving structure.  No
global regularity claim is made.
\end{assessment}

% =============================================================================
\section{Integrated V53 conclusion and next acquisition order}
\label{sec:c2v53-conclusion}
% =============================================================================

\begin{theorem}[What V53 proves]
\label{thm:c2v53-conclusion}
Relative to the first cycle and V1--V52: the window stretching bound
\eqref{eq:c2v53-stretching} and growth bound \eqref{eq:c2v53-growth},
the absence of seeding on uniformly spread windows, and the moving
concentration on travelling windows.  The full Navier--Stokes
stability/global-regularity theorem remains unproved.
\end{theorem}

\begin{proof}
The cited results.
\end{proof}

\begin{programme}[V54 acquisition order]
\label{prog:c2v54}
\begin{enumerate}[label=\textup{(AC\arabic*)},leftmargin=4.8em]
\item Sweeping speed.  On a travelling window, bound the speed at which
      the enstrophy concentration moves through the cells by the local
      velocity: the enstrophy flux through a cell boundary is bounded by
      \(\sup_{\partial Q}|U|\) times the local enstrophy, so a concentration
      moving \(n\) cells in time \(\ell_M\) needs \(\sup|U|\gtrsim n/\ell_M\) along
      its path; combine with the vanishing energy density and the
      sup-norm bound on \(H^2\)-bounded cells (V52) to determine whether
      the path must lie in the exceptional cells, and record the
      resulting constraint on \(n\).
\item Prepare V55: the audit and the consolidation of V51--V54 as the
      window-structure chapter, with the position after fifty-five
      versions.
\item The next compulsory companion audit is V55.
\end{enumerate}
\end{programme}

\begin{assessment}[Position after V53]
\label{ass:c2v53-position}
Spread windows are uniformly weak (no stretching) or travelling (a
sweeping concentration); neither is excluded.  No full stability or
global-regularity claim is made.
\end{assessment}

% =============================================================================
\section*{Version V53 (second cycle) append-only record}
\addcontentsline{toc}{section}{Version V53 (second cycle) append-only record}
% =============================================================================

The complete V52 source of the second cycle was retained before this
continuation.  Its SHA--256 digest was
\begin{center}\scriptsize\ttfamily
c8a8a07f3a4dc92d1af1c3c547eed11bd772c21d96781e178e1e97fdedb5f120.
\end{center}
V53 proved the window stretching bound and distinguished uniformly
spread from travelling windows.  No full regularity claim is made.

% =============================================================================
\section{V54: there are no travelling windows}
\label{sec:c2v54-entry}
% =============================================================================

Item (AC1) of Programme~\ref{prog:c2v54} is carried out and closes the
travelling alternative of V53.  The local enstrophy identity on a cell
with a cutoff shows that the enstrophy of a cell whose neighbourhood has
bounded window \(H^2\) norm can increase over a time interval by at most
an explicit constant times the fourth root of the interval's length:
stretching, transport, diffusion, and both pressure terms are each
bounded on such a cell by powers of the interval's local enstrophy
integral and of the window budgets.  Hence a cell that carries a fixed
enstrophy at some time has carried at least half of it for a fixed
time before, so its window-integrated enstrophy is bounded below; on a
spread window, where every cell's window-integrated enstrophy tends to
zero, no cell with bounded neighbourhood can ever carry a fixed
enstrophy, and the cells with unbounded neighbourhood are boundedly
many and also have vanishing integrals, so the maximal local enstrophy
exceeds any fixed level only on a set of times of vanishing measure.
Every spread window is therefore uniformly spread, and by V53 no
enstrophy growth occurs within it.  No global regularity claim is made.
The next version is the compulsory companion audit.

% =============================================================================
\section{The local enstrophy influx bound}
\label{sec:c2v54-influx}
% =============================================================================

Let \(u\) be in the intermediate class with constant \(M\), \(U\), \(P\) the
rescaled velocity and pressure at a window scale, \(J_M\) the window,
\(Q_\alpha\) the unit cells, \(3Q_\alpha\) the concentric cube of side three,
\(\phi_\alpha\) a cutoff equal to \(1\) on \(Q_\alpha\) and supported in \(2Q_\alpha\),
and
\begin{equation}
 y_\alpha(s):=\int|\nabla U(s)|^2\phi_\alpha,\qquad
 H_\alpha:=\int_{J_M}\|U(s)\|_{H^2(3Q_\alpha)}^2\,\dd s .
 \label{eq:c2v54-notation}
\end{equation}
Call the cell \(\alpha\) \emph{\(H\)-regular} if \(H_\alpha\le H\).  Since
\(\sum_\alpha H_\alpha\le27\int_{J_M}\|U\|_{H^2}^2\le C_H(M,\nu)\) by the window
energy, enstrophy, and palinstrophy budgets, at most \(C_H/H\) cells are
not \(H\)-regular.

\begin{lemma}[Influx bound on an \(H\)-regular cell]
\label{lem:c2v54-influx}
There is \(C_\flat=C_\flat(M,\nu,H)\) such that for every \(H\)-regular cell
\(\alpha\) and every interval \([s_1,s_2]\subset J_M\) of length \(\tau\),
\begin{equation}
 y_\alpha(s_2)-y_\alpha(s_1)\ \le\ C_\flat\,\tau^{1/4} .
 \label{eq:c2v54-influx}
\end{equation}
\end{lemma}

\begin{proof}
Multiply the equation for \(\nabla U\) by \(\nabla U\phi_\alpha\):
\[
 \frac{\dd}{\dd s}\frac12y_\alpha=-\nu\int|\Delta U|^2\phi_\alpha-\nu\int\Delta U\cdot(\nabla U\nabla\phi_\alpha)
 -\int\partial_kU_i\partial_kU_j\partial_jU_i\phi_\alpha+\frac12\int|\nabla U|^2U\cdot\nabla\phi_\alpha
 +\int P\,\Delta U\cdot\nabla\phi_\alpha-\int\partial_iP\,\partial_kU_i\partial_k\phi_\alpha ,
\]
where the transport term used \(\operatorname{div}U=0\) and the two pressure
terms come from \(-\int\partial_kU_i\partial_k\partial_iP\phi_\alpha\) by two integrations
by parts with \(\operatorname{div}U=0\).  Drop the dissipation.  Over
\([s_1,s_2]\), with \(y:=\int_{s_1}^{s_2}\int_{2Q_\alpha}|\nabla U|^2\le2M\tau\) and
\(h:=\int_{s_1}^{s_2}\|U\|_{H^2(3Q_\alpha)}^2\le H\): the diffusion cutoff term is at
most \(C\nu h^{1/2}y^{1/2}\); the stretching term at most
\(C\int\|\nabla U\|_{L^3(2Q_\alpha)}^3\le C\int y_\alpha'^{3/4}h_\alpha'^{3/4}\le C(2M\tau)^{1/4}(2M)^{1/2}H^{3/4}\)
(cellwise Gagliardo--Nirenberg as in Theorem~\ref{thm:c2v53-stretching},
then H\"older in time with \(\int y_\alpha'^3\le(2M)^2\int y_\alpha'\)); the
transport term at most \(C\int\|U\|_{L^\infty(2Q_\alpha)}y_\alpha'\le C\sup_{J_M}\|U\|_{L^\infty(2Q_\alpha)}\cdot2M\tau\)
with \(\|U\|_{L^\infty(2Q_\alpha)}\le C\|U\|_{L^2(3Q_\alpha)}^{1/4}\|U\|_{H^2(3Q_\alpha)}^{3/4}\), so that
\(\int_{s_1}^{s_2}\|U\|_{L^\infty(2Q_\alpha)}\le C\sup\|U\|_{L^2(3Q_\alpha)}^{1/4}h^{3/8}\tau^{5/8}\le C(M)H^{3/8}\tau^{5/8}\)
by H\"older with exponents \(\frac83\) and \(\frac85\), giving at most
\(C(M,H)\tau^{5/8}\le C(M,H)\tau^{1/4}\); the first pressure term at
most \((\int\|P\|_{L^2(2Q_\alpha)}^2)^{1/2}h^{1/2}\le C(M)\tau^{1/2}H^{1/2}\) by the
pressure bounds of Proposition~\ref{prop:c2v47-pressure} (far part
bounded, local part bounded by \(\|U\|_{L^4(3Q_\alpha)}^2\le C(M)\) through
the \(H^1\) bound); the second pressure term at most
\(\int\|\nabla P\|_{L^2(2Q_\alpha)}y_\alpha'^{1/2}\) with
\(\|\nabla P\|_{L^2(2Q_\alpha)}\le C\|U\|_{L^\infty(3Q_\alpha)}\|\nabla U\|_{L^2(3Q_\alpha)}+C(M)\)
(local part by the Riesz bound, far part harmonic and bounded), giving
at most \(C(M,H)\tau^{1/2}\).  Each term is bounded by \(C(M,\nu,H)\tau^{1/4}\)
since \(\tau\le\ell_M\).
\end{proof}

% =============================================================================
\section{No travelling windows}
\label{sec:c2v54-no-travelling}
% =============================================================================

\begin{theorem}[Every spread window is uniformly spread]
\label{thm:c2v54-uniform}
Let \(\lambda_j\) be spread window scales (every cell's window-integrated
enstrophy \(e_\alpha^{(j)}\to0\) uniformly in \(\alpha\)).  Then
\(\int_{J_M}\delta_j(s)^2\,\dd s\to0\), where \(\delta_j(s)=\max_\alpha y_\alpha(s)\):
there are no travelling windows, and every spread window is uniformly
spread.  Consequently, on every spread window, \(\sup_{J_M}Y\le M+o(1)\)
and no enstrophy growth occurs within the window.
\end{theorem}

\begin{proof}
Fix \(H\) and \(\eta>0\).  Let \(\alpha\) be \(H\)-regular and suppose
\(y_\alpha(s)\ge\eta\) for some \(s\in J_M\).  By \eqref{eq:c2v54-influx}, for
\(s'\in[s-\tau_\eta,s]\cap J_M\) with \(\tau_\eta:=(\eta/(2C_\flat))^4\),
\(y_\alpha(s')\ge y_\alpha(s)-C_\flat(s-s')^{1/4}\ge\eta/2\); hence
\(e^{(j)}_\alpha\ge\int|\nabla U|^2\phi_\alpha\ge\frac\eta2\min(\tau_\eta,s+1)\).  For \(j\)
large, \(e^{(j)}_\alpha<\frac\eta2\tau_\eta\cdot\frac12\), so an \(H\)-regular cell
can satisfy \(y_\alpha(s)\ge\eta\) only for \(s\in[-1,-1+\tau_\eta/2)\), a set of
measure at most \(\tau_\eta/2\); on that initial set the same argument
forward (from \(s\) to \(s+\tau_\eta\)) gives \(e_\alpha\ge\frac\eta2\tau_\eta\) as
well, contradicting \(e_\alpha\to0\); so for \(j\) large no \(H\)-regular cell
reaches level \(\eta\) at any time.  The remaining cells number at most
\(C_H/H\), and for each of them \(|\{s:y_\alpha(s)\ge\eta\}|\le e^{(j)}_\alpha/\eta\to0\).
Hence \(|\{s\in J_M:\delta_j(s)\ge\eta\}|\le(C_H/H)\max_\alpha e^{(j)}_\alpha/\eta\to0\), and
\(\int_{J_M}\delta_j^2\le\eta^2\ell_M+(2M)^2|\{\delta_j\ge\eta\}|\), whose limsup is at
most \(\eta^2\ell_M\) for every \(\eta\): the integral tends to zero.  The
last statement is Corollary~\ref{cor:c2v53-no-seed}.
\end{proof}

\begin{assessment}[What has been closed]
\label{ass:c2v54-closed}
The travelling alternative of V53 is impossible: a concentration of
enstrophy cannot sweep through the parabolic cells of a spread window,
because entering a cell whose neighbourhood is under \(H^2\) control
takes a fixed time, during which the cell accumulates a fixed
window-integrated enstrophy, contradicting spreading; and the cells
outside \(H^2\) control are too few to host the concentration for a
positive fraction of the window.  The spreading configuration of the
intermediate class is therefore exactly the uniformly weak state of
V51--V53: enstrophy of order one at parabolic and larger scales, spread
over a diverging volume, with the maximal local enstrophy vanishing in
time-average and no growth within the window.  What remains open is
whether such a state can, after the window, grow a spike; the
enstrophy inequality permits it, and the influx bound of this version,
which needs the window's \(H^2\) budget, lapses with the window.
\end{assessment}

\begin{firewall}[Item (AC1): the sweeping-speed bound is subsumed]
\label{fw:c2v54-AC1}
The programme's sweeping-speed argument (a concentration moving \(n\)
cells in time \(\ell_M\) needs local velocity of order \(n/\ell_M\)) is
subsumed by Theorem~\ref{thm:c2v54-uniform}: the influx bound shows that
no fixed enstrophy can enter an \(H\)-regular cell in a short time by
any mechanism, transport included, so the question of the speed does
not arise.  The exceptional cells could host a concentration only for a
vanishing fraction of the window.
\end{firewall}

% =============================================================================
\section{Integrated V54 conclusion and next acquisition order}
\label{sec:c2v54-conclusion}
% =============================================================================

\begin{theorem}[What V54 proves]
\label{thm:c2v54-conclusion}
Relative to the first cycle and V1--V53, modulo (E2): the influx bound
\eqref{eq:c2v54-influx} on \(H\)-regular cells and the theorem that every
spread window is uniformly spread, with no enstrophy growth within it
(Theorem~\ref{thm:c2v54-uniform}).  The full Navier--Stokes
stability/global-regularity theorem remains unproved.
\end{theorem}

\begin{proof}
The cited results.
\end{proof}

\begin{programme}[V55 acquisition order, with compulsory companion audit]
\label{prog:c2v55}
\begin{enumerate}[label=\textup{(AC\arabic*)},leftmargin=4.8em]
\item Perform the compulsory review of the current SM and YM notes;
      record digests; import nothing numerical.
\item Consolidate V51--V54 as the window-structure chapter: the
      enstrophy band, the Fourier description of spread windows, the
      sharp enstrophy inequality and its optimal palinstrophy, the
      window stretching bound, the influx bound, and the theorem that
      spread windows are uniformly spread.
\item Record the position after fifty-five versions and the direction
      for V56--V60: the stretch after a uniformly spread window, where
      the \(H^2\) control lapses, and whether the influx bound can be
      continued past the window by a bootstrap on the maximal local
      enstrophy (the local enstrophy identity closes if the neighbourhood
      \(H^2\) norm is controlled by the local enstrophy through the
      local dissipation).
\item The next compulsory companion audit after V55 is V60.
\end{enumerate}
\end{programme}

\begin{assessment}[Position after V54]
\label{ass:c2v54-position}
Spread windows are uniformly spread; a concentration of enstrophy
cannot travel through the parabolic cells of a type-I window.  No full
stability or global-regularity claim is made.
\end{assessment}

% =============================================================================
\section*{Version V54 (second cycle) append-only record}
\addcontentsline{toc}{section}{Version V54 (second cycle) append-only record}
% =============================================================================

The complete V53 source of the second cycle was retained before this
continuation.  Its SHA--256 digest was
\begin{center}\scriptsize\ttfamily
8bcbb7c9fa894ab8a92727aaed29732966d4158ba5426a1b295dfedbbdee1bf5.
\end{center}
V54 proved the local enstrophy influx bound on \(H^2\)-regular cells
and that every spread window is uniformly spread.  No full regularity
claim is made.

% =============================================================================
\section{V55: compulsory companion audit, the window-structure chapter
consolidated, and the position after fifty-five versions}
\label{sec:c2v55-entry}
% =============================================================================

This is the fifty-fifth version of the second cycle.  The compulsory
review of the two companion programmes is performed first.  V51--V54
are then consolidated as the window-structure chapter of the
intermediate class.  The position after fifty-five versions is
recorded with the direction for V56--V60.  No global regularity claim
is made.  The next compulsory companion audit is V60.

% =============================================================================
\section{Compulsory V55 review of the current companion notes}
\label{sec:c2v55-companion-audit}
% =============================================================================

The current companion files in \texttt{latex\_notes} were inspected
read-only.  The frozen checkpoints are
\begin{align}
 &\texttt{Renewal\_SM\_Einstein\_Continuum\_Research\_Notes\_V2.tex},
 \notag\\[-2pt]
 &\qquad
 \texttt{24497db4ecd645e16d2b1e686c183af4597df94823be73d0d07f6838a176a5af},
 \label{eq:c2v55-SM-checkpoint}\\
 &\texttt{Renewal\_Yang\_Mills\_Mass\_Gap\_Research\_Notes\_V78.tex},
 \notag\\[-2pt]
 &\qquad
 \texttt{d9416f63f155cc90a99bbc6c684301368998ab361683bd9ca71b29782a83d224}.
 \label{eq:c2v55-YM-checkpoint}
\end{align}
The SM programme has restarted its numbering (its file is a new first
series, ``seventh series'', of 953 lines at V2); a YM file named
\(\ldots\)\texttt{V77.tex} with the same digest as \eqref{eq:c2v55-YM-checkpoint}
was also present.

\emph{SM V2 (new series).}  For a gapped Wilson kernel, a field
variation enters the overlap sign once between two resolvents and its
trace norm costs one inverse gap; a finite link or Higgs support gives a
finite-cutoff trace-class relative determinant connection; on a fixed
physical four-volume the raw bound retains the site count; the next
operation is a local heat/symbol subtraction.

\emph{YM V77.}  The YM programme expanded local child-star powers into
weighted Cayley trees, glued them to one global tree with exact
preservation of the Pr\"ufer monomial, identified the gluing fibre with
constrained connected edge partitions with a sharp fibre bound, proved
an exact global-tree representation, and closed its rooted-hypertree
conditions and every incidence-hypertree centre core.

\begin{theorem}[V55 companion-audit decision]
\label{thm:c2v55-companion-decision}
No SM V2 trace-norm or gap constant and no YM V77 tree or fibre bound is
imported.  Neither bears on the window-structure chapter or on the
gates.
\end{theorem}

\begin{proof}
The SM results concern overlap determinants of lattice Dirac operators;
the YM results concern tree expansions of a lattice clock.  None is a
Navier--Stokes statement.
\end{proof}

% =============================================================================
\section{The window-structure chapter, consolidated}
\label{sec:c2v55-consolidated}
% =============================================================================

\begin{theorem}[Window structure of the intermediate class, V51--V54]
\label{thm:c2v55-consolidated}
Let \(u\) be in the intermediate class with constant \(M\), with window
scales \(\lambda_j\), windows \(J_M\) of length \(\ell_M=3\nu^3/(8C_AM^2)\),
\(C_A=\frac{27}{16}C_B^4\), and unit cells.  Modulo (E2):
\begin{enumerate}[label=\textup{(\arabic*)},leftmargin=3.2em]
\item \emph{Band.}  \(\nu\int_{J_M}\|\Delta U\|_2^2\le4M\); on window average
      the enstrophy above wavenumber \(K\) is at most \(\frac{8C_AM^2\cdot4M}{3\nu^4K^2}\),
      and at least half of Leray's floor sits below the explicit
      \(K_M\); at a spread window the band carries fixed enstrophy and
      energy while the energy density and every local quantity tend to
      zero.
\item \emph{Sharp growth.}  \(Y'\le\frac{27C_B^4}{128\nu^3}Y^3\), attained at
      palinstrophy \((3C_B/(4\nu))^4Y^3\), which is within the window
      budget; delocalization does not reduce the available palinstrophy.
\item \emph{Stretching.}  The window-integrated stretching is at most
      \(C(2M)^{1/4}H_M^{3/4}(\int_{J_M}\delta^2)^{1/4}\) with \(\delta\) the maximal
      local enstrophy; uniformly spread windows have no enstrophy
      growth.
\item \emph{No travelling windows.}  On an \(H\)-regular cell the
      enstrophy increases by at most \(C_\flat\tau^{1/4}\) over any interval
      of length \(\tau\); hence every spread window is uniformly spread,
      \(\int_{J_M}\delta^2\to0\), \(\sup_{J_M}Y\le M+o(1)\), and a spike after a
      spread window is seeded only after the window.
\end{enumerate}
\end{theorem}

\begin{proof}
(1) Proposition~\ref{prop:c2v51-palinstrophy}, Theorem~\ref{thm:c2v51-band},
Proposition~\ref{prop:c2v51-spread}.  (2) Propositions~\ref{prop:c2v52-sharp},
\ref{prop:c2v52-budget}.  (3) Theorem~\ref{thm:c2v53-stretching},
Corollary~\ref{cor:c2v53-no-seed}.  (4) Lemma~\ref{lem:c2v54-influx},
Theorem~\ref{thm:c2v54-uniform}.
\end{proof}

% =============================================================================
\section{Position after fifty-five versions and direction}
\label{sec:c2v55-position}
% =============================================================================

\begin{assessment}[Position after fifty-five versions of the second cycle]
\label{ass:c2v55-position-fifty-five}
The intermediate class is now described at its window scales in full:
a type-I window is either concentrated (nonzero strong limits at
boundedly many centres, with the Gaussian ledgers bounded and the
window flux signed for small sequence constant) or uniformly spread (a
weak delocalized state in an explicit band of parabolic and larger
scales, with no growth and no travelling structure), and the transition
to the next spike happens strictly after the window, in the stretch
where the \(H^2\) control lapses.  The theorem that spread windows are
uniformly spread is the first exclusion of a configuration in the
spreading chapter.  For V56--V60 the direction is the stretch after a
uniformly spread window: whether the influx bound can be continued
past the window by a bootstrap in which the neighbourhood \(H^2\) norm
is controlled by the local enstrophy through the local dissipation,
so that a uniformly weak state cannot grow a spike at all within a
time comparable to the window, and, if so, how this interacts with the
backward cone's parabolic delay and the gap summability.  No global
regularity claim is made.
\end{assessment}

% =============================================================================
\section{Integrated V55 conclusion and next acquisition order}
\label{sec:c2v55-conclusion}
% =============================================================================

\begin{theorem}[What V55 establishes]
\label{thm:c2v55-conclusion}
Relative to the first cycle and V1--V54: the companion-audit decision
(Theorem~\ref{thm:c2v55-companion-decision}) and the consolidated
window-structure chapter (Theorem~\ref{thm:c2v55-consolidated}).  The
full Navier--Stokes stability/global-regularity theorem remains
unproved.
\end{theorem}

\begin{proof}
The cited results.
\end{proof}

\begin{programme}[V56 acquisition order]
\label{prog:c2v56}
\begin{enumerate}[label=\textup{(AC\arabic*)},leftmargin=4.8em]
\item Bootstrap after the window.  For \(s\) beyond the window, with
      \(\delta(s)\) the maximal local enstrophy and \(Y(s)\) the global
      enstrophy, derive from the local enstrophy identity a
      differential inequality for \(\delta\) in which the stretching is
      bounded by \(C\delta^{1/2}Y^{1/4}\mathcal E_2^{3/4}\)-type terms and the
      palinstrophy is controlled by the dissipation term of the same
      identity; determine whether a closed inequality
      \(\delta'\le F(\delta,Y)\) with \(F(\delta,Y)\le C\delta^{a}Y^{b}\), \(a>1\), follows, so
      that a uniformly weak state (\(\delta\) small) stays weak for a time
      inversely proportional to a power of \(\delta\), longer than the
      window.
\item Combine with the global inequality \(Y'\le C_A\nu^{-3}Y^3\) to
      determine whether a spike (unbounded \(Y\)) after a uniformly
      spread window requires \(\delta\) to reach order one first, and
      whether that requires a time exceeding the cone's parabolic delay.
\item The next compulsory companion audit is V60.
\end{enumerate}
\end{programme}

\begin{assessment}[Position after V55]
\label{ass:c2v55-position}
The window-structure chapter is consolidated; spread windows are
uniformly spread.  No full stability or global-regularity claim is
made.
\end{assessment}

% =============================================================================
\section*{Version V55 (second cycle) append-only record}
\addcontentsline{toc}{section}{Version V55 (second cycle) append-only record}
% =============================================================================

The complete V54 source of the second cycle was retained before this
continuation.  Its SHA--256 digest was
\begin{center}\scriptsize\ttfamily
c868b62be1ab2942dbcaf15e29edeec3e6884a6f08deb894a93d6a62a2dc075c.
\end{center}
V55 performed the compulsory companion audit and consolidated the
window-structure chapter.  No full regularity claim is made.

% =============================================================================
\section{V56: the post-window bootstrap and the delay of a spike after a
uniformly spread window}
\label{sec:c2v56-entry}
% =============================================================================

Items (AC1) and (AC2) of Programme~\ref{prog:c2v56} are carried out.  A
closed differential system is derived for the global enstrophy and the
maximal local enstrophy of the rescaled solution, valid at all times
and without any \(H^2\) control: the local enstrophy identity on a cell,
with its own dissipation absorbing the stretching, transport, and local
pressure terms through Gagliardo--Nirenberg and Young inequalities,
bounds the growth of the square root of the maximal local enstrophy by
a constant times the global enstrophy plus lower-order terms, the
far-pressure gradient being controlled by the global enstrophy; and
the global enstrophy identity, with the stretching bounded through the
maximal local enstrophy, bounds the growth rate of the global
enstrophy by a constant times the enstrophy times a power of the
maximal local enstrophy.  Consequently a uniformly weak state cannot
spike before a time of the order of the inverse of the global enstrophy
in rescaled units, an explicit multiple of the parabolic scale that is
longer than the backward cone's delay for large sequence constants.
No global regularity claim is made.  The next compulsory companion
audit is V60.

% =============================================================================
\section{The local inequality without \(H^2\) control}
\label{sec:c2v56-local}
% =============================================================================

Let \(U\), \(P\) be a smooth mean-zero solution on a torus of side \(L\ge3\)
(the rescaled solution at any scale, for any time before the singular
time), \(Q_\alpha\) unit cells, \(\phi_\alpha\) cutoffs as in V54,
\(y_\alpha=\int|\nabla U|^2\phi_\alpha\), \(y'_\alpha:=\int_{3Q_\alpha}|\nabla U|^2\le27\delta\) with
\(\delta:=\max_\alpha y_\alpha\), \(Y=\|\nabla U\|_2^2\), \(\mathcal E_\alpha:=\int|\Delta U|^2\phi_\alpha\),
and \(K:=C_S^2Y\) the bound on \(\|U\|_{L^2(3Q)}^2\le|3Q|^{2/3}\|U\|_6^2\le C_S^2Y\cdot9\)
(absorbed into constants).

\begin{proposition}[Closed local enstrophy inequality]
\label{prop:c2v56-local}
There is an absolute constant \(C_0\) such that for every cell \(\alpha\),
\begin{equation}
 \frac{\dd}{\dd s}y_\alpha\ \le\ -\frac\nu2\mathcal E_\alpha+C_0\Bigl[\nu^{-3}y'^3_\alpha+\nu\,y'_\alpha+\nu^{-3}y'^{5/3}_\alpha
 +\nu^{-1}Y^{1/2}y'^{5/4}_\alpha+\nu^{-1/3}Y^{2/3}y'^{4/3}_\alpha+Y\,y'^{1/2}_\alpha\Bigr],
 \label{eq:c2v56-local}
\end{equation}
where the powers of \(\nu\) are those of the Young inequalities used.
\end{proposition}

\begin{proof}
Start from the local enstrophy identity of Lemma~\ref{lem:c2v54-influx}
with the pressure terms rewritten with \(P-c_\alpha\) (\(c_\alpha\) the mean of
\(P\) over \(2Q_\alpha\)), which is allowed since the cutoff terms involve
\(\nabla\phi_\alpha\) only through integrals of \(\nabla U\), and use
Poincar\'e \(\|P-c_\alpha\|_{L^2(2Q_\alpha)}\le C\|\nabla P\|_{L^2(2Q_\alpha)}\).  Each
term is bounded with the local \(H^1\) norm \(\|\nabla U\|_{H^1(2Q_\alpha)}^2\le y'_\alpha+\mathcal E'_\alpha\),
\(\mathcal E'_\alpha:=\|\Delta U\|_{L^2(2Q_\alpha)}^2\), and the dissipation
\(\frac\nu2\mathcal E_\alpha\) of the cell absorbs the \(\mathcal E'\)-dependence
after summing the inequality over a bounded number of overlapping
cutoffs (the standard covering: replace \(\phi_\alpha\) by a partition of
unity subordinate to the cells, so that \(\sum_\alpha\mathcal E'_\alpha\le8\sum_\alpha\mathcal E_\alpha\);
the inequality \eqref{eq:c2v56-local} is then understood for the
maximizing cell with \(\mathcal E_\alpha\) replaced by the dissipation of its
neighbourhood, which only weakens it).  Stretching:
\(\|\nabla U\|_{L^3(2Q)}^3\le Cy'^{3/4}(y'+\mathcal E')^{3/4}\le\frac\nu8\mathcal E'+C(\nu^{-3}y'^3+y'^{3/2})\).
Transport: \(\int|\nabla U|^2|U||\nabla\phi|\le\|U\|_{L^6(2Q)}\|\nabla U\|_{L^{12/5}(2Q)}^2\le CY^{1/2}y'^{3/4}(y'+\mathcal E')^{1/4}\le\frac\nu8\mathcal E'+C\nu^{-1/3}Y^{2/3}y'+CY^{1/2}y'\)
(Young with exponents \(4\) and \(\frac43\)); the displayed
\(\nu^{-1}Y^{1/2}y'^{5/4}\) and \(\nu^{-1/3}Y^{2/3}y'^{4/3}\) collect this and the
local-pressure gradient term
\(\|\nabla P^{\rm loc}\|_{L^2(2Q)}y'^{1/2}\le C\|U\|_{L^6(3Q)}\|\nabla U\|_{L^3(3Q)}y'^{1/2}\le CY^{1/2}y'^{3/4}(y'+\mathcal E')^{1/4}y'^{1/2}\),
absorbed in the same way.  Far pressure:
\(\|\nabla P^{\rm far}\|_{L^\infty(2Q)}\le C\operatorname{osc}_{3Q}P^{\rm far}\le CC_S^2Y\)
by the dyadic argument of Proposition~\ref{prop:c2v47-pressure}(i)
(\(\|U\|_{L^2(2^kQ)}^2\le C4^kC_S^2Y\)), giving the term \(CYy'^{1/2}\).
Diffusion cutoff: \(\nu\int|\Delta U||\nabla U||\nabla\phi|\le\frac\nu8\mathcal E'+C\nu y'\).
The pressure term with \(P-c\) against \(\nabla U\nabla^2\phi\) is bounded by
\(\|\nabla P\|_{L^2(2Q)}y'^{1/2}\), already counted.  Collect, with
\(y'^{3/2}\le y'+y'^3\) absorbed.
\end{proof}

\begin{corollary}[Growth of the maximal local enstrophy]
\label{cor:c2v56-delta}
For \(\delta(s)=\max_\alpha y_\alpha(s)\) (absolutely continuous, with derivative
equal a.e.\ to that of a maximizing cell),
\begin{equation}
 \frac{\dd}{\dd s}\delta^{1/2}\ \le\ C_1\Bigl[Y+\nu\delta^{1/2}+\nu^{-3}\delta^{5/2}+\nu^{-3}\delta^{7/6}+\nu^{-1}Y^{1/2}\delta^{3/4}+\nu^{-1/3}Y^{2/3}\delta^{5/6}\Bigr]
 \label{eq:c2v56-delta}
\end{equation}
with \(C_1\) absolute (the factor \(27\) from \(y'\le27\delta\) absorbed).  In
particular, as long as \(\delta\le1\) and \(Y\ge1\),
\(\frac{\dd}{\dd s}\delta^{1/2}\le C_2(Y+\nu^{-3})\) with \(C_2\) absolute.
\end{corollary}

\begin{proof}
Divide \eqref{eq:c2v56-local} by \(2\delta^{1/2}\) at a maximizing cell,
dropping the dissipation; each term \(y'^{a}\) becomes \(\delta^{a-1/2}\); the
simplified form uses \(\delta\le1\le Y\).
\end{proof}

% =============================================================================
\section{The global inequality with the maximal local enstrophy}
\label{sec:c2v56-global}
% =============================================================================

\begin{proposition}[Global growth controlled by the maximal local enstrophy]
\label{prop:c2v56-global}
\begin{equation}
 \frac{\dd}{\dd s}Y\ \le\ -\nu\,\mathcal E_2+C_3\bigl(\nu^{-3}\delta^2+\delta^{1/2}\bigr)\,Y ,
 \label{eq:c2v56-global}
\end{equation}
with \(C_3\) absolute.  Hence as long as \(\delta\le\eta\),
\(Y(s)\le Y(s_0)\exp\bigl(C_3(\nu^{-3}\eta^2+\eta^{1/2})(s-s_0)\bigr)\).
\end{proposition}

\begin{proof}
The stretching is at most \(\sum_\alpha\|\nabla U\|_{L^3(Q_\alpha)}^3\le C\sum_\alpha y'^{3/4}_\alpha h'^{3/4}_\alpha\)
with \(h'_\alpha=\|\nabla U\|_{H^1(2Q_\alpha)}^2\); H\"older over the cells with
exponents \(4\), \(\frac43\) and \(\sum_\alpha y'^3_\alpha\le(27\delta)^2\sum_\alpha y'_\alpha\le C\delta^2Y\),
\(\sum_\alpha h'_\alpha\le C(Y+\mathcal E_2)\), give stretching
\(\le C\delta^{1/2}Y^{1/4}(Y+\mathcal E_2)^{3/4}\le\frac\nu2\mathcal E_2+C\nu^{-3}\delta^2Y+C\delta^{1/2}Y\)
by Young (exponents \(4\), \(\frac43\)) and \(Y^{1/4}Y^{3/4}=Y\).  Insert
into \(Y'=-2\nu\mathcal E_2+2\cdot\)stretching.
\end{proof}

% =============================================================================
\section{The delay of a spike after a uniformly spread window}
\label{sec:c2v56-delay}
% =============================================================================

\begin{theorem}[Delay theorem]
\label{thm:c2v56-delay}
Let \(u\) be in the intermediate class with constant \(M\ge1\), and let
\(s_0\) be a time (in the rescaled variables at a window scale
\(\lambda_j\)) at which \(Y(s_0)\le2M\) and \(\delta(s_0)\le\delta_0\le\frac14\).  Then for
\begin{equation}
 s-s_0\ \le\ T_{\rm delay}:=\frac{1}{4C_2\,(2M+\nu^{-3})}\qquad\text{(and while }s<0),
 \label{eq:c2v56-delay}
\end{equation}
one has \(\delta(s)\le1\) and \(Y(s)\le2Me^{C_3(\nu^{-3}+1)T_{\rm delay}}\le C_4M\):
no spike occurs, and the enstrophy stays bounded by an explicit
multiple of the sequence constant, for a rescaled time \(T_{\rm delay}\)
of order \(1/M\).  In the original variables the delay after the window
time \(t_j\) is \(T_{\rm delay}\lambda_j^2\), which exceeds the backward cone's
delay \(\nu^3\lambda_j^2/(32C_1M^2)\) by the factor of order \(M\).
Consequently, after a uniformly spread window (\(\delta_0\to0\)), the
window's enstrophy control persists for an additional rescaled time of
order \(1/M\) beyond the window length \(\ell_M\) of order \(1/M^2\).
\end{theorem}

\begin{proof}
Bootstrap on \([s_0,s]\): while \(\delta\le1\) and \(Y\le C_4M\),
\eqref{eq:c2v56-delta} gives \(\delta^{1/2}(s)\le\delta_0^{1/2}+C_2(C_4M+\nu^{-3})(s-s_0)\),
and \eqref{eq:c2v56-global} gives \(Y(s)\le2Me^{C_3(\nu^{-3}+1)(s-s_0)}\); with
\(T_{\rm delay}\) as stated (and \(C_4\) chosen as \(2e^{C_3(\nu^{-3}+1)/(4C_2)}\),
consistent), \(\delta^{1/2}\le\frac12+\frac14<1\) and the bounds close by
continuity.  The comparison with the cone is \eqref{eq:c2v46-distance}.
\end{proof}

\begin{assessment}[Items (AC1)--(AC2)]
\label{ass:c2v56-AC}
The bootstrap closes: a uniformly weak state cannot develop a spike
until its maximal local enstrophy has grown to order one, and the
growth of the maximal local enstrophy is driven by the global enstrophy
(through the far pressure and the diffusion into the cell) at a rate
that keeps it small for a time of order \(1/M\).  This extends the
type-I window of a spread state from length \(\ell_M\sim M^{-2}\) to
length of order \(M^{-1}\) in the rescaled units, with an explicit
constant, and it is the first place where the post-window stretch is
constrained by a local mechanism.  It does not exclude spikes: after
the delay the maximal local enstrophy may be of order one, the
inequalities revert to the global cubic bound, and a spike may follow.
What it says is that between a uniformly spread window and the next
spike there is a quiet stretch of parabolic duration of order \(1/M\),
during which the solution remains a strong solution with enstrophy at
most \(C_4M\); on this stretch all the tools of the type-I windows apply,
so the windows of the intermediate class can be taken of length
\(T_{\rm delay}\) after spread windows.  No global regularity claim is
made.
\end{assessment}

% =============================================================================
\section{Integrated V56 conclusion and next acquisition order}
\label{sec:c2v56-conclusion}
% =============================================================================

\begin{theorem}[What V56 proves]
\label{thm:c2v56-conclusion}
Relative to the first cycle and V1--V55: the closed local inequality
\eqref{eq:c2v56-local}, the growth bound \eqref{eq:c2v56-delta} for the
maximal local enstrophy, the global inequality \eqref{eq:c2v56-global},
and the delay theorem \eqref{eq:c2v56-delay}.  The full Navier--Stokes
stability/global-regularity theorem remains unproved.
\end{theorem}

\begin{proof}
The cited results.
\end{proof}

\begin{programme}[V57 acquisition order]
\label{prog:c2v57}
\begin{enumerate}[label=\textup{(AC\arabic*)},leftmargin=4.8em]
\item Iterate.  At the end of the delay stretch the state has enstrophy
      at most \(C_4M\) and maximal local enstrophy at most one; apply the
      window trichotomy at the new scale: if the new window is spread,
      the delay theorem applies again with \(M\) replaced by \(C_4M\);
      determine whether the iteration produces, from a spread window, a
      sequence of quiet stretches whose total parabolic duration
      diverges (excluding a spike at all), or whether the growth of the
      constant \(M\to C_4M\to C_4^2M\ldots\) makes the stretches shrink
      summably (allowing a spike after a finite time).
\item Compute the two possibilities explicitly: with \(T_k\sim1/(C_4^kM)\)
      the total delay \(\sum_kT_k\) converges, so the iteration alone does
      not exclude a spike; determine whether the maximal local enstrophy
      bound \(\delta\le1\) can be maintained across the iteration (it is the
      global enstrophy that grows, not \(\delta\), in the bootstrap), in
      which case the global inequality \eqref{eq:c2v56-global} gives
      exponential rather than cubic growth and no finite-time spike
      while \(\delta\le1\).
\item The next compulsory companion audit is V60.
\end{enumerate}
\end{programme}

\begin{assessment}[Position after V56]
\label{ass:c2v56-position}
A uniformly spread window is followed by a quiet stretch of parabolic
duration of order \(1/M\), during which no spike can occur.  No full
stability or global-regularity claim is made.
\end{assessment}

% =============================================================================
\section*{Version V56 (second cycle) append-only record}
\addcontentsline{toc}{section}{Version V56 (second cycle) append-only record}
% =============================================================================

The complete V55 source of the second cycle was retained before this
continuation.  Its SHA--256 digest was
\begin{center}\scriptsize\ttfamily
e46b3ceaaa668cea75543823746a283a38f0b5f39ad523452cf000ff7ec065f2.
\end{center}
V56 proved the closed local enstrophy inequality, the growth bound for
the maximal local enstrophy, the global inequality controlled by it,
and the delay theorem.  No full regularity claim is made.

% =============================================================================
\section{V57: the cell trap and the exclusion of spread windows}
\label{sec:c2v57-entry}
% =============================================================================

Items (AC1) and (AC2) of Programme~\ref{prog:c2v57} are carried out, and
the answer is stronger than the programme anticipated.  The far
pressure gradient on a cell is bounded by the maximal energy of a unit
cell (not by the global enstrophy), the local enstrophy of a cell is
bounded by the geometric mean of its energy and its dissipation, and
the two facts together turn the local enstrophy inequality of V56 into
a trap: a weighted maximum of the cell enstrophies decreases whenever
it exceeds a fixed multiple of the weighted maximum of the cell
energies, while the weighted maximum of the cell energies grows at
most exponentially at a rate depending only on the viscosity and the
sequence constant.  Consequently a state whose energy per parabolic
cell is small enough stays weak for a rescaled time that diverges as
the cell energy tends to zero, longer than the time remaining before
the singular time, and its global enstrophy stays bounded up to the
singular time, which contradicts the singularity.  The theorem that
follows is that the intermediate class has no spread windows: every
window is concentrated, with a cell of energy at least an explicit
function of the viscosity and the sequence constant.  The V56 delay
theorem is thereby superseded for spread windows.  No global
regularity claim is made.  The next compulsory companion audit is V60.

\begin{remark}[Two constants named \(C_1\)]
\label{rem:c2v57-constants}
In Theorem~\ref{thm:c2v56-delay} the cone delay
\(\nu^3\lambda_j^2/(32C_1M^2)\) uses the first-cycle cone constant \(C_1\) of
\eqref{eq:c2v46-distance}, while \(C_1\) in \eqref{eq:c2v56-delta} is the
absolute constant of the corollary.  They are distinct; the comparison
in the theorem is between \(T_{\rm delay}\) and the cone delay with the
cone constant.
\end{remark}

% =============================================================================
\section{Cells, partitions of unity, and weighted maxima}
\label{sec:c2v57-setup}
% =============================================================================

Throughout, \(U\), \(P\) is a smooth mean-zero solution on the torus
\(\mathbb T_L\) of side \(L\ge8\) (the rescaled solution \(U_j\) on
\([-1,0)\), \(L=2\pi/\lambda_j\)).  The cells are the unit cubes \(Q_\alpha\)
centred at the points \(\alpha\) of the integer lattice of the torus
(\(L\) is taken an integer by a harmless adjustment of \(\lambda_j\)), and
\(rQ_\alpha\) is the concentric cube of side \(r\).  Let \(\theta\in C_c^\infty(-1,1)\),
\(\theta\ge0\), with \(\sum_{n\in\mathbb Z}\theta(t-n)=1\), and
\(\phi_\alpha(x):=\theta(x_1-\alpha_1)\theta(x_2-\alpha_2)\theta(x_3-\alpha_3)\), so that
\(\phi_\alpha\in C_c^\infty(2Q_\alpha)\), \(0\le\phi_\alpha\le1\), and \(\sum_\alpha\phi_\alpha=1\).
The lattice distance is \(|\alpha-\beta|:=\max_i|\alpha_i-\beta_i|\) (on the
torus, the periodic distance).  The cell quantities are
\[
 e_\alpha:=\int|U|^2\phi_\alpha,\qquad y_\alpha:=\int|\nabla U|^2\phi_\alpha,\qquad
 d_\alpha:=\int|\Delta U|^2\phi_\alpha ,
\]
so that \(\sum_\alpha e_\alpha=\|U\|_2^2\), \(\sum_\alpha y_\alpha=Y\), \(\sum_\alpha d_\alpha=\|\Delta U\|_2^2\),
and for a cube \(rQ_\alpha\) with \(r\le5\) one has
\(\|U\|_{L^2(rQ_\alpha)}^2\le\sum_{|\beta-\alpha|\le2}e_\beta\), and likewise for the
other two quantities, since \(5Q_\alpha\) is covered by the supports of
the \(\phi_\beta\) with \(|\beta-\alpha|\le2\).  Fix \(\kappa\in(0,1)\) (say
\(\kappa=\frac12\)) and define the weighted quantities and their maxima
\begin{equation}
 w_\alpha:=\sum_\beta\kappa^{|\beta-\alpha|}y_\beta,\qquad
 \varepsilon_\alpha:=\sum_\beta\kappa^{|\beta-\alpha|}e_\beta,\qquad
 W:=\max_\alpha w_\alpha,\qquad \varepsilon:=\max_\alpha\varepsilon_\alpha .
 \label{eq:c2v57-weighted}
\end{equation}
Since the number of \(\beta\) at distance \(k\ge1\) from \(\alpha\) is at most
\(26k^2\), \(S_\kappa:=\sum_\beta\kappa^{|\beta-\alpha|}\le1+26\sum_kk^2\kappa^k<\infty\), and
\begin{equation}
 \max_\alpha y_\alpha\le W\le\min(Y,\,S_\kappa\max_\alpha y_\alpha),\qquad
 \max_\alpha e_\alpha\le\varepsilon\le S_\kappa\max_\alpha e_\alpha .
 \label{eq:c2v57-comparison}
\end{equation}
The maximal local enstrophy \(\delta=\max_\alpha y_\alpha\) of V53--V56 (there
defined with \(\int_{2Q_\alpha}\) or with a cutoff; the versions are
equivalent up to the factor \(27\)) is thus at most \(W\).  Both \(W\)
and \(\varepsilon\) are maxima of finitely many smooth functions of \(s\),
hence Lipschitz, with \(W'(s)=w'_{\alpha}(s)\) for almost every \(s\) and any
\(\alpha\) attaining the maximum at \(s\) (Danskin), and likewise for
\(\varepsilon\).

% =============================================================================
\section{Two lemmas: the far pressure at small cell energy, and the
enstrophy--energy interpolation}
\label{sec:c2v57-lemmas}
% =============================================================================

For a cell \(\alpha\) let \(\chi_\alpha\in C_c^\infty(4Q_\alpha)\) be a cutoff equal to \(1\)
on \(3Q_\alpha\), and \(P=P^{\rm loc}_\alpha+P^{\rm far}_\alpha\),
\(P^{\rm loc}_\alpha:=\partial_i\partial_j(-\Delta)^{-1}(\chi_\alpha U_iU_j)\) (on the torus, the
zero Fourier mode is killed by \(\partial_i\partial_j\)).

\begin{lemma}[Far pressure gradient at small cell energy]
\label{lem:c2v57-far}
There is an absolute constant \(C_{\rm far}'\) such that for every cell
\[
 \|\nabla P^{\rm far}_\alpha\|_{L^\infty(2Q_\alpha)}\ \le\ C_{\rm far}'\max_\beta e_\beta\ \le\ C_{\rm far}'\varepsilon .
\]
\end{lemma}

\begin{proof}
On \(2Q_\alpha\), \(\nabla P^{\rm far}_\alpha=K_L*((1-\chi_\alpha)U\otimes U)\) with
\(K_L\) the periodization of the kernel \(K=\nabla\partial_i\partial_j\Gamma\) of
\(\nabla\partial_i\partial_j(-\Delta)^{-1}\) on \(\R^3\), which is smooth away from the
origin and homogeneous of degree \(-4\): \(K_L(x)=\sum_{n\in\mathbb Z^3}K(x+nL)\),
the sum over \(n\neq0\) converging absolutely and uniformly for
\(|x|_\infty\le L/2\) with \(\sum_{n\ne0}|K(x+nL)|\le C\sum_{n\ne0}|nL|^{-4}\le CL^{-4}\)
(the periodization of a kernel homogeneous of degree \(-4\) in three
dimensions has Fourier coefficients equal to the symbol of the
operator at the lattice frequencies, which identifies it with the
torus kernel).  Hence \(|K_L(x)|\le C|x|^{-4}+CL^{-4}\) for
\(|x|_\infty\le L/2\).  Since \(1-\chi_\alpha\) vanishes on \(3Q_\alpha\), for
\(x\in2Q_\alpha\) the integrand is supported at distance at least \(\frac12\)
from \(x\).  Decompose the complement of \(2Q_\alpha\) into the dyadic
shells \(S_k:=2^{k+1}Q_\alpha\setminus2^kQ_\alpha\), \(k\ge1\), up to the torus: for
\(x\in2Q_\alpha\) and \(y\in S_k\) one has \(|x-y|\ge2^{k-2}\) for \(k\ge2\) and
\(|x-y|\ge\frac12\) on the support for \(k=1\); \(2^{k+1}Q_\alpha\) is covered by
\(8^{k+1}\) cells, so \(\int_{S_k}|U|^2\le8^{k+1}\max_\beta e_\beta\); the \(k\)-th
shell contributes at most \(C2^{-4k}8^{k+1}\max_\beta e_\beta=C'2^{-k}\max_\beta e_\beta\),
summable in \(k\).  The \(L^{-4}\) part of the kernel contributes at most
\(CL^{-4}\int_{\mathbb T_L}|U|^2=CL^{-4}\sum_\beta e_\beta\le CL^{-4}L^3\max_\beta e_\beta\le C\max_\beta e_\beta\),
the torus having \(L^3\) cells.
\end{proof}

\begin{lemma}[Interpolation: local enstrophy is bounded by the geometric
mean of the local energy and the local dissipation]
\label{lem:c2v57-interp}
For every cell \(\gamma\),
\begin{equation}
 y_\gamma^2\ \le\ 2\Bigl(\sum_{|\beta-\gamma|\le1}e_\beta\Bigr)\Bigl(\sum_{|\beta-\gamma|\le1}d_\beta+C_\theta^2\sum_{|\beta-\gamma|\le1}y_\beta\Bigr),
 \qquad C_\theta:=3\|\theta'\|_\infty .
 \label{eq:c2v57-interp}
\end{equation}
\end{lemma}

\begin{proof}
\(y_\gamma=\int\partial_kU_i\partial_kU_i\phi_\gamma=-\int U_i\Delta U_i\phi_\gamma-\int U_i\partial_kU_i\partial_k\phi_\gamma\le\|U\|_{L^2(2Q_\gamma)}(\|\Delta U\|_{L^2(2Q_\gamma)}+\|\nabla\phi_\gamma\|_\infty\|\nabla U\|_{L^2(2Q_\gamma)})\),
then \((a+b)^2\le2a^2+2b^2\) and the covering of \(2Q_\gamma\) by the
supports of the \(\phi_\beta\), \(|\beta-\gamma|\le1\), together with
\(\sum_\beta\phi_\beta=1\), which gives \(\|f\|_{L^2(2Q_\gamma)}^2\le\sum_{|\beta-\gamma|\le1}\int|f|^2\phi_\beta\)
for every \(f\).
\end{proof}

% =============================================================================
\section{The weighted local inequalities}
\label{sec:c2v57-weighted}
% =============================================================================

\begin{proposition}[Weighted local enstrophy inequality]
\label{prop:c2v57-enstrophy}
Let \(\Lambda_0\ge1\) and suppose \(W(s)\le\Lambda_0\) and \(\varepsilon(s)\le1\).  There are
\(a=a(\nu,\kappa)>0\) and \(b=b(\nu,\kappa,\Lambda_0)\) such that, at almost
every such \(s\),
\begin{equation}
 W'\ \le\ -\,a\,\frac{W^2}{\varepsilon}+b\,(W+\varepsilon) .
 \label{eq:c2v57-W}
\end{equation}
\end{proposition}

\begin{proof}
\emph{Step 1: the cell inequality.}  Fix a cell \(\alpha\) and write
\(N_\alpha:=\{\beta:|\beta-\alpha|\le2\}\), \(y^N:=\sum_{\beta\in N_\alpha}y_\beta\),
\(e^N:=\sum_{\beta\in N_\alpha}e_\beta\), \(d^N:=\sum_{\beta\in N_\alpha}d_\beta\), so that the
\(L^2\) norms of \(U\), \(\nabla U\), \(\Delta U\) on any \(rQ_\alpha\), \(r\le5\), are
bounded by \((e^N)^{1/2}\), \((y^N)^{1/2}\), \((d^N)^{1/2}\).  By the interior
estimate \(\|\nabla^2U\|_{L^2(4Q_\alpha)}^2\le C(\|\Delta U\|_{L^2(5Q_\alpha)}^2+\|\nabla U\|_{L^2(5Q_\alpha)}^2)\le C(d^N+y^N)\),
the local Sobolev and Gagliardo--Nirenberg inequalities on the cubes
\(rQ_\alpha\), \(r\le4\), read \(\|U\|_{L^6}\le C(e^N+y^N)^{1/2}\),
\(\|U\|_{L^4}^2\le C(e^N+y^N)\), \(\|U\|_{L^3}^3\le C(e^N)^{3/4}(e^N+y^N)^{3/4}\),
\(\|\nabla U\|_{L^3}^3\le C(y^N)^{3/4}(y^N+d^N)^{3/4}\),
\(\|\nabla U\|_{L^{12/5}}^2\le C(y^N)^{3/4}(y^N+d^N)^{1/4}\).  Start from the
identity of Lemma~\ref{lem:c2v54-influx} with \(P\) replaced by
\(P-c_\alpha\), \(c_\alpha:=P^{\rm far}_\alpha(\alpha)\), in the two pressure terms
(allowed: \(\nabla c_\alpha=0\), and
\(\int\Delta U\cdot\nabla\phi_\alpha=\int U\cdot\nabla\Delta\phi_\alpha=-\int(\operatorname{div}U)\Delta\phi_\alpha=0\);
moreover \(\int P\Delta U\cdot\nabla\phi_\alpha=-\int\nabla P\cdot(\nabla U\nabla\phi_\alpha)-\int P\,\partial_kU_i\partial_k\partial_i\phi_\alpha\)
with \(\int\partial_kU_i\partial_{ki}\phi_\alpha=0\) by \(\operatorname{div}U=0\), so both
pressure terms are integrals of \(\nabla P\) against \(\nabla U\) times
derivatives of \(\phi_\alpha\), and the constant enters only through
\(\nabla P\)).  Each term is bounded, with a Young
parameter \(\eta\in(0,1)\) to be fixed:
\begin{itemize}
\item diffusion cutoff: \(\nu\int|\Delta U||\nabla U||\nabla\phi_\alpha|\le\eta\nu d^N+C_\eta\nu y^N\);
\item stretching: \(\int|\nabla U|^3\phi_\alpha\le C(y^N)^{3/4}(y^N+d^N)^{3/4}\le\eta\nu d^N+C_\eta\nu^{-3}(y^N)^3+Cy^N\cdot(y^N)^{1/2}\);
\item transport: \(\int|\nabla U|^2|U||\nabla\phi_\alpha|\le C\|U\|_{L^6}\|\nabla U\|_{L^{12/5}}^2\le C(e^N+y^N)^{1/2}(y^N)^{3/4}(y^N+d^N)^{1/4}\le\eta\nu d^N+C_\eta\nu^{-1/3}(e^N+y^N)^{2/3}y^N+C(e^N+y^N)^{1/2}y^N\);
\item local pressure: \(\|\nabla P^{\rm loc}_\alpha\|_{L^2(\mathbb T_L)}\le C\|\nabla(\chi_\alpha U\otimes U)\|_2\le C\|U\|_{L^6(4Q_\alpha)}\|\nabla U\|_{L^3(4Q_\alpha)}+C\|U\|_{L^4(4Q_\alpha)}^2\)
      by the Riesz bound, so the two pressure terms with \(P^{\rm loc}\)
      are at most \(C\|\nabla P^{\rm loc}\|_2(y^N)^{1/2}\le\eta\nu d^N+C_\eta\nu^{-1/3}(e^N+y^N)^{2/3}y^N+C(e^N+y^N)(y^N)^{1/2}\)
      (Young on the factor \((y^N+d^N)^{1/4}\) of \(\|\nabla U\|_{L^3}\));
\item far pressure: by Lemma~\ref{lem:c2v57-far}, the two pressure
      terms with \(P^{\rm far}-c_\alpha\) are at most
      \(C\|\nabla P^{\rm far}_\alpha\|_{L^\infty(2Q_\alpha)}(y^N)^{1/2}\le C\varepsilon(y^N)^{1/2}\le C(\varepsilon^{3/2}+\varepsilon^{1/2}y^N)\).
\end{itemize}
Since \(y^N\le125\Lambda_0\) and \(e^N\le125\varepsilon\le125\), every factor
\((e^N+y^N)^{q}\), \((y^N)^{q}\) with \(q>0\) is bounded by \(C(\Lambda_0)\) (and
\((y^N)^3\le(125\Lambda_0)^2y^N\)), so the positive terms are at most
\(4\eta\nu d^N+b_1(\nu,\eta,\Lambda_0)(y^N+\varepsilon^{3/2})\).  Hence
\begin{equation}
 y'_\alpha\ \le\ -\nu d_\alpha+4\eta\nu\sum_{\beta\in N_\alpha}d_\beta+b_1\Bigl(\sum_{\beta\in N_\alpha}y_\beta+\varepsilon^{3/2}\Bigr).
 \label{eq:c2v57-cell}
\end{equation}

\emph{Step 2: the weighted sum.}  Let \(\alpha^*\) attain \(W\) at \(s\).
Multiply \eqref{eq:c2v57-cell} by \(\kappa^{|\alpha-\alpha^*|}\) and sum over
\(\alpha\).  The coefficient of \(d_\gamma\) is
\(-\nu\kappa^{|\gamma-\alpha^*|}+4\eta\nu\sum_{\alpha:\gamma\in N_\alpha}\kappa^{|\alpha-\alpha^*|}\le\nu\kappa^{|\gamma-\alpha^*|}(-1+500\eta\kappa^{-2})\),
since \(\gamma\in N_\alpha\) means \(|\alpha-\gamma|\le2\), there are \(125\) such \(\alpha\),
and \(|\alpha-\alpha^*|\ge|\gamma-\alpha^*|-2\).  With \(\eta:=\kappa^2/1000\) the
coefficient is at most \(-\frac\nu2\kappa^{|\gamma-\alpha^*|}\).  The coefficient of
\(y_\gamma\) in the \(b_1\)-sum is at most \(125b_1\kappa^{-2}\kappa^{|\gamma-\alpha^*|}\), and
\(\sum_\alpha\kappa^{|\alpha-\alpha^*|}=S_\kappa\).  Hence, with
\(D:=\sum_\gamma\kappa^{|\gamma-\alpha^*|}d_\gamma\),
\[
 W'=w'_{\alpha^*}\ \le\ -\frac\nu2D+125b_1\kappa^{-2}W+b_1S_\kappa\varepsilon^{3/2}.
\]

\emph{Step 3: the lower bound on \(D\).}  By
Lemma~\ref{lem:c2v57-interp}, with \(e^{(1)}_\gamma:=\sum_{|\beta-\gamma|\le1}e_\beta\le27\max e_\beta\le27\varepsilon\),
\[
 \sum_\gamma\kappa^{|\gamma-\alpha^*|}y_\gamma^2\ \le\ 54\varepsilon\sum_\gamma\kappa^{|\gamma-\alpha^*|}\Bigl(\sum_{|\beta-\gamma|\le1}d_\beta+C_\theta^2\sum_{|\beta-\gamma|\le1}y_\beta\Bigr)
 \ \le\ 54\varepsilon\cdot27\kappa^{-1}\bigl(D+C_\theta^2W\bigr),
\]
while by Cauchy--Schwarz
\(W^2=(\sum_\gamma\kappa^{|\gamma-\alpha^*|}y_\gamma)^2\le S_\kappa\sum_\gamma\kappa^{|\gamma-\alpha^*|}y_\gamma^2\).
Hence \(D\ge\frac{\kappa}{1458S_\kappa}\frac{W^2}{\varepsilon}-C_\theta^2W\), and
\[
 W'\ \le\ -\frac{\nu\kappa}{2916S_\kappa}\frac{W^2}{\varepsilon}+\Bigl(\frac\nu2C_\theta^2+125b_1\kappa^{-2}\Bigr)W+b_1S_\kappa\varepsilon^{3/2},
\]
which is \eqref{eq:c2v57-W} with \(a:=\nu\kappa/(2916S_\kappa)\) and
\(b:=\frac\nu2C_\theta^2+125b_1\kappa^{-2}+b_1S_\kappa\) (using \(\varepsilon^{3/2}\le\varepsilon\)).
\end{proof}

\begin{proposition}[Weighted local energy inequality]
\label{prop:c2v57-energy}
Under the same hypotheses (\(W\le\Lambda_0\), \(\varepsilon\le1\)), at almost every \(s\),
\begin{equation}
 \varepsilon'\ \le\ b'\,(\varepsilon+W),\qquad b'=b'(\nu,\kappa,\Lambda_0).
 \label{eq:c2v57-eps}
\end{equation}
\end{proposition}

\begin{proof}
The local energy identity,
\(e'_\alpha=-2\nu y_\alpha+\nu\int|U|^2\Delta\phi_\alpha+\int|U|^2U\cdot\nabla\phi_\alpha+2\int(P-c_\alpha)U\cdot\nabla\phi_\alpha\)
(the constant is allowed since \(\int U\cdot\nabla\phi_\alpha=0\)), gives, with the
bounds of Step 1 of the previous proof,
\(e'_\alpha\le C\nu e^N+C(e^N)^{3/4}(e^N+y^N)^{3/4}+C\|P^{\rm loc}_\alpha\|_{L^{3/2}}\|U\|_{L^3(2Q_\alpha)}+C\varepsilon(e^N)^{1/2}\),
with \(\|P^{\rm loc}_\alpha\|_{L^{3/2}}\le C\|U\|_{L^3(4Q_\alpha)}^2\) by the
Calder\'on--Zygmund bound; every term is at most
\(C(\Lambda_0)(e^N+y^N)+C\varepsilon\) for \(\varepsilon\le1\).  Take the weighted sum at
a maximizer of \(\varepsilon\) as in Step 2.
\end{proof}

% =============================================================================
\section{Capture and trapping}
\label{sec:c2v57-trap}
% =============================================================================

\begin{theorem}[Capture and trapping]
\label{thm:c2v57-trap}
Let \(W\), \(\varepsilon\) be nonnegative Lipschitz functions on \([s_0,s_1]\)
satisfying \eqref{eq:c2v57-W} and \eqref{eq:c2v57-eps} almost
everywhere as long as \(W\le\Lambda_0\) and \(\varepsilon\le1\), with
\(W(s_0)\le\Lambda_0\) and \(\varepsilon(s_0)=:\varepsilon_0\le1\).  Put
\[
 K:=\frac{4b}{a},\qquad p:=\frac{4b'}{a},\qquad \tau:=\frac2{aK},\qquad
 \Lambda:=b'(1+K),\qquad
 \varepsilon_1:=\Bigl(\varepsilon_0\Lambda_0^{\,p}K^{-p}\Bigr)^{1/(1+p)} .
\]
Then, provided \(K\varepsilon_1e^{\Lambda(s_1-s_0)}\le1\),
\begin{equation}
 W(s)\le\Lambda_0\ \text{ on }[s_0,s_1],\qquad
 W(s)\le K\varepsilon_1e^{\Lambda(s-s_0)},\quad \varepsilon(s)\le\varepsilon_1e^{\Lambda(s-s_0)}\ \text{ on }[s_0+\tau,s_1]\cap[s_0,s_1] .
 \label{eq:c2v57-trap}
\end{equation}
\end{theorem}

\begin{proof}
\emph{The trap region.}  Whenever \(W\ge K\varepsilon\) (and \(W\le\Lambda_0\),
\(\varepsilon\le1\)): \(bW\le\frac bK\frac{W^2}{\varepsilon}=\frac a4\frac{W^2}{\varepsilon}\) and
\(b\varepsilon\le\frac bK W\le\frac a4\frac{W^2}{\varepsilon}\), so \eqref{eq:c2v57-W} gives
\(W'\le-\frac a2\frac{W^2}{\varepsilon}<0\), and \eqref{eq:c2v57-eps} gives
\(\varepsilon'\le b'(1+\frac1K)W\le2b'W\).

\emph{Capture.}  Let \(Z:=\varepsilon W^p\).  In the trap region,
\(Z'=\varepsilon'W^p+p\varepsilon W^{p-1}W'\le2b'W^{p+1}-\frac{pa}2W^{p+1}=0\).  Let
\(s_c:=\inf\{s\ge s_0:W(s)\le K\varepsilon(s)\}\) (\(=s_1\) if there is no such
\(s\)).  On \([s_0,s_c)\), \(W\) decreases (so \(W\le\Lambda_0\)) and
\(Z\) is nonincreasing, so \(\varepsilon\le\varepsilon_0\Lambda_0^{\,p}W^{-p}\) and, with
\(W>K\varepsilon\), \(\varepsilon^{1+p}\le\varepsilon_0\Lambda_0^{\,p}K^{-p}\): \(\varepsilon\le\varepsilon_1\) on
\([s_0,s_c)\) (in particular \(\varepsilon\le1\) there).  Moreover
\(W'\le-\frac a2W^2/\varepsilon_1\) gives \(W(s)\le2\varepsilon_1/(a(s-s_0))\), which is
at most \(K\varepsilon_1\) for \(s\ge s_0+\tau\).

\emph{Trapping.}  Let \(m:=\max(\varepsilon,W/K)\), Lipschitz.  Almost
everywhere: where \(W<K\varepsilon\), \(m=\varepsilon\) and \(m'=\varepsilon'\le b'(\varepsilon+W)\le b'(1+K)\varepsilon=\Lambda m\);
where \(W>K\varepsilon\), \(m=W/K\) and \(m'=W'/K<0\le\Lambda m\); the set
\(\{W=K\varepsilon\}\) contributes nothing to a Lipschitz function's derivative
beyond one of the two one-sided values.  Hence \(m(s)\le m(s')e^{\Lambda(s-s')}\)
for \(s\ge s'\) as long as the hypotheses \(W\le\Lambda_0\), \(\varepsilon\le1\) hold.
Let \(s_2:=\min(s_c,s_0+\tau)\).  At \(s_2\), either \(W(s_2)\le K\varepsilon(s_2)\)
with \(\varepsilon(s_2)\le\varepsilon_1\), or \(W(s_2)\le K\varepsilon_1\) and
\(\varepsilon(s_2)\le\varepsilon_1\); in both cases \(m(s_2)\le\varepsilon_1\).  Therefore
\(m(s)\le\varepsilon_1e^{\Lambda(s-s_2)}\le\varepsilon_1e^{\Lambda(s-s_0)}\) on \([s_2,s_1]\), and
the proviso \(K\varepsilon_1e^{\Lambda(s_1-s_0)}\le1\) keeps \(W\le Km\le1\le\Lambda_0\)
and \(\varepsilon\le m\le1\) throughout, so the hypotheses never fail (a
continuity argument on the first time they would).  On \([s_0,s_2]\),
\(W\le\Lambda_0\) by monotonicity and \(\varepsilon\le\varepsilon_1\).  This is
\eqref{eq:c2v57-trap}, since \([s_0+\tau,s_1]\subset[s_2,s_1]\).
\end{proof}

% =============================================================================
\section{No spread windows}
\label{sec:c2v57-main}
% =============================================================================

\begin{theorem}[The intermediate class has no spread windows]
\label{thm:c2v57-no-spread}
Let \(u\) be in the intermediate class with constant \(M\)
(Definition~\ref{def:c2v42-intermediate}), with window times \(t_j\),
scales \(\lambda_j\), rescalings \(U_j\) on \([-1,0)\).  There is
\(\varepsilon_*=\varepsilon_*(\nu,M)>0\) such that for every \(j\) with \(2\pi/\lambda_j\ge8\),
\begin{equation}
 \max_\alpha\int|U_j(-1)|^2\phi_\alpha\ \ge\ \varepsilon_* ,
 \label{eq:c2v57-no-spread}
\end{equation}
that is, at every window time some parabolic cell carries energy at
least \(\varepsilon_*\).  In particular no subsequence of the windows is
spread in the sense of Corollary~\ref{cor:c2v48-collapse}: every
window is concentrated, and the two-case description of the
intermediate class at window scales reduces to the concentrated case.
\end{theorem}

\begin{proof}
Fix \(\kappa=\frac12\), \(\Lambda_0:=M\), and let \(a,b,b'\) be the constants of
Propositions~\ref{prop:c2v57-enstrophy}--\ref{prop:c2v57-energy} for
this \(\Lambda_0\), and \(K,p,\tau,\Lambda\) as in Theorem~\ref{thm:c2v57-trap}.
Define \(\varepsilon_*\) by \(K\varepsilon_1e^{\Lambda}=1\) with \(\varepsilon_0=S_\kappa\varepsilon_*\), that is
\(\varepsilon_*:=S_\kappa^{-1}K^{p}M^{-p}(Ke^{\Lambda})^{-(1+p)}\).  Suppose
\eqref{eq:c2v57-no-spread} fails for some \(j\): then \(\varepsilon(-1)\le S_\kappa\max_\alpha e_\alpha(-1)<S_\kappa\varepsilon_*=\varepsilon_0\),
while \(W(-1)\le Y_j(-1)\le M=\Lambda_0\) by \eqref{eq:c2v57-comparison} and
the definition of the class.  Theorem~\ref{thm:c2v57-trap} on
\([s_0,s_1]=[-1,-\sigma]\) for every \(\sigma\in(0,1)\), whose proviso holds
since \(s_1-s_0\le1\), gives \(W\le\max(\Lambda_0,1)=M\) on \([-1,-1+\tau]\) and
\(W\le K\varepsilon_1e^{\Lambda}\le1\) on \([-1+\tau,0)\) (if \(\tau\ge1\) the first
bound alone applies).  Hence the maximal local enstrophy satisfies
\(\delta\le W\le M\) on \([-1,0)\) and \(\delta\le1\) on \([-1+\tau,0)\).  By the
global inequality \eqref{eq:c2v56-global}, with \(\delta\le M\),
\[
 Y_j(s)\ \le\ Y_j(-1)\exp\Bigl(C_3\bigl(\nu^{-3}M^2+M^{1/2}\bigr)(s+1)\Bigr)\ \le\ M\,e^{C_3(\nu^{-3}M^2+M^{1/2})}
 \qquad(-1\le s<0):
\]
the enstrophy of the rescaled solution is bounded on \([-1,0)\).  In
the original variables, \(\|\Lambda u(t)\|_2^2=\lambda_j^{-1}Y_j(s)\) is bounded on
\([t_j,T_*)\), so \(u\in L^\infty(t_j,T_*;H^1)\), and by the classical
criterion (bounded enstrophy up to \(T_*\) gives \(u\in L^2(t_j,T_*;H^2)\)
and a strong continuation past \(T_*\)) \(T_*\) is not a singular time,
contradicting \(T_*<\infty\) being the maximal existence time of the
smooth solution.  The identification of spread windows uses
Corollary~\ref{cor:c2v48-collapse}: a spread subsequence has
\(\sup_xE^x_j\to0\), hence \(\max_\alpha e_\alpha(-1)\le\sup_xE^x_j\cdot8\to0\)
(a cell is covered by eight unit balls), which violates
\eqref{eq:c2v57-no-spread} for large \(j\).
\end{proof}

\begin{corollary}[Consequences]
\label{cor:c2v57-consequences}
In the intermediate class with constant \(M\):
\begin{enumerate}[label=\textup{(\roman*)},leftmargin=3.2em]
\item Every window is concentrated: there is a centre \(x_j\) with
      \(E^{x_j}_j\ge\varepsilon_*/8\), hence by
      Theorem~\ref{thm:c2v48-energy-active} window dissipation at least
      \(\varepsilon'(\varepsilon_*/8,M,\nu)\) in the ball of radius
      \(R'(\varepsilon_*/8,M,\nu)\) around \(x_j\), a nonzero strong window limit
      \(U'\) with \(\int_{B_1}|U'(-1)|^2\ge\varepsilon_*/8\) along every subsequence,
      and a coherent active centre map.
\item The both-spread case (Proposition~\ref{prop:c2v48-both}), the
      uniformly spread and travelling windows of V53--V54, and the
      spread-window delay of V56 are vacuous in the intermediate
      class; Theorem~\ref{thm:c2v55-consolidated}(4) is superseded by
      the exclusion.
\item Quantitatively, for every window time, the energy of the most
      energetic parabolic cell is at least \(\varepsilon_*(\nu,M)\), an
      explicit (in principle) function of the viscosity and the
      sequence constant through the absolute constants of the local
      Sobolev, Gagliardo--Nirenberg, Calder\'on--Zygmund and
      kernel bounds; the corresponding scale-invariant statement in
      the original variables is
      \(\sup_x\lambda_j^{-1}\int_{B_{\lambda_j}(x)}|u(t_j)|^2\ge c\,\varepsilon_*\): the
      critical Morrey quantity of the solution at the window time is
      bounded below.
\end{enumerate}
\end{corollary}

\begin{proof}
(i) and (iii) are restatements with the covering constants; (ii) lists
the statements whose hypothesis is now known to be empty.
\end{proof}

\begin{firewall}[What the exclusion does and does not say]
\label{fw:c2v57-scope}
The theorem is an a priori statement inside the intermediate class:
it uses the type-I bound at the window time (\(W(-1)\le M\)) to run the
capture phase with constants depending on \(M\), and the small cell
energy to run the trap; a small critical Morrey norm alone, without
the enstrophy bound at the window time, is not shown to suffice (the
capture phase would have no a priori bound \(\Lambda_0\)), which is
consistent with the endpoint Morrey space being outside the known
well-posedness theory.  The exclusion does not bound the number of
energetic cells, does not say that the concentrated structure persists
after the window, and does not exclude the intermediate class: it
identifies its windows as concentrated.  The trap mechanism is
local, and its localized form (quiet cells stay quiet unless energy
flows in from energetic neighbours) is the next acquisition.  No
global regularity claim is made.
\end{firewall}

% =============================================================================
\section{Integrated V57 conclusion and next acquisition order}
\label{sec:c2v57-conclusion}
% =============================================================================

\begin{theorem}[What V57 proves]
\label{thm:c2v57-conclusion}
Relative to the first cycle and V1--V56: the far-pressure bound
(Lemma~\ref{lem:c2v57-far}), the interpolation
\eqref{eq:c2v57-interp}, the weighted local inequalities
\eqref{eq:c2v57-W}--\eqref{eq:c2v57-eps}, the capture-and-trapping
theorem \eqref{eq:c2v57-trap}, and the exclusion of spread windows
\eqref{eq:c2v57-no-spread} with its consequences.  The full
Navier--Stokes stability/global-regularity theorem remains unproved.
\end{theorem}

\begin{proof}
The cited results.
\end{proof}

\begin{programme}[V58 acquisition order]
\label{prog:c2v58}
\begin{enumerate}[label=\textup{(AC\arabic*)},leftmargin=4.8em]
\item Localize the trap.  For a set \(\mathcal Q\) of cells (the quiet
      set) define the weighted maxima over \(\mathcal Q\) only, with the
      weights cut off at the boundary of \(\mathcal Q\); determine the
      inequalities satisfied by them when the cells outside \(\mathcal Q\)
      carry arbitrary (window-bounded) energy and enstrophy, in
      particular the far-pressure gradient on a quiet cell at distance
      \(r\) from the energetic set (a sum of \(r^{-4}\) times the energies
      of the energetic cells), and derive the rate at which the quiet
      region can be invaded.
\item Count.  At a window the total energy of the rescaled solution
      is \(\lambda_j^{-1}\|u(t_j)\|_2^2\), unbounded; but the cells with energy
      at least \(\varepsilon_*\) and the cells with dissipation at least a
      threshold are counted by the window dissipation budget and by
      the \(L^6\) bound; determine whether the energetic set at a window
      is finite uniformly in \(j\), or whether the energy \(\varepsilon_*\) can be
      carried by a diverging number of cells, each with its own
      dissipation, and what the localized trap says about the
      interaction of far-apart energetic cells.
\item The next compulsory companion audit is V60.
\end{enumerate}
\end{programme}

\begin{assessment}[Position after V57]
\label{ass:c2v57-position}
The intermediate class has no spread windows: every window has a
parabolic cell with energy at least \(\varepsilon_*(\nu,M)\), and every window
is concentrated with a nonzero strong limit.  No full stability or
global-regularity claim is made.
\end{assessment}

% =============================================================================
\section*{Version V57 (second cycle) append-only record}
\addcontentsline{toc}{section}{Version V57 (second cycle) append-only record}
% =============================================================================

The complete V56 source of the second cycle was retained before this
continuation.  Its SHA--256 digest was
\begin{center}\scriptsize\ttfamily
56b16436e25ecf817ea38c9ddecc943d1f1b81c33ecc9ab34c425aef8792fb08.
\end{center}
V57 proved the cell trap and excluded spread windows in the
intermediate class.  No full regularity claim is made.

% =============================================================================
\section{V58: the energetic set is finite, and the localized trap on the window}
\label{sec:c2v58-entry}
% =============================================================================

Items (AC1) and (AC2) of Programme~\ref{prog:c2v58} are carried out.
Item (AC2) closes cleanly: at a window the number of parabolic cells
with energy above a threshold is bounded by an explicit multiple of the
cube of the sequence constant over the cube of the threshold, by the
\(L^6\) bound alone, so the energetic set of a window is finite
uniformly in the window index; with V57, a window has between one and
\(N_M\) energetic cells.  Item (AC1) is carried out on the window: the
far-pressure gradient on a cell is bounded by a polynomially weighted
sum of all cell energies, hence on a quiet cell at distance \(r\) from
the energetic set by the quiet energy plus an explicit multiple of
\(N_M M/r\); the weighted local inequalities of V57 hold for the quiet
region with two corrections, a polynomial halo term from the energetic
set and an exponentially small boundary term integrable over the
window.  The trap for the quiet region is conditional: its effective
energy at distance \(r\) involves the quiet energies at distance
\(r/4\), a hierarchy that the window length does not close in general,
and beyond the window the boundary term is uncontrolled.  This is
recorded as the firewall, with the precise missing estimate.  No
global regularity claim is made.  The next compulsory companion audit
is V60.

% =============================================================================
\section{The energetic set is finite}
\label{sec:c2v58-count}
% =============================================================================

Notation as in V57 (cells, partition of unity, \(e_\alpha\), \(y_\alpha\),
\(d_\alpha\)).  Let \(u\) be in the intermediate class with constant \(M\ge1\),
\(U=U_j\) a window rescaling, \(J_M\) the window, on which \(Y\le2M\).

\begin{proposition}[Counting energetic cells]
\label{prop:c2v58-count}
On the window, every cell has
\(e_\alpha\le4\|U\|_{L^6(2Q_\alpha)}^2\le a_M:=8C_S^2M\) and \(y_\alpha\le Y\le2M\), and for
every threshold \(\varepsilon_{\rm th}>0\)
\begin{equation}
 \#\{\alpha:\ e_\alpha(s)\ge\varepsilon_{\rm th}\}\ \le\ N_M(\varepsilon_{\rm th}):=\frac{1728\,C_S^6\,(2M)^3}{\varepsilon_{\rm th}^3}
 \qquad(s\in J_M).
 \label{eq:c2v58-count}
\end{equation}
In particular, with \(\varepsilon_{\rm th}:=\varepsilon_*/2\) (\(\varepsilon_*\) of
Theorem~\ref{thm:c2v57-no-spread}), the energetic set
\(\mathcal A(s):=\{\alpha:e_\alpha(s)\ge\varepsilon_*/2\}\) at the window time \(s=-1\) is
nonempty and has at most \(N_M:=N_M(\varepsilon_*/2)\) elements, uniformly in
\(j\).
\end{proposition}

\begin{proof}
H\"older on \(2Q_\alpha\) (volume \(8\)) gives
\(e_\alpha\le\|U\|_{L^2(2Q_\alpha)}^2\le8^{2/3}\|U\|_{L^6(2Q_\alpha)}^2\), and Sobolev
\(\|U\|_6^2\le C_S^2Y\le2C_S^2M\).  Cubing and summing over the cells,
\(\sum_\alpha e_\alpha^3\le64\sum_\alpha\|U\|_{L^6(2Q_\alpha)}^6\le64\cdot27\|U\|_6^6\le1728C_S^6(2M)^3\),
since each point lies in at most \(27\) cubes \(2Q_\alpha\).  Chebyshev
gives \eqref{eq:c2v58-count}.  Nonemptiness is
\eqref{eq:c2v57-no-spread}.
\end{proof}

\begin{corollary}[Windows are finite configurations]
\label{cor:c2v58-configuration}
At every window time of the intermediate class, the rescaled solution
has between \(1\) and \(N_M\) parabolic cells of energy at least
\(\varepsilon_*/2\), all other cells having energy below \(\varepsilon_*/2\); by
Theorem~\ref{thm:c2v48-energy-active}, each energetic cell is
dissipation-active at radius \(R'(\varepsilon_*/16,M,\nu)\) with level
\(\varepsilon'(\varepsilon_*/16,M,\nu)\) on the window, and along every
subsequence the rescalings centred at an energetic cell converge to a
nonzero strong solution on the window.
\end{corollary}

\begin{proof}
Proposition~\ref{prop:c2v58-count} and Theorem~\ref{thm:c2v48-energy-active}
(a cell of energy \(\varepsilon_*/2\) contains a unit ball of energy at least
\(\varepsilon_*/16\), by the covering of the cell by eight unit balls).
\end{proof}

% =============================================================================
\section{The far pressure with polynomial weights}
\label{sec:c2v58-far}
% =============================================================================

\begin{lemma}[Far pressure gradient as a weighted sum of cell energies]
\label{lem:c2v58-far}
With the notation of Lemma~\ref{lem:c2v57-far}, for every cell \(\beta\),
\begin{equation}
 \|\nabla P^{\rm far}_\beta\|_{L^\infty(2Q_\beta)}\ \le\ C_{\rm far}''\sum_\gamma(1+|\gamma-\beta|)^{-4}e_\gamma ,
 \label{eq:c2v58-far}
\end{equation}
with \(C_{\rm far}''\) absolute.  Consequently, if \(\mathcal A\) is any set
of at most \(N\) cells with \(e_\gamma\le a\) for all \(\gamma\), and
\(\beta\) is at lattice distance \(\rho\ge2\) from \(\mathcal A\), then
\begin{equation}
 \|\nabla P^{\rm far}_\beta\|_{L^\infty(2Q_\beta)}\ \le\ C_{\rm far}''\Bigl(S_4\max_{\gamma:\,{\rm dist}(\gamma,\mathcal A)\ge\rho/2}e_\gamma+\frac{1024\,Na}{\rho}\Bigr),
 \qquad S_4:=\sum_{\gamma}(1+|\gamma|)^{-4}<\infty .
 \label{eq:c2v58-halo}
\end{equation}
\end{lemma}

\begin{proof}
In the proof of Lemma~\ref{lem:c2v57-far}, the \(k\)-th shell
contribution is \(C2^{-4k}\int_{S_k}|U|^2\le C2^{-4k}\sum_{\gamma:\,Q_\gamma\cap S_k\ne\emptyset}e_\gamma\),
and for such \(\gamma\), \(1+|\gamma-\beta|\le2^{k+2}\), so \(2^{-4k}\le256(1+|\gamma-\beta|)^{-4}\);
each cell meets at most two shells.  The \(L^{-4}\) part is bounded
by \(CL^{-4}\sum_\gamma e_\gamma\le C\sum_\gamma(1+|\gamma-\beta|)^{-4}e_\gamma\) since
\(1+|\gamma-\beta|\le L\) on the torus.  For \eqref{eq:c2v58-halo}: the cells
\(\gamma\) with \({\rm dist}(\gamma,\mathcal A)\ge\rho/2\) contribute at most
\(S_4\) times their maximal energy; the remaining cells lie within
distance \(\rho/2\) of one of the \(N\) cells of \(\mathcal A\), so they
number at most \(N(\rho+1)^3\le8N\rho^3\), and each is at distance more
than \(\rho/2\) from \(\beta\), contributing at most \((\rho/2)^{-4}a\).
\end{proof}

% =============================================================================
\section{The localized weighted inequalities on the window}
\label{sec:c2v58-localized}
% =============================================================================

Fix the window rescaling \(U\) on \(J_M\), \(\mathcal A:=\mathcal A(-1)\) the
energetic set at the window time, \(N:=|\mathcal A|\le N_M\), and for
\(r\ge3\) the quiet regions \(\mathcal Q_r:=\{\alpha:{\rm dist}(\alpha,\mathcal A)\ge r\}\).
Define, with \(\kappa=\frac12\),
\begin{equation}
 w^{\mathcal Q}_\alpha:=\sum_{\beta\in\mathcal Q_3}\kappa^{|\beta-\alpha|}y_\beta,\quad
 \varepsilon^{\mathcal Q}_\alpha:=\sum_{\beta\in\mathcal Q_3}\kappa^{|\beta-\alpha|}e_\beta,\quad
 W_r:=\max_{\alpha\in\mathcal Q_r}w^{\mathcal Q}_\alpha,\quad
 \varepsilon_r:=\max_{\alpha\in\mathcal Q_r}\varepsilon^{\mathcal Q}_\alpha,\quad
 \bar e_r:=\max_{\alpha\in\mathcal Q_r}e_\alpha .
 \label{eq:c2v58-localized}
\end{equation}

\begin{proposition}[Localized weighted inequalities on the window]
\label{prop:c2v58-localized}
There are constants \(a=a(\nu)>0\), \(b,b',C_M\) depending on
\((\nu,M)\) only, such that for every \(r\ge8\) and almost every
\(s\in J_M\) with \(\varepsilon_r(s)\le1\),
\begin{align}
 W_r'\ &\le\ -\,a\,\frac{W_r^2}{\hat\varepsilon_r}+b\,(W_r+\hat\varepsilon_r)+\pi_r(s),
 \label{eq:c2v58-W}\\
 \varepsilon_r'\ &\le\ b'\,(\varepsilon_r+W_r+\hat\varepsilon_r)+\pi_r(s),
 \label{eq:c2v58-eps}
\end{align}
where
\begin{equation}
 \hat\varepsilon_r:=\varepsilon_r+S_4\,\bar e_{r/4}+\frac{C_M N}{r}+\kappa^{r/2}C_M,\qquad
 \pi_r(s):=\kappa^{r/2}C_M\bigl(1+\nu\|\Delta U(s)\|_2^2\bigr),\qquad
 \int_{J_M}\pi_r\le\kappa^{r/2}C_M(\ell_M+4M).
 \label{eq:c2v58-eff}
\end{equation}
\end{proposition}

\begin{proof}
Let \(\alpha^*\in\mathcal Q_r\) attain \(W_r\) at \(s\).  Then
\(W_r'=\sum_{\beta\in\mathcal Q_3}\kappa^{|\beta-\alpha^*|}y'_\beta\).  For each
\(\beta\in\mathcal Q_3\), the neighbourhood \(N_\beta\) lies in \(\mathcal Q_1\), so
the cell inequality \eqref{eq:c2v57-cell} holds with \(\Lambda_0=2M\)
(all cell enstrophies are at most \(2M\) and all cell energies at most
\(a_M\) on the window) and with the far-pressure term bounded by
Lemma~\ref{lem:c2v58-far}: for \(\beta\) at distance \(\rho\) from
\(\mathcal A\), the far-pressure contribution is at most
\(C(S_4\bar e_{\rho/2}+1024Na_M/\rho)(y^N_\beta)^{1/2}\).  Split the sum over
\(\beta\) according to \(\rho\ge r/2\) or \(\rho<r/2\).

For \(\rho\ge r/2\): \(\bar e_{\rho/2}\le\bar e_{r/4}\) and \(Na_M/\rho\le2Na_M/r\),
so the far-pressure energy of \(\beta\) is at most
\(f_r:=C(S_4\bar e_{r/4}+2Na_M/r)\), and the cell inequality reads
\(y'_\beta\le-\nu d_\beta+4\eta\nu\sum_{N_\beta}d_\gamma+b_1(\sum_{N_\beta}y_\gamma+f_r^{3/2}+f_r^{1/2}\sum_{N_\beta}y_\gamma)\)
(the far-pressure term \(Cf_r(y^N)^{1/2}\le C(f_r^{3/2}+f_r^{1/2}y^N)\)).
The neighbours \(\gamma\in N_\beta\) lie in \(\mathcal Q_{r/2-2}\subset\mathcal Q_3\), so
their \(y_\gamma\), \(d_\gamma\) are among the summed quantities.

For \(\rho<r/2\): \(|\beta-\alpha^*|\ge r/2\), so the weight is at most
\(\kappa^{r/2}\); the cell inequality with the crude far-pressure bound
\(Ca_MS_4\) (all energies at most \(a_M\)) gives
\(y'_\beta\le-\nu d_\beta+4\eta\nu\sum_{N_\beta}d_\gamma+C_M\), and the number of such
\(\beta\) is at most \(N r^3\), while their neighbours' dissipations
\(d_\gamma\) (\(\gamma\) possibly in \(\mathcal Q_1\setminus\mathcal Q_3\), not among the
summed quantities) total at most \(\|\Delta U\|_2^2\).

Weighted sum: the dissipation coefficients for \(\gamma\in\mathcal Q_3\) are
handled as in Step 2 of Proposition~\ref{prop:c2v57-enstrophy}, giving
\(-\frac\nu2D\) with \(D:=\sum_{\gamma\in\mathcal Q_3}\kappa^{|\gamma-\alpha^*|}d_\gamma\); the
dissipations of cells outside \(\mathcal Q_3\) appear only through
\(\beta\) with \(\rho<5\), hence with weight at most \(\kappa^{r-5}\), and
contribute at most \(4\eta\nu\kappa^{r-5}\cdot125\|\Delta U\|_2^2\).  The
\(b_1\)-terms with \(\rho\ge r/2\) give
\(b_1(125\kappa^{-2}(1+f_r^{1/2})W_r+S_\kappa f_r^{3/2})\).  The \(\rho<r/2\) terms
give at most \(\kappa^{r/2}Nr^3C_M\le\kappa^{r/4}C_M'\) (absorbing
\(Nr^3\kappa^{r/4}\le C_M\)).  The lower bound on \(D\) is Step 3 of
Proposition~\ref{prop:c2v57-enstrophy} restricted to \(\gamma\in\mathcal Q_3\):
the interpolation \eqref{eq:c2v57-interp} for \(\gamma\in\mathcal Q_3\) involves
\(\sum_{|\beta-\gamma|\le1}e_\beta\), which is at most \(27\bar e_{r/4}\) for
\(\gamma\in\mathcal Q_{r/4+1}\) and at most \(27a_M\) otherwise, the latter
\(\gamma\) having weight at most \(\kappa^{3r/4-1}\); and the terms
\(\sum_{|\beta-\gamma|\le1}d_\beta\) with \(\beta\notin\mathcal Q_3\) occur only for
\(\gamma\) at distance at most \(4\) from \(\mathcal A\), weight at most
\(\kappa^{r-4}\), contributing \(\kappa^{r-4}C\|\Delta U\|_2^2\) to the
comparison.  Altogether
\(W_r^2\le S_\kappa\sum\kappa^{|\gamma-\alpha^*|}y_\gamma^2\le C(\bar e_{r/4}+\kappa^{3r/4-1}a_M)(D+C_\theta^2W_r)+\kappa^{r-4}C(D+\|\Delta U\|_2^2)\cdot a_M\),
which gives \(D\ge c\,W_r^2/(\bar e_{r/4}+\kappa^{r/2}C_M)-C_\theta^2W_r-\kappa^{r/2}C_M\|\Delta U\|_2^2\)
with \(c\) absolute (the denominator is at most \(C_M\), since
\(\bar e_{r/4}\le a_M\)).  Since \(\varepsilon_r\ge\bar e_r\) and \(\hat\varepsilon_r\ge\bar e_{r/4}S_4+\kappa^{r/2}C_M\),
collecting the terms yields \eqref{eq:c2v58-W} with
\(a=\nu c/2\), the \(b\)-terms absorbing \(f_r\le C\hat\varepsilon_r\), and
\(\pi_r\) collecting the \(\kappa^{r/2}\)-terms.  The energy inequality
\eqref{eq:c2v58-eps} is obtained in the same way from the local energy
identity, the far-pressure work term being \(C f_r(e^N_\beta)^{1/2}\le C(f_r+e^N_\beta)\).
The integral of \(\pi_r\) uses \(\int_{J_M}\|\Delta U\|_2^2\le4M/\nu\)
(Proposition~\ref{prop:c2v51-palinstrophy}).
\end{proof}

\begin{theorem}[Conditional trapping of the quiet region on the window]
\label{thm:c2v58-conditional}
Let \(r\ge8\), and suppose that on an interval \([-1,s_1]\subset J_M\) the
quiet energies at distance \(r/4\) satisfy \(\bar e_{r/4}\le\bar e\) and the
quiet enstrophy at the window time satisfies \(W_r(-1)\le\Lambda_0\le2M\).
Put \(g_r:=S_4\bar e+C_MN/r+\kappa^{r/2}C_M\) (the non-dynamical part of
\(\hat\varepsilon_r\)) and \(\Pi_r:=\int_{J_M}\pi_r\le\kappa^{r/2}C_M(\ell_M+4M)\).  Then,
with \(K=4b/a\), \(\Lambda=b'(1+K)+b'\), on \([-1,s_1]\):
\begin{equation}
 \max\bigl(\varepsilon_r(s)+g_r,\,W_r(s)/K\bigr)\ \le\ e^{\Lambda(s+1)}\Bigl(\max\bigl(\varepsilon_r(-1)+g_r,\,\Lambda_0/K\bigr)+2\Pi_r\Bigr).
 \label{eq:c2v58-conditional}
\end{equation}
In particular, if \(\varepsilon_r(-1)\), \(\bar e\), \(N/r\), \(\kappa^{r/2}\) and
\(\Lambda_0/K\) are all small, the quiet region at distance \(r\) stays
quiet (small energy and small enstrophy) through the window.
\end{theorem}

\begin{proof}
Let \(\hat\varepsilon:=\varepsilon_r+g_r\); on \([-1,s_1]\) one has \(\hat\varepsilon_r\le\hat\varepsilon\),
so \eqref{eq:c2v58-W}--\eqref{eq:c2v58-eps} hold with \(\hat\varepsilon_r\)
replaced by \(\hat\varepsilon\) (a larger denominator weakens the negative
term, and larger positive terms weaken the inequality).  Let
\(m:=\max(\hat\varepsilon,W_r/K)\).  Where \(W_r<K\hat\varepsilon\):
\(m'=\varepsilon_r'\le b'(\varepsilon_r+W_r+\hat\varepsilon)+\pi_r\le b'(2+K)m+\pi_r=\Lambda m+\pi_r\).
Where \(W_r>K\hat\varepsilon\): as in Theorem~\ref{thm:c2v57-trap},
\(W_r'\le-\frac a2W_r^2/\hat\varepsilon+\pi_r\le\pi_r\), so \(m'\le\pi_r/K\le\pi_r\).
Hence \(m'\le\Lambda m+\pi_r\) almost everywhere, and Gr\"onwall gives
\(m(s)\le e^{\Lambda(s+1)}(m(-1)+\int_{-1}^s\pi_r)\), which is
\eqref{eq:c2v58-conditional} (with a factor \(2\) to spare).  The
hypothesis \(\varepsilon_r\le1\) of the proposition is maintained by
continuity when the right side is at most \(1\).
\end{proof}

\begin{firewall}[Item (AC1): what the localized trap does not close]
\label{fw:c2v58-AC1}
Three gaps separate Theorem~\ref{thm:c2v58-conditional} from an
unconditional statement.
\begin{enumerate}[label=\textup{(\alph*)},leftmargin=3.2em]
\item \emph{The hierarchy.}  The effective energy at distance \(r\)
      contains the quiet energy at distance \(r/4\), through the far
      pressure of the cells between \(r/4\) and \(r/2\) from the energetic
      set; controlling \(\bar e_{r/4}\) requires the statement at
      distance \(r/4\), and so on down to distance \(8\), where the
      cells adjacent to the buffer are not controlled.  The regress
      terminates only if the total energy of the non-quiet region is
      bounded, which requires the energy identity on that region with a
      quiet boundary, itself the conclusion sought.
\item \emph{The capture time.}  Quiet cells at the window time may
      carry enstrophy up to \(2M\) with small energy; the capture
      phase of Theorem~\ref{thm:c2v57-trap} takes a time \(\tau=2/(aK)\)
      of order \(\nu/C(M)\), which need not be shorter than the window
      \(\ell_M\) of order \(\nu^3/M^2\).  Without capture within the window
      the bound \eqref{eq:c2v58-conditional} is the crude
      \(\Lambda_0/K\), of order \(M\nu/C(M)\), not small.
\item \emph{Beyond the window.}  The boundary term \(\pi_r\) and the
      constants \(\Lambda_0=2M\), \(a_M\) use the window bounds \(Y\le2M\),
      \(\int\|\Delta U\|^2\le4M/\nu\); after the window nothing bounds the
      buffer cells, and the localized inequalities have no content.
\end{enumerate}
The clean part of the position is Corollary~\ref{cor:c2v58-configuration}:
a window is a finite configuration of at most \(N_M\) energetic cells
with a quiet complement.  No global regularity claim is made.
\end{firewall}

% =============================================================================
\section{Integrated V58 conclusion and next acquisition order}
\label{sec:c2v58-conclusion}
% =============================================================================

\begin{theorem}[What V58 proves]
\label{thm:c2v58-conclusion}
Relative to the first cycle and V1--V57: the counting bound
\eqref{eq:c2v58-count} and the finite-configuration corollary, the
weighted far-pressure bound \eqref{eq:c2v58-far}--\eqref{eq:c2v58-halo},
the localized inequalities \eqref{eq:c2v58-W}--\eqref{eq:c2v58-eps}, and
the conditional trapping \eqref{eq:c2v58-conditional}.  The full
Navier--Stokes stability/global-regularity theorem remains unproved.
\end{theorem}

\begin{proof}
The cited results.
\end{proof}

\begin{programme}[V59 acquisition order]
\label{prog:c2v59}
\begin{enumerate}[label=\textup{(AC\arabic*)},leftmargin=4.8em]
\item Go upstream of the hierarchy: instead of localizing in space,
      use the finite configuration globally.  At a window the energetic
      set has at most \(N_M\) cells; apply the window compactness at each
      energetic cell and determine whether the \(N_M\) window limits
      (nonzero strong solutions on \(J_M\)) can be assembled into a
      single limit configuration on the big torus, with the quiet
      complement converging to zero in \(L^2_{\rm loc}\) uniformly, and
      whether the trap of V57 applied to the limit (which has zero
      energy far from the configuration) shows that the limit solution
      is regular past the window: the limit of the rescalings on
      \([-1,0)\) is a suitable weak solution whose cell energies at
      \(s=-1\) vanish outside a bounded set.
\item Determine whether the V57 trap can be run for the limit
      solution on \([-1,0)\) with the energetic set as a finite
      perturbation: the quiet complement of the limit has zero energy
      at \(s=-1\), so its effective energy is the halo \(C_MN/r\) alone,
      and the hierarchy of Firewall~\ref{fw:c2v58-AC1}(a) reduces to
      the halo's decay; record the resulting statement for the limit,
      and whether it transfers to the rescalings by the compactness of
      V42.
\item The next compulsory companion audit is V60.
\end{enumerate}
\end{programme}

\begin{assessment}[Position after V58]
\label{ass:c2v58-position}
A window is a finite configuration of at most \(N_M\) energetic cells
with a quiet complement; the localized trap on the window is
conditional on a hierarchy of quiet energies.  No full stability or
global-regularity claim is made.
\end{assessment}

% =============================================================================
\section*{Version V58 (second cycle) append-only record}
\addcontentsline{toc}{section}{Version V58 (second cycle) append-only record}
% =============================================================================

The complete V57 source of the second cycle was retained before this
continuation.  Its SHA--256 digest was
\begin{center}\scriptsize\ttfamily
8f2468f73d5d5759fba38976a755c14827bf8ac1f1ef198938ebe3af4ab8b9dd.
\end{center}
(The V57 file was corrected in place after its assembly: the pressure
parenthesis, the torus term of the far-pressure lemma, the smallness
hypothesis \(\varepsilon\le1\) in the weighted inequalities, and the exponent
of the local-pressure Young inequality; the digest above is that of
the corrected file.)  V58 proved the finiteness of the energetic set
and the conditional localized trap on the window.  No full regularity
claim is made.

% =============================================================================
\section{V59: the Morrey floor at every time before a singularity, and the
cluster limits of a window}
\label{sec:c2v59-entry}
% =============================================================================

Programme~\ref{prog:c2v59} is carried out with a change of order.  The
trap of V57 is first restated as a theorem about an arbitrary singular
solution at an arbitrary time: the critical Morrey energy at the
parabolic scale is bounded below by an explicit function of the scaled
enstrophy at that time, for every time before the singular time, with
no class hypothesis.  This is the upstream form of the exclusion of
spread windows, and it gives, for the intermediate class, the same
floor at window times, and for a type-I singularity a floor at all
times.  Item (AC1) is then carried out: the energetic cells of a window
fall into finitely many clusters of bounded diameter, each cluster
producing a nonzero strong window limit on the whole space whose far
cells have small but not zero energy (the programme's ``zero energy''
is corrected).  Item (AC2) is closed by a firewall: the limit
inherits the energetic cells, so the trap cannot be run on it for the
same reason as on the rescalings.  No global regularity claim is made.
The next compulsory companion audit is V60.

% =============================================================================
\section{The Morrey floor}
\label{sec:c2v59-floor}
% =============================================================================

Let \(u\) be a smooth solution on the torus of side \(2\pi\) with first
singular time \(T_*<\infty\).  For \(t<T_*\) with \(T_*-t\le\pi^2/16\) put
\begin{equation}
 \lambda(t):=(T_*-t)^{1/2},\qquad
 \mathfrak Y(t):=\lambda(t)\|\Lambda u(t)\|_2^2,\qquad
 \mathfrak M(t):=\sup_x\lambda(t)^{-1}\int_{B_{\lambda(t)}(x)}|u(t)|^2 ,
 \label{eq:c2v59-quantities}
\end{equation}
the scaled enstrophy and the critical Morrey energy at the parabolic
scale.  Both are scale-invariant.  H\"older and Sobolev give the upper
bound \(\mathfrak M(t)\le C\,C_S^2\,\mathfrak Y(t)\).

\begin{theorem}[Morrey floor at every time before the singularity]
\label{thm:c2v59-floor}
Let \(\varepsilon_*(\nu,\Lambda)\) be the function of
Theorem~\ref{thm:c2v57-no-spread} (defined there for the sequence
constant \(M=\Lambda\ge1\)).  There is an absolute covering constant
\(c_0>0\) such that for every \(t<T_*\) with \(T_*-t\le\pi^2/16\),
\begin{equation}
 \mathfrak M(t)\ \ge\ c_0\,\varepsilon_*\bigl(\nu,\max(\mathfrak Y(t),1)\bigr) .
 \label{eq:c2v59-floor}
\end{equation}
The function \(\varepsilon_*(\nu,\Lambda)\) is nonincreasing in \(\Lambda\), and
tracing the constants of V57 gives \(\varepsilon_*(\nu,\Lambda)\ge\exp(-C(\nu)\Lambda^{3})\)
for \(\Lambda\ge1\), so that
\begin{equation}
 \mathfrak M(t)\,\exp\bigl(C(\nu)\max(\mathfrak Y(t),1)^3\bigr)\ \ge\ c_0
 \qquad(t<T_*,\ T_*-t\le\pi^2/16):
 \label{eq:c2v59-tradeoff}
\end{equation}
a small parabolic-scale energy at any time forces a large scaled
enstrophy at that time.
\end{theorem}

\begin{proof}
Fix \(t_0=t\), \(\lambda=\lambda(t_0)\), and rescale:
\(U(y,s):=\lambda u(\lambda y,T_*+\lambda^2s)\) on the torus of side \(L=2\pi/\lambda\ge8\)
and \(s\in[-1,0)\), so that \(Y(-1)=\|\nabla U(-1)\|_2^2=\mathfrak Y(t_0)\).  Put
\(\Lambda_0:=\max(\mathfrak Y(t_0),1)\).  With the cells of V57, every cell has
\(e_\alpha(-1)\le\int_{2Q_\alpha}|U(-1)|^2\le c_0^{-1}\sup_x\int_{B_1(x)}|U(-1)|^2=c_0^{-1}\mathfrak M(t_0)\),
where \(c_0^{-1}\) is the number of unit balls needed to cover a cube of
side \(2\), and the last equality is the scaling of the Morrey
quantity.  Suppose \eqref{eq:c2v59-floor} fails.  Then
\(\max_\alpha e_\alpha(-1)<\varepsilon_*(\nu,\Lambda_0)\) and \(W(-1)\le Y(-1)\le\Lambda_0\), and
the proof of Theorem~\ref{thm:c2v57-no-spread} applies verbatim (it
used only these two facts, the torus size, and the global inequality
\eqref{eq:c2v56-global}): \(Y(s)\) is bounded on \([-1,0)\), hence
\(\|\Lambda u\|_2^2\) is bounded on \([t_0,T_*)\), and \(T_*\) is not singular,
a contradiction.  Monotonicity in \(\Lambda\): the constants \(b,b'\) of
Propositions~\ref{prop:c2v57-enstrophy}--\ref{prop:c2v57-energy} are
nondecreasing in \(\Lambda_0\), hence \(K,p,\Lambda\) are, and \(\varepsilon_*\) is
nonincreasing.  The exponential bound: \(a\) is of order \(\nu\);
\(b_1\), hence \(b\) and \(K\), is at most \(C(\nu)\Lambda_0^2\) (the largest power
comes from \((y^N)^3\le(125\Lambda_0)^2y^N\)); \(b'\) is at most
\(C(\nu)\Lambda_0^{1/2}\) (from \((e^N)^{3/4}(e^N+y^N)^{3/4}\le(e^N+y^N)^{3/2}\)),
so \(p\le C(\nu)\Lambda_0^{1/2}\) and \(\Lambda=b'(1+K)\le C(\nu)\Lambda_0^{5/2}\); then
\(\varepsilon_*=S_\kappa^{-1}K^pM^{-p}(Ke^\Lambda)^{-(1+p)}\ge\exp(-C(\nu)\Lambda_0^{5/2}\cdot\Lambda_0^{1/2})\cdot\exp(-C(\nu)\Lambda_0^{1/2}\log\Lambda_0)\ge\exp(-C(\nu)\Lambda_0^3)\).
\end{proof}

\begin{corollary}[Floors in the two rate classes]
\label{cor:c2v59-classes}
\begin{enumerate}[label=\textup{(\roman*)},leftmargin=3.2em]
\item (Intermediate class.)  At the window times,
      \(\mathfrak M(t_j)\ge c_0\varepsilon_*(\nu,M)\): Theorem~\ref{thm:c2v57-no-spread}.
\item (Type-I enstrophy rate.)  If \(\mathfrak Y(t)\le M\) on \([T_*-\delta,T_*)\),
      then \(\mathfrak M(t)\ge c_0\varepsilon_*(\nu,M)\) on that interval: the
      parabolic-scale energy concentration is persistent in time, with
      an explicit floor; this is the explicit form of the lower pinching
      \(c_\Phi\le\Phi_W\) of the first cycle (Theorem V91-pinched, whose
      constant came from compactness).
\item (Any singularity.)  Along any sequence \(t_k\uparrow T_*\) with
      \(\mathfrak M(t_k)\to0\) one has \(\mathfrak Y(t_k)\ge(C(\nu)^{-1}\log(c_0/\mathfrak M(t_k)))^{1/3}\to\infty\):
      a singularity whose parabolic energy vanishes along a sequence
      is of type II along that sequence, at least at the logarithmic
      rate displayed.
\end{enumerate}
\end{corollary}

\begin{proof}
\eqref{eq:c2v59-floor} and \eqref{eq:c2v59-tradeoff}.
\end{proof}

\begin{remark}[The floor is a statement about the energy, not the dissipation]
\label{rem:c2v59-energy-vs-dissipation}
The classical (E1)/(E7) criteria bound the solution when the
parabolic-cylinder \emph{dissipation} \(\lambda^{-1}\iint|\nabla u|^2\) is small.
The floor \eqref{eq:c2v59-floor} bounds it when the parabolic-ball
\emph{energy} at a single time is small, given the scaled enstrophy at
that time; the two are independent inputs (Proposition~\ref{prop:c2v26-independent}
recorded that the dissipation condition does not imply the energy
condition at the parabolic order).  The floor is not a regularity
criterion in a critical space alone: it needs \(\mathfrak Y(t)\) finite,
which is a subcritical quantity, and its dependence on \(\mathfrak Y\) is
exponential.
\end{remark}

% =============================================================================
\section{Cluster limits of a window}
\label{sec:c2v59-clusters}
% =============================================================================

Let \(u\) be in the intermediate class with constant \(M\), and at the
window time of \(U_j\) let \(\mathcal A_j\) be the energetic set
(\(1\le|\mathcal A_j|\le N_M\), Proposition~\ref{prop:c2v58-count}).

\begin{proposition}[Clusters and their limits]
\label{prop:c2v59-clusters}
Along a subsequence, the energetic sets \(\mathcal A_j\) split into
\(n\le N_M\) clusters \(\mathcal A_j^{(1)},\dots,\mathcal A_j^{(n)}\) with the
following properties: the diameter of each cluster is bounded
uniformly in \(j\); the distances between distinct clusters tend to
infinity; the number of cells in each cluster is constant along the
subsequence; and, for each cluster \(c\), with \(x^{(c)}_j\) the centre of
one of its cells, the rescalings centred at \(x^{(c)}_j\) converge on
\(B_R\times J_M\) for every \(R\), in the sense of
Theorem~\ref{thm:c2v42-compactness}, to a strong solution \(U'_c\) on
\(\R^3\times J_M\) with:
\begin{enumerate}[label=\textup{(\roman*)},leftmargin=3.2em]
\item \(\|\nabla U'_c(s)\|_{L^2(\R^3)}^2\le2M\) on \(J_M\), and \(U'_c(-1)\in\dot H^1(\R^3)\);
\item \(\int_{Q}|U'_c(-1)|^2\ge\varepsilon_*/2\) for the cell \(Q\) at the origin,
      and \(\int_{2Q_\alpha}|U'_c(-1)|^2\le27\cdot4\|U'_c(-1)\|_{L^6(2Q_\alpha)}^2\) for
      every cell;
\item (small far energy) \(\int_{Q_\alpha}|U'_c(-1)|^2\le\varepsilon_*/2\) for every
      cell \(Q_\alpha\) outside the bounded set \(D_c\) occupied by the
      cluster's cells: the limit has energy below the threshold in
      every far cell, but not zero energy;
\item the limit is nonzero, dissipation-active on the window in
      \(B_{R'}\) with level \(\varepsilon'\) (Theorem~\ref{thm:c2v48-energy-active}),
      and the limits of distinct clusters are obtained from rescalings
      whose centres separate: no relation between \(U'_c\) and \(U'_{c'}\)
      is asserted.
\end{enumerate}
\end{proposition}

\begin{proof}
Clustering: since \(|\mathcal A_j|\le N_M\), pass to a subsequence with
\(|\mathcal A_j|=N\) constant, and for each pair of cells decide, along a
further subsequence, whether their distance stays bounded or tends to
infinity; ``bounded distance'' is an equivalence relation along the
subsequence (a chain of bounded distances is bounded, with at most
\(N\) links), and its classes are the clusters.  Convergence:
Theorem~\ref{thm:c2v42-compactness} at the moving centre \(x^{(c)}_j\)
(the rescaling at a moving centre is the same field translated, and
the compactness is uniform in the centre).  (i) is weak lower
semicontinuity of the \(H^1\) bound \(Y_j\le2M\).  (ii): at the fixed
time \(-1\), \(U_j(-1)\) is bounded in \(H^1(B_R)\), so a further
subsequence converges strongly in \(L^2(B_R)\) (Rellich), and the energy
of the central cell passes to the limit; the cellwise \(L^6\) bound
passes by lower semicontinuity of the \(L^6\) norm.  (iii): the same
strong \(L^2_{\rm loc}\) convergence at time \(-1\) and the fact that every
non-energetic cell of \(U_j\) has energy below \(\varepsilon_*/2\), together
with the divergence of the distance to the other clusters, so that for
any fixed \(R\) and large \(j\) all cells of \(B_R\) outside \(D_c\) are
non-energetic.  (iv) is Theorem~\ref{thm:c2v48-energy-active} applied
to the centre map \(x^{(c)}_j\), which is energy-active with level
\(\varepsilon_*/16\).
\end{proof}

\begin{firewall}[Item (AC2): the trap does not run on the limit]
\label{fw:c2v59-AC2}
The limit \(U'_c\) is a strong solution on \(\R^3\times J_M\) only (no
extension to \([-1,0)\) is available without parabolic bounds beyond the
window, Definition~\ref{def:c2v26-parabolic}), and it has, in the
bounded set \(D_c\), cells with energy at least \(\varepsilon_*/2\), exactly
as the rescalings do; the far cells have energy at most \(\varepsilon_*/2\),
which is below the trap threshold, but the trap of V57 is global and
the localized trap of V58 has the hierarchy, the capture time and the
window limitation.  Nothing is gained by passing to the limit except
the removal of the torus and of the subsequence; item (AC2) is closed
without a new statement.  The statement that survives is
Corollary~\ref{cor:c2v59-classes}: the floor is a property of the
solution at every time, and its persistence under a type-I enstrophy
rate is explicit.  No global regularity claim is made.
\end{firewall}

% =============================================================================
\section{Integrated V59 conclusion and next acquisition order}
\label{sec:c2v59-conclusion}
% =============================================================================

\begin{theorem}[What V59 proves]
\label{thm:c2v59-conclusion}
Relative to the first cycle and V1--V58: the Morrey floor
\eqref{eq:c2v59-floor} with the trade-off \eqref{eq:c2v59-tradeoff}
and its consequences (Corollary~\ref{cor:c2v59-classes}), and the
cluster limits (Proposition~\ref{prop:c2v59-clusters}).  The full
Navier--Stokes stability/global-regularity theorem remains unproved.
\end{theorem}

\begin{proof}
The cited results.
\end{proof}

\begin{programme}[V60 acquisition order]
\label{prog:c2v60}
\begin{enumerate}[label=\textup{(AC\arabic*)},leftmargin=4.8em]
\item The compulsory companion audit: read the current SM and YM
      notes and record their digests and decision.
\item Consolidate V56--V59 as the trap chapter: the closed local
      inequality and delay theorem (V56), the cell trap and the
      exclusion of spread windows (V57), the finite energetic set and
      the conditional localized trap (V58), the Morrey floor and the
      cluster limits (V59); record the position after sixty versions
      and the direction for V61--V65.
\item Direction to examine at V60: the floor \eqref{eq:c2v59-floor}
      holds at every time; combine it with the time evolution of the
      concentrating ball.  The energy of the ball \(B_{\lambda(t)}(x(t))\)
      is at least \(c_0\varepsilon_*(\mathfrak Y(t))\lambda(t)\), and the local energy
      identity of the first cycle controls how the energy in a
      shrinking ball changes; determine whether the floor along a
      window sequence, together with the vanishing of the energy in
      shrinking balls (\(\int_{B_\lambda}|u|^2\le E_0\), and
      \(\lambda^{-1}\int_{B_\lambda}|u|^2\) bounded by \(C_S^2\mathfrak Y\)), yields an
      inequality between successive windows, and whether the gap
      summability of the first cycle (Theorem V37-cone-share) enters.
\end{enumerate}
\end{programme}

\begin{assessment}[Position after V59]
\label{ass:c2v59-position}
At every time before a singularity the parabolic-scale energy is at
least an explicit function of the scaled enstrophy; the windows of the
intermediate class are finite configurations of clusters with nonzero
strong limits.  No full stability or global-regularity claim is made.
\end{assessment}

% =============================================================================
\section*{Version V59 (second cycle) append-only record}
\addcontentsline{toc}{section}{Version V59 (second cycle) append-only record}
% =============================================================================

The complete V58 source of the second cycle was retained before this
continuation.  Its SHA--256 digest was
\begin{center}\scriptsize\ttfamily
0b086b4b10ab08edc26082fe0728d001f62d11ecd38073b1189095df45da7cfe.
\end{center}
V59 proved the Morrey floor at every time before a singularity and
the cluster limits of a window.  No full regularity claim is made.

% =============================================================================
\section{V60: compulsory companion audit, the trap chapter consolidated,
and the position after sixty versions}
\label{sec:c2v60-entry}
% =============================================================================

This is the sixtieth version of the second cycle.  The compulsory
review of the two companion programmes is performed first.  V56--V59
are then consolidated as the trap chapter.  The position after sixty
versions is recorded with the direction for V61--V65, and the first
step of the direction of Programme~\ref{prog:c2v60}(AC3) is taken: the
floor is combined with the energy of the concentrating ball.  No
global regularity claim is made.  The next compulsory companion audit
is V65.

% =============================================================================
\section{Compulsory V60 review of the current companion notes}
\label{sec:c2v60-companion-audit}
% =============================================================================

The current companion files in \texttt{latex\_notes} were inspected
read-only.  The frozen checkpoints are
\begin{align}
 &\texttt{Renewal\_SM\_Einstein\_Continuum\_Research\_Notes\_V9.tex},
 \notag\\[-2pt]
 &\qquad
 \texttt{7fbe1d2664eb4bc6995c96554b1f852625c0b54d4fa30f3e1de7b6a0c66f4f41},
 \label{eq:c2v60-SM-checkpoint}\\
 &\texttt{Renewal\_Yang\_Mills\_Mass\_Gap\_Research\_Notes\_V82.tex},
 \notag\\[-2pt]
 &\qquad
 \texttt{3bf7c95adc8dcd03bfef48aaff1cd5f277ed9724ab9e90e4442eb7047366ca2c}.
 \label{eq:c2v60-YM-checkpoint}
\end{align}
(The SM file advanced from V8 to V9 during the inspection; the digest
is that of V9, of 3991 lines.  A YM file named \(\ldots\)\texttt{V81.tex}
with the same digest as \eqref{eq:c2v60-YM-checkpoint} was also
present.)

\emph{SM V9 (seventh series).}  After closing the free
principal-symbol and doubler part of its main theorem at V8, the SM
programme derived the exact Hessian of the overlap heat trace with
respect to the link connection on an analytic gauge chart, proved the
vanishing of the free one-point functional and an exact Hessian Ward
identity, identified the unique leading local tensor, and recorded
that a termwise operator-norm bound of the polarization loses
transversality and yields neither a mass-term exclusion nor a
mesoscopic continuum bound.

\emph{YM V82.}  The YM programme constructed a shared-cluster
small-clock law with a percolation representation of its block
covariances (positive semidefinite connectivity matrices, endpoint
identities), proved a full squarefree logarithm and an exact
block-tree representation, characterised the exact percolation gates
seen by label roots, showed that conditioning all differentiated edges
at once makes the acyclic quotient an exact spanning gate, and
recorded by a counterexample that edgewise pivotality is false.

\begin{theorem}[V60 companion-audit decision]
\label{thm:c2v60-companion-decision}
No SM V9 Hessian, Ward-identity or tensor constant and no YM V82
percolation or block-tree bound is imported.  Neither bears on the
trap chapter or on the gates.
\end{theorem}

\begin{proof}
The SM results concern second variations of lattice heat traces; the
YM results concern percolation representations of a lattice clock.
None is a Navier--Stokes statement.  The structural analogy that the
YM programme's ``condition all differentiated edges at once'' resolves
a failure of termwise (edgewise) estimates is noted as parallel to the
trap chapter's replacement of cellwise inequalities by weighted maxima
(V57), but nothing quantitative transfers.
\end{proof}

% =============================================================================
\section{The trap chapter, consolidated}
\label{sec:c2v60-consolidated}
% =============================================================================

\begin{theorem}[The trap chapter, V56--V59]
\label{thm:c2v60-consolidated}
Let \(u\) be a smooth solution on the torus of side \(2\pi\) with first
singular time \(T_*\), and \(\mathfrak Y\), \(\mathfrak M\) the scaled
enstrophy and the parabolic Morrey energy \eqref{eq:c2v59-quantities}.
\begin{enumerate}[label=\textup{(\arabic*)},leftmargin=3.2em]
\item \emph{Closed local inequalities (V56).}  In the rescaled
      variables, the maximal local enstrophy \(\delta\) and the global
      enstrophy \(Y\) satisfy \eqref{eq:c2v56-delta} and
      \eqref{eq:c2v56-global}; a state with \(\delta\le\frac14\) and \(Y\le2M\)
      has no spike for a rescaled time \(T_{\rm delay}\) of order \(1/M\).
\item \emph{The cell trap (V57).}  With cells, a partition of unity and
      geometric weights, the weighted maxima \(W\) (enstrophy) and
      \(\varepsilon\) (energy) satisfy \(W'\le-aW^2/\varepsilon+b(W+\varepsilon)\),
      \(\varepsilon'\le b'(\varepsilon+W)\) while \(W\le\Lambda_0\), \(\varepsilon\le1\); the
      capture-and-trapping theorem bounds \(W\le K\varepsilon_1e^{\Lambda(s-s_0)}\)
      after the capture time, with \(\varepsilon_1\to0\) as \(\varepsilon(s_0)\to0\).
\item \emph{Morrey floor (V57, V59).}  For every \(t<T_*\) with
      \(T_*-t\le\pi^2/16\): \(\mathfrak M(t)\ge c_0\varepsilon_*(\nu,\max(\mathfrak Y(t),1))\),
      \(\varepsilon_*(\nu,\Lambda)\ge e^{-C(\nu)\Lambda^3}\).  In the intermediate class
      no window is spread; under a type-I enstrophy rate the floor is
      persistent; a singularity with vanishing parabolic energy along
      a sequence is of type II along it.
\item \emph{Finite configurations (V58, V59).}  At a window of the
      intermediate class, between \(1\) and \(N_M=1728C_S^6(2M)^3(\varepsilon_*/2)^{-3}\)
      cells carry energy at least \(\varepsilon_*/2\); they form finitely many
      clusters with nonzero strong window limits on the whole space
      whose far cells have energy below the threshold.
\item \emph{Localized trap on the window (V58).}  The weighted maxima
      over the quiet region at distance \(r\) from the energetic set
      satisfy the trap inequalities with the halo correction
      \(C_MN/r\), the hierarchy term \(S_4\bar e_{r/4}\), and an exponentially
      small boundary term integrable over the window; the trapping is
      conditional (Firewall~\ref{fw:c2v58-AC1}).
\end{enumerate}
\end{theorem}

\begin{proof}
(1) Corollary~\ref{cor:c2v56-delta}, Proposition~\ref{prop:c2v56-global},
Theorem~\ref{thm:c2v56-delay}.  (2) Propositions~\ref{prop:c2v57-enstrophy},
\ref{prop:c2v57-energy}, Theorem~\ref{thm:c2v57-trap}.  (3)
Theorems~\ref{thm:c2v57-no-spread}, \ref{thm:c2v59-floor},
Corollary~\ref{cor:c2v59-classes}.  (4) Proposition~\ref{prop:c2v58-count},
Corollary~\ref{cor:c2v58-configuration}, Proposition~\ref{prop:c2v59-clusters}.
(5) Proposition~\ref{prop:c2v58-localized}, Theorem~\ref{thm:c2v58-conditional}.
\end{proof}

% =============================================================================
\section{The floor and the energy of the concentrating ball}
\label{sec:c2v60-ball}
% =============================================================================

The first step of Programme~\ref{prog:c2v60}(AC3).  For \(t<T_*\) let
\(x(t)\) be a centre attaining the supremum in \(\mathfrak M(t)\) (the
supremum of a continuous function on the torus), and
\(\mathcal B(t):=\int_{B_{\lambda(t)}(x(t))}|u(t)|^2\), the energy of the
concentrating ball.

\begin{proposition}[The concentrating ball]
\label{prop:c2v60-ball}
For \(t<T_*\) with \(T_*-t\le\pi^2/16\),
\begin{equation}
 c_0\,\varepsilon_*\bigl(\nu,\max(\mathfrak Y(t),1)\bigr)\,\lambda(t)\ \le\ \mathcal B(t)\ \le\ C\,C_S^2\,\mathfrak Y(t)\,\lambda(t) ,
 \label{eq:c2v60-ball}
\end{equation}
so that \(\mathcal B(t)\to0\) at the rate \(\lambda(t)\) exactly when
\(\mathfrak Y(t)\) is bounded, and in general \(\mathcal B(t)\) lies between
\(e^{-C(\nu)\mathfrak Y(t)^3}\lambda(t)\) and \(C\mathfrak Y(t)\lambda(t)\).  Moreover the
energy of the concentrating ball is at least its parabolic share of
the total energy: with \(E(t)=\|u(t)\|_2^2\le E_0\),
\(\mathcal B(t)/E(t)\ge c_0\varepsilon_*\lambda(t)/E_0\), and the number of disjoint
parabolic balls needed to carry the total energy at the floor level is
at most \(E_0/(c_0\varepsilon_*\lambda(t))\), which diverges as \(t\uparrow T_*\); the
floor concerns one ball, not the distribution of the energy over the
parabolic balls.
\end{proposition}

\begin{proof}
\eqref{eq:c2v59-floor} and the H\"older--Sobolev upper bound of
Section~\ref{sec:c2v59-floor}, multiplied by \(\lambda(t)\).
\end{proof}

\begin{firewall}[What the ball's energy identity would need]
\label{fw:c2v60-ball}
The local energy identity on the ball \(B_{\lambda}(x)\) over a time interval
of length \(\lambda^2\) (the first cycle's Theorem V66-local-energy) bounds
the change of the ball's energy by the dissipation in the ball, the
cubic flux through its boundary, and the pressure work; along a
window sequence the dissipation in a ball of radius \(\lambda_j\) over the
window is at most \(\nu M\lambda_j\) (the type-I dissipation share), the
cubic flux is bounded by the local \(L^3\) norm, which the window
\(H^1\) bound controls, and the pressure work by the far-pressure and
local-pressure bounds of V47.  The identity gives that the ball's
energy changes by at most \(C(M,\nu)\lambda_j\) across the window, which is
of the same order as the floor \(c_0\varepsilon_*\lambda_j\) itself, so the
identity does not by itself show that the concentration persists
across the window without the compactness of V48, nor does it relate
successive windows, between which the enstrophy is unbounded.  The
inequality between successive windows sought by
Programme~\ref{prog:c2v60}(AC3) therefore requires either a bound on
the enstrophy between windows (unavailable in the intermediate class)
or a use of the gap summability (Theorem V37-cone-share) that the
present chapter does not provide; the item is recorded as open.
\end{firewall}

% =============================================================================
\section{Position after sixty versions and direction}
\label{sec:c2v60-position}
% =============================================================================

\begin{assessment}[Position after sixty versions of the second cycle]
\label{ass:c2v60-position-sixty}
The second cycle has: the self-similar profile and statistical
frameworks (V1--V15), the closed moment ledger (V16--V22), the
statistical dichotomy and coherence (V23--V25), the parabolic-scale
compactness and the invisibility of frequency escape (V26--V32), the
push budget and the constant-chasing gate with its explicit upper side
(V33--V41), the intermediate class with its type-I windows (V42--V45),
the spreading chapter (V46--V50), the window-structure chapter
(V51--V55), and now the trap chapter (V56--V59): at every time before
any singularity the parabolic-scale energy is bounded below by an
explicit function of the scaled enstrophy, and in the intermediate
class the windows are finite configurations of concentrated clusters.
The exclusion of spread windows is the strongest structural result of
the cycle.  What remains open is the same as at V55, now sharper: the
evolution between windows, where the enstrophy of the intermediate
class is unbounded, and the fate of a concentrated cluster after its
window.  For V61--V65 the direction is the cluster after the window:
the trap is a statement about quiet cells, and the cluster is the
complement; the tools to apply to the cluster itself are the first
cycle's Gaussian ledgers at the cluster's centre (bounded on the
window by V42), the backward cone (a spike after the window is at
parabolic distance at least \(\nu^3/(32C_1M^2)\) in rescaled time), and
the local energy identity with the far-pressure bound of V58, which
holds at all times with the polynomial weights: the far pressure at the
cluster is bounded by the weighted energies of all cells, an
\(L^1\)-type quantity in the energies that is bounded by the total energy
\(\lambda_j^{-1}E_0\) times \((1+r)^{-4}\) summed, hence by
\(C\lambda_j^{-1}E_0\), a bound that degenerates but is available beyond the
window.  Whether a cluster's energy identity with this far-pressure
bound gives a lower bound on the time needed to move energy of order
\(\varepsilon_*\) into or out of a parabolic ball, hence a lower bound on the
parabolic duration of the concentration beyond the window, is the
question for V61.  No global regularity claim is made.
\end{assessment}

% =============================================================================
\section{Integrated V60 conclusion and next acquisition order}
\label{sec:c2v60-conclusion}
% =============================================================================

\begin{theorem}[What V60 establishes]
\label{thm:c2v60-conclusion}
Relative to the first cycle and V1--V59: the companion-audit decision
(Theorem~\ref{thm:c2v60-companion-decision}), the consolidated trap
chapter (Theorem~\ref{thm:c2v60-consolidated}), and the concentrating
ball bounds \eqref{eq:c2v60-ball}.  The full Navier--Stokes
stability/global-regularity theorem remains unproved.
\end{theorem}

\begin{proof}
The cited results.
\end{proof}

\begin{programme}[V61 acquisition order]
\label{prog:c2v61}
\begin{enumerate}[label=\textup{(AC\arabic*)},leftmargin=4.8em]
\item Cluster energy identity beyond the window.  For a cluster cell
      \(\alpha\) (energy at least \(\varepsilon_*/2\) at \(s=-1\)), write the local
      energy identity on \(2Q_\alpha\) for \(s\in[-1,0)\) with the far pressure
      bounded by Lemma~\ref{lem:c2v58-far} (the weighted sum of all cell
      energies, at most \(C\lambda_j^{-1}E_0\), or better if the energies at
      distance \(r\) are bounded by the window bound \(a_M\) for \(r\) large)
      and the cubic and local-pressure terms by the local \(L^3\) norm;
      determine the minimal rescaled time in which the cell's energy
      can drop below \(\varepsilon_*/4\), and compare with the backward cone's
      delay.
\item The reverse: the minimal rescaled time in which a quiet cell can
      reach energy \(\varepsilon_*/2\) beyond the window, using the same
      identity with the quiet cell's trapped enstrophy while the trap
      holds; determine whether the energetic set can change only on a
      time scale bounded below, so that the finite configuration is
      stable for a parabolic time beyond the window.
\item The next compulsory companion audit is V65.
\end{enumerate}
\end{programme}

\begin{assessment}[Position after V60]
\label{ass:c2v60-position}
The trap chapter is consolidated; the concentrating ball's energy is
pinched between explicit multiples of the parabolic scale.  No full
stability or global-regularity claim is made.
\end{assessment}

% =============================================================================
\section*{Version V60 (second cycle) append-only record}
\addcontentsline{toc}{section}{Version V60 (second cycle) append-only record}
% =============================================================================

The complete V59 source of the second cycle was retained before this
continuation.  Its SHA--256 digest was
\begin{center}\scriptsize\ttfamily
284f666998334c6581ac1b40950cb64701409a585c05eabf3a158749769be32e.
\end{center}
V60 performed the compulsory companion audit and consolidated the trap
chapter.  No full regularity claim is made.

% =============================================================================
\section{V61: finite propagation of energy activity on a window}
\label{sec:c2v61-entry}
% =============================================================================

Programme~\ref{prog:c2v61} asked for the cluster energy identity beyond
the window and its reverse.  Beyond the window neither closes (the
firewall below records why in one paragraph), and the version goes
upstream to what the same identity gives \emph{on} the window, where
it closes the hierarchy of Firewall~\ref{fw:c2v58-AC1}(a) in a form the
localized trap did not reach: energy activity has a finite propagation
speed.  On a window, a cell that is regular in the sense of V54 needs a
time bounded below, depending only on \((\nu,M)\), to gain a fixed
fraction of the energy threshold, and a cell far from every non-quiet
cell gains, over the whole window, less than half the threshold.
Consequently the non-quiet set at any window time lies within an
explicit distance of a seed set of boundedly many cells fixed at the
window time, and every cell outside that neighbourhood is quiet
throughout the window.  A counting statement of V52/V54 is corrected in
passing (the homogeneous \(H^2\) seminorm must be used).  No global
regularity claim is made.  The next compulsory companion audit is V65.

\begin{remark}[Correction to the counting of regular cells in V52 and V54]
\label{rem:c2v61-correction}
Proposition~\ref{prop:c2v52-supnorm} and the definition of \(H\)-regular
cells in V54 count the cells with \(\int_{J_M}\|U\|_{H^2(2Q_\alpha)}^2\le H\)
using the bound \(\sum_\alpha\int_{J_M}\|U\|_{H^2(2Q_\alpha)}^2\le C(M,\nu)\).  With
the inhomogeneous norm this bound is false: the \(L^2\) part sums to
\(8\ell_M\|U\|_2^2=8\ell_M\lambda_j^{-1}\|u\|_2^2\), unbounded in \(j\).  The
correct statement uses the homogeneous seminorm
\(|U|_{H^2(rQ_\alpha)}^2:=\|\nabla U\|_{L^2(rQ_\alpha)}^2+\|\nabla^2U\|_{L^2(rQ_\alpha)}^2\):
\(\sum_\alpha\int_{J_M}|U|_{H^2(5Q_\alpha)}^2\le125(2M\ell_M+C_{\rm int}(4M/\nu+2M\ell_M))=:C_H\)
by the palinstrophy budget (Proposition~\ref{prop:c2v51-palinstrophy})
and the interior estimate \(\|\nabla^2U\|_{L^2(5Q)}^2\le C_{\rm int}(\|\Delta U\|_2^2+\|\nabla U\|_2^2)\)
summed over the cells.  Henceforth a cell is \emph{\(H\)-regular} if
\(\int_{J_M}|U|_{H^2(5Q_\alpha)}^2\le H\); all but \(C_H/H\) cells are
\(H\)-regular, and on an \(H\)-regular cell
\(\int_{J_M}\|U\|_{H^2(5Q_\alpha)}^2\le2(125a_M\ell_M+H)\) since
\(\|U\|_{L^2(5Q_\alpha)}^2\le125a_M\) on the window.  Every use in V52--V54
(the sup-norm bound and the influx bound) holds with \(H\) replaced by
\(2(125a_M\ell_M+H)\) in the constants; no conclusion changes.
\end{remark}

% =============================================================================
\section{Setting and two growth lemmas}
\label{sec:c2v61-lemmas}
% =============================================================================

Let \(u\) be in the intermediate class with constant \(M\ge1\), \(U\) a
window rescaling on \(J_M\), with the cells and notation of V57--V58,
\(a_M=8C_S^2M\) the window bound on cell energies, \(H:=1\).  Fix a
threshold
\begin{equation}
 \varepsilon_{\rm th}:=\min\bigl(\varepsilon_*/2,\ \varepsilon_\star(\nu,M)\bigr),
 \label{eq:c2v61-threshold}
\end{equation}
with \(\varepsilon_\star(\nu,M)\) the smallness required in
\eqref{eq:c2v61-far-small} below, \(N_M:=N_M(\varepsilon_{\rm th})\) the count of
Proposition~\ref{prop:c2v58-count}, \(\mathcal N(s):=\{\beta:e_\beta(s)\ge\varepsilon_{\rm th}\}\)
the non-quiet set (\(|\mathcal N(s)|\le N_M\) for all \(s\in J_M\)), and
\begin{equation}
 F_M:=S_4\varepsilon_{\rm th}+N_Ma_M ,\qquad
 \Gamma(\tau):=C_\sharp\Bigl[a_M^{9/8}\bigl(2(125a_M\ell_M+H)\bigr)^{3/8}\tau^{5/8}+\bigl(\nu a_M+F_Ma_M^{1/2}\bigr)\tau\Bigr],
 \label{eq:c2v61-crude}
\end{equation}
with \(C_\sharp\) the absolute constant of Lemma~\ref{lem:c2v61-crude};
\(\Gamma\) is increasing with \(\Gamma(0)=0\).  Define the \emph{contact time}
\(\tau_c=\tau_c(\nu,M)>0\) by \(\Gamma(\tau_c)=\frac38\varepsilon_{\rm th}\).

\begin{lemma}[Crude growth bound on regular cells]
\label{lem:c2v61-crude}
Let \(\beta\) be \(H\)-regular and \(\mathcal T\subset J_M\) measurable.  Then
\(\int_{\mathcal T}(e_\beta)'_+\,\dd s\le\Gamma(|\mathcal T|)\).  In particular an
\(H\)-regular cell gains at most \(\Gamma(\tau)\) of energy over any set
of times of measure \(\tau\).
\end{lemma}

\begin{proof}
The local energy identity gives
\(e'_\beta\le\nu\int|U|^2|\Delta\phi_\beta|+\int|U|^3|\nabla\phi_\beta|+2\int|P-c_\beta||U||\nabla\phi_\beta|\),
\(c_\beta:=P^{\rm far}_\beta(\beta)\).  The first term is at most
\(C\nu\cdot27a_M\).  The cubic term is at most
\(C\|U\|_{L^\infty(2Q_\beta)}\cdot27a_M\).  The local pressure term is at most
\(C\|P^{\rm loc}_\beta\|_{L^{3/2}}\|U\|_{L^3(2Q_\beta)}\le C\|U\|_{L^3(4Q_\beta)}^3\le C\|U\|_{L^\infty(4Q_\beta)}\cdot125a_M\).
The far pressure term is at most
\(2\sup_{2Q_\beta}|P^{\rm far}_\beta-c_\beta|\cdot\|U\|_{L^1(2Q_\beta)}\le C\|\nabla P^{\rm far}_\beta\|_{L^\infty(2Q_\beta)}(27a_M)^{1/2}\),
and \(\|\nabla P^{\rm far}_\beta\|_{L^\infty(2Q_\beta)}\le C(S_4\varepsilon_{\rm th}+N_Ma_M)=CF_M\)
by \eqref{eq:c2v58-far}, the quiet cells contributing at most \(S_4\varepsilon_{\rm th}\)
and the at most \(N_M\) non-quiet cells at most \(a_M\) each.  Hence
\(e'_\beta\le Ca_M\|U\|_{L^\infty(4Q_\beta)}+C(\nu a_M+F_Ma_M^{1/2})\).  Integrate
over \(\mathcal T\): by Gagliardo--Nirenberg on \(5Q_\beta\),
\(\|U\|_{L^\infty(4Q_\beta)}\le C\|U\|_{L^2(5Q_\beta)}^{1/4}\|U\|_{H^2(5Q_\beta)}^{3/4}\le Ca_M^{1/8}\|U\|_{H^2(5Q_\beta)}^{3/4}\),
and H\"older over \(\mathcal T\) with exponents \(\frac83,\frac85\) gives
\(\int_{\mathcal T}\|U\|_{L^\infty(4Q_\beta)}\le Ca_M^{1/8}(\int_{J_M}\|U\|_{H^2(5Q_\beta)}^2)^{3/8}|\mathcal T|^{5/8}\le Ca_M^{1/8}(2(125a_M\ell_M+H))^{3/8}|\mathcal T|^{5/8}\)
by Remark~\ref{rem:c2v61-correction}.
\end{proof}

\begin{lemma}[Growth bound far from the non-quiet set]
\label{lem:c2v61-far}
Let \(\beta\) be \(H\)-regular, \(\rho\ge3\), and \(\mathcal T\subset J_M\) a
measurable set of times at each of which \({\rm dist}(\beta,\mathcal N(s))\ge\rho\).
Then
\begin{equation}
 \int_{\mathcal T}(e_\beta)'_+\,\dd s\ \le\ C_\flat'\Bigl[\varepsilon_{\rm th}^{9/8}\bigl(2(125a_M\ell_M+H)\bigr)^{3/8}\ell_M^{5/8}+\nu\varepsilon_{\rm th}\ell_M+S_4\varepsilon_{\rm th}^{3/2}\ell_M\Bigr]
 +C_\flat'\,N_Ma_M\,\varepsilon_{\rm th}^{1/2}\,\ell_M\,(1+\rho)^{-4},
 \label{eq:c2v61-far}
\end{equation}
with \(C_\flat'\) absolute.  Let \(\varepsilon_\star(\nu,M)\) be such that for
\(\varepsilon_{\rm th}\le\varepsilon_\star\) the first bracket is at most
\(\varepsilon_{\rm th}/8\), that is
\begin{equation}
 C_\flat'\Bigl[\varepsilon_{\rm th}^{1/8}\bigl(2(125a_M\ell_M+H)\bigr)^{3/8}\ell_M^{5/8}+\nu\ell_M+S_4\varepsilon_{\rm th}^{1/2}\ell_M\Bigr]\le\tfrac18
 \label{eq:c2v61-far-small}
\end{equation}
(possible since \(\nu\ell_M=3\nu^4/(8C_AM^2)\) can be made at most
\(1/(24C_\flat')\) by taking \(M\) large, or otherwise by shrinking the
window to a sub-window \(\ell\le\ell_M\), which is harmless everywhere
below; the two \(\varepsilon_{\rm th}\)-terms are small for \(\varepsilon_{\rm th}\)
small), and define the \emph{contact radius}
\begin{equation}
 \rho_0=\rho_0(\nu,M):=\Bigl(4C_\flat'N_Ma_M\ell_M\,\varepsilon_{\rm th}^{-1/2}\Bigr)^{1/4} .
 \label{eq:c2v61-rho0}
\end{equation}
Then for \(\rho\ge\rho_0\) the total gain over \(\mathcal T\) is at most
\(\frac38\varepsilon_{\rm th}\).
\end{lemma}

\begin{proof}
At the times of \(\mathcal T\) every cell within distance \(2\) of \(\beta\) is
quiet, so \(\|U\|_{L^2(5Q_\beta)}^2\le125\varepsilon_{\rm th}\) there; the terms of the
previous proof become: first term \(\le C\nu\varepsilon_{\rm th}\); cubic and local
pressure \(\le C\varepsilon_{\rm th}\|U\|_{L^\infty(4Q_\beta)}\) with
\(\|U\|_{L^\infty(4Q_\beta)}\le C\varepsilon_{\rm th}^{1/8}\|U\|_{H^2(5Q_\beta)}^{3/4}\), integrated
as before over \(\mathcal T\) (\(|\mathcal T|\le\ell_M\)); far pressure
\(\le C(S_4\varepsilon_{\rm th}+N_Ma_M(1+\rho)^{-4})(27\varepsilon_{\rm th})^{1/2}\) by
\eqref{eq:c2v58-halo} with the current non-quiet set in the role of
\(\mathcal A\) (at most \(N_M\) cells, each of energy at most \(a_M\), all at
distance at least \(\rho\)), integrated over \(\mathcal T\).  The choice of
\(\rho_0\) makes the last term of \eqref{eq:c2v61-far} at most
\(\varepsilon_{\rm th}/4\).
\end{proof}

% =============================================================================
\section{Finite propagation}
\label{sec:c2v61-propagation}
% =============================================================================

\begin{theorem}[Finite propagation of energy activity on a window]
\label{thm:c2v61-propagation}
Let the \emph{seed set} be
\[
 \mathcal S:=\{\beta:\ e_\beta(-1)\ge\varepsilon_{\rm th}/4\}\ \cup\ \{\beta:\ \beta\text{ not \(H\)-regular}\},
 \qquad |\mathcal S|\le64N_M+C_H .
\]
Then for every \(s\in J_M\), with \(k(s):=\lfloor(s+1)/\tau_c\rfloor\),
\begin{equation}
 \mathcal N(s)\ \subset\ \{\beta:\ {\rm dist}(\beta,\mathcal S)\le k(s)\,\rho_0\},
 \qquad\text{hence}\qquad
 \mathcal N(s)\subset B(\mathcal S,R_{\rm inv}),\quad R_{\rm inv}:=\lfloor\ell_M/\tau_c\rfloor\rho_0 ,
 \label{eq:c2v61-propagation}
\end{equation}
and every cell at distance more than \(R_{\rm inv}\) from \(\mathcal S\) has
energy below \(\varepsilon_{\rm th}\) at every window time.  The constants
\(\tau_c,\rho_0,R_{\rm inv},N_M,C_H\) depend only on \((\nu,M)\).
\end{theorem}

\begin{proof}
The bound on \(|\mathcal S|\) is Proposition~\ref{prop:c2v58-count} with
threshold \(\varepsilon_{\rm th}/4\) (count \(64N_M\)) and
Remark~\ref{rem:c2v61-correction} (\(H=1\)).  We prove by induction on
\(k\ge1\) the statement
\[
 (\mathrm P_k)\colon\quad\text{every cell that is non-quiet at some time }s<-1+k\tau_c\text{ satisfies }{\rm dist}(\beta,\mathcal S)\le(k-1)\rho_0 .
\]
\((\mathrm P_1)\): let \(\beta\notin\mathcal S\) be non-quiet at some
\(s<-1+\tau_c\).  Then \(\beta\) is \(H\)-regular with \(e_\beta(-1)<\varepsilon_{\rm th}/4\),
so its energy gained at least \(\frac34\varepsilon_{\rm th}\) on \([-1,s]\), while
Lemma~\ref{lem:c2v61-crude} bounds the gain by \(\Gamma(s+1)<\Gamma(\tau_c)=\frac38\varepsilon_{\rm th}\):
impossible.  \((\mathrm P_k)\Rightarrow(\mathrm P_{k+1})\): let
\(\beta\notin\mathcal S\) be non-quiet first at \(s_\beta\in[-1+k\tau_c,-1+(k+1)\tau_c)\).
Split \([-1,s_\beta]=\mathcal T_{\rm near}\cup\mathcal T_{\rm far}\) according to
\({\rm dist}(\beta,\mathcal N(s))<\rho_0\) or \(\ge\rho_0\) (measurable sets, the
non-quiet set being determined by finitely many continuous functions).
The gain on \(\mathcal T_{\rm far}\) is at most \(\frac38\varepsilon_{\rm th}\)
(Lemma~\ref{lem:c2v61-far}, \(\rho_0\ge3\) by the size of the constants,
or by replacing \(\rho_0\) with \(\max(\rho_0,3)\)), and the total gain is
at least \(\frac34\varepsilon_{\rm th}\), so the gain on \(\mathcal T_{\rm near}\) is at
least \(\frac38\varepsilon_{\rm th}\), whence \(|\mathcal T_{\rm near}|\ge\tau_c\) by
Lemma~\ref{lem:c2v61-crude}.  Therefore \(\mathcal T_{\rm near}\), a subset of \([-1,s_\beta]\) of measure
at least \(\tau_c\), contains a time \(s'\le s_\beta-\tau_c<-1+k\tau_c\); at
\(s'\) there is
\(\gamma\in\mathcal N(s')\) with \(|\gamma-\beta|<\rho_0\).  By \((\mathrm P_k)\),
\({\rm dist}(\gamma,\mathcal S)\le(k-1)\rho_0\), so \({\rm dist}(\beta,\mathcal S)<k\rho_0\).
This is \((\mathrm P_{k+1})\), and \eqref{eq:c2v61-propagation} follows
with \(k(s)+1\) in the role of \(k\).
\end{proof}

\begin{corollary}[Windows are finite configurations at all window times]
\label{cor:c2v61-configuration}
At every window time \(s\in J_M\) the non-quiet set has at most \(N_M\)
cells and lies within distance \(R_{\rm inv}\) of the \(64N_M+C_H\) seed
cells fixed at the window time; the quiet exterior, all cells at
distance more than \(R_{\rm inv}\) from \(\mathcal S\), has energy below
\(\varepsilon_{\rm th}\) in every cell throughout the window.  Consequently, in
Proposition~\ref{prop:c2v58-localized} with \(\mathcal A\) replaced by
\(\mathcal S\), the hierarchy term satisfies \(\bar e_{r/4}\le\varepsilon_{\rm th}\)
for \(r\ge4(R_{\rm inv}+1)\), and the effective energy of the quiet region
at distance \(r\) is at most \(\varepsilon_r+S_4\varepsilon_{\rm th}+C_M(64N_M+C_H)/r+\kappa^{r/2}C_M\):
the hierarchy of Firewall~\ref{fw:c2v58-AC1}(a) is closed on the
window.
\end{corollary}

\begin{proof}
Theorem~\ref{thm:c2v61-propagation} and Proposition~\ref{prop:c2v58-count};
the localized inequalities of V58 were derived for an arbitrary set
\(\mathcal A\) of at most \(N\) cells with the window bounds, so \(\mathcal S\)
may replace \(\mathcal A\) with \(N=64N_M+C_H\).
\end{proof}

\begin{firewall}[Items (AC1)--(AC2) beyond the window]
\label{fw:c2v61-beyond}
Beyond the window the two lemmas lose their inputs: the crude bound
needs the cell energy bound \(a_M\) and the \(H^2\) budget, both from
\(Y\le2M\); the far bound needs the count \(N_M\), from the \(L^6\) bound.
Without them the far-pressure gradient on a cell is bounded only by
\(C\lambda_j^{-1}E_0\) (the total energy times the summed weights), the cubic
flux only by the local \(L^3\) norm, which nothing bounds, and the
non-quiet set can be the whole torus.  The cluster energy identity of
Programme~\ref{prog:c2v61}(AC1) therefore gives no minimal time for the
loss of a cluster's energy beyond the window, and (AC2) no minimal time
for a quiet cell to become energetic.  What the version establishes is
the finite propagation on the window and the closure of the hierarchy
there.  Two limitations of the localized trap remain: the capture time
(Firewall~\ref{fw:c2v58-AC1}(b)), and the window itself.  No global
regularity claim is made.
\end{firewall}

% =============================================================================
\section{Integrated V61 conclusion and next acquisition order}
\label{sec:c2v61-conclusion}
% =============================================================================

\begin{theorem}[What V61 proves]
\label{thm:c2v61-conclusion}
Relative to the first cycle and V1--V60: the corrected counting of
regular cells (Remark~\ref{rem:c2v61-correction}), the crude and far
growth bounds (Lemmas~\ref{lem:c2v61-crude}--\ref{lem:c2v61-far}), the
finite propagation theorem \eqref{eq:c2v61-propagation}, and the closure
of the hierarchy on the window (Corollary~\ref{cor:c2v61-configuration}).
The full Navier--Stokes stability/global-regularity theorem remains
unproved.
\end{theorem}

\begin{proof}
The cited results.
\end{proof}

\begin{programme}[V62 acquisition order]
\label{prog:c2v62}
\begin{enumerate}[label=\textup{(AC\arabic*)},leftmargin=4.8em]
\item The capture time on the window.  With the hierarchy closed, the
      localized trap at distance \(r\ge4(R_{\rm inv}+1)\) is limited by
      the capture phase: quiet cells may carry enstrophy up to \(2M\) at
      the window time.  Use the influx bound of V54 in reverse (an
      \(H\)-regular cell's enstrophy is \(\frac14\)-H\"older in time) and the
      dissipation lower bound of Lemma~\ref{lem:c2v57-interp}
      (\(d\ge y^2/(2e)-C^2y\) with \(e\le27\varepsilon_{\rm th}\)) to determine
      whether a quiet regular cell with enstrophy \(y\ge K\varepsilon_{\rm th}\)
      at the window time loses it within a time \(\ll\ell_M\), so that
      the capture phase is shorter than the window for the quiet
      exterior; if so, state the unconditional trapping of the quiet
      exterior on the window: enstrophy at most \(K\varepsilon_{\rm th}\) in every
      exterior cell after the capture time.
\item Consequence for the dissipation: if the exterior is trapped,
      the window dissipation \(\nu\int_{J_M}Y\ge c_L^2\nu^{3/2}\ell_M\) (the
      Leray floor of V43) is carried, up to \(K\varepsilon_{\rm th}\ell_M\) per
      exterior cell, by the \(O(N_MR_{\rm inv}^3)\) cells of the seed
      neighbourhood; record the resulting lower bound on the
      dissipation of the seed neighbourhood and its consequence for
      the active centres of V43.
\item The next compulsory companion audit is V65.
\end{enumerate}
\end{programme}

\begin{assessment}[Position after V61]
\label{ass:c2v61-position}
Energy activity on a window propagates at a finite speed from a
bounded seed set; the hierarchy of the localized trap is closed on the
window.  No full stability or global-regularity claim is made.
\end{assessment}

% =============================================================================
\section*{Version V61 (second cycle) append-only record}
\addcontentsline{toc}{section}{Version V61 (second cycle) append-only record}
% =============================================================================

The complete V60 source of the second cycle was retained before this
continuation.  Its SHA--256 digest was
\begin{center}\scriptsize\ttfamily
2c13df20c489b7df7cc3c88654156e3ac5d0edc39660ee61ef180722e71c8305.
\end{center}
V61 proved finite propagation of energy activity on a window and
corrected the counting of regular cells.  No full regularity claim is
made.

% =============================================================================
\section{V62: the quiet exterior is trapped on the window}
\label{sec:c2v62-entry}
% =============================================================================

Items (AC1) and (AC2) of Programme~\ref{prog:c2v62} are carried out.
(AC1) closes: with the hierarchy closed by V61, the effective energy of
the quiet exterior at distance \(r\) from the seed set is a
time-independent constant on the window, and the capture phase of the
weighted trap for the exterior lasts a time \(1/b\), which can be taken
at most half the window by enlarging the constant \(b\) (which only
weakens the inequality); hence after the capture time every exterior
cell has enstrophy at most an explicit multiple of the threshold plus
a halo decaying like the fourth power of the distance to the seed set.
This is unconditional.  (AC2) does not close, for a reason recorded as
a firewall: the trap bounds the enstrophy per cell uniformly, and the
exterior has unboundedly many cells, so its total enstrophy and
dissipation are not controlled by the trap.  A proposition records the
resulting form of the enstrophy inequality on the window (linear
growth outside the seed neighbourhood) and why it does not lengthen
the window.  No global regularity claim is made.  The next compulsory
companion audit is V65.

% =============================================================================
\section{The exterior trap}
\label{sec:c2v62-trap}
% =============================================================================

Notation of V58 and V61: window rescaling \(U\) on \(J_M\), seed set
\(\mathcal S\) with \(N':=|\mathcal S|\le64N_M+C_H\), invasion radius
\(R_{\rm inv}\), threshold \(\varepsilon_{\rm th}\), quiet regions
\(\mathcal Q_r=\{\alpha:{\rm dist}(\alpha,\mathcal S)\ge r\}\), weighted maxima \(W_r\),
\(\varepsilon_r\) of \eqref{eq:c2v58-localized} with \(\mathcal A=\mathcal S\), and
the constants \(a=a(\nu)\), \(b,b',C_M\) of Proposition~\ref{prop:c2v58-localized}
(with \(N=N'\)); replace \(b\) by \(\max(b,2/\ell_M)\), which weakens
\eqref{eq:c2v58-W}, so that \(\tau_{\rm cap}:=1/b\le\ell_M/2\).  Put
\(K:=4b/a\).

\begin{lemma}[Time-independent effective energy of the exterior]
\label{lem:c2v62-effective}
For \(r\ge2R_{\rm inv}+8\) and every \(s\in J_M\),
\begin{equation}
 \hat\varepsilon_r(s)\ \le\ \hat\varepsilon(r):=(S_\kappa+S_4)\varepsilon_{\rm th}+S_\kappa\kappa^{r-R_{\rm inv}}a_M+\frac{C_MN'}{r}+\kappa^{r/2}C_M ,
 \label{eq:c2v62-effective}
\end{equation}
and the far-pressure gradient on any cell of \(\mathcal Q_r\) is at most
\(C(S_4\varepsilon_{\rm th}+N_Ma_M(1+r-R_{\rm inv})^{-4})\).
\end{lemma}

\begin{proof}
By Theorem~\ref{thm:c2v61-propagation}, at every window time every
cell at distance more than \(R_{\rm inv}\) from \(\mathcal S\) has energy
below \(\varepsilon_{\rm th}\).  Hence \(\bar e_{r/4}\le\varepsilon_{\rm th}\) for
\(r/4>R_{\rm inv}\), and \(\varepsilon_r\le S_\kappa\varepsilon_{\rm th}+S_\kappa\kappa^{r-R_{\rm inv}}a_M\),
the cells of \(\mathcal Q_3\) within distance \(R_{\rm inv}\) of \(\mathcal S\)
being at distance at least \(r-R_{\rm inv}\) from \(\mathcal Q_r\) and having
energy at most \(a_M\).  Insert into \eqref{eq:c2v58-eff}.  The
far-pressure bound is \eqref{eq:c2v58-halo} with the current non-quiet
set (at most \(N_M\) cells within \(B(\mathcal S,R_{\rm inv})\)) in the role of
\(\mathcal A\).
\end{proof}

\begin{theorem}[Unconditional trapping of the quiet exterior on the window]
\label{thm:c2v62-exterior}
There is \(r_1=r_1(\nu,M)\ge2R_{\rm inv}+8\) such that for every \(r\ge r_1\)
and every \(s\in[-1+\tau_{\rm cap},-1+\ell_M]\),
\begin{equation}
 \max_{\alpha\in\mathcal Q_r}y_\alpha(s)\ \le\ W_r(s)\ \le\ 2K\,\hat\varepsilon(r)\ \le\ C(\nu,M)\Bigl(\varepsilon_{\rm th}+\frac{N'}{r}\Bigr).
 \label{eq:c2v62-exterior}
\end{equation}
In words: after the capture time, which is at most half the window,
every cell at distance at least \(r\) from the seed set carries
enstrophy at most an explicit multiple of the threshold plus a halo of
order \(N'/r\); with the finite propagation, for \(r\ge2R_{\rm inv}+8\) the
far pressure on such cells is of order \(N_Ma_M r^{-4}\), and the halo
term \(C_MN'/r\) in \eqref{eq:c2v62-effective} can be replaced by
\(C_MN_M(1+r-R_{\rm inv})^{-4}\) by rerunning Proposition~\ref{prop:c2v58-localized}
with the sharper far-pressure bound of Lemma~\ref{lem:c2v62-effective}.
\end{theorem}

\begin{proof}
By Lemma~\ref{lem:c2v62-effective}, \eqref{eq:c2v58-W} holds with
\(\hat\varepsilon_r\) replaced by the constant \(\hat\varepsilon=\hat\varepsilon(r)\)
(enlarging the denominator and the positive terms only weakens it) as
long as \(\varepsilon_r\le1\), which holds for \(r\ge r_1\) since
\(\hat\varepsilon(r)\to(S_\kappa+S_4)\varepsilon_{\rm th}<1\).  Let \(\Pi_r:=\int_{J_M}\pi_r\le\kappa^{r/2}C_M(\ell_M+4M)\)
and choose \(r_1\) so that \(\Pi_r\le\frac12K\hat\varepsilon(r)\) for \(r\ge r_1\).

\emph{Capture.}  On any interval where \(W_r\ge K\hat\varepsilon\):
\(b(W_r+\hat\varepsilon)\le\frac a4W_r^2/\hat\varepsilon+\frac a4W_r^2/\hat\varepsilon\) (as in
Theorem~\ref{thm:c2v57-trap}), so \(W_r'\le-\frac a2W_r^2/\hat\varepsilon+\pi_r\).
Let \(V:=W_r-\int_{-1}^s\pi_r\); then \(V\le W_r\) and
\(V'\le-\frac a2W_r^2/\hat\varepsilon\le-\frac a2V^2/\hat\varepsilon\) where \(V>0\), so
\(1/V(s)\ge1/V(-1)+\frac a{2\hat\varepsilon}(s+1)\) and
\(W_r(s)\le\frac{2\hat\varepsilon}{a(s+1)}+\Pi_r\) as long as \(W_r\ge K\hat\varepsilon\) on
\([-1,s]\).  At \(s=-1+\tau_{\rm cap}\), \(\frac{2\hat\varepsilon}{a\tau_{\rm cap}}=\frac{2b\hat\varepsilon}{a}=\frac K2\hat\varepsilon\),
so either \(W_r\) has dropped below \(K\hat\varepsilon\) before \(-1+\tau_{\rm cap}\)
or \(W_r(-1+\tau_{\rm cap})\le\frac K2\hat\varepsilon+\Pi_r\le K\hat\varepsilon\).

\emph{Trapping.}  After the first time \(s_c\le-1+\tau_{\rm cap}\) at which
\(W_r\le K\hat\varepsilon\): on \(\{W_r>K\hat\varepsilon\}\) one has \(W_r'\le\pi_r\), so
\(W_r(s)\le K\hat\varepsilon+\int_{s_c}^s\pi_r\le K\hat\varepsilon+\Pi_r\le\frac32K\hat\varepsilon\)
for \(s\ge s_c\) (a Lipschitz function that can increase only where it
exceeds \(K\hat\varepsilon\), and there at rate \(\pi_r\)).  The first
inequality of \eqref{eq:c2v62-exterior} is \eqref{eq:c2v57-comparison}
restricted to \(\mathcal Q_r\) (the weights include \(\kappa^0=1\) on the
cell itself); the last is the form of \(\hat\varepsilon(r)\) with
\(K=4b/a\) depending on \((\nu,M)\).
\end{proof}

\begin{corollary}[The exterior after capture]
\label{cor:c2v62-exterior}
On \([-1+\tau_{\rm cap},-1+\ell_M]\), outside \(B(\mathcal S,r_1)\), every cell has
energy below \(\varepsilon_{\rm th}\) and enstrophy at most
\(C(\nu,M)(\varepsilon_{\rm th}+N_M(1+r-R_{\rm inv})^{-4})\) at distance \(r\) from
\(\mathcal S\); every cell at distance at least \(r_2(\nu,M):=R_{\rm inv}+(N_M/\varepsilon_{\rm th})^{1/4}\)
has enstrophy at most \(2C(\nu,M)\varepsilon_{\rm th}\).  The seed neighbourhood
\(B(\mathcal S,r_2)\) has at most \(N_s:=N'(2r_2+1)^3\) cells.
\end{corollary}

\begin{proof}
Theorem~\ref{thm:c2v62-exterior} with the sharper halo, and counting.
\end{proof}

% =============================================================================
\section{The enstrophy inequality on the window with a trapped exterior}
\label{sec:c2v62-enstrophy}
% =============================================================================

\begin{proposition}[Linear growth outside the seed neighbourhood]
\label{prop:c2v62-linear}
On \([-1+\tau_{\rm cap},-1+\ell_M]\), with \(\delta_{\rm ext}:=2C(\nu,M)\varepsilon_{\rm th}\) the
exterior enstrophy level of Corollary~\ref{cor:c2v62-exterior} and
\(N_s\) the size of the seed neighbourhood,
\begin{equation}
 \frac{\dd Y}{\dd s}\ \le\ -\nu\,\mathcal E_2+C\Bigl(\nu^{-3}\delta_{\rm ext}^2+\delta_{\rm ext}^{1/2}\Bigr)Y+C\nu^{-3}N_s(2M)^3+C\bigl(N_s(2M)^3\bigr)^{1/4}Y^{3/4} .
 \label{eq:c2v62-linear}
\end{equation}
The growth is at most linear in \(Y\) up to the additive constant
\(A:=C\nu^{-3}N_s(2M)^3\); the time for \(Y\) to double from \(M\) under
\eqref{eq:c2v62-linear} is at least \(c\min(1/B,M/A)\) with
\(B:=C(\nu^{-3}\delta_{\rm ext}^2+\delta_{\rm ext}^{1/2})+C(N_s(2M)^3)^{1/4}M^{-1/4}\), which is
shorter than \(\ell_M\) because \(N_s\) is large (of order
\(N_M(N_M/\varepsilon_{\rm th})^{3/4}\)); the inequality does not lengthen the
window.
\end{proposition}

\begin{proof}
Split the stretching \(\sum_\alpha\|\nabla U\|_{L^3(Q_\alpha)}^3\le C\sum_\alpha y'^{3/4}_\alpha h'^{3/4}_\alpha\)
(proof of Proposition~\ref{prop:c2v56-global}) into the cells of
\(B(\mathcal S,r_2)\) and the rest.  For the rest, \(y'_\alpha\le27\delta_{\rm ext}\)
and H\"older over the cells give
\(C\delta_{\rm ext}^{1/2}Y^{1/4}(Y+\mathcal E_2)^{3/4}\), absorbed as in
Proposition~\ref{prop:c2v56-global}.  For the seed neighbourhood,
\(\sum y'^3\le N_s(27\cdot2M)^3\) crudely (each \(y'_\alpha\le Y\le2M\)), so
its stretching is at most \(C(N_s(2M)^3)^{1/4}(Y+\mathcal E_2)^{3/4}\le\frac\nu4\mathcal E_2+C\nu^{-3}N_s(2M)^3+C(N_s(2M)^3)^{1/4}Y^{3/4}\).
The doubling-time statement is the comparison ODE.
\end{proof}

\begin{firewall}[Item (AC2): the trap does not control sums over the exterior]
\label{fw:c2v62-AC2}
The exterior trap bounds \(y_\alpha\) uniformly by \(C\varepsilon_{\rm th}\) on
unboundedly many cells (the rescaled torus has \((2\pi/\lambda_j)^3\) cells),
so the exterior's total enstrophy \(\sum_{\rm ext}y_\alpha\le Y\le2M\) and
its dissipation are not bounded by the trap: the sum of \(\varepsilon_{\rm th}\)
over the cells diverges.  The trap is a bound on the ratio of
enstrophy to the \emph{maximal} energy, not to the cell's own energy;
a per-cell ratio bound \(y_\alpha\le Ke_\alpha\) would give
\(\sum_{\rm ext}y_\alpha\le K\sum_{\rm ext}e_\alpha\), still unbounded since the
total energy \(\lambda_j^{-1}\|u\|_2^2\) is.  Hence the Leray floor of the
window dissipation cannot be attributed to the seed neighbourhood, and
(AC2) is closed negatively.  What survives is
Corollary~\ref{cor:c2v62-exterior}: after the capture time the exterior
is uniformly weak in energy and enstrophy, and the enstrophy
inequality is linear outside the seed neighbourhood.  No global
regularity claim is made.
\end{firewall}

% =============================================================================
\section{Integrated V62 conclusion and next acquisition order}
\label{sec:c2v62-conclusion}
% =============================================================================

\begin{theorem}[What V62 proves]
\label{thm:c2v62-conclusion}
Relative to the first cycle and V1--V61: the time-independent
effective energy \eqref{eq:c2v62-effective}, the unconditional exterior
trap \eqref{eq:c2v62-exterior} with its corollary, and the linear
enstrophy inequality \eqref{eq:c2v62-linear}.  The full Navier--Stokes
stability/global-regularity theorem remains unproved.
\end{theorem}

\begin{proof}
The cited results.
\end{proof}

\begin{programme}[V63 acquisition order]
\label{prog:c2v63}
\begin{enumerate}[label=\textup{(AC\arabic*)},leftmargin=4.8em]
\item The seed neighbourhood from inside.  The crude bound \(y'_\alpha\le2M\)
      on the \(N_s\) cells of \(B(\mathcal S,r_2)\) is the obstruction in
      Proposition~\ref{prop:c2v62-linear}.  Determine whether the trap
      can be run on the quiet cells inside the seed neighbourhood at
      their own halo level, using the finite propagation in the form
      \(\mathcal N(s)\subset B(\mathcal S,k(s)\rho_0)\) (shrinking regions
      \(\mathcal Q(s)\) at whose jump times the maximum can only drop),
      so that the cells at distance \(\rho\) from the current non-quiet
      set have enstrophy at most \(C(\varepsilon_{\rm th}+N_Ma_M\rho^{-4})\) after
      capture; if so, sum the halo (\(\sum_\rho26\rho^2\rho^{-4}<\infty\)) to
      bound the enstrophy of the seed neighbourhood minus the non-quiet
      cells by \(C(\nu,M)\), and the stretching of the seed
      neighbourhood by a constant depending on \(N_M\) but not on \(N_s\).
\item If (AC1) succeeds, redo Proposition~\ref{prop:c2v62-linear} with
      the improved constant and determine whether the window can be
      lengthened, and by how much, relative to \(\ell_M\).
\item The next compulsory companion audit is V65.
\end{enumerate}
\end{programme}

\begin{assessment}[Position after V62]
\label{ass:c2v62-position}
After the capture time the quiet exterior of a window is uniformly
weak in energy and enstrophy, unconditionally; the enstrophy
inequality is linear outside the seed neighbourhood.  No full
stability or global-regularity claim is made.
\end{assessment}

% =============================================================================
\section*{Version V62 (second cycle) append-only record}
\addcontentsline{toc}{section}{Version V62 (second cycle) append-only record}
% =============================================================================

The complete V61 source of the second cycle was retained before this
continuation.  Its SHA--256 digest was
\begin{center}\scriptsize\ttfamily
bb40a1b4bb905afafc97435934ac713729849adc8eb0225cb144d0c40e2daa2d.
\end{center}
V62 proved the unconditional trapping of the quiet exterior on the
window.  No full regularity claim is made.

% =============================================================================
\section{V63: the seed neighbourhood from inside, and the trap chapter
under a type-I enstrophy rate}
\label{sec:c2v63-entry}
% =============================================================================

Item (AC1) of Programme~\ref{prog:c2v63} is examined and closed
negatively, with the obstruction identified: a trap relative to the
current non-quiet set needs regions that shrink in time (so that the
maximum can only drop at the jump times), but the non-quiet set moves
non-monotonically, and a region defined by the distance to it can
enlarge, admitting cells of enstrophy up to the window bound; the
monotone envelope \(B(\mathcal S,k(s)\rho_0)\) of V61 gives regions that
shrink, but then the halo is measured from the envelope and the result
is V62 again.  Item (AC2) is therefore moot.  The version then goes
upstream and records what the trap chapter (V56--V62) says under a
type-I enstrophy rate on an interval, where every time is a window
time: at every time a finite configuration of at most \(N_M\) energetic
parabolic cells, finite propagation of energy activity between
consecutive parabolic times, and a uniformly weak exterior after the
capture time; and the corresponding statement for the ancient limit of
the first cycle, whose far field is uniformly diffuse at every time.
No global regularity claim is made.  The next compulsory companion
audit is V65.

% =============================================================================
\section{The obstruction to a trap relative to the current non-quiet set}
\label{sec:c2v63-obstruction}
% =============================================================================

\begin{proposition}[Two regions, two outcomes]
\label{prop:c2v63-regions}
On a window, let \(\rho\ge\rho_{\min}(\nu,M)\) (the radius beyond which the
boundary and interpolation pollution terms are below the threshold
level, of order \(\log(1/\varepsilon_{\rm th})\)).
\begin{enumerate}[label=\textup{(\roman*)},leftmargin=3.2em]
\item (Envelope regions.)  For \(\mathcal Q^{\rm env}_\rho(s):=\{\alpha:{\rm dist}(\alpha,\mathcal S)\ge\rho+k(s)\rho_0\}\),
      \(k(s)=\lfloor(s+1)/\tau_c\rfloor\), the weighted maximum
      \(W^{\rm env}_\rho(s):=\max_{\alpha\in\mathcal Q^{\rm env}_\rho(s)}w_\alpha\) is
      Lipschitz between the jump times \(-1+k\tau_c\), drops (weakly)
      at the jump times, and satisfies the trap inequality with the
      constant effective energy \(\hat\varepsilon(\rho)\) of
      Lemma~\ref{lem:c2v62-effective} (with \(R_{\rm inv}\) replaced by
      \(k(s)\rho_0\)); the capture-and-trapping argument of
      Theorem~\ref{thm:c2v62-exterior} applies and gives
      \(W^{\rm env}_\rho(s)\le2K\hat\varepsilon(\rho)\) for \(s\ge-1+\tau_{\rm cap}\).
      Since \(k(s)\rho_0\le R_{\rm inv}\), this is Theorem~\ref{thm:c2v62-exterior}
      with \(r=\rho+k(s)\rho_0\): no cell closer than \(k(s)\rho_0+\rho\) to
      \(\mathcal S\) is reached.
\item (Current regions.)  For \(\mathcal Q^{\rm cur}_\rho(s):=\{\alpha:{\rm dist}(\alpha,\mathcal N(s))\ge\rho\}\),
      the weighted maximum \(W^{\rm cur}_\rho\) can jump upward: when a
      non-quiet cell becomes quiet, the cells within distance \(\rho\) of
      it enter the region with enstrophy bounded only by the window
      bound \(2M\); each such entry restarts the capture phase, which
      lasts \(\tau_{\rm cap}\).  The non-quiet set changes at the crossing
      times of the finitely many continuous functions \(e_\beta-\varepsilon_{\rm th}\),
      whose number on the window is not bounded by the series (a cell
      may oscillate about the threshold).  Hence no trapping statement
      follows for \(W^{\rm cur}_\rho\) on the window.
\end{enumerate}
\end{proposition}

\begin{proof}
(i) The regions are nested decreasing in \(s\), so the maximum over
them is nonincreasing at the jump times; between jumps the region is
fixed and Proposition~\ref{prop:c2v58-localized} applies with
\(\mathcal A=\mathcal S\) and, for the far pressure, the current non-quiet set
inside \(B(\mathcal S,k(s)\rho_0)\) (Theorem~\ref{thm:c2v61-propagation}),
at distance at least \(\rho\) from the region; the pollution terms are
as in V62.  The capture-and-trapping argument uses only the
differential inequality and the fact that downward jumps do not
increase the maximum.  (ii) is the description of the region's
evolution; the bound \(2M\) on entering cells' enstrophy is the window
bound, and the capture phase of Theorem~\ref{thm:c2v62-exterior} needs
\(\tau_{\rm cap}\) from any level up to \(2M\).
\end{proof}

\begin{firewall}[Items (AC1)--(AC2)]
\label{fw:c2v63-AC}
A halo-level trap relative to the current non-quiet set would require
a hysteresis (two thresholds, so that a cell leaving the non-quiet set
is counted as non-quiet until its energy drops below a lower threshold,
which takes a time bounded below by the reverse of
Lemma~\ref{lem:c2v61-crude}) and a bound on the number of threshold
crossings on the window, which the series does not have: the energy of
a cell can oscillate about a threshold arbitrarily often within the
window bounds.  Without it, the seed neighbourhood's enstrophy is
bounded only by the window bound on each of its \(N_s\) cells, and
Proposition~\ref{prop:c2v62-linear} is not improved.  (AC1) and (AC2)
are closed negatively.
\end{firewall}

% =============================================================================
\section{The trap chapter under a type-I enstrophy rate}
\label{sec:c2v63-typeI}
% =============================================================================

Let now \(u\) satisfy \(\mathfrak Y(t)=(T_*-t)^{1/2}\|\Lambda u(t)\|_2^2\le M\) on
\([T_*-\delta_0,T_*)\) (a type-I enstrophy rate; the first cycle's class,
Theorem V46-upgrade of the first cycle gives then the \(L^\infty\)
type-I rate as well).  Every \(t\) in the interval is a window time
with the same constant, and the windows \([t,t+\ell_M(T_*-t)]\) cover the
interval.  For \(t\) in the interval let \(\lambda=\lambda(t)\), \(U_t\) the
rescaling at scale \(\lambda\) (time \(s=-1\) corresponding to \(t\)), and use
the cells of side \(\lambda\) in the original variables.

\begin{theorem}[Type-I structure at every time]
\label{thm:c2v63-typeI}
Under a type-I enstrophy rate with constant \(M\), for every
\(t\in[T_*-\delta_0,T_*)\) with \(\lambda(t)\le\pi/4\):
\begin{enumerate}[label=\textup{(\roman*)},leftmargin=3.2em]
\item (Finite configuration.)  The number of parabolic cells (side
      \(\lambda(t)\)) with energy at least \(\varepsilon_{\rm th}\lambda(t)\) is between
      \(1\) and \(N_M\); those with energy at least \(\varepsilon_{\rm th}\lambda(t)/4\)
      number at most \(64N_M\).
\item (Finite propagation.)  On \([t,t+\ell_M\lambda^2]\) the non-quiet set
      (cells with energy at least \(\varepsilon_{\rm th}\lambda\)) stays within
      distance \(R_{\rm inv}\lambda\) of the seed set fixed at time \(t\)
      (the at most \(64N_M+C_H\) cells that are energetic at level
      \(\varepsilon_{\rm th}\lambda/4\) or not \(H\)-regular on the window), and moves
      at most \(\rho_0\lambda\) per time \(\tau_c\lambda^2\).
\item (Weak exterior.)  On \([t+\tau_{\rm cap}\lambda^2,t+\ell_M\lambda^2]\) every cell at
      distance at least \(r_2\lambda\) from the seed set has energy below
      \(\varepsilon_{\rm th}\lambda\) and enstrophy \(\int_{\rm cell}|\nabla u|^2\le2C(\nu,M)\varepsilon_{\rm th}\lambda^{-1}\).
\item (Persistent floor.)  \(\mathfrak M(t)\ge c_0\varepsilon_*(\nu,M)\) for all such \(t\).
\end{enumerate}
All constants depend only on \((\nu,M)\).
\end{theorem}

\begin{proof}
Each statement is the corresponding window statement
(Proposition~\ref{prop:c2v58-count}, Theorem~\ref{thm:c2v61-propagation},
Corollary~\ref{cor:c2v62-exterior}, Theorem~\ref{thm:c2v59-floor})
applied to the rescaling at \(t\), which is a window rescaling with
constant \(M\) because \(\mathfrak Y(t)\le M\); the scalings of energy
(\(\lambda\)), enstrophy (\(\lambda^{-1}\)), length (\(\lambda\)) and time (\(\lambda^2\))
convert the cell statements to the original variables.
\end{proof}

\begin{corollary}[The ancient limit is uniformly diffuse at infinity]
\label{cor:c2v63-ancient}
Let \(W\) be an ancient limit of the first cycle at a type-I singular
point (a nonzero element of the class \(\mathcal A_\nu(C)\), with
\(\|\nabla W(s)\|_2^2\le C(-s)^{-1/2}\)).  Then for every \(s<0\), with
\(\lambda=(-s)^{1/2}\) and the cells of side \(\lambda\): at most \(N_C\) cells carry
energy at least \(\varepsilon_{\rm th}\lambda\), and every cell at distance at least
\(r_2\lambda\) from the seed set of the time \(s\) has, on
\([s+\tau_{\rm cap}\lambda^2,s+\ell_C\lambda^2]\), energy below \(\varepsilon_{\rm th}\lambda\) and
enstrophy at most \(2C(\nu,C)\varepsilon_{\rm th}\lambda^{-1}\).  In particular the
far field of the ancient limit carries, per parabolic cell, at most
the fraction \(2C(\nu,C)\varepsilon_{\rm th}/C\) of the total enstrophy bound
\(C\lambda^{-1}\), at every time, on the corresponding sub-window.
\end{corollary}

\begin{proof}
The rescalings converging to \(W\) satisfy the window bounds at every
scale with the constant \(C\); the statements of
Theorem~\ref{thm:c2v63-typeI} hold for them and pass to the limit:
the count by strong \(L^2_{\rm loc}\) convergence at fixed times, the
energy bounds likewise, and the enstrophy bounds by weak lower
semicontinuity.
\end{proof}

\begin{assessment}[What the type-I structure adds to the gate]
\label{ass:c2v63-gate}
The type-I gate of the first cycle (Theorem V41-gate and its
restatements) asks whether a nonzero element of \(\mathcal A_\nu(C)\) can
exist.  The trap chapter adds to the description of such an element:
at every time it is a finite configuration of at most \(N_C\)
energetic parabolic cells, its energy activity propagates at a finite
parabolic speed, and its far field is uniformly diffuse (small energy
and small enstrophy per parabolic cell) after a parabolic capture
time.  These are consistent with the first cycle's description
(vorticity nonzero on every exterior region at every time, V63 of the
first cycle; uniform concentration at homogeneous points, V82; satellites
at lacunary points, V77--V79): the far field is nonzero but weak.  The
gate is not closed: a finite configuration with a weak far field is
exactly what a self-similar-like ancient solution would look like,
and the exclusion of self-similar elements (E9) does not extend to
this description.  No global regularity claim is made.
\end{assessment}

% =============================================================================
\section{Integrated V63 conclusion and next acquisition order}
\label{sec:c2v63-conclusion}
% =============================================================================

\begin{theorem}[What V63 proves]
\label{thm:c2v63-conclusion}
Relative to the first cycle and V1--V62: the two-region proposition
(Proposition~\ref{prop:c2v63-regions}) with the negative closure of
(AC1)--(AC2), the type-I structure theorem
(Theorem~\ref{thm:c2v63-typeI}), and the uniform diffuseness of the
ancient limit's far field (Corollary~\ref{cor:c2v63-ancient}).  The
full Navier--Stokes stability/global-regularity theorem remains
unproved.
\end{theorem}

\begin{proof}
The cited results.
\end{proof}

\begin{programme}[V64 acquisition order]
\label{prog:c2v64}
\begin{enumerate}[label=\textup{(AC\arabic*)},leftmargin=4.8em]
\item Hysteresis.  Define the non-quiet set with two thresholds
      (entry at \(\varepsilon_{\rm th}\), exit at \(\varepsilon_{\rm th}/2\)) and determine
      whether the reverse of Lemma~\ref{lem:c2v61-crude} (an
      \(H\)-regular cell loses at most \(\Gamma(\tau)\) of energy in time
      \(\tau\), by the same identity with the dissipation term retained
      and bounded by the \(H^2\) budget) bounds the number of
      entry--exit cycles of a cell on the window by \(\ell_M/\tau_c\), so that
      the current-region trap of Proposition~\ref{prop:c2v63-regions}(ii)
      restarts at most boundedly many times; if so, state the
      halo-level trap and the improved seed-neighbourhood enstrophy
      bound.
\item Prepare V65: the audit and the consolidation of V61--V64 as the
      propagation chapter, with the position after sixty-five versions.
\item The next compulsory companion audit is V65.
\end{enumerate}
\end{programme}

\begin{assessment}[Position after V63]
\label{ass:c2v63-position}
The trap relative to the current non-quiet set is blocked by
threshold oscillation; under a type-I enstrophy rate the trap chapter
gives a finite configuration with finite propagation and a uniformly
diffuse far field at every time.  No full stability or
global-regularity claim is made.
\end{assessment}

% =============================================================================
\section*{Version V63 (second cycle) append-only record}
\addcontentsline{toc}{section}{Version V63 (second cycle) append-only record}
% =============================================================================

The complete V62 source of the second cycle was retained before this
continuation.  Its SHA--256 digest was
\begin{center}\scriptsize\ttfamily
a62a90e16b2e1e4543f75be0de05af034c4c751c8095e1466a0a1d946838406c.
\end{center}
V63 closed the current-region trap negatively and recorded the type-I
structure at every time.  No full regularity claim is made.

% =============================================================================
\section{V64: persistence of energetic cells, bounded oscillation, and the
failure of the hysteresis budget}
\label{sec:c2v64-entry}
% =============================================================================

Item (AC1) of Programme~\ref{prog:c2v64} is carried out.  The reverse
of the crude growth lemma holds: a regular cell loses at most an
explicit amount of energy in a given time, so an energetic cell stays
above half its level for a time bounded below, and the number of
threshold crossings of a regular cell on a window is bounded.  The
hysteresis count is therefore finite, but the budget fails: the number
of cells that can ever be marked is of the order of the seed
neighbourhood, each recapture costs the capture time, and the product
exceeds the window by many orders of magnitude for every admissible
choice of the constants, because the invasion radius is a power of the
inverse threshold.  The halo-level trap relative to the current
non-quiet set is closed negatively, and with it the improvement of the
seed-neighbourhood enstrophy bound.  What survives is the persistence
statement, which is recorded as the time-regularity of the finite
configuration.  No global regularity claim is made.  The next
compulsory companion audit is V65.

% =============================================================================
\section{Reverse growth bound and persistence}
\label{sec:c2v64-persistence}
% =============================================================================

Notation of V61.

\begin{lemma}[Crude loss bound on regular cells]
\label{lem:c2v64-loss}
Let \(\beta\) be \(H\)-regular and \(\mathcal T\subset J_M\) measurable.  Then
\(\int_{\mathcal T}(e_\beta)'_-\,\dd s\le\Gamma(|\mathcal T|)+4\nu M|\mathcal T|=:\Gamma^-(|\mathcal T|)\).
\end{lemma}

\begin{proof}
The local energy identity has, besides the terms bounded in
Lemma~\ref{lem:c2v61-crude} (which bound \((e_\beta)'\) from above and,
with the same bounds on their absolute values, from below), the
dissipation \(-2\nu y_\beta\ge-2\nu Y\ge-4\nu M\) on the window.
\end{proof}

\begin{proposition}[Persistence of energetic cells and bounded oscillation]
\label{prop:c2v64-persistence}
Let \(\tau'\) and \(\tau''\) be defined by \(\Gamma(\tau')=\varepsilon_{\rm th}/2\) and
\(\Gamma^-(\tau'')=\varepsilon_{\rm th}/2\), and \(\tau_h:=\min(\tau',\tau'')\).  For an
\(H\)-regular cell \(\beta\):
\begin{enumerate}[label=\textup{(\roman*)},leftmargin=3.2em]
\item if \(e_\beta(s_0)\ge\varepsilon_{\rm th}\) then \(e_\beta\ge\varepsilon_{\rm th}/2\) on
      \([s_0,s_0+\tau'']\cap J_M\): an energetic cell stays half-energetic
      for the time \(\tau''\);
\item if \(e_\beta(s_0)<\varepsilon_{\rm th}/2\) then \(e_\beta<\varepsilon_{\rm th}\) on
      \([s_0,s_0+\tau']\cap J_M\);
\item the number of disjoint intervals on which \(e_\beta\) rises from
      below \(\varepsilon_{\rm th}/2\) to at least \(\varepsilon_{\rm th}\), or falls from at
      least \(\varepsilon_{\rm th}\) to below \(\varepsilon_{\rm th}/2\), is at most
      \(\ell_M/\tau_h+1\).
\end{enumerate}
In particular the finite configuration of a window is Lipschitz in
time in the following sense: the set of cells with energy at least
\(\varepsilon_{\rm th}/2\) at time \(s\) contains the energetic set of every time in
\([s-\tau'',s]\) and is contained in the set of cells energetic at
level \(\varepsilon_{\rm th}/4\) at every time in \([s,s+\tau''']\), \(\Gamma^-(\tau''')=\varepsilon_{\rm th}/4\),
up to the \(C_H\) non-regular cells.
\end{proposition}

\begin{proof}
(i) and (ii) are Lemmas~\ref{lem:c2v64-loss} and \ref{lem:c2v61-crude}
with \(\mathcal T\) the interval.  (iii): each rise or fall takes time at
least \(\tau_h\), and the intervals are disjoint within \(J_M\).  The last
statement is (i) applied at levels \(\varepsilon_{\rm th}\) and \(\varepsilon_{\rm th}/2\).
\end{proof}

% =============================================================================
\section{The hysteresis count and its budget}
\label{sec:c2v64-hysteresis}
% =============================================================================

\begin{proposition}[Hysteresis: finitely many restarts, but too many]
\label{prop:c2v64-hysteresis}
Mark a cell at the first time its energy reaches \(\varepsilon_{\rm th}\), unmark
it when its energy falls below \(\varepsilon_{\rm th}/2\), and let the non-regular
cells be permanently marked.  Then on the window:
\begin{enumerate}[label=\textup{(\roman*)},leftmargin=3.2em]
\item the marked set has at most \(8N_M+C_H\) cells at every time;
\item the marked set changes at most \(E_{\max}:=N'(2R_{\rm inv}+1)^3(\ell_M/\tau_h+1)\)
      times on the window;
\item the current-region trap of Proposition~\ref{prop:c2v63-regions}(ii),
      run with the marked set in place of the non-quiet set, is
      restarted at most \(E_{\max}\) times, each restart costing the
      capture time \(\tau_{\rm cap}\); a trapped stretch exists only if
      \(E_{\max}\tau_{\rm cap}<\ell_M\).
\end{enumerate}
The inequality in (iii) fails: with \(\tau_{\rm cap}\le\ell_M/2\) at best of
order \(\nu^3/(C M^2)\), \(R_{\rm inv}=\lfloor\ell_M/\tau_c\rfloor\rho_0\), and
\(\tau_c,\tau_h\) of order \(\varepsilon_{\rm th}^{8/5}\) (from \(\Gamma(\tau)\approx C a_M^{9/8}(\ldots)^{3/8}\tau^{5/8}\)
for small \(\tau\)), one has \(E_{\max}\tau_{\rm cap}/\ell_M\) of order
\(N'\ell_M^{3}\rho_0^3\tau_c^{-3}\tau_h^{-1}\tau_{\rm cap}/\ell_M\gtrsim\varepsilon_{\rm th}^{-6}\), which is
astronomically large.
\end{proposition}

\begin{proof}
(i) Marked regular cells have energy at least \(\varepsilon_{\rm th}/2\), and
Proposition~\ref{prop:c2v58-count} at that threshold gives \(8N_M\).
(ii) A cell that is ever marked is non-quiet at some window time,
hence within \(B(\mathcal S,R_{\rm inv})\) (Theorem~\ref{thm:c2v61-propagation}),
at most \(N'(2R_{\rm inv}+1)^3\) cells; each changes state at most
\(\ell_M/\tau_h+1\) times (Proposition~\ref{prop:c2v64-persistence}(iii)).
(iii) At each change the current region may enlarge and the maximum
may jump to the window bound, after which the capture phase of
Theorem~\ref{thm:c2v62-exterior} is needed again.  The size estimate is
the substitution of the constants.
\end{proof}

\begin{firewall}[Item (AC1) closed negatively]
\label{fw:c2v64-AC1}
The halo-level trap relative to the current non-quiet set, with or
without hysteresis, does not close on the window: the number of
restarts is finite but exceeds the window's capacity by a factor that
is a large power of the inverse threshold, and the threshold is
exponentially small in \(M^3\).  The obstruction is not the counting
of events but the recapture cost: each restart resets the region's
maximum to the window bound because the entering cells' enstrophy is
unknown.  A restart-free argument would need a bound on the enstrophy
of cells adjacent to a cell leaving the marked set, which the series
does not have.  The seed-neighbourhood enstrophy bound of
Proposition~\ref{prop:c2v62-linear} stands as the crude \(N_s\cdot2M\).
No global regularity claim is made.
\end{firewall}

% =============================================================================
\section{Integrated V64 conclusion and next acquisition order}
\label{sec:c2v64-conclusion}
% =============================================================================

\begin{theorem}[What V64 proves]
\label{thm:c2v64-conclusion}
Relative to the first cycle and V1--V63: the crude loss bound
(Lemma~\ref{lem:c2v64-loss}), the persistence and bounded-oscillation
proposition (Proposition~\ref{prop:c2v64-persistence}), and the
hysteresis count with its negative budget
(Proposition~\ref{prop:c2v64-hysteresis}).  The full Navier--Stokes
stability/global-regularity theorem remains unproved.
\end{theorem}

\begin{proof}
The cited results.
\end{proof}

\begin{programme}[V65 acquisition order]
\label{prog:c2v65}
\begin{enumerate}[label=\textup{(AC\arabic*)},leftmargin=4.8em]
\item The compulsory companion audit: read the current SM and YM
      notes and record their digests and decision.
\item Consolidate V61--V64 as the propagation chapter: finite
      propagation and the corrected counting (V61), the unconditional
      exterior trap and the linear enstrophy inequality (V62), the
      two-region proposition and the type-I structure at every time
      (V63), persistence, bounded oscillation and the hysteresis budget
      (V64); record the position after sixty-five versions.
\item Direction for V66--V70, to be fixed at V65 after the audit.  The
      candidates are: (a) the type-I gate with the new structure
      (finite configuration, finite propagation, diffuse far field at
      every time), in particular whether the moment ledger of V16--V22
      or the Rayleigh chapter of V31--V34, evaluated on an ancient
      limit whose far field is uniformly diffuse, yields a new signed
      row; (b) the intermediate class between windows, using the
      Morrey floor and the count at every time (both hold for any
      \(\mathfrak Y(t)\), with constants degenerating as
      \(\mathfrak Y\to\infty\)), to describe the spike stretch as a
      configuration whose cells are counted by \(\mathfrak Y(t)^3\); (c)
      going upstream of the cell decomposition to the Gaussian
      ledgers at moving centres, where the far-pressure bound of V58
      (a weighted sum of energies) has a Gaussian analogue.
\end{enumerate}
\end{programme}

\begin{assessment}[Position after V64]
\label{ass:c2v64-position}
Energetic cells persist for a time bounded below and oscillate
boundedly often about the threshold; the hysteresis budget for a
current-set trap fails by a large power of the inverse threshold.  No
full stability or global-regularity claim is made.
\end{assessment}

% =============================================================================
\section*{Version V64 (second cycle) append-only record}
\addcontentsline{toc}{section}{Version V64 (second cycle) append-only record}
% =============================================================================

The complete V63 source of the second cycle was retained before this
continuation.  Its SHA--256 digest was
\begin{center}\scriptsize\ttfamily
c77caf405002c6d50a11fc994b552d90927e49b8005314930796b7c6272ff64c.
\end{center}
V64 proved persistence and bounded oscillation of energetic cells and
closed the hysteresis budget negatively.  No full regularity claim is
made.

% =============================================================================
\section{V65: compulsory companion audit, the propagation chapter
consolidated, and the position after sixty-five versions}
\label{sec:c2v65-entry}
% =============================================================================

This is the sixty-fifth version of the second cycle.  The compulsory
review of the two companion programmes is performed first.  V61--V64
are then consolidated as the propagation chapter.  The position after
sixty-five versions is recorded, the three candidate directions of
Programme~\ref{prog:c2v65}(AC3) are weighed, and the direction for
V66--V70 is fixed, with one obstruction (the linear decay of Morrey
levels under refinement) recorded so that it is not rediscovered.  No
global regularity claim is made.  The next compulsory companion audit
is V70.

% =============================================================================
\section{Compulsory V65 review of the current companion notes}
\label{sec:c2v65-companion-audit}
% =============================================================================

The current companion files in \texttt{latex\_notes} were inspected
read-only.  The frozen checkpoints are
\begin{align}
 &\texttt{Renewal\_SM\_Einstein\_Continuum\_Research\_Notes\_V11.tex},
 \notag\\[-2pt]
 &\qquad
 \texttt{b6140b9259bc0e156e408e3f57fd6a597497df1185b87d4ba9e6192743a5502b},
 \label{eq:c2v65-SM-checkpoint}\\
 &\texttt{Renewal\_Yang\_Mills\_Mass\_Gap\_Research\_Notes\_V85.tex},
 \notag\\[-2pt]
 &\qquad
 \texttt{4d46962277ff998d08b99f8c6e3686cfb03e80e32c22d082524c9b6f955deab7}.
 \label{eq:c2v65-YM-checkpoint}
\end{align}
(SM V11 has 4893 lines; YM V85 has 38707 lines.)

\emph{SM V11 (seventh series).}  After reducing its local polarization
coefficient to a second spatial moment and isolating the torus-holonomy
Hessian (V10), the SM programme proved the free overlap heat-kernel
estimate with rapid off-diagonal decay (not a pointwise Gaussian
bound), uniform zeroth and second moments, and an exact periodization
with the zero-momentum holonomy term separated, suppressing the
finite-torus winding contributions to the local heat coefficients.

\emph{YM V85.}  The YM programme recorded that its termwise block-tree
target is stronger than necessary, introduced a Kruskal partition of
unity, showed that normalized percolation configurations must be
retained, introduced labeled external ports with anchor maps, gave a
tree-valued form of its V77 core and an exact remaining allocation
card, performed its own companion audit, and showed that centre-only
output trees survive its V78 skeleton while the restoration of
noncentral labels cannot be inferred.

\begin{theorem}[V65 companion-audit decision]
\label{thm:c2v65-companion-decision}
No SM V11 heat-kernel moment or periodization constant and no YM V85
partition, port or allocation bound is imported.  Neither bears on the
propagation chapter or on the gates.
\end{theorem}

\begin{proof}
The SM results concern lattice heat kernels; the YM results concern
tree and percolation representations of a lattice clock.  None is a
Navier--Stokes statement.
\end{proof}

% =============================================================================
\section{The propagation chapter, consolidated}
\label{sec:c2v65-consolidated}
% =============================================================================

\begin{theorem}[The propagation chapter, V61--V64]
\label{thm:c2v65-consolidated}
Let \(u\) be in the intermediate class with constant \(M\ge1\), \(U\) a
window rescaling on \(J_M\), \(\varepsilon_{\rm th}\) the threshold
\eqref{eq:c2v61-threshold}, \(H=1\).
\begin{enumerate}[label=\textup{(\arabic*)},leftmargin=3.2em]
\item \emph{Counting corrected (V61).}  Regular cells are counted with the
      homogeneous \(H^2\) seminorm; all but \(C_H(\nu,M)\) cells are regular.
\item \emph{Finite propagation (V61).}  With the seed set \(\mathcal S\)
      (cells energetic at level \(\varepsilon_{\rm th}/4\) at the window time,
      and non-regular cells; \(|\mathcal S|\le64N_M+C_H\)), the non-quiet
      set at time \(s\) lies within distance \(\lfloor(s+1)/\tau_c\rfloor\rho_0\) of
      \(\mathcal S\); every cell beyond \(R_{\rm inv}\) from \(\mathcal S\) is quiet
      throughout the window.
\item \emph{Exterior trap (V62).}  After the capture time
      \(\tau_{\rm cap}\le\ell_M/2\), every cell at distance \(r\ge r_1\) from
      \(\mathcal S\) has enstrophy at most \(C(\nu,M)(\varepsilon_{\rm th}+N_M(1+r-R_{\rm inv})^{-4})\),
      unconditionally; the enstrophy inequality is linear outside the
      seed neighbourhood \(B(\mathcal S,r_2)\), with an additive constant
      of order \(N_sM^3\nu^{-3}\) that prevents lengthening the window.
\item \emph{Two regions (V63).}  Envelope regions reproduce (3); current
      regions admit upward jumps and give no trap.
\item \emph{Type-I structure (V63).}  Under a type-I enstrophy rate,
      (1)--(3) hold at every time, and the ancient limit's far field is
      uniformly diffuse at every time on the corresponding sub-window.
\item \emph{Persistence and oscillation (V64).}  A regular energetic
      cell stays half-energetic for \(\tau''\); a regular cell crosses
      the threshold band at most \(\ell_M/\tau_h+1\) times; the hysteresis
      budget for a current-set trap fails by a factor of order
      \(\varepsilon_{\rm th}^{-6}\).
\end{enumerate}
\end{theorem}

\begin{proof}
(1)--(2) Remark~\ref{rem:c2v61-correction}, Theorem~\ref{thm:c2v61-propagation}.
(3) Theorem~\ref{thm:c2v62-exterior}, Corollary~\ref{cor:c2v62-exterior},
Proposition~\ref{prop:c2v62-linear}.  (4) Proposition~\ref{prop:c2v63-regions}.
(5) Theorem~\ref{thm:c2v63-typeI}, Corollary~\ref{cor:c2v63-ancient}.
(6) Propositions~\ref{prop:c2v64-persistence}, \ref{prop:c2v64-hysteresis}.
\end{proof}

% =============================================================================
\section{An obstruction recorded: Morrey levels decay linearly under refinement}
\label{sec:c2v65-obstruction}
% =============================================================================

\begin{proposition}[No coarse-to-fine tree of energetic cells at a fixed level]
\label{prop:c2v65-refinement}
Let \(Q'\) be a cell of side \(\lambda'\le\lambda\) contained in a cube \(Q\) of
side \(\lambda\), and let \(Q'\) be energetic at Morrey level \(\varepsilon\)
(\(\int_{Q'}|u|^2\ge\varepsilon\lambda'\)).  Then \(Q\) is energetic at Morrey level
\(\varepsilon\lambda'/\lambda\) and, in general, at no higher level.  Consequently a
chain of finite-propagation statements from a scale \(\lambda_0\) to the
scales \(\lambda_0 2^{-k}\), at a fixed Morrey threshold, would require
coarse thresholds \(\varepsilon2^{-k}\), whose invasion radii
\(R_{\rm inv}(\varepsilon2^{-k})\lambda_0\) diverge: the energetic sets at
successive scales cannot be nested into a tree of bounded radius by
the propagation chapter alone, and the finiteness of the singular set
of a type-I singularity does not follow from
Theorem~\ref{thm:c2v63-typeI} by chaining.
\end{proposition}

\begin{proof}
\(\int_Q|u|^2\ge\int_{Q'}|u|^2\ge\varepsilon\lambda'=(\varepsilon\lambda'/\lambda)\lambda\); a field
supported in \(Q'\) shows the level is sharp.  The rest is the
dependence of \(R_{\rm inv}\) on the threshold (a negative power).
\end{proof}

% =============================================================================
\section{Position after sixty-five versions and direction}
\label{sec:c2v65-position}
% =============================================================================

\begin{assessment}[Position after sixty-five versions of the second cycle]
\label{ass:c2v65-position-sixty-five}
The trap and propagation chapters (V56--V64) give the intermediate
class at its windows a complete qualitative description: a finite
configuration of energetic parabolic cells, finite propagation, a
uniformly weak exterior after a capture time; and they give, for any
singularity, the Morrey floor at every time.  Three negative closures
mark the chapter's limits: the trap controls maxima, not sums (V62);
a trap relative to the moving non-quiet set restarts too often (V63,
V64); the Morrey level decays under refinement, so the configuration
cannot be chained across scales (V65).  Of the three candidates of
Programme~\ref{prog:c2v65}(AC3): (a) the type-I gate with the new
structure is where the description is richest (every time is a
window), and the finite configuration at every scale, together with
the first cycle's active-scale analysis, is the natural place to look
for a count of singular points at a common active scale; (b) the
intermediate class between windows has only the degenerate constants;
(c) the Gaussian analogue of the far-pressure bound is a technical
transfer with no new mechanism.  The direction for V66--V70 is (a),
in two steps: first remove the non-regular cells from the seed set by
running the propagation and the trap jointly with a shortened capture
time (enlarging \(b\) to \(\varepsilon_{\rm th}^{-1/2}\), which makes the
pre-capture energy gain of a non-regular exterior cell negligible), so
that the seed set is the energetic set alone; second, for a type-I
singularity, relate singular points to energetic cells at their active
scales (a singular point has a nontrivial ancient limit, hence active
scales at which its parabolic ball is energetic), and count the
singular points sharing a common active scale by \(N_M\).  No global
regularity claim is made.
\end{assessment}

% =============================================================================
\section{Integrated V65 conclusion and next acquisition order}
\label{sec:c2v65-conclusion}
% =============================================================================

\begin{theorem}[What V65 establishes]
\label{thm:c2v65-conclusion}
Relative to the first cycle and V1--V64: the companion-audit decision
(Theorem~\ref{thm:c2v65-companion-decision}), the consolidated
propagation chapter (Theorem~\ref{thm:c2v65-consolidated}), and the
refinement obstruction (Proposition~\ref{prop:c2v65-refinement}).  The
full Navier--Stokes stability/global-regularity theorem remains
unproved.
\end{theorem}

\begin{proof}
The cited results.
\end{proof}

\begin{programme}[V66 acquisition order]
\label{prog:c2v66}
\begin{enumerate}[label=\textup{(AC\arabic*)},leftmargin=4.8em]
\item Energetic seeds only.  With \(b:=\max(b,\varepsilon_{\rm th}^{-1/2})\)
      (so \(\tau_{\rm cap}\le\varepsilon_{\rm th}^{1/2}\) and the trapped level is
      \(2K\hat\varepsilon\le C\varepsilon_{\rm th}^{1/2}\)), run the finite propagation and the
      exterior trap jointly by a first-time argument: let \(s^*\) be the
      first time a cell outside \(B(\mathcal S',R)\) becomes non-quiet,
      \(\mathcal S'\) the energetic seeds; before \(s^*\) the far-pressure
      bounds hold for the exterior, the trap captures within
      \(\varepsilon_{\rm th}^{1/2}\), and a non-regular exterior cell's energy gain
      (cubic term bounded by \((e^N)^{3/4}(e^N+y^N)^{3/4}\) with \(y^N\) the
      trapped enstrophy after capture and \(2M\) before) is at most
      \(C\varepsilon_{\rm th}^{5/4}M^{3/4}+C\varepsilon_{\rm th}^{3/2}\ell_M\), below
      \(\varepsilon_{\rm th}/4\); determine whether the induction of
      Theorem~\ref{thm:c2v61-propagation} then closes with
      \(\mathcal S'\) and \(R=R_{\rm inv}'\), and state the theorem.
\item Singular points at active scales.  Under a type-I enstrophy
      rate, show that at every active scale of a singular point \(x_s\)
      (a scale at which the rescaling at \(x_s\) is close to a nonzero
      element of the ancient family), the parabolic cell of \(x_s\) is
      energetic at an explicit level, and deduce that the singular
      points sharing a common active scale number at most \(N_M\) at
      that level.
\item The next compulsory companion audit is V70.
\end{enumerate}
\end{programme}

\begin{assessment}[Position after V65]
\label{ass:c2v65-position}
The propagation chapter is consolidated; Morrey levels decay linearly
under refinement, so the configuration cannot be chained across
scales.  No full stability or global-regularity claim is made.
\end{assessment}

% =============================================================================
\section*{Version V65 (second cycle) append-only record}
\addcontentsline{toc}{section}{Version V65 (second cycle) append-only record}
% =============================================================================

The complete V64 source of the second cycle was retained before this
continuation.  Its SHA--256 digest was
\begin{center}\scriptsize\ttfamily
e6cf001d868c726e34ef4c8ca359b017398d5b78f672613df69db807efca2413.
\end{center}
V65 performed the compulsory companion audit and consolidated the
propagation chapter.  No full regularity claim is made.

% =============================================================================
\section{V66: energetic seeds only, and singular points at active scales}
\label{sec:c2v66-entry}
% =============================================================================

Items (AC1) and (AC2) of Programme~\ref{prog:c2v66} are carried out.
(AC1): the finite propagation theorem holds with the energetic cells
alone as seeds.  The non-regular cells, which V61 had to include in the
seed set because the crude growth lemma needs a sup-norm bound, are
handled by a joint first-time argument: with the capture time
shortened to \(\varepsilon_{\rm th}^{1/2}\) (by enlarging the constant \(b\), which
weakens the trap inequality but keeps the trapped level of order
\(\varepsilon_{\rm th}^{1/2}\)), a non-regular cell in the trapped exterior
gains less than a quarter of the threshold over the window, so it can
become non-quiet only inside a band around the current localization
radius, and each such event (at most \(C_H\) of them) enlarges the
radius by a fixed amount.  (AC2): under a type-I rate, at an active
scale of a singular point (in the sense of the first cycle's Gaussian
activity) a parabolic cell within bounded distance of the point is
energetic at an explicit level, so the singular points active at a
common scale form at most \(N_M\) clusters of bounded diameter at that
scale, and the homogeneous singular points number at most \(N_M\).  No
global regularity claim is made.  The next compulsory companion audit
is V70.

% =============================================================================
\section{Finite propagation from the energetic seeds}
\label{sec:c2v66-seeds}
% =============================================================================

Notation of V61--V62 on a window.  Let \(\mathcal S':=\{\beta:e_\beta(-1)\ge\varepsilon_{\rm th}/4\}\),
\(|\mathcal S'|\le64N_M\), and let \(\mathcal R\) be the set of non-regular cells,
\(|\mathcal R|\le C_H\).  Replace \(b\) by
\(\tilde b:=\max(b,\varepsilon_{\rm th}^{-1/2},2/\ell_M)\), so that
\(\tau_{\rm cap}=1/\tilde b\le\varepsilon_{\rm th}^{1/2}\) and \(\tilde K=4\tilde b/a\); the
trapped level of the exterior at distance \(\rho\ge\rho_{\min}\) from the
non-quiet set is then
\(2\tilde K\hat\varepsilon(\rho)\le C(\nu)\varepsilon_{\rm th}^{-1/2}\bigl(C\varepsilon_{\rm th}+N_Ma_M(1+\rho)^{-4}+\kappa^{\rho/2}C_M\bigr)\le C_1(\nu,M)\varepsilon_{\rm th}^{1/2}\)
for \(\rho\ge\rho_{\min}:=\max\bigl(\rho_0,\ (N_Ma_M/\varepsilon_{\rm th})^{1/4},\ \rho_{\rm pol}\bigr)\),
with \(\rho_{\rm pol}\) the radius beyond which the pollution terms are
below \(\varepsilon_{\rm th}\) (of order \(\log(1/\varepsilon_{\rm th})\)).

\begin{lemma}[A non-regular cell in the trapped exterior stays quiet]
\label{lem:c2v66-nonregular}
Let \(\beta\in\mathcal R\) with \(e_\beta(-1)<\varepsilon_{\rm th}/4\), and let
\([-1,s_1]\subset J_M\) be an interval on which every cell within distance
\(2\) of \(\beta\) is at distance at least \(\rho_{\min}\) from the non-quiet
set and, after \(-1+\tau_{\rm cap}\), has enstrophy at most
\(C_1\varepsilon_{\rm th}^{1/2}\).  Then \(e_\beta<\varepsilon_{\rm th}/2\) on \([-1,s_1]\), provided
\(\varepsilon_{\rm th}\le\varepsilon_{\star\star}(\nu,M)\), an explicit smallness.
\end{lemma}

\begin{proof}
The local energy identity for \(\beta\), without any sup-norm bound:
\(e'_\beta\le C\nu e^N+C(e^N)^{3/4}(e^N+y^N)^{3/4}+C\|P^{\rm loc}\|_{3/2}\|U\|_{L^3(2Q)}+C\|\nabla P^{\rm far}\|_\infty(e^N)^{1/2}\),
with the local pressure term bounded by \(C\|U\|_{L^3(4Q_\beta)}^3\le C(e^N)^{3/4}(e^N+y^N)^{3/4}\)
and \(e^N\le125\varepsilon_{\rm th}\) as long as \(e_\beta<\varepsilon_{\rm th}\) and the
neighbours are quiet (they are at distance at least \(\rho_{\min}\ge3\)
from the non-quiet set).  Before \(-1+\tau_{\rm cap}\): \(y^N\le250M\), so
the gain is at most \(C\varepsilon_{\rm th}^{3/4}M^{3/4}\tau_{\rm cap}\le C\varepsilon_{\rm th}^{5/4}M^{3/4}\).
After: \(y^N\le125C_1\varepsilon_{\rm th}^{1/2}\), so the cubic rate is at most
\(C\varepsilon_{\rm th}^{3/4}\varepsilon_{\rm th}^{3/8}=C\varepsilon_{\rm th}^{9/8}\), and the gain over the
window at most \(C\varepsilon_{\rm th}^{9/8}\ell_M\).  The far pressure, by
\eqref{eq:c2v58-halo} with \(\rho\ge\rho_{\min}\), is at most
\(C(S_4\varepsilon_{\rm th}+N_Ma_M\rho_{\min}^{-4})\le C\varepsilon_{\rm th}\), giving a gain at most
\(C\varepsilon_{\rm th}^{3/2}\ell_M\).  The diffusion term gives \(C\nu\varepsilon_{\rm th}\ell_M\).
The total is below \(\varepsilon_{\rm th}/4\) once
\(C(\varepsilon_{\rm th}^{1/4}M^{3/4}+\varepsilon_{\rm th}^{1/8}\ell_M+\varepsilon_{\rm th}^{1/2}\ell_M+\nu\ell_M)\le\frac14\),
which defines \(\varepsilon_{\star\star}\) (the \(\nu\ell_M\) term as in
\eqref{eq:c2v61-far-small}); the energy then stays below
\(\varepsilon_{\rm th}/4+\varepsilon_{\rm th}/4\).
\end{proof}

\begin{theorem}[Finite propagation from the energetic seeds]
\label{thm:c2v66-propagation}
Let \(\varepsilon_{\rm th}\le\min(\varepsilon_\star,\varepsilon_{\star\star})\) and
\(\Delta:=\rho_{\min}+3\).  Define the localization radius
\begin{equation}
 L(s):=\lfloor(s+1)/\tau_c\rfloor\,\rho_0+n(s)\,\Delta ,
 \label{eq:c2v66-radius}
\end{equation}
where \(n(s)\le C_H\) is the number of non-regular cells that have been
non-quiet at some time before \(s\).  Then for every \(s\in J_M\),
\begin{equation}
 \mathcal N(s)\ \subset\ B(\mathcal S',L(s))\ \subset\ B(\mathcal S',R'_{\rm inv}),\qquad
 R'_{\rm inv}:=\lfloor\ell_M/\tau_c\rfloor\rho_0+C_H\Delta ,
 \label{eq:c2v66-propagation}
\end{equation}
and every cell at distance more than \(R'_{\rm inv}\) from the energetic
seeds is quiet throughout the window.  Moreover, for \(r\ge R'_{\rm inv}+\rho_{\min}\)
and \(s\ge-1+\tau_{\rm cap}\), every cell at distance at least \(r\) from
\(\mathcal S'\) has enstrophy at most \(C_1\varepsilon_{\rm th}^{1/2}\).
\end{theorem}

\begin{proof}
Let \(s^*\) be the first time at which \(\mathcal N(s)\not\subset B(\mathcal S',L(s))\)
(if none, there is nothing to prove), and \(\beta\) a cell at distance
more than \(L(s^*)\) from \(\mathcal S'\) that becomes non-quiet at \(s^*\);
\(\beta\notin\mathcal S'\), so \(e_\beta(-1)<\varepsilon_{\rm th}/4\).  On \([-1,s^*)\) the
non-quiet set is within \(B(\mathcal S',L(s))\).

\emph{The exterior trap on \([-1,s^*)\).}  The envelope regions
\(\mathcal Q(s):=\{\alpha:{\rm dist}(\alpha,\mathcal S')\ge L(s)+\rho_{\min}\}\) shrink in
time; on each interval of constancy of \(L\) the localized inequality
of Proposition~\ref{prop:c2v58-localized} holds for the weighted
maximum over \(\mathcal Q(s)\) with the constant effective energy
\(\hat\varepsilon(\rho_{\min})\) (the non-quiet set being within \(B(\mathcal S',L(s))\),
at distance at least \(\rho_{\min}\) from the region) and with \(\tilde b\);
the maximum drops at the jump times of \(L\); the capture-and-trapping
argument of Theorem~\ref{thm:c2v62-exterior} gives, for
\(s\in[-1+\tau_{\rm cap},s^*)\), enstrophy at most \(2\tilde K\hat\varepsilon(\rho_{\min})\le C_1\varepsilon_{\rm th}^{1/2}\)
on \(\mathcal Q(s)\).

\emph{Case \(\beta\) regular.}  As in the proof of
Theorem~\ref{thm:c2v61-propagation}: the near times
\(\mathcal T_{\rm near}=\{s<s^*:{\rm dist}(\beta,\mathcal N(s))<\rho_0\}\) satisfy
\({\rm dist}(\beta,\mathcal S')<L(s)+\rho_0\), hence \(L(s)>L(s^*)-\rho_0\), which
forces \(\lfloor(s+1)/\tau_c\rfloor=\lfloor(s^*+1)/\tau_c\rfloor\) and \(n(s)=n(s^*)\)
(each unit of either counter raises \(L\) by at least \(\rho_0\), as
\(\Delta\ge\rho_0\)); so \(\mathcal T_{\rm near}\subset[-1+\lfloor(s^*+1)/\tau_c\rfloor\tau_c,s^*)\),
an interval of length less than \(\tau_c\), and the gain on it is less
than \(\frac38\varepsilon_{\rm th}\);
the far gain is at most \(\frac38\varepsilon_{\rm th}\) (Lemma~\ref{lem:c2v61-far});
total less than \(\frac34\varepsilon_{\rm th}\): \(\beta\) cannot reach \(\varepsilon_{\rm th}\).

\emph{Case \(\beta\) non-regular.}  Read \(L\) as right-continuous, so
that the event ``\(\beta\) becomes non-quiet at \(s^*\)'' is counted in
\(n(s^*)\): \(L(s^*)\ge L(s^*-)+\Delta\), and \({\rm dist}(\beta,\mathcal S')>L(s^*)\ge L(s)+\rho_{\min}+3\)
for every \(s<s^*\).  Hence every cell within distance \(2\) of \(\beta\) is
at distance more than \(L(s)+\rho_{\min}\) from \(\mathcal S'\), that is, lies in
\(\mathcal Q(s)\), for every \(s<s^*\); so the hypotheses of
Lemma~\ref{lem:c2v66-nonregular} hold on \([-1,s^*)\), giving
\(e_\beta<\varepsilon_{\rm th}/2\) there and, by continuity, \(e_\beta(s^*)\le\varepsilon_{\rm th}/2<\varepsilon_{\rm th}\):
\(\beta\) is not non-quiet at \(s^*\), a contradiction.

Thus there is no first failure time, and the localization holds on the
whole window.  Since \(n\le C_H\) and \(\lfloor(s+1)/\tau_c\rfloor\le\lfloor\ell_M/\tau_c\rfloor\),
the radius is at most \(R'_{\rm inv}\).  The last statement is the
exterior trap on \(\mathcal Q(s)\supset\{{\rm dist}(\cdot,\mathcal S')\ge R'_{\rm inv}+\rho_{\min}\}\).
\end{proof}

\begin{corollary}[The configuration of a window]
\label{cor:c2v66-configuration}
At every window time the non-quiet set (at most \(N_M\) cells) lies
within \(R'_{\rm inv}(\nu,M)\) of the at most \(64N_M\) cells energetic at the
window time; all other cells are quiet throughout, and beyond a
further \(\rho_{\min}\) their enstrophy is at most \(C_1\varepsilon_{\rm th}^{1/2}\) after
the capture time \(\varepsilon_{\rm th}^{1/2}\).  Under a type-I enstrophy rate this
holds at every time.
\end{corollary}

\begin{proof}
Theorem~\ref{thm:c2v66-propagation}, Proposition~\ref{prop:c2v58-count},
and Theorem~\ref{thm:c2v63-typeI}.
\end{proof}

% =============================================================================
\section{Singular points at their active scales}
\label{sec:c2v66-singular}
% =============================================================================

Let \(u\) satisfy a type-I enstrophy rate with constant \(M\) on
\([T_*-\delta_0,T_*)\), and let \(x_s\) be a singular point.  Recall from the
first cycle (V97) that a scale \(\lambda\) is \emph{active at level \(c\)} at
\(x_s\) if the Gaussian ledger satisfies \(\Phi_{u,x_s}(\lambda)\ge c\), that is
\(\lambda^{-1}\int|u(T_*-\lambda^2)|^2\psi_\lambda(\cdot-x_s)\ge c\) with the Gaussian
weight \(\psi_\lambda(y)=e^{-|y|^2/(4\nu\lambda^2)}\); every singular point has active
scales (its ancient limits are nonzero), and the homogeneous points of
the first cycle are those with all scales active at a common level.

\begin{proposition}[Active scales are energetic cells]
\label{prop:c2v66-active}
If \(\lambda\) is active at level \(c\) at \(x_s\), then some parabolic cell of
side \(\lambda\) within lattice distance \(R_c:=\lceil(4\nu\log(2S_G/c))^{1/2}\rceil\)
of the cell of \(x_s\) has energy at least
\(c'\lambda\), \(c':=c/(2(2R_c+1)^3)\), at time \(T_*-\lambda^2\), where
\(S_G:=\sum_\alpha e^{-(|\alpha|-1)_+^2/(4\nu)}<\infty\) is the lattice Gaussian sum.
Consequently the singular points that are active at level \(c\) at a
common scale \(\lambda\) are covered by at most \(N_M(c')\) balls of radius
\((2R_c+2)\lambda\) (the cells energetic at level \(c'\) at time \(T_*-\lambda^2\),
enlarged), and any family of such points pairwise at distance more
than \(2(2R_c+2)\lambda\) has at most \(N_M(c')\) members.  In particular the
homogeneous singular points (all scales active at level \(c\)) number at
most \(N_M(c')\).
\end{proposition}

\begin{proof}
In the rescaled variables at \(x_s\) and scale \(\lambda\), \(\Phi\ge c\) reads
\(\sum_\alpha\int_{Q_\alpha}|U(-1)|^2e^{-|y|^2/4\nu}\ge c\); the cells at lattice
distance more than \(R_c\) from the origin contribute at most
\(\sum_{|\alpha|>R_c}a_Me^{-(|\alpha|-1)^2/4\nu}\), which is at most \(c/2\) for
\(R_c\) as chosen (the tail of the Gaussian lattice sum against the
uniform cell bound \(a_M\); absorb \(a_M\) into \(S_G\)); so the
\((2R_c+1)^3\) cells within distance \(R_c\) carry at least \(c/2\), and one
of them at least \(c'\).  The covering statement is
Proposition~\ref{prop:c2v58-count} at level \(c'\) (valid at every time
under the type-I rate); the counting of separated points follows since
each covering ball contains at most one of them.  For homogeneous
points, apply the covering at scales \(\lambda\to0\): two distinct
homogeneous points are separated at small scales, so their number is
at most \(N_M(c')\).
\end{proof}

\begin{assessment}[Items (AC1)--(AC2)]
\label{ass:c2v66-AC}
(AC1) closes: the seed set is the energetic set.  (AC2) closes in the
form stated: singular points active at a common scale and level form
at most \(N_M(c')\) clusters at that scale, and the homogeneous points
number at most \(N_M(c')\), an energy-based count parallel to the
first cycle's dissipation-based count of uniformly homogeneous points
(Theorem V81-count of the first cycle).  The general singular set is
not counted: distinct singular points may have disjoint sets of active
scales (the lacunary points of the first cycle), and the refinement
obstruction of V65 prevents chaining.  No global regularity claim is
made.
\end{assessment}

% =============================================================================
\section{Integrated V66 conclusion and next acquisition order}
\label{sec:c2v66-conclusion}
% =============================================================================

\begin{theorem}[What V66 proves]
\label{thm:c2v66-conclusion}
Relative to the first cycle and V1--V65: the non-regular-cell lemma
(Lemma~\ref{lem:c2v66-nonregular}), finite propagation from the
energetic seeds \eqref{eq:c2v66-propagation} with its corollary, and
the active-scale proposition (Proposition~\ref{prop:c2v66-active}).
The full Navier--Stokes stability/global-regularity theorem remains
unproved.
\end{theorem}

\begin{proof}
The cited results.
\end{proof}

\begin{programme}[V67 acquisition order]
\label{prog:c2v67}
\begin{enumerate}[label=\textup{(AC\arabic*)},leftmargin=4.8em]
\item Active scales of distinct singular points.  Under the type-I
      rate, determine whether two singular points at distance \(d\)
      apart must share an active scale of order \(d\): at the scale
      \(\lambda\sim d\) both points lie in a bounded number of cells, and
      the ancient limits of the pair (rescaled at the midpoint) are
      nonzero if either point is active there; use the first cycle's
      minimal log-period of the lacunary alternation (Theorem
      V97-log-period) to show that the active scales of a point have
      bounded gaps in log-scale, hence that every point is active at
      some scale in \([d,d\,\Theta]\) with \(\Theta=\Theta(\nu,M)\); if so, count
      the \(d\)-separated singular points by \(N_M(c')\) uniformly in \(d\),
      and conclude that the singular set of a type-I singularity is
      finite with at most \(N_M(c')\) points.
\item If (AC1) succeeds, record the consequences for the type-I gate:
      the ancient limit at each of the finitely many singular points,
      and the relation between the points' ancient limits.
\item The next compulsory companion audit is V70.
\end{enumerate}
\end{programme}

\begin{assessment}[Position after V66]
\label{ass:c2v66-position}
Energy activity propagates from the energetic seeds alone; at an
active scale a singular point sits in an energetic cell, and the
homogeneous singular points number at most \(N_M(c')\).  No full
stability or global-regularity claim is made.
\end{assessment}

% =============================================================================
\section*{Version V66 (second cycle) append-only record}
\addcontentsline{toc}{section}{Version V66 (second cycle) append-only record}
% =============================================================================

The complete V65 source of the second cycle was retained before this
continuation.  Its SHA--256 digest was
\begin{center}\scriptsize\ttfamily
06d72be06cb270a41d26574afcc6e660100cb05830e972d25ca1cf3ab3ae1037.
\end{center}
V66 proved finite propagation from the energetic seeds and the
active-scale count.  No full regularity claim is made.

% =============================================================================
\section{V67: active scales of distinct singular points, and the quiet-scale
dichotomy}
\label{sec:c2v67-entry}
% =============================================================================

Item (AC1) of Programme~\ref{prog:c2v67} is examined and closed
negatively: the first cycle's minimal log-period bounds the alternation
of active and intermediate scales from below, not the gaps between
active scales from above, and the lacunary points of the first cycle
(whose activity at intermediate scales sits at satellites at rescaled
distance tending to infinity) are exactly points whose active scales
may have unbounded gaps in log-scale; two singular points at distance
\(d\) need not share an active scale of order \(d\), and the singular set
of a type-I singularity is not counted by the active-scale argument.
What the trap chapter adds is a dichotomy at every scale of every
singular point: either an energetic cell lies within a bounded
distance of the point, or the point's parabolic neighbourhood is
trapped on that window and carries a dissipation below an explicit
fraction of the threshold, so that the dissipation of the CKN cylinder
at the point is carried by the later windows.  Item (AC2) is moot.
No global regularity claim is made.  The next compulsory companion
audit is V70.

% =============================================================================
\section{Why active scales need not be shared}
\label{sec:c2v67-gaps}
% =============================================================================

\begin{proposition}[Gaps between active scales are not bounded by the series]
\label{prop:c2v67-gaps}
Under a type-I enstrophy rate, let \(x_s\) be a singular point.
\begin{enumerate}[label=\textup{(\roman*)},leftmargin=3.2em]
\item Every scale is dissipation-active at \(x_s\) in the CKN sense:
      \(\lambda^{-1}\iint_{Q_\lambda(x_s)}|\nabla u|^2\ge\varepsilon_{\rm CKN}\) for every
      \(\lambda\le\lambda_0\) (import (E1)); every rescaling limit at \(x_s\) is a
      nonzero ancient solution.
\item Gaussian activity at \(x_s\) (\(\Phi_{u,x_s}(\lambda)\ge c\)) can fail along
      scale intervals of unbounded log-length: the first cycle's
      lacunary points have intermediate scales on which the Gaussian
      ledger at \(x_s\) tends to zero while the dissipation activity of
      (i) is carried by satellites at rescaled distance tending to
      infinity (Theorems V77-satellite and V97-log-period of the first
      cycle; the log-period bound is a lower bound on the length of
      each phase of the alternation, and no upper bound on the
      intermediate phases is proved).
\item Consequently, for two singular points \(x_s\ne x'_s\) at distance
      \(d\), the scale interval \([d,\Theta d]\) need not contain a scale
      active at both, for any \(\Theta=\Theta(\nu,M)\); the covering count of
      Proposition~\ref{prop:c2v66-active} applies only to points sharing
      a common active scale and level.
\end{enumerate}
\end{proposition}

\begin{proof}
(i) is the CKN criterion in its contrapositive form, together with the
strong convergence of rescalings under the type-I rate (first cycle,
Theorem V41-gate's construction), which passes the dissipation lower
bound to the limit.  (ii) is the content of the first cycle's satellite
theorems: at a lacunary point the intermediate scales exist by
definition, and their union may have unbounded log-length since only
the minimal period of the alternation is bounded.  (iii) follows: if
\(x_s\) is lacunary with an intermediate phase covering \([d,\Theta d]\), the
covering at any scale in that interval need not contain \(x_s\).
\end{proof}

\begin{firewall}[Item (AC1) closed negatively]
\label{fw:c2v67-AC1}
The finiteness of the singular set of a type-I singularity does not
follow from the trap and propagation chapters: the active-scale count
is scale-by-scale, the Morrey level decays under refinement
(Proposition~\ref{prop:c2v65-refinement}), and the active scales of
distinct points need not align.  A proof of finiteness would need
either an upper bound on the log-length of intermediate phases (a
statement about the ancient family, not available) or a scale-free
count.  No global regularity claim is made.
\end{firewall}

% =============================================================================
\section{The quiet-scale dichotomy}
\label{sec:c2v67-dichotomy}
% =============================================================================

Under the type-I rate, fix a singular point \(x_s\) and a scale \(\lambda\)
with \(t=T_*-\lambda^2\) in the type-I interval; let \(\mathcal S'_\lambda\) be the
energetic seeds at time \(t\) (cells of side \(\lambda\) with energy at least
\(\varepsilon_{\rm th}\lambda/4\)), \(R'_{\rm inv}\), \(\rho_{\min}\), \(\tau_{\rm cap}\), \(C_1\) the
constants of Theorem~\ref{thm:c2v66-propagation}, and
\(r_\ast:=R'_{\rm inv}+\rho_{\min}+3\).

\begin{theorem}[Quiet-scale dichotomy at a singular point]
\label{thm:c2v67-dichotomy}
For every such \(\lambda\), exactly one of the following holds.
\begin{enumerate}[label=\textup{(\alph*)},leftmargin=3.2em]
\item (Active at scale \(\lambda\).)  Some energetic seed lies within lattice
      distance \(r_\ast\) of the cell of \(x_s\): there is a cell of side
      \(\lambda\) within distance \((r_\ast+1)\lambda\) of \(x_s\) with energy at least
      \(\varepsilon_{\rm th}\lambda/4\) at time \(t\).
\item (Quiet at scale \(\lambda\).)  No such seed exists, and then every cell
      within distance \(2\lambda\) of \(x_s\) is quiet on the window
      \([t,t+\ell_M\lambda^2]\) and carries enstrophy at most
      \(C_1\varepsilon_{\rm th}^{1/2}\lambda^{-1}\) on \([t+\tau_{\rm cap}\lambda^2,t+\ell_M\lambda^2]\); in
      particular the dissipation of the parabolic neighbourhood on the
      window satisfies
      \begin{equation}
       \int_t^{t+\ell_M\lambda^2}\int_{B_{2\lambda}(x_s)}|\nabla u|^2
       \ \le\ 27\bigl(2M\tau_{\rm cap}+C_1\varepsilon_{\rm th}^{1/2}\ell_M\bigr)\lambda
       \ \le\ C_2(\nu,M)\,\varepsilon_{\rm th}^{1/2}\,\lambda ,
       \label{eq:c2v67-quiet}
      \end{equation}
      which is below \(\varepsilon_{\rm CKN}\lambda\) by the factor \(C_2\varepsilon_{\rm th}^{1/2}/\varepsilon_{\rm CKN}\);
      the CKN dissipation \(\varepsilon_{\rm CKN}\lambda\) of the cylinder
      \(Q_\lambda(x_s)\) is then carried by \(Q_\lambda(x_s)\) minus
      \(B_{2\lambda}(x_s)\times[t,t+\ell_M\lambda^2]\), that is by the later part of
      the cylinder.
\end{enumerate}
\end{theorem}

\begin{proof}
In (b) the cells within distance \(2\) of the cell of \(x_s\) are at
distance more than \(r_\ast-3=R'_{\rm inv}+\rho_{\min}\) from \(\mathcal S'_\lambda\);
Theorem~\ref{thm:c2v66-propagation} (valid at every time under the
type-I rate, Corollary~\ref{cor:c2v66-configuration}) gives quietness
throughout the window and the trapped enstrophy after the capture
time; before the capture time the enstrophy of a cell is at most the
global \(Y\le2M\) (rescaled), giving the first term of
\eqref{eq:c2v67-quiet} after the \(27\)-cell covering of \(B_{2\lambda}\) and the
rescaling of dissipation (\(\lambda\)).  The comparison with
\(\varepsilon_{\rm CKN}\) uses \(\tau_{\rm cap}\le\varepsilon_{\rm th}^{1/2}\).
\end{proof}

\begin{corollary}[Trap-active scales need not accumulate]
\label{cor:c2v67-active}
Call \(\lambda\) \emph{trap-active} at \(x_s\) if (a) holds.  If every scale
below \(\lambda_0\) is quiet at \(x_s\), then, summing \eqref{eq:c2v67-quiet}
over the windows tiling \((T_*-\lambda_0^2,T_*)\) at the scales
\(\lambda_k=\lambda_0(1-\ell_M)^{k/2}\), the dissipation of the parabolic cone
\(\bigcup_kB_{2\lambda_k}(x_s)\times[t_k,t_{k+1}]\) is at most
\(C_2\varepsilon_{\rm th}^{1/2}\sum_k\lambda_k\le2C_2\varepsilon_{\rm th}^{1/2}\lambda_0/\ell_M\), while
the CKN dissipation of each cylinder \(Q_{\lambda_k}(x_s)\) is at least
\(\varepsilon_{\rm CKN}\lambda_k\): every rescaling limit at \(x_s\) along these
scales carries its dissipation outside the cone
\(\{(y,s):|y|\le2(-s)^{1/2}\}\) up to the small amount above.  This is the
satellite configuration of the first cycle (activity at rescaled
distance tending to infinity at every scale), which the series does
not exclude (Gaussian invisibility at all small scales is admitted by
V23 and V27 of this cycle).  Hence trap-active scales need not
accumulate at \(0\), and when they do they need not have bounded
log-gaps.
\end{corollary}

\begin{proof}
The cone bound is \eqref{eq:c2v67-quiet} at each scale, using that (b)
at every \(\lambda_k\) places the seeds at distance at least \(r_\ast\lambda_k\)
from \(x_s\); the CKN bound is Proposition~\ref{prop:c2v67-gaps}(i); the
identification with the satellite configuration is the first cycle's
Theorem V77-satellite read at the cell level.
\end{proof}

\begin{assessment}[Position of the singular-point analysis]
\label{ass:c2v67-position-singular}
Under the type-I rate a singular point alternates between trap-active
scales (an energetic cell within \(r_\ast\) parabolic units) and quiet
scales (its parabolic neighbourhood trapped, with dissipation below
\(C_2\varepsilon_{\rm th}^{1/2}\) per window); trap-active scales need not
accumulate at \(0\) (the fully quiet case is the satellite
configuration), and when they do they need not have bounded log-gaps.  This is the cell-level form
of the first cycle's active/intermediate alternation, with the new
information that at a quiet scale the neighbourhood is trapped rather
than merely Gaussian-invisible.  No global regularity claim is made.
\end{assessment}

% =============================================================================
\section{Integrated V67 conclusion and next acquisition order}
\label{sec:c2v67-conclusion}
% =============================================================================

\begin{theorem}[What V67 proves]
\label{thm:c2v67-conclusion}
Relative to the first cycle and V1--V66: the gap proposition
(Proposition~\ref{prop:c2v67-gaps}) with the negative closure of
(AC1), the quiet-scale dichotomy \eqref{eq:c2v67-quiet}, and the
cone bound in the fully quiet case (Corollary~\ref{cor:c2v67-active}).
The full Navier--Stokes stability/global-regularity theorem remains
unproved.
\end{theorem}

\begin{proof}
The cited results.
\end{proof}

\begin{programme}[V68 acquisition order]
\label{prog:c2v68}
\begin{enumerate}[label=\textup{(AC\arabic*)},leftmargin=4.8em]
\item Internal audit of the trap chapter.  Re-derive, with the
      hypotheses stated in full, the five load-bearing statements:
      the far-pressure kernel bound on the torus
      (Lemma~\ref{lem:c2v57-far}), the interpolation
      (Lemma~\ref{lem:c2v57-interp}), the weighted enstrophy inequality
      (Proposition~\ref{prop:c2v57-enstrophy}, in particular the
      absorption of the neighbourhood dissipation with the geometric
      weights and the Young parameter), the capture-and-trapping ODE
      argument (Theorem~\ref{thm:c2v57-trap}), and the global
      inequality with the maximal local enstrophy
      (Proposition~\ref{prop:c2v56-global}); record any gap, and the
      exact dependence of \(\varepsilon_*\) on \((\nu,M)\).
\item Internal audit of the propagation chapter: the crude and far
      growth lemmas, the first-time induction with the two counters,
      and the non-regular-cell lemma; check the claim
      \(\tau_{\rm cap}\le\ell_M/2\) after enlarging \(b\), and the use of
      right-continuity of the localization radius.
\item The next compulsory companion audit is V70.
\end{enumerate}
\end{programme}

\begin{assessment}[Position after V67]
\label{ass:c2v67-position}
Active scales of distinct singular points need not align; at every
scale a singular point is either trap-active or has a trapped
neighbourhood with small dissipation; the fully quiet case is the
satellite configuration.  No full stability or global-regularity claim is
made.
\end{assessment}

% =============================================================================
\section*{Version V67 (second cycle) append-only record}
\addcontentsline{toc}{section}{Version V67 (second cycle) append-only record}
% =============================================================================

The complete V66 source of the second cycle was retained before this
continuation.  Its SHA--256 digest was
\begin{center}\scriptsize\ttfamily
6ca28a37bacd238b9bfa426b3f44b9eddd11fc207e254183a253cfe91266dbca.
\end{center}
(The V66 file was corrected in place after its assembly: a legacy
citation name.  The digest above is that of the corrected file.)  V67
closed the shared-active-scale question negatively and proved the
quiet-scale dichotomy.  No full regularity claim is made.

% =============================================================================
\section{V68: internal audit of the trap and propagation chapters}
\label{sec:c2v68-entry}
% =============================================================================

Programme~\ref{prog:c2v68} is carried out: the load-bearing statements
of V56--V67 are re-derived with their hypotheses in full.  The audit
finds one substantive error, in V56: the single-cell local enstrophy
inequality absorbed the neighbourhood palinstrophy into the cell's own
dissipation, which have different supports, so the corollary bounding
the growth of the maximal local enstrophy and the delay theorem built
on it are not proved as stated; the correct instrument is the weighted
maximum of V57, which gives a repaired delay of a shorter, explicit
length.  The exclusion of spread windows (V57), the Morrey floor (V59),
and the propagation chapter do not use the V56 corollary and stand.
Two justifications are corrected without change to any inequality (a
sign display in the V54 identity, and a false vanishing claim for a
pressure term in V57), the lattice adjustment for non-integer torus
sides is spelled out, and the V58 localized inequality is marked as
proved at sketch level with its structure identical to V57.  No global
regularity claim is made.  The next compulsory companion audit is V70.

% =============================================================================
\section{Audit findings}
\label{sec:c2v68-findings}
% =============================================================================

\begin{theorem}[Audit of the trap chapter]
\label{thm:c2v68-audit}
\begin{enumerate}[label=\textup{(A\arabic*)},leftmargin=3.6em]
\item \emph{V54 identity, signs.}  In the display of
      Lemma~\ref{lem:c2v54-influx} the two pressure terms have the wrong
      signs: the correct form is
      \(-\int(P-c)\Delta U\cdot\nabla\phi_\alpha+\int\partial_iP\,\partial_kU_i\partial_k\phi_\alpha\).
      All uses bound these terms in absolute value; nothing changes.
\item \emph{V57 Step 1, a false vanishing.}  The parenthesis in the
      proof of Proposition~\ref{prop:c2v57-enstrophy} asserts that
      \(\int P\,\partial_kU_i\partial_k\partial_i\phi_\alpha\) vanishes by \(\operatorname{div}U=0\); this
      holds without the factor \(P\), not with it.  The term is bounded
      instead by \(\|P^{\rm far}_\alpha-c_\alpha\|_{L^\infty(2Q_\alpha)}\|\nabla U\|_{L^1(2Q_\alpha)}\|\nabla^2\phi_\alpha\|_\infty\le C\|\nabla P^{\rm far}_\alpha\|_{L^\infty(2Q_\alpha)}(y^N)^{1/2}\)
      for the far part and by
      \(\|P^{\rm loc}_\alpha\|_{L^2}(y^N)^{1/2}\le C\|U\|_{L^4(4Q_\alpha)}^2(y^N)^{1/2}\le C(e^N+y^N)(y^N)^{1/2}\)
      for the local part, both of the forms already present in the
      list; the cell inequality \eqref{eq:c2v57-cell} and everything
      after it are unchanged.
\item \emph{V56, the single-cell absorption (substantive).}  The proof
      of Proposition~\ref{prop:c2v56-local} absorbs
      \(\|\Delta U\|_{L^2(2Q_\alpha)}^2\) into \(-\frac\nu2\int|\Delta U|^2\phi_\alpha\); the
      supports differ, and the remark about a partition of unity does
      not produce a single-cell inequality (this is exactly the
      regress that V57 resolved with geometric weights).  Consequently
      Corollary~\ref{cor:c2v56-delta} is not proved (its right side
      should contain the neighbourhood palinstrophy, which is not
      bounded pointwise in time), and Theorem~\ref{thm:c2v56-delay}
      is not proved as stated.  The repaired statement is
      Proposition~\ref{prop:c2v68-delay} below, with a shorter delay.
      Proposition~\ref{prop:c2v56-global} (the global inequality with
      the maximal local enstrophy) is correct: it uses only the global
      palinstrophy.  Theorem~\ref{thm:c2v57-no-spread},
      Theorem~\ref{thm:c2v59-floor} and the propagation chapter use
      \eqref{eq:c2v56-global} and the weighted inequalities of V57 only,
      never Corollary~\ref{cor:c2v56-delta}; they stand.  The
      consolidated Theorem~\ref{thm:c2v60-consolidated}(1) is to be read
      with the repaired delay.
\item \emph{V57, lattice on a non-integer torus.}  The cells were taken
      on the integer lattice ``after a harmless adjustment of \(\lambda_j\)''.
      Explicitly: with \(L=2\pi/\lambda_j\ge8\), use the lattice of spacing
      \(h:=L/\lfloor L\rfloor\in[1,\frac98)\) and cells of side \(h\); every
      constant of V57--V67 changes by a factor bounded by a power of
      \(\frac98\), and no statement changes.
\item \emph{V57, the hypotheses of the weighted inequalities.}
      Proposition~\ref{prop:c2v57-enstrophy} requires \(W\le\Lambda_0\) and
      \(\varepsilon\le1\); in Theorem~\ref{thm:c2v57-no-spread} both hold along
      the argument (\(\varepsilon\le\varepsilon_1\le1\) in the capture phase, \(m\le1\)
      after), and the continuity argument of Theorem~\ref{thm:c2v57-trap}
      maintains them.  On windows, where cell energies reach \(a_M\),
      the localized inequalities of V58 and V62 apply the cell
      inequality to quiet cells (energy below \(\varepsilon_{\rm th}\le1\)) and
      absorb the energetic cells' energies into the constants \(C_M\);
      consistent.
\item \emph{V58, proof status.}  The proof of
      Proposition~\ref{prop:c2v58-localized} is at sketch level: it
      lists the pollution terms and their weights without displaying
      each sum.  Its structure is that of Steps 2--3 of
      Proposition~\ref{prop:c2v57-enstrophy} with three additional
      bookkeeping items (cells at distance less than \(r/2\) from the
      seeds, weight at most \(\kappa^{r/2}\); dissipations of cells outside
      \(\mathcal Q_3\), weight at most \(\kappa^{r-5}\); interpolation for cells
      near the seeds, weight at most \(\kappa^{3r/4-1}\)), each of which is a
      finite sum bounded as stated.  The statement stands; the
      write-up is marked as a sketch.
\item \emph{V57 capture and trapping, checked.}  In the trap region
      \(W\ge K\varepsilon\): \(bW\le\frac a4W^2/\varepsilon\) needs \(W/\varepsilon\ge4b/a=K\), and
      \(b\varepsilon\le\frac a4W^2/\varepsilon\) needs \(W\ge K^{1/2}\varepsilon\), implied by
      \(W\ge K\varepsilon\) for \(K\ge1\).  \(Z=\varepsilon W^p\) is nonincreasing with
      \(p=4b'/a\): \(Z'\le2b'W^{p+1}-\frac{pa}2W^{p+1}=0\).  The a.e.\ derivative
      of \(m=\max(\varepsilon,W/K)\) on the set \(\{W=K\varepsilon\}\) coincides with one
      of the one-sided values.  Correct.
\item \emph{V56 global inequality, checked.}  From
      \(Y'=-2\nu\|\Delta U\|_2^2-2\int\partial_kU_j\partial_jU_i\partial_kU_i\), the cellwise
      Gagliardo--Nirenberg bound, H\"older over the cells with
      \(\sum_\alpha y'^3_\alpha\le(27\delta)^2\cdot8Y\) and \(\sum_\alpha h'_\alpha\le8(Y+\|\Delta U\|_2^2)\),
      and Young, one gets \eqref{eq:c2v56-global}.  Correct.
\item \emph{V61 smallness condition.}  \eqref{eq:c2v61-far-small}
      requires \(\nu\ell_M\le1/(24C_\flat')\); when it fails, every window
      statement of V61--V67 is to be read on the sub-window of length
      \(\ell:=\min(\ell_M,1/(24C_\flat'\nu))\), with \(\ell\) in place of \(\ell_M\)
      in the constants.  The type-I statements at every time are
      unaffected (windows of length \(\ell\) still cover the interval).
\item \emph{V67 in-place correction.}  The original
      Corollary~\ref{cor:c2v67-active} claimed that trap-active scales
      accumulate at \(0\) at every singular point, by comparing the
      parabolic cone's dissipation with the CKN cylinder's; the
      cylinder is larger than the cone, and the claim does not follow
      (the fully quiet case is the satellite configuration).  The
      corollary was corrected in place in the V67 file before this
      version; the digest in the record below is that of the corrected
      file.
\end{enumerate}
\end{theorem}

\begin{proof}
Each item is verified in the text of the cited statement; (A2) and
(A3) are the only ones affecting a proof, and (A3) the only one
affecting a statement.
\end{proof}

% =============================================================================
\section{The repaired delay statement}
\label{sec:c2v68-delay}
% =============================================================================

\begin{proposition}[Delay after a weak state, repaired]
\label{prop:c2v68-delay}
Let \(U\) be a smooth mean-zero solution on the torus of side \(L\ge8\),
\(W\), \(\varepsilon\) the weighted maxima of V57, \(\Lambda_0\ge1\), and suppose
\(W(s_0)\le\frac12\), \(\varepsilon(s_0)\le\frac12\), \(Y(s_0)\le Y_0\).  Then with
\(b^\sharp:=\max(b,b')\) the constants of
Propositions~\ref{prop:c2v57-enstrophy}--\ref{prop:c2v57-energy} for
this \(\Lambda_0=1\), on
\begin{equation}
 s-s_0\ \le\ T^\sharp:=\frac{\log 2}{2b^\sharp}
 \label{eq:c2v68-delay}
\end{equation}
one has \(W\le1\), \(\varepsilon\le1\), and
\(Y(s)\le Y_0\exp\bigl(C_3(\nu^{-3}+1)(s-s_0)\bigr)\).  In the rescaled
variables of a window with \(Y_0=2M\), \(T^\sharp\) is of order
\(\nu^3/(C\,\cdot)\) with an absolute constant, independent of \(M\); it is
shorter than the cone delay \(\nu^3/(32C_1M^2)\) for \(M\) below an absolute
constant and longer above it.
\end{proposition}

\begin{proof}
While \(W\le1=\Lambda_0\) and \(\varepsilon\le1\), \eqref{eq:c2v57-W} gives
\(W'\le b(W+\varepsilon)\) (dropping the negative term) and \eqref{eq:c2v57-eps}
gives \(\varepsilon'\le b'(\varepsilon+W)\); so \(m:=\max(W,\varepsilon)\) satisfies
\(m'\le2b^\sharp m\) a.e.\ and \(m(s)\le\frac12e^{2b^\sharp(s-s_0)}\le1\) for
\(s-s_0\le T^\sharp\); the hypotheses persist by continuity.  The bound on
\(Y\) is \eqref{eq:c2v56-global} with \(\delta\le W\le1\).  The order of
\(T^\sharp\): \(b,b'\) with \(\Lambda_0=1\) depend on \(\nu\) only; in the rescaled
variables \(\nu\) is the physical viscosity, and \(b\sim C\nu^{-3}\) from the
stretching Young term, so \(T^\sharp\sim\nu^3/C\).
\end{proof}

\begin{remark}[What the repaired delay says and does not say]
\label{rem:c2v68-delay}
The repaired statement bounds the time a uniformly weak state (weighted
maxima at most \(\frac12\)) needs to stop being weak by an absolute
multiple of \(\nu^3\) in rescaled time, and the global enstrophy stays
at most \(Y_0e^{C(\nu^{-3}+1)T^\sharp}\) meanwhile.  With the exclusion of
spread windows this has no application in the intermediate class (a
window is never uniformly weak); it is recorded as the correct form of
V56 and as a general statement.  The original claim that the delay is
of order \(1/M\) and exceeds the cone delay by a factor of order \(M\) is
withdrawn.
\end{remark}

% =============================================================================
\section{Integrated V68 conclusion and next acquisition order}
\label{sec:c2v68-conclusion}
% =============================================================================

\begin{theorem}[What V68 establishes]
\label{thm:c2v68-conclusion}
Relative to the first cycle and V1--V67: the audit
(Theorem~\ref{thm:c2v68-audit}) with the withdrawal of
Corollary~\ref{cor:c2v56-delta} and Theorem~\ref{thm:c2v56-delay} as
stated, and the repaired delay (Proposition~\ref{prop:c2v68-delay}).
The full Navier--Stokes stability/global-regularity theorem remains
unproved.
\end{theorem}

\begin{proof}
The cited results.
\end{proof}

\begin{programme}[V69 acquisition order]
\label{prog:c2v69}
\begin{enumerate}[label=\textup{(AC\arabic*)},leftmargin=4.8em]
\item Write out Proposition~\ref{prop:c2v58-localized} in full (the
      three bookkeeping sums), so that the propagation chapter rests
      on displayed inequalities; record the constants \(C_M\) and the
      minimal \(r\).
\item Prepare V70: the audit and the consolidation of V66--V69 with the
      position after seventy versions, and the direction for
      V71--V75, which should go upstream of the cell decomposition:
      the two remaining open fronts are the type-I gate (Liouville for
      the ancient family with the finite-configuration description)
      and the intermediate class between windows; determine at V70
      which of the first cycle's tools (the Gaussian ledgers at moving
      centres, the backward cone, the gap summability) can be combined
      with the finite configuration.
\item The next compulsory companion audit is V70.
\end{enumerate}
\end{programme}

\begin{assessment}[Position after V68]
\label{ass:c2v68-position}
The trap chapter is audited: V56's delay theorem is withdrawn and
repaired with a shorter delay; the exclusion of spread windows, the
Morrey floor and the propagation chapter stand.  No full stability or
global-regularity claim is made.
\end{assessment}

% =============================================================================
\section*{Version V68 (second cycle) append-only record}
\addcontentsline{toc}{section}{Version V68 (second cycle) append-only record}
% =============================================================================

The complete V67 source of the second cycle was retained before this
continuation.  Its SHA--256 digest was
\begin{center}\scriptsize\ttfamily
5ae2e187107b27048ca2aecc00851cd147ccdf5fc33b1f1554e7d65d087ead59.
\end{center}
(The V67 file was corrected in place after its assembly: the corollary
on trap-active scales, see (A10); the digest above is that of the
corrected file.)  V68 audited the trap and propagation chapters and
repaired the delay statement.  No full regularity claim is made.

% =============================================================================
\section{V69: the localized weighted inequalities written out in full}
\label{sec:c2v69-entry}
% =============================================================================

Item (AC1) of Programme~\ref{prog:c2v69} is carried out: the localized
weighted inequalities of V58 are proved with every sum displayed.  The
write-out finds that the sketch's lower bound on the weighted
dissipation was organised wrongly (it divided a pollution term by the
quiet energy), and replaces it by a far/near split of the weighted
maximum itself: the cells near the seeds contribute to the maximum an
amount exponentially small in the distance, and the interpolation is
applied only to the far cells, whose neighbourhoods carry the quiet
energy.  The statement of V58 is unchanged, with the minimal distance
\(r\ge12\) made explicit and the constants named.  No global regularity
claim is made.  The next compulsory companion audit is V70.

% =============================================================================
\section{Setting}
\label{sec:c2v69-setting}
% =============================================================================

Window rescaling \(U\) on \(J_M\) (or the sub-window of (A9) of V68), cells,
partition of unity, \(e_\alpha,y_\alpha,d_\alpha\) as in V57; window bounds
\(y_\alpha\le Y\le2M\), \(e_\alpha\le a_M\), \(\int_{J_M}\|\Delta U\|_2^2\le4M/\nu\).  A set
\(\mathcal A\) of at most \(N\) cells; \(\rho_\beta:={\rm dist}(\beta,\mathcal A)\);
\(\mathcal Q_r=\{\rho\ge r\}\); \(\kappa=\frac12\); for \(\alpha^*\) fixed,
\(\omega_\beta:=\kappa^{|\beta-\alpha^*|}\); \(N_\beta=\{\gamma:|\gamma-\beta|\le2\}\),
\(y^N_\beta=\sum_{N_\beta}y_\gamma\), etc.; \(\bar e_\rho:=\max_{\mathcal Q_\rho}e\);
\(W_r=\max_{\alpha\in\mathcal Q_r}\sum_{\beta\in\mathcal Q_3}\kappa^{|\beta-\alpha|}y_\beta\),
\(\varepsilon_r=\max_{\alpha\in\mathcal Q_r}\sum_{\beta\in\mathcal Q_3}\kappa^{|\beta-\alpha|}e_\beta\).
The cell inequality of V57 (Step 1 of Proposition~\ref{prop:c2v57-enstrophy}
with \(\Lambda_0=2M\), energies at most \(a_M\), Young parameter
\(\eta=\kappa^2/1000\)), for every cell \(\beta\):
\begin{equation}
 y'_\beta\ \le\ -\nu d_\beta+4\eta\nu\,d^N_\beta+b_1\bigl(y^N_\beta+f_\beta^{3/2}+f_\beta^{1/2}y^N_\beta\bigr),
 \qquad f_\beta:=\|\nabla P^{\rm far}_\beta\|_{L^\infty(2Q_\beta)},
 \label{eq:c2v69-cell}
\end{equation}
with \(b_1=b_1(\nu,M)\); by \eqref{eq:c2v58-halo},
\(f_\beta\le C_{\rm far}''(S_4\bar e_{\rho_\beta/2}+1024Na_M/\rho_\beta)\) for \(\rho_\beta\ge2\),
and \(f_\beta\le F:=C_{\rm far}''(S_4+N)a_M\) always.  Put
\begin{equation}
 f_r:=C_{\rm far}''\Bigl(S_4\bar e_{r/4}+\frac{2048Na_M}{r}\Bigr),\qquad
 \hat\varepsilon_r:=\varepsilon_r+S_4\bar e_{r/4}+\frac{C_MN}{r}+\kappa^{r/2}C_M ,
 \label{eq:c2v69-eff}
\end{equation}
so that \(f_\beta\le f_r\le C\hat\varepsilon_r\) for \(\rho_\beta\ge r/2\).

% =============================================================================
\section{The enstrophy inequality in full}
\label{sec:c2v69-enstrophy}
% =============================================================================

\begin{proposition}[V58 enstrophy inequality, full proof]
\label{prop:c2v69-W}
For \(r\ge12\) and almost every \(s\in J_M\) with \(\varepsilon_r(s)\le1\),
\begin{equation}
 W_r'\ \le\ -\,a\,\frac{W_r^2}{\hat\varepsilon_r}+b\,(W_r+\hat\varepsilon_r)+\pi_r(s),\qquad
 a:=\frac{\nu\kappa}{8\cdot1458\,S_\kappa},\quad
 \pi_r(s):=\kappa^{r/4}C_M\bigl(1+\nu\|\Delta U(s)\|_2^2\bigr),
 \label{eq:c2v69-W}
\end{equation}
with \(b=b(\nu,M)\) and \(C_M=C_M(\nu,M,N)\) explicit in the proof.
\end{proposition}

\begin{proof}
Let \(\alpha^*\in\mathcal Q_r\) attain \(W_r\) at \(s\); \(W_r'=\sum_{\beta\in\mathcal Q_3}\omega_\beta y'_\beta\).
Insert \eqref{eq:c2v69-cell}.

\emph{Dissipation.}  With \(D:=\sum_{\gamma\in\mathcal Q_3}\omega_\gamma d_\gamma\),
\[
 4\eta\nu\sum_{\beta\in\mathcal Q_3}\omega_\beta d^N_\beta
 =4\eta\nu\sum_\gamma d_\gamma\!\!\sum_{\beta\in\mathcal Q_3:\,|\beta-\gamma|\le2}\!\!\omega_\beta
 \le4\eta\nu\Bigl(125\kappa^{-2}\sum_{\gamma\in\mathcal Q_3}\omega_\gamma d_\gamma+125\kappa^{r-4}\sum_{\gamma\notin\mathcal Q_3}d_\gamma\Bigr)
 \le\frac\nu2D+\frac{\nu}{8}\kappa^{r-4}\|\Delta U\|_2^2 ,
\]
using \(\omega_\beta\le\kappa^{-2}\omega_\gamma\) for \(|\beta-\gamma|\le2\), \(\omega_\beta\le\kappa^{r-4}\)
when \(\beta\in\mathcal Q_3\) is within \(2\) of a cell outside \(\mathcal Q_3\)
(then \(\rho_\beta\le4\), so \(|\beta-\alpha^*|\ge r-4\)), and \(4\eta\cdot125\kappa^{-2}=\frac12\).
Hence the dissipation terms total at most \(-\frac\nu2D+\pi^{(1)}\),
\(\pi^{(1)}:=\frac\nu8\kappa^{r-4}\|\Delta U\|_2^2\).

\emph{Linear terms.}  Split \(\mathcal Q_3\) into \(\mathcal B_{\rm far}:=\{\rho_\beta\ge r/2\}\)
and \(\mathcal B_{\rm near}:=\mathcal Q_3\setminus\mathcal B_{\rm far}\); on \(\mathcal B_{\rm near}\),
\(|\beta-\alpha^*|\ge r/2\), \(\omega_\beta\le\kappa^{r/2}\), and
\(|\mathcal B_{\rm near}|\le N(r+1)^3\).  Then
\[
 \sum_{\beta\in\mathcal Q_3}\omega_\beta y^N_\beta
 \le125\kappa^{-2}\sum_{\gamma\in\mathcal Q_3}\omega_\gamma y_\gamma+125\kappa^{r-4}\sum_{\gamma\in\mathcal Q_1\setminus\mathcal Q_3}y_\gamma
 \le125\kappa^{-2}W_r+125^2\kappa^{r-4}N\cdot2M ,
\]
and therefore
\[
 b_1\sum_{\beta\in\mathcal Q_3}\omega_\beta\bigl(1+f_\beta^{1/2}\bigr)y^N_\beta
 \le b_1(1+f_r^{1/2})\bigl(125\kappa^{-2}W_r+125^2\kappa^{r-4}\cdot2NM\bigr)
 +b_1(1+F^{1/2})\kappa^{r/2}\cdot125\cdot N(r+1)^3\cdot2M ,
\]
the last term collecting \(\mathcal B_{\rm near}\) (where \(f_\beta\le F\) only).
Similarly
\(b_1\sum_{\beta}\omega_\beta f_\beta^{3/2}\le b_1S_\kappa f_r^{3/2}+b_1F^{3/2}\kappa^{r/2}N(r+1)^3\).
Since \((r+1)^3\kappa^{r/4}\le C\), all the \(\kappa^{r/2}\)- and \(\kappa^{r-4}\)-terms are
at most \(\kappa^{r/4}C_M\) with
\(C_M:=C\,b_1(1+F^{1/2}+F^{3/2})NM\), and, with \(f_r\le C\hat\varepsilon_r\),
\[
 W_r'\ \le\ -\frac\nu2D+b_2W_r+b_3\hat\varepsilon_r^{3/2}+\pi^{(1)}+\kappa^{r/4}C_M,
 \qquad b_2:=125\kappa^{-2}b_1(1+CF^{1/2}),\ \ b_3:=Cb_1S_\kappa .
\]

\emph{Far/near split of the maximum.}  Write
\(W_r=W^{\rm far}+W^{\rm near}\) with \(W^{\rm far}:=\sum_{\gamma\in\mathcal Q_{r/4+1}}\omega_\gamma y_\gamma\);
then \(W^{\rm near}\le\kappa^{3r/4-1}\cdot2M\cdot N(r/2+3)^3\le\kappa^{r/2}C_M\).
For \(\gamma\in\mathcal Q_{r/4+1}\) every cell within distance \(1\) of \(\gamma\)
lies in \(\mathcal Q_{r/4}\subset\mathcal Q_3\) (as \(r\ge12\)), so
Lemma~\ref{lem:c2v57-interp} gives
\(y_\gamma^2\le54\,\bar e_{r/4}\bigl(d^{(1)}_\gamma+C_\theta^2y^{(1)}_\gamma\bigr)\) with
\(d^{(1)}_\gamma=\sum_{|\beta-\gamma|\le1}d_\beta\), and
\[
 \sum_{\gamma\in\mathcal Q_{r/4+1}}\omega_\gamma d^{(1)}_\gamma\le27\kappa^{-1}D,\qquad
 \sum_{\gamma\in\mathcal Q_{r/4+1}}\omega_\gamma y^{(1)}_\gamma\le27\kappa^{-1}W_r ,
\]
all the cells involved being in \(\mathcal Q_3\).  By Cauchy--Schwarz,
\((W^{\rm far})^2\le S_\kappa\sum_{\gamma\in\mathcal Q_{r/4+1}}\omega_\gamma y_\gamma^2\le1458\kappa^{-1}S_\kappa\,\bar e_{r/4}\,(D+C_\theta^2W_r)\),
that is
\begin{equation}
 D\ \ge\ \frac{\kappa}{1458S_\kappa}\,\frac{(W^{\rm far})^2}{\bar e_{r/4}}-C_\theta^2W_r .
 \label{eq:c2v69-D}
\end{equation}

\emph{Conclusion.}  If \(W_r\ge2\kappa^{r/2}C_M\), then \(W^{\rm far}\ge W_r/2\)
and \eqref{eq:c2v69-D} gives
\(\frac\nu2D\ge\frac{\nu\kappa}{8\cdot1458S_\kappa}\frac{W_r^2}{\bar e_{r/4}}-\frac\nu2C_\theta^2W_r\ge a\frac{W_r^2}{\hat\varepsilon_r}-\frac\nu2C_\theta^2W_r\)
since \(\hat\varepsilon_r\ge S_4\bar e_{r/4}\ge\bar e_{r/4}\) (\(S_4\ge1\)); hence
\(W_r'\le-aW_r^2/\hat\varepsilon_r+(b_2+\frac\nu2C_\theta^2)W_r+b_3\hat\varepsilon_r^{3/2}+\pi^{(1)}+\kappa^{r/4}C_M\).
If \(W_r<2\kappa^{r/2}C_M\le2\hat\varepsilon_r\), then \(aW_r^2/\hat\varepsilon_r\le4a\hat\varepsilon_r\), and
the same inequality holds with \(b\) enlarged by \(4a\), the dissipation
being dropped.  With \(\hat\varepsilon_r^{3/2}\le\hat\varepsilon_r\cdot\max(1,\hat\varepsilon_r^{1/2})\)
and \(\hat\varepsilon_r\le C_M\) on the window, \(b:=b_2+\frac\nu2C_\theta^2+b_3C_M^{1/2}+4a\)
gives \eqref{eq:c2v69-W} with \(\pi_r:=\pi^{(1)}+\kappa^{r/4}C_M\le\kappa^{r/4}C_M(1+\nu\|\Delta U\|_2^2)\)
after enlarging \(C_M\).
\end{proof}

\begin{proposition}[V58 energy inequality, full proof]
\label{prop:c2v69-eps}
Under the same hypotheses,
\(\varepsilon_r'\le b'(\varepsilon_r+W_r+\hat\varepsilon_r)+\pi_r(s)\) with \(b'=b'(\nu,M)\).
\end{proposition}

\begin{proof}
Let \(\alpha^{**}\in\mathcal Q_r\) attain \(\varepsilon_r\), \(\omega_\beta:=\kappa^{|\beta-\alpha^{**}|}\).
The cell energy inequality (proof of Proposition~\ref{prop:c2v57-energy})
with the window bounds reads
\(e'_\beta\le C\nu e^N_\beta+C(e^N_\beta+y^N_\beta)^{3/2}+Cf_\beta(e^N_\beta)^{1/2}\le C_M'(e^N_\beta+y^N_\beta)+C(f_\beta+e^N_\beta)\),
using \((e^N+y^N)^{1/2}\le(125a_M+250M)^{1/2}\).  Summing with the
weights as in the previous proof:
\(\sum_{\beta\in\mathcal Q_3}\omega_\beta e^N_\beta\le125\kappa^{-2}\varepsilon_r+125^2\kappa^{r-4}Na_M\),
\(\sum_{\beta\in\mathcal Q_3}\omega_\beta y^N_\beta\le125\kappa^{-2}W_r+125^2\kappa^{r-4}\cdot2NM\)
(as \(\alpha^{**}\in\mathcal Q_r\), the weighted enstrophy at \(\alpha^{**}\) is at
most \(W_r\)), and
\(\sum_{\beta\in\mathcal Q_3}\omega_\beta f_\beta\le S_\kappa f_r+F\kappa^{r/2}N(r+1)^3\).
Collecting, with \(f_r\le C\hat\varepsilon_r\), gives the claim with
\(b':=C(C_M'+1)125\kappa^{-2}+CS_\kappa\) and the pollution absorbed in \(\pi_r\).
\end{proof}

\begin{remark}[What changed relative to the V58 sketch]
\label{rem:c2v69-change}
The sketch bounded \(W_r^2\) by the weighted interpolation over all of
\(\mathcal Q_3\), which brought the near cells' dissipation into the
comparison with a pollution term that had to be divided by the quiet
energy; the far/near split above avoids the division: the near cells
contribute at most \(\kappa^{r/2}C_M\) to the maximum, and when the maximum
is below twice that amount no lower bound on the dissipation is
needed.  The statements of V58 and their uses in V62--V67 are unchanged,
with \(\kappa^{r/4}\) in place of \(\kappa^{r/2}\) in the pollution term (which
only changes the minimal radii \(r_1,\rho_{\min}\) by a factor \(2\)) and
\(r\ge12\).
\end{remark}

% =============================================================================
\section{Integrated V69 conclusion and next acquisition order}
\label{sec:c2v69-conclusion}
% =============================================================================

\begin{theorem}[What V69 proves]
\label{thm:c2v69-conclusion}
Relative to the first cycle and V1--V68: the full proofs of the
localized weighted inequalities (Propositions~\ref{prop:c2v69-W} and
\ref{prop:c2v69-eps}) with the far/near split \eqref{eq:c2v69-D}.  The
full Navier--Stokes stability/global-regularity theorem remains
unproved.
\end{theorem}

\begin{proof}
The cited results.
\end{proof}

\begin{programme}[V70 acquisition order]
\label{prog:c2v70}
\begin{enumerate}[label=\textup{(AC\arabic*)},leftmargin=4.8em]
\item The compulsory companion audit: read the current SM and YM
      notes and record their digests and decision.
\item Consolidate V66--V69 (energetic seeds, active scales and the
      quiet-scale dichotomy, the internal audit with the repaired
      delay, the full localized inequalities) and record the position
      after seventy versions.
\item Fix the direction for V71--V75 upstream of the cell
      decomposition, choosing between the type-I gate with the
      finite-configuration description and the intermediate class
      between windows; the decision criterion is whether a first
      cycle tool (Gaussian ledgers at moving centres, backward cone,
      gap summability) combines with a trap-chapter statement to give
      a new signed row or a new exclusion.
\end{enumerate}
\end{programme}

\begin{assessment}[Position after V69]
\label{ass:c2v69-position}
The propagation chapter rests on displayed inequalities.  No full
stability or global-regularity claim is made.
\end{assessment}

% =============================================================================
\section*{Version V69 (second cycle) append-only record}
\addcontentsline{toc}{section}{Version V69 (second cycle) append-only record}
% =============================================================================

The complete V68 source of the second cycle was retained before this
continuation.  Its SHA--256 digest was
\begin{center}\scriptsize\ttfamily
d2d92e111f5c3863fc1d5bd6493c9098d6abe78b135fa2a76e7be6fda18f2068.
\end{center}
V69 wrote out the localized weighted inequalities in full.  No full
regularity claim is made.

% =============================================================================
\section{V70: compulsory companion audit, V66--V69 consolidated, and the
position after seventy versions}
\label{sec:c2v70-entry}
% =============================================================================

This is the seventieth version of the second cycle.  The compulsory
review of the two companion programmes is performed first.  V66--V69
are then consolidated.  The position after seventy versions is
recorded, and the direction for V71--V75 is fixed by the criterion of
Programme~\ref{prog:c2v70}(AC3).  No global regularity claim is made.
The next compulsory companion audit is V75.

% =============================================================================
\section{Compulsory V70 review of the current companion notes}
\label{sec:c2v70-companion-audit}
% =============================================================================

The current companion files in \texttt{latex\_notes} were inspected
read-only.  The frozen checkpoints are
\begin{align}
 &\texttt{Renewal\_SM\_Einstein\_Continuum\_Research\_Notes\_V14.tex},
 \notag\\[-2pt]
 &\qquad
 \texttt{291948294930b1efb7bd6a154b7037a05145abda0e4920a858c4bf80b912f12d},
 \label{eq:c2v70-SM-checkpoint}\\
 &\texttt{Renewal\_Yang\_Mills\_Mass\_Gap\_Research\_Notes\_V87.tex},
 \notag\\[-2pt]
 &\qquad
 \texttt{ac1148fa6b3fab623fcfb66641b2c67ccb1d9a78fdf9d2ccdcb90d5cb71a8de5}.
 \label{eq:c2v70-YM-checkpoint}
\end{align}
(SM V14 has 6123 lines; YM V87 has 39681 lines.)

\emph{SM V14 (seventh series).}  After factoring its local two-point
heat response through the linearized lattice curvature and reducing
the parity-even part to one scalar (V13), the SM programme compared
that scalar with the continuum Dirac-square coefficient: continuum and
lattice two-point densities, physical-branch consistency of the
overlap vertices, the twice-differentiated assembled integrand,
convergence of the curvature scalar, and a logarithmic integrability
of the coefficient error over the mesoscopic heat-time range, for the
zero-winding free-background vectorlike overlap modulus.

\emph{YM V87.}  The YM programme identified its exact remaining
coarsening operation (insertions indexed by support and edge set,
assembled before the type-factorial estimate), proved two independent
polarizations, exact coefficients of its normalized tree gate, a
product bound for forest coefficients, the coefficientwise collapse of
repeated supports, and defined the remaining normalized-anchor
conversion card with a fibre bound.

\begin{theorem}[V70 companion-audit decision]
\label{thm:c2v70-companion-decision}
No SM V14 curvature-coefficient or integrability constant and no YM
V87 coarsening, polarization or anchor bound is imported.  Neither
bears on the trap and propagation chapters or on the gates.
\end{theorem}

\begin{proof}
The SM results concern lattice heat-kernel coefficients; the YM
results concern cluster expansions of a lattice clock.  None is a
Navier--Stokes statement.
\end{proof}

% =============================================================================
\section{V66--V69 consolidated}
\label{sec:c2v70-consolidated}
% =============================================================================

\begin{theorem}[Energetic seeds, active scales, audit, full inequalities]
\label{thm:c2v70-consolidated}
Let \(u\) be in the intermediate class with constant \(M\ge1\), with
threshold \(\varepsilon_{\rm th}\le\min(\varepsilon_*/2,\varepsilon_\star,\varepsilon_{\star\star})\).
\begin{enumerate}[label=\textup{(\arabic*)},leftmargin=3.2em]
\item \emph{Energetic seeds (V66).}  On a window the non-quiet set lies
      within \(R'_{\rm inv}(\nu,M)\) of the at most \(64N_M\) cells energetic at
      the window time; beyond a further \(\rho_{\min}\), after the capture
      time \(\varepsilon_{\rm th}^{1/2}\), every cell has enstrophy at most
      \(C_1\varepsilon_{\rm th}^{1/2}\).  Under a type-I enstrophy rate this holds at
      every time.
\item \emph{Active scales (V66--V67).}  Under the type-I rate, at an
      active scale (Gaussian level \(c\)) of a singular point some cell
      within \(R_c\) is energetic at level \(c'\); singular points sharing
      an active scale and level form at most \(N_M(c')\) clusters at
      that scale; homogeneous points number at most \(N_M(c')\); active
      scales of distinct points need not align, trap-active scales
      need not accumulate, and at a quiet scale the parabolic
      neighbourhood carries window dissipation at most
      \(C_2\varepsilon_{\rm th}^{1/2}\lambda\).
\item \emph{Audit (V68).}  V56's single-cell inequality, its corollary
      and the delay theorem are withdrawn as stated and repaired
      (Proposition~\ref{prop:c2v68-delay}); two justifications
      corrected; the exclusion of spread windows, the Morrey floor and
      the propagation chapter stand.
\item \emph{Full inequalities (V69).}  The localized weighted
      inequalities hold for \(r\ge12\) with the far/near split
      \eqref{eq:c2v69-D} and pollution \(\kappa^{r/4}C_M(1+\nu\|\Delta U\|_2^2)\).
\end{enumerate}
\end{theorem}

\begin{proof}
(1) Theorem~\ref{thm:c2v66-propagation}, Corollary~\ref{cor:c2v66-configuration}.
(2) Proposition~\ref{prop:c2v66-active}, Proposition~\ref{prop:c2v67-gaps},
Theorem~\ref{thm:c2v67-dichotomy}, Corollary~\ref{cor:c2v67-active}.
(3) Theorem~\ref{thm:c2v68-audit}.  (4) Propositions~\ref{prop:c2v69-W},
\ref{prop:c2v69-eps}.
\end{proof}

% =============================================================================
\section{Position after seventy versions and direction}
\label{sec:c2v70-position}
% =============================================================================

\begin{assessment}[Position after seventy versions of the second cycle]
\label{ass:c2v70-position-seventy}
The second cycle's trap and propagation chapters (V56--V69) are
complete and audited.  Their content, in one sentence each: at every
time before any singularity the parabolic-scale energy is bounded
below by an explicit function of the scaled enstrophy (V57, V59); in
the intermediate class every window is a finite configuration of at
most \(N_M\) energetic parabolic cells whose activity propagates at a
finite parabolic speed and outside of which the flow is uniformly weak
after a capture time (V58, V61, V62, V66); under a type-I rate this
holds at every time and singular points are counted at their active
scales (V63, V66); the limits of the method are the maxima-not-sums
nature of the trap (V62), the non-monotone motion of the non-quiet
set (V63--V64), the linear decay of Morrey levels under refinement
(V65), and the non-alignment of active scales (V67).

The criterion of Programme~\ref{prog:c2v70}(AC3) is applied.  For the
type-I gate, the first cycle's tools evaluated on a finite
configuration reproduce known rows (the count of active centres per
scale, the satellite alternation) and no new signed row was found
in V63 or V67.  For the intermediate class between windows, the trap
chapter's statements at a general time have constants degenerating
polynomially in \(\mathfrak Y(t)\) (the count) and exponentially (the
floor); the backward cone (V46) gives the spike times a parabolic
delay after each window and summable square-root gaps; these can be
combined into a description of the state at a spike time as a finite
configuration at the spike's own scale, counted by \(\mathfrak Y^3\), which
is new in form though not yet in consequence.  One general statement
is also available cheaply: the floor at a sub-parabolic scale chosen
so that the scaled enstrophy is one, with a threshold exponential in
\(\mathfrak Y^2\) rather than \(\mathfrak Y^3\).  The direction for V71--V75 is
therefore the intermediate class at spike times: V71 the sub-parabolic
floor and count at a general time; V72 the state at a spike time
relative to the preceding window (which cells are energetic at the
spike's scale, and their distance from the window's seeds, using the
backward cone's delay and the finite propagation on the window); V73
whether the energy needed at the spike's scale can be supplied from
the window's configuration in the available time; V74 consolidation;
V75 audit.  No global regularity claim is made.
\end{assessment}

% =============================================================================
\section{Integrated V70 conclusion and next acquisition order}
\label{sec:c2v70-conclusion}
% =============================================================================

\begin{theorem}[What V70 establishes]
\label{thm:c2v70-conclusion}
Relative to the first cycle and V1--V69: the companion-audit decision
(Theorem~\ref{thm:c2v70-companion-decision}) and the consolidation
(Theorem~\ref{thm:c2v70-consolidated}).  The full Navier--Stokes
stability/global-regularity theorem remains unproved.
\end{theorem}

\begin{proof}
The cited results.
\end{proof}

\begin{programme}[V71 acquisition order]
\label{prog:c2v71}
\begin{enumerate}[label=\textup{(AC\arabic*)},leftmargin=4.8em]
\item Sub-parabolic floor.  For a general time \(t\) with
      \(\Lambda=\mathfrak Y(t)\ge1\), rescale at the scale \(\lambda'=\lambda(t)/\Lambda\) (so that
      the scaled enstrophy at scale \(\lambda'\) is \(1\)) and run the trap
      with \(\Lambda_0=1\): determine the required smallness of the
      \(\lambda'\)-cell energies for the trapping time to reach \(T_*\)
      (rescaled time \(\Lambda^2\)), and state the resulting floor
      \(\sup_x\lambda'^{-1}\int_{B_{\lambda'}(x)}|u(t)|^2\ge\exp(-C(\nu)\Lambda^2)\), comparing
      with \eqref{eq:c2v59-tradeoff}.
\item Count at a general time.  State the count of energetic
      \(\lambda(t)\)-cells and \(\lambda'\)-cells at level \(\varepsilon\) at a general time
      in terms of \(\Lambda\), and the finite propagation on the interval
      \([t,t+\ell_\Lambda\lambda(t)^2]\) with the constants of the class with
      constant \(\Lambda\) (every time is a window time for its own
      constant).
\item The next compulsory companion audit is V75.
\end{enumerate}
\end{programme}

\begin{assessment}[Position after V70]
\label{ass:c2v70-position}
The trap and propagation chapters are consolidated and audited; the
direction for V71--V75 is the intermediate class at spike times.  No
full stability or global-regularity claim is made.
\end{assessment}

% =============================================================================
\section*{Version V70 (second cycle) append-only record}
\addcontentsline{toc}{section}{Version V70 (second cycle) append-only record}
% =============================================================================

The complete V69 source of the second cycle was retained before this
continuation.  Its SHA--256 digest was
\begin{center}\scriptsize\ttfamily
6b4de333c0b7bb981381d7da657bbf7a826ee2fbf677a8219afcf55f9c77e3e5.
\end{center}
V70 performed the compulsory companion audit and consolidated
V66--V69.  No full regularity claim is made.

% =============================================================================
\section{V71: the sub-parabolic floor and the configuration at a general time}
\label{sec:c2v71-entry}
% =============================================================================

Items (AC1) and (AC2) of Programme~\ref{prog:c2v71} are carried out.
At a general time before the singularity, rescaling at the
sub-parabolic scale at which the scaled enstrophy equals one makes
the trap's constants depend on the viscosity only and the remaining
time equal to the square of the scaled enstrophy; the resulting floor
on the sub-parabolic Morrey energy is exponential in the square of the
scaled enstrophy, and it improves the parabolic floor of V59 from an
exponential in the cube to an exponential in the square.  At the
sub-parabolic scale the counts, the finite propagation and the
exterior trap hold with constants depending on the viscosity only,
and a dichotomy follows: either some sub-parabolic cell is energetic
at a viscosity-only level, or the whole flow at that scale is trapped
for a viscosity-only time.  No global regularity claim is made.  The
next compulsory companion audit is V75.

% =============================================================================
\section{The sub-parabolic rescaling}
\label{sec:c2v71-rescaling}
% =============================================================================

Let \(u\) be a smooth solution on the torus of side \(2\pi\) with first
singular time \(T_*\), \(t<T_*\) with \(\lambda:=\lambda(t)=(T_*-t)^{1/2}\le\pi/4\), and
\(\Lambda:=\max(\mathfrak Y(t),1)\), \(\mathfrak Y(t)=\lambda\|\Lambda u(t)\|_2^2\).  Put
\begin{equation}
 \lambda':=\lambda/\Lambda,\qquad
 U'(y,s):=\lambda'\,u\bigl(\lambda'y,\ t+\lambda'^2(s+1)\bigr),\quad s\in[-1,-1+\Lambda^2),
 \label{eq:c2v71-rescaling}
\end{equation}
a smooth mean-zero solution on the torus of side \(L'=2\pi\Lambda/\lambda\ge8\)
with \(Y'(-1)=\|\nabla U'(-1)\|_2^2=\lambda'\|\Lambda u(t)\|_2^2=\mathfrak Y(t)/\Lambda\le1\), the
singular time corresponding to \(s=-1+\Lambda^2\).  The cells of V57 at
scale \(\lambda'\) are the \emph{sub-parabolic cells} of the time \(t\).  Let
\(a,b,b',K,p,\Lambda_{\rm tr}:=b'(1+K)\) be the constants of
Theorem~\ref{thm:c2v57-trap} for \(\Lambda_0=1\); they depend on \(\nu\) only.

\begin{theorem}[Sub-parabolic floor]
\label{thm:c2v71-floor}
With
\(\varepsilon'_*(\nu,\Lambda):=S_\kappa^{-1}K^{-1}\exp\bigl(-(1+p)\Lambda_{\rm tr}\Lambda^2\bigr)\ge\exp(-C(\nu)\Lambda^2)\),
\begin{equation}
 \mathfrak M'(t):=\sup_x\frac1{\lambda'}\int_{B_{\lambda'}(x)}|u(t)|^2\ \ge\ c_0\,\varepsilon'_*(\nu,\Lambda)
 \qquad(t<T_*,\ \lambda(t)\le\pi/4).
 \label{eq:c2v71-floor}
\end{equation}
Consequently the parabolic floor improves to
\begin{equation}
 \mathfrak M(t)\ \ge\ \frac{c_0\,\varepsilon'_*(\nu,\Lambda)}{\Lambda}\ \ge\ \exp\bigl(-C'(\nu)\,\Lambda^2\bigr),
 \label{eq:c2v71-parabolic}
\end{equation}
an exponential in \(\mathfrak Y(t)^2\) in place of \eqref{eq:c2v59-tradeoff}'s
\(\mathfrak Y(t)^3\).
\end{theorem}

\begin{proof}
Suppose \eqref{eq:c2v71-floor} fails.  Then every sub-parabolic cell
has \(e_\alpha(-1)<c_0^{-1}c_0\varepsilon'_*=\varepsilon'_*\) (covering as in
Theorem~\ref{thm:c2v59-floor}), so \(\varepsilon(-1)\le S_\kappa\varepsilon'_*=:\varepsilon_0\), and
\(W(-1)\le Y'(-1)\le1=\Lambda_0\).  Apply Theorem~\ref{thm:c2v57-trap} on
\([-1,-1+\Lambda^2-\sigma]\): the proviso \(K\varepsilon_1e^{\Lambda_{\rm tr}\Lambda^2}\le1\) holds with
\(\varepsilon_1=(\varepsilon_0\Lambda_0^pK^{-p})^{1/(1+p)}=(\varepsilon_0K^{-p})^{1/(1+p)}\) by the choice
of \(\varepsilon'_*\): \(K^{1+p}\varepsilon_1^{1+p}e^{(1+p)\Lambda_{\rm tr}\Lambda^2}=K\varepsilon_0e^{(1+p)\Lambda_{\rm tr}\Lambda^2}=K S_\kappa\varepsilon'_*e^{(1+p)\Lambda_{\rm tr}\Lambda^2}=1\).
Hence \(W\le1\) on \([-1,-1+\Lambda^2)\) (the capture phase has \(W\le\Lambda_0=1\),
the trapped phase \(W\le K\varepsilon_1e^{\Lambda_{\rm tr}(s+1)}\le1\)), so
\(\delta\le1\) there and \eqref{eq:c2v56-global} gives
\(Y'(s)\le Y'(-1)\exp(C_3(\nu^{-3}+1)\Lambda^2)\) on \([-1,-1+\Lambda^2)\): the enstrophy
of \(u\) is bounded on \([t,T_*)\), and \(T_*\) is not singular, a
contradiction.  The exponential bound on \(\varepsilon'_*\) is immediate
since \(K,p,\Lambda_{\rm tr},S_\kappa\) depend on \(\nu\) only.  For
\eqref{eq:c2v71-parabolic}: \(\int_{B_\lambda(x)}|u|^2\ge\int_{B_{\lambda'}(x)}|u|^2\ge c_0\varepsilon'_*\lambda'=c_0\varepsilon'_*\lambda/\Lambda\)
at the maximizing \(x\).
\end{proof}

\begin{remark}[Comparison with V59]
\label{rem:c2v71-comparison}
V59 ran the trap at the parabolic scale with \(\Lambda_0=\Lambda\), paying the
polynomial growth of the trap constants in \(\Lambda_0\) (\(K\sim\Lambda_0^2\),
\(\Lambda_{\rm tr}\sim\Lambda_0^{5/2}\)) over a remaining time \(1\); V71 runs it at
the scale \(\lambda/\Lambda\) with \(\Lambda_0=1\) over a remaining time \(\Lambda^2\), paying
only the exponential in the remaining time.  Rescaling to a scale
\(\lambda/\Lambda^\theta\) with \(\theta\in(0,1)\) gives \(\Lambda_0=\Lambda^{1-\theta}\) and remaining
time \(\Lambda^{2\theta}\), hence an exponent of order
\(\Lambda^{(1-\theta)5/2+(1-\theta)1/2+2\theta}=\Lambda^{3-\theta}\) (from \((1+p)\Lambda_{\rm tr}\)
times the remaining time), minimized at \(\theta=1\): the sub-parabolic
choice is optimal within the method.
\end{remark}

% =============================================================================
\section{Counts and configuration at a general time}
\label{sec:c2v71-configuration}
% =============================================================================

\begin{proposition}[Counts at a general time]
\label{prop:c2v71-counts}
For every \(t<T_*\) with \(\lambda(t)\le\pi/4\) and every level \(\varepsilon>0\):
\begin{enumerate}[label=\textup{(\roman*)},leftmargin=3.2em]
\item the number of parabolic cells (side \(\lambda\)) with energy at least
      \(\varepsilon\lambda\) is at most \(1728C_S^6\mathfrak Y(t)^3/\varepsilon^3\);
\item the number of sub-parabolic cells (side \(\lambda'=\lambda/\Lambda\)) with
      energy at least \(\varepsilon\lambda'\) is at most \(1728C_S^6/\varepsilon^3\),
      independently of \(\mathfrak Y(t)\);
\item at level \(\varepsilon=c_0\varepsilon'_*(\nu,\Lambda)\) the count in (ii) is at least \(1\).
\end{enumerate}
\end{proposition}

\begin{proof}
Proposition~\ref{prop:c2v58-count} applied to the rescaling at scale
\(\lambda\) (with \(Y=\mathfrak Y(t)\) in place of \(2M\)) and to
\eqref{eq:c2v71-rescaling} (with \(Y'\le1\)); (iii) is
\eqref{eq:c2v71-floor}.
\end{proof}

\begin{theorem}[Configuration at the sub-parabolic scale, with viscosity-only constants]
\label{thm:c2v71-configuration}
Let \(\varepsilon_{\rm th}^{(1)}:=\varepsilon_{\rm th}(\nu,1)\) be the threshold
\eqref{eq:c2v61-threshold} of the class with constant \(M=1\), and
\(N_1,R'_{\rm inv,1},\rho_{\min,1},\tau_{{\rm cap},1},C_{1,1},\ell_1,\tau_1\) the
corresponding constants, all depending on \(\nu\) only.  At every
\(t<T_*\) with \(\lambda(t)\le\pi/4\), in the sub-parabolic cells of \(t\), on the
interval \([t,t+\ell_1\lambda'^2]\subset[t,T_*)\):
\begin{enumerate}[label=\textup{(\roman*)},leftmargin=3.2em]
\item (Configuration.)  At most \(N_1\) cells have energy at least
      \(\varepsilon^{(1)}_{\rm th}\lambda'\) at any time of the interval; those at level
      \(\varepsilon^{(1)}_{\rm th}\lambda'/4\) at time \(t\) (at most \(64N_1\)) are the
      seeds; the non-quiet set stays within \(R'_{\rm inv,1}\lambda'\) of the
      seeds; beyond a further \(\rho_{\min,1}\lambda'\) every cell is quiet and,
      after \(t+\tau_{{\rm cap},1}\lambda'^2\), has enstrophy at most
      \(C_{1,1}\varepsilon^{(1)}_{\rm th}{}^{1/2}\lambda'^{-1}\).
\item (Dichotomy.)  Either some sub-parabolic cell has energy at least
      \(\varepsilon^{(1)}_{\rm th}\lambda'\) at time \(t\), or every sub-parabolic cell is
      quiet at time \(t\) and then the trap of Theorem~\ref{thm:c2v57-trap}
      with \(\Lambda_0=1\), \(\varepsilon_0\le S_\kappa\varepsilon^{(1)}_{\rm th}\), keeps the maximal
      sub-parabolic enstrophy at most \(1\) on \([t,t+\tau_1\lambda'^2]\), where
      \(\tau_1:=\log(1/(K\varepsilon_1^{(1)}))/\Lambda_{\rm tr}>0\) depends on \(\nu\) only; in
      the second case the enstrophy of \(u\) satisfies
      \(\|\Lambda u(t')\|_2^2\le\|\Lambda u(t)\|_2^2e^{C_3(\nu^{-3}+1)\tau_1}\) on that
      interval.
\end{enumerate}
The physical length of the interval in (i) is \(\ell_1\lambda^2/\Lambda^2\), and of
the trapped interval in (ii) \(\tau_1\lambda^2/\Lambda^2\), both a fraction
\(\Lambda^{-2}\) of the remaining time.
\end{theorem}

\begin{proof}
The rescaling \eqref{eq:c2v71-rescaling} has \(Y'(-1)\le1\), so it is a
window rescaling of the class with constant \(1\) on
\([-1,-1+\ell_1]\) (Proposition~\ref{prop:c2v42-window} with \(M=1\)), and
Theorem~\ref{thm:c2v66-propagation} and
Corollary~\ref{cor:c2v66-configuration} apply with \(M=1\): this is (i).
(ii): in the second case all cell energies are below \(\varepsilon^{(1)}_{\rm th}\)
at \(s=-1\), so \(\varepsilon(-1)\le S_\kappa\varepsilon^{(1)}_{\rm th}\), and
Theorem~\ref{thm:c2v57-trap} with \(\Lambda_0=1\) gives \(W\le1\) as long as
\(K\varepsilon_1e^{\Lambda_{\rm tr}(s+1)}\le1\), that is for \(s+1\le\tau_1\); the enstrophy
bound is \eqref{eq:c2v56-global} with \(\delta\le1\), rescaled.
\end{proof}

\begin{assessment}[Items (AC1)--(AC2)]
\label{ass:c2v71-AC}
The sub-parabolic scale \(\lambda(t)/\mathfrak Y(t)\) is the natural scale of the
trap chapter at a general time: there the constants are those of the
class with constant one, and the price of a large scaled enstrophy is
paid only through the remaining time \(\mathfrak Y^2\).  The floor
\eqref{eq:c2v71-floor} is the sharpest form of the floor within the
method (Remark~\ref{rem:c2v71-comparison}).  The configuration
theorem says that the flow, seen at the sub-parabolic scale, is at
every time a finite configuration with viscosity-only constants on a
time interval that is the fraction \(\mathfrak Y^{-2}\) of the remaining
time.  For the intermediate class at a spike time \(t'\) with
\(\mathfrak Y(t')=\Lambda\) large, this describes the spike at its own scale
\(\lambda(t')/\Lambda\); relating it to the preceding window is the next item.
No global regularity claim is made.
\end{assessment}

% =============================================================================
\section{Integrated V71 conclusion and next acquisition order}
\label{sec:c2v71-conclusion}
% =============================================================================

\begin{theorem}[What V71 proves]
\label{thm:c2v71-conclusion}
Relative to the first cycle and V1--V70: the sub-parabolic floor
\eqref{eq:c2v71-floor} with the improved parabolic floor
\eqref{eq:c2v71-parabolic}, the counts at a general time
(Proposition~\ref{prop:c2v71-counts}), and the sub-parabolic
configuration theorem with its dichotomy
(Theorem~\ref{thm:c2v71-configuration}).  The full Navier--Stokes
stability/global-regularity theorem remains unproved.
\end{theorem}

\begin{proof}
The cited results.
\end{proof}

\begin{programme}[V72 acquisition order]
\label{prog:c2v72}
\begin{enumerate}[label=\textup{(AC\arabic*)},leftmargin=4.8em]
\item The spike relative to the window.  In the intermediate class,
      let \(t_j\) be a window time (\(\mathfrak Y(t_j)\le M\)) and \(t'>t_j\) a
      spike time with \(\mathfrak Y(t')=\Lambda\ge4M\); by the backward cone,
      \(t'-t_j\ge\nu^3\lambda_j^2/(32C_1M^2)\).  Determine where the
      sub-parabolic energetic cells of \(t'\) (side \(\lambda(t')/\Lambda\), energy at
      least \(\varepsilon'_*(\nu,\Lambda)\lambda(t')/\Lambda\)) can lie relative to the window's
      seeds: use the finite propagation on the window to bound the
      non-quiet set at the window's end, then the energy identity on
      the stretch from the window's end to \(t'\) with the only
      available bounds (total energy, far pressure by the total
      energy), and determine whether a lower bound on \(t'-t_j\) beyond
      the cone's follows from the need to move energy of order
      \(\varepsilon'_*\lambda(t')/\Lambda\) into a ball of radius \(\lambda(t')/\Lambda\) far from
      the window's configuration.
\item If (AC1) yields nothing, determine the reverse: whether the
      spike's energetic cells must be within a bounded number of
      parabolic cells (at scale \(\lambda_j\)) of the window's seeds, using
      that outside the seed neighbourhood the exterior is trapped up to
      the window's end and that the spike's sub-parabolic cells carry
      energy \(\varepsilon'_*\lambda(t')/\Lambda\), which is below the window's quiet
      level \(\varepsilon_{\rm th}\lambda_j\) times \(\lambda(t')/(\Lambda\lambda_j)\).
\item The next compulsory companion audit is V75.
\end{enumerate}
\end{programme}

\begin{assessment}[Position after V71]
\label{ass:c2v71-position}
At every time before a singularity the sub-parabolic Morrey energy
is at least an exponential in the square of the scaled enstrophy, and
the sub-parabolic configuration has viscosity-only constants.  No full
stability or global-regularity claim is made.
\end{assessment}

% =============================================================================
\section*{Version V71 (second cycle) append-only record}
\addcontentsline{toc}{section}{Version V71 (second cycle) append-only record}
% =============================================================================

The complete V70 source of the second cycle was retained before this
continuation.  Its SHA--256 digest was
\begin{center}\scriptsize\ttfamily
48ee68e9b6ee6ed74bc0901635a5c3a92887de859e529cbfcf4dedc552be629a.
\end{center}
V71 proved the sub-parabolic floor and the configuration at a general
time.  No full regularity claim is made.

% =============================================================================
\section{V72: the pre-spike floor, and spikes are seeded in the window's
configuration}
\label{sec:c2v72-entry}
% =============================================================================

Programme~\ref{prog:c2v72} is carried out with a change of tool.  The
energy accounting of (AC1)--(AC2) gives nothing: a spike's
sub-parabolic energetic cell needs an energy far below the quiet level
of a parabolic cell, so quiet cells can host it, and beyond the window
no bound localizes anything.  The trap, however, can be run with the
spike time as reference instead of the singular time: if at an earlier
time all cells at the backward-parabolic scale of the spike are quiet,
the trap keeps the enstrophy bounded up to the spike time, which
bounds the spike.  Contrapositively, every sufficiently high spike is
preceded, at every earlier time, by an energetic ball at the
backward-parabolic scale of that time.  In the intermediate class,
with the window time as the earlier time and the backward cone's
lower bound on the gap, the seed ball lies in a parabolic cell that is
energetic at a fixed \((\nu,M)\)-level at the window time: every high
spike after a window is seeded in one of the window's finitely many
energetic cells.  No global regularity claim is made.  The next
compulsory companion audit is V75.

% =============================================================================
\section{The pre-spike floor}
\label{sec:c2v72-floor}
% =============================================================================

\begin{theorem}[Pre-spike floor]
\label{thm:c2v72-prespike}
Let \(u\) be a smooth solution on the torus of side \(2\pi\) on
\([0,T)\), and \(t<t_2<T\) with \(t_2-t\le\pi^2/16\).  Put
\(\lambda_g:=(t_2-t)^{1/2}\), \(\Lambda_g:=\max(1,\lambda_g\|\Lambda u(t)\|_2^2)\), and
\begin{equation}
 \mathfrak M_g(t):=\sup_x\lambda_g^{-1}\int_{B_{\lambda_g}(x)}|u(t)|^2,\qquad
 \mathfrak M'_g(t):=\sup_x\frac{\Lambda_g}{\lambda_g}\int_{B_{\lambda_g/\Lambda_g}(x)}|u(t)|^2 .
 \label{eq:c2v72-quantities}
\end{equation}
If \(\mathfrak M'_g(t)<c_0\varepsilon'_*(\nu,\Lambda_g)\) (the sub-parabolic threshold of
Theorem~\ref{thm:c2v71-floor}), or if \(\mathfrak M_g(t)<c_0\varepsilon_*(\nu,\Lambda_g)\)
(the parabolic threshold of Theorem~\ref{thm:c2v57-no-spread}), then
\begin{equation}
 \lambda_g\,\|\Lambda u(t_2)\|_2^2\ \le\ \Lambda_g\,e^{C_4(\nu)\Lambda_g^2}
 \qquad\text{respectively}\qquad
 \lambda_g\,\|\Lambda u(t_2)\|_2^2\ \le\ \Lambda_g\,e^{C_4'(\nu)} ,
 \label{eq:c2v72-bound}
\end{equation}
with \(C_4,C_4'\) depending on \(\nu\) only.  Contrapositively: if
\(\lambda_g\|\Lambda u(t_2)\|_2^2>\Lambda_ge^{C_4\Lambda_g^2}\) then \(\mathfrak M'_g(t)\ge c_0\varepsilon'_*(\nu,\Lambda_g)\),
and if \(\lambda_g\|\Lambda u(t_2)\|_2^2>\Lambda_ge^{C_4'}\) then
\(\mathfrak M_g(t)\ge c_0\varepsilon_*(\nu,\Lambda_g)\): a high spike relative to the gap
is preceded by an energetic ball at the gap scale.
\end{theorem}

\begin{proof}
Rescale with reference time \(t_2\): \(U(y,s):=\lambda_gu(\lambda_gy,t_2+\lambda_g^2s)\),
\(s\in[-1,0]\), on the torus of side \(2\pi/\lambda_g\ge8\), with
\(Y(-1)=\lambda_g\|\Lambda u(t)\|_2^2\le\Lambda_g\) and \(Y(0)=\lambda_g\|\Lambda u(t_2)\|_2^2\).  The
solution is smooth on the closed interval, so all identities hold up
to \(s=0\).  Parabolic version: under \(\mathfrak M_g(t)<c_0\varepsilon_*(\nu,\Lambda_g)\)
the proof of Theorem~\ref{thm:c2v57-no-spread} gives \(\delta\le W\le\Lambda_g\)
on the capture phase \([-1,-1+\tau]\), \(\tau=2/(aK)\), and \(\delta\le1\) on
\([-1+\tau,0]\); \eqref{eq:c2v56-global} integrated gives
\(Y(0)\le\Lambda_g\exp\bigl(C_3(\nu^{-3}\Lambda_g^2+\Lambda_g^{1/2})\tau+C_3(\nu^{-3}+1)\bigr)\), and
\(\tau=1/b\le C\nu^3\Lambda_g^{-2}\) (the constant \(b\) of
Proposition~\ref{prop:c2v57-enstrophy} is at least \(C\nu^{-3}\Lambda_g^2\) from
the stretching Young term; enlarging \(b\) only weakens the inequality,
so \(b\ge\nu^{-3}\Lambda_g^2\) may be assumed), so the first exponent is at
most an absolute multiple of \(1+\Lambda_g^{-3/2}\nu^3\le1+\nu^3\); this is the
second bound in \eqref{eq:c2v72-bound}.  Sub-parabolic version: rescale
at \(\lambda_g/\Lambda_g\) with reference \(t_2\) (remaining rescaled time
\(\Lambda_g^2\), \(Y\le1\) at the start), and the proof of
Theorem~\ref{thm:c2v71-floor} gives \(\delta\le1\) up to \(s=0\) and
\(Y\le\exp(C_3(\nu^{-3}+1)\Lambda_g^2)\) in that rescaling, which in the
\(\lambda_g\)-rescaling reads \(Y(0)\le\Lambda_g\exp(C_3(\nu^{-3}+1)\Lambda_g^2)\).
\end{proof}

\begin{remark}[The reference time is arbitrary]
\label{rem:c2v72-reference}
Theorems~\ref{thm:c2v57-no-spread}, \ref{thm:c2v59-floor} and
\ref{thm:c2v71-floor} are the case \(t_2=T_*\), where the conclusion
(bounded enstrophy up to \(t_2\)) is contradictory; for a regular
\(t_2\) the conclusion is a bound on the enstrophy at \(t_2\), which is
the content of \eqref{eq:c2v72-bound}.  Nothing in the trap uses that
\(t_2\) is singular.
\end{remark}

% =============================================================================
\section{Spikes after a window are seeded in the window's configuration}
\label{sec:c2v72-seeded}
% =============================================================================

Let \(u\) be in the intermediate class with constant \(M\), \(t_j\) a
window time, \(\lambda_j=\lambda(t_j)\), and \(t'\in(t_j,T_*)\) a \emph{spike time of
height \(\Lambda\)}: \(\mathfrak Y(t')=\lambda(t')\|\Lambda u(t')\|_2^2=\Lambda\).  By
Proposition~\ref{prop:c2v46-cone} (the backward cone), if \(\Lambda\ge4M\) then
\(t'-t_j\ge c_g^2\lambda_j^2\) with \(c_g:=\bigl(\nu^3/(32C_1M^2)\bigr)^{1/2}\) (\(C_1\) the
first-cycle cone constant).  Put \(\lambda_g:=(t'-t_j)^{1/2}\in[c_g\lambda_j,\lambda_j)\).

\begin{theorem}[Spikes are seeded]
\label{thm:c2v72-seeded}
Let \(\Lambda_{\rm seed}:=\max\bigl(4M,\ Me^{C_4'(\nu)}/c_g\bigr)\).  Every spike time
\(t'>t_j\) of height \(\Lambda\ge\Lambda_{\rm seed}\) satisfies:
\begin{enumerate}[label=\textup{(\roman*)},leftmargin=3.2em]
\item (Seed ball.)  There is \(x_{\rm seed}\) with
      \(\int_{B_{\lambda_g}(x_{\rm seed})}|u(t_j)|^2\ge c_0\varepsilon_*(\nu,M)\lambda_g\): an energetic
      ball at the gap scale at the window time.
\item (Seed cell.)  The parabolic cell (side \(\lambda_j\)) containing
      \(x_{\rm seed}\), enlarged to its \(27\)-cell neighbourhood, has energy
      at least \(\varepsilon_2\lambda_j\), \(\varepsilon_2:=c_0c_g\varepsilon_*(\nu,M)\), at the window
      time: it is one of the at most \(N_M(\varepsilon_2)=1728C_S^6(2M)^3/\varepsilon_2^3\)
      cells energetic at the level \(\varepsilon_2\).
\item (Sub-parabolic seed.)  With \(\Lambda_g=\max(1,\lambda_g\|\Lambda u(t_j)\|_2^2)\le M\),
      if moreover \(c_g\Lambda>\Lambda_ge^{C_4(\nu)\Lambda_g^2}\), there is a ball of
      radius \(\lambda_g/\Lambda_g\ge\lambda_g/M\) at time \(t_j\) with energy at least
      \(c_0\varepsilon'_*(\nu,\Lambda_g)\lambda_g/\Lambda_g\).
\end{enumerate}
Consequently every spike of height at least \(\Lambda_{\rm seed}\) after the
window \(t_j\) is seeded, at the window time, inside the finite
configuration of the window at level \(\varepsilon_2\): a family of at most
\(N_M(\varepsilon_2)\) parabolic cells, fixed at the window time, contains a
seed ball for every later high spike.
\end{theorem}

\begin{proof}
(i) Apply Theorem~\ref{thm:c2v72-prespike} with \(t=t_j\), \(t_2=t'\):
\(\Lambda_g\le\lambda_g\|\Lambda u(t_j)\|_2^2\le\lambda_j\|\Lambda u(t_j)\|_2^2\le M\) (or \(1\le M\)), and
\(\lambda_g\|\Lambda u(t')\|_2^2=(\lambda_g/\lambda(t'))\Lambda\ge c_g\Lambda\) since
\(\lambda(t')^2=\lambda_j^2-\lambda_g^2\le(1-c_g^2)\lambda_j^2\) gives
\(\lambda_g/\lambda(t')\ge c_g/(1-c_g^2)^{1/2}\ge c_g\); so
\(\lambda_g\|\Lambda u(t')\|_2^2\ge c_g\Lambda\ge Me^{C_4'}\ge\Lambda_ge^{C_4'}\) for
\(\Lambda\ge\Lambda_{\rm seed}\), and the contrapositive of the parabolic bound
gives \(\mathfrak M_g(t_j)\ge c_0\varepsilon_*(\nu,\Lambda_g)\ge c_0\varepsilon_*(\nu,M)\) (monotonicity of
\(\varepsilon_*\) in its second argument).  (ii) \(B_{\lambda_g}(x_{\rm seed})\) has radius
at most \(\lambda_j\), so it lies in the \(27\)-cell neighbourhood of the cell
of \(x_{\rm seed}\), whose energy is then at least
\(c_0\varepsilon_*\lambda_g\ge c_0c_g\varepsilon_*\lambda_j\); the count is
Proposition~\ref{prop:c2v58-count} at level \(\varepsilon_2\) for the
\(27\)-cell neighbourhoods (a cell whose neighbourhood carries \(\varepsilon_2\lambda_j\)
has a neighbour carrying \(\varepsilon_2\lambda_j/27\); the count at level
\(\varepsilon_2/27\) is \(27^3N_M(\varepsilon_2)\), absorbed by redefining \(\varepsilon_2\)).
(iii) is the sub-parabolic contrapositive.
\end{proof}

\begin{firewall}[What is not localized]
\label{fw:c2v72-location}
The seed is localized at the window time; the spike itself, at time
\(t'\), is not: between the window's end and \(t'\) no bound of the series
controls where the enstrophy grows, and the energy required by the
spike's sub-parabolic cells (\(c_0\varepsilon'_*(\Lambda)\lambda(t')/\Lambda\)) is far below
the quiet level \(\varepsilon_{\rm th}\lambda_j\) of a parabolic cell, so the energy
accounting of Programme~\ref{prog:c2v72}(AC1)--(AC2) gives no
constraint.  The seed ball's future is tracked by the finite
propagation only to the window's end; after that the seed may move
by any amount compatible with the total energy.  No global regularity
claim is made.
\end{firewall}

% =============================================================================
\section{Integrated V72 conclusion and next acquisition order}
\label{sec:c2v72-conclusion}
% =============================================================================

\begin{theorem}[What V72 proves]
\label{thm:c2v72-conclusion}
Relative to the first cycle and V1--V71: the pre-spike floor
(Theorem~\ref{thm:c2v72-prespike}) and the seeding of spikes in the
window's configuration (Theorem~\ref{thm:c2v72-seeded}).  The full
Navier--Stokes stability/global-regularity theorem remains unproved.
\end{theorem}

\begin{proof}
The cited results.
\end{proof}

\begin{programme}[V73 acquisition order]
\label{prog:c2v73}
\begin{enumerate}[label=\textup{(AC\arabic*)},leftmargin=4.8em]
\item The backward cone of seeds.  Apply the pre-spike floor with
      \(t_2=t'\) fixed and \(t\) varying in \((t_j,t')\): for every \(t\) with
      \((t'-t)^{1/2}\|\Lambda u(t')\|_2^2>\Lambda_g(t)e^{C_4'}\), a ball of radius
      \((t'-t)^{1/2}\) at time \(t\) carries energy at least
      \(c_0\varepsilon_*(\nu,\Lambda_g(t))(t'-t)^{1/2}\); determine the range of \(t\)
      (down to \(t'-t\) of order \(\|\Lambda u(t')\|_2^{-4}\)) and the dependence
      of the level on \(\Lambda_g(t)=(t'-t)^{1/2}\|\Lambda u(t)\|_2^2\), which is
      unbounded in general between windows; state the resulting
      backward cone of seeds and its relation to the first cycle's
      backward Leray cone (Theorem V37-cone-share).
\item Energy supply along the cone: the seed balls at successive
      scales \((t'-t)^{1/2}\) are nested in space only if their centres
      converge; determine whether the local energy identity between
      two seed times, with the far pressure bounded by the total
      energy (the only bound available between windows), constrains
      the distance between successive seed centres, and record the
      result.
\item The next compulsory companion audit is V75.
\end{enumerate}
\end{programme}

\begin{assessment}[Position after V72]
\label{ass:c2v72-position}
Every high spike after a window is preceded, at the window time, by
an energetic ball at the gap scale inside the window's finite
configuration; the spike's own location is not localized.  No full
stability or global-regularity claim is made.
\end{assessment}

% =============================================================================
\section*{Version V72 (second cycle) append-only record}
\addcontentsline{toc}{section}{Version V72 (second cycle) append-only record}
% =============================================================================

The complete V71 source of the second cycle was retained before this
continuation.  Its SHA--256 digest was
\begin{center}\scriptsize\ttfamily
91b7b07cacd6910c0d89ea400ceef62505cbe724dffcc1209a730c2152ea3d38.
\end{center}
V72 proved the pre-spike floor and the seeding of spikes.  No full
regularity claim is made.

% =============================================================================
\section{V73: the backward cone of seeds}
\label{sec:c2v73-entry}
% =============================================================================

Items (AC1) and (AC2) of Programme~\ref{prog:c2v73} are carried out.
(AC1): for a fixed spike time, the pre-spike floor applied at every
earlier time whose enstrophy is a fixed factor below the spike's
yields an energetic ball at the backward-parabolic scale of that time,
at a level depending on the scaled enstrophy of that time relative to
the spike; the set of such seed times and scales is the backward cone
of seeds, which ranges from the spike's own scale to the largest gap
over which the spike dominates, and which is the energy counterpart
of the first cycle's backward Leray cone.  (AC2): consecutive seeds
whose times lie within one window of each other (at the earlier gap
scale) are localized, the later seed lying within the invasion radius
of the earlier time's configuration; but the configuration itself is
not propagated across the refinement of scales (the Morrey level of a
configuration cell decays by a fixed factor per step), so the seeds
may hop among the configuration's cells and the cone need not
converge to a point.  No global regularity claim is made.  The next
compulsory companion audit is V75.

% =============================================================================
\section{The cone of seeds}
\label{sec:c2v73-cone}
% =============================================================================

Let \(u\) be a smooth solution on the torus of side \(2\pi\) on \([0,T)\),
\(t'<T\) a fixed time with \(Y_2:=\|\Lambda u(t')\|_2^2\), and for \(t<t'\) with
\(t'-t\le\pi^2/16\): \(\lambda_g(t):=(t'-t)^{1/2}\), \(\Lambda_g(t):=\max(1,\lambda_g(t)\|\Lambda u(t)\|_2^2)\).

\begin{definition}[Seed times]
\label{def:c2v73-seed-times}
\(t<t'\) is a \emph{seed time of \(t'\)} if
\(\lambda_g(t)Y_2>\Lambda_g(t)e^{C_4'(\nu)}\), that is, if either
\(\lambda_g(t)\|\Lambda u(t)\|_2^2\ge1\) and \(\|\Lambda u(t)\|_2^2<e^{-C_4'}Y_2\), or
\(\lambda_g(t)\|\Lambda u(t)\|_2^2<1\) and \(t'-t>e^{2C_4'}Y_2^{-2}\).  The set of seed
times is \(\mathcal T(t')\).
\end{definition}

\begin{theorem}[Backward cone of seeds]
\label{thm:c2v73-cone}
For every seed time \(t\in\mathcal T(t')\) there is \(x(t)\) with
\begin{equation}
 \int_{B_{\lambda_g(t)}(x(t))}|u(t)|^2\ \ge\ c_0\,\varepsilon_*\bigl(\nu,\Lambda_g(t)\bigr)\,\lambda_g(t),
 \label{eq:c2v73-seed}
\end{equation}
and, if moreover \(\lambda_g(t)Y_2>\Lambda_g(t)e^{C_4(\nu)\Lambda_g(t)^2}\), a ball of radius
\(\lambda_g(t)/\Lambda_g(t)\) at time \(t\) with energy at least
\(c_0\varepsilon'_*(\nu,\Lambda_g(t))\lambda_g(t)/\Lambda_g(t)\).  The seed times satisfy:
\begin{enumerate}[label=\textup{(\roman*)},leftmargin=3.2em]
\item (Smallest scale.)  \(t'-t>e^{2C_4'}Y_2^{-2}\) for every seed time:
      the cone starts at the spike's own scale \(e^{C_4'}/Y_2\).
\item (Backward Leray cone.)  For \(t'-t\ge\nu^3/(4C_1Y_2^2)\) the first
      cycle's cone gives \(\|\Lambda u(t)\|_2^2\ge(\nu^3/(8C_1(t'-t)))^{1/2}\), hence
      \(\Lambda_g(t)\ge\min(1,(\nu^3/8C_1)^{1/2})\): the level in
      \eqref{eq:c2v73-seed} is at most \(c_0\varepsilon_*(\nu,\max(1,(\nu^3/8C_1)^{1/2}))\),
      and it is at least \(c_0\varepsilon_*(\nu,\bar\Lambda)\) on any set of seed times
      where \(\Lambda_g\le\bar\Lambda\).
\item (Type-I approach to the spike.)  If \(\Lambda_g(t)\le\bar\Lambda\) on an
      interval \((t_1,t')\) (the spike is approached in a type-I manner
      relative to itself), then every \(t\in(t_1,t')\) with
      \(t'-t>(\bar\Lambda e^{C_4'}/Y_2)^2\) is a seed time with level at least
      \(c_0\varepsilon_*(\nu,\bar\Lambda)\), and the seed scales fill
      \((\bar\Lambda e^{C_4'}/Y_2,\ (t'-t_1)^{1/2})\).
\end{enumerate}
\end{theorem}

\begin{proof}
\eqref{eq:c2v73-seed} and the sub-parabolic version are the
contrapositives in Theorem~\ref{thm:c2v72-prespike}.  (i) Since
\(\Lambda_g\ge1\), \(\lambda_gY_2>e^{C_4'}\) is necessary.  (ii) is
Proposition~\ref{prop:c2v46-cone}'s input (the first cycle's cone) and
the monotonicity of \(\varepsilon_*\).  (iii) With \(\Lambda_g(t)\le\bar\Lambda\), the seed
condition \(\lambda_g(t)Y_2>\Lambda_g(t)e^{C_4'}\) follows from
\(\lambda_g(t)Y_2>\bar\Lambda e^{C_4'}\).
\end{proof}

\begin{remark}[Two cones]
\label{rem:c2v73-two-cones}
The first cycle's backward Leray cone bounds the enstrophy from below
at all times within \(\nu^3/(4C_1Y_2^2)\)-parabolic reach of a spike; the
cone of seeds bounds the energy concentration from below, at the
backward-parabolic scale, at all seed times.  The Leray cone is a
consequence of the cubic enstrophy inequality alone; the seed cone is
a consequence of the trap, and it says where (in scale, not yet in
space) the spike's energy sits beforehand.
\end{remark}

% =============================================================================
\section{One-step localization and the hopping of seeds}
\label{sec:c2v73-hopping}
% =============================================================================

\begin{proposition}[Consecutive seeds within one window]
\label{prop:c2v73-one-step}
Let \(t_1<t_2\) be seed times of \(t'\) with \(\Lambda_g(t_1),\Lambda_g(t_2)\le\bar\Lambda\) and
\(t_2-t_1\le\ell\,\lambda_g(t_1)^2\), \(\ell:=\ell_{\bar\Lambda}\) the window length of the
class with constant \(\bar\Lambda\) (or the sub-window of (A9) of V68).  Let
the parabolic cells at scale \(\lambda_g(t_1)\), the threshold
\(\varepsilon''':=\varepsilon_{\rm th}(\nu,\bar\Lambda)\) lowered if necessary to at most
\(c_0\varepsilon_*(\nu,\bar\Lambda)(1-\ell)^{1/2}/27\), and the constants \(R'_{\rm inv}=R'_{\rm inv}(\nu,\bar\Lambda;\varepsilon''')\)
be those of Theorem~\ref{thm:c2v66-propagation} for the class with
constant \(\bar\Lambda\) at that threshold.  Then the seed ball of \(t_2\) lies
within distance \((R'_{\rm inv}+2)\lambda_g(t_1)\) of a cell that is energetic at
level \(\varepsilon'''/4\) at time \(t_1\) (at scale \(\lambda_g(t_1)\)).
\end{proposition}

\begin{proof}
The rescaling at time \(t_1\) and scale \(\lambda_g(t_1)\) has scaled
enstrophy \(\Lambda_g(t_1)\le\bar\Lambda\), so it is a window rescaling of the class
with constant \(\bar\Lambda\) on \([t_1,t_1+\ell\lambda_g(t_1)^2]\ni t_2\).  The seed ball
\(B_{\lambda_g(t_2)}(x(t_2))\) at time \(t_2\) has radius at most \(\lambda_g(t_1)\) and
energy at least \(c_0\varepsilon_*(\nu,\bar\Lambda)\lambda_g(t_2)\ge c_0\varepsilon_*(\nu,\bar\Lambda)(1-\ell)^{1/2}\lambda_g(t_1)\)
(as \(\lambda_g(t_2)^2=\lambda_g(t_1)^2-(t_2-t_1)\ge(1-\ell)\lambda_g(t_1)^2\)), so one of the
\(27\) cells of side \(\lambda_g(t_1)\) meeting it has energy at least
\(\varepsilon'''\lambda_g(t_1)\) at time \(t_2\): it is non-quiet at \(t_2\), hence within
\(R'_{\rm inv}\lambda_g(t_1)\) of the seeds of \(t_1\) by
Theorem~\ref{thm:c2v66-propagation}.
\end{proof}

\begin{firewall}[Item (AC2): the seeds hop, the cone need not converge]
\label{fw:c2v73-hopping}
Iterating Proposition~\ref{prop:c2v73-one-step} along a geometric
sequence of seed times \(t_{k+1}=t_k+\frac\ell2\lambda_g(t_k)^2\) (which stays in
\(\mathcal T(t')\) under the type-I approach hypothesis of
Theorem~\ref{thm:c2v73-cone}(iii)) localizes the seed of \(t_{k+1}\)
within \((R'_{\rm inv}+2)\lambda_g(t_k)\) of \emph{some} cell energetic at level
\(\varepsilon'''/4\) at \(t_k\), of which there are up to \(N_{\bar\Lambda}(\varepsilon'''/4)\);
it does not localize it near the seed of \(t_k\).  To follow the whole
configuration from \(t_k\) to \(t_{k+1}\) one would need every cell
energetic at \(t_{k+1}\) at scale \(\lambda_g(t_{k+1})\) to be non-quiet at scale
\(\lambda_g(t_k)\), which fails by the refinement decay of
Proposition~\ref{prop:c2v65-refinement}: the level drops by the factor
\(27\cdot4/(1-\ell/2)^{1/2}\) per step, and the number of steps from the
first gap to the spike's scale is of order \(\ell^{-1}\log(\lambda_g(t_1)Y_2)\),
unbounded.  Hence the seeds may hop among the cells of successive
configurations, and the cone of seeds is not shown to converge to a
point; the spike's location at \(t'\) remains unlocalized
(Firewall~\ref{fw:c2v72-location}).  The energy identity between seed
times, with the far pressure bounded only by the total energy, adds
nothing (Firewall~\ref{fw:c2v61-beyond}).  No global regularity claim
is made.
\end{firewall}

\begin{assessment}[What the cone of seeds adds]
\label{ass:c2v73-adds}
For any spike of any solution, at every seed time the energy is
concentrated at the backward-parabolic scale at an explicit level;
under a type-I approach relative to the spike the seed scales fill
the whole range from the spike's own scale to the first gap.  In the
intermediate class this holds for the spikes after each window with
the window time as a seed time (V72), and for all later seed times
with levels depending on the enstrophy relative to the spike.  The
first cycle's picture of a spike (Theorem V37-cone-share: a backward
Leray cone with summable square-root gaps) is thereby completed on
the energy side; the spatial convergence of the cone is the open
item, and its obstruction is the refinement decay.
\end{assessment}

% =============================================================================
\section{Integrated V73 conclusion and next acquisition order}
\label{sec:c2v73-conclusion}
% =============================================================================

\begin{theorem}[What V73 proves]
\label{thm:c2v73-conclusion}
Relative to the first cycle and V1--V72: the backward cone of seeds
(Theorem~\ref{thm:c2v73-cone}) and the one-step localization
(Proposition~\ref{prop:c2v73-one-step}), with the hopping firewall.
The full Navier--Stokes stability/global-regularity theorem remains
unproved.
\end{theorem}

\begin{proof}
The cited results.
\end{proof}

\begin{programme}[V74 acquisition order]
\label{prog:c2v74}
\begin{enumerate}[label=\textup{(AC\arabic*)},leftmargin=4.8em]
\item Consolidate V71--V73 as the spike chapter (sub-parabolic floor
      and configuration, pre-spike floor and seeding, cone of seeds),
      and state the position of the intermediate class after it: the
      window is a finite configuration, every high spike is seeded in
      it, the spike's seeds exist at every backward-parabolic scale,
      and the spike's location is open.
\item Determine whether the refinement decay can be avoided by
      choosing, at each step, the threshold as a fixed fraction of the
      floor level (which is scale-free) rather than as a fraction of
      the previous step's threshold: the seed at \(t_{k+1}\) has level
      \(c_0\varepsilon_*\) at its own scale by the floor, so the localization
      of Proposition~\ref{prop:c2v73-one-step} applies at every step
      with the same threshold; record precisely why this does not give
      convergence (the hopping among up to \(N\) cells) and what
      additional statement would (that the configuration at \(t_k\) has
      a single cell, or that the seed ball is unique up to bounded
      distance).
\item The next compulsory companion audit is V75.
\end{enumerate}
\end{programme}

\begin{assessment}[Position after V73]
\label{ass:c2v73-position}
Every spike is preceded by energetic balls at all backward-parabolic
seed scales; consecutive seeds within a window are localized to the
earlier configuration, but the seeds may hop and the cone need not
converge.  No full stability or global-regularity claim is made.
\end{assessment}

% =============================================================================
\section*{Version V73 (second cycle) append-only record}
\addcontentsline{toc}{section}{Version V73 (second cycle) append-only record}
% =============================================================================

The complete V72 source of the second cycle was retained before this
continuation.  Its SHA--256 digest was
\begin{center}\scriptsize\ttfamily
be07fdbe80f0fb12672f7a176da6ccd6ab536ca765a6f6a95e45919d29646f18.
\end{center}
V73 proved the backward cone of seeds and the one-step localization.
No full regularity claim is made.

% =============================================================================
\section{V74: the spike chapter consolidated, and what convergence of the
cone would need}
\label{sec:c2v74-entry}
% =============================================================================

Items (AC1) and (AC2) of Programme~\ref{prog:c2v74} are carried out.
V71--V73 are consolidated as the spike chapter, and the position of
the intermediate class after it is stated.  The proposal of (AC2), to
use at each step of the cone the floor level rather than the previous
step's threshold, is examined: it is exactly what
Proposition~\ref{prop:c2v73-one-step} already does, and it localizes
each seed to the previous time's configuration at a fixed level; what
it cannot do is propagate the configuration, whose cells are counted
at a level that must decrease by a fixed factor per step, so the
seeds may hop among up to \(N\) cells.  The two statements that would
give convergence are recorded: that the configuration at every seed
time is a single cluster, or that the seed of each time is unique up
to a bounded distance.  No global regularity claim is made.  The next
compulsory companion audit is V75.

% =============================================================================
\section{The spike chapter, consolidated}
\label{sec:c2v74-consolidated}
% =============================================================================

\begin{theorem}[The spike chapter, V71--V73]
\label{thm:c2v74-consolidated}
Let \(u\) be a smooth solution on the torus of side \(2\pi\) with first
singular time \(T_*\).
\begin{enumerate}[label=\textup{(\arabic*)},leftmargin=3.2em]
\item \emph{Sub-parabolic floor (V71).}  For every \(t<T_*\) with
      \(\lambda(t)\le\pi/4\) and \(\Lambda=\max(\mathfrak Y(t),1)\), the Morrey energy at
      the scale \(\lambda(t)/\Lambda\) is at least \(c_0\varepsilon'_*(\nu,\Lambda)\ge e^{-C(\nu)\Lambda^2}\),
      and the parabolic Morrey energy at least \(e^{-C'(\nu)\Lambda^2}\); at the
      sub-parabolic scale the counts, finite propagation and exterior
      trap hold with viscosity-only constants on a time interval that
      is the fraction \(\Lambda^{-2}\) of the remaining time.
\item \emph{Pre-spike floor (V72).}  For any \(t<t_2<T_*\) with
      \(\lambda_g=(t_2-t)^{1/2}\le\pi/4\) and \(\Lambda_g=\max(1,\lambda_g\|\Lambda u(t)\|_2^2)\): if the
      parabolic (sub-parabolic) Morrey energy at scale \(\lambda_g\)
      (\(\lambda_g/\Lambda_g\)) at time \(t\) is below \(c_0\varepsilon_*(\nu,\Lambda_g)\)
      (\(c_0\varepsilon'_*(\nu,\Lambda_g)\)), then
      \(\lambda_g\|\Lambda u(t_2)\|_2^2\le\Lambda_ge^{C_4'}\) (\(\le\Lambda_ge^{C_4\Lambda_g^2}\)).
\item \emph{Seeding (V72).}  In the intermediate class, every spike of
      height at least \(\Lambda_{\rm seed}(\nu,M)\) after a window time \(t_j\) is
      seeded at \(t_j\) by a ball at the gap scale lying in one of at
      most \(N_M(\varepsilon_2)\) parabolic cells energetic at level
      \(\varepsilon_2=c_0c_g\varepsilon_*(\nu,M)\).
\item \emph{Cone of seeds (V73).}  For any \(t'\) and every seed time
      \(t\in\mathcal T(t')\), a ball of radius \((t'-t)^{1/2}\) at time \(t\) carries
      energy at least \(c_0\varepsilon_*(\nu,\Lambda_g(t))(t'-t)^{1/2}\); the seed scales
      range from \(e^{C_4'}/\|\Lambda u(t')\|_2^2\) upward, and under a type-I
      approach relative to \(t'\) they fill the range up to the first
      gap; consecutive seeds within one window are within
      \((R'_{\rm inv}+2)\) gap-lengths of the earlier configuration.
\end{enumerate}
\end{theorem}

\begin{proof}
(1) Theorem~\ref{thm:c2v71-floor}, Proposition~\ref{prop:c2v71-counts},
Theorem~\ref{thm:c2v71-configuration}.  (2) Theorem~\ref{thm:c2v72-prespike}.
(3) Theorem~\ref{thm:c2v72-seeded}.  (4) Theorem~\ref{thm:c2v73-cone},
Proposition~\ref{prop:c2v73-one-step}.
\end{proof}

\begin{assessment}[The intermediate class after the spike chapter]
\label{ass:c2v74-intermediate}
At a window time the flow is a finite configuration of at most \(N_M\)
energetic parabolic cells (V57--V58, V66); on the window its activity
propagates at finite speed from the energetic seeds, and outside the
seed neighbourhood it is uniformly weak after a capture time (V61,
V62, V66); every high spike after the window is seeded at the window
time inside the configuration at the gap scale (V72), and at every
seed time before the spike a ball at the backward-parabolic scale is
energetic (V73); at the spike time the flow at its sub-parabolic
scale is again a finite configuration with viscosity-only constants
(V71).  Open: the location of the spike relative to the window's
configuration, equivalently the convergence of the cone of seeds; the
evolution between the window's end and the spike, where the enstrophy
is unbounded and no local bound holds.  No global regularity claim
is made.
\end{assessment}

% =============================================================================
\section{What convergence of the cone would need}
\label{sec:c2v74-convergence}
% =============================================================================

\begin{proposition}[Fixed-level localization does not propagate the configuration]
\label{prop:c2v74-fixed-level}
Let \(t_k\uparrow t'\) be seed times with \(\Lambda_g(t_k)\le\bar\Lambda\) and
\(t_{k+1}-t_k=\frac\ell2\lambda_g(t_k)^2\).  With the fixed threshold \(\varepsilon'''\) of
Proposition~\ref{prop:c2v73-one-step}:
\begin{enumerate}[label=\textup{(\roman*)},leftmargin=3.2em]
\item for every \(k\), the seed ball of \(t_{k+1}\) is within
      \((R'_{\rm inv}+2)\lambda_g(t_k)\) of the set \(\mathcal C_k\) of cells (scale
      \(\lambda_g(t_k)\)) energetic at level \(\varepsilon'''/4\) at \(t_k\), and
      \(|\mathcal C_k|\le N_{\bar\Lambda}(\varepsilon'''/4)\);
\item the set \(\mathcal C_{k+1}\) (scale \(\lambda_g(t_{k+1})\), same level) is
      contained in the \((R'_{\rm inv}(\varepsilon^{(k)})+2)\lambda_g(t_k)\)-neighbourhood of
      the cells energetic at \(t_k\) at the level \(\varepsilon^{(k)}/4\) with
      \(\varepsilon^{(k)}:=\varepsilon'''(1-\ell/2)^{1/2}/108\), and not, in general, in a
      neighbourhood of \(\mathcal C_k\) itself;
\item iterating (ii) requires levels decreasing by the factor
      \((1-\ell/2)^{1/2}/108\) per step, so no fixed-level chain contains
      all the configurations, and the sequence of seed balls need not
      converge.
\end{enumerate}
\end{proposition}

\begin{proof}
(i) is Proposition~\ref{prop:c2v73-one-step} with the count of
Proposition~\ref{prop:c2v58-count}.  (ii) A cell of scale \(\lambda_g(t_{k+1})\)
energetic at level \(\varepsilon'''/4\) lies in a cube of side \(2\lambda_g(t_k)\)
covered by \(27\) cells of scale \(\lambda_g(t_k)\), one of which then has
energy at least \((\varepsilon'''/4)\lambda_g(t_{k+1})/27\ge\varepsilon^{(k)}\lambda_g(t_k)\); this cell
is non-quiet at \(t_{k+1}\) for the threshold \(\varepsilon^{(k)}\), and
Theorem~\ref{thm:c2v66-propagation} at that threshold places it near
the cells energetic at level \(\varepsilon^{(k)}/4\) at \(t_k\), a superset of
\(\mathcal C_k\).  (iii) is the iteration.
\end{proof}

\begin{theorem}[Two sufficient conditions for convergence]
\label{thm:c2v74-sufficient}
Under the hypotheses of Proposition~\ref{prop:c2v74-fixed-level},
either of the following implies that the seed balls \(B_{\lambda_g(t_k)}(x(t_k))\)
converge to a point \(x_\infty\) with
\(|x(t_k)-x_\infty|\le C(R'_{\rm inv}+2)\lambda_g(t_k)/\ell\) for all \(k\):
\begin{enumerate}[label=\textup{(\alph*)},leftmargin=3.2em]
\item (Single cluster.)  For every \(k\), \(\mathcal C_k\) has diameter at most
      \(D\lambda_g(t_k)\) for a fixed \(D\).
\item (Unique seed.)  For every \(k\), any two balls of radius
      \(\lambda_g(t_k)\) with energy at least \(c_0\varepsilon_*(\nu,\bar\Lambda)(1-\ell)^{1/2}\lambda_g(t_k)/27\)
      at time \(t_k\) have centres within \(D\lambda_g(t_k)\) of each other.
\end{enumerate}
In case (b) the spike's sub-parabolic energetic cells at \(t'\), which
are seeds of \(t'\) at the smallest scale, lie within
\(C(R'_{\rm inv}+2)e^{C_4'}/(\ell\,\|\Lambda u(t')\|_2^2)\) of \(x_\infty\): the spike is
localized.
\end{theorem}

\begin{proof}
(a): by Proposition~\ref{prop:c2v74-fixed-level}(i) the seed of \(t_{k+1}\)
is within \((R'_{\rm inv}+2+D)\lambda_g(t_k)\) of the seed of \(t_k\) (both are
within the neighbourhood of the single cluster), and
\(\sum_k\lambda_g(t_k)\le\lambda_g(t_1)\sum_k(1-\ell/2)^{k/2}\le4\lambda_g(t_1)/\ell\), a Cauchy
sequence.  (b): the seed of \(t_{k+1}\), of radius \(\lambda_g(t_{k+1})\le\lambda_g(t_k)\),
carries energy at least \(c_0\varepsilon_*(1-\ell)^{1/2}\lambda_g(t_k)\), so the ball of
radius \(\lambda_g(t_k)\) with the same centre carries at least that energy
and, by (b), its centre is within \(D\lambda_g(t_k)\) of \(x(t_k)\); sum as
before.  The last statement: at \(t'\) the sub-parabolic cells of the
spike are the seeds at the smallest scale \(e^{C_4'}/\|\Lambda u(t')\|_2^2\)
(Theorem~\ref{thm:c2v73-cone}(i)), and the chain ends there.
\end{proof}

\begin{firewall}[Neither condition is available]
\label{fw:c2v74-conditions}
Condition (a) fails in general: the configuration may have several
clusters (Proposition~\ref{prop:c2v59-clusters}).  Condition (b) is a
uniqueness of the concentrating ball up to bounded distance, which
the floor does not give (the floor is a supremum).  Both are
statements about the multiplicity of energy concentration at the
backward-parabolic scale; the series has the count \(N\) as an upper
bound and \(1\) as a lower bound, and nothing between.  No global
regularity claim is made.
\end{firewall}

% =============================================================================
\section{Integrated V74 conclusion and next acquisition order}
\label{sec:c2v74-conclusion}
% =============================================================================

\begin{theorem}[What V74 establishes]
\label{thm:c2v74-conclusion}
Relative to the first cycle and V1--V73: the consolidated spike
chapter (Theorem~\ref{thm:c2v74-consolidated}), the fixed-level
proposition (Proposition~\ref{prop:c2v74-fixed-level}), and the two
sufficient conditions for convergence (Theorem~\ref{thm:c2v74-sufficient}).
The full Navier--Stokes stability/global-regularity theorem remains
unproved.
\end{theorem}

\begin{proof}
The cited results.
\end{proof}

\begin{programme}[V75 acquisition order]
\label{prog:c2v75}
\begin{enumerate}[label=\textup{(AC\arabic*)},leftmargin=4.8em]
\item The compulsory companion audit: read the current SM and YM
      notes and record their digests and decision.
\item Record the position after seventy-five versions and fix the
      direction for V76--V80.  The candidate is the multiplicity
      question of Firewall~\ref{fw:c2v74-conditions}: whether the number
      of \(D\)-separated energetic balls at the backward-parabolic scale
      of a spike can be bounded by something better than \(N\), or
      whether multiple clusters at every seed scale are compatible
      with the energy inequality; and, upstream, whether the first
      cycle's satellite theorems (activity at rescaled distance
      tending to infinity at intermediate scales) already describe the
      multi-cluster case, so that the cone of seeds converges at
      homogeneous points.
\end{enumerate}
\end{programme}

\begin{assessment}[Position after V74]
\label{ass:c2v74-position}
The spike chapter is consolidated; the cone of seeds converges under
a single-cluster or a unique-seed condition, neither of which is
available.  No full stability or global-regularity claim is made.
\end{assessment}

% =============================================================================
\section*{Version V74 (second cycle) append-only record}
\addcontentsline{toc}{section}{Version V74 (second cycle) append-only record}
% =============================================================================

The complete V73 source of the second cycle was retained before this
continuation.  Its SHA--256 digest was
\begin{center}\scriptsize\ttfamily
1bc6091b85fa0b78b557f1f9b21676506188a53507c75da6742cde93c9e1040c.
\end{center}
V74 consolidated the spike chapter and recorded the conditions for
convergence of the cone of seeds.  No full regularity claim is made.

% =============================================================================
\section{V75: compulsory companion audit and the position after seventy-five
versions}
\label{sec:c2v75-entry}
% =============================================================================

This is the seventy-fifth version of the second cycle.  The compulsory
review of the two companion programmes is performed first.  The
position after seventy-five versions is recorded, the multiplicity
question of V74 is examined against the first cycle's satellite
theorems, and the direction for V76--V80 is fixed.  No global
regularity claim is made.  The next compulsory companion audit is V80.

% =============================================================================
\section{Compulsory V75 review of the current companion notes}
\label{sec:c2v75-companion-audit}
% =============================================================================

The current companion files in \texttt{latex\_notes} were inspected
read-only.  The frozen checkpoints are
\begin{align}
 &\texttt{Renewal\_SM\_Einstein\_Continuum\_Research\_Notes\_V17.tex},
 \notag\\[-2pt]
 &\qquad
 \texttt{7775ce7aa7fe9273bfac0762eb2c4a467c396af635bccccc430b8c256ff02c51},
 \label{eq:c2v75-SM-checkpoint}\\
 &\texttt{Renewal\_Yang\_Mills\_Mass\_Gap\_Research\_Notes\_V90.tex},
 \notag\\[-2pt]
 &\qquad
 \texttt{6091681f973aa01b43b85f62f94e74201f6ba33d3e4bd320fba6efadd6e687a7}.
 \label{eq:c2v75-YM-checkpoint}
\end{align}
(SM V17 has 7372 lines; YM V90 has 40965 lines; a YM file
\(\ldots\)\texttt{V89.tex} of 40534 lines was also present.)

\emph{SM V17 (seventh series).}  After closing the deterministic
smooth-link transfer of its vectorlike gauge coefficient (V16), the SM
programme computed the continuum coefficients of the Dirac--Yukawa
\(A_4\) owners in its fixed heat convention (the Dirac--Yukawa square,
the integrated \(A_4\) formula, equivariant Higgs--Yukawa invariants,
a relative Yukawa subtraction) and isolated the still open
overlap--Yukawa transfer gate (lattice and chiral-projector transfer).

\emph{YM V89--V90.}  The YM programme proved an exact hierarchical
Cayley identity, showed that the complete full-label tree sum is
insufficient, identified the compatibility defect as an exact cycle
excess, closed one defect with a fresh background tree (the one-defect
exchange and the exact fresh-coalescence card), performed its own
companion audit, showed that freshness is not forced by connectedness,
recorded a proxy-level no-go for a universal allocation card, and
stated a corrected hybrid target with local and transversal root
occurrences.

\begin{theorem}[V75 companion-audit decision]
\label{thm:c2v75-companion-decision}
No SM V17 Yukawa coefficient or subtraction constant and no YM V90
exchange, coalescence or hybrid-target bound is imported.  Neither
bears on the spike chapter or on the gates.
\end{theorem}

\begin{proof}
The SM results concern local Standard Model tensors of lattice
operators; the YM results concern tree allocations of a lattice
clock.  None is a Navier--Stokes statement.  The YM programme's
recording of a no-go for a universal card and its replacement by a
hybrid target is structurally parallel to this cycle's replacement of
the chained localization by fixed-level seeds, but nothing
transfers.
\end{proof}

% =============================================================================
\section{The multiplicity question against the satellite theorems}
\label{sec:c2v75-multiplicity}
% =============================================================================

\begin{proposition}[Homogeneous and satellite spikes]
\label{prop:c2v75-dichotomy}
Let \(t'\) be a spike time approached in a type-I manner relative to
itself (\(\Lambda_g\le\bar\Lambda\) on \((t_1,t')\)), with the geometric seed times
\(t_k\) of Proposition~\ref{prop:c2v74-fixed-level} and seed balls
\(B_{\lambda_g(t_k)}(x(t_k))\).  Exactly one of the following holds.
\begin{enumerate}[label=\textup{(\alph*)},leftmargin=3.2em]
\item (Homogeneous spike.)  There are a point \(x_\infty\) and \(D\) with
      \(|x(t_k)-x_\infty|\le D\lambda_g(t_k)\) for all large \(k\): the cone of seeds
      converges, and the spike's sub-parabolic energetic cells at \(t'\)
      lie within \(D'\lambda_g(t_K)\) of \(x_\infty\) at the smallest seed scale.
\item (Satellite spike.)  \(\limsup_k|x(t_{k+1})-x(t_k)|/\lambda_g(t_k)=\infty\)
      along some choice of seed balls, or the seeds' distance to every
      fixed point exceeds \(D\lambda_g(t_k)\) for every \(D\) along a
      subsequence: the energy concentration at the backward-parabolic
      scale moves away from any fixed centre at rescaled distance
      tending to infinity, which is the first cycle's satellite
      configuration (Theorems V77--V79 of the first cycle) read at a
      spike instead of the singular time.
\end{enumerate}
In case (a) the seeds' distances are summable and the spike is
localized as in Theorem~\ref{thm:c2v74-sufficient}; in case (b) the
one-step localization (Proposition~\ref{prop:c2v73-one-step}) shows
that the seeds hop among the cells of successive configurations, each
hop bounded by \((R'_{\rm inv}+2)\lambda_g(t_k)\) to some configuration cell but
not to the previous seed, so the configurations at successive seed
times must contain cells at mutual distance exceeding any fixed
multiple of the scale: multiple clusters at infinitely many seed
scales.
\end{proposition}

\begin{proof}
The dichotomy is a tautology on the sequence \(x(t_k)/\lambda_g(t_k)\)-relative
positions; the consequences are Theorem~\ref{thm:c2v74-sufficient}(a)
in case (a) and, in case (b), the observation that if every
configuration \(\mathcal C_k\) had diameter at most \(D\lambda_g(t_k)\) then
case (a) would hold.
\end{proof}

\begin{assessment}[Position after seventy-five versions of the second cycle]
\label{ass:c2v75-position-seventy-five}
The spike chapter reproduces at a spike, with the trap as the engine,
the first cycle's active/satellite alternation at a singular point:
a spike is either homogeneous (its seeds converge to a point, near
which its sub-parabolic energetic cells lie) or a satellite spike
(its seeds hop among clusters at rescaled distance tending to
infinity).  The intermediate class is therefore described, at the
level of energy concentration, as: finite configurations at window
times; every high spike seeded in them; every spike homogeneous or
satellite.  What the series does not have is a mechanism that
excludes either kind or relates the spike's clusters to the window's
clusters across the refinement of scales.  Two upstream facts frame
V76--V80: the first cycle's satellite theorems were proved for the
singular time with the type-I rate on the whole interval, and their
proofs (compactness of the ancient family, moving-centre ledgers)
are available at a spike under the type-I approach relative to the
spike; and the multiplicity of energy concentration (the number of
\(D\)-separated energetic balls at a scale) is bounded by \(N\) only,
with no lower-order structure.  The direction for V76--V80 is the
multiplicity: V76 transfers the first cycle's moving-centre Gaussian
ledgers to a spike under the type-I approach and asks whether the
satellite alternative has the first cycle's minimal log-period at a
spike; V77 asks whether the energy inequality on the window at the
seed scale bounds the number of clusters that can each carry energy
of order \(\varepsilon_*\lambda\) after a capture time (the exterior trap says
the clusters are the only places where enstrophy is not small, and
the Leray floor of the window dissipation must be carried by them);
V78--V79 consolidate; V80 audit.  No global regularity claim is made.
\end{assessment}

% =============================================================================
\section{Integrated V75 conclusion and next acquisition order}
\label{sec:c2v75-conclusion}
% =============================================================================

\begin{theorem}[What V75 establishes]
\label{thm:c2v75-conclusion}
Relative to the first cycle and V1--V74: the companion-audit decision
(Theorem~\ref{thm:c2v75-companion-decision}) and the homogeneous/satellite
dichotomy at a spike (Proposition~\ref{prop:c2v75-dichotomy}).  The
full Navier--Stokes stability/global-regularity theorem remains
unproved.
\end{theorem}

\begin{proof}
The cited results.
\end{proof}

\begin{programme}[V76 acquisition order]
\label{prog:c2v76}
\begin{enumerate}[label=\textup{(AC\arabic*)},leftmargin=4.8em]
\item Gaussian ledgers at a spike.  Under the type-I approach relative
      to \(t'\) (\(\Lambda_g\le\bar\Lambda\) on \((t_1,t')\)), define the Gaussian
      ledger \(\Phi^{(t')}_{x}(\lambda):=\lambda^{-1}\int|u(t'-\lambda^2)|^2\psi_\lambda(\cdot-x)\) with
      reference time \(t'\); show that the first cycle's window
      compactness (V42 of this cycle, with \(t'\) in place of \(T_*\))
      and the ledger bounds hold on the windows \([t'-\lambda^2,t'-(1-\ell)\lambda^2]\);
      determine whether the seed of scale \(\lambda\) gives
      \(\Phi^{(t')}_{x(t)}(\lambda)\ge c_0\varepsilon_*e^{-1/4\nu}\) and whether the
      first cycle's minimal log-period of the active/intermediate
      alternation (Theorem V97-log-period) transfers to the
      homogeneous/satellite alternation at the spike.
\item Multiplicity on a window.  On a window at the seed scale, after
      the capture time, the exterior of the configuration is trapped
      (enstrophy at most \(C_1\varepsilon_{\rm th}^{1/2}\) per cell); the window's
      dissipation is at least the Leray floor \(c_L^2\nu^{3/2}\ell\) when the
      reference is the singular time, or at least the backward-cone
      floor when the reference is a spike; determine whether these
      floors, carried by the configuration's clusters, bound the
      number of clusters from below or above in terms of the per-cluster
      dissipation, and record the result.
\item The next compulsory companion audit is V80.
\end{enumerate}
\end{programme}

\begin{assessment}[Position after V75]
\label{ass:c2v75-position}
A spike is homogeneous or a satellite spike; the intermediate class is
described at the level of energy concentration, and the multiplicity
of concentration is the open quantity.  No full stability or
global-regularity claim is made.
\end{assessment}

% =============================================================================
\section*{Version V75 (second cycle) append-only record}
\addcontentsline{toc}{section}{Version V75 (second cycle) append-only record}
% =============================================================================

The complete V74 source of the second cycle was retained before this
continuation.  Its SHA--256 digest was
\begin{center}\scriptsize\ttfamily
27141cb24272c67bccc20ea6527d9cbf77b7364beef92a8741d7136edcce7f76.
\end{center}
V75 performed the compulsory companion audit and recorded the
homogeneous/satellite dichotomy at a spike.  No full regularity claim
is made.

% =============================================================================
\section{V76: Gaussian ledgers at a spike, seed persistence, and the first
spike after a window}
\label{sec:c2v76-entry}
% =============================================================================

Items (AC1) and (AC2) of Programme~\ref{prog:c2v76} are carried out.
(AC1) closes: with the spike time as reference, under a type-I
approach relative to the spike, the window compactness and the
Gaussian ledger bounds of V42 hold verbatim, a seed of scale \(\lambda\)
makes the ledger at its centre at least an explicit constant, and the
first cycle's Lipschitz bound of the ledger in log-scale transfers, so
that a seed centre remains active at half level over a log-scale
interval of explicit length: hops of the cone of seeds are separated
by an explicit number of geometric steps.  A lemma shows that the
first spike of a given height after a window, if it occurs in the
first three quarters of the remaining time, is approached in a type-I
manner relative to itself with twice its height as constant, so the
spike chapter applies to it unconditionally.  (AC2) closes negatively
for the reason of V62: the window's dissipation floor is not
attributable to the clusters, since the trap bounds the exterior's
enstrophy per cell and not in sum.  No global regularity claim is
made.  The next compulsory companion audit is V80.

% =============================================================================
\section{Ledgers with a spike as reference}
\label{sec:c2v76-ledgers}
% =============================================================================

Let \(u\) be a smooth solution on the torus of side \(2\pi\), \(t'\) a
time, and suppose the type-I approach \(\Lambda_g(t)=(t'-t)^{1/2}\|\Lambda u(t)\|_2^2\le\bar\Lambda\)
on \((t_1,t')\), with \((t'-t_1)^{1/2}\le\pi/4\).  For \(x\) and \(\lambda\le(t'-t_1)^{1/2}\)
define the reference-\(t'\) rescaling \(U^{x,\lambda}(y,s):=\lambda u(x+\lambda y,t'+\lambda^2s)\)
on \([-1,0]\), and the ledger
\begin{equation}
 \Phi^{(t')}_x(\lambda):=\lambda^{-1}\int|u(t'-\lambda^2)|^2\,\psi_\lambda(\cdot-x),\qquad
 \psi_\lambda(y)=e^{-|y|^2/(4\nu\lambda^2)}
 \label{eq:c2v76-ledger}
\end{equation}
(periodized on the torus as in the first cycle's V87), and likewise
the dissipation and flux ledgers \(\lambda D^{(t')}_x(\lambda)\), \(\lambda J^{(t')}_x(\lambda)\) of
the first cycle's exact scaling identity.

\begin{theorem}[Windows, compactness and ledger bounds with reference \(t'\)]
\label{thm:c2v76-transfer}
Under the type-I approach, for every \(t\in(t_1,t')\) with
\(\lambda=(t'-t)^{1/2}\):
\begin{enumerate}[label=\textup{(\roman*)},leftmargin=3.2em]
\item \(U^{x,\lambda}\) is a window rescaling of the class with constant
      \(\bar\Lambda\) on \([-1,-1+\ell]\), \(\ell=\ell_{\bar\Lambda}\): \(Y\le2\bar\Lambda\) there, and
      \(\int_{-1}^{-1+\ell}\|\Delta U^{x,\lambda}\|_2^2\le4\bar\Lambda/\nu\);
\item along any sequence of such rescalings a subsequence converges
      on \(B_R\times[-1,-1+\ell]\) to a strong solution with the bounds of
      Theorem~\ref{thm:c2v42-compactness};
\item the ledgers satisfy \eqref{eq:c2v42-ledgers} with \(M=\bar\Lambda\), and
      the exact scaling identity \(\frac{\dd}{\dd\lambda}\Phi^{(t')}_x=D^{(t')}_x+2\Phi^{(t')}_x/\lambda+J^{(t')}_x\)
      of the first cycle (Theorem V87-derivative-u) holds with reference
      \(t'\) (it is an identity for smooth solutions, independent of the
      reference being singular), with \(|\lambda D|\le8\nu\bar\Lambda\) and
      \(|\lambda J|\le K_J(\nu,\bar\Lambda)\) on the window scales;
\item consequently \(\lambda\mapsto\Phi^{(t')}_x(\lambda)\) is Lipschitz in \(\log\lambda\)
      with constant \(K:=8\nu\bar\Lambda+2C_S^2\cdot2\bar\Lambda\|G\|_{3/2}+K_J\) on
      \((0,(t'-t_1)^{1/2}]\).
\end{enumerate}
\end{theorem}

\begin{proof}
(i) is Proposition~\ref{prop:c2v42-window} applied to the rescaling,
whose enstrophy at \(s=-1\) is \(\Lambda_g(t)\le\bar\Lambda\); (ii) is
Theorem~\ref{thm:c2v42-compactness}, whose proof uses only the window
bounds; (iii) the ledger bounds are the same computation, and the
scaling identity is the first cycle's identity for the rescaled
solution with the Gaussian weight, which holds for every smooth
solution on \([-1,0]\); the flux bound \(K_J\) is the first cycle's bound
of the cubic and pressure fluxes by the window \(H^1\) bound (Lemma V88-order), which uses only \(Y\le2\bar\Lambda\) and the Sobolev
inequality.  (iv) is \(|\lambda\,\dd\Phi/\dd\lambda|\le|\lambda D|+2\Phi+|\lambda J|\) with the
bounds of (iii).
\end{proof}

\begin{corollary}[Seeds as ledger activity, and their persistence in scale]
\label{cor:c2v76-persistence}
Let \(t\in\mathcal T(t')\) be a seed time with \(\Lambda_g(t)\le\bar\Lambda\), \(\lambda=(t'-t)^{1/2}\),
and \(x(t)\) the seed centre of Theorem~\ref{thm:c2v73-cone}.  Then
\[
 \Phi^{(t')}_{x(t)}(\lambda)\ \ge\ c_1:=e^{-1/(4\nu)}c_0\varepsilon_*(\nu,\bar\Lambda),
\]
and for every \(\lambda'\) with \(|\log(\lambda'/\lambda)|\le c_1/(2K)\),
\(\Phi^{(t')}_{x(t)}(\lambda')\ge c_1/2\): at the time \(t'-\lambda'^2\) some cell of side
\(\lambda'\) within lattice distance \(R_{c_1/2}\) of \(x(t)\) carries energy at
least \(c_1'\lambda'\), \(c_1':=c_1/(4(2R_{c_1/2}+1)^3)\) (Proposition~\ref{prop:c2v66-active}).
In the geometric seed sequence \(t_{k+1}=t_k+\frac\ell2\lambda_g(t_k)^2\), the
log-step is \(\frac12|\log(1-\ell/2)|\le\ell/2\), so the seed centre \(x(t_k)\)
remains active at level \(c_1'\) for at least \(n_0:=\lfloor c_1/(K\ell)\rfloor\)
consecutive steps: a hop of the cone of seeds to a centre at distance
more than \((R_{c_1/2}+1)\lambda_g\) from \(x(t_k)\) cannot be forced by the floor
alone before \(n_0\) steps, and any sequence of seed centres can be
chosen so that hops are separated by at least \(n_0\) steps.
\end{corollary}

\begin{proof}
\(\psi_\lambda\ge e^{-1/(4\nu)}\) on \(B_\lambda(x(t))\) gives the ledger bound from
\eqref{eq:c2v73-seed}; the Lipschitz bound of Theorem~\ref{thm:c2v76-transfer}(iv)
gives the persistence; the cell statement is
Proposition~\ref{prop:c2v66-active} with the reference \(t'\)
(its proof uses only the cell energy bound \(a_{\bar\Lambda}\) on the window
and the Gaussian tail); the step count is arithmetic.
\end{proof}

\begin{remark}[What persistence does and does not give]
\label{rem:c2v76-persistence}
Persistence bounds the frequency of hops, not their number: over the
cone from the first gap to the spike's scale there are of order
\(\ell^{-1}\log(\lambda_g(t_1)\|\Lambda u(t')\|_2^2)\) steps, hence up to that number
divided by \(n_0\) hops, unbounded as the spike's height grows.  The
convergence of the cone (Theorem~\ref{thm:c2v74-sufficient}) is not
obtained.  Persistence also shows that the level at a fixed centre
decays at most linearly in log-scale, at rate \(K\), so a seed centre
that ceases to be active does so over a log-interval of length at
least \(c_1/(2K)\): the alternation of active and inactive scales at a
fixed centre has the first cycle's minimal log-period, now at a spike
(Theorem V97-log-period transferred).
\end{remark}

% =============================================================================
\section{The first spike after a window}
\label{sec:c2v76-first-spike}
% =============================================================================

\begin{lemma}[The first spike of a given height is approached type-I relative to itself]
\label{lem:c2v76-first-spike}
Let \(u\) be in the intermediate class with constant \(M\), \(t_j\) a
window time, \(\Lambda\ge4M\), and \(t'\) the first time after \(t_j\) at which
\(\mathfrak Y(t')=\Lambda\).  If \(t'-t_j\le\frac34\lambda_j^2\), then
\(\Lambda_g(t)=(t'-t)^{1/2}\|\Lambda u(t)\|_2^2\le2\Lambda\) for all \(t\in[t_j,t')\): the
spike is approached in a type-I manner relative to itself with
\(\bar\Lambda=2\Lambda\), and the whole spike chapter (V72--V76) applies to it
with that constant, on \((t_j,t')\).
\end{lemma}

\begin{proof}
For \(t\in[t_j,t')\), \(\mathfrak Y(t)\le\Lambda\) by the choice of \(t'\) (and
continuity), so \(\|\Lambda u(t)\|_2^2\le\Lambda/\lambda(t)\le\Lambda/\lambda(t')\), and
\((t'-t)^{1/2}\le(t'-t_j)^{1/2}\le\lambda_j\); since \(\lambda(t')^2=\lambda_j^2-(t'-t_j)\ge\frac14\lambda_j^2\),
\(\Lambda_g(t)\le\lambda_j\Lambda/\lambda(t')\le2\Lambda\).
\end{proof}

\begin{theorem}[Description of the first spike after a window]
\label{thm:c2v76-first-spike}
In the setting of Lemma~\ref{lem:c2v76-first-spike} with
\(\Lambda\ge\Lambda_{\rm seed}(\nu,M)\):
\begin{enumerate}[label=\textup{(\roman*)},leftmargin=3.2em]
\item the spike is seeded at \(t_j\) in the window's configuration at
      level \(\varepsilon_2\) (Theorem~\ref{thm:c2v72-seeded});
\item every \(t\in(t_j,t')\) with \((t'-t)^{1/2}\|\Lambda u(t')\|_2^2>2\Lambda e^{C_4'}\)
      is a seed time, with seed level at least \(c_0\varepsilon_*(\nu,2\Lambda)\)
      (Theorem~\ref{thm:c2v73-cone}(iii)); the seed scales fill
      \((2\Lambda e^{C_4'}\lambda(t')/\Lambda,\ (t'-t_j)^{1/2})=(2e^{C_4'}\lambda(t'),(t'-t_j)^{1/2})\);
\item the seed centres persist over \(n_0(\nu,2\Lambda)\) geometric steps
      (Corollary~\ref{cor:c2v76-persistence});
\item at \(t'\) the flow at the sub-parabolic scale \(\lambda(t')/\Lambda\) is a
      finite configuration with viscosity-only constants
      (Theorem~\ref{thm:c2v71-configuration}).
\end{enumerate}
The item not established is the position of the configuration in
(iv) relative to the seed of (i).
\end{theorem}

\begin{proof}
Each item is the cited statement with \(\bar\Lambda=2\Lambda\); in (ii) the seed
condition \(\lambda_g\|\Lambda u(t')\|_2^2>\Lambda_ge^{C_4'}\) is implied by
\(\lambda_g\|\Lambda u(t')\|_2^2>2\Lambda e^{C_4'}\), and \(\|\Lambda u(t')\|_2^2=\Lambda/\lambda(t')\) gives the
smallest seed scale \(\lambda_g=2e^{C_4'}\lambda(t')\): note that the spike's own
sub-parabolic scale \(\lambda(t')/\Lambda\) is below the smallest seed scale by
the factor \(2\Lambda e^{C_4'}\), so the cone of seeds ends \(\log(2\Lambda e^{C_4'})\)
log-units above the sub-parabolic configuration.
\end{proof}

\begin{firewall}[Item (AC2): multiplicity is not bounded by the dissipation floors]
\label{fw:c2v76-AC2}
On a window the exterior of the configuration has enstrophy at most
\(C_1\varepsilon_{\rm th}^{1/2}\) per cell but unboundedly many cells, so its total
dissipation is not small and the window's dissipation floor (the
Leray floor with reference \(T_*\), or the backward-cone floor with
reference a spike) cannot be attributed to the clusters
(Firewall~\ref{fw:c2v62-AC2}).  A per-cluster dissipation lower bound
exists (each energetic cell is dissipation-active at radius \(R'\),
Theorem~\ref{thm:c2v48-energy-active}), giving \(\ge\varepsilon'\) per cluster and
hence at most \(4M/(\nu\varepsilon')\) clusters per window from the palinstrophy
budget; this is the count \(N_M\) in another form, not an improvement.
No global regularity claim is made.
\end{firewall}

% =============================================================================
\section{Integrated V76 conclusion and next acquisition order}
\label{sec:c2v76-conclusion}
% =============================================================================

\begin{theorem}[What V76 proves]
\label{thm:c2v76-conclusion}
Relative to the first cycle and V1--V75: the transfer of windows,
compactness and ledgers to a spike reference
(Theorem~\ref{thm:c2v76-transfer}), seed persistence
(Corollary~\ref{cor:c2v76-persistence}), the type-I approach to the
first spike (Lemma~\ref{lem:c2v76-first-spike}) and its description
(Theorem~\ref{thm:c2v76-first-spike}).  The full Navier--Stokes
stability/global-regularity theorem remains unproved.
\end{theorem}

\begin{proof}
The cited results.
\end{proof}

\begin{programme}[V77 acquisition order]
\label{prog:c2v77}
\begin{enumerate}[label=\textup{(AC\arabic*)},leftmargin=4.8em]
\item The gap between the smallest seed scale and the sub-parabolic
      configuration.  The cone of seeds ends at scale \(2e^{C_4'}\lambda(t')\),
      the configuration of \(t'\) lives at scale \(\lambda(t')/\Lambda\); on the
      intermediate scales the reference-\(t'\) rescalings have enstrophy
      \(\Lambda_g\) between \(1\) and \(2\Lambda e^{C_4'}\), and the pre-spike floor at
      those scales with reference \(t'\) is trivial (the spike is not
      high relative to the gap).  Determine instead what the
      reference-\(T_*\) statements give at \(t'\) at scales between
      \(\lambda(t')/\Lambda\) and \(\lambda(t')\): the count at scale \(\lambda(t')\) is
      \(C\Lambda^3/\varepsilon^3\), at scale \(\lambda(t')/\Lambda\) it is \(C/\varepsilon^3\); record
      the multiscale count profile \(N(\rho)\) for \(\rho\in[\lambda(t')/\Lambda,\lambda(t')]\)
      (from the \(L^6\) bound at each scale) and whether the energetic
      cells at scale \(\rho\) are nested across \(\rho\) (upward nesting is
      automatic by the refinement decay in reverse: a fine energetic
      cell at level \(\varepsilon\) makes its coarse cell energetic at level
      \(\varepsilon\rho/\rho'\)).
\item Prepare V78--V79: consolidation of V76--V77 and the position
      before the V80 audit; and a decision on the remaining twenty
      versions of the cycle, in particular whether to spend V81--V95
      on a second internal audit of the whole cycle in preparation
      for the restart manifest at V100.
\item The next compulsory companion audit is V80.
\end{enumerate}
\end{programme}

\begin{assessment}[Position after V76]
\label{ass:c2v76-position}
The spike chapter's ledgers and persistence transfer to any spike
under a type-I approach relative to it, which the first spike of a
given height after a window satisfies; the position of the spike's
configuration relative to its seed remains open.  No full stability
or global-regularity claim is made.
\end{assessment}

% =============================================================================
\section*{Version V76 (second cycle) append-only record}
\addcontentsline{toc}{section}{Version V76 (second cycle) append-only record}
% =============================================================================

The complete V75 source of the second cycle was retained before this
continuation.  Its SHA--256 digest was
\begin{center}\scriptsize\ttfamily
1858ecb9bbc25242d3afe32f45ea2dec2114db755961ee4c479c9b9c283c68ae.
\end{center}
V76 transferred the ledgers to a spike reference, proved seed
persistence, and described the first spike after a window.  No full
regularity claim is made.

% =============================================================================
\section{V77: the multiscale profile at a spike time, and the plan for the
rest of the cycle}
\label{sec:c2v77-entry}
% =============================================================================

Items (AC1) and (AC2) of Programme~\ref{prog:c2v77} are carried out.
(AC1): at a spike time, across the scales between the sub-parabolic
and the parabolic one, the count of energetic cells and the Morrey
floor are computed as functions of the scale; the count grows like the
cube of the scale ratio and the floor's exponent like the scale ratio,
both extremal at the sub-parabolic scale, and the energetic cells nest
upward with levels decreasing linearly in the scale ratio.  (AC2): the
plan for V78--V99 is fixed: V78--V79 consolidate the spike chapter's
last block and record the position before the V80 audit; V81--V95
carry out a second internal audit of the whole cycle, block by block,
so that the restart manifest at V100 rests on audited statements;
V96--V99 write the manifest.  No global regularity claim is made.  The
next compulsory companion audit is V80.

% =============================================================================
\section{The multiscale profile}
\label{sec:c2v77-profile}
% =============================================================================

Let \(u\) be a smooth solution on the torus with first singular time
\(T_*\), \(t'<T_*\) with \(\lambda:=\lambda(t')\le\pi/4\) and \(\Lambda:=\max(\mathfrak Y(t'),1)\).
For \(\rho\in[\lambda/\Lambda,\lambda]\) put \(\Lambda_\rho:=\rho\|\Lambda u(t')\|_2^2=\rho\Lambda/\lambda\in[1,\Lambda]\), the
scaled enstrophy at scale \(\rho\), and
\(\mathfrak M_\rho(t'):=\sup_x\rho^{-1}\int_{B_\rho(x)}|u(t')|^2\),
\(N_\rho(\varepsilon):=\#\{\text{cells of side }\rho\text{ with energy}\ge\varepsilon\rho\}\).

\begin{proposition}[Count and floor profiles across scales]
\label{prop:c2v77-profile}
For \(\rho\in[\lambda/\Lambda,\lambda]\):
\begin{enumerate}[label=\textup{(\roman*)},leftmargin=3.2em]
\item (Count.)  \(N_\rho(\varepsilon)\le1728C_S^6\Lambda_\rho^3\varepsilon^{-3}=1728C_S^6(\rho\Lambda/\lambda)^3\varepsilon^{-3}\).
\item (Floor.)  \(\mathfrak M_\rho(t')\ge c_0\,\varepsilon_*^{(\rho)}\) with
      \(\varepsilon_*^{(\rho)}\ge\exp\bigl(-C(\nu)\Lambda_\rho^{3}(\Lambda/\Lambda_\rho)^2\bigr)=\exp\bigl(-C(\nu)\Lambda^2\Lambda_\rho\bigr)\),
      the exponent decreasing from \(C\Lambda^3\) at \(\rho=\lambda\) to \(C\Lambda^2\) at
      \(\rho=\lambda/\Lambda\).
\item (Upward nesting.)  A cell of side \(\rho\) energetic at level
      \(\varepsilon\) lies in a cell of side \(\rho'\ge\rho\) (of a coarser lattice
      containing it, or in one of \(8\) cells of a translate) energetic
      at level \(\varepsilon\rho/\rho'\); in particular the sub-parabolic
      configuration (at most \(1728C_S^6\varepsilon^{-3}\) cells at level \(\varepsilon\))
      lies inside parabolic cells energetic at level \(\varepsilon/\Lambda\), of
      which there are at most \(1728C_S^6\Lambda^6\varepsilon^{-3}\).
\item (Extremality.)  Both the count at a fixed level and the floor's
      exponent are smallest at \(\rho=\lambda/\Lambda\); at no scale in the range
      does the trap give a floor independent of \(\Lambda\).
\end{enumerate}
\end{proposition}

\begin{proof}
(i) Proposition~\ref{prop:c2v58-count} for the rescaling at scale \(\rho\),
whose enstrophy is \(\Lambda_\rho\).  (ii) Run the trap at scale \(\rho\) with
\(\Lambda_0=\Lambda_\rho\) over the remaining rescaled time \((\lambda/\rho)^2=(\Lambda/\Lambda_\rho)^2\):
the proviso \(K\varepsilon_1e^{\Lambda_{\rm tr}(\Lambda/\Lambda_\rho)^2}\le1\) with \(K\le C\Lambda_\rho^2\),
\(p\le C\Lambda_\rho^{1/2}\), \(\Lambda_{\rm tr}\le C\Lambda_\rho^{5/2}\) (proof of
Theorem~\ref{thm:c2v59-floor}) is met by
\(\varepsilon_0\le K^{-1}\Lambda_\rho^{-p}\exp(-(1+p)\Lambda_{\rm tr}(\Lambda/\Lambda_\rho)^2)\), whose logarithm
is at least \(-C\Lambda_\rho^{3}(\Lambda/\Lambda_\rho)^2-C\Lambda_\rho^{1/2}\log\Lambda_\rho\ge-C'\Lambda^2\Lambda_\rho\); the
contradiction argument is that of Theorem~\ref{thm:c2v71-floor}.
(iii) \(\int_{Q'}|u|^2\ge\int_Q|u|^2\ge\varepsilon\rho=(\varepsilon\rho/\rho')\rho'\); the count is
(i) at scale \(\lambda\) and level \(\varepsilon/\Lambda\).  (iv) is the monotonicity of the
expressions in (i)--(ii) in \(\rho\).
\end{proof}

\begin{assessment}[Item (AC1)]
\label{ass:c2v77-AC1}
The profile confirms that the sub-parabolic scale is the scale at which
the trap sees the spike best, and that the gap between the smallest
seed scale \(2e^{C_4'}\lambda\) and the sub-parabolic scale \(\lambda/\Lambda\) (a factor
\(2\Lambda e^{C_4'}\) in scale) is not bridged by the floors, which hold at every
scale but with exponents growing with the scale.  The nesting (iii)
places the sub-parabolic configuration inside a bounded-count family
of parabolic cells energetic at level \(\varepsilon/\Lambda\), which is the
refinement decay once more.  The item is closed with these
statements.
\end{assessment}

% =============================================================================
\section{The plan for V78--V99}
\label{sec:c2v77-plan}
% =============================================================================

\begin{programme}[Plan for the remainder of the second cycle]
\label{prog:c2v77-plan}
\begin{enumerate}[label=\textup{(P\arabic*)},leftmargin=3.6em]
\item V78: consolidate V76--V77 (ledgers at a spike, persistence, the
      first spike, the multiscale profile) and close the spike chapter.
\item V79: the position before the V80 audit, with the list of the
      cycle's theorems that are candidates for the restart manifest's
      summary (M-list), each with its hypotheses and the versions that
      corrected it.
\item V80: compulsory companion audit; delivery.
\item V81--V95: second internal audit of the whole cycle, in blocks:
      V81--V83 the profile and statistical frameworks (V1--V15),
      V84--V85 the moment ledger and statistical dichotomy (V16--V25),
      V86--V87 the rate-free and Rayleigh chapters (V26--V35), V88--V89
      the constant-chasing chapter (V36--V41), V90 audit, V91--V92 the
      intermediate class, spreading and window-structure chapters
      (V42--V55), V93--V94 the trap, propagation and spike chapters
      (V56--V78) as a whole, V95 audit and the audited M-list.
\item V96--V99: the restart manifest for the third cycle (context,
      summary of audited results, imports, gates), written so that the
      new V1 can be assembled from it; V100 closes the cycle.
\end{enumerate}
The audit versions may also contain new results when an audit item
suggests one; each audit version records corrections in place as the
series has done, with the corrected digest in the next record.
\end{programme}

% =============================================================================
\section{Integrated V77 conclusion and next acquisition order}
\label{sec:c2v77-conclusion}
% =============================================================================

\begin{theorem}[What V77 establishes]
\label{thm:c2v77-conclusion}
Relative to the first cycle and V1--V76: the multiscale profile
(Proposition~\ref{prop:c2v77-profile}) and the plan
(Programme~\ref{prog:c2v77-plan}).  The full Navier--Stokes
stability/global-regularity theorem remains unproved.
\end{theorem}

\begin{proof}
The cited results.
\end{proof}

\begin{programme}[V78 acquisition order]
\label{prog:c2v78}
\begin{enumerate}[label=\textup{(AC\arabic*)},leftmargin=4.8em]
\item Consolidate V76--V77 and close the spike chapter with a single
      statement of what is known about a high spike after a window in
      the intermediate class.
\item Record, for the spike chapter, the exact list of hypotheses
      under which each statement holds (reference time, type-I
      approach, first-spike condition, window sub-length), in the form
      the M-list of V79 will use.
\item The next compulsory companion audit is V80.
\end{enumerate}
\end{programme}

\begin{assessment}[Position after V77]
\label{ass:c2v77-position}
The multiscale profile at a spike is computed; the remainder of the
cycle is planned as consolidation, a second internal audit, and the
restart manifest.  No full stability or global-regularity claim is
made.
\end{assessment}

% =============================================================================
\section*{Version V77 (second cycle) append-only record}
\addcontentsline{toc}{section}{Version V77 (second cycle) append-only record}
% =============================================================================

The complete V76 source of the second cycle was retained before this
continuation.  Its SHA--256 digest was
\begin{center}\scriptsize\ttfamily
7df8fd653d733a9f364a246a8167a90264772c9702755d9b6d581a550a72d659.
\end{center}
(The V76 file was corrected in place after its assembly: two legacy
citation names; the digest above is that of the corrected file.)  V77
computed the multiscale profile at a spike and fixed the plan for the
rest of the cycle.  No full regularity claim is made.

% =============================================================================
\section{V78: the spike chapter closed}
\label{sec:c2v78-entry}
% =============================================================================

Items (AC1) and (AC2) of Programme~\ref{prog:c2v78} are carried out:
V76--V77 are consolidated, the spike chapter (V71--V77) is closed with
a single statement about a high spike after a window in the
intermediate class, and the hypotheses of each statement of the
chapter are listed in the form the M-list of V79 will use.  No global
regularity claim is made.  The next compulsory companion audit is V80.

% =============================================================================
\section{V76--V77 consolidated and the chapter closed}
\label{sec:c2v78-consolidated}
% =============================================================================

\begin{theorem}[A high spike after a window, everything known]
\label{thm:c2v78-spike}
Let \(u\) be in the intermediate class with constant \(M\ge1\), \(t_j\) a
window time with scale \(\lambda_j\), \(\Lambda\ge\Lambda_{\rm seed}(\nu,M)\), and \(t'\) the
first time after \(t_j\) with \(\mathfrak Y(t')=\Lambda\), assumed to satisfy
\(t'-t_j\le\frac34\lambda_j^2\).  Write \(\lambda=\lambda(t')\), \(\lambda_g=(t'-t_j)^{1/2}\in[c_g\lambda_j,\lambda_j)\).
\begin{enumerate}[label=\textup{(\arabic*)},leftmargin=3.2em]
\item (Window.)  At \(t_j\) the flow is a finite configuration of at
      most \(N_M\) energetic parabolic cells; on the window
      \([t_j,t_j+\ell\lambda_j^2]\) its activity stays within \(R'_{\rm inv}\lambda_j\) of
      the energetic seeds and the exterior is trapped after
      \(\varepsilon_{\rm th}^{1/2}\lambda_j^2\).
\item (Seed.)  A ball of radius \(\lambda_g\) at time \(t_j\), inside one of at
      most \(N_M(\varepsilon_2)\) parabolic cells energetic at level \(\varepsilon_2\),
      carries energy at least \(c_0\varepsilon_*(\nu,M)\lambda_g\).
\item (Approach.)  \(\Lambda_g\le2\Lambda\) on \([t_j,t')\); every \(t\in(t_j,t')\) with
      \(t'-t>(2e^{C_4'}\lambda)^2\) is a seed time with a ball of radius
      \((t'-t)^{1/2}\) carrying energy at least \(c_0\varepsilon_*(\nu,2\Lambda)(t'-t)^{1/2}\);
      the ledgers with reference \(t'\) are bounded and Lipschitz in
      log-scale with constant \(K(\nu,2\Lambda)\); seed centres persist over
      \(n_0(\nu,2\Lambda)\) geometric steps; consecutive seeds within one
      window are localized to the earlier configuration.
\item (Spike.)  At \(t'\), at the sub-parabolic scale \(\lambda/\Lambda\), the flow
      is a finite configuration with viscosity-only constants; the
      Morrey energy at scale \(\rho\in[\lambda/\Lambda,\lambda]\) is at least
      \(\exp(-C(\nu)\Lambda^2\rho\Lambda/\lambda)\), and the energetic cells nest upward
      with linearly decaying levels.
\item (Open.)  The position of the configuration in (4) relative to
      the seed in (2); the convergence of the cone of seeds
      (sufficient: a single cluster or a unique seed at every scale);
      the evolution between the window's end and \(t'\).
\end{enumerate}
\end{theorem}

\begin{proof}
(1) Theorems~\ref{thm:c2v57-no-spread}, \ref{thm:c2v66-propagation},
Corollary~\ref{cor:c2v66-configuration}.  (2) Theorem~\ref{thm:c2v72-seeded}.
(3) Lemma~\ref{lem:c2v76-first-spike}, Theorem~\ref{thm:c2v73-cone}(iii),
Theorem~\ref{thm:c2v76-transfer}, Corollary~\ref{cor:c2v76-persistence},
Proposition~\ref{prop:c2v73-one-step}.  (4) Theorem~\ref{thm:c2v71-configuration},
Proposition~\ref{prop:c2v77-profile}.  (5) Firewalls~\ref{fw:c2v72-location},
\ref{fw:c2v73-hopping}, \ref{fw:c2v74-conditions}.
\end{proof}

% =============================================================================
\section{Hypotheses of the spike chapter}
\label{sec:c2v78-hypotheses}
% =============================================================================

\begin{proposition}[Hypothesis ledger of the spike chapter]
\label{prop:c2v78-hypotheses}
The statements of V71--V77 hold under the following hypotheses, and
under no weaker ones proved in the series.
\begin{enumerate}[label=\textup{(H\arabic*)},leftmargin=3.6em]
\item (General singularity.)  Theorem~\ref{thm:c2v71-floor} (sub-parabolic
      floor), Proposition~\ref{prop:c2v71-counts}, Theorem~\ref{thm:c2v71-configuration},
      Proposition~\ref{prop:c2v77-profile}: a smooth solution on the
      torus of side \(2\pi\) with first singular time \(T_*\), at any
      \(t<T_*\) with \(\lambda(t)\le\pi/4\); constants depend on \(\nu\) and on
      \(\max(\mathfrak Y(t),1)\) where stated.
\item (Any two times.)  Theorem~\ref{thm:c2v72-prespike} (pre-spike
      floor) and Theorem~\ref{thm:c2v73-cone} (cone of seeds): a smooth
      solution and two times \(t<t_2\) with \(t_2-t\le\pi^2/16\); no
      singularity is needed.
\item (Intermediate class, window and spike.)  Theorem~\ref{thm:c2v72-seeded}:
      constant \(M\), window time \(t_j\), spike height \(\Lambda\ge\Lambda_{\rm seed}(\nu,M)\);
      the cone constant \(C_1\) enters through \(c_g\).
\item (Type-I approach relative to the spike.)  Theorem~\ref{thm:c2v73-cone}(iii),
      Proposition~\ref{prop:c2v73-one-step}, Theorem~\ref{thm:c2v76-transfer},
      Corollary~\ref{cor:c2v76-persistence}: \(\Lambda_g\le\bar\Lambda\) on \((t_1,t')\);
      supplied by Lemma~\ref{lem:c2v76-first-spike} for the first spike
      of height \(\Lambda\) after a window with \(t'-t_j\le\frac34\lambda_j^2\), with
      \(\bar\Lambda=2\Lambda\).
\item (Window sub-length.)  All window statements use the sub-window
      of length \(\min(\ell_M,1/(24C_\flat'\nu))\) when \(\nu\ell_M\) is not small
      ((A9) of V68).
\item (Threshold.)  \(\varepsilon_{\rm th}\le\min(\varepsilon_*/2,\varepsilon_\star,\varepsilon_{\star\star})\), and
      in V73 the lowered \(\varepsilon'''\); all constants of the propagation
      chapter are evaluated at the threshold in use.
\end{enumerate}
\end{proposition}

\begin{proof}
Inspection of the statements; each cites its hypotheses.
\end{proof}

% =============================================================================
\section{Integrated V78 conclusion and next acquisition order}
\label{sec:c2v78-conclusion}
% =============================================================================

\begin{theorem}[What V78 establishes]
\label{thm:c2v78-conclusion}
Relative to the first cycle and V1--V77: the closing statement of the
spike chapter (Theorem~\ref{thm:c2v78-spike}) and its hypothesis
ledger (Proposition~\ref{prop:c2v78-hypotheses}).  The full
Navier--Stokes stability/global-regularity theorem remains unproved.
\end{theorem}

\begin{proof}
The cited results.
\end{proof}

\begin{programme}[V79 acquisition order]
\label{prog:c2v79}
\begin{enumerate}[label=\textup{(AC\arabic*)},leftmargin=4.8em]
\item The M-list: the cycle's theorems that are candidates for the
      restart manifest's summary, each with its hypotheses, its
      constants' dependence, the versions that corrected or withdrew
      parts of it, and its status (proved in full, proved at sketch
      level, conditional).
\item The position before the V80 audit, including the three gates
      as they stand after the trap, propagation and spike chapters.
\item The next compulsory companion audit is V80.
\end{enumerate}
\end{programme}

\begin{assessment}[Position after V78]
\label{ass:c2v78-position}
The spike chapter is closed with a single statement and a hypothesis
ledger.  No full stability or global-regularity claim is made.
\end{assessment}

% =============================================================================
\section*{Version V78 (second cycle) append-only record}
\addcontentsline{toc}{section}{Version V78 (second cycle) append-only record}
% =============================================================================

The complete V77 source of the second cycle was retained before this
continuation.  Its SHA--256 digest was
\begin{center}\scriptsize\ttfamily
140564ff15e702fd4b2ae03ac9470a4945b448e27fbd1e520cbadd55ddaa9354.
\end{center}
V78 closed the spike chapter.  No full regularity claim is made.

% =============================================================================
\section{V79: the M-list and the position before the V80 audit}
\label{sec:c2v79-entry}
% =============================================================================

Items (AC1) and (AC2) of Programme~\ref{prog:c2v79} are carried out:
the M-list of the cycle's candidate results for the restart manifest,
with hypotheses, constant dependence, corrections and status; and the
position of the three gates after the trap, propagation and spike
chapters.  No global regularity claim is made.  The next compulsory
companion audit is V80.

% =============================================================================
\section{The M-list}
\label{sec:c2v79-mlist}
% =============================================================================

\begin{theorem}[M-list: candidate results of the second cycle for the restart manifest]
\label{thm:c2v79-mlist}
Status codes: F = proved in full in the series; S = proved at sketch
level with structure displayed; C = conditional on a stated
hypothesis; W = withdrawn or corrected (with the correcting version).
\begin{enumerate}[label=\textup{(M\arabic*)},leftmargin=3.6em]
\item \emph{Profile framework} (V1--V10, F): the self-similar profile
      \(V\), the operator \(L_-\) with Gaussian weight \(G\), the second
      identity \(\dd\mathcal L/\dd\sigma=-\|\partial_\sigma V\|_G^2-\langle N,\partial_\sigma V\rangle_G\),
      weak lower semicontinuity of the defect, the dichotomy (A)/(B),
      and the import (E9) excluding (A) at homogeneous type-I points.
      Hypotheses: type-I rate at a homogeneous point; imports
      (E1)--(E9).
\item \emph{Statistical framework} (V11--V15, F): recurrence, invariant
      and blow-up measures, exact mean identities
      \(\int j=-8\int\mathcal L\), \(\int\langle N,Y\rangle=-\int\|Y\|^2\le-\underline\delta\),
      rate-free second identity for any smooth solution.
\item \emph{Moment ledger} (V16--V22, F; V17 and V21 corrected sentences
      of V16 and V20): centre mode, momentum, tensor, inertia,
      enstrophy identities; \(R=8\nu\mathcal L-4\nu^2\mathcal E\); the ledger
      search closed without an incompatible row.
\item \emph{Statistical dichotomy and coherence} (V23--V25, F):
      \(0\in\mathcal F(x_0)\) iff intermediate scales; compact dissipation
      \(\Leftrightarrow\) Gaussian activity at moving centres; coherence
      needed for moving-centre measures.
\item \emph{Rate-free chapter} (V26--V30, F): parabolic scales,
      suitable weak limits (E2), frequency escape, Gaussian
      invisibility, the paraboloid dissipation seen by the Dirichlet
      ledger, the bridge \(\|L_-V\|_G\ge2\mathcal L/\varphi^{1/2}\), the trichotomy.
\item \emph{Rayleigh chapter} (V31--V35, F): \(r=2\mathcal L/\varphi\), the
      Rayleigh identity, forced high-mode escape, push budget,
      \(\mathrm{Var}\ge v_0>0\) on the family (non-explicit constant).
\item \emph{Constant-chasing gate} (V36--V41, F for the upper side; the
      lower side open): Gaussian Hardy, flux bound
      \(|j|\le\kappa(d+4\varphi)+K_P(d+3\varphi)^{1/2}\), sharpened threshold
      \(\kappa<\sqrt2\), modewise identities \(\langle n_k\rangle=-\mu_k\langle a_k\rangle\) and
      \(\langle a_k\rangle=O(1/k)\).
\item \emph{Intermediate class and windows} (V42--V45, F): type-I windows
      of length \(\ell_M=3\nu^3/(8C_AM^2)\), compactness to strong solutions,
      ledger bounds, window trichotomy, window flux signed for
      \(M<M_*\).
\item \emph{Spreading and window-structure chapters} (V46--V55, F; the
      regular-cell counting corrected in V61 to the homogeneous
      seminorm): backward cone at windows, far/local pressure on
      cells, energy-active \(\Rightarrow\) dissipation-active, palinstrophy
      budget and enstrophy band, sharp enstrophy inequality, window
      stretching bound, influx bound, no travelling windows.
\item \emph{Cell trap and exclusion of spread windows} (V57, F; V68
      corrected a justification): far-pressure gradient bounded by the
      maximal cell energy; interpolation \(y^2\le2e(d+C^2y)\); weighted
      inequalities \(W'\le-aW^2/\varepsilon+b(W+\varepsilon)\), \(\varepsilon'\le b'(\varepsilon+W)\);
      capture and trapping; \emph{the intermediate class has no spread
      windows}, \(\max_\alpha e_\alpha(-1)\ge\varepsilon_*(\nu,M)\).
\item \emph{Delay statements} (V56, W: Corollary~\ref{cor:c2v56-delta}
      and Theorem~\ref{thm:c2v56-delay} withdrawn in V68; repaired as
      Proposition~\ref{prop:c2v68-delay} with delay \(\log2/(2b^\sharp)\)).
      The global inequality \eqref{eq:c2v56-global} stands (F).
\item \emph{Morrey floors} (V59, V71, F): for any singularity and every
      \(t<T_*\), \(\mathfrak M(t)\ge e^{-C'(\nu)\max(\mathfrak Y,1)^2}\) and the
      sub-parabolic floor \(\mathfrak M'(t)\ge e^{-C(\nu)\max(\mathfrak Y,1)^2}\); type-I
      \(\Rightarrow\) persistent floor; vanishing parabolic energy along a
      sequence \(\Rightarrow\) type II along it.
\item \emph{Finite configurations and clusters} (V58--V59, F): at a window
      between \(1\) and \(N_M=1728C_S^6(2M)^3(\varepsilon_*/2)^{-3}\) energetic cells;
      cluster limits with small far energy; counts at any time by
      \(\mathfrak Y^3\).
\item \emph{Finite propagation} (V61, V66, F; V58's localized
      inequalities S in V58, F in V69): non-quiet set within
      \(R'_{\rm inv}\) of the energetic seeds through the window; crude and
      far growth lemmas; non-regular cells handled by the joint
      argument.
\item \emph{Exterior trap} (V62, F): unconditional after the capture
      time; linear enstrophy inequality outside the seed
      neighbourhood; negative: maxima not sums.
\item \emph{Type-I structure at every time} (V63, F) and \emph{singular
      points at active scales} (V66--V67, F; V67 corrected in place):
      active scales are energetic cells, homogeneous points number at
      most \(N_M(c')\), quiet-scale dichotomy; negative: no alignment of
      active scales, no finiteness of the singular set, no accumulation
      of trap-active scales.
\item \emph{Persistence and oscillation} (V64, F) and the hysteresis
      no-go; \emph{refinement obstruction} (V65, F).
\item \emph{Pre-spike floor and seeding} (V72, F): any two times; every
      high spike after a window is seeded in the window's
      configuration at level \(\varepsilon_2\).
\item \emph{Cone of seeds} (V73--V76, F with hypothesis (H4) for the
      type-I approach; C for convergence): seed times, one-step
      localization, hopping; two sufficient conditions for convergence;
      ledgers with a spike as reference, persistence over \(n_0\) steps;
      the first spike after a window satisfies (H4).
\item \emph{Multiscale profile at a spike} (V77, F).
\end{enumerate}
\end{theorem}

\begin{proof}
Each item cites the versions in which it is proved and the versions
that corrected it; the status codes follow the audits of V68 and the
proofs' form.
\end{proof}

% =============================================================================
\section{The three gates after the trap, propagation and spike chapters}
\label{sec:c2v79-gates}
% =============================================================================

\begin{assessment}[The gates before the V80 audit]
\label{ass:c2v79-gates}
\emph{G1 (type-I Liouville).}  A nonzero element of the ancient family
at a type-I point is, at every time, a finite configuration of at
most \(N_C\) energetic parabolic cells with a uniformly diffuse far
field, its homogeneous points are at most \(N_C(c')\), and its
self-similar elements are excluded (E9); the missing statement is
still a single mean inequality for one blow-up measure (V15), or an
independent exclusion of the finite-configuration ancient solutions.
\emph{G2 (intermediate class).}  Windows are finite configurations;
spread windows are excluded; every high spike is seeded in the
window's configuration and approached in a type-I manner relative to
itself (first spike), with a cone of seeds whose convergence is the
open item; the evolution between the window's end and the spike is
unconstrained by the series.  \emph{G3 (strong type II).}  The floors
degenerate exponentially in \(\mathfrak Y^2\); the counts grow like
\(\mathfrak Y^3\); a singularity with vanishing parabolic energy along a
sequence is of type II along it; no new constraint.  The cycle's
strongest results are the exclusion of spread windows and the Morrey
floors, both valid for any singularity in their general forms, and the
finite propagation of energy activity on windows.  No global
regularity claim is made.
\end{assessment}

% =============================================================================
\section{Integrated V79 conclusion and next acquisition order}
\label{sec:c2v79-conclusion}
% =============================================================================

\begin{theorem}[What V79 establishes]
\label{thm:c2v79-conclusion}
Relative to the first cycle and V1--V78: the M-list
(Theorem~\ref{thm:c2v79-mlist}) and the position of the gates
(Assessment~\ref{ass:c2v79-gates}).  The full Navier--Stokes
stability/global-regularity theorem remains unproved.
\end{theorem}

\begin{proof}
The cited results.
\end{proof}

\begin{programme}[V80 acquisition order]
\label{prog:c2v80}
\begin{enumerate}[label=\textup{(AC\arabic*)},leftmargin=4.8em]
\item The compulsory companion audit: read the current SM and YM
      notes and record their digests and decision.
\item Record the position after eighty versions and begin the second
      internal audit per Programme~\ref{prog:c2v77-plan}(P4) with the
      first block, V1--V5 (the profile framework: the dictionary, the
      profile equation, the second identity and its rigid case, the
      Lamb--Bernoulli cross term), re-deriving each identity.
\end{enumerate}
\end{programme}

\begin{assessment}[Position after V79]
\label{ass:c2v79-position}
The M-list is recorded with statuses; the gates stand as described.
No full stability or global-regularity claim is made.
\end{assessment}

% =============================================================================
\section*{Version V79 (second cycle) append-only record}
\addcontentsline{toc}{section}{Version V79 (second cycle) append-only record}
% =============================================================================

The complete V78 source of the second cycle was retained before this
continuation.  Its SHA--256 digest was
\begin{center}\scriptsize\ttfamily
65c70628c1aad471c6a79498589807fe9898ea1363db289e1fc17fe2175e21ef.
\end{center}
V79 recorded the M-list and the position of the gates.  No full
regularity claim is made.

% =============================================================================
\section{V80: compulsory companion audit, the position after eighty
versions, and the second internal audit begun (V1--V5)}
\label{sec:c2v80-entry}
% =============================================================================

This is the eightieth version of the second cycle.  The compulsory
review of the two companion programmes is performed first.  The
position after eighty versions is recorded.  The second internal
audit of the cycle (Programme~\ref{prog:c2v77-plan}(P4)) begins with
the profile framework, V1--V5: the profile equation, the dictionary,
the symmetry of the drift operator, the second identity and its rigid
case, and the Lamb--Bernoulli form of the cross term are re-derived;
no correction is needed.  No global regularity claim is made.  The
next compulsory companion audit is V85.

% =============================================================================
\section{Compulsory V80 review of the current companion notes}
\label{sec:c2v80-companion-audit}
% =============================================================================

The current companion files in \texttt{latex\_notes} were inspected
read-only.  The frozen checkpoints are
\begin{align}
 &\texttt{Renewal\_SM\_Einstein\_Continuum\_Research\_Notes\_V18.tex},
 \notag\\[-2pt]
 &\qquad
 \texttt{7bfd396f174142755956647463fa89f246316c50a76b280cbe848d4429ac07c5},
 \label{eq:c2v80-SM-checkpoint}\\
 &\texttt{Renewal\_Yang\_Mills\_Mass\_Gap\_Research\_Notes\_V91.tex},
 \notag\\[-2pt]
 &\qquad
 \texttt{65a24353e41176b134b9d084b908abfe94af2b92a5765e2d17efd5d89395f38f}.
 \label{eq:c2v80-YM-checkpoint}
\end{align}
(SM V18 has 7824 lines; YM V91 has 41379 lines.)

\emph{SM V18 (seventh series).}  The SM programme supplied the finite
Standard Model Yukawa block for an arbitrary number of generations
(the Higgs-dependent finite matrix left unspecified by V17), proved
its gauge equivariance, computed its quadratic and quartic traces,
distinguished the doubled heat trace from the physical chiral weight,
and described its improved-overlap realization.

\emph{YM V91.}  The YM programme analysed the component structure with
one transversal root, proved an exact rooted-forest polynomial, glued
rooted attachments to the transversal tree, closed the single-anchor
row of its allocation card, and recorded that with several anchors in
one local component the acyclic replacement is a forest whose
partition count is not yet controlled.

\begin{theorem}[V80 companion-audit decision]
\label{thm:c2v80-companion-decision}
No SM V18 trace constant and no YM V91 forest polynomial or gluing
bound is imported.  Neither bears on the second cycle's chapters or
on the gates.
\end{theorem}

\begin{proof}
The SM results concern finite internal traces of lattice Dirac
operators; the YM results concern rooted forests of a lattice clock.
None is a Navier--Stokes statement.
\end{proof}

% =============================================================================
\section{Position after eighty versions}
\label{sec:c2v80-position}
% =============================================================================

\begin{assessment}[Position after eighty versions of the second cycle]
\label{ass:c2v80-position-eighty}
The cycle's mathematical work is complete in the sense of
Programme~\ref{prog:c2v77-plan}: the profile, statistical, ledger,
rate-free, Rayleigh, constant-chasing, intermediate-class, spreading,
window-structure, trap, propagation and spike chapters are written,
consolidated, and listed with statuses in the M-list
(Theorem~\ref{thm:c2v79-mlist}).  The remaining twenty versions are
the second internal audit (V80--V95) and the restart manifest
(V96--V99).  The audit re-derives the load-bearing identities and
inequalities of each block, corrects in place where needed, and marks
each M-list item as audited; new results are admitted when an audit
item suggests one.  No global regularity claim is made.
\end{assessment}

% =============================================================================
\section{Second internal audit, block 1: the profile framework (V1--V5)}
\label{sec:c2v80-audit-block1}
% =============================================================================

\begin{theorem}[Audit of V1--V5]
\label{thm:c2v80-audit}
The following are re-derived and found correct as stated.
\begin{enumerate}[label=\textup{(B1.\arabic*)},leftmargin=4.2em]
\item \emph{Profile equation} (Lemma~\ref{lem:c2v2-dictionary}).  With
      \(W(s,y)=(-s)^{-1/2}V(y(-s)^{-1/2},\sigma)\), \(\sigma=-\log(-s)\):
      \(\partial_sW=(-s)^{-3/2}[\frac12V+\frac12z\cdot\nabla V+\partial_\sigma V]\) (from
      \(\partial_s(-s)^{-1/2}=\frac12(-s)^{-3/2}\), \(\partial_sz=\frac12z(-s)^{-1}\),
      \(\partial_s\sigma=(-s)^{-1}\)), \(\nu\Delta_yW=(-s)^{-3/2}\nu\Delta_zV\),
      \((W\cdot\nabla_y)W=(-s)^{-3/2}(V\cdot\nabla_z)V\), \(\nabla_yP=(-s)^{-3/2}\nabla_zP_V\);
      hence \(\partial_\sigma V=\nu\Delta V-\frac12z\cdot\nabla V-\frac12V-N(V)=L_-V-N(V)\).
\item \emph{Dictionary} \eqref{eq:c2v2-dictionary}.  At \(s=-\lambda^2\),
      \(W=\lambda^{-1}V(y/\lambda)\): \(\Phi_W=\lambda^{-1}\lambda^{-2}\lambda^3\|V\|_G^2=\|V\|_G^2\);
      \(\int|\nabla_yW|^2\psi_\lambda\,\dd y=\lambda^{-4}\lambda^3\|\nabla V\|_G^2\), so with the first
      cycle's \(D_W=4\nu\int|\nabla W|^2\psi_\lambda\) one has \(\lambda D_W=4\nu\|\nabla V\|_G^2\);
      and \(\langle V,N(V)\rangle_G=-\frac12\int|V|^2V\cdot\nabla G-\int P_VV\cdot\nabla G=\frac1{4\nu}\int(|V|^2+2P_V)(z\cdot V)G\)
      (\(\nabla G=-\frac z{2\nu}G\), \(\operatorname{div}V=0\)), matching
      \(\lambda\mathcal J_W\).  The first identity in the form
      \(\varphi'=d+2\varphi+j\) (derivative in \(\log\lambda\), \(\dd\sigma=-2\dd\log\lambda\))
      follows from \(\frac{\dd}{\dd\sigma}\|V\|_G^2=-2\nu\|\nabla V\|_G^2-\|V\|_G^2-2\langle V,N\rangle_G\).
\item \emph{Symmetry of \(L_-\).}  \(\operatorname{div}(G\nabla V)=G(\Delta V-\frac z{2\nu}\cdot\nabla V)\),
      so \(\nu\Delta V-\frac12z\cdot\nabla V=\nu G^{-1}\operatorname{div}(G\nabla V)\) and
      \(\langle L_-V,\phi\rangle_G=-\nu\langle\nabla V,\nabla\phi\rangle_G-\frac12\langle V,\phi\rangle_G\).
\item \emph{Second identity} \eqref{eq:c2v2-second}.
      \(\frac{\dd}{\dd\sigma}\mathcal L=\nu\langle\nabla V,\nabla\partial_\sigma V\rangle_G+\frac12\langle V,\partial_\sigma V\rangle_G=-\langle L_-V,\partial_\sigma V\rangle_G\)
      by (B1.3); inserting \(L_-V=\partial_\sigma V+N(V)\) gives
      \(-\|\partial_\sigma V\|_G^2-\langle N(V),\partial_\sigma V\rangle_G\), inserting
      \(\partial_\sigma V=L_-V-N(V)\) gives \(-\|L_-V\|_G^2+\langle L_-V,N(V)\rangle_G\).
      The finiteness and the cutoff justification rely on the first
      cycle's bounds on the class (Corollary V56-class of the first
      cycle: \(W,\nabla W,\nabla^2W\) bounded on half-lines) and the
      Calder\'on--Zygmund bound for \(\nabla P\); correct as stated.
      The rigid case (Theorem~\ref{thm:c2v2-rigidity}) follows.
\item \emph{Lamb--Bernoulli cross term} \eqref{eq:c2v3-cross}.
      \((V\cdot\nabla)V=\nabla\frac12|V|^2-V\times\Omega=\nabla\frac12|V|^2+\Omega\times V\), so
      \(N(V)=\Omega\times V+\nabla\Pi\); \((\Omega\times V)\cdot\partial_\sigma V=\det(\Omega,V,\partial_\sigma V)\);
      \(\langle\nabla\Pi,\partial_\sigma V\rangle_G=-\int\Pi\,\partial_\sigma V\cdot\nabla G=\frac1{2\nu}\int\Pi(z\cdot\partial_\sigma V)G\)
      using \(\operatorname{div}\partial_\sigma V=0\).  Correct.
\item \emph{V1 and V4--V5 at statement level.}  The scale-invariant form
      \(\varphi'=2\varphi-\xi\) with \(\xi:=-(d+j)\) is the first identity
      rewritten; the log-average law is its integration over a log-scale
      interval with the bounds of the class; the weak lower
      semicontinuity of the defect on the compact family and the
      small-defect statement (via (E8) and (E6)) are as stated; V5's
      original-variables form of alternative (B) is the dictionary
      applied to \(\partial_\sigma V\).  No correction.
\end{enumerate}
M-list items (M1) and the first half of (M2) are marked audited.
\end{theorem}

\begin{proof}
The displayed computations.
\end{proof}

% =============================================================================
\section{Integrated V80 conclusion and next acquisition order}
\label{sec:c2v80-conclusion}
% =============================================================================

\begin{theorem}[What V80 establishes]
\label{thm:c2v80-conclusion}
Relative to the first cycle and V1--V79: the companion-audit decision
(Theorem~\ref{thm:c2v80-companion-decision}) and the audit of V1--V5
(Theorem~\ref{thm:c2v80-audit}).  The full Navier--Stokes
stability/global-regularity theorem remains unproved.
\end{theorem}

\begin{proof}
The cited results.
\end{proof}

\begin{programme}[V81 acquisition order]
\label{prog:c2v81}
\begin{enumerate}[label=\textup{(AC\arabic*)},leftmargin=4.8em]
\item Audit block 2, V6--V10: the minimal windowed defect \(m(\ell)\)
      (attainment, superadditivity, the limit \(\underline\delta\)), the
      regularity of the defect (\(\nabla^2W,\nabla^3W\in L^2_tL^2_x\) on
      half-lines), the defect equation \(\partial_\sigma Y=L_-Y-N'(V)Y\) with its
      Gaussian energy identity, the OU spectrum of \(L_-\) and the
      Gaussian Poincar\'e inequality, the three averaged signs, and
      the import (E9) with its consequence; re-derive the defect
      equation and the spectrum, and check the sign conventions of
      the three averaged inequalities.
\item Audit block 3, V11--V15: the rescaling flow on the compact
      family, Birkhoff minimal sets, Krylov--Bogolyubov measures, and
      the exact mean identities; check that the mean identities are
      the integrals of the first and second identities with the
      boundary terms vanishing by invariance.
\item The next compulsory companion audit is V85.
\end{enumerate}
\end{programme}

\begin{assessment}[Position after V80]
\label{ass:c2v80-position}
The profile framework is audited without correction; the second
internal audit continues.  No full stability or global-regularity
claim is made.
\end{assessment}

% =============================================================================
\section*{Version V80 (second cycle) append-only record}
\addcontentsline{toc}{section}{Version V80 (second cycle) append-only record}
% =============================================================================

The complete V79 source of the second cycle was retained before this
continuation.  Its SHA--256 digest was
\begin{center}\scriptsize\ttfamily
b73a529326e6655913b71c4f9e6d143b296fb647f3b82b825596ced017feafed.
\end{center}
V80 performed the compulsory companion audit and audited V1--V5.  No
full regularity claim is made.

% =============================================================================
\section{V81: second internal audit, blocks 2--3 (V6--V15)}
\label{sec:c2v81-entry}
% =============================================================================

Items (AC1) and (AC2) of Programme~\ref{prog:c2v81} are carried out.
The defect equation, the Gaussian energy identity for the defect with
its Hardy inequality, the Ornstein--Uhlenbeck spectrum and Gaussian
Poincar\'e inequality, the invariance of the divergence-free subspace,
the three averaged signs, and the exact mean identities of the
invariant measures are re-derived; no correction is needed.  The
import ledger (E1)--(E9) is unchanged.  No global regularity claim is
made.  The next compulsory companion audit is V85.

% =============================================================================
\section{Block 2: the defect and its identities (V6--V10)}
\label{sec:c2v81-block2}
% =============================================================================

\begin{theorem}[Audit of V6--V10]
\label{thm:c2v81-block2}
The following are re-derived and found correct as stated.
\begin{enumerate}[label=\textup{(B2.\arabic*)},leftmargin=4.2em]
\item \emph{Defect equation} \eqref{eq:c2v7-defect-equation}.
      Differentiating \(\partial_\sigma V=L_-V-N(V)\) in \(\sigma\) with
      \(Y=\partial_\sigma V\): \(\partial_\sigma Y=L_-Y-N'(V)Y\),
      \(N'(V)Y=(Y\cdot\nabla)V+(V\cdot\nabla)Y+\nabla P'\), \(P'=\partial_\sigma P_V\); the
      integrability statements rest on the first cycle's half-line
      bounds and the Calder\'on--Zygmund bound, as in (B1.4).
\item \emph{Gaussian energy identity} \eqref{eq:c2v7-defect-energy}.
      \(\langle Y,L_-Y\rangle_G=-\nu\|\nabla Y\|_G^2-\frac12\|Y\|_G^2\) by (B1.3);
      \(\langle Y,(V\cdot\nabla)Y\rangle_G=\int V\cdot\nabla\frac12|Y|^2\,G=-\frac12\int|Y|^2V\cdot\nabla G=\frac1{4\nu}\int|Y|^2(z\cdot V)G\)
      (\(\operatorname{div}V=0\), \(\nabla G=-\frac z{2\nu}G\));
      \(\langle Y,\nabla P'\rangle_G=-\int P'\,Y\cdot\nabla G=\frac1{2\nu}\int P'(z\cdot Y)G\)
      (\(\operatorname{div}Y=0\)); the three nonlinear terms enter with the
      minus sign of \(-\langle Y,N'(V)Y\rangle_G\).  Correct.
\item \emph{Gaussian Hardy inequality.}  From
      \(z\cdot\nabla G=-\frac{|z|^2}{2\nu}G\) and
      \(\int|Y|^2z\cdot\nabla G=-\int\operatorname{div}(|Y|^2z)G=-\int(3|Y|^2+2Y\cdot(z\cdot\nabla)Y)G\)
      one gets \(\frac1{2\nu}\int|z|^2|Y|^2G=\int(3|Y|^2+2Y\cdot(z\cdot\nabla)Y)G\),
      hence \(\||z|Y\|_G^2\le2\nu(3\|Y\|_G^2+2\||z|Y\|_G\|\nabla Y\|_G)\) and
      \(\||z|Y\|_G\le C(\nu)(\|Y\|_G+\|\nabla Y\|_G)\).  Correct.
\item \emph{Spectrum and Poincar\'e} (Proposition~\ref{prop:c2v8-spectrum}).
      With \(z=(2\nu)^{1/2}x\): \(\nu\Delta_z=\frac12\Delta_x\), \(\frac12z\cdot\nabla_z=\frac12x\cdot\nabla_x\),
      \(G=e^{-|x|^2/2}\); \(A=\frac12(\Delta_x-x\cdot\nabla_x)\) has eigenvalues
      \(-k/2\) on Hermite polynomials, so \(L_-=A-\frac12\) has spectrum
      \(\{-\frac12-\frac k2\}\); the Gaussian Poincar\'e inequality with
      constant \(1\) gives \(\nu\|\nabla_zY\|_G^2=\frac12\|\nabla_xY\|^2\ge\frac12\|Y-\bar Y\|^2\).
      Invariance of the divergence-free subspace:
      \(\operatorname{div}(z\cdot\nabla V)=\operatorname{div}V+z\cdot\nabla\operatorname{div}V\),
      so \(\operatorname{div}L_-V=(\nu\Delta-\frac12z\cdot\nabla-1)\operatorname{div}V\).
      Correct.
\item \emph{Three averaged signs} (Theorem~\ref{thm:c2v8-three-signs}).
      (sign 1) integrates the second identity:
      \(\int_\sigma^{\sigma+L}\langle N,\partial_\sigma V\rangle_G=-\int\|\partial_\sigma V\|_G^2-\mathcal L\big|_\sigma^{\sigma+L}\le-m(L)+\operatorname{osc}\mathcal L\).
      (sign 2) integrates the defect energy identity with the
      Poincar\'e inequality:
      \(\int\langle Y,N'(V)Y\rangle_G=-\int(\nu\|\nabla Y\|^2+\frac12\|Y\|^2)-\frac12\|Y\|^2\big|_\sigma^{\sigma+L}\).
      (sign 3) is the product rule
      \(\frac{\dd}{\dd\sigma}\langle N,Y\rangle=\langle N'Y,Y\rangle+\langle N,\partial_\sigma Y\rangle\)
      integrated, with \(|\langle N,Y\rangle|\le C_N^{1/2}C_Y\).  The
      inequalities' directions and the constants
      \(\frac{C_Y^2}{2L}\), \(\frac{C_Y^2+4C_N^{1/2}C_Y}{2L}\) are as
      stated.  Correct.
\item \emph{Minimal windowed defect} (V6).  \(m(\ell)=\inf_{\mathcal F}\int_0^\ell\|\partial_\sigma V\|_G^2\)
      is attained by compactness of the family and lower
      semicontinuity of the defect (V4), superadditive by
      concatenation of windows and the invariance of the family under
      the rescaling flow, so \(m(\ell)/\ell\to\underline\delta=\sup m(\ell)/\ell\le C_N\)
      (Fekete).  Correct.
\item \emph{Import (E9) and alternative (A)} (V9).  The self-similar
      element produced by alternative (A) is a Leray profile in
      \(L^q\), \(q\in(3,\infty]\), by the \(L^\infty\) rate of the class; (E9)
      makes it zero, contradicting nontriviality.  The import is
      recorded as an import, not a proof.  Correct.
\end{enumerate}
\end{theorem}

\begin{proof}
The displayed computations.
\end{proof}

% =============================================================================
\section{Block 3: the statistical framework (V11--V15)}
\label{sec:c2v81-block3}
% =============================================================================

\begin{theorem}[Audit of V11--V15]
\label{thm:c2v81-block3}
The following are re-derived and found correct as stated.
\begin{enumerate}[label=\textup{(B3.\arabic*)},leftmargin=4.2em]
\item \emph{The rescaling flow.}  \(\Theta_\tau W:=W^{e^\tau}\) (the rescaling
      by \(e^\tau\)) maps the family to itself, is continuous in the
      metric of the family (which is that of local strong convergence
      on compact space-time sets, dominated by the class bounds), and
      corresponds to the shift \(\sigma\mapsto\sigma+2\tau\) of the profile;
      the family is compact (first cycle, V83); Birkhoff's theorem
      gives minimal sets, and Krylov--Bogolyubov gives invariant
      probability measures \(\mathfrak m\) on the family.  Correct.
\item \emph{Exact mean identities} (V12).  For a bounded continuous
      functional \(f\) on the family, invariance gives
      \(\int(f\circ\Theta_\tau-f)\,\dd\mathfrak m=0\), and for \(f\) absolutely
      continuous along the flow with bounded derivative,
      \(\int\frac{\dd}{\dd\sigma}f\,\dd\mathfrak m=0\).  Applied to
      \(f=\varphi=\|V\|_G^2\) with the first identity
      \(\frac{\dd\varphi}{\dd\sigma}=-\frac12(d+2\varphi)-\frac12j=-4\mathcal L-\frac12j\)
      (from \(-2\nu\|\nabla V\|^2-\varphi-2\langle V,N\rangle\), \(d=4\nu\|\nabla V\|_G^2\),
      \(j=4\langle V,N\rangle_G\), \(8\mathcal L=d+2\varphi\)): \(\int j\,\dd\mathfrak m=-8\int\mathcal L\,\dd\mathfrak m\).
      Applied to \(f=\mathcal L\) with the second identity:
      \(\int\langle N,Y\rangle_G\,\dd\mathfrak m=-\int\|Y\|_G^2\,\dd\mathfrak m\le-\underline\delta\)
      (the last by the definition of \(\underline\delta\) and the ergodic
      average).  Applied to \(f=\frac12\|Y\|_G^2\) with the defect energy
      identity: \(\int\langle Y,N'(V)Y\rangle_G\,\dd\mathfrak m=-\int(\nu\|\nabla Y\|_G^2+\frac12\|Y\|_G^2)\,\dd\mathfrak m\le-\frac12\underline\delta\).
      The boundedness and continuity of \(\varphi,\mathcal L,\|Y\|_G^2\) on
      the family are the class bounds and the strong convergence.
      Correct.
\item \emph{Blow-up measures} (V13).  Weak limits of the empirical
      scale-measures \(\frac1L\int_\sigma^{\sigma+L}\delta_{\Theta_{\sigma'/2}W}\,\dd\sigma'\)
      are invariant (the shift by a fixed \(\tau\) changes the empirical
      measure by at most \(2\tau/L\) in total variation), and carry the
      identities of (B3.2).  Correct.
\item \emph{Rate-free second identity} (V14).  For any smooth solution
      on the torus, \(\lambda\frac{\dd}{\dd\lambda}\mathcal L_u=2\|Y_u\|^2+2\langle N,Y_u\rangle\)
      with the periodized Gaussian weight is the second identity read
      in the original variables; \(\mathcal L_u\ge\frac14\Phi_u\) is the
      definition.  Correct.
\item \emph{Gates restated} (V15).  The missing statement at a
      homogeneous type-I point, a single mean inequality
      \(\int\langle N,\partial_\sigma V\rangle\,\dd\mathfrak m\ge-\theta\int\|\partial_\sigma V\|^2\,\dd\mathfrak m\)
      with \(\theta<1\) for one blow-up measure, is equivalent by
      (B3.2) to \(\int\|Y\|^2\,\dd\mathfrak m=0\), hence to a self-similar element
      in the support, excluded by (E9): the gate is correctly stated
      as the negation of (B).  Correct.
\end{enumerate}
M-list items (M2) and (M4)'s statistical part are marked audited.
\end{theorem}

\begin{proof}
The displayed computations.
\end{proof}

% =============================================================================
\section{Integrated V81 conclusion and next acquisition order}
\label{sec:c2v81-conclusion}
% =============================================================================

\begin{theorem}[What V81 establishes]
\label{thm:c2v81-conclusion}
Relative to the first cycle and V1--V80: the audits of V6--V10
(Theorem~\ref{thm:c2v81-block2}) and V11--V15
(Theorem~\ref{thm:c2v81-block3}), without correction.  The full
Navier--Stokes stability/global-regularity theorem remains unproved.
\end{theorem}

\begin{proof}
The cited results.
\end{proof}

\begin{programme}[V82 acquisition order]
\label{prog:c2v82}
\begin{enumerate}[label=\textup{(AC\arabic*)},leftmargin=4.8em]
\item Audit block 4, V16--V22 (the moment ledger): re-derive the
      centre-mode ODE, the momentum identity, the tensor and inertia
      identities, the Gaussian enstrophy identity, the relation
      \(\mathcal E=\|\nabla V\|_G^2+\varphi/(2\nu)-\frac1{4\nu^2}\int(z\cdot V)^2G\), and
      \(R=8\nu\mathcal L-4\nu^2\mathcal E\); check the V17 and V21 corrections.
\item Audit block 5, V23--V25 (statistical dichotomy, coherence):
      check the equivalence \(0\in\mathcal F(x_0)\Leftrightarrow\) intermediate
      scales and the moving-centre statements.
\item The next compulsory companion audit is V85.
\end{enumerate}
\end{programme}

\begin{assessment}[Position after V81]
\label{ass:c2v81-position}
Blocks 2--3 are audited without correction.  No full stability or
global-regularity claim is made.
\end{assessment}

% =============================================================================
\section*{Version V81 (second cycle) append-only record}
\addcontentsline{toc}{section}{Version V81 (second cycle) append-only record}
% =============================================================================

The complete V80 source of the second cycle was retained before this
continuation.  Its SHA--256 digest was
\begin{center}\scriptsize\ttfamily
755ef1d0c8477df09b7f085d4a07dee445f2d88339863a36864f61ee3f8bb470.
\end{center}
V81 audited V6--V15.  No full regularity claim is made.

% =============================================================================
\section{V82: second internal audit, blocks 4--5 (V16--V25)}
\label{sec:c2v82-entry}
% =============================================================================

Items (AC1) and (AC2) of Programme~\ref{prog:c2v82} are carried out.
The moment ledger (centre-mode equation, mean profile, tensor and
inertia identities, Gaussian enstrophy identity, the
enstrophy--Dirichlet relation and the closure \(R=8\nu\mathcal L-4\nu^2\mathcal E\))
and the statistical dichotomy and coherence statements are re-derived;
no correction is needed, and the V17 and V21 corrections are confirmed
as the correct forms.  No global regularity claim is made.  The next
compulsory companion audit is V85.

% =============================================================================
\section{Block 4: the moment ledger (V16--V22)}
\label{sec:c2v82-block4}
% =============================================================================

\begin{theorem}[Audit of V16--V22]
\label{thm:c2v82-block4}
The following are re-derived and found correct as stated.
\begin{enumerate}[label=\textup{(B4.\arabic*)},leftmargin=4.2em]
\item \emph{Gaussian mean of \(L_-\)} (V16).  With
      \(\Delta G=(\frac{|z|^2}{4\nu^2}-\frac3{2\nu})G\) and
      \(\operatorname{div}(zG)=(3-\frac{|z|^2}{2\nu})G\):
      \(\int\nu\Delta Y\,G-\frac12\int(z\cdot\nabla Y)G=\int Y\bigl[\nu\Delta G+\frac12\operatorname{div}(zG)\bigr]=\int Y\bigl[\frac{|z|^2}{4\nu}-\frac32+\frac32-\frac{|z|^2}{4\nu}\bigr]G=0\),
      so \(\overline{L_-Y}=-\frac12\bar Y\); the centre-mode equation
      \eqref{eq:c2v16-centre-equation} follows with
      \(\overline{(V\cdot\nabla)Y}=\frac1{2\nu|G|}\int(z\cdot V)YG\),
      \(\overline{(Y\cdot\nabla)V}=\frac1{2\nu|G|}\int(z\cdot Y)VG\) (integration
      by parts against \(\nabla G=-\frac z{2\nu}G\) with
      \(\operatorname{div}V=\operatorname{div}Y=0\)), and \(\overline{\nabla P'}=\frac1{2\nu|G|}\int P'zG\).
      Correct.
\item \emph{Mean profile} (V17).  \(\overline{N(V)}=\frac1{2\nu|G|}\int[(z\cdot V)V+P_Vz]G\)
      by the same integrations by parts, and
      \(\frac{\dd}{\dd\sigma}\bar V=-\frac12\bar V-\overline{N(V)}\); against an
      invariant measure \(\int\bar V=-2\int\overline{N(V)}\).  The V17
      correction (the dipole matrix of an odd profile need not vanish)
      is right: oddness of \(V\) makes \(\bar V=0\) but not the matrix
      \(\int V_i z_jG\), which is even.  Correct.
\item \emph{Tensor identity} (V18).  With \(T_{ij}=\langle V_i,V_j\rangle_G\),
      \(D_{ij}=2\nu\langle\nabla V_i,\nabla V_j\rangle_G\), \(S_{ij}=\langle N_i,V_j\rangle_G+\langle N_j,V_i\rangle_G\):
      \(\frac{\dd}{\dd\sigma}T_{ij}=\langle L_-V_i,V_j\rangle_G+\langle V_i,L_-V_j\rangle_G-S_{ij}=-D_{ij}-T_{ij}-S_{ij}\)
      by (B1.3) componentwise; the trace is
      \(\frac{\dd\varphi}{\dd\sigma}=-\frac12d-\varphi-\frac12j\).  Correct.
\item \emph{Inertia identity} (V18).  Linear part:
      \(2\nu\int|z|^2V\cdot\operatorname{div}(G\nabla V)=-2\nu I_\nabla-2\nu\int z\cdot\nabla|V|^2G=-2\nu I_\nabla+2\nu\int|V|^2\operatorname{div}(zG)=-2\nu I_\nabla+6\nu\varphi-I\),
      plus \(-I\) from \(-\frac12V\): \(-2\nu I_\nabla+6\nu\varphi-2I\).  Nonlinear
      part: \(\int|z|^2V\cdot(V\cdot\nabla)VG=-\int|V|^2(z\cdot V)G+\frac1{4\nu}\int|z|^2|V|^2(z\cdot V)G\)
      and \(\int|z|^2V\cdot\nabla P_VG=-2\int P_V(z\cdot V)G+\frac1{2\nu}\int|z|^2P_V(z\cdot V)G\),
      so \(-2\int|z|^2V\cdot NG=2\int(|V|^2+2P_V)(z\cdot V)G-\frac1{2\nu}\int|z|^2(|V|^2+2P_V)(z\cdot V)G=2\nu j-\frac12j_2\)
      with \(j=\frac1\nu\int(|V|^2+2P_V)(z\cdot V)G\) (the dictionary) and
      \(j_2:=\frac1\nu\int|z|^2(|V|^2+2P_V)(z\cdot V)G\).  Correct.
\item \emph{Enstrophy--Dirichlet relation} (V19).
      \(|\Omega|^2=|\nabla V|^2-\partial_iV_j\partial_jV_i\) (as \(|\nabla V|^2=|S|^2+|A|^2\),
      \(\partial_iV_j\partial_jV_i=|S|^2-|A|^2\), \(|\Omega|^2=2|A|^2\));
      \(\int\partial_iV_j\partial_jV_iG=-\int V_j\partial_jV_i\partial_iG=\frac1{2\nu}\int z_iV_j\partial_jV_iG\)
      (\(\operatorname{div}V=0\)), and
      \(\int z_iV_j\partial_jV_iG=-\int V_i\partial_j(z_iV_jG)=-\|V\|_G^2+\frac1{2\nu}\int(z\cdot V)^2G\);
      hence \(\mathcal E=\|\nabla V\|_G^2+\frac1{2\nu}\|V\|_G^2-\frac1{4\nu^2}\int(z\cdot V)^2G\).
      Correct.
\item \emph{Closure} (V21).  \(8\nu\mathcal L=4\nu^2\|\nabla V\|_G^2+2\nu\|V\|_G^2\) and
      \(4\nu^2\mathcal E=4\nu^2\|\nabla V\|_G^2+2\nu\|V\|_G^2-\int(z\cdot V)^2G\), so
      \(R=\int(z\cdot V)^2G=8\nu\mathcal L-4\nu^2\mathcal E\) and \(\mathcal E\le\frac2\nu\mathcal L\)
      with equality iff \(z\cdot V\equiv0\).  The V21 correction of the
      programme identity is this form.  The mean of
      \(\frac{\dd}{\dd\sigma}R\) vanishes identically by the second and
      enstrophy identities, so the ledger search closes without a new
      row (V22).  Correct.
\item \emph{Gaussian enstrophy identity} (V19).  From the vorticity
      form of the profile equation, multiplying by \(\Omega G\):
      \(\int\Omega\cdot(\nu\Delta\Omega-\frac12z\cdot\nabla\Omega)G=-\nu\|\nabla\Omega\|_G^2\) (by
      (B1.3) applied to \(\Omega\) without the zero-order term), the
      zero-order term of the vorticity equation being \(-\Omega\) (the
      vorticity of the profile scales with one more power), giving
      \(-\mathcal E\); the stretching \(\mathcal S=\int\Omega\cdot(\Omega\cdot\nabla)VG\)
      and the transport \(\mathcal T=\int\Omega\cdot(V\cdot\nabla)\Omega G=\frac1{4\nu}\int|\Omega|^2(z\cdot V)G\);
      hence \(\frac12\frac{\dd}{\dd\sigma}\mathcal E=-\mathcal E_\nabla-\mathcal E+\mathcal S-\mathcal T\)
      with \(\mathcal E_\nabla=\nu\|\nabla\Omega\|_G^2\).  Correct.
\end{enumerate}
M-list item (M3) is marked audited.
\end{theorem}

\begin{proof}
The displayed computations.
\end{proof}

% =============================================================================
\section{Block 5: statistical dichotomy and coherence (V23--V25)}
\label{sec:c2v82-block5}
% =============================================================================

\begin{theorem}[Audit of V23--V25]
\label{thm:c2v82-block5}
The following are checked and found correct as stated.
\begin{enumerate}[label=\textup{(B5.\arabic*)},leftmargin=4.2em]
\item \emph{Family at any singular point} (V23).  \(\mathcal F(x_0)\) is the
      set of local limits of rescalings and their rescaling limits;
      \(0\in\mathcal F(x_0)\) iff there are scales along which the
      rescalings tend to zero locally, that is, intermediate scales
      (\(\Phi_{u,x_0}\to0\) along them, by the strong local convergence
      and the Gaussian tail); \(\delta_0\) is invariant, and any invariant
      measure with zero mean squared defect is supported on
      self-similar elements, hence (E9) on \(\{0\}\).  Correct.
\item \emph{Statistical dichotomy} (V23).  (I)/(II) are exclusive and
      exhaustive; (II) is equivalent to zero logarithmic density of the
      active scales at every level by the argument with the truncated
      observable \(\min(\varphi,\varepsilon)\) (its log-average is at most
      \(\varepsilon\) plus the class bound times the density), and conversely
      positive density of active scales at some level makes some
      blow-up measure charge \(\{\varphi\ge\varepsilon/2\}\).  A homogeneous
      point has all scales active, hence (I).  Correct.
\item \emph{Compact dissipation and Gaussian activity} (V24).  The
      equivalence rests on the first cycle's moving-centre lower bound
      and its quantitative converse (Propositions V98-moving-lower and
      V98-converse of the first cycle), with the pigeonhole over the
      \(N\) centres and the slab lemma; the scale shift by \(S^{1/2}\)
      preserves logarithmic density.  Correct as a transfer.
\item \emph{Coherent centre maps} (V24).  Invariance of the
      moving-centre empirical measures needs the recentring
      translation \((x(\mu e^{\tau_0})-x(\mu))/(\mu e^{\tau_0})\to0\), which is the
      definition of coherence; without it the limits need not be
      invariant, and the series correctly records that statistically
      invisible non-compact blow-up is not excluded.  Correct.
\item \emph{V25} is the audit and consolidation; nothing to re-derive.
\end{enumerate}
M-list item (M4) is marked audited.
\end{theorem}

\begin{proof}
Inspection against (B1)--(B4).
\end{proof}

% =============================================================================
\section{Integrated V82 conclusion and next acquisition order}
\label{sec:c2v82-conclusion}
% =============================================================================

\begin{theorem}[What V82 establishes]
\label{thm:c2v82-conclusion}
Relative to the first cycle and V1--V81: the audits of V16--V22
(Theorem~\ref{thm:c2v82-block4}) and V23--V25
(Theorem~\ref{thm:c2v82-block5}), without correction.  The full
Navier--Stokes stability/global-regularity theorem remains unproved.
\end{theorem}

\begin{proof}
The cited results.
\end{proof}

\begin{programme}[V83 acquisition order]
\label{prog:c2v83}
\begin{enumerate}[label=\textup{(AC\arabic*)},leftmargin=4.8em]
\item Audit block 6, V26--V30 (the rate-free chapter): the three
      parabolic conditions and their independence, the suitable weak
      limits (E2), frequency escape and Gaussian invisibility, the
      paraboloid dissipation bound
      \(\int_0^\lambda\mathcal L_u\,\dd\lambda'/\lambda'\ge\frac\nu4e^{-1/\nu}\varepsilon'\), the bridge
      \(\|L_-V\|_G\ge2\mathcal L/\varphi^{1/2}\) and the trichotomy; re-derive the
      bridge and the paraboloid bound.
\item Audit block 7, V31--V35 (the Rayleigh chapter): the Rayleigh
      identity \(\frac{\dd r}{\dd\sigma}=-2\mathrm{Var}+2\varphi^{1/2}\langle\mathcal R,\hat N\rangle\),
      the closed identity of V32, the push budget, and
      \(\mathrm{Var}\ge v_0\); re-derive the Rayleigh identity.
\item The next compulsory companion audit is V85.
\end{enumerate}
\end{programme}

\begin{assessment}[Position after V82]
\label{ass:c2v82-position}
Blocks 4--5 are audited without correction.  No full stability or
global-regularity claim is made.
\end{assessment}

% =============================================================================
\section*{Version V82 (second cycle) append-only record}
\addcontentsline{toc}{section}{Version V82 (second cycle) append-only record}
% =============================================================================

The complete V81 source of the second cycle was retained before this
continuation.  Its SHA--256 digest was
\begin{center}\scriptsize\ttfamily
15676debc754d8057455ff9e912dd623ff02bca6191ea379271bf6aed824ddfd.
\end{center}
V82 audited V16--V25.  No full regularity claim is made.

% =============================================================================
\section{V83: second internal audit, blocks 6--7 (V26--V35)}
\label{sec:c2v83-entry}
% =============================================================================

Items (AC1) and (AC2) of Programme~\ref{prog:c2v83} are carried out.
The rate-free chapter (parabolic scales, suitable limits, frequency
escape and Gaussian invisibility, the paraboloid dissipation bound, the
bridge and the trichotomy) and the Rayleigh chapter (the normalized
profile equation, the Rayleigh identity, the closed identity, the push
budget, the uniform positivity of the spectral variance) are
re-derived; no correction is needed.  One notational point is
recorded: the normalized flux \(\widehat\jmath=j/\varphi\) of V33 equals
\(\varphi^{1/2}\widehat j\) with \(\widehat j=4\langle\widehat V,\widehat N\rangle_G\) of V29, and the
two versions are consistent.  No global regularity claim is made.  The
next compulsory companion audit is V85.

% =============================================================================
\section{Block 6: the rate-free chapter (V26--V30)}
\label{sec:c2v83-block6}
% =============================================================================

\begin{theorem}[Audit of V26--V30]
\label{thm:c2v83-block6}
The following are re-derived and found correct as stated.
\begin{enumerate}[label=\textup{(B6.\arabic*)},leftmargin=4.2em]
\item \emph{Parabolic scales and their independence} (V26).  The three
      conditions (energy, dissipation, pressure of parabolic order)
      rescale to the bounds \eqref{eq:c2v26-bounds}; the independence
      proposition uses the first cycle's rate-free local energy bound
      (Theorem V66-local-energy) and the Ladyzhenskaya interpolation,
      giving the orders \(\lambda^{1/4}\), \(\lambda^{3/4}\), \(\lambda^{5/4}\) displayed;
      the arithmetic of the exponents is checked
      (\(\lambda^{-1/2}\lambda^{3/4}=\lambda^{1/4}\); \(\lambda^{-1}\lambda^{1/2}\lambda^{3/4}\lambda^{1/2}=\lambda^{3/4}\);
      \(\lambda^{3/4}\lambda^{1/2}=\lambda^{5/4}\)).  Correct.
\item \emph{Suitable weak limits} (V26).  Aubin--Lions (E2) with
      \(\partial_sU_k\) bounded in \(L^{3/2}H^{-2}\) gives strong \(L^2\) and, by
      interpolation with \(L^\infty L^2\cap L^2L^6\), strong \(L^3\)
      convergence on compact subcylinders; the local energy inequality
      passes to the limit by lower semicontinuity of the dissipation
      and strong convergence of the cubic and pressure terms.  Correct.
\item \emph{Paraboloid dissipation} \eqref{eq:c2v27-dirichlet-sees}.
      \(\mathcal L_u(\lambda')\ge\frac18\lambda'D_u(\lambda')=\frac\nu2\lambda'\int|\nabla u|^2\Psi\), and
      \(\Psi\ge e^{-1/\nu}\) on \(|x-x_0|<2\lambda'\); the log-integral over
      \((0,\lambda]\) of \(\frac\nu2\lambda'^2\int_{B_{2\lambda'}}|\nabla u(T_*-\lambda'^2)|^2\,\dd\lambda'/\lambda'\)
      is \(\frac\nu4\int_{T_*-\lambda^2}^{T_*}\int_{B_{2\lambda'(t)}}|\nabla u|^2\) with
      \(\lambda'(t)=(T_*-t)^{1/2}\) (\(\dd t=2\lambda'\dd\lambda'\)), the paraboloid
      dissipation, at least \(\varepsilon'\) by hypothesis.  Correct, with
      the constant \(\frac\nu4e^{-1/\nu}\).
\item \emph{The bridge} (V28).  \(2\mathcal L=-\langle L_-V,V\rangle_G\le\|L_-V\|_G\varphi^{1/2}\)
      by Cauchy--Schwarz, so \(\|L_-V\|_G\ge2\mathcal L/\varphi^{1/2}\); with
      \(\|L_-V\|_G\le\|Y\|_G+\|N\|_G\) and \(\varphi\le\eta\) on the window,
      \(\int_a^b(\|Y\|+\|N\|)\dd\lambda/\lambda\ge2\eta^{-1/2}\int_a^b\mathcal L\,\dd\lambda/\lambda\ge2c_0\eta^{-1/2}\).
      The trichotomy follows by cases on the integrability of \(\|Y\|\)
      and \(\|N\|\).  Correct.  (The V28 assembly repair recorded in
      V32 is confirmed present.)
\item \emph{Normalized profile} \eqref{eq:c2v29-Vhat}--\eqref{eq:c2v29-logphi}.
      With \(\widehat V=V/\varphi^{1/2}\), \(\widehat N=N(\widehat V)=N(V)/\varphi\),
      \(r=-\langle L_-\widehat V,\widehat V\rangle_G=2\mathcal L/\varphi\):
      \(\partial_\sigma\widehat V=\varphi^{-1/2}Y-\frac12\varphi^{-3/2}(\partial_\sigma\varphi)V\) with
      \(Y=\varphi^{1/2}L_-\widehat V-\varphi\widehat N\) and
      \(\partial_\sigma\varphi=2\langle V,Y\rangle_G=-2r\varphi-2\varphi^{3/2}\langle\widehat V,\widehat N\rangle_G\)
      gives \(\partial_\sigma\widehat V=L_-\widehat V+r\widehat V-\varphi^{1/2}(\widehat N-\langle\widehat N,\widehat V\rangle_G\widehat V)\)
      and \(\frac{\dd}{\dd\sigma}\log\varphi=-2r-2\varphi^{1/2}\langle\widehat V,\widehat N\rangle_G=-2r-\frac12\varphi^{1/2}\widehat j\).
      Correct.
\end{enumerate}
M-list item (M5) is marked audited.
\end{theorem}

\begin{proof}
The displayed computations.
\end{proof}

% =============================================================================
\section{Block 7: the Rayleigh chapter (V31--V35)}
\label{sec:c2v83-block7}
% =============================================================================

\begin{theorem}[Audit of V31--V35]
\label{thm:c2v83-block7}
The following are re-derived and found correct as stated.
\begin{enumerate}[label=\textup{(B7.\arabic*)},leftmargin=4.2em]
\item \emph{Rayleigh identity} \eqref{eq:c2v31-identity}.
      \(\frac{\dd r}{\dd\sigma}=-2\langle L_-\widehat V,\partial_\sigma\widehat V\rangle_G\) by symmetry;
      inserting (B6.5): \(-2\langle L_-\widehat V,L_-\widehat V+r\widehat V\rangle_G=-2(\|L_-\widehat V\|_G^2-r^2)=-2\operatorname{Var}\)
      (and \(\operatorname{Var}=\|L_-\widehat V+r\widehat V\|_G^2=\|\mathcal R\|_G^2\) since
      \(\|L_-\widehat V+r\widehat V\|^2=\|L_-\widehat V\|^2-2r^2+r^2\)), and
      \(2\varphi^{1/2}\langle L_-\widehat V,P^\perp\widehat N\rangle_G=2\varphi^{1/2}\langle\mathcal R,\widehat N\rangle_G=\frac2\varphi\langle L_-V+rV,N\rangle_G\).
      Correct.
\item \emph{Push decomposition and closed identity} (V32).
      \(\langle L_-V+rV,N\rangle_G=\int\det(\Omega,V,L_-V)G+\frac1{2\nu}\int\Pi(z\cdot L_-V)G+\frac r4j\)
      by the Lamb--Bernoulli form (B1.5) with \(\partial_\sigma V\) replaced by
      \(L_-V\) (divergence-free by (B2.4)) and \(\langle V,N\rangle_G=j/4\);
      then \(-2\|L_-\widehat V\|^2+2r^2+\frac2\varphi\frac r4j=-2\|L_-\widehat V\|^2-r\frac{\dd}{\dd\sigma}\log\varphi\)
      since \(-r\frac{\dd}{\dd\sigma}\log\varphi=2r^2+\frac r2\widehat\jmath\) and
      \(\widehat\jmath=j/\varphi\): the closed identity \eqref{eq:c2v32-closed}.
      Correct.
\item \emph{Notation.}  V29's \(\widehat j=4\langle\widehat V,\widehat N\rangle_G=j/\varphi^{3/2}\) and
      V33's \(\widehat\jmath=j/\varphi\) satisfy \(\widehat\jmath=\varphi^{1/2}\widehat j\); both
      log-derivative formulas \(-2r-\frac12\varphi^{1/2}\widehat j\) and
      \(-2r-\frac12\widehat\jmath\) agree.  Consistent.
\item \emph{Push budget} \eqref{eq:c2v33-budget}.  Integrating the closed
      identity with \(\|L_-\widehat V\|^2=\operatorname{Var}+r^2\) and
      \(-r\frac{\dd}{\dd\sigma}\log\varphi=2r^2+\frac r2\widehat\jmath\):
      \(\int\frac{2\mathcal P}\varphi=\Delta r+2\int\operatorname{Var}+2\int r^2-2\int r^2-\int\frac r2\widehat\jmath=\Delta r+2\int\operatorname{Var}+\int\frac{r(-\widehat\jmath)}2\).
      Correct.
\item \emph{Uniform variance} (V34).  \(\operatorname{Var}(W)=0\) at a scale means
      \(\widehat V\) is an eigenfunction of \(L_-\), a Hermite-mode profile,
      which is not a Navier--Stokes profile with the class bounds
      (the argument of V34 via the profile equation); lower
      semicontinuity on the compact family and invariance under the
      flow give \(v_0>0\), non-explicit.  Correct as stated (the
      non-explicitness is recorded in the M-list).
\item \emph{V30 and V35} are audits and consolidations; the direction
      recorded at V35 (the flux bound by Gaussian Hardy) was carried
      out in V36--V37.
\end{enumerate}
M-list item (M6) is marked audited.
\end{theorem}

\begin{proof}
The displayed computations.
\end{proof}

% =============================================================================
\section{Integrated V83 conclusion and next acquisition order}
\label{sec:c2v83-conclusion}
% =============================================================================

\begin{theorem}[What V83 establishes]
\label{thm:c2v83-conclusion}
Relative to the first cycle and V1--V82: the audits of V26--V30
(Theorem~\ref{thm:c2v83-block6}) and V31--V35
(Theorem~\ref{thm:c2v83-block7}), without correction.  The full
Navier--Stokes stability/global-regularity theorem remains unproved.
\end{theorem}

\begin{proof}
The cited results.
\end{proof}

\begin{programme}[V84 acquisition order]
\label{prog:c2v84}
\begin{enumerate}[label=\textup{(AC\arabic*)},leftmargin=4.8em]
\item Audit block 8, V36--V41 (the constant-chasing chapter): the
      Gaussian Hardy inequality \(X^2\le4\nu(d+3\varphi)\), the flux bound
      \(|j|\le\kappa(d+4\varphi)+K_P(d+3\varphi)^{1/2}\) with its constants, the
      sharpened energy part \(2\kappa(\varphi d)^{1/2}\) and threshold
      \(\kappa<\sqrt2\), the modewise identities
      \(\langle n_k\rangle=-\mu_k\langle a_k\rangle\) and the \(O(1/k)\) bound; re-derive
      the Hardy inequality and the modewise identity.
\item Audit block 9, V42--V45 (the intermediate class): the window
      length \(\ell_M\) from the cubic inequality, the window
      compactness, the ledger bounds \eqref{eq:c2v42-ledgers}, the
      count of active centres, and the window flux sign for \(M<M_*\).
\item The next compulsory companion audit is V85.
\end{enumerate}
\end{programme}

\begin{assessment}[Position after V83]
\label{ass:c2v83-position}
Blocks 6--7 are audited without correction.  No full stability or
global-regularity claim is made.
\end{assessment}

% =============================================================================
\section*{Version V83 (second cycle) append-only record}
\addcontentsline{toc}{section}{Version V83 (second cycle) append-only record}
% =============================================================================

The complete V82 source of the second cycle was retained before this
continuation.  Its SHA--256 digest was
\begin{center}\scriptsize\ttfamily
baae70ac212eabb1791d7454ab02880b1c33984d814093023a0800c24567041b.
\end{center}
V83 audited V26--V35.  No full regularity claim is made.

% =============================================================================
\section{V84: second internal audit, blocks 8--9 (V36--V45)}
\label{sec:c2v84-entry}
% =============================================================================

Items (AC1) and (AC2) of Programme~\ref{prog:c2v84} are carried out.
The constant-chasing chapter (Gaussian Hardy inequality, flux bound,
sharpened energy part, modewise identities and modewise bound) and the
intermediate-class chapter (window length, compactness, ledger bounds,
count of active centres, window flux sign) are re-derived; no
correction is needed.  No global regularity claim is made.  The next
compulsory companion audit is V85.

% =============================================================================
\section{Block 8: the constant-chasing chapter (V36--V41)}
\label{sec:c2v84-block8}
% =============================================================================

\begin{theorem}[Audit of V36--V41]
\label{thm:c2v84-block8}
The following are re-derived and found correct as stated.
\begin{enumerate}[label=\textup{(B8.\arabic*)},leftmargin=4.2em]
\item \emph{Gaussian Hardy} \eqref{eq:c2v36-hardy}.  \(|z|^2G=-2\nu z\cdot\nabla G\), so
      \(X^2=2\nu\int\operatorname{div}(|V|^2z)G=2\nu\int(3|V|^2+2V\cdot(z\cdot\nabla)V)G\le6\nu\varphi+4\nu X\|\nabla V\|_G\le6\nu\varphi+\frac12X^2+8\nu^2\|\nabla V\|_G^2\),
      hence \(X^2\le12\nu\varphi+16\nu^2\|\nabla V\|_G^2=12\nu\varphi+4\nu d=4\nu(d+3\varphi)\)
      (\(d=4\nu\|\nabla V\|_G^2\)); and \(d+3\varphi=\varphi+8\mathcal L\).  Correct.
\item \emph{Flux bound} \eqref{eq:c2v36-flux-bound}.  Energy part:
      \(\frac1\nu\int|V|^2|z\cdot V|G\le\frac{C_\infty}\nu X\varphi^{1/2}\le\frac{C_\infty}\nu\cdot2\nu^{1/2}(\varphi+8\mathcal L)^{1/2}\varphi^{1/2}\le\kappa(2\varphi+8\mathcal L)=\kappa(d+4\varphi)\)
      with \(\kappa=C_\infty\nu^{-1/2}\) and \(2ab\le a^2+b^2\).  Pressure part:
      \(\frac2\nu\int|P_V||z\cdot V|G\le\frac2\nu\|P_V\|_{L^3}\||z||V|G\|_{L^{3/2}}\), with the
      Calder\'on--Zygmund bound \(\|P_V\|_3\le C_{\rm CZ}\|V\|_6^2\le C_{\rm CZ}K_6^2\)
      (\(K_6\) the class bound on \(\|V\|_{L^6}\)) and
      \(\||z||V|G\|_{3/2}\le\||z|V G^{1/2}\|_2\|G^{1/2}\|_6=X\,(\int G^3)^{1/6}\),
      \((\int e^{-3|z|^2/4\nu})^{1/6}=(4\pi\nu/3)^{1/4}\); with \(X\le2\nu^{1/2}(d+3\varphi)^{1/2}\)
      this is \(K_P(d+3\varphi)^{1/2}\), \(K_P=4\nu^{-1/2}(4\pi\nu/3)^{1/4}C_{\rm CZ}K_6^2\).
      Correct.
\item \emph{Sharpened energy part} \eqref{eq:c2v37-energy}.
      \(\int|V|^2(z\cdot V)G=2\nu\int V\cdot\nabla|V|^2G\) (by \((z\cdot V)G=-2\nu V\cdot\nabla G\)
      and \(\operatorname{div}V=0\)), bounded by \(4\nu C_\infty\|V\|_G\|\nabla V\|_G=2\nu^{1/2}C_\infty(\varphi d)^{1/2}\);
      divided by \(\nu\): \(2\kappa(\varphi d)^{1/2}\le\frac\kappa2(d+4\varphi)\).  The
      threshold: the mean identity \(\langle j\rangle=-\langle d+2\varphi\rangle\) with
      \(|j|\le2\kappa(\varphi d)^{1/2}+K_P(d+3\varphi)^{1/2}\) closes when
      \(2\kappa(\varphi d)^{1/2}<d+2\varphi\) uniformly, that is
      \(\kappa<\inf_{\varphi,d>0}(d+2\varphi)/(2(\varphi d)^{1/2})=\sqrt2\) (attained at
      \(d=2\varphi\)), with \(m'_\kappa=(3-\sqrt{1+4\kappa^2})/2\) the resulting
      margin.  Correct.
\item \emph{Modewise identities} \eqref{eq:c2v40-modewise}.  With \(V_k\) the
      \(G\)-orthogonal projection of \(V\) on the \(k\)-th
      Ornstein--Uhlenbeck eigenspace of \(L_-\) (eigenvalue \(-\mu_k\),
      \(\mu_k=\frac12+\frac k2\)), \(a_k=\|V_k\|_G^2\), \(n_k=\langle V_k,N\rangle_G\):
      \(\frac{\dd a_k}{\dd\sigma}=2\langle V_k,L_-V-N\rangle_G=-2\mu_ka_k-2n_k\); summing,
      \(\varphi=\sum a_k\) and \(2\mathcal L=\sum\mu_ka_k\) give the displayed
      forms; against an invariant measure \(\langle n_k\rangle=-\mu_k\langle a_k\rangle\).
      Correct.
\item \emph{Modewise bound} \eqref{eq:c2v41-bound}.  The advective part
      after integration by parts is
      \(-\int(V\otimes V):\nabla V_kG+\frac1{2\nu}\int(V\cdot V_k)(z\cdot V)G\); the first is
      bounded by \(C_\infty\varphi^{1/2}\|\nabla V_k\|_G\) with \(\nu\|\nabla V_k\|_G^2=\frac k2a_k\)
      (the eigenvalue relation of \(A=\nu\Delta-\frac12z\cdot\nabla\)), giving
      \(\kappa(k/2)^{1/2}(\varphi a_k)^{1/2}\); the second by
      \(\frac{C_\infty}{2\nu}a_k^{1/2}X\le\kappa a_k^{1/2}(d+3\varphi)^{1/2}\); the pressure part
      by \(\frac{K_P}4a_k^{1/2}(3+2k)^{1/2}\) via the same Calder\'on--Zygmund
      step with \(\||z|V_kG\|_{3/2}\) and the Hardy inequality for
      \(V_k\) (\(\||z|V_k\|_G^2\le2\nu(3a_k+2a_k^{1/2}\||z|V_k\|_G\|\nabla V_k\|_G)\),
      whence \(\||z|V_k\|_G^2\le4\nu(3+2k)a_k/2\cdot\) an absolute factor
      absorbed in \(K_P/4\)).  Correct up to the absolute factor, which
      the statement absorbs.
\item \emph{V38--V39.}  The lower side's obstruction (the rate-free
      bound is of order \(\rho^{1/4}\), quantitative unique continuation
      needed) and the consolidation are records, not computations.
\end{enumerate}
M-list item (M7) is marked audited.
\end{theorem}

\begin{proof}
The displayed computations.
\end{proof}

% =============================================================================
\section{Block 9: the intermediate class (V42--V45)}
\label{sec:c2v84-block9}
% =============================================================================

\begin{theorem}[Audit of V42--V45]
\label{thm:c2v84-block9}
The following are re-derived and found correct as stated.
\begin{enumerate}[label=\textup{(B9.\arabic*)},leftmargin=4.2em]
\item \emph{Window length} \eqref{eq:c2v42-window}.  From
      \(Y'\le C_A\nu^{-3}Y^3\): \((Y^{-2})'\ge-2C_A\nu^{-3}\), so
      \(Y(s)^{-2}\ge M^{-2}-2C_A\nu^{-3}(s+1)\ge\frac14M^{-2}\) for
      \(s+1\le\frac{3\nu^3}{8C_AM^2}=\ell_M\); the \(H^2\) bound integrates
      \(Y'\le-\nu\|\Delta U\|_2^2+C_A\nu^{-3}Y^3\).  Correct.
\item \emph{Compactness and ledgers} (Theorem~\ref{thm:c2v42-compactness}).
      The window bounds \(Y\le2M\), \(\int\|\Delta U\|^2\le C\) give
      \(L^\infty H^1\cap L^2H^2\) bounds, Aubin--Lions (E2) gives strong
      \(L^2(H^1)\) convergence on compact sets, and the limit is a strong
      solution; the ledger bounds are H\"older and Sobolev with the
      Gaussian weight: \(\Phi\le\|U\|_6^2\|G\|_{3/2}\le C_S^2\cdot2M\|G\|_{3/2}\),
      \(\lambda D\le4\nu Y\le8\nu M\), \(\mathcal L=\frac18(\lambda D+2\Phi)\le\nu M+\frac14C_S^2\,2M\|G\|_{3/2}\).
      Correct.
\item \emph{Few active centres} (Lemma~\ref{lem:c2v43-few}).  Disjoint
      cylinders of radius \(\lambda_j\) over the window, each with dissipation
      at least \(\varepsilon\lambda_j\), against the total \(\lambda_j\int_{J_M}Y_j\le2M\ell_M\lambda_j\)
      (the scaling of the dissipation is \(\lambda\)): \(m\le2M\ell_M/\varepsilon\).
      Correct.
\item \emph{Window flux} \eqref{eq:c2v44-window}.  \(\int_{a_j}^{b_j}D_j\,\dd\lambda=4\nu\int|\nabla u|^2\Psi\,\dd\lambda=2\nu\iint|\nabla u|^2\Psi\,\dd t/\lambda\)
      (\(\dd t=2\lambda\dd\lambda\)) and \(\Psi\ge e^{-1/(4\nu(1-\ell_M))}\) on the unit
      ball at the smallest window scale \(a_j=\lambda_j(1-\ell_M)^{1/2}\), giving
      \(2\nu e^{-1/(4\nu(1-\ell_M))}\varepsilon\); the flux integral is
      \(\Phi(b_j)-\Phi(a_j)-\int D-\int2\Phi/\lambda\le\Phi_M-\int D\).  The
      threshold \(M_*\) is the arithmetic of \(\varepsilon\le2M\ell_M\) against
      \(\varepsilon_M\).  Correct.
\item \emph{V45} is the audit and consolidation.
\end{enumerate}
M-list item (M8) is marked audited.
\end{theorem}

\begin{proof}
The displayed computations.
\end{proof}

% =============================================================================
\section{Integrated V84 conclusion and next acquisition order}
\label{sec:c2v84-conclusion}
% =============================================================================

\begin{theorem}[What V84 establishes]
\label{thm:c2v84-conclusion}
Relative to the first cycle and V1--V83: the audits of V36--V41
(Theorem~\ref{thm:c2v84-block8}) and V42--V45
(Theorem~\ref{thm:c2v84-block9}), without correction.  The full
Navier--Stokes stability/global-regularity theorem remains unproved.
\end{theorem}

\begin{proof}
The cited results.
\end{proof}

\begin{programme}[V85 acquisition order]
\label{prog:c2v85}
\begin{enumerate}[label=\textup{(AC\arabic*)},leftmargin=4.8em]
\item The compulsory companion audit: read the current SM and YM
      notes and record their digests and decision.
\item Audit block 10, V46--V50 (the spreading chapter): the cell
      accounting, the backward cone at windows, the far/local pressure
      proposition, the cellwise energy budget, the static enstrophy
      bound and the energy-active/dissipation-active theorem, the
      both-spread proposition, and the spectral face; re-derive the
      static bound and the far-pressure bound.
\item The next compulsory companion audit is V90.
\end{enumerate}
\end{programme}

\begin{assessment}[Position after V84]
\label{ass:c2v84-position}
Blocks 8--9 are audited without correction.  No full stability or
global-regularity claim is made.
\end{assessment}

% =============================================================================
\section*{Version V84 (second cycle) append-only record}
\addcontentsline{toc}{section}{Version V84 (second cycle) append-only record}
% =============================================================================

The complete V83 source of the second cycle was retained before this
continuation.  Its SHA--256 digest was
\begin{center}\scriptsize\ttfamily
7f60723920e0fc6b8d8658808347f1b48aea5fde64bfecff15dd1436c481616b.
\end{center}
V84 audited V36--V45.  No full regularity claim is made.

% =============================================================================
\section{V85: compulsory companion audit and the second internal audit,
block 10 (V46--V50)}
\label{sec:c2v85-entry}
% =============================================================================

This is the eighty-fifth version of the second cycle.  The compulsory
review of the two companion programmes is performed first.  Block 10
of the second internal audit, the spreading chapter, is then carried
out: the cell accounting, the backward cone at windows, the far and
local pressure on a cell, the cellwise energy budget, the static
enstrophy bound and the energy-active/dissipation-active theorem, the
both-spread proposition and the spectral face are re-derived; no
correction is needed.  No global regularity claim is made.  The next
compulsory companion audit is V90.

% =============================================================================
\section{Compulsory V85 review of the current companion notes}
\label{sec:c2v85-companion-audit}
% =============================================================================

The current companion files in \texttt{latex\_notes} were inspected
read-only.  The frozen checkpoints are
\begin{align}
 &\texttt{Renewal\_SM\_Einstein\_Continuum\_Research\_Notes\_V19.tex},
 \notag\\[-2pt]
 &\qquad
 \texttt{78d16412440aeb8daa6d127f1364ce9de67d6a05af9aa13e2e509b79f2a683a7},
 \label{eq:c2v85-SM-checkpoint}\\
 &\texttt{Renewal\_Yang\_Mills\_Mass\_Gap\_Research\_Notes\_V93.tex},
 \notag\\[-2pt]
 &\qquad
 \texttt{ecb8a8145a4022dc25662956fac775ccf320cb8edd1c0a04dbcececec6cca47f}.
 \label{eq:c2v85-YM-checkpoint}
\end{align}
(SM V19 has 8252 lines; YM V93 has 42180 lines.)

\emph{SM V19 (seventh series).}  After closing its finite Standard
Model Yukawa algebra (V18), the SM programme expanded the assembled
improved-Yukawa operator rather than the square of the chiral block,
proved an exact commutator estimate, and identified the finite set of
heat words that can contribute to its flat-space weight-four gate
(factorial-free joint gauge--Yukawa words).

\emph{YM V93.}  The YM programme proved a stable Duhamel contraction
in its intermediate wedge, showed that primitive root insertions
retain the reserve, proved an all-label incidence-hypertree theorem,
and identified the precise remaining cycle sector through an
incidence excess of the occurrence hypergraph.

\begin{theorem}[V85 companion-audit decision]
\label{thm:c2v85-companion-decision}
No SM V19 commutator or heat-word constant and no YM V93 contraction
or hypertree bound is imported.  Neither bears on the second cycle's
chapters or on the gates.
\end{theorem}

\begin{proof}
The SM results concern heat-kernel expansions of lattice Dirac--Yukawa
operators; the YM results concern cluster expansions of a lattice
clock.  None is a Navier--Stokes statement.
\end{proof}

% =============================================================================
\section{Block 10: the spreading chapter (V46--V50)}
\label{sec:c2v85-block10}
% =============================================================================

\begin{theorem}[Audit of V46--V50]
\label{thm:c2v85-block10}
The following are re-derived and found correct as stated.
\begin{enumerate}[label=\textup{(B10.\arabic*)},leftmargin=4.6em]
\item \emph{Cell accounting} \eqref{eq:c2v46-cells}.  The window
      dissipation \(\sum_\alpha e_\alpha=\int_{J_M}Y_j\) is at most \(2M\ell_M\) by the
      window bound and at least \(c_L^2\nu^{3/2}\ell_M\) by Leray's rate
      (\(Y_j(s)\ge c_L^2\nu^{3/2}(-s)^{-1/2}\ge c_L^2\nu^{3/2}\)); the counts are the
      pigeonhole principle.  Correct.
\item \emph{Backward cone at windows} \eqref{eq:c2v46-distance}.  From
      \(Y'\le C_1\nu^{-3}Y^3\) backward from \(t_2\): \(Y(t)^{-2}\le Y(t_2)^{-2}+2C_1\nu^{-3}(t_2-t)\le6C_1\nu^{-3}(t_2-t)\)
      once \(t_2-t\ge\nu^3/(4C_1Y(t_2)^2)\), so \(Y(t)\ge(\nu^3/(6C_1(t_2-t)))^{1/2}\),
      which implies the stated \((\nu^3/(8C_1(t_2-t)))^{1/2}\); with
      \(Y(t_j)\le M\lambda_j^{-1}\) and a spike \(Y(t_2)\ge4M\lambda_j^{-1}\) this gives
      \(t_2-t_j\ge\nu^3\lambda_j^2/(32C_1M^2)\) after the arithmetic
      (\(\nu^3/(8C_1(t_2-t_j))\le M^2\lambda_j^{-2}\), and the reach condition
      holds since \(\nu^3/(4C_1\cdot16M^2\lambda_j^{-2})\le t_2-t_j\) is implied).
      Summability of square-root gaps: consecutive spikes with
      diverging heights have gaps bounded by the cone, and the sum of
      \((t'_{k+1}-t'_k)^{1/2}\) is at most the sum of the reach lengths,
      finite since the heights diverge.  Correct as stated.
\item \emph{Far pressure at a window} (Proposition~\ref{prop:c2v47-pressure}).
      The kernel of \(\partial_i\partial_j(-\Delta)^{-1}\) is bounded by \(C|x-y|^{-3}\)
      off the diagonal; the dyadic shells give
      \(C2^{-3k}\int_{2^kQ}|U|^2\le C2^{-3k}(2^k)^{2}C_S^2Y\) by H\"older on the
      cube of side \(2^k\) (\(\|U\|_{L^2(2^kQ)}^2\le|2^kQ|^{2/3}\|U\|_6^2=4^kC_S^2Y\)),
      summable; the local part is Calder\'on--Zygmund in \(L^{3/2}\).
      (V57's Lemma~\ref{lem:c2v57-far} sharpened this for the gradient
      with the cell energies; both are correct.)  Correct.
\item \emph{Cellwise energy budget} (Theorem~\ref{thm:c2v47-budget}).
      The local energy inequality with the cutoff, the cubic term by
      the local \(L^3\) norm, the local pressure by
      \(\|P^{\rm loc}\|_{3/2}\|U\|_3\), the far pressure by its sup-norm
      times \(\|U\|_{L^1}\) (H\"older against the cubic norm and the
      cylinder volume), the diffusion cutoff by the energy; the
      interpolation of the cubic norm between \(L^\infty L^2\) and
      \(L^2H^1\) on a unit cube over the window.  Correct.
\item \emph{Static enstrophy bound} (Proposition~\ref{prop:c2v48-static}).
      With \(c\) the mean over \(B_R\): \(|c|^2|B_R|\le|B_R|^{2/3}K^2\) so
      \(|c|^2\le C'K^2/R\); Poincar\'e \(\|U-c\|_{L^2(B_R)}^2\le C_P^2R^2\delta\); then
      \(\int_{B_1}|U|^2\le2|c|^2|B_1|+2C_P^2R^2\delta<\varepsilon\) when
      \(R\ge16C'K^2/\varepsilon\) and \(\delta<\varepsilon/(4C_P^2R^2)\), contradiction.
      Correct.
\item \emph{Energy-active implies dissipation-active}
      (Theorem~\ref{thm:c2v48-energy-active}).  Strong \(L^2(B_R)\)
      convergence at the fixed time by Rellich from the \(H^1\) bound
      passes the energy lower bound to the limit; the static bound
      gives enstrophy of the limit at time \(-1\), continuous in time
      for the strong limit, hence positive dissipation on an initial
      part of the window; the uniform version by contradiction and
      compactness.  Correct.
\item \emph{Both-spread and spectral face} (V48, V50).  At a spread
      window all window limits vanish, the Gaussian ledgers tend to
      zero by the strong local convergence with the Gaussian tail, the
      dissipation is carried by a diverging number of cells; the
      spectral face is Parseval on the torus with the enstrophy fixed
      and the energy density tending to zero.  Correct.
\item \emph{V49} is the consolidation and the restart manifest
      (M1)--(M5), superseded by the M-list of V79.
\end{enumerate}
M-list item (M9)'s first half (V46--V50) is marked audited.
\end{theorem}

\begin{proof}
The displayed computations.
\end{proof}

% =============================================================================
\section{Integrated V85 conclusion and next acquisition order}
\label{sec:c2v85-conclusion}
% =============================================================================

\begin{theorem}[What V85 establishes]
\label{thm:c2v85-conclusion}
Relative to the first cycle and V1--V84: the companion-audit decision
(Theorem~\ref{thm:c2v85-companion-decision}) and the audit of V46--V50
(Theorem~\ref{thm:c2v85-block10}), without correction.  The full
Navier--Stokes stability/global-regularity theorem remains unproved.
\end{theorem}

\begin{proof}
The cited results.
\end{proof}

\begin{programme}[V86 acquisition order]
\label{prog:c2v86}
\begin{enumerate}[label=\textup{(AC\arabic*)},leftmargin=4.8em]
\item Audit block 11, V51--V55 (the window-structure chapter): the
      palinstrophy budget \(\nu\int_{J_M}\mathcal E_2\le4M\) with
      \(C_A=\frac{27}{16}C_B^4\), the enstrophy band, the sharp enstrophy
      inequality with its optimal palinstrophy, the window stretching
      bound, the influx bound and the no-travelling-windows theorem
      (with the V61 correction of the regular-cell counting); re-derive
      the sharp inequality and the influx bound.
\item Audit block 12, V56--V60 (the trap chapter's first block), in
      the light of V68: confirm the withdrawal of the V56 delay
      statements, re-derive the global inequality and the cell trap
      (V57) once more in full, and the counting and cluster statements
      of V58--V59.
\item The next compulsory companion audit is V90.
\end{enumerate}
\end{programme}

\begin{assessment}[Position after V85]
\label{ass:c2v85-position}
Block 10 is audited without correction.  No full stability or
global-regularity claim is made.
\end{assessment}

% =============================================================================
\section*{Version V85 (second cycle) append-only record}
\addcontentsline{toc}{section}{Version V85 (second cycle) append-only record}
% =============================================================================

The complete V84 source of the second cycle was retained before this
continuation.  Its SHA--256 digest was
\begin{center}\scriptsize\ttfamily
88fbe0823bb660d8c72e56bb8747fd10b2cdbeba083b9c5696545197cbd5beae.
\end{center}
V85 performed the compulsory companion audit and audited V46--V50.
No full regularity claim is made.

% =============================================================================
\section{V86: second internal audit, blocks 11--12 (V51--V60)}
\label{sec:c2v86-entry}
% =============================================================================

Items (AC1) and (AC2) of Programme~\ref{prog:c2v86} are carried out.
The window-structure chapter (palinstrophy budget, band, sharp
enstrophy inequality, window stretching bound, influx bound, no
travelling windows) and the first block of the trap chapter (the
global inequality, the cell trap and the exclusion of spread windows,
counting, cluster limits, the Morrey floor, the concentrating ball)
are re-derived, the V68 withdrawal of the V56 delay statements and the
V61 correction of the regular-cell counting are confirmed, and one
constant is sharpened in passing (the Young step of the palinstrophy
budget is exact).  No further correction is needed.  No global
regularity claim is made.  The next compulsory companion audit is V90.

% =============================================================================
\section{Block 11: the window-structure chapter (V51--V55)}
\label{sec:c2v86-block11}
% =============================================================================

\begin{theorem}[Audit of V51--V55]
\label{thm:c2v86-block11}
The following are re-derived and found correct as stated.
\begin{enumerate}[label=\textup{(B11.\arabic*)},leftmargin=4.6em]
\item \emph{Palinstrophy budget} \eqref{eq:c2v51-palinstrophy}.  Young with
      \(a=2C_BY^{3/4}\), \(b=\mathcal E_2^{3/4}\), exponents \(4\) and \(\frac43\), and the
      scaling \(\theta=(4\nu/3)^{3/4}\): \(ab\le\frac{a^4}{4\theta^4}+\frac34\theta^{4/3}b^{4/3}=\frac{27}{16}\nu^{-3}C_B^4Y^3+\nu\mathcal E_2\);
      hence \(Y'\le-\nu\mathcal E_2+C_A\nu^{-3}Y^3\) with \(C_A=\frac{27}{16}C_B^4\), and
      integrating over the window with \(Y\le2M\), \(Y(-1)\le M\):
      \(\nu\int\mathcal E_2\le M+C_A\nu^{-3}8M^3\ell_M=4M\).  Correct, and the Young
      step is sharp.
\item \emph{Band} \eqref{eq:c2v51-tail}--\eqref{eq:c2v51-band}.  Parseval:
      \(\sum_{|k|\ge K}|k|^2|\widehat U|^2\le K^{-2}\sum|k|^4|\widehat U|^2=K^{-2}\mathcal E_2\);
      average over the window with the budget; the band statement is
      Leray's floor minus the tail at \(K_M\) chosen so that the tail is
      at most half the floor.  Correct.
\item \emph{Sharp enstrophy inequality} \eqref{eq:c2v52-sharp}.
      \(f(\mathcal E)=-2\nu\mathcal E+2C_BY^{3/4}\mathcal E^{3/4}\) has \(f'=0\) at
      \(\mathcal E_*^{1/4}=\frac{3C_B}{4\nu}Y^{3/4}\) and
      \(f(\mathcal E_*)=\mathcal E_*^{3/4}(2C_BY^{3/4}-2\nu\mathcal E_*^{1/4})=\bigl(\frac{3C_B}{4\nu}\bigr)^3Y^{9/4}\cdot\frac12C_BY^{3/4}=\frac{27C_B^4}{128\nu^3}Y^3\).
      The \(\theta\)-form and the budget comparison
      (\(\frac{81}{32\cdot18}\approx0.14\)) are arithmetic.  Correct.
\item \emph{Window stretching bound} \eqref{eq:c2v53-stretching}.
      Cellwise Gagliardo--Nirenberg \(\|\nabla U\|_{L^3(Q)}^3\le C_{\rm GN}y^{3/4}h^{3/4}\),
      H\"older over cells with \(\sum y^3\le\delta^2\sum y\le8\delta^2Y\), H\"older in
      time.  Correct.
\item \emph{Influx bound} \eqref{eq:c2v54-influx} and no travelling
      windows.  The identity (with the sign display corrected in (A1)
      of V68) and the term-by-term bounds on \(H\)-regular cells, now
      with the homogeneous \(H^2\) seminorm (Remark~\ref{rem:c2v61-correction}),
      give \(y_\alpha(s_2)-y_\alpha(s_1)\le C_\flat\tau^{1/4}\); the no-travelling
      argument (a cell at level \(\eta\) stays at level \(\eta/2\) for time
      \(\tau_\eta\), contradicting vanishing window-integrated enstrophy)
      is correct; the count of non-regular cells is \(C_H/H\) with the
      corrected budget.  Correct after the V61 correction.
\item \emph{V55} is the audit and consolidation.
\end{enumerate}
M-list item (M9)'s second half is marked audited.
\end{theorem}

\begin{proof}
The displayed computations.
\end{proof}

% =============================================================================
\section{Block 12: the trap chapter, first block (V56--V60)}
\label{sec:c2v86-block12}
% =============================================================================

\begin{theorem}[Audit of V56--V60]
\label{thm:c2v86-block12}
The following are re-derived and found correct as stated, with the
V68 withdrawal confirmed.
\begin{enumerate}[label=\textup{(B12.\arabic*)},leftmargin=4.6em]
\item \emph{V56.}  Proposition~\ref{prop:c2v56-local},
      Corollary~\ref{cor:c2v56-delta} and Theorem~\ref{thm:c2v56-delay}
      remain withdrawn as stated ((A3) of V68), with
      Proposition~\ref{prop:c2v68-delay} as the repaired form.  The
      global inequality \eqref{eq:c2v56-global} is re-derived in (A8) of
      V68 and is correct.
\item \emph{Cell trap} (V57), re-derived once more in full outline.
      (a) Far pressure: the periodized kernel of \(\nabla\partial_i\partial_j(-\Delta)^{-1}\)
      is bounded by \(C|x|^{-4}+CL^{-4}\); shells give
      \(\sum_kC2^{-4k}8^{k+1}\max e=C'\max e\), the torus part
      \(CL^{-4}L^3\max e\); so \(\|\nabla P^{\rm far}\|_{L^\infty(2Q)}\le C\max_\beta e_\beta\).
      (b) Interpolation: \(y_\gamma=-\int U\cdot\Delta U\phi_\gamma-\int U\cdot(\nabla U\nabla\phi_\gamma)\le\|U\|_{L^2(2Q)}(\|\Delta U\|_{L^2(2Q)}+C\|\nabla U\|_{L^2(2Q)})\),
      squared with the covering.  (c) Cell inequality: each term of
      the local enstrophy identity bounded by local Sobolev and
      Gagliardo--Nirenberg on cubes with the Young parameter \(\eta\),
      the far-pressure terms by (a) times \((y^N)^{1/2}\) (with the (A2)
      correction), giving \eqref{eq:c2v57-cell}.  (d) Weighted sum
      with \(\kappa^{|\beta-\alpha^*|}\): the dissipation coefficient
      \(-\nu\kappa^{|\gamma-\alpha^*|}(1-500\eta\kappa^{-2})\le-\frac\nu2\kappa^{|\gamma-\alpha^*|}\) for
      \(\eta=\kappa^2/1000\); the \(b_1\)-terms give \(125\kappa^{-2}b_1W+b_1S_\kappa\varepsilon^{3/2}\).
      (e) Lower bound on \(D\) by (b), Cauchy--Schwarz and the covering:
      \(D\ge\frac{\kappa}{1458S_\kappa}\frac{W^2}\varepsilon-C_\theta^2W\).  (f) The
      capture-and-trapping ODE argument ((A7) of V68).  (g) The main
      theorem: small cell energies and \(W(-1)\le M\) give \(\delta\le1\) up to
      \(T_*\), the global inequality bounds \(Y\), contradiction with the
      singularity.  Correct.
\item \emph{Counting} \eqref{eq:c2v58-count}.  \(e_\alpha\le8^{2/3}\|U\|_{L^6(2Q_\alpha)}^2\),
      cubed and summed with the \(27\)-fold overlap: \(\sum e_\alpha^3\le1728C_S^6Y^3\);
      Chebyshev.  Correct.
\item \emph{Clusters} (Proposition~\ref{prop:c2v59-clusters}).  Bounded
      distance is an equivalence relation along a subsequence with at
      most \(N\) links; Rellich at the fixed time gives strong
      \(L^2_{\rm loc}\) convergence; the far cells' energies pass to the
      limit below the threshold.  Correct (with the ``small, not
      zero'' reading of V59).
\item \emph{Morrey floor} (Theorem~\ref{thm:c2v59-floor}).  The proof of
      (B12.2)(g) with an arbitrary time \(t\) and \(\Lambda_0=\max(\mathfrak Y(t),1)\);
      the exponent bound traces \(K\le C\Lambda_0^2\), \(p\le C\Lambda_0^{1/2}\),
      \(\Lambda\le C\Lambda_0^{5/2}\), giving \(\varepsilon_*\ge e^{-C\Lambda_0^3}\); improved in
      V71 to \(e^{-C\Lambda_0^2}\) by the sub-parabolic rescaling.  Correct.
\item \emph{Concentrating ball} (V60): the floor times \(\lambda\) and the
      H\"older--Sobolev upper bound.  Correct.
\end{enumerate}
M-list items (M10), (M11), (M12) and (M13) are marked audited.
\end{theorem}

\begin{proof}
The displayed computations and the cited audit items of V68.
\end{proof}

% =============================================================================
\section{Integrated V86 conclusion and next acquisition order}
\label{sec:c2v86-conclusion}
% =============================================================================

\begin{theorem}[What V86 establishes]
\label{thm:c2v86-conclusion}
Relative to the first cycle and V1--V85: the audits of V51--V55
(Theorem~\ref{thm:c2v86-block11}) and V56--V60
(Theorem~\ref{thm:c2v86-block12}), with the earlier corrections
confirmed and no new one.  The full Navier--Stokes
stability/global-regularity theorem remains unproved.
\end{theorem}

\begin{proof}
The cited results.
\end{proof}

\begin{programme}[V87 acquisition order]
\label{prog:c2v87}
\begin{enumerate}[label=\textup{(AC\arabic*)},leftmargin=4.8em]
\item Audit block 13, V61--V70 (the propagation chapter and the
      audit versions): the crude and far growth lemmas, the first-time
      induction with two counters, the exterior trap, the two-region
      proposition, the type-I structure, persistence and hysteresis,
      the refinement obstruction, the energetic-seeds theorem, the
      quiet-scale dichotomy (as corrected), and the full localized
      inequalities of V69; re-derive the crude lemma's exponents and
      the far/near split.
\item Audit block 14, V71--V78 (the spike chapter): the sub-parabolic
      rescaling and floor, the counts, the pre-spike floor with its
      two forms, the seeding theorem's arithmetic (\(c_g\), \(\Lambda_{\rm seed}\),
      \(\varepsilon_2\)), the cone of seeds, the one-step localization, the
      transfer of ledgers to a spike reference, the persistence
      corollary, the first-spike lemma, and the multiscale profile.
\item The next compulsory companion audit is V90.
\end{enumerate}
\end{programme}

\begin{assessment}[Position after V86]
\label{ass:c2v86-position}
Blocks 11--12 are audited; the earlier corrections stand and no new
one is needed.  No full stability or global-regularity claim is made.
\end{assessment}

% =============================================================================
\section*{Version V86 (second cycle) append-only record}
\addcontentsline{toc}{section}{Version V86 (second cycle) append-only record}
% =============================================================================

The complete V85 source of the second cycle was retained before this
continuation.  Its SHA--256 digest was
\begin{center}\scriptsize\ttfamily
0bba3dace78258526751ba3b27dbd85ebf6e01a499ac78e09f6d7e9551be21b9.
\end{center}
V86 audited V51--V60.  No full regularity claim is made.

% =============================================================================
\section{V87: second internal audit, blocks 13--14 (V61--V78)}
\label{sec:c2v87-entry}
% =============================================================================

Items (AC1) and (AC2) of Programme~\ref{prog:c2v87} are carried out.
The propagation chapter (crude and far growth lemmas, first-time
induction with two counters, exterior trap, two-region proposition,
type-I structure, persistence and hysteresis, refinement obstruction,
energetic seeds, quiet-scale dichotomy as corrected, full localized
inequalities) and the spike chapter (sub-parabolic floor and counts,
pre-spike floor, seeding arithmetic, cone of seeds, one-step
localization, ledgers at a spike reference, persistence, first-spike
lemma, multiscale profile) are re-derived; no correction is needed
beyond those already recorded in V67--V69.  No global regularity claim
is made.  The next compulsory companion audit is V90.

% =============================================================================
\section{Block 13: the propagation chapter (V61--V70)}
\label{sec:c2v87-block13}
% =============================================================================

\begin{theorem}[Audit of V61--V70]
\label{thm:c2v87-block13}
The following are re-derived and found correct as stated.
\begin{enumerate}[label=\textup{(B13.\arabic*)},leftmargin=4.6em]
\item \emph{Crude growth lemma} (Lemma~\ref{lem:c2v61-crude}).  The
      terms of the local energy identity: diffusion cutoff
      \(\le C\nu\cdot27a_M\); cubic \(\le C\|U\|_{L^\infty(2Q)}\cdot27a_M\); local
      pressure \(\le C\|U\|_{L^3(4Q)}^3\le C\|U\|_{L^\infty(4Q)}\cdot125a_M\); far
      pressure \(\le C\|\nabla P^{\rm far}\|_{L^\infty(2Q)}(27a_M)^{1/2}\) with
      \(\|\nabla P^{\rm far}\|\le C(S_4\varepsilon_{\rm th}+N_Ma_M)\).  Time integration:
      \(\|U\|_{L^\infty(4Q)}\le C\|U\|_{L^2(5Q)}^{1/4}\|U\|_{H^2(5Q)}^{3/4}\) (three
      dimensions, \(\theta=\frac34\)), and H\"older over \(\mathcal T\) with
      exponents \(\frac83,\frac85\) (\(\frac38+\frac58=1\)) gives
      \(\int_{\mathcal T}\|U\|_{H^2}^{3/4}\le(\int\|U\|_{H^2}^2)^{3/8}|\mathcal T|^{5/8}\).  Correct.
\item \emph{Far growth lemma} (Lemma~\ref{lem:c2v61-far}).  With all
      cells within distance \(2\) quiet, \(\|U\|_{L^2(5Q)}^2\le125\varepsilon_{\rm th}\), and
      the far pressure at distance \(\rho\) from the current non-quiet
      set at most \(C(S_4\varepsilon_{\rm th}+N_Ma_M(1+\rho)^{-4})\) (the quiet cells with
      \(S_4\), the at most \(N_M\) non-quiet cells at distance at least
      \(\rho\)); the smallness condition \eqref{eq:c2v61-far-small} and
      the contact radius \(\rho_0=(4C'_\flat N_Ma_M\ell_M\varepsilon_{\rm th}^{-1/2})^{1/4}\)
      make the total gain at most \(\frac38\varepsilon_{\rm th}\).  Correct.
\item \emph{First-time induction} (Theorems~\ref{thm:c2v61-propagation},
      \ref{thm:c2v66-propagation}).  The near times have measure at
      least \(\tau_c\) (else the near gain is below \(\frac38\varepsilon_{\rm th}\) and
      the total below \(\frac34\varepsilon_{\rm th}\)), hence contain a time at least
      \(\tau_c\) before the event, at which the localization was in force
      with a smaller counter; the non-regular case uses the
      right-continuous counter \(n\) and Lemma~\ref{lem:c2v66-nonregular}
      on \([-1,s^*)\).  Correct.
\item \emph{Exterior trap} (Theorem~\ref{thm:c2v62-exterior}).  The
      effective energy is time-independent after the finite
      propagation; capture with the forcing handled by
      \(V=W-\int\pi\); trapping by the ``increase only above \(K\hat\varepsilon\), at
      rate \(\pi\)'' argument; \(\tau_{\rm cap}=1/b\le\ell_M/2\) by enlarging \(b\).
      Correct.
\item \emph{Two regions, type-I structure, persistence, hysteresis,
      refinement} (V63--V65).  The envelope regions shrink (maxima drop
      at jumps); the current regions may enlarge; the type-I
      statements are the window statements at every time; the loss
      lemma adds \(4\nu M|\mathcal T|\); the hysteresis count and its budget
      are arithmetic; the refinement decay \(\varepsilon\lambda'/\lambda\) is the
      inclusion of cells.  Correct.
\item \emph{Non-regular cells and energetic seeds} (V66).  With
      \(\tilde b\ge\varepsilon_{\rm th}^{-1/2}\): pre-capture gain
      \(\le C\varepsilon_{\rm th}^{3/4}M^{3/4}\varepsilon_{\rm th}^{1/2}\), post-capture cubic rate
      \(\le C\varepsilon_{\rm th}^{3/4}(\varepsilon_{\rm th}^{1/2})^{3/4}=C\varepsilon_{\rm th}^{9/8}\), far and
      diffusion terms as before; total below \(\varepsilon_{\rm th}/4\) under
      \(\varepsilon_{\star\star}\).  The trapped level \(2\tilde K\hat\varepsilon\le C\varepsilon_{\rm th}^{1/2}\)
      for \(\rho\ge\rho_{\min}\) since \(\tilde K\sim\varepsilon_{\rm th}^{-1/2}\) and
      \(\hat\varepsilon\le C\varepsilon_{\rm th}\).  Correct.
\item \emph{Active scales and the quiet-scale dichotomy} (V66--V67).
      The Gaussian tail against the uniform cell bound places at least
      \(c/2\) of the ledger within \(R_c\); the count at level \(c'\); the
      dichotomy from the energetic-seeds theorem; the corrected
      corollary (cone versus cylinder).  Correct as corrected.
\item \emph{Full localized inequalities} (V69).  The three sums
      (dissipation coefficients, linear terms, far/near split of the
      maximum) are displayed; the near part of the maximum is at most
      \(\kappa^{3r/4-1}\cdot2M\cdot N(r/2+3)^3\le\kappa^{r/2}C_M\); the interpolation
      on the far cells uses only cells of \(\mathcal Q_3\) for \(r\ge12\); the
      conclusion handles \(W_r<2\kappa^{r/2}C_M\) by enlarging \(b\) by \(4a\).
      Correct.
\item \emph{V68, V70} are the audit and the consolidation.
\end{enumerate}
M-list items (M14), (M15), (M16) and (M17) are marked audited.
\end{theorem}

\begin{proof}
The displayed computations.
\end{proof}

% =============================================================================
\section{Block 14: the spike chapter (V71--V78)}
\label{sec:c2v87-block14}
% =============================================================================

\begin{theorem}[Audit of V71--V78]
\label{thm:c2v87-block14}
The following are re-derived and found correct as stated.
\begin{enumerate}[label=\textup{(B14.\arabic*)},leftmargin=4.6em]
\item \emph{Sub-parabolic floor} (Theorem~\ref{thm:c2v71-floor}).  At scale
      \(\lambda'=\lambda/\Lambda\): \(Y'(-1)=\lambda'\|\Lambda u\|^2=\mathfrak Y/\Lambda\le1\), torus side
      \(2\pi\Lambda/\lambda\ge8\), remaining rescaled time \(\Lambda^2\); the proviso
      \(K\varepsilon_1e^{\Lambda_{\rm tr}\Lambda^2}\le1\) with \(\varepsilon_1=(\varepsilon_0K^{-p})^{1/(1+p)}\) is met by
      \(\varepsilon_0=S_\kappa\varepsilon'_*\), \(\varepsilon'_*=S_\kappa^{-1}K^{-1}e^{-(1+p)\Lambda_{\rm tr}\Lambda^2}\):
      \(K^{1+p}\varepsilon_1^{1+p}e^{(1+p)\Lambda_{\rm tr}\Lambda^2}=K\varepsilon_0e^{(1+p)\Lambda_{\rm tr}\Lambda^2}=1\).
      The parabolic consequence divides by \(\Lambda\).  The optimality
      remark: exponent \(\Lambda^{(1-\theta)3+2\theta}=\Lambda^{3-\theta}\), minimal at
      \(\theta=1\).  Correct.
\item \emph{Counts and configuration at a general time} (V71).  The
      count at scale \(\lambda'\) uses \(Y'\le1\); the configuration theorem is
      the window theorem for the class with constant \(1\); the physical
      lengths are \(\ell_1\lambda^2/\Lambda^2\) and \(\tau_1\lambda^2/\Lambda^2\).  Correct.
\item \emph{Pre-spike floor} (Theorem~\ref{thm:c2v72-prespike}).  The
      reference time \(t_2\) is regular, the solution smooth on the
      closed interval; the trap gives \(\delta\le\Lambda_g\) on the capture
      phase and \(\delta\le1\) after; the global inequality integrates to
      \(Y(0)\le\Lambda_g\exp(C_3(\nu^{-3}\Lambda_g^2+\Lambda_g^{1/2})\tau+C_3(\nu^{-3}+1))\) with
      \(\tau=1/b\le C\nu^3\Lambda_g^{-2}\) (assuming \(b\ge\nu^{-3}\Lambda_g^2\), which
      only weakens the inequality), so the first exponent is bounded
      by an absolute multiple of \(1+\nu^3\).  The sub-parabolic form is
      the same at scale \(\lambda_g/\Lambda_g\).  Correct.
\item \emph{Seeding arithmetic} (Theorem~\ref{thm:c2v72-seeded}).
      \(\lambda(t')^2=\lambda_j^2-\lambda_g^2\le(1-c_g^2)\lambda_j^2\) and \(\lambda_g\ge c_g\lambda_j\) give
      \(\lambda_g/\lambda(t')\ge c_g/(1-c_g^2)^{1/2}\ge c_g\); so
      \(\lambda_g\|\Lambda u(t')\|_2^2=(\lambda_g/\lambda(t'))\Lambda\ge c_g\Lambda\ge Me^{C_4'}\ge\Lambda_ge^{C_4'}\)
      for \(\Lambda\ge\Lambda_{\rm seed}\); the seed cell's level \(c_0c_g\varepsilon_*\) and
      the count at level \(\varepsilon_2\) (with the \(27^3\) covering factor
      absorbed).  Correct.
\item \emph{Cone of seeds and one-step localization} (V73).  The seed
      condition and its two cases; the smallest scale \(e^{C_4'}/Y_2\);
      the type-I approach filling the range; the one-step
      localization with \(\lambda_g(t_2)^2\ge(1-\ell)\lambda_g(t_1)^2\) and the lowered
      threshold \(\varepsilon'''\).  Correct.
\item \emph{Ledgers at a spike and persistence} (V76).  The window
      statements need only \(Y\le2\bar\Lambda\) on the window and the
      scaling identity for smooth solutions; the Lipschitz constant
      \(K=8\nu\bar\Lambda+4C_S^2\bar\Lambda\|G\|_{3/2}+K_J\); \(\psi_\lambda\ge e^{-1/(4\nu)}\) on the
      unit ball gives \(c_1\); the log-step of the geometric sequence is
      \(\frac12|\log(1-\ell/2)|\le\ell/2\) for \(\ell\le1\); \(n_0=\lfloor c_1/(K\ell)\rfloor\).
      Correct.
\item \emph{First spike} (Lemma~\ref{lem:c2v76-first-spike}).
      \(\mathfrak Y(t)\le\Lambda\) before the first time it equals \(\Lambda\), so
      \(\|\Lambda u(t)\|^2\le\Lambda/\lambda(t)\le\Lambda/\lambda(t')\) (\(\lambda\) decreasing), and
      \((t'-t)^{1/2}\le\lambda_j\le2\lambda(t')\) when \(t'-t_j\le\frac34\lambda_j^2\); hence
      \(\Lambda_g\le2\Lambda\).  Correct.
\item \emph{Multiscale profile} (V77).  Count \(1728C_S^6\Lambda_\rho^3\varepsilon^{-3}\)
      with \(\Lambda_\rho=\rho\Lambda/\lambda\); floor exponent
      \(\Lambda_\rho^{5/2}\Lambda_\rho^{1/2}(\Lambda/\Lambda_\rho)^2=\Lambda^2\Lambda_\rho\); upward nesting by
      inclusion.  Correct.
\item \emph{V74, V75, V78} are consolidations, the dichotomy, and the
      hypothesis ledger; nothing to re-derive.
\end{enumerate}
M-list items (M18), (M19) and (M20) are marked audited.
\end{theorem}

\begin{proof}
The displayed computations.
\end{proof}

% =============================================================================
\section{Integrated V87 conclusion and next acquisition order}
\label{sec:c2v87-conclusion}
% =============================================================================

\begin{theorem}[What V87 establishes]
\label{thm:c2v87-conclusion}
Relative to the first cycle and V1--V86: the audits of V61--V70
(Theorem~\ref{thm:c2v87-block13}) and V71--V78
(Theorem~\ref{thm:c2v87-block14}), without new correction.  The full
Navier--Stokes stability/global-regularity theorem remains unproved.
\end{theorem}

\begin{proof}
The cited results.
\end{proof}

\begin{programme}[V88 acquisition order]
\label{prog:c2v88}
\begin{enumerate}[label=\textup{(AC\arabic*)},leftmargin=4.8em]
\item The audited M-list: restate Theorem~\ref{thm:c2v79-mlist} with
      every item marked audited (blocks 1--14), the withdrawn items
      marked, and the statuses updated (the V58 localized inequalities
      are F after V69).
\item The import ledger (E1)--(E9) restated for the manifest, with
      the version of first use and the statements that depend on
      each import.
\item The next compulsory companion audit is V90.
\end{enumerate}
\end{programme}

\begin{assessment}[Position after V87]
\label{ass:c2v87-position}
All fourteen blocks of the cycle are audited; the corrections are
those of V17, V21, V61, V67, V68 and V69, and no new one was needed.
No full stability or global-regularity claim is made.
\end{assessment}

% =============================================================================
\section*{Version V87 (second cycle) append-only record}
\addcontentsline{toc}{section}{Version V87 (second cycle) append-only record}
% =============================================================================

The complete V86 source of the second cycle was retained before this
continuation.  Its SHA--256 digest was
\begin{center}\scriptsize\ttfamily
6e1403a3d15cbc12cdec09f24cde360f0fdb4c30d1f6bf8f208d5e4576c6e802.
\end{center}
V87 audited V61--V78.  No full regularity claim is made.

% =============================================================================
\section{V88: the audited M-list and the import ledger for the manifest}
\label{sec:c2v88-entry}
% =============================================================================

Items (AC1) and (AC2) of Programme~\ref{prog:c2v88} are carried out:
the M-list of V79 is restated with every item marked audited by the
second internal audit (blocks 1--14, V80--V87) and with the statuses
updated; and the import ledger (E1)--(E9) is restated for the restart
manifest, with the versions of first use and the statements that
depend on each import.  No global regularity claim is made.  The next
compulsory companion audit is V90.

% =============================================================================
\section{The audited M-list}
\label{sec:c2v88-mlist}
% =============================================================================

\begin{theorem}[Audited M-list]
\label{thm:c2v88-mlist}
Each item of Theorem~\ref{thm:c2v79-mlist} has been re-derived in the
second internal audit; the audit block and the final status are:
\begin{center}
\small
\begin{tabular}{p{0.07\textwidth}p{0.42\textwidth}p{0.12\textwidth}p{0.28\textwidth}}
\toprule
Item & Content & Audited in & Status\\
\midrule
(M1) & profile framework, dichotomy (A)/(B), (E9) & V80 (B1), V81 (B2) & F\\
(M2) & statistical framework, mean identities & V81 (B3) & F\\
(M3) & moment ledger, \(R=8\nu\mathcal L-4\nu^2\mathcal E\) & V82 (B4) & F\\
(M4) & statistical dichotomy, coherence & V82 (B5) & F (transfer of first-cycle bounds)\\
(M5) & rate-free chapter, bridge, trichotomy & V83 (B6) & F\\
(M6) & Rayleigh chapter, \(\mathrm{Var}\ge v_0\) & V83 (B7) & F (\(v_0\) non-explicit)\\
(M7) & constant-chasing gate, upper side & V84 (B8) & F upper side; lower side open\\
(M8) & intermediate class, windows & V84 (B9) & F\\
(M9) & spreading and window-structure chapters & V85 (B10), V86 (B11) & F (counting corrected in V61)\\
(M10) & cell trap, no spread windows & V86 (B12) & F ((A2) of V68 applied)\\
(M11) & delay statements & V86 (B12) & W (V56) / F (repaired, V68)\\
(M12) & Morrey floors & V86 (B12), V87 (B14) & F\\
(M13) & finite configurations, clusters & V86 (B12) & F\\
(M14) & finite propagation & V87 (B13) & F (V69 for the localized inequalities)\\
(M15) & exterior trap & V87 (B13) & F\\
(M16) & type-I structure, active scales & V87 (B13) & F (V67 corrected in place)\\
(M17) & persistence, hysteresis, refinement & V87 (B13) & F\\
(M18) & pre-spike floor, seeding & V87 (B14) & F\\
(M19) & cone of seeds, ledgers at a spike & V87 (B14) & F under (H4); C for convergence\\
(M20) & multiscale profile & V87 (B14) & F\\
\bottomrule
\end{tabular}
\end{center}
The corrections made during the cycle are: V17 (a sentence of V16),
V21 (a programme identity of V20), V61 (the regular-cell counting of
V52/V54), V67 (in place, the trap-active corollary), V68 (V56's delay
statements withdrawn; two justifications of V54 and V57 corrected;
lattice adjustment; sub-window condition), V69 (the V58 sketch's
dissipation bound reorganised), and the legacy-citation names of V66
and V76 (in place).
\end{theorem}

\begin{proof}
Theorems~\ref{thm:c2v80-audit}, \ref{thm:c2v81-block2},
\ref{thm:c2v81-block3}, \ref{thm:c2v82-block4}, \ref{thm:c2v82-block5},
\ref{thm:c2v83-block6}, \ref{thm:c2v83-block7}, \ref{thm:c2v84-block8},
\ref{thm:c2v84-block9}, \ref{thm:c2v85-block10}, \ref{thm:c2v86-block11},
\ref{thm:c2v86-block12}, \ref{thm:c2v87-block13}, \ref{thm:c2v87-block14}.
\end{proof}

% =============================================================================
\section{The import ledger}
\label{sec:c2v88-imports}
% =============================================================================

\begin{theorem}[Import ledger of the second cycle]
\label{thm:c2v88-imports}
Every result of the second cycle is proved modulo the following
external theorems, used as stated and none re-proved; the classical
periodic Sobolev, Gagliardo--Nirenberg, Ladyzhenskaya, Agmon,
Poincar\'e, Calder\'on--Zygmund and Riesz-transform inequalities,
Leray's existence theorem and its enstrophy blow-up rate, the
Ornstein--Uhlenbeck spectral theory and the Gaussian Poincar\'e
inequality, and the Krylov--Bogolyubov and Birkhoff theorems are
used without listing.
\begin{center}
\small
\begin{tabular}{p{0.06\textwidth}p{0.44\textwidth}p{0.42\textwidth}}
\toprule
& Statement & Used in (second cycle)\\
\midrule
(E1) & CKN dissipation criterion & V67 (singular points are dissipation-active), V72--V73 (regularity of \(t_2\) not needed), the type-I statements\\
(E2) & Aubin--Lions compactness & V26 (suitable limits), V42 (window compactness), V48, V59, V76\\
(E3) & Fujita--Kato small data & first cycle only (class construction)\\
(E4) & Serrin regularity in \(L^\infty_{\rm loc}L^6\) & first cycle only\\
(E5) & CKN cubic criterion with \(L^\infty\) and derivative bounds & V4 (uniform derivative bounds on the family), V7 (half-line bounds)\\
(E6) & ESS backward uniqueness and unique continuation \cite{ESS} & V4 (small defect \(\Rightarrow\) self-similar element), V10 (discrete dichotomy)\\
(E7) & quantitative form of (E5) & V59 remark, V67 (dissipation-active at every scale)\\
(E8) & forward uniqueness of mild solutions in \(C_tL^6\) & V4, V10, V24 (moving-centre statements via the first cycle)\\
(E9) & Leray profiles in \(L^q\), \(q\in(3,\infty]\), vanish (Ne\v cas--R\accent23u\v zi\v cka--\v Sver\'ak, Tsai) & V9 (alternative (A) impossible), V15, V23 (\(\delta_0\) the only zero-defect measure)\\
\bottomrule
\end{tabular}
\end{center}
The trap and propagation chapters (V56--V78) use none of (E1)--(E9)
except (E1)/(E7) in the statements about singular points; their
theorems for smooth solutions (the cell trap, the Morrey floors, the
pre-spike floor, finite propagation) are self-contained modulo the
classical inequalities and the standard continuation criterion
(bounded enstrophy up to \(T\) extends the smooth solution past \(T\)).
\end{theorem}

\begin{proof}
Inspection of the proofs; each cites its import at the point of use.
\end{proof}

% =============================================================================
\section{Integrated V88 conclusion and next acquisition order}
\label{sec:c2v88-conclusion}
% =============================================================================

\begin{theorem}[What V88 establishes]
\label{thm:c2v88-conclusion}
Relative to the first cycle and V1--V87: the audited M-list
(Theorem~\ref{thm:c2v88-mlist}) and the import ledger
(Theorem~\ref{thm:c2v88-imports}).  The full Navier--Stokes
stability/global-regularity theorem remains unproved.
\end{theorem}

\begin{proof}
The cited results.
\end{proof}

\begin{programme}[V89 acquisition order]
\label{prog:c2v89}
\begin{enumerate}[label=\textup{(AC\arabic*)},leftmargin=4.8em]
\item Manifest part 1: the context paragraph for the third cycle's
      V1 (what the programme is, the two cycles, the rules of the
      series), and the summary of the first cycle as it will appear
      (condensed from Section~\ref{sec:c2v1-summary}).
\item Manifest part 2: the summary of the second cycle from the
      audited M-list, in the verbal form required for cross-file
      citation (``Theorem C2-V57-no-spread''), with hypotheses.
\item The next compulsory companion audit is V90.
\end{enumerate}
\end{programme}

\begin{assessment}[Position after V88]
\label{ass:c2v88-position}
The M-list is audited and the import ledger recorded.  No full
stability or global-regularity claim is made.
\end{assessment}

% =============================================================================
\section*{Version V88 (second cycle) append-only record}
\addcontentsline{toc}{section}{Version V88 (second cycle) append-only record}
% =============================================================================

The complete V87 source of the second cycle was retained before this
continuation.  Its SHA--256 digest was
\begin{center}\scriptsize\ttfamily
e7b1b353b7bb10178a5c322a531c79b587b15db55295a972e42821f04830ba4c.
\end{center}
V88 recorded the audited M-list and the import ledger.  No full
regularity claim is made.

% =============================================================================
\section{V89: the restart manifest, parts 1--2 (context and summaries)}
\label{sec:c2v89-entry}
% =============================================================================

Items (AC1) and (AC2) of Programme~\ref{prog:c2v89} are carried out.
The restart manifest for the third cycle's first version is drafted
in two parts: the context paragraph and the condensed summary of the
first cycle (part 1), and the summary of the second cycle from the
audited M-list, in the verbal cross-file citation form (part 2).  The
manifest is a draft to be assembled into the new file at the
restart; nothing mathematical is added.  No global regularity claim
is made.  The next compulsory companion audit is V90.

% =============================================================================
\section{Manifest part 1: context and the first cycle}
\label{sec:c2v89-part1}
% =============================================================================

\begin{programme}[Manifest part 1: context paragraph for the third cycle's V1]
\label{prog:c2v89-context}
The third cycle's V1 opens with the following context, adapted from
Section~\ref{sec:c2v1-context}.  \emph{What is being attempted:} for a
smooth mean-zero divergence-free solution \(u\) of the Navier--Stokes
equations \eqref{eq:NS} on \([0,T_*)\times\Torus^3\) with maximal time
\(T_*\), rule out \(T_*<\infty\) for smooth periodic data
\cite{FeffermanClay}; notation \(E_*\), \(Y(t)=\|\Lambda u(t)\|_2^2\), \(\Sigma\),
the type-I enstrophy rate \eqref{eq:type-I}, the scaled enstrophy
\(\mathfrak Y(t)=(T_*-t)^{1/2}Y(t)\) and the parabolic Morrey energy
\(\mathfrak M(t)\) of the second cycle.  \emph{Rules:} append-only single
file with digests, companion audits every fifth version, cycle closure
every hundredth version with the suffix \texttt{\_f}\(n\), instructions
in files treated as content, no version claims global regularity.
\emph{Legacy files:} the first cycle
\texttt{Renewal\_NS\_Stability\_Research\_Notes\_V100\_f1.tex} (digest
\texttt{4d4f4c22\ldots}) and the second cycle
\texttt{Renewal\_NS\_Stability\_Research\_Notes\_V100\_f2.tex} (digest to
be recorded at closure); references ``Theorem V41-gate'' name labelled
statements of the first cycle's file and ``Theorem C2-V57-no-spread''
those of the second cycle's file; neither is a cross-reference within
the new file.  \emph{Imports:} the ledger of
Theorem~\ref{thm:c2v88-imports}, (E1)--(E9), reproduced as a table.
\end{programme}

\begin{programme}[Manifest part 1: condensed summary of the first cycle]
\label{prog:c2v89-cycle1}
Reproduce, condensed to their statements, the five theorems of
Section~\ref{sec:c2v1-summary} of this file: the ledger results and
the two-regime classification (Theorem~\ref{thm:c2v1-ledgers}), the
ancient class and the type-I gate (Definition~\ref{def:c2v1-class},
Theorem~\ref{thm:c2v1-ancient}, with the V93 correction), the
rate-free results and the three gates (Theorem~\ref{thm:c2v1-rate-free}),
the shapes of singular points (Theorem~\ref{thm:c2v1-shapes}), and the
Gaussian-weighted energy with its pinching, representation, density
of inward scales, active/intermediate alternation and satellite
Gaussians (Theorem~\ref{thm:c2v1-gaussian}).  The condensation keeps
each statement's hypotheses and drops the proofs and the assessments.
\end{programme}

% =============================================================================
\section{Manifest part 2: the second cycle}
\label{sec:c2v89-part2}
% =============================================================================

\begin{programme}[Manifest part 2: summary of the second cycle]
\label{prog:c2v89-cycle2}
The third cycle's V1 states the following, each as a theorem with its
hypotheses, cited verbally as indicated.
\begin{enumerate}[label=\textup{(S\arabic*)},leftmargin=3.6em]
\item \emph{Profile and statistics} (Theorems C2-V2-second,
      C2-V4-dichotomy, C2-V9-B, C2-V12-means).  At a homogeneous type-I
      point, with \(V\) the backward self-similar profile of an ancient
      limit, \(L_-=\nu\Delta-\frac12z\cdot\nabla-\frac12\) and \(G=e^{-|z|^2/4\nu}\):
      \(\partial_\sigma V=L_-V-N(V)\); \(\frac{\dd}{\dd\sigma}\mathcal L=-\|\partial_\sigma V\|_G^2-\langle N,\partial_\sigma V\rangle_G\);
      every rescaling family satisfies alternative (B) (defect rate
      \(\underline\delta>0\), anti-alignment), alternative (A) being excluded
      by (E9); every invariant measure on the family satisfies
      \(\int j=-8\int\mathcal L\), \(\int\langle N,Y\rangle=-\int\|Y\|^2\le-\underline\delta\),
      \(\int\langle Y,N'Y\rangle\le-\frac12\underline\delta\), and the moment ledger's
      rows (centre mode, momentum, tensor, inertia, enstrophy, closure
      \(R=8\nu\mathcal L-4\nu^2\mathcal E\)); the ledger search closed without an
      incompatible row.
\item \emph{Dichotomies at any singular point} (C2-V23-dichotomy,
      C2-V24-equivalence, C2-V28-dichotomy, C2-V29-spectral).  Statistical
      nontriviality iff positive log-density of active scales; compact
      dissipation iff Gaussian activity at moving centres; at
      Dirichlet-active Gaussian-invisible points the defect plus
      nonlinearity is not log-integrable; invisibility is high-mode
      escape or outward normalized flux.
\item \emph{Constant-chasing gate} (C2-V36-flux-bound,
      C2-V37-constraint, C2-V40-modewise, C2-V41-bound).
      \(|j|\le2\kappa(\varphi d)^{1/2}+K_P(d+3\varphi)^{1/2}\), \(\kappa=C_\infty\nu^{-1/2}\); for
      \(\kappa<\sqrt2\) the mean ledger closes on the upper side; the
      lower side needs a quantitative unique continuation not
      available; modewise \(\langle n_k\rangle=-\mu_k\langle a_k\rangle\), \(\langle a_k\rangle=O(1/k)\).
\item \emph{Intermediate class} (C2-V42-window, C2-V42-compactness,
      C2-V54-uniform).  Type-I windows of length \(\ell_M=3\nu^3/(8C_AM^2)\),
      compactness to strong solutions, ledger bounds; the palinstrophy
      budget \(4M\); no travelling windows.
\item \emph{Cell trap and Morrey floors} (C2-V57-no-spread,
      C2-V59-floor, C2-V71-floor, C2-V72-prespike).  For a smooth
      solution on a torus of side \(L\ge8\) with cells of side one and
      geometric weights, the weighted maxima of cell enstrophy and
      energy satisfy \(W'\le-aW^2/\varepsilon+b(W+\varepsilon)\), \(\varepsilon'\le b'(\varepsilon+W)\);
      hence: (i) the intermediate class has no spread windows,
      \(\max_\alpha e_\alpha(-1)\ge\varepsilon_*(\nu,M)\); (ii) for any singularity and
      every \(t<T_*\), \(\mathfrak M(t)\ge e^{-C(\nu)\max(\mathfrak Y(t),1)^2}\), and the
      sub-parabolic floor at scale \(\lambda/\mathfrak Y\); (iii) for any two
      times, a spike high relative to the gap is preceded by an
      energetic ball at the gap scale.
\item \emph{Finite configurations and propagation}
      (C2-V58-count, C2-V61-propagation, C2-V62-exterior,
      C2-V66-propagation).  At a window at most \(N_M=1728C_S^6(2M)^3(\varepsilon_*/2)^{-3}\)
      energetic cells; energy activity propagates from the energetic
      seeds at speed \(\rho_0/\tau_c\); the exterior is trapped after the
      capture time; the trap bounds maxima, not sums; the localized
      inequalities hold for \(r\ge12\).
\item \emph{Type-I structure and singular points} (C2-V63-typeI,
      C2-V66-active, C2-V67-dichotomy).  Under a type-I rate the
      configuration, propagation and trap hold at every time; active
      scales are energetic cells; homogeneous points number at most
      \(N_M(c')\); at every scale a singular point is trap-active or has
      a trapped neighbourhood; active scales of distinct points need
      not align; Morrey levels decay linearly under refinement.
\item \emph{Spikes} (C2-V72-seeded, C2-V73-cone, C2-V74-sufficient,
      C2-V76-first-spike, C2-V77-profile).  Every high spike after a
      window is seeded in the window's configuration; the cone of
      seeds; convergence under a single-cluster or unique-seed
      condition; the first spike of a given height is approached
      type-I relative to itself; the multiscale profile.
\item \emph{Withdrawn and corrected.}  C2-V56-delay (withdrawn, repaired
      as C2-V68-delay); the corrections of V17, V21, V61, V67, V69.
\item \emph{The gates after the second cycle}
      (Assessment~\ref{ass:c2v79-gates}): G1 with the finite-configuration
      description; G2 with windows as finite configurations, spikes
      seeded, the between-window stretch open; G3 with the degenerate
      floors.
\end{enumerate}
\end{programme}

% =============================================================================
\section{Integrated V89 conclusion and next acquisition order}
\label{sec:c2v89-conclusion}
% =============================================================================

\begin{theorem}[What V89 establishes]
\label{thm:c2v89-conclusion}
Relative to the first cycle and V1--V88: the manifest parts 1--2
(Programmes~\ref{prog:c2v89-context}, \ref{prog:c2v89-cycle1},
\ref{prog:c2v89-cycle2}).  The full Navier--Stokes
stability/global-regularity theorem remains unproved.
\end{theorem}

\begin{proof}
The manifest is a record.
\end{proof}

\begin{programme}[V90 acquisition order]
\label{prog:c2v90}
\begin{enumerate}[label=\textup{(AC\arabic*)},leftmargin=4.8em]
\item The compulsory companion audit: read the current SM and YM
      notes and record their digests and decision.
\item Record the position after ninety versions; manifest part 3: the
      open gates and the first acquisitions proposed for the third
      cycle, ordered by the criterion of V70 (a first-cycle tool
      combined with a trap-chapter statement yielding a new signed
      row or exclusion).
\end{enumerate}
\end{programme}

\begin{assessment}[Position after V89]
\label{ass:c2v89-position}
The manifest's context and summaries are drafted.  No full stability
or global-regularity claim is made.
\end{assessment}

% =============================================================================
\section*{Version V89 (second cycle) append-only record}
\addcontentsline{toc}{section}{Version V89 (second cycle) append-only record}
% =============================================================================

The complete V88 source of the second cycle was retained before this
continuation.  Its SHA--256 digest was
\begin{center}\scriptsize\ttfamily
c3466ab4bb2b58f2efa1991f8655c2c75551a6831441bc5b4a63b38253f61026.
\end{center}
V89 drafted manifest parts 1--2.  No full regularity claim is made.

% =============================================================================
\section{V90: compulsory companion audit, the position after ninety versions,
and manifest part 3}
\label{sec:c2v90-entry}
% =============================================================================

This is the ninetieth version of the second cycle.  The compulsory
review of the two companion programmes is performed first.  The
position after ninety versions is recorded, and the third part of the
restart manifest is drafted: the open gates as they stand and the
first acquisitions proposed for the third cycle, ordered by the
criterion of V70.  No global regularity claim is made.  The next
compulsory companion audit is V95.

% =============================================================================
\section{Compulsory V90 review of the current companion notes}
\label{sec:c2v90-companion-audit}
% =============================================================================

The current companion files in \texttt{latex\_notes} were inspected
read-only.  The frozen checkpoints are
\begin{align}
 &\texttt{Renewal\_SM\_Einstein\_Continuum\_Research\_Notes\_V20.tex},
 \notag\\[-2pt]
 &\qquad
 \texttt{4737f516e6586f92c8651784bac453085027c134f04a70d5f4d53a02a489cf71},
 \label{eq:c2v90-SM-checkpoint}\\
 &\texttt{Renewal\_Yang\_Mills\_Mass\_Gap\_Research\_Notes\_V94.tex},
 \notag\\[-2pt]
 &\qquad
 \texttt{111e389c2fe4ca66d3688b642dc6f3c5ec5b18c70f007fd7bad580258654e345}.
 \label{eq:c2v90-YM-checkpoint}
\end{align}
(SM V20 has 8656 lines; YM V94 has 42562 lines.)

\emph{SM V20 (seventh series).}  The SM programme performed its own
companion audit and closed the simplest surviving joint
gauge--Yukawa word: an exact constant-Higgs quartic transfer computed
on the constant broken-phase free symbol, where its V18 block gives an
exact orthogonal sum of massive-overlap operators, with the
lower-order subtraction kept at the lattice level until the quartic
coefficient is isolated.

\emph{YM V94.}  The YM programme proved a two-root union--intersection
inequality, showed that local overlaps form a small residual gas,
treated the locally overlapped fundamental cycle, and defined the
remaining transversal-overlap contraction card.

\begin{theorem}[V90 companion-audit decision]
\label{thm:c2v90-companion-decision}
No SM V20 quartic-transfer constant and no YM V94 overlap or
contraction bound is imported.  Neither bears on the second cycle's
chapters, the audit, or the manifest.
\end{theorem}

\begin{proof}
The SM results concern heat-kernel coefficients of lattice
Dirac--Yukawa operators; the YM results concern cluster expansions of
a lattice clock.  None is a Navier--Stokes statement.
\end{proof}

% =============================================================================
\section{Position after ninety versions}
\label{sec:c2v90-position}
% =============================================================================

\begin{assessment}[Position after ninety versions of the second cycle]
\label{ass:c2v90-position-ninety}
The cycle's mathematics is written (V1--V78), consolidated (V79),
audited in full (V80--V87) with the corrections listed in
Theorem~\ref{thm:c2v88-mlist}, and summarized for the restart
(V88--V89).  The strongest results of the cycle are, for any
singularity, the Morrey floors at every time (V59, V71) and the
pre-spike floor (V72); for the intermediate class, the exclusion of
spread windows (V57), the finite configuration with finite
propagation and a trapped exterior at every window (V58--V66), and the
seeding of every high spike in the window's configuration (V72); for
type I, the same at every time (V63) with the counts of homogeneous
points (V66).  The three gates remain open.  Ten versions remain:
V90 (this), V91--V94 (manifest parts 3--4 and the new V1 assembled in
draft), V95 (audit), V96--V99 (final checks of the manifest against
the file, the closing theorem of the cycle), V100 (closure).  No
global regularity claim is made.
\end{assessment}

% =============================================================================
\section{Manifest part 3: the gates and the first acquisitions of the third cycle}
\label{sec:c2v90-part3}
% =============================================================================

\begin{programme}[Manifest part 3: the open gates]
\label{prog:c2v90-gates}
The third cycle's V1 records the three gates in the following form.
\begin{enumerate}[label=\textup{(G\arabic*)},leftmargin=3.6em]
\item \emph{Type-I Liouville.}  A nonzero element \(W\) of the ancient
      class \(\mathcal A_\nu(C)\) with the properties of Theorem V41-gate
      would have to satisfy, at every time, the finite-configuration
      description (at most \(N_C\) energetic parabolic cells, finite
      propagation, uniformly diffuse far field after a capture time;
      Theorem C2-V63-typeI, Corollary C2-V63-ancient), the
      profile identities and the invariant-measure mean identities of
      the second cycle at its homogeneous points, and the exclusion of
      self-similar elements (E9).  The missing statement is a single
      mean inequality for one blow-up measure (C2-V15), or an
      independent exclusion of finite-configuration ancient solutions.
\item \emph{Intermediate class.}  Windows are finite configurations
      (C2-V57, C2-V58, C2-V66); every high spike is seeded in the
      window's configuration (C2-V72) and, for the first spike of a
      given height, approached type-I relative to itself (C2-V76);
      the cone of seeds converges under a single-cluster or
      unique-seed condition (C2-V74).  The missing statement is the
      localization of the spike relative to its seed, equivalently the
      convergence of the cone, or any bound between the window's end
      and the spike.
\item \emph{Strong type II.}  The floors degenerate as
      \(e^{-C\mathfrak Y^2}\) and the counts grow as \(\mathfrak Y^3\); a
      singularity with vanishing parabolic energy along a sequence is
      of type II along it (C2-V59).  No constraint beyond the first
      cycle's price sheet.
\end{enumerate}
\end{programme}

\begin{programme}[Manifest part 3: first acquisitions proposed for the third cycle]
\label{prog:c2v90-acquisitions}
Ordered by the criterion of V70 (a first-cycle or second-cycle tool
combined with a trap-chapter statement that could yield a new signed
row or a new exclusion):
\begin{enumerate}[label=\textup{(T\arabic*)},leftmargin=3.6em]
\item \emph{Multiplicity of energy concentration at a scale.}  The
      number of \(D\)-separated energetic balls at the backward-parabolic
      scale of a spike is between \(1\) and \(N\); whether the energy
      inequality between two seed times, on the window at the seed
      scale (where all bounds hold), forces the number of clusters
      carrying energy of order \(\varepsilon_*\lambda\) to be one for a spike
      approached type-I relative to itself, or whether two clusters can
      each persist; a positive answer gives convergence of the cone by
      Theorem C2-V74-sufficient(a).
\item \emph{The mean inequality at homogeneous points from the finite
      configuration.}  On the support of a blow-up measure every profile
      is a finite configuration with a diffuse far field; whether the
      cross term \(\langle N,\partial_\sigma V\rangle_G\), written in the Lamb--Bernoulli
      form, has a sign on such profiles on average, using that the
      defect of a diffuse far field is small in \(L^2(G)\) outside the
      configuration.
\item \emph{A dissipation-based propagation.}  The trap chapter is
      energy-based; the first cycle's per-scale equivalence
      (terminal density \(\Leftrightarrow\) dissipation activity) suggests
      a propagation statement for dissipation-active cells, which would
      have time-integrated control and might survive beyond the window
      where the energy-based one does not.
\item \emph{The sub-parabolic configuration across time.}  At every
      time the sub-parabolic configuration has viscosity-only
      constants on a time interval that is the fraction \(\mathfrak Y^{-2}\) of
      the remaining time; whether the configurations at consecutive
      such intervals are related by finite propagation at the
      sub-parabolic scale, with the level loss of refinement
      controlled because the scale ratio between consecutive intervals
      is \(1-O(\mathfrak Y^{-2})\), not a fixed fraction.
\end{enumerate}
\end{programme}

% =============================================================================
\section{Integrated V90 conclusion and next acquisition order}
\label{sec:c2v90-conclusion}
% =============================================================================

\begin{theorem}[What V90 establishes]
\label{thm:c2v90-conclusion}
Relative to the first cycle and V1--V89: the companion-audit decision
(Theorem~\ref{thm:c2v90-companion-decision}) and manifest part 3
(Programmes~\ref{prog:c2v90-gates}, \ref{prog:c2v90-acquisitions}).
The full Navier--Stokes stability/global-regularity theorem remains
unproved.
\end{theorem}

\begin{proof}
The cited records.
\end{proof}

\begin{programme}[V91 acquisition order]
\label{prog:c2v91}
\begin{enumerate}[label=\textup{(AC\arabic*)},leftmargin=4.8em]
\item Examine (T4) of Programme~\ref{prog:c2v90-acquisitions} now,
      since it is a trap-chapter question: at time \(t\) with
      \(\Lambda=\mathfrak Y(t)\), the sub-parabolic interval is
      \([t,t+\ell_1\lambda'^2]\), \(\lambda'=\lambda/\Lambda\); at its end \(t_1\) the scaled
      enstrophy \(\mathfrak Y(t_1)\) may differ from \(\Lambda\), and the new
      sub-parabolic scale is \(\lambda(t_1)/\mathfrak Y(t_1)\); determine the ratio
      of consecutive sub-parabolic scales in terms of the enstrophy
      growth over the interval (bounded by the cubic inequality at the
      class-\(1\) constant), and whether the refinement loss per step is
      \(1-O(\Lambda^{-2})\) or worse.
\item Manifest part 4: the list of tools and scripts of the series
      (the assembly and check discipline), and the form of the new V1's
      first acquisition.
\end{enumerate}
\end{programme}

\begin{assessment}[Position after V90]
\label{ass:c2v90-position}
The gates and the first acquisitions of the third cycle are drafted.
No full stability or global-regularity claim is made.
\end{assessment}

% =============================================================================
\section*{Version V90 (second cycle) append-only record}
\addcontentsline{toc}{section}{Version V90 (second cycle) append-only record}
% =============================================================================

The complete V89 source of the second cycle was retained before this
continuation.  Its SHA--256 digest was
\begin{center}\scriptsize\ttfamily
1bc5147ee2a6a0aa53cc0a143154fc110ce9d1a5e0908e3d5cb352c501d5d455.
\end{center}
V90 performed the compulsory companion audit and drafted manifest part
3.  No full regularity claim is made.

% =============================================================================
\section{V91: the sub-parabolic chain, and manifest part 4}
\label{sec:c2v91-entry}
% =============================================================================

Items (AC1) and (AC2) of Programme~\ref{prog:c2v91} are carried out.
(AC1): the ratio of consecutive sub-parabolic scales along the chain
of sub-parabolic intervals is at least one half (the class-one cubic
inequality allows at most a doubling of the enstrophy per interval),
and close to one when the scaled enstrophy does not grow; but a step
loses a fixed factor of the Morrey level through the re-thresholding
of the seeds, and the number of steps needed to reach the singular
time (or the first spike) is of the order of the square of the scaled
enstrophy, so the chained level is exponentially small: (T4) closes
negatively, and the refinement obstruction of V65 is confirmed in this
form.  (AC2): manifest part 4 records the assembly discipline of the
series and the form of the third cycle's first acquisition.  No global
regularity claim is made.  The next compulsory companion audit is
V95.

% =============================================================================
\section{The sub-parabolic chain}
\label{sec:c2v91-chain}
% =============================================================================

\begin{proposition}[Consecutive sub-parabolic scales]
\label{prop:c2v91-ratio}
Let \(t<T_*\), \(\Lambda=\max(\mathfrak Y(t),1)\), \(\lambda=\lambda(t)\), \(\lambda'=\lambda/\Lambda\),
\(t_1:=t+\ell_1\lambda'^2\) (\(\ell_1=3\nu^3/(8C_A)\)), and \(\Lambda_1=\max(\mathfrak Y(t_1),1)\),
\(\lambda'_1=\lambda(t_1)/\Lambda_1\).  Then
\begin{equation}
 \frac{\lambda'_1}{\lambda'}\ \ge\ \frac{(1-\ell_1/\Lambda^2)^{1/2}}{2}\ \ge\ \frac12\Bigl(1-\frac{\ell_1}{\Lambda^2}\Bigr),
 \qquad
 \frac{\lambda'_1}{\lambda'}=\frac{(1-\ell_1/\Lambda^2)^{1/2}\,\Lambda}{\Lambda_1},
 \label{eq:c2v91-ratio}
\end{equation}
so the ratio is close to one when \(\Lambda_1\approx\Lambda\) and at least one half
always; if the interval is shortened to \(\theta\ell_1\lambda'^2\), the enstrophy
growth factor is at most \((1-\frac34\theta)^{-1/2}\) and the ratio at least
\((1-\theta\ell_1/\Lambda^2)^{1/2}(1-\frac34\theta)^{1/2}\).
\end{proposition}

\begin{proof}
\(\lambda(t_1)^2=\lambda^2-\ell_1\lambda^2/\Lambda^2\); in the rescaling at scale \(\lambda'\) the
enstrophy satisfies \(Y'(-1)\le1\) and, by the class-one cubic
inequality (Proposition~\ref{prop:c2v42-window} with \(M=1\)),
\(Y'(s)\le(1-2C_A\nu^{-3}(s+1))^{-1/2}\), which is \(2\) at \(s+1=\ell_1\) and
\((1-\frac34\theta)^{-1/2}\) at \(s+1=\theta\ell_1\); hence
\(\|\Lambda u(t_1)\|_2^2\le2/\lambda'=2\Lambda/\lambda\) and
\(\mathfrak Y(t_1)\le2\Lambda(1-\ell_1/\Lambda^2)^{1/2}\), so \(\Lambda_1\le2\Lambda\).
\end{proof}

\begin{theorem}[The chained level is exponentially small]
\label{thm:c2v91-chain}
Along the chain \(t_0<t_1<\dots\) of sub-parabolic intervals (with
\(\theta\le1\)), let the seeds at step \(k\) be the \(\lambda'_k\)-cells energetic at
a level \(\varepsilon^{(k)}\).  Finite propagation at scale \(\lambda'_k\) with
class-one constants localizes the seeds of step \(k+1\) within
\(R'_{\rm inv}(\nu;\varepsilon^{(k)}/(4\cdot216))\,\lambda'_k\) of the cells energetic at level
\(\varepsilon^{(k+1)}(\lambda'_{k+1}/\lambda'_k)/(4\cdot27)\) at \(t_k\), so that a chain of
seed sets nested by finite propagation requires
\begin{equation}
 \varepsilon^{(k)}\ \le\ \frac{\lambda'_{k+1}}{\lambda'_k}\cdot\frac{\varepsilon^{(k+1)}}{108}\quad\text{for all }k,
 \qquad\text{hence}\qquad
 \varepsilon^{(0)}\ \le\ 108^{-n}\frac{\lambda'_n}{\lambda'_0}\,\varepsilon^{(n)}
 \label{eq:c2v91-loss}
\end{equation}
after \(n\) steps.  The number of steps from \(t_0\) to a time \(t'\) is at
least \((t'-t_0)\Lambda_{\min}^2/(\theta\ell_1\lambda(t_0)^2)\), where \(\Lambda_{\min}\) is the
least scaled enstrophy along the chain; for the first spike of height
\(\Lambda\) after a window (\(t'-t_0\ge c_g^2\lambda_j^2\), \(\Lambda_{\min}\ge1\)) it is at least
\(c_g^2/(\theta\ell_1)\), and for the approach to \(T_*\) it is unbounded.
Consequently no fixed level survives the chain: (T4) of
Programme~\ref{prog:c2v90-acquisitions} closes negatively, in agreement
with the refinement obstruction (Proposition~\ref{prop:c2v65-refinement});
the improvement of the scale ratio to \(1-O(\Lambda^{-2})\) per step is
cancelled by the \(\Lambda^2\) steps per unit of physical time.
\end{theorem}

\begin{proof}
The seed re-thresholding is Proposition~\ref{prop:c2v74-fixed-level}(ii)
at the sub-parabolic scales, with the covering factor \(27\) and the
level factor \(4\) of the seed definition; \eqref{eq:c2v91-loss} is its
iteration, using \(\lambda'_{k+1}/\lambda'_k\le1+O(\Lambda_k^{-2})\) in the favourable
direction only to say that the scale ratio does not compensate the
factor \(108\).  The step count is the sum of the interval lengths
\(\theta\ell_1\lambda'^2_k=\theta\ell_1\lambda(t_k)^2/\Lambda_k^2\le\theta\ell_1\lambda(t_0)^2/\Lambda_{\min}^2\).
\end{proof}

% =============================================================================
\section{Manifest part 4: discipline and the first acquisition}
\label{sec:c2v91-part4}
% =============================================================================

\begin{programme}[Manifest part 4: the discipline of the series]
\label{prog:c2v91-discipline}
The third cycle's V1 records the discipline in one paragraph: each
version is drafted as a separate source ending with
\texttt{\textbackslash end\{document\}}, appended before the previous file's
\texttt{\textbackslash end\{document\}} with the title, PDF title and date retitled;
the previous file's SHA--256 digest is recorded in the version's
append-only record; unresolved references, duplicate labels,
environment balance, brace balance and citation keys are checked
before the previous file is deleted; corrections to earlier versions
are made in place and announced in the next record with the corrected
digest; legacy results are cited verbally, never by cross-reference;
companion notes are audited every fifth version and delivered every
fifth version; no version claims global regularity.  Labels of the
third cycle are prefixed \texttt{c3vNN-}.
\end{programme}

\begin{programme}[Manifest part 4: the first acquisition of the third cycle]
\label{prog:c2v91-first}
The third cycle's V1, after the context, the summaries, the imports
and the gates, makes its first acquisition on (T1) of
Programme~\ref{prog:c2v90-acquisitions}: on a window at the seed
scale with reference a spike approached type-I relative to itself,
write the local energy identities of two clusters each carrying
energy at least \(c_0\varepsilon_*\lambda\) after the capture time, with the
exterior trapped between them (so that the far pressure at each is
the other's halo plus the quiet level), and determine whether the
two can persist through the window or whether the identity forces
one to lose its energy; the outcome, positive or negative, is the
first record of the third cycle.
\end{programme}

% =============================================================================
\section{Integrated V91 conclusion and next acquisition order}
\label{sec:c2v91-conclusion}
% =============================================================================

\begin{theorem}[What V91 establishes]
\label{thm:c2v91-conclusion}
Relative to the first cycle and V1--V90: the ratio of consecutive
sub-parabolic scales \eqref{eq:c2v91-ratio}, the negative closure of
(T4) (Theorem~\ref{thm:c2v91-chain}), and manifest part 4.  The full
Navier--Stokes stability/global-regularity theorem remains unproved.
\end{theorem}

\begin{proof}
The cited results.
\end{proof}

\begin{programme}[V92 acquisition order]
\label{prog:c2v92}
\begin{enumerate}[label=\textup{(AC\arabic*)},leftmargin=4.8em]
\item Draft the third cycle's V1 source in the scratch space of the
      series (not appended): preamble copied from the present file,
      the context paragraph (Programme~\ref{prog:c2v89-context}), the
      import table, the condensed first-cycle summary
      (Programme~\ref{prog:c2v89-cycle1}), the second-cycle summary
      (Programme~\ref{prog:c2v89-cycle2}) as theorems with verbal
      citations, the gates (Programme~\ref{prog:c2v90-gates}), the
      discipline paragraph, and a placeholder for the first
      acquisition; check it with the series' checker; record its
      line count and digest in V92.
\item Record in V92 the list of statements of the second cycle whose
      constants are explicit in principle (all of the trap chapter)
      versus those with compactness constants (\(c_\Phi\), \(v_0\),
      \(\varepsilon'(\varepsilon,M,\nu)\), \(R'\)), for the manifest.
\end{enumerate}
\end{programme}

\begin{assessment}[Position after V91]
\label{ass:c2v91-position}
The sub-parabolic chain loses the level exponentially; the manifest's
discipline and first acquisition are drafted.  No full stability or
global-regularity claim is made.
\end{assessment}

% =============================================================================
\section*{Version V91 (second cycle) append-only record}
\addcontentsline{toc}{section}{Version V91 (second cycle) append-only record}
% =============================================================================

The complete V90 source of the second cycle was retained before this
continuation.  Its SHA--256 digest was
\begin{center}\scriptsize\ttfamily
b02acb15f3f600e2bf6af5589e34a2cc9e340aa40563e0260904d6a674c27c7b.
\end{center}
V91 closed (T4) negatively and drafted manifest part 4.  No full
regularity claim is made.

% =============================================================================
\section{V92: the third cycle's V1 drafted, and the constants ledger}
\label{sec:c2v92-entry}
% =============================================================================

Items (AC1) and (AC2) of Programme~\ref{prog:c2v92} are carried out.
The third cycle's V1 has been drafted as a separate source in the
scratch space of the series and checked with the series' checker; its
line count and digest are recorded here, and it will be assembled as
the new file at the closure of this cycle, with the digest of the
closed second-cycle file inserted in place of its placeholder.  The
constants ledger lists which statements of the second cycle have
constants explicit in principle and which rest on compactness.  No
global regularity claim is made.  The next compulsory companion audit
is V95.

% =============================================================================
\section{The drafted V1 of the third cycle}
\label{sec:c2v92-draft}
% =============================================================================

\begin{theorem}[Record of the drafted V1]
\label{thm:c2v92-draft}
The draft of the third cycle's V1 consists of: the preamble of the
present file with the title retitled; an abstract; the context and
rules (Section~1, with the placeholder \texttt{C2DIGEST} for the
digest of the closed second-cycle file); the import table (E1)--(E9)
(Section~2); the audited summary of the first cycle in five theorems
(Section~3, ledgers, ancient limit and gate, rate-free results and
gates, shapes, Gaussian ledger); the audited summary of the second
cycle in six theorems and a remark (Section~4, profile and
statistics (S1)--(S2), constant-chasing (S3), windows (S4), trap and
floors (S5), configurations and propagation (S6), type-I structure
(S7), spikes (S8), withdrawn and corrected statements with the
constants note); the three gates (Section~5); the first acquisition
programmed on two clusters on a window (Section~6); the V1 record
and the bibliography.  Its checker output is: no missing references,
no duplicate labels, environments and braces balanced, citation keys
present; it has \(548\) lines, and its SHA--256 digest with the
placeholder in place is
\begin{center}\scriptsize\ttfamily
9c7de8c64137cceaefbac34780a2a1b29802ded8442cfbda44334b5a39bfa746.
\end{center}
At closure the placeholder is replaced by the digest of
\texttt{Renewal\_NS\_Stability\_Research\_Notes\_V100\_f2.tex}, the file
is saved as \texttt{Renewal\_NS\_Stability\_Research\_Notes\_V1.tex}, and
its final digest is recorded in the third cycle's V2.
\end{theorem}

\begin{proof}
The record of the draft's construction and check.
\end{proof}

% =============================================================================
\section{The constants ledger}
\label{sec:c2v92-constants}
% =============================================================================

\begin{theorem}[Explicit and compactness constants of the second cycle]
\label{thm:c2v92-constants}
\begin{enumerate}[label=\textup{(\roman*)},leftmargin=3.2em]
\item \emph{Explicit in principle} (absolute constants of the periodic
      Sobolev, Gagliardo--Nirenberg, Poincar\'e and Calder\'on--Zygmund
      inequalities, the kernel bound of the far pressure, the Young
      parameters, and the displayed dependence on \(\nu\), \(M\), the
      threshold): the window length \(\ell_M\) and the palinstrophy
      budget; the Gaussian Hardy and flux bounds \(\kappa,K_P\) and the
      modewise bounds; the cell trap's \(a,b,b',K,p,\Lambda_{\rm tr}\) and hence
      \(\varepsilon_*(\nu,\Lambda)\), \(\varepsilon'_*(\nu,\Lambda)\), the Morrey floors and the
      pre-spike floor; the counts \(N_M(\varepsilon)\); the propagation constants
      \(\tau_c,\rho_0,R_{\rm inv},R'_{\rm inv},\rho_{\min},\tau_{\rm cap},C_1\); the
      seeding constants \(c_g\), \(\Lambda_{\rm seed}\), \(\varepsilon_2\) (given the first
      cycle's cone constant \(C_1\), itself the constant of the cubic
      enstrophy inequality); the ledger Lipschitz constant \(K\) at a
      spike; the multiscale profile.
\item \emph{From compactness} (existence proved, value not
      computable from the proof): the pinching constants \(c_\Phi,c_\Phi'\)
      of the first cycle; the variance floor \(v_0\) (C2-V34); the
      dissipation level \(\varepsilon'(\varepsilon,M,\nu)\) and radius \(R'(\varepsilon,M,\nu)\)
      of energy-active centres (C2-V48-energy-active), hence the
      per-cluster dissipation used in the third cycle's first
      acquisition; the constants of the window limits' nontriviality;
      the defect rate \(\underline\delta\) and \(m(\ell)\) (defined by infima on
      the family).
\item \emph{Imported} (external theorems): \(\varepsilon_{\rm CKN}\) of (E1)/(E7),
      the constants of (E5) and (E8).
\end{enumerate}
The third cycle's first acquisition mixes (i) (trap, halo, quiet
level) with one constant of (ii) (the per-cluster dissipation); if the
outcome is a bound on the number of windows over which two clusters
coexist, that bound is explicit in \(\varepsilon'\) and \(R'\) but not in
\((\nu,M)\) alone.
\end{theorem}

\begin{proof}
Inspection of the statements and of the audit records (V80--V87).
\end{proof}

% =============================================================================
\section{Integrated V92 conclusion and next acquisition order}
\label{sec:c2v92-conclusion}
% =============================================================================

\begin{theorem}[What V92 establishes]
\label{thm:c2v92-conclusion}
Relative to the first cycle and V1--V91: the record of the drafted V1
(Theorem~\ref{thm:c2v92-draft}) and the constants ledger
(Theorem~\ref{thm:c2v92-constants}).  The full Navier--Stokes
stability/global-regularity theorem remains unproved.
\end{theorem}

\begin{proof}
The cited records.
\end{proof}

\begin{programme}[V93 acquisition order]
\label{prog:c2v93}
\begin{enumerate}[label=\textup{(AC\arabic*)},leftmargin=4.8em]
\item Cross-check the drafted V1 against the present file: every
      verbal citation ``Theorem C2-VNN-name'' in the draft must
      correspond to a label \texttt{thm:c2vNN-name} (or
      \texttt{prop:}, \texttt{cor:}, \texttt{lem:}) present in this file;
      list the correspondences and correct the draft where a name
      differs; record the corrected draft's digest.
\item Cross-check the draft's first-cycle citations (``Theorem
      V41-gate'' and the others) against the legacy file's labels in
      the same way.
\end{enumerate}
\end{programme}

\begin{assessment}[Position after V92]
\label{ass:c2v92-position}
The third cycle's V1 is drafted and checked; the constants ledger is
recorded.  No full stability or global-regularity claim is made.
\end{assessment}

% =============================================================================
\section*{Version V92 (second cycle) append-only record}
\addcontentsline{toc}{section}{Version V92 (second cycle) append-only record}
% =============================================================================

The complete V91 source of the second cycle was retained before this
continuation.  Its SHA--256 digest was
\begin{center}\scriptsize\ttfamily
536ab24ade9528d4bbc210dc725367474cceeec0e4a8bb6dcee695a6a057f4fd.
\end{center}
(The V91 file was corrected in place after its assembly: two verbatim
occurrences of the closing-line text in the discipline paragraph were
replaced by typewriter text so that the series' checker counts the
document environment correctly; the digest above is that of the
corrected file.)  V92 recorded the drafted V1 and the constants
ledger.  No full regularity claim is made.

% =============================================================================
\section{V93: cross-check of the drafted V1 against the two files}
\label{sec:c2v93-entry}
% =============================================================================

Items (AC1) and (AC2) of Programme~\ref{prog:c2v93} are carried out.
Every verbal citation of the drafted third-cycle V1 was matched by a
script against the labels of the present file (second-cycle
citations) and of the legacy file (first-cycle citations): of the
thirty-six second-cycle citations one name was wrong and was
corrected in the draft; all eleven first-cycle citations resolve; the
corrected draft's digest is recorded.  No global regularity claim is
made.  The next compulsory companion audit is V95.

% =============================================================================
\section{The cross-check}
\label{sec:c2v93-xcheck}
% =============================================================================

\begin{theorem}[Citation correspondences of the drafted V1]
\label{thm:c2v93-xcheck}
\begin{enumerate}[label=\textup{(\roman*)},leftmargin=3.2em]
\item \emph{Second cycle.}  The draft cites, in the form
      ``C2-VNN-name'', the following labelled statements of this file,
      each present with the prefix \texttt{thm:}, \texttt{prop:},
      \texttt{cor:} or \texttt{lem:} and the suffix \texttt{c2vNN-name}:
      V2-second, V9-B, V12-mean-identities, V23-dichotomy,
      V24-equivalence, V28-divergence, V29-spectral, V37-constraint,
      V48-energy-active, V54-uniform, V57-enstrophy, V57-energy,
      V57-trap, V57-no-spread, V58-count, V59-clusters, V59-floor,
      V61-propagation, V62-exterior, V63-typeI, V63-ancient,
      V65-refinement, V66-propagation, V66-active, V67-dichotomy,
      V68-delay, V71-floor, V72-prespike, V72-seeded, V73-cone,
      V73-one-step, V74-sufficient, V76-first-spike, V76-persistence,
      V77-profile, V91-chain.  The draft had ``C2-V12-means'' for the
      mean identities, whose label is \texttt{thm:c2v12-mean-identities};
      the draft was corrected.  Citations by version number alone
      (C2-V22, C2-V31, C2-V34, C2-V61) are verbal and need no label.
\item \emph{First cycle.}  The draft cites the following labelled
      statements of the legacy file, each present: V41-gate,
      V66-local-energy, V70-gates, V81-count, V81-equivalence,
      V87-derivative-u, V91-pinched, V93-G1, V97-log-period, V98-count;
      the phrase ``Theorems V77--V79'' is a range, not a citation.
\item \emph{Corrected draft.}  After the correction the draft has \(548\)
      lines, passes the series' checker, and has SHA--256 digest
      \begin{center}\scriptsize\ttfamily
      cbe3817dfebe45f24c7e30b8b407cab0c1337cef4fa4fd53872e558b68676c0c
      \end{center}
      with the placeholder \texttt{C2DIGEST} still in place.
\end{enumerate}
\end{theorem}

\begin{proof}
The script's output, recorded at the time of the check.
\end{proof}

% =============================================================================
\section{Integrated V93 conclusion and next acquisition order}
\label{sec:c2v93-conclusion}
% =============================================================================

\begin{theorem}[What V93 establishes]
\label{thm:c2v93-conclusion}
Relative to the first cycle and V1--V92: the citation cross-check
(Theorem~\ref{thm:c2v93-xcheck}).  The full Navier--Stokes
stability/global-regularity theorem remains unproved.
\end{theorem}

\begin{proof}
The cited record.
\end{proof}

\begin{programme}[V94 acquisition order]
\label{prog:c2v94}
\begin{enumerate}[label=\textup{(AC\arabic*)},leftmargin=4.8em]
\item The closing theorem of the second cycle: a single statement,
      with proof by citation, of what the cycle proved (the audited
      M-list in one paragraph), what it withdrew, and what it leaves
      open, in the form the first cycle used at its V100.
\item The closure procedure for V100: the present file is retained
      as \texttt{Renewal\_NS\_Stability\_Research\_Notes\_V100\_f2.tex}
      after the V100 append; its digest is computed; the draft V1 has
      its placeholder replaced by that digest and is saved as
      \texttt{Renewal\_NS\_Stability\_Research\_Notes\_V1.tex}; the
      third cycle's V2 records the digest of the new V1.
\end{enumerate}
\end{programme}

\begin{assessment}[Position after V93]
\label{ass:c2v93-position}
The drafted V1 cites correctly.  No full stability or
global-regularity claim is made.
\end{assessment}

% =============================================================================
\section*{Version V93 (second cycle) append-only record}
\addcontentsline{toc}{section}{Version V93 (second cycle) append-only record}
% =============================================================================

The complete V92 source of the second cycle was retained before this
continuation.  Its SHA--256 digest was
\begin{center}\scriptsize\ttfamily
3bec6db7df0a6f86c2c95037da686d40cc055161b5d6c39a10b6a74656af5624.
\end{center}
V93 cross-checked the drafted V1.  No full regularity claim is made.

% =============================================================================
\section{V94: the closing theorem of the second cycle, and the closure procedure}
\label{sec:c2v94-entry}
% =============================================================================

Items (AC1) and (AC2) of Programme~\ref{prog:c2v94} are carried out:
the closing theorem of the second cycle, in the form the first cycle
used at its V100, and the closure procedure for V100.  No global
regularity claim is made.  The next compulsory companion audit is V95.

% =============================================================================
\section{The closing theorem}
\label{sec:c2v94-closing}
% =============================================================================

\begin{theorem}[What the second cycle proved]
\label{thm:c2v94-closing}
Let \(u\) be a smooth mean-zero solution of the Navier--Stokes
equations on \([0,T_*)\times\Torus^3\) with maximal time \(T_*<\infty\).
Modulo the imports (E1)--(E9) of Theorem~\ref{thm:c2v88-imports}, the
second cycle proved the following, and nothing beyond it.
\begin{enumerate}[label=\textup{(C\arabic*)},leftmargin=3.8em]
\item \emph{Profile framework and statistics} (V1--V25).  At a
      homogeneous type-I point, the backward self-similar profile
      satisfies \(\partial_\sigma V=L_-V-N(V)\) with the \(G\)-symmetric drift
      operator \(L_-\) of spectrum \(\{-\frac12-\frac k2\}\); the second
      identity \(\frac{\dd}{\dd\sigma}\mathcal L=-\|\partial_\sigma V\|_G^2-\langle N,\partial_\sigma V\rangle_G\)
      with its rigid equality case and Lamb--Bernoulli cross term;
      the dichotomy (A)/(B) with (A) excluded by (E9), so every
      rescaling family is uniformly non-self-similar with the three
      averaged signs; the rescaling flow on the compact family, its
      invariant and blow-up measures, and the exact mean identities;
      the moment ledger (centre mode, momentum, tensor, inertia,
      enstrophy, closure \(R=8\nu\mathcal L-4\nu^2\mathcal E\)), searched and
      closed without an incompatible row; the statistical dichotomy
      at any singular point and the equivalence of compact
      dissipation with Gaussian activity at moving centres.
\item \emph{Rate-free, Rayleigh and constant-chasing chapters}
      (V26--V41).  Parabolic scales and suitable limits; frequency
      escape and Gaussian invisibility with the paraboloid dissipation
      seen by the Dirichlet ledger; the bridge and the trichotomy; the
      Rayleigh identity, the closed identity, the push budget and the
      uniform variance floor; the Gaussian Hardy inequality, the flux
      bound with the threshold \(\kappa<\sqrt2\), the explicit upper side of
      the constant-chasing gate, and the modewise identities and
      bounds; the lower side recorded as needing a quantitative unique
      continuation.
\item \emph{The intermediate class at its windows} (V42--V55).  Type-I
      windows with compactness and explicit ledger bounds; few active
      centres; the backward cone at windows; the far and local pressure
      on cells; energy-active implies dissipation-active; the
      palinstrophy budget, band and sharp enstrophy inequality; the
      window stretching and influx bounds; no travelling windows.
\item \emph{The cell trap} (V56--V60).  The far-pressure gradient
      bounded by the maximal cell energy, the enstrophy--energy
      interpolation, the weighted local inequalities, capture and
      trapping; \emph{the intermediate class has no spread windows};
      the Morrey floor at every time before any singularity; finite
      configurations and cluster limits; the V56 delay statements
      withdrawn and repaired (V68).
\item \emph{Propagation} (V61--V70).  Finite propagation of energy
      activity on a window from the energetic seeds; the unconditional
      exterior trap after the capture time; the type-I structure at
      every time and the diffuse far field of the ancient limit;
      active scales as energetic cells and the count of homogeneous
      points; the quiet-scale dichotomy; persistence and bounded
      oscillation; the negative records (maxima not sums, hysteresis
      budget, refinement decay, non-alignment of active scales); the
      full localized inequalities.
\item \emph{Spikes} (V71--V78).  The sub-parabolic floor
      \(e^{-C(\nu)\mathfrak Y^2}\) and configuration with viscosity-only
      constants; the pre-spike floor for any two times; every high
      spike after a window seeded in the window's configuration; the
      backward cone of seeds with one-step localization and hopping;
      the two sufficient conditions for convergence; the ledgers with a
      spike as reference, persistence, the first-spike lemma, the
      multiscale profile; the homogeneous/satellite dichotomy at a
      spike.
\item \emph{Consolidation and audit} (V79--V93).  The M-list with
      statuses, the second internal audit of all fourteen blocks with
      the corrections of V17, V21, V61, V67, V68, V69 confirmed and no
      new one, the import ledger, the constants ledger, the negative
      closure of the sub-parabolic chain, and the drafted and
      cross-checked first version of the third cycle.
\end{enumerate}
The cycle did not prove the full stability/global-regularity theorem.
The open statements are the three gates of
Assessment~\ref{ass:c2v79-gates}: the type-I Liouville statement for
finite-configuration ancient solutions, the localization of a spike
relative to its seed in the intermediate class, and any constraint on
strong type II beyond the degenerate floors.
\end{theorem}

\begin{proof}
By citation: (C1) Theorems~\ref{thm:c2v80-audit}, \ref{thm:c2v81-block2},
\ref{thm:c2v81-block3}, \ref{thm:c2v82-block4}, \ref{thm:c2v82-block5};
(C2) Theorems~\ref{thm:c2v83-block6}, \ref{thm:c2v83-block7},
\ref{thm:c2v84-block8}; (C3) Theorems~\ref{thm:c2v84-block9},
\ref{thm:c2v85-block10}, \ref{thm:c2v86-block11}; (C4)
Theorem~\ref{thm:c2v86-block12}; (C5)--(C6) Theorems~\ref{thm:c2v87-block13},
\ref{thm:c2v87-block14}; (C7) Theorems~\ref{thm:c2v88-mlist},
\ref{thm:c2v88-imports}, \ref{thm:c2v92-constants}, \ref{thm:c2v91-chain},
\ref{thm:c2v93-xcheck}.
\end{proof}

% =============================================================================
\section{The closure procedure}
\label{sec:c2v94-closure}
% =============================================================================

\begin{programme}[Closure at V100]
\label{prog:c2v94-closure}
\begin{enumerate}[label=\textup{(\arabic*)},leftmargin=3.2em]
\item V95: compulsory companion audit; delivery of the file.
\item V96--V99: final checks (the checker on the whole file; the list
      of in-place corrections and their records; the count of
      versions, theorems and firewalls), and any last acquisition an
      audit item suggests.
\item V100: the closing version appends the companion audit, the
      final record, and the statement that the cycle is closed; the
      file is then renamed
      \texttt{Renewal\_NS\_Stability\_Research\_Notes\_V100\_f2.tex} and
      its digest computed.
\item The drafted V1 (digest of Theorem~\ref{thm:c2v93-xcheck}(iii))
      receives the digest of the closed file in place of its
      placeholder and is saved as
      \texttt{Renewal\_NS\_Stability\_Research\_Notes\_V1.tex}; the third
      cycle's V2 records the digest of that V1 and makes the first
      acquisition programmed in it.
\end{enumerate}
\end{programme}

% =============================================================================
\section{Integrated V94 conclusion and next acquisition order}
\label{sec:c2v94-conclusion}
% =============================================================================

\begin{theorem}[What V94 establishes]
\label{thm:c2v94-conclusion}
Relative to the first cycle and V1--V93: the closing theorem
(Theorem~\ref{thm:c2v94-closing}) and the closure procedure.  The
full Navier--Stokes stability/global-regularity theorem remains
unproved.
\end{theorem}

\begin{proof}
The cited records.
\end{proof}

\begin{programme}[V95 acquisition order]
\label{prog:c2v95}
\begin{enumerate}[label=\textup{(AC\arabic*)},leftmargin=4.8em]
\item The compulsory companion audit: read the current SM and YM
      notes and record their digests and decision.
\item The count of the cycle: versions, theorems, propositions,
      lemmas, corollaries, firewalls, programmes, and in-place
      corrections, for the final record.
\end{enumerate}
\end{programme}

\begin{assessment}[Position after V94]
\label{ass:c2v94-position}
The closing theorem and the closure procedure are recorded.  No full
stability or global-regularity claim is made.
\end{assessment}

% =============================================================================
\section*{Version V94 (second cycle) append-only record}
\addcontentsline{toc}{section}{Version V94 (second cycle) append-only record}
% =============================================================================

The complete V93 source of the second cycle was retained before this
continuation.  Its SHA--256 digest was
\begin{center}\scriptsize\ttfamily
ad6682f9a9d75e93a45fbac80b5daee2e04d4d70748eca3b25bb41d61be191c3.
\end{center}
V94 recorded the closing theorem and the closure procedure.  No full
regularity claim is made.

% =============================================================================
\section{V95: compulsory companion audit and the count of the cycle}
\label{sec:c2v95-entry}
% =============================================================================

This is the ninety-fifth version of the second cycle.  The compulsory
review of the two companion programmes is performed first.  The count
of the cycle (versions, statements, firewalls, programmes, in-place
corrections) is then recorded for the final record.  No global
regularity claim is made.  The next compulsory companion audit is
V100, the closing version.

% =============================================================================
\section{Compulsory V95 review of the current companion notes}
\label{sec:c2v95-companion-audit}
% =============================================================================

The current companion files in \texttt{latex\_notes} were inspected
read-only.  The frozen checkpoints are
\begin{align}
 &\texttt{Renewal\_SM\_Einstein\_Continuum\_Research\_Notes\_V20.tex},
 \notag\\[-2pt]
 &\qquad
 \texttt{4737f516e6586f92c8651784bac453085027c134f04a70d5f4d53a02a489cf71},
 \label{eq:c2v95-SM-checkpoint}\\
 &\texttt{Renewal\_Yang\_Mills\_Mass\_Gap\_Research\_Notes\_V95.tex},
 \notag\\[-2pt]
 &\qquad
 \texttt{c43dc30eb18f14c95a31c5e052870087a575e1ffc0360e8e7a27d426cb4bf1a1}.
 \label{eq:c2v95-YM-checkpoint}
\end{align}
(SM V20 is unchanged since the V90 audit, digest identical to
\eqref{eq:c2v90-SM-checkpoint}; YM V95 has 43026 lines.)

\emph{SM V20.}  Unchanged: the companion audit and the exact
constant-Higgs quartic transfer recorded at V90.

\emph{YM V95.}  The YM programme treated fundamental cycles with
submacroscopic overlap, proved a reserve-preserving two-laminar
uncrossing, and defined the remaining high-overlap nested-chain card
on the canonical positive-excess transversal core.

\begin{theorem}[V95 companion-audit decision]
\label{thm:c2v95-companion-decision}
Nothing is imported from SM V20 or YM V95.  Neither bears on the
closing of the second cycle or on the drafted first version of the
third.
\end{theorem}

\begin{proof}
The SM file is unchanged; the YM results concern cluster expansions
of a lattice clock.  None is a Navier--Stokes statement.
\end{proof}

% =============================================================================
\section{The count of the cycle}
\label{sec:c2v95-count}
% =============================================================================

\begin{theorem}[Count at V95]
\label{thm:c2v95-count}
The file at V94 contains, from V1 to V94: ninety-four version
sections; \(239\) theorem environments, \(77\) propositions, \(25\)
lemmas, \(24\) corollaries, \(8\) definitions and \(19\) remarks; \(48\)
firewalls; \(105\) programmes; \(170\) assessments; and eleven records
announcing in-place corrections.  Of the theorem environments, those
withdrawn as stated are the two of V56 (Corollary~\ref{cor:c2v56-delta}
and Theorem~\ref{thm:c2v56-delay}, repaired in V68); every other
statement stands after the second internal audit.
\end{theorem}

\begin{proof}
Counted on the V94 file by the series' tools; the withdrawal list is
Theorem~\ref{thm:c2v88-mlist}.
\end{proof}

% =============================================================================
\section{Integrated V95 conclusion and next acquisition order}
\label{sec:c2v95-conclusion}
% =============================================================================

\begin{theorem}[What V95 establishes]
\label{thm:c2v95-conclusion}
Relative to the first cycle and V1--V94: the companion-audit decision
(Theorem~\ref{thm:c2v95-companion-decision}) and the count
(Theorem~\ref{thm:c2v95-count}).  The full Navier--Stokes
stability/global-regularity theorem remains unproved.
\end{theorem}

\begin{proof}
The cited records.
\end{proof}

\begin{programme}[V96 acquisition order]
\label{prog:c2v96}
\begin{enumerate}[label=\textup{(AC\arabic*)},leftmargin=4.8em]
\item Final check of the whole file with the series' checker, and the
      list of all in-place corrections with the record that announced
      each.
\item The list of negative records of the cycle (statements shown not
      to follow from the series' tools), for the third cycle not to
      repeat them.
\end{enumerate}
\end{programme}

\begin{assessment}[Position after V95]
\label{ass:c2v95-position}
The companion audit is recorded and the cycle counted.  No full
stability or global-regularity claim is made.
\end{assessment}

% =============================================================================
\section*{Version V95 (second cycle) append-only record}
\addcontentsline{toc}{section}{Version V95 (second cycle) append-only record}
% =============================================================================

The complete V94 source of the second cycle was retained before this
continuation.  Its SHA--256 digest was
\begin{center}\scriptsize\ttfamily
88d7022fd1fa060e4a498a8dcec11496fe6a33dd10a26c27401d568062255511.
\end{center}
V95 performed the compulsory companion audit and counted the cycle.
No full regularity claim is made.

% =============================================================================
\section{V96: final check of the file and the list of in-place corrections}
\label{sec:c2v96-entry}
% =============================================================================

Items (AC1) and (AC2) of Programme~\ref{prog:c2v96} are carried out:
the whole file passes the series' checker, the in-place corrections
made during the cycle are listed with the records that announced them,
and the negative records of the cycle are listed for the third cycle
not to repeat them.  No global regularity claim is made.  The next
compulsory companion audit is V100.

% =============================================================================
\section{Final check and in-place corrections}
\label{sec:c2v96-check}
% =============================================================================

\begin{theorem}[State of the file at V96]
\label{thm:c2v96-check}
The file at V95 has no unresolved reference, no duplicate label,
balanced environments and braces, and both citation keys present.
The in-place corrections made to earlier versions after their
assembly, each announced in the record of the following version with
the corrected digest, are:
\begin{enumerate}[label=\textup{(\roman*)},leftmargin=3.2em]
\item V28: a lost firewall restored by a splice (announced in V32);
\item V51: a factor in the spectral tail display (announced in V52);
\item V57: the pressure parenthesis, the torus term of the
      far-pressure lemma, the smallness hypothesis of the weighted
      inequalities, and the exponent of the local-pressure Young
      inequality (announced in V58);
\item V66: a legacy citation name (announced in V67);
\item V67: the corollary on trap-active scales (announced in V68);
\item V76: two legacy citation names (announced in V77);
\item V91: two verbatim occurrences of the closing-line text
      (announced in V92).
\end{enumerate}
The corrections of statements made by later versions without editing
the earlier text are those of Theorem~\ref{thm:c2v88-mlist}: V17, V21,
V61, V68, V69.
\end{theorem}

\begin{proof}
The checker's output on the V95 file, and the records cited.
\end{proof}

% =============================================================================
\section{Negative records of the cycle}
\label{sec:c2v96-negative}
% =============================================================================

\begin{theorem}[What the cycle showed does not follow from its tools]
\label{thm:c2v96-negative}
\begin{enumerate}[label=\textup{(N\arabic*)},leftmargin=3.6em]
\item The moment ledger contains no incompatible row (V22).
\item The third averaged sign is not a strengthening of the first two
      (V8).
\item The inertia identity gives no unconditional sign (V18).
\item The lower side of the constant-chasing gate is not reached by
      the rate-free local energy bound (V38).
\item Delocalization does not bound the palinstrophy per enstrophy;
      spreading does not slow enstrophy growth (V52).
\item The cellwise energy budget does not exclude spreading; only
      ``small stays small'' (V47).
\item The trap bounds maxima, not sums: the exterior's total
      enstrophy and dissipation are not controlled (V62).
\item A trap relative to the moving non-quiet set restarts too often,
      with or without hysteresis (V63, V64).
\item Morrey levels decay linearly under refinement: no coarse-to-fine
      tree, no chained finiteness of the singular set (V65), no
      sub-parabolic chain (V91).
\item Active scales of distinct singular points need not align;
      trap-active scales need not accumulate (V67).
\item The window dissipation floors are not attributable to the
      clusters (V76).
\item The spike's location relative to its seed is not localized; the
      cone of seeds need not converge (V72--V74).
\item The V56 single-cell inequality does not hold as stated; the
      delay is of order \(\nu^3\), not \(1/M\) (V68).
\end{enumerate}
\end{theorem}

\begin{proof}
The cited firewalls and audit items.
\end{proof}

% =============================================================================
\section{Integrated V96 conclusion and next acquisition order}
\label{sec:c2v96-conclusion}
% =============================================================================

\begin{theorem}[What V96 establishes]
\label{thm:c2v96-conclusion}
Relative to the first cycle and V1--V95: the state of the file
(Theorem~\ref{thm:c2v96-check}) and the negative records
(Theorem~\ref{thm:c2v96-negative}).  The full Navier--Stokes
stability/global-regularity theorem remains unproved.
\end{theorem}

\begin{proof}
The cited records.
\end{proof}

\begin{programme}[V97 acquisition order]
\label{prog:c2v97}
\begin{enumerate}[label=\textup{(AC\arabic*)},leftmargin=4.8em]
\item The index of the cycle's firewalls grouped by the gate they
      concern, for the third cycle's use.
\item The list of the cycle's explicit constants with their defining
      versions, as a table.
\end{enumerate}
\end{programme}

\begin{assessment}[Position after V96]
\label{ass:c2v96-position}
The file is checked and its corrections and negative records are
listed.  No full stability or global-regularity claim is made.
\end{assessment}

% =============================================================================
\section*{Version V96 (second cycle) append-only record}
\addcontentsline{toc}{section}{Version V96 (second cycle) append-only record}
% =============================================================================

The complete V95 source of the second cycle was retained before this
continuation.  Its SHA--256 digest was
\begin{center}\scriptsize\ttfamily
a395057782f1adf72afd7aecaa586e1668721a2f9396fabfd238805bd25013cd.
\end{center}
V96 recorded the final check, the in-place corrections and the
negative records.  No full regularity claim is made.

% =============================================================================
\section{V97: index of firewalls by gate, and the table of explicit constants}
\label{sec:c2v97-entry}
% =============================================================================

Items (AC1) and (AC2) of Programme~\ref{prog:c2v97} are carried out.
No global regularity claim is made.  The next compulsory companion
audit is V100.

% =============================================================================
\section{Firewalls by gate}
\label{sec:c2v97-firewalls}
% =============================================================================

\begin{theorem}[Index of the cycle's firewalls]
\label{thm:c2v97-firewalls}
The forty-eight firewalls of V1--V96 concern the gates as follows
(each listed by version and short name).
\begin{enumerate}[label=\textup{(G\arabic*)},leftmargin=3.6em]
\item \emph{Type-I Liouville.}  V2 (the averaged form does not close
      pointwise), V3 (trivial cross terms), V8 (third sign not a
      strengthening), V9 (E9 is an import), V16--V22 (the ledger rows
      without signs: centre normalization, inertia, stress split,
      closure), V28 (the trichotomy's third branch), V31 (forced
      escape is not exclusion), V33--V34 (push budget, variance
      non-explicit), V36--V38 (the gate's lower side), V41 (linear-mode
      law needs the third moment), V57 (what the exclusion of spread
      windows does not say), V59 (the trap does not run on the limit),
      V63 (the finite-configuration description does not close the
      gate).
\item \emph{Intermediate class.}  V42--V44 (nonzero window limits open
      without spreading control; window flux signed only for small
      \(M\)), V47 (the budget does not exclude spreading), V48 (precision
      of the nearly inviscid cell), V52 (delocalization does not slow
      growth), V58 (three gaps of the localized trap), V60 (the ball's
      energy identity), V61 (beyond the window), V62 (maxima not sums),
      V63--V64 (current-region trap, hysteresis), V66 (what the
      exclusion does not localize), V67 (no finiteness of the singular
      set), V72--V74 (spike not localized, hopping, conditions not
      available), V76 (multiplicity not bounded by dissipation floors).
\item \emph{Strong type II.}  V24 (statistically invisible non-compact
      blow-up not excluded), V27--V29 (invisibility mechanisms not
      excluded), V71 (the floor degenerates), V77 (the profile's
      exponents grow with the scale).
\end{enumerate}
\end{theorem}

\begin{proof}
Inspection of the firewall environments.
\end{proof}

% =============================================================================
\section{Explicit constants}
\label{sec:c2v97-constants}
% =============================================================================

\begin{theorem}[Table of the cycle's explicit constants]
\label{thm:c2v97-constants}
\begin{center}
\small
\begin{tabular}{p{0.22\textwidth}p{0.52\textwidth}p{0.14\textwidth}}
\toprule
Constant & Definition & Version\\
\midrule
\(\ell_M\) & \(3\nu^3/(8C_AM^2)\), window length & V42\\
\(C_A\) & \(\frac{27}{16}C_B^4\), cubic enstrophy constant & V51\\
\(P_M\) & \(4M\), palinstrophy budget & V51\\
\(\kappa\), \(K_P\) & \(C_\infty\nu^{-1/2}\); \(4\nu^{-1/2}(4\pi\nu/3)^{1/4}C_{\rm CZ}K_6^2\) & V36\\
\(m'_\kappa\) & \((3-\sqrt{1+4\kappa^2})/2\), margin for \(\kappa<\sqrt2\) & V37\\
\(a\), \(b\), \(b'\), \(K\), \(p\), \(\Lambda_{\rm tr}\) & trap constants, \(a=\nu\kappa/(2916S_\kappa)\), \(K=4b/a\), \(p=4b'/a\), \(\Lambda_{\rm tr}=b'(1+K)\) & V57\\
\(\varepsilon_*(\nu,M)\) & \(S_\kappa^{-1}K^pM^{-p}(Ke^{\Lambda_{\rm tr}})^{-(1+p)}\ge e^{-C(\nu)M^3}\) & V57, V59\\
\(\varepsilon'_*(\nu,\Lambda)\) & \(S_\kappa^{-1}K^{-1}e^{-(1+p)\Lambda_{\rm tr}\Lambda^2}\ge e^{-C(\nu)\Lambda^2}\) (class-one constants) & V71\\
\(N_M(\varepsilon)\) & \(1728C_S^6(2M)^3\varepsilon^{-3}\) & V58\\
\(a_M\) & \(8C_S^2M\), window cell-energy bound & V58\\
\(\varepsilon_{\rm th}\) & \(\min(\varepsilon_*/2,\varepsilon_\star,\varepsilon_{\star\star})\) & V61, V66\\
\(\Gamma(\tau)\), \(\tau_c\), \(\rho_0\), \(R_{\rm inv}\) & crude growth function; \(\Gamma(\tau_c)=\frac38\varepsilon_{\rm th}\); \((4C'_\flat N_Ma_M\ell_M\varepsilon_{\rm th}^{-1/2})^{1/4}\); \(\lfloor\ell_M/\tau_c\rfloor\rho_0\) & V61\\
\(\tau_{\rm cap}\), \(r_1\), \(r_2\) & \(1/b\le\ell_M/2\); pollution radius; \(R_{\rm inv}+(N_M/\varepsilon_{\rm th})^{1/4}\) & V62\\
\(\rho_{\min}\), \(\Delta\), \(R'_{\rm inv}\), \(C_1\) & energetic-seeds constants, \(\tau_{\rm cap}\le\varepsilon_{\rm th}^{1/2}\) & V66\\
\(C_4'(\nu)\), \(C_4(\nu)\) & pre-spike exponents & V72\\
\(c_g\), \(\Lambda_{\rm seed}\), \(\varepsilon_2\) & \((\nu^3/(32C_1M^2))^{1/2}\); \(\max(4M,Me^{C_4'}/c_g)\); \(c_0c_g\varepsilon_*\) & V72\\
\(K\) (ledger), \(c_1\), \(n_0\) & Lipschitz constant at a spike; \(e^{-1/4\nu}c_0\varepsilon_*\); \(\lfloor c_1/(K\ell)\rfloor\) & V76\\
\bottomrule
\end{tabular}
\end{center}
All are explicit in principle in \((\nu,M)\) and the absolute constants
of the classical inequalities (Theorem~\ref{thm:c2v92-constants}).
\end{theorem}

\begin{proof}
The cited definitions.
\end{proof}

% =============================================================================
\section{Integrated V97 conclusion and next acquisition order}
\label{sec:c2v97-conclusion}
% =============================================================================

\begin{theorem}[What V97 establishes]
\label{thm:c2v97-conclusion}
Relative to the first cycle and V1--V96: the index of firewalls
(Theorem~\ref{thm:c2v97-firewalls}) and the table of constants
(Theorem~\ref{thm:c2v97-constants}).  The full Navier--Stokes
stability/global-regularity theorem remains unproved.
\end{theorem}

\begin{proof}
The cited records.
\end{proof}

\begin{programme}[V98 acquisition order]
\label{prog:c2v98}
\begin{enumerate}[label=\textup{(AC\arabic*)},leftmargin=4.8em]
\item The dependency graph of the cycle's principal theorems (which
      statement uses which), so that a future correction can be
      traced.
\item The pre-closure state: confirm that the drafted V1 and the
      closure procedure of Programme~\ref{prog:c2v94-closure} are
      ready.
\end{enumerate}
\end{programme}

\begin{assessment}[Position after V97]
\label{ass:c2v97-position}
Firewalls are indexed and constants tabulated.  No full stability or
global-regularity claim is made.
\end{assessment}

% =============================================================================
\section*{Version V97 (second cycle) append-only record}
\addcontentsline{toc}{section}{Version V97 (second cycle) append-only record}
% =============================================================================

The complete V96 source of the second cycle was retained before this
continuation.  Its SHA--256 digest was
\begin{center}\scriptsize\ttfamily
53de22e8767893b384e16b10ca56891ac02d7f4c467799447ab247dfb401c916.
\end{center}
V97 indexed the firewalls and tabulated the constants.  No full
regularity claim is made.

% =============================================================================
\section{V98: the dependency graph of the principal theorems, and the
pre-closure state}
\label{sec:c2v98-entry}
% =============================================================================

Items (AC1) and (AC2) of Programme~\ref{prog:c2v98} are carried out.
No global regularity claim is made.  The next compulsory companion
audit is V100.

% =============================================================================
\section{Dependencies}
\label{sec:c2v98-dependencies}
% =============================================================================

\begin{theorem}[Dependency graph of the principal theorems]
\label{thm:c2v98-dependencies}
Arrows read ``is used by''.
\begin{enumerate}[label=\textup{(D\arabic*)},leftmargin=3.6em]
\item Profile equation and dictionary (V2) \(\to\) second identity (V2)
      \(\to\) dichotomy (V4), defect equation (V7), Rayleigh identity
      (V31), mean identities (V12), moment ledger (V16--V22).
\item Defect regularity (V7) and (E5), (E6), (E8) \(\to\) small-defect
      theorem (V4) \(\to\) dichotomy (A)/(B); (E9) \(\to\) exclusion of (A)
      (V9) \(\to\) three signs (V8), gates (V15), statistical dichotomy
      (V23).
\item Window proposition (V42, cubic inequality) \(\to\) window
      compactness (V42) \(\to\) active/spread (V43--V44), energy-active
      implies dissipation-active (V48), cluster limits (V59), ledgers at
      a spike (V76).
\item Palinstrophy budget (V51) \(\to\) band (V51), sharp inequality
      (V52), stretching bound (V53), influx bound (V54), regular-cell
      counting (V52/V54, corrected V61), no travelling windows (V54).
\item Far-pressure lemma (V57), interpolation (V57), cell inequality
      (V57) \(\to\) weighted inequalities (V57) \(\to\) capture and trapping
      (V57) \(\to\) no spread windows (V57), Morrey floor (V59),
      sub-parabolic floor (V71), pre-spike floor (V72), repaired delay
      (V68); the global inequality (V56) is used in all four.
\item Count (V58, \(L^6\) bound) \(\to\) configurations (V58, V71),
      seeding count (V72), halo bound (V58) \(\to\) localized inequalities
      (V58, full in V69) \(\to\) exterior trap (V62), energetic seeds
      (V66).
\item Crude and far growth lemmas (V61) \(\to\) finite propagation (V61)
      \(\to\) hierarchy closure (V61), exterior trap (V62), type-I
      structure (V63), persistence (V64), energetic seeds (V66),
      quiet-scale dichotomy (V67), one-step localization (V73).
\item Pre-spike floor (V72) \(\to\) seeding (V72, with the backward cone
      of V46 and the first cycle's cone constant), cone of seeds (V73)
      \(\to\) sufficient conditions (V74), persistence at a spike (V76),
      first-spike description (V76).
\item Everything in the trap, propagation and spike chapters depends
      on the classical inequalities and the continuation criterion
      only; the singular-point statements (V63, V66, V67) depend in
      addition on (E1)/(E7) and on the first cycle's active-scale
      analysis.
\end{enumerate}
A correction to a statement propagates along the arrows from it; the
two withdrawals of V56 propagated nowhere, since nothing used
Corollary~\ref{cor:c2v56-delta} or Theorem~\ref{thm:c2v56-delay}.
\end{theorem}

\begin{proof}
Inspection of the proofs' citations.
\end{proof}

% =============================================================================
\section{Pre-closure state}
\label{sec:c2v98-preclosure}
% =============================================================================

\begin{theorem}[Ready for closure]
\label{thm:c2v98-ready}
The drafted first version of the third cycle (Theorem~\ref{thm:c2v93-xcheck}(iii))
is complete except for its placeholder digest; the closure procedure
(Programme~\ref{prog:c2v94-closure}) is fixed; the file passes the
checker (Theorem~\ref{thm:c2v96-check}); the audited M-list, the
import ledger, the constants ledger, the negative records, the
firewall index and the dependency graph are recorded.  V99 records
the last state before closure and V100 closes.
\end{theorem}

\begin{proof}
The cited records.
\end{proof}

% =============================================================================
\section{Integrated V98 conclusion and next acquisition order}
\label{sec:c2v98-conclusion}
% =============================================================================

\begin{theorem}[What V98 establishes]
\label{thm:c2v98-conclusion}
Relative to the first cycle and V1--V97: the dependency graph
(Theorem~\ref{thm:c2v98-dependencies}) and the pre-closure state
(Theorem~\ref{thm:c2v98-ready}).  The full Navier--Stokes
stability/global-regularity theorem remains unproved.
\end{theorem}

\begin{proof}
The cited records.
\end{proof}

\begin{programme}[V99 acquisition order]
\label{prog:c2v99}
\begin{enumerate}[label=\textup{(AC\arabic*)},leftmargin=4.8em]
\item The last state before closure: the statement of the three
      gates in their final second-cycle form, one paragraph each, and
      the sentence that will open the third cycle's first acquisition.
\item Nothing else; V100 performs the companion audit and closes.
\end{enumerate}
\end{programme}

\begin{assessment}[Position after V98]
\label{ass:c2v98-position}
The dependency graph is recorded and the cycle is ready for closure.
No full stability or global-regularity claim is made.
\end{assessment}

% =============================================================================
\section*{Version V98 (second cycle) append-only record}
\addcontentsline{toc}{section}{Version V98 (second cycle) append-only record}
% =============================================================================

The complete V97 source of the second cycle was retained before this
continuation.  Its SHA--256 digest was
\begin{center}\scriptsize\ttfamily
3cc54af403f20cf35f9522f929802f19b8fba886beae52d85178f45a3529f18a.
\end{center}
V98 recorded the dependency graph and the pre-closure state.  No full
regularity claim is made.

% =============================================================================
\section{V99: the last state before closure}
\label{sec:c2v99-entry}
% =============================================================================

Item (AC1) of Programme~\ref{prog:c2v99} is carried out: the three
gates in their final second-cycle form, one paragraph each, and the
sentence that opens the third cycle's first acquisition.  No global
regularity claim is made.  V100 performs the compulsory companion
audit and closes the cycle.

% =============================================================================
\section{The gates in their final second-cycle form}
\label{sec:c2v99-gates}
% =============================================================================

\begin{assessment}[Gate G1, type-I Liouville]
\label{ass:c2v99-G1}
A type-I singularity produces, at each singular point, a nonzero
ancient solution in \(\mathcal A_\nu(C_I^2\nu^{3/2})\) with the first cycle's
properties; the second cycle adds that, at every time, it is a finite
configuration of at most \(N_C\) energetic parabolic cells with finite
propagation and a uniformly diffuse far field after a capture time;
that at its homogeneous points the profile satisfies the second
identity, the mean identities, the moment ledger and the Rayleigh
identity, with alternative (B) (defect rate \(\underline\delta>0\),
anti-alignment) forced by (E9); that the homogeneous points number at
most \(N_C(c')\); and that the constant-chasing gate closes on its
upper side for \(\kappa<\sqrt2\).  Missing: one mean inequality for one
blow-up measure, or an independent exclusion of finite-configuration
ancient solutions.
\end{assessment}

\begin{assessment}[Gate G2, the intermediate class]
\label{ass:c2v99-G2}
A solution that is type I along a sequence of window times without a
rate has, at each window, a finite configuration of between one and
\(N_M\) energetic parabolic cells (no spread windows), with the
activity confined to the invasion radius of the energetic seeds
through the window and the exterior uniformly weak after the capture
time; every spike of height at least \(\Lambda_{\rm seed}\) after the window is
seeded at the window time in the configuration at the gap scale, the
first spike of a given height is approached type-I relative to itself,
and its cone of seeds is localized step by step but may hop among the
clusters of successive configurations.  Missing: the localization of
a spike relative to its seed (convergence of the cone), or any bound
in the stretch between the window's end and the spike, where the
enstrophy is unbounded and no local bound holds.
\end{assessment}

\begin{assessment}[Gate G3, strong type II]
\label{ass:c2v99-G3}
Along a sequence of times with \(\mathfrak Y\to\infty\), the second cycle
gives only the degenerate floors \(\mathfrak M\ge e^{-C\mathfrak Y^2}\) and
counts \(\le C\mathfrak Y^3\), the sub-parabolic configuration with
viscosity-only constants on a fraction \(\mathfrak Y^{-2}\) of the remaining
time, and the fact that a singularity with vanishing parabolic energy
along a sequence is of type II along it.  Missing: any constraint
beyond the first cycle's price sheet.
\end{assessment}

% =============================================================================
\section{The opening sentence of the third cycle's first acquisition}
\label{sec:c2v99-opening}
% =============================================================================

\begin{programme}[Opening sentence]
\label{prog:c2v99-opening}
``On a window at the seed scale of a spike approached type-I relative
to itself, two clusters each carrying energy at least \(c_0\varepsilon_*\lambda\)
after the capture time are separated by a trapped exterior, and the
local energy identity of each cluster's neighbourhood, with the far
pressure bounded by the other's halo and the quiet level and the
dissipation bounded below by the per-cluster level of energy-active
centres, is examined for the number of window lengths over which both
can persist.''
\end{programme}

% =============================================================================
\section{Integrated V99 conclusion and next acquisition order}
\label{sec:c2v99-conclusion}
% =============================================================================

\begin{theorem}[What V99 establishes]
\label{thm:c2v99-conclusion}
Relative to the first cycle and V1--V98: the final second-cycle form
of the three gates (Assessments~\ref{ass:c2v99-G1}--\ref{ass:c2v99-G3})
and the opening sentence of the third cycle.  The full Navier--Stokes
stability/global-regularity theorem remains unproved.
\end{theorem}

\begin{proof}
The cited records.
\end{proof}

\begin{programme}[V100 acquisition order]
\label{prog:c2v100}
\begin{enumerate}[label=\textup{(AC\arabic*)},leftmargin=4.8em]
\item The compulsory companion audit.
\item The closing record of the second cycle and the closure
      (Programme~\ref{prog:c2v94-closure}).
\end{enumerate}
\end{programme}

\begin{assessment}[Position after V99]
\label{ass:c2v99-position}
The gates are stated in their final form; the cycle closes at V100.
No full stability or global-regularity claim is made.
\end{assessment}

% =============================================================================
\section*{Version V99 (second cycle) append-only record}
\addcontentsline{toc}{section}{Version V99 (second cycle) append-only record}
% =============================================================================

The complete V98 source of the second cycle was retained before this
continuation.  Its SHA--256 digest was
\begin{center}\scriptsize\ttfamily
9618b7b29ec7e31edf604544c6f5b675dd326366d4722e6abec77060a308fe29.
\end{center}
V99 recorded the last state before closure.  No full regularity claim
is made.

% =============================================================================
\section{V100: compulsory companion audit and the closure of the second cycle}
\label{sec:c2v100-entry}
% =============================================================================

This is the hundredth and closing version of the second cycle.  The
compulsory review of the two companion programmes is performed, the
closing record is made, and the cycle is closed according to
Programme~\ref{prog:c2v94-closure}: this file is retained as the legacy
file of the second cycle, and the third cycle opens with the first
version drafted in V92--V93.  No global regularity claim is made.

% =============================================================================
\section{Compulsory V100 review of the current companion notes}
\label{sec:c2v100-companion-audit}
% =============================================================================

The current companion files in \texttt{latex\_notes} were inspected
read-only.  The frozen checkpoints are
\begin{align}
 &\texttt{Renewal\_SM\_Einstein\_Continuum\_Research\_Notes\_V21.tex},
 \notag\\[-2pt]
 &\qquad
 \texttt{385e336333afaedc86b5005d7602fde51e9a6b7ed1324141ddd82d4ab0d89f64},
 \label{eq:c2v100-SM-checkpoint}\\
 &\texttt{Renewal\_Yang\_Mills\_Mass\_Gap\_Research\_Notes\_V95.tex},
 \notag\\[-2pt]
 &\qquad
 \texttt{c43dc30eb18f14c95a31c5e052870087a575e1ffc0360e8e7a27d426cb4bf1a1}.
 \label{eq:c2v100-YM-checkpoint}
\end{align}
(SM V21 has 9044 lines; YM V95 is unchanged since the V95 audit,
digest identical to \eqref{eq:c2v95-YM-checkpoint}.)

\emph{SM V21 (seventh series).}  After closing the constant-Higgs
quartic word (V20), the SM programme treated the remaining pure-Higgs
weight-four word, the covariant gradient pair, by a plane-wave Higgs
perturbation and differentiation of the complete heat trace in its
amplitude, obtaining the Higgs-gradient two-point symbol and its
physical coefficient.

\emph{YM V95.}  Unchanged: the submacroscopic-overlap cycles, the
two-laminar uncrossing and the high-overlap nested-chain card recorded
at V95.

\begin{theorem}[V100 companion-audit decision]
\label{thm:c2v100-companion-decision}
Nothing is imported from SM V21 or YM V95.  Neither bears on the
closure of the second cycle or on the third cycle's first version.
\end{theorem}

\begin{proof}
The SM results concern heat-kernel coefficients of lattice
Dirac--Yukawa operators; the YM file is unchanged.  None is a
Navier--Stokes statement.
\end{proof}

% =============================================================================
\section{Closure of the second cycle}
\label{sec:c2v100-closure}
% =============================================================================

\begin{theorem}[The second cycle is closed]
\label{thm:c2v100-closed}
The second cycle consists of the versions V1--V100 of this file.  Its
mathematical content is Theorem~\ref{thm:c2v94-closing}; its audited
statements are those of Theorem~\ref{thm:c2v88-mlist}; its imports are
those of Theorem~\ref{thm:c2v88-imports}; its withdrawn statements are
the two of V56; its negative records are those of
Theorem~\ref{thm:c2v96-negative}; its open gates are those of
Assessments~\ref{ass:c2v99-G1}--\ref{ass:c2v99-G3}.  The full
Navier--Stokes stability/global-regularity theorem is not proved by
this cycle.  This file is retained as
\begin{center}
\texttt{Renewal\_NS\_Stability\_Research\_Notes\_V100\_f2.tex},
\end{center}
and the third cycle opens with the file
\texttt{Renewal\_NS\_Stability\_Research\_Notes\_V1.tex} assembled from
the draft of Theorem~\ref{thm:c2v93-xcheck}(iii) with the digest of
the present file in place of its placeholder; the third cycle's V2
records the digest of that V1.
\end{theorem}

\begin{proof}
The cited records and Programme~\ref{prog:c2v94-closure}.
\end{proof}

\begin{assessment}[Position at the closure of the second cycle]
\label{ass:c2v100-position}
Two hundred versions of the series are complete.  The strongest
results of the second cycle, valid for any singularity, are the Morrey
floors at every time and the pre-spike floor; for the intermediate
class, the exclusion of spread windows, the finite configuration with
finite propagation and a trapped exterior at every window, and the
seeding of every high spike in the window's configuration; for type
I, the same at every time with the count of homogeneous points.  The
three gates remain open.  No full stability or global-regularity
claim is made.
\end{assessment}

% =============================================================================
\section*{Version V100 (second cycle) append-only record}
\addcontentsline{toc}{section}{Version V100 (second cycle) append-only record}
% =============================================================================

The complete V99 source of the second cycle was retained before this
continuation.  Its SHA--256 digest was
\begin{center}\scriptsize\ttfamily
8cd69f634a3c5df3ed07055120df1457e25a99da99a7f1708f4fc378555fd732.
\end{center}
V100 performed the compulsory companion audit and closed the second
cycle.  No full regularity claim is made.

\end{document}
